\documentclass[UTF8,12pt,a4paper,twoside,openany]{ctexbook}

\usepackage[top=2.7cm,bottom=2.8cm,inner=2.8cm,outer=2.4cm,headheight=15pt]{geometry}
\usepackage{amsmath,amssymb,mathtools}
\usepackage{mathrsfs}
\usepackage{slashed}
\usepackage{booktabs,tabularx,array,longtable}
\usepackage{enumitem}
\usepackage[table]{xcolor}
\usepackage{graphicx}
\usepackage{float}
\usepackage{seqsplit}
\usepackage{tikz}
\usetikzlibrary{arrows.meta,positioning,calc,decorations.pathmorphing,shapes.geometric}
\usepackage{microtype}
\usepackage[numbers,sort&compress]{natbib}
\usepackage{caption}
\usepackage[most]{tcolorbox}
\usepackage{titlesec}
\usepackage{fancyhdr}
\usepackage{emptypage}
\usepackage{algorithm}
\usepackage{algpseudocode}
\usepackage{hyperref}

\definecolor{qcdnavy}{HTML}{173B57}
\definecolor{qcdmist}{HTML}{EDF3F7}
\definecolor{qcdrust}{HTML}{A65332}
\definecolor{qcdsand}{HTML}{F8F3E9}
\definecolor{qcdsage}{HTML}{426B5A}
\definecolor{qcdsagebg}{HTML}{EEF4F0}

\hypersetup{
  colorlinks=true,
  linkcolor=blue!55!black,
  citecolor=blue!55!black,
  urlcolor=blue!55!black,
  pdftitle={格点QCD中的核子胶子TMD-PDF：理论、算法与可复现实现},
  pdfauthor={MyQCD 理论物理研究组}
}

\captionsetup{
  font=small,
  labelfont={bf,color=qcdnavy},
  textfont={small},
  skip=5pt
}

\titleformat{\part}[display]
  {\normalfont\Huge\bfseries\color{qcdnavy}\centering}
  {\Large\color{qcdrust}\partname\ \thepart}
  {1.2ex}
  {\vspace{0.4ex}{\color{qcdnavy!28}\titlerule[0.8pt]}\vspace{1.2ex}}
\titlespacing*{\part}{0pt}{2.4cm}{2.0cm}
\titleformat{\chapter}[display]
  {\normalfont\Huge\bfseries\color{qcdnavy}}
  {\filleft\Large\color{qcdrust}\chaptertitlename\ \thechapter}
  {0.8ex}
  {\vspace{0.25ex}{\color{qcdnavy!28}\titlerule[0.8pt]}\vspace{1.2ex}\filright}
\titlespacing*{\chapter}{0pt}{-1.0cm}{2.1cm}
\titleformat{\section}
  {\normalfont\Large\bfseries\color{qcdnavy}}
  {\color{qcdrust}\thesection}{0.72em}{}
  [\vspace{0.18em}{\color{qcdnavy!28}\titlerule[0.7pt]}]
\titlespacing*{\section}{0pt}{2.4ex plus .35ex minus .15ex}{1.15ex}
\titleformat{\subsection}
  {\normalfont\large\bfseries\color{qcdnavy!92!black}}
  {\color{qcdrust}\thesubsection}{0.62em}{}
\titlespacing*{\subsection}{0pt}{1.8ex plus .25ex minus .1ex}{0.75ex}
\titleformat{\subsubsection}
  {\normalfont\normalsize\bfseries\color{qcdnavy!78!black}}
  {\thesubsubsection}{0.55em}{}

\pagestyle{fancy}
\fancyhf{}
\fancyhead[LE,RO]{\small\thepage}
\fancyhead[RE]{\small\nouppercase{\leftmark}}
\fancyhead[LO]{\small\nouppercase{\rightmark}}
\renewcommand{\headrulewidth}{0.4pt}
\renewcommand{\chaptermark}[1]{\markboth{\thechapter\quad #1}{}}
\renewcommand{\sectionmark}[1]{\markright{\thesection\quad #1}}
\fancypagestyle{plain}{
  \fancyhf{}
  \fancyhead[LE,RO]{\small\thepage}
  \renewcommand{\headrulewidth}{0pt}
}

\tcbset{
  enhanced,
  breakable,
  boxrule=0.55pt,
  arc=1.6mm,
  left=7pt,
  right=7pt,
  top=5pt,
  bottom=5pt,
  before skip=7pt,
  after skip=7pt,
  fonttitle=\bfseries
}
\newtcolorbox{keyformula}[1]{
  colback=qcdmist,
  colframe=qcdnavy!68!black,
  colbacktitle=qcdnavy!88!black,
  coltitle=white,
  title={#1}
}
\newtcolorbox{physicalbox}[1]{
  colback=qcdsand,
  colframe=qcdrust!72!black,
  colbacktitle=qcdrust!86!black,
  coltitle=white,
  title={#1}
}
\newtcolorbox{checkpointbox}[1]{
  colback=qcdsagebg,
  colframe=qcdsage!75!black,
  colbacktitle=qcdsage!88!black,
  coltitle=white,
  title={#1}
}

\linespread{1.58}
\setlength{\parindent}{2em}
\setlength{\parskip}{0.32em}
\setlength{\emergencystretch}{3em}
\setlength{\abovedisplayskip}{7pt plus 2pt minus 2pt}
\setlength{\belowdisplayskip}{7pt plus 2pt minus 2pt}
\renewcommand{\arraystretch}{1.12}
\allowdisplaybreaks
\setlist[itemize]{leftmargin=2.2em,itemsep=0.15em,topsep=0.25em}
\setlist[enumerate]{leftmargin=2.5em,itemsep=0.18em,topsep=0.25em}
\setcounter{secnumdepth}{3}
\setcounter{tocdepth}{2}
\raggedbottom

\newcommand{\MSbar}{\overline{\mathrm{MS}}}
\newcommand{\dd}{\mathrm{d}}
\newcommand{\Tr}{\operatorname{Tr}}
\newcommand{\Ord}{\mathcal{O}}
\newcommand{\Pz}{P_z}
\newcommand{\bT}{\boldsymbol b_T}
\newcommand{\kT}{\boldsymbol k_T}
\newcommand{\pT}{\boldsymbol p_T}
\newcommand{\bperp}{\boldsymbol b_\perp}
\newcommand{\Ocal}{\mathcal O}
\newcommand{\LQCD}{\Lambda_{\rm QCD}}
\newcommand{\fone}{f_{1}^{g/N}}
\newcommand{\Oq}{\widetilde{\mathcal O}_{g}}

\floatname{algorithm}{算法}
\renewcommand{\algorithmicrequire}{\textbf{输入：}}
\renewcommand{\algorithmicensure}{\textbf{输出：}}

\newcolumntype{Y}{>{\raggedright\arraybackslash}X}
\newcolumntype{L}[1]{>{\raggedright\arraybackslash}p{#1}}

\title{\heiti 格点QCD中的核子胶子TMD-PDF\\[0.2em]\Large 理论、算法与可复现实现}
\author{MyQCD 理论物理研究组}
\date{2026 年 10 月 23 日}

\begin{document}

\frontmatter
\maketitle

\chapter*{摘要与关键词}
\addcontentsline{toc}{chapter}{摘要与关键词}
\markboth{摘要与关键词}{}

\begin{tcolorbox}[
  enhanced,
  breakable,
  colback=white,
  colframe=qcdnavy!46!black,
  colbacktitle=qcdnavy!90!black,
  coltitle=white,
  title={摘要 | Abstract},
  boxrule=0.7pt,
  arc=1.8mm,
  left=9pt,
  right=9pt,
  top=7pt,
  bottom=7pt,
  fonttitle=\bfseries
]
\small
\noindent
胶子的横动量依赖部分子分布函数（transverse-momentum-dependent parton distribution function，TMD-PDF）同时描述核子内胶子的纵向动量份额 \(x\)、横向动量 \(\kT\)、自旋关联以及非微扰软胶子效应，是理解质子三维结构、小 \(x\) 动力学和未来电子离子对撞机（EIC）物理的重要输入。本书以使用格点 QCD 从第一性原理计算核子胶子 TMD-PDF 为主线，从光锥胶子关联函数及其规范链定义出发，系统建立快度（rapidity）发散、软函数、Collins--Soper 演化、准 TMD 算符、LaMET 因子化、Wilson 圈减除、短距比率与混合重正化、梯度流以及连续极限之间的逻辑关系。特别强调，梯度流提供的是受控的 UV 正规化与短流时匹配工具，而不是 rapidity 发散、表示相关软因子或完整单态混合的替代品。本书进一步给出格点两点与三点关联函数、Clover 场强、核子激发态、有限体积、Dirac 传播子求逆、Krylov 算法、Fourier 反演、统计推断和外推的系统误差矩阵，并以库内文献与实现审计为证据，区分共线胶子 PDF 的已发表先例、TMD 软函数与 CS 核的格点结果、夸克 TMD 的方法学基准，以及尚待闭合的直接核子胶子 TMD 计算。书中还把 BiCGSTAB 从“线性求解器”还原为规范背景中的传播子构造，说明近零模、条件数、breakdown、残差和求解容差如何进入物理误差预算。作为全书证据层，本书同时吸收固定提交下 PyQCD 参考快照中的 163 个内容文件，以 87 个知识单元、逐文件哈希账本、历史验证矩阵和来源图谱保证每项实质内容均可追溯；对重复版本作去重和演进解释，而不是简单拼接。完整计算必须同时满足算符几何、表示、软因子、重正化、匹配和大动量尺度窗口的一致性；缺少任一环节时，只能报告诊断量或中间矩阵元，不能称为物理胶子 TMD-PDF。

\vspace{0.7em}
\noindent\textbf{关键词：}格点 QCD；胶子 TMD-PDF；大动量有效理论；Collins--Soper 核；梯度流；重正化；核子结构

\vspace{1.2em}
\noindent\textbf{Abstract.}
This book develops the theoretical, algorithmic, and reproducible-implementation framework for calculating the gluon transverse-momentum-dependent parton distribution function (TMD-PDF) of the nucleon in lattice QCD. Starting from the light-cone gluon correlator and its gauge link, it connects rapidity divergences, the soft function, Collins--Soper evolution, quasi-TMD operators, LaMET matching, Wilson-loop subtraction, short-distance ratio and hybrid renormalization, gradient flow, Dirac inversions, statistical inference, and continuum extrapolation. Gradient flow is presented as a controlled ultraviolet regulator and short-flow-time matching tool, not as a replacement for the representation-dependent soft factor, rapidity renormalization, or complete singlet mixing. The book distinguishes published lattice calculations of collinear gluon PDFs, determinations of the soft function and Collins--Soper kernel, methodological benchmarks from quark TMDs, and the remaining requirements for a physical nucleon gluon TMD. A calculation is physically complete only when operator geometry, color representation, soft and rapidity factors, renormalization, matching, and the large-momentum scale hierarchy are mutually consistent. It also explains how near-zero modes, conditioning, iterative-solver breakdown, true residuals, and tolerance choices enter the physical error budget rather than remaining hidden implementation details. The evidence layer integrates all 163 content files in a fixed PyQCD reference snapshot through 87 knowledge units, a per-file hash ledger, historical validation matrices, and a source atlas. Foundations, worked derivations, numerical algorithms, implementation audits, a glossary, a formula-and-algorithm handbook, and executable reproduction laboratories are included to make the monograph usable as both a research reference and a course text.

\vspace{0.7em}
\noindent\textbf{Keywords:} lattice QCD; gluon TMD; LaMET; Collins--Soper kernel; gradient flow; renormalization; nucleon structure
\end{tcolorbox}

\clearpage
\tableofcontents

\clearpage
\chapter*{阅读路线与课程设计}
\addcontentsline{toc}{chapter}{阅读路线与课程设计}
\markboth{阅读路线与课程设计}{}

\section*{本书解决的三个层次问题}

本书面向三类读者：希望理解胶子TMD物理定义的场论研究者，准备实现格点矩阵元与匹配链的数值研究者，以及需要评估结果可靠性的现象学或实验研究者。全文围绕三个层次展开：
\begin{enumerate}
  \item \textbf{对象层：}胶子TMD是什么，为什么同时依赖 \(x,\bT,\mu,\zeta\)，它和共线PDF、TMD波函数、准TMD有什么区别；
  \item \textbf{方法层：}如何从Euclidean矩阵元经过重正化、软减除、匹配、Fourier变换和极限外推得到物理量；
  \item \textbf{证据层：}如何证明一个数值结果不是算法假象、方案混用或外推产物。
\end{enumerate}
只回答第一层会得到概念综述；只回答第二层会得到代码说明；只有第三层闭合，结果才可作为可复核的物理输入。

\section*{四条阅读路径}

\begin{table}[htbp]
\centering
\caption{面向不同读者的阅读路径。章节号为逻辑依赖顺序，不要求逐字读完。}
\label{tab:reading-paths}
\small
\begin{tabularx}{\textwidth}{L{2.6cm}L{3.6cm}Y}
\toprule
读者 & 建议主线 & 可先略读 \\
\midrule
场论与TMD理论 & 引言、光锥定义、因子化、高级专题、关键推导 & HMC、实现审计、算法索引 \\
格点实现 & 格点基础、Euclidean表述、关联函数、Krylov算法、实现审计 & 高级小 \(x\) 专题、部分场论推导 \\
数值与统计 & 线性代数、统计推断、误差协议、算法手册 & 历史证据矩阵、EIC细节 \\
现象学与实验 & 光锥定义、匹配、误差协议、高级专题、数据产品 & 大规模系综生成与求解器细节 \\
\bottomrule
\end{tabularx}
\end{table}

\section*{符号阅读法}

同一符号在不同文献中可能有不同含义。本书优先使用以下规则：
\begin{itemize}
  \item 坐标空间的横向分离用粗体 \(\bT\)，其模长写作 \(b_T\)；
  \item 动量空间的横向动量用 \(\kT\)，TMD分布可写成 \(f(x,\kT)\) 或Fourier共轭 \(f(x,\bT)\)；
  \item 准TMD用波浪号 \(\widetilde f\)，光锥TMD不加波浪号；
  \item 重正化标度用 \(\mu\)，rapidity标度用 \(\zeta\)，大动量用 \(P_z\)；
  \item 表示标签 \(q,g\) 只在必要处写出，未写时不表示两者相同。
\end{itemize}
遇到公式时，先看指标属于色空间、Lorentz空间、Dirac空间还是统计空间，再判断能否代入数值。

\section*{算法阅读法}

本书中的算法框不是代码库替代品，而是用于检查“输入、状态、循环、分支、停止条件和输出”是否完整。阅读一个算法时，建议依次问：
\begin{enumerate}
  \item 它求解的对象是物理方程、矩阵方程还是统计估计；
  \item 哪些量必须来自同一组态或同一replica；
  \item 停止条件是否检查真残差、协方差或物理窗口；
  \item 失败时是返回诊断状态、切换算法还是停止流水线；
  \item 复杂度和内存是否与报告中的实现一致。
\end{enumerate}
算法索引表~\ref{tab:algorithm-index} 可作为快速入口。

\section*{课程与练习建议}

若把本书作为一学期专题课，可按十六周安排：
\begin{enumerate}
  \item 第1--2周：QCD记号、Wilson线与光锥定义；
  \item 第3--4周：TMD因子化、rapidity与CS演化；
  \item 第5--6周：格点路径积分、链接与Wilson/Clover费米子；
  \item 第7--8周：两点/三点函数、源汇与谱学；
  \item 第9--10周：Krylov求解器、预条件与deflation；
  \item 第11--12周：重正化、软减除、hybrid与匹配；
  \item 第13--14周：Fourier逆问题、bootstrap与外推；
  \item 第15--16周：实现审计、复现包和高级专题。
\end{enumerate}
每两周至少完成一个“小系统可手算、代码可验证”的练习，例如二乘二CG、指数尾部Fourier、相关比值bootstrap和plaquette纯规范测试。

\section*{本书不替代什么}

本书不替代原始论文、格点系综文档或生产代码。它提供的是概念链、公式约定、算法结构与验收标准。遇到具体数值时，仍应回到原始文献、代码版本、组态元数据和完整误差预算。若正文与原始论文在约定上冲突，应以论文和可复现实验证据为准，并在项目文档中记录差异。

\mainmatter
\part{场论、格点与胶子结构基础}
\chapter{引言}

\section{从一维 PDF 到胶子 TMD}

共线部分子分布函数 \(f_{i/h}(x,\mu)\) 描述强子中携带动量份额 \(x\) 的部分子数密度，已在深度非弹性散射、Drell--Yan 过程和高 \(p_T\) 强子产生中得到系统应用。共线 PDF 对横向动量积分，因此不保留横向动量守恒和软辐射造成的方位角关联。当实验观测到硬过程中末态粒子或弱规范玻色子的横向动量 \(q_T\) 远小于硬标度 \(Q\) 时，TMD 因子化成为组织多尺度 QCD 的计算框架\cite{Collins2011}。此时截面不仅依赖 \(x\)，还依赖横向分离 \(\bT\) 或横向动量 \(\kT\)，从而直接编码核子的三维部分子图像。

胶子在核子动量中占有重要份额，也是 Higgs 产生、Drell--Yan、重夸克产生和小 \(x\) 演化中的主导部分子之一。胶子 TMD 还与胶子饱和、横向极化胶子、质子自旋分解和 EIC 中的方位角不对称性相关\cite{JiReview2021,Meissner2007}。核子非极化胶子 TMD-PDF \(\fone(x,\kT;\mu,\zeta)\) 可作为最基础的基准量：它不依赖核子极化，却仍然包含 UV 重正化标度 \(\mu\)、rapidity 标度 \(\zeta\)、规范链方向和软胶子减除等全部关键理论结构。

本文讨论的是\textbf{理论综述}，目标是建立一条可以逐项审计的计算链，而不是宣布某个现有数值为最终物理结果。文章采用以下证据优先级：
\begin{enumerate}
  \item 原始理论论文和库内全文转排用于定义算符、因子化公式、重正化方案与发表结果；
  \item 库内教科书和专题笔记用于补充 Wilson 线、格点场强、关联函数和统计方法；
  \item \texttt{refer/git-rep/PyQCD/} 源码与 \texttt{myqcd/} 公式审计仅作为实现和代数一致性的证据；
  \item 历史报告和诊断运行不是新的物理测量，不能替代多格距、多流时、同方案软因子和真实系综验证。
\end{enumerate}

\section{物理对象与记号}

采用 Minkowski 度规 \(g^{\mu\nu}=\mathrm{diag}(1,-1,-1,-1)\)，并定义光锥坐标
\begin{equation}
  v^\pm=\frac{v^0\pm v^3}{\sqrt{2}},
  \qquad
  v\cdot w=v^+w^-+v^-w^+-v_T\cdot w_T .
  \label{eq:lightcone}
\end{equation}
核子沿 \(z\) 方向以大动量 \(P_z\) 运动，\(P^\mu=(P^0,0,0,P_z)\)。胶子的纵向动量份额为 \(x=k^+/P^+\)，横向动量与其 Fourier 共轭变量分别为
\begin{equation}
  \bT=(b^1,b^2),\qquad \kT=(k^1,k^2),\qquad
  f(x,\bT)=\int\frac{\dd^2\kT}{(2\pi)^2} e^{i\kT\cdot\bT}f(x,\kT).
  \label{eq:Fourier2d}
\end{equation}
逆变换为
\begin{equation}
  f(x,\kT)=\int\dd^2\bT\,e^{-i\kT\cdot\bT}f(x,\bT).
  \label{eq:Fourier2dinverse}
\end{equation}
式~\eqref{eq:Fourier2d} 与式~\eqref{eq:Fourier2dinverse} 表明，格点上的有限 \(b_T\) 截断和实验中的 \(k_T\) 分布不是两个独立对象，而是同一关联函数在不同基底下的表示。任何相位、归一化或截断变化都必须成对处理。
表~\ref{tab:symbols} 汇总了后文反复使用的标度。\textbf{符号本身不是方案}：只有当 UV 方案、rapidity 方案、规范链方向、色表示和硬匹配核同时固定时，同一个符号才代表同一个物理量。

\begin{table}[htbp]
\centering
\caption{胶子 TMD 计算中的主要变量与物理作用。}
\label{tab:symbols}
\small
\begin{tabularx}{\textwidth}{L{2.3cm}L{3.2cm}Y}
\toprule
变量 & 物理意义 & 验收要点 \\
\midrule
\(x\) & 胶子纵向动量份额 & 应注明是否使用 \(G(x)=xg(x)\) 的动量密度约定 \\
\(\kT,\bT\) & 横向动量及共轭空间分离 & TMD 在 \(b_T\) 空间更便于建立规范不变算符和因子化 \\
\(\mu\) & UV 重正化标度 & 必须与 Wilson 系数、CS 核和软函数处于同一方案 \\
\(\zeta\) & rapidity 重正化标度 & 由规范链和软因子定义，不应与 \(P_z^2\) 混为一谈 \\
\(\Pz\) & 核子 boost & 控制 LaMET 幂修正，大动量通常要求更小格距和更高统计量 \\
\(a,L_s,T_s\) & 格距、空间与时间长度 & 决定连续极限、有限体积和激发态控制 \\
\(L_\parallel\) & staple 纵向臂长 & 用于实现 rapidity 正则化和 pinch-pole 减除；必须检查平台或极限 \\
\(\tau,r_f\) & 流时间与流半径 \(r_f=\sqrt{8\tau}\) & 形成独立于 \(a\) 的 UV 尺度，需满足短流时窗口 \\
\bottomrule
\end{tabularx}
\end{table}

\section{文献范围与证据映射}

本文不是对一个封闭数据库作 PRISMA 式系统综述，而是在 MyQCD 库内进行尽可能完整的理论核对。库内 \texttt{refer/papers/INDEX.md} 收录 50 篇原文转排及其中文译本，另有 38 篇评注导读、163 项 PyQCD 文档快照、若干原始 PDF 和可执行公式审计。表~\ref{tab:source-map} 给出与胶子 TMD 直接相关的来源分层。凡文献只讨论共线胶子 PDF、夸克 TMD 或纯梯度流方法时，本文只把它作为方法学先例，不把其结论跨通道升级。

\begin{table}[htbp]
\centering
\caption{本文使用的核心文献与证据类型。}
\label{tab:source-map}
\scriptsize
\begin{tabularx}{\textwidth}{L{3.0cm}L{3.2cm}Y}
\toprule
文献组 & 代表工作 & 在本文中的用途与边界 \\
\midrule
LaMET 基础 & Ji 2013, 2014；Ji--Zhang--Zhao 2018；Izubuchi 等 2018；Ji 等 2021 & 定义准算符、因子化、重正化和大动量展开；不提供胶子 TMD 数值结论\cite{Ji2013,Ji2014,JiZhangZhao2018,Izubuchi2018,JiReview2021} \\
胶子准 PDF & Zhang 等 2019；Li--Ma--Qiu 2019；Balitsky 等 2020 & 给出胶子算符分类、可乘性重正化和匹配语言\cite{ZhangGluonLaMET2019,LiMaQiu2019,Balitsky2020} \\
胶子 PDF 格点 & Fan 等 2018；Fan--Lin 2021；Salas-Chavira 等 2022；Good--Yao--Lin 2026 & 验证场强、HYP/flow、Ioffe-time、hybrid-ratio 与核子先例；均为共线 PDF\cite{FanQuasiGluon2018,FanLinPion2021,SalasChavira2022,GoodYaoLin2026} \\
TMD 软结构 & Zhang 等 2020；Chu 等 2022, 2023；Zhang 等 2022 & 软函数、CS 核、Wilson 圈减除和五格距验证\cite{ZhangSoft2020,ChuCS2022,ChuSoft2023,Zhang2022TMD} \\
夸克 TMD 基准 & He 等 2024；Ma 等 2025 & 展示完整 TMD 链和核子应用；不能直接替代胶子伴随表示\cite{HeTMD2024,MaBM2025} \\
梯度流 & Lüscher 2010；Lüscher--Weisz 2011；Monahan--Orginos 2017；Monahan 2018 & 流方程、UV 有限性、短流时展开和准 PDF 匹配\cite{Luscher2010,LuscherWeisz2011,MonahanOrginos2017,Monahan2018} \\
格点算法 & Peardon 等 2009；Bali 等 2016；CLQCD/PyQCD 实现 & distillation、momentum smearing、关联函数和统计流程\cite{Peardon2009,Bali2016} \\
\bottomrule
\end{tabularx}
\end{table}

\section{核心公式地图}

为避免后文的长推导淹没主线，先用一个公式卡片概括四个必须同时闭合的层次：光锥定义、rapidity 演化、Euclidean 匹配和数值误差。这里只给出结构性摘要，具体约定分别在后文展开。

\begin{keyformula}{从定义到误差的四层公式}
\begin{align}
  \fone(x,\bT;\mu,\zeta)
  &=\frac12\delta_T^{ij}\Phi_{g,ij}(x,\bT;\mu,\zeta),
  \label{eq:formulamap-def}\\
  2\zeta\frac{\partial}{\partial\zeta}
  \ln f^{R}(x,\bT;\mu,\zeta)
  &=K_R(\bT;\mu),
  \label{eq:formulamap-cs}\\
  \widetilde f^{g}(x,\bT;P_z;\mu)
  &=\left[H_g\otimes f_g\right](x,\bT;\mu;\zeta_z)
  \frac{\exp\left[\frac12K_g\ln(\zeta_z/\zeta)\right]}
  {\sqrt{S_{I,g}(\bT;\mu)}}+\text{高扭度},
  \label{eq:formulamap-match}\\
  \left|\Delta O_{\rm solve}\right|
  &\lesssim\kappa_O\,\rho_i^{\rm true}\|\eta\|_2.
  \label{eq:formulamap-error}
\end{align}
\end{keyformula}

式~\eqref{eq:formulamap-def} 说明“对象是什么”，式~\eqref{eq:formulamap-cs} 说明它在快度方向如何演化，式~\eqref{eq:formulamap-match} 说明格点准 TMD 如何回到光锥物理，式~\eqref{eq:formulamap-error} 则把 Dirac 求解误差接入最终观测量。四者缺一不可；任一层的符号或方案变化都不能由另一层的拟合补偿。

\clearpage
\chapter{数学结构、场论语言与量纲分析}
\label{sec:toolbox}

\section{为什么需要一层方法工具箱}

后文出现的许多困难并不来自单个公式，而来自公式之间使用了不同的数学对象：场强是Lie代数取值张量，Wilson线是路径有序指数，准TMD是Euclidean矩阵元，TMD分布是含Dirac delta与光锥投影的分布函数，重正化因子是方案依赖矩阵，格点关联函数又是随机采样估计。若只记住最终公式，容易把“算符相等”“矩阵元相等”“分布函数相等”混为一谈。本节先建立一套共同语言，使后续每个公式都能回答四个问题：对象属于哪个空间，指标如何缩并，量纲由什么标度补偿，取极限时哪一层展开先于另一层。

本文采用的主要约定如下：
\begin{enumerate}
  \item Minkowski度规取 \(g^{\mu\nu}=\mathrm{diag}(1,-1,-1,-1)\)，光锥分量用 \(v^\pm=(v^0\pm v^3)/\sqrt2\) 时可保持 \(v^+w^-+v^-w^+=v^0w^0-v^3w^3\)；
  \item 四维动量积分记作 \(\int_p\equiv\int\dd^4p/(2\pi)^4\)，三维格点动量只取离散值 \(\mathbf p=2\pi\mathbf n/L\)；
  \item 色指标在基本表示中记作 \(i,j\)，在伴随表示中记作 \(a,b\)；重复指标按所在空间缩并；
  \item 未特别说明时，裸量用下标 \(0\)，重正化量用上标或下标 \(R\)，格点量保留显式 \(a\) 依赖；
  \item \(P_z\) 是强子沿 \(z\) 方向的动量，\(\zeta=(2xP_z)^2\) 或文中有明确定义的rapidity标度，不能与 \(\mu^2\) 混用。
\end{enumerate}

\section{自然单位、量纲与主导项}

在自然单位 \(\hbar=c=1\) 下，作用量无量纲，因此场和算符的工程量纲由动能项固定。规范场满足
\begin{equation}
  [A_\mu]=1,
  \qquad
  [F_{\mu\nu}]=2,
  \qquad
  [g]=0,
  \qquad
  [D_\mu]=1.
  \label{eq:toolbox-engineeringdim}
\end{equation}
四维中 \(g\) 的裸量无量纲，但重正化耦合随标度演化。一个局部算符 \(O_d\) 若有工程量纲 \(d\)，其矩阵元在标度 \(Q\) 下通常写成
\begin{equation}
  \langle O_d\rangle
  =Q^{d}\sum_{n}c_n(\mu/Q)\left(\frac{\Lambda}{Q}\right)^n,
  \label{eq:toolbox-powercounting}
\end{equation}
其中 \(d\) 由自由场量纲确定，\(n\) 来自非微扰或高扭度修正。TMD展开中的 \(1/P_z^2\)、\(1/b_T^2\) 与 \(M_N^2/P_z^2\) 都是在比较同一物理量与相应硬标度。

判断一项是否“小”不能只看符号。若 \(\Lambda_{\rm QCD}\sim0.2~\mathrm{GeV}\)，则 \(P_z=2.0~\mathrm{GeV}\) 时 \((0.2/2.0)^2=10^{-2}\)，但在 \(P_z=0.8~\mathrm{GeV}\) 时已经达到 \(6.25\times10^{-2}\)。因此同一 \(1/P_z^2\) 修正能否忽略取决于目标精度、拟合窗口和系数大小。量纲分析只告诉缩放形式，不替代系数计算。

\begin{physicalbox}{尺度的三层含义}
\(\mu\) 是重正化或因子化标度，\(\zeta\) 是rapidity标度，\(P_z\) 是 Euclidean 大动量。三者不能用同一个变量代替。改变 \(\mu\) 不应改变物理截面，但会改变算符矩阵元与Wilson系数的分配；改变 \(\zeta\) 由Collins--Soper核演化；改变 \(P_z\) 同时改变幂修正、统计噪声和微扰匹配的适用范围。
\end{physicalbox}

\section{SU(3)色代数与表示}

夸克场属于基本表示，胶子场属于伴随表示。生成元取Hermitian矩阵 \(T^a\)，满足
\begin{equation}
  [T^a,T^b]=if^{abc}T^c,
  \qquad
  \Tr(T^aT^b)=T_F\delta^{ab},
  \qquad
  T_F=\frac12.
  \label{eq:toolbox-generators}
\end{equation}
基本表示与伴随表示的二次Casimir分别为
\begin{equation}
  C_F=\frac{N_c^2-1}{2N_c}=\frac43,
  \qquad
  C_A=N_c=3.
  \label{eq:toolbox-casimir}
\end{equation}
它们直接控制软胶子辐射的强度。若把伴随Wilson线的自能近似成有效质量项，线性发散系数与 \(C_A\) 成正比；若讨论夸克传播子，则软辐射强度由 \(C_F\) 控制。把这两种Casimir互换是胶子TMD计算中最常见的量纲正确但物理错误的来源之一。

色单态算符的矩阵元必须与表示匹配。胶子场强双线性在规范链连接后整体属于伴随表示，但在TMD定义的横向指标与Dirac/Lorentz投影中，最终外部态是色单态。实际计算要先用色代数收缩内部指标，再讨论剩余的Lorentz结构。一个看似仅含 \(F_{\mu\nu}\) 的表达式，若未与正确伴随链相乘，通常并不规范不变。

\section{规范协变性、平行移动与Wilson线}

格点链接定义为
\begin{equation}
  U_\mu(x)=\mathcal P\exp\left[ig\int_x^{x+a\hat\mu}\dd s\,A_\mu(s)\right],
  \qquad
  U_{-\mu}(x)=U_\mu^\dagger(x-a\hat\mu).
  \label{eq:toolbox-linkdef}
\end{equation}
在规范变换 \(\Omega(x)\) 下，基本表示链接满足
\begin{equation}
  U_\mu(x)\to \Omega(x)U_\mu(x)\Omega^\dagger(x+a\hat\mu).
  \label{eq:toolbox-linkgauge}
\end{equation}
伴随表示只需把 \(\Omega\) 换成 \(A_{ab}=2\Tr(T^a\Omega T^b\Omega^\dagger)\)。沿路径 \(\Gamma\) 的Wilson线把两端规范变换重新接起来：
\begin{equation}
  W_\Gamma(x,y)=\mathcal P\exp\left[ig\int_\Gamma A_\mu\,\dd x^\mu\right],
  \qquad
  W_\Gamma\to\Omega(x)W_\Gamma\Omega^\dagger(y).
  \label{eq:toolbox-wilsonline}
\end{equation}
若路径从 \(x\) 到无穷远再返回，必须区分两条腿的表示和方向；TMD中的未来指向与过去指向不是同一个算符，它们对应不同的软函数和过程。

\section{Dirac代数、旋量与投影}

Clifford代数为
\begin{equation}
  \{\gamma^\mu,\gamma^\nu\}=2g^{\mu\nu},
  \qquad
  \gamma_5=i\gamma^0\gamma^1\gamma^2\gamma^3,
  \qquad
  \gamma_5^2=1.
  \label{eq:toolbox-clifford}
\end{equation}
常用投影算符为
\begin{equation}
  P_\pm=\frac{1\pm\gamma_5}{2},
  \qquad
  \Sigma^{\mu\nu}=\frac{i}{2}[\gamma^\mu,\gamma^\nu].
  \label{eq:toolbox-projectors}
\end{equation}
核子矩阵元中出现的 \(\gamma^\mu\)、\(\gamma^\mu\gamma_5\) 和 \(\sigma^{\mu\nu}\) 分别探测不同自旋结构。胶子TMD的非极化投影通常先对场强指标做横向对称无迹组合，再与核子自旋平均。不能因为胶子本身不带电荷就省略Dirac结构；外部核子态仍由旋量描述，投影的选择决定矩阵元是否为实数、是否含反称张量以及在高扭度中如何混合。

对Euclidean计算，常用Hermitian矩阵 \(\gamma_5\) 与 \(\gamma_\mu\) 的性质会随具体表示改变。代码中的 \(\gamma_5\)-Hermiticity、Clover系数和手征投影必须与文章约定一致。论文只写一个 \(\gamma_5\) 而不说明表示，读者无法复核传播子或三点函数。

\section{路径积分、相关函数与矩阵元}

格点QCD的路径积分写成
\begin{equation}
  Z=\int\mathcal D U\,\mathcal D\bar\psi\,\mathcal D\psi\,
  e^{-S_g[U]-S_f[\bar\psi,\psi,U]},
  \label{eq:toolbox-pathintegral}
\end{equation}
费米子场是Grassmann数，积分后得到
\begin{equation}
  \int\mathcal D\bar\psi\,\mathcal D\psi\,
  e^{-\bar\psi D[U]\psi}
  =\det D[U].
  \label{eq:toolbox-fermiondet}
\end{equation}
一个规范不变算符 \(O\) 的真空期望值为
\begin{equation}
  \langle O\rangle
  =\frac{1}{Z}\int\mathcal DU\,O[U,D^{-1}[U]]\det D[U]e^{-S_g[U]}.
  \label{eq:toolbox-vev}
\end{equation}
传播子 \(D^{-1}\) 不是可观测量，只有闭合收缩后的规范不变量才是。手征极限下 \(\det D\) 可能接近零，导致HMC力增大、拓扑冻结和条件数恶化；因此“能否求逆”与“采样是否覆盖正确拓扑扇区”是两个独立问题。

路径积分中的算符插入位置必须保持时序与规范链。对于非局域Wilson线算符，传播子收缩不能只看图形；端点的伴随指标、路径方向和反粒子传播子的 \(\gamma_5\)-Hermiticity都要在同一约定下处理。

\section{谱分解、正定性与物理可解释性}

两点函数在固定动量下插入完备态：
\begin{equation}
  C_2(t,\mathbf p)
  =\sum_n \frac{\langle 0|O|\mathbf p,n\rangle\langle\mathbf p,n|O^\dagger|0\rangle}{2E_n(\mathbf p)}
  e^{-E_n(\mathbf p)t}.
  \label{eq:toolbox-spectral}
\end{equation}
在Euclidean时间 \(0<t<T\) 的有限环面上，还应包括反向传播项。忽略周期性项会在大 \(t\) 处制造假平台。基态能量为 \(E_0=\sqrt{M^2+\mathbf p^2}\)，若拟合出的 \(E_0\) 与色散关系不符，后续矩阵元即使数值稳定也不应直接采信。

谱分解提供三种诊断：早期由高激发态控制，中期可能出现基态平台，晚期信噪比下降。一个“平台”若只在两三个时间片出现且随源、汇算符变化，通常是拟合自由度过多或统计涨落，而不是物理信号。TMD矩阵元必须沿用与两点函数相同的源、汇和态投影，避免用不同态定义比值。

\section{傅里叶约定、分布与有限体积}

连续Fourier变换采用
\begin{equation}
  \widetilde f(x)=\int\frac{\dd z}{2\pi}\,e^{ixP_zz}f(z),
  \qquad
  f(z)=\int\dd x\,e^{-ixP_zz}\widetilde f(x),
  \label{eq:toolbox-fourierconv}
\end{equation}
有限体积中 \(z\) 只能取离散格点，\(P_z\) 也只能取离散动量，积分变成求和加插值。若代码使用 \(e^{ipz}\) 而非 \(e^{-ipz}\)，则会得到复共轭或反射后的分布；若同时缺失 \(2\pi\)，归一化看起来“差一个因子”，但端点与正定性判据也会随之改变。

有限体积导致两类误差：一是动量量子化使 \(P_z\) 不能连续靠近目标值；二是大 \(z\) 的边界效应。后者在准TMD中常表现为 \(z\) 接近 \(L/2\) 时两端Wilson链与周期性镜像相遇，若未检查 \(z_{\max}<L/2\) 与 \(L\to\infty\) 趋势，就不能把尾部当作物理衰减。

\begin{figure}[htbp]
\centering
\begin{tikzpicture}[
  box/.style={draw=qcdnavy!70!black,rounded corners=2pt,fill=qcdmist,align=center,minimum width=2.8cm,minimum height=0.9cm},
  arrow/.style={-{Latex[length=2mm]},thick,qcdrust!80!black}
]
\node[box] (hard) {硬标度 \(Q\)\\微扰匹配};
\node[box,right=1.0cm of hard] (coll) {共线标度\\\(xP_z,\;1/b_T\)};
\node[box,right=1.0cm of coll] (soft) {软/rapidity\\非微扰核};
\node[box,below=0.8cm of coll] (lat) {格点矩阵元\\\(a,P_z,L\)};
\node[box,below=0.8cm of hard] (ren) {UV重正化\\\(\mu,\mu_0\)};
\draw[arrow] (hard) -- (coll);
\draw[arrow] (coll) -- (soft);
\draw[arrow] (lat) -- (coll);
\draw[arrow] (ren) -- (hard);
\end{tikzpicture}
\caption{TMD计算中四类标度的不可替代关系。箭头表示演化或匹配方向，不表示数值大小。}
\label{fig:toolbox-scales}
\end{figure}

\section{重正化群、方案依赖与方案无关性}

一个重正化算符满足
\begin{equation}
  O_R(\mu)=Z_O(\mu,\mu_0)O_R(\mu_0),
  \qquad
  \mu\frac{\dd}{\dd\mu}O_R(\mu)=\gamma_O(g)O_R(\mu).
  \label{eq:toolbox-rge}
\end{equation}
物理矩阵元还应乘上相应的Wilson系数。若匹配公式写成 \(F=C(\mu)O_R(\mu)\)，则有
\begin{equation}
  \mu\frac{\dd}{\dd\mu}C=-\gamma_O C,
  \label{eq:toolbox-wilsonrge}
\end{equation}
使得乘积与 \(\mu\) 无关。TMD匹配中的 \(H\)、CS核与软函数必须使用同一光锥或rapidity方案；把一个在 \(\MSbar\) 下计算的系数乘到一个在ratio方案下定义的矩阵元，通常会留下人为的 \(\mu\) 依赖。

\section{有效算符、可乘性重正化与混合矩阵}

有效场论的语言不是“把高能自由度删掉”，而是把其效应压进局域算符的Wilson系数。若低能算符 \(O_i\) 在重正化后互相混合，则
\begin{equation}
  O_{i,R}(\mu)=Z_{ij}(\mu,\mu_0)O_{j,R}(\mu_0),
  \qquad
  \mu\frac{\dd}{\dd\mu}O_{i,R}=\gamma_{ij}O_{j,R}.
  \label{eq:toolbox-mixing}
\end{equation}
若 \(Z\) 可对角化且本征算符有确定反常量纲，则称该基底可乘性重正化；否则应保留矩阵结构。对于胶子TMD，纯胶子准算符与稀味匹配层的混合是两个不同问题：前者是UV算符空间的基底选择，后者是把准分布转成光锥分布时的味单态硬匹配。把两者合并会让“单态混合”一词失去可操作含义。

\section{极限交换与奇异展开}

物理量常写成
\begin{equation}
  F(a,P_z)=\sum_{n\ge0}\frac{c_n(P_z)}{a^n}
  +\sum_{m\ge1}\frac{d_m(P_z)}{P_z^m},
  \label{eq:toolbox-expansions}
\end{equation}
但这只是示意。正确顺序取决于先取哪个极限。如果 \(aP_z\) 固定，则 \(a^{-1}\) 与 \(P_z\) 的幂次不能分别处理；若先取 \(a\to0\) 再取 \(P_z\to\infty\)，重正化与匹配的定义可能与格点可实现的顺序不同。LaMET的核心正是用匹配核组织两类极限的差异，并在
\begin{equation}
  \Lambda_{\rm QCD},M_N\ll P_z\ll a^{-1}
  \label{eq:toolbox-lametwindow}
\end{equation}
的窗口内做外推。若 \(P_z\) 接近截断，微扰匹配与格点幂修正的展开参数都变大，外推不是“多拟合几个点”能修复的。

\section{一致性检查清单}

在进入光锥定义以前，应能独立检查以下项目：
\begin{checkpointbox}{数学约定自检}
\begin{enumerate}
  \item 每个算符的色表示、Lorentz指标、Dirac指标和路径方向是否明确；
  \item 量纲是否平衡，是否存在未说明的 \(a^{-1}\)、\(P_z\) 或 \(\mu\)；
  \item 传播子、\(\gamma_5\)-Hermiticity和Fourier相位是否使用同一约定；
  \item 两个极限的取法是否写清，是否能退化为已知的局域算符或自由场结果；
  \item 重正化因子是否为标量、向量或矩阵，混合是否发生在正确的物理层；
  \item 端点、零模、周期镜像与大 \(z\) 边界是否被显式排除或纳入误差。
\end{enumerate}
\end{checkpointbox}

这些检查并不增加物理内容，却能显著减少把方案差异误判为数值信号的概率。下面的光锥定义将对每个项目给出具体答案。

\clearpage
\chapter{光锥胶子 TMD 的理论定义}

\section{场强双线性与伴随规范链}

QCD 场强张量取为
\begin{equation}
  F_{\mu\nu}=\frac{i}{g}[D_\mu,D_\nu],
  \qquad
  D_\mu=\partial_\mu-igA_\mu,
  \qquad
  A_\mu=A_\mu^aT^a,
  \label{eq:fieldstrength}
\end{equation}
其中基本表示生成元取厄米约定，并统一采用
\begin{equation}
  \Tr(T^aT^b)=\frac{1}{2}\delta^{ab},
  \qquad
  C_F=\frac{N_c^2-1}{2N_c}=\frac{4}{3},
  \qquad
  C_A=N_c=3 .
  \label{eq:generatornormalization}
\end{equation}
胶子 TMD 使用场强而不是胶子势 \(A_\mu\)，原因是非局域算符必须在规范变换下保持协变，而 \(F_{\mu\nu}\) 在 \(b_T=0\) 的局域极限直接连接到 twist-2 胶子算符和胶子动量分数\cite{Collins2011,JiReview2021}。

在伴随表示中，规范链为
\begin{equation}
  U_{\Gamma,ab}(\bT)=\left[\mathcal P\exp\left(-ig\int_\Gamma \dd z^\mu A_\mu(z)\right)\right]_{ab},
  \label{eq:adjlink}
\end{equation}
把指数逐阶展开，可以清楚看到 Wilson 线的物理内容：
\begin{equation}
  U_{\Gamma,ab}
  =\delta_{ab}
  -ig\int_\Gamma \dd z^\mu A^c_\mu(z)(T^c)_{ab}
  +\frac{(-ig)^2}{2!}\mathcal P\int_\Gamma \dd z^\mu\dd w^\nu
  [A_\mu(z)A_\nu(w)]_{ab}
  +\cdots .
  \label{eq:adjlinkexpansion}
\end{equation}
第 \(n\) 阶项对应交换 \(n\) 个软规范量子，路径顺序记录颜色算符如何沿传播路径排列；指数形式则把这些交换重求和为规范协变的传播相位。
路径 \(\Gamma\) 连接场强算符的两个端点。若算符放在 \(0\) 与 \(z\hat z+\bT\)，准 TMD 的典型 staple 路径可写成
\begin{equation}
  \Gamma_\pm:\quad
  0\longrightarrow \pm L_\parallel\hat z
  \longrightarrow \pm L_\parallel\hat z+\bT
  \longrightarrow z\hat z+\bT .
  \label{eq:staplepath}
\end{equation}
正负方向对应初态或末态相互作用相反的规范链。对非极化 \(f_1^g\)，在严格 leading twist 下两种方向可属于同一普适类；对 T-odd 或自旋相关分布，方向改变会产生物理效应\cite{Collins2011,JiReview2021}。因此，格子计算不能只比较数值曲线，还必须声明路径方向。

\begin{physicalbox}{规范链的物理读法}
Wilson 线不是把两个场强点“连接起来”的数值工具，而是对软胶子交换的协变重求和。改变路径端点、纵向臂长或表示，等同于改变被重求和的软模式集合；因此裸矩阵元、Wilson 圈减除和软函数必须共享同一几何。路径方向对非极化 leading-twist 胶子可能落入同一普适类，但对 T-odd 与自旋相关分布则可以改变符号或相位。
\end{physicalbox}

\begin{figure}[htbp]
\centering
\begin{tikzpicture}[>=Stealth,font=\small,scale=0.88,transform shape]
  \tikzset{
    link/.style={line width=1.35pt,blue!60!black},
    soft/.style={line width=1.1pt,red!65!black,decorate,decoration={snake,amplitude=1.2pt,segment length=5pt}},
    field/.style={circle,fill=orange!80!black,inner sep=1.8pt},
    label/.style={align=center}
  }
  % Light-cone panel
  \node[font=\bfseries] at (0,3.15) {(a) 光锥 TMD};
  \draw[->] (-1.8,0)--(1.8,0) node[right] {\(z^-\)};
  \draw[->] (0,-0.3)--(0,2.3) node[above] {\(z^+\)};
  \draw[link] (-1.1,0.55)--(1.1,0.55);
  \draw[link] (-1.1,1.55)--(1.1,1.55);
  \draw[soft] (0,0.55)--(0,1.55);
  \node[field] at (-1.1,0.55) {};
  \node[field] at (1.1,1.55) {};
  \node[label,below] at (0,-0.55) {两条光样 Wilson 线\\+ 软函数};

  % Quasi panel
  \begin{scope}[xshift=5.5cm]
  \node[font=\bfseries] at (0,3.15) {(b) Euclidean 准 TMD};
  \draw[->] (-2.0,0)--(2.2,0) node[right] {\(z\)};
  \draw[->] (0,-0.3)--(0,2.35) node[above] {\(b_T\)};
  \draw[link] (-0.9,0.35)--(-0.9,1.75)--(1.3,1.75)--(1.3,1.0);
  \node[field] at (-0.9,0.35) {};
  \node[field] at (1.3,1.0) {};
  \draw[dashed,red!60!black] (-0.9,1.75)--(1.3,1.75) node[midway,above] {rapidity 臂};
  \node[label,below] at (0,-0.55) {有限 \(L_\parallel\) staple\\+ Wilson 圈减除};
  \end{scope}

  % Matching panel
  \begin{scope}[xshift=10.8cm]
  \node[font=\bfseries] at (0,3.15) {(c) 匹配与极限};
  \node[draw,rounded corners,fill=blue!7,minimum width=2.2cm,minimum height=0.7cm] (q) at (0,2.15) {准 TMD};
  \node[draw,rounded corners,fill=green!8,minimum width=2.2cm,minimum height=0.7cm] (r) at (0,1.05) {重正化矩阵};
  \node[draw,rounded corners,fill=orange!10,minimum width=2.2cm,minimum height=0.7cm] (m) at (0,-0.05) {硬核/软因子};
  \node[draw,rounded corners,fill=red!7,minimum width=2.2cm,minimum height=0.7cm] (p) at (0,-1.15) {光锥 TMD};
  \draw[->,thick] (q)--(r); \draw[->,thick] (r)--(m); \draw[->,thick] (m)--(p);
  \draw[->,dashed] (2.0,2.15) to[bend left=25] node[right,align=left] {\(a\to0\)\\\(L\to\infty\)} (2.0,0.95);
  \draw[->,dashed] (2.0,0.95) to[bend left=25] node[right,align=left] {\(P_z\to\infty\)\\\(\tau\to0\) 窗口} (2.0,-1.05);
  \end{scope}
\end{tikzpicture}
\caption{胶子 TMD 的三个理论层次。光锥算符给出物理定义；Euclidean staple 算符使格点 QCD 可以直接计算；因子化、软/rapidity 减除和多尺度极限把格点矩阵元连接到光锥分布。}
\label{fig:three-levels}
\end{figure}

图~\ref{fig:three-levels} 概括了本文反复使用的三层结构：光锥算符提供定义，Euclidean staple 提供可计算对象，soft/rapidity、重正化、匹配和极限把两者连接起来。后续所有公式和流程都应能放回图中的某一条边，不能把不同层次的对象直接等同。

\section{非极化投影与共线极限}

定义裸胶子关联函数
\begin{equation}
  \Phi^{ij}_{g}(x,\bT)
  =\frac{1}{xP^+}\int\frac{\dd\xi^-}{2\pi}\,
  e^{-ixP^+\xi^-}
  \langle P|\,
  F^{+i}_a(\xi^-,\bT)\,
  U^{ab}_{\Gamma}(\xi^-,\bT;0)\,F^{+j}_b(0)\,|P\rangle ,
  \label{eq:gluoncorrelator}
\end{equation}
其中横向指标 \(i,j=1,2\)，色指标在伴随表示中缩并。非极化胶子分布由横向无迹投影给出
\begin{equation}
  \fone(x,\bT;\mu,\zeta)
  =\frac{1}{2}\delta^{ij}_{T}\,
  \Phi_{g,ij}(x,\bT;\mu,\zeta),
  \label{eq:f1projection}
\end{equation}
其中整体归一化应随场强和生成元约定一起固定。若需要研究线性极化胶子，可定义无迹对称投影
\begin{equation}
  h_{1}^{\perp g}(x,\bT;\mu,\zeta)
  \sim
  \left(\frac{b_T^ib_T^j}{b_T^2}-\frac{1}{2}\delta_T^{ij}\right)
  \Phi_{g,ij}(x,\bT;\mu,\zeta),
  \label{eq:h1perpprojection}
\end{equation}
其具体归一化和符号必须与所用 TMD 分解约定一致。本文的主线是非极化 \(f_1^g\)；其它极化胶子分布的原则相同，但会增加自旋投影、手征结构和规范链方向相关的复杂性\cite{Meissner2007}。

在 \(\bT\to0\) 且对横向动量积分后，TMD 应退化为共线胶子 PDF：
\begin{equation}
  f_{g/h}(x,\mu)
  =\int\dd^2\kT\,\fone(x,\kT;\mu,\zeta)
  \quad\text{(leading twist)},
  \label{eq:collinearlimit}
\end{equation}
并满足动量求和
\begin{equation}
  \int_0^1\dd x\,x\left[f_{g/h}(x,\mu)+\sum_q f_{q/h}(x,\mu)\right]=1 .
  \label{eq:momentumsum}
\end{equation}
这两个关系是重要的物理一致性检查，但不能替代数值验证：在格点上，积分截断、单态混合、零模减除和重正化方案都可能改变端点行为。

\section{rapidity 发散、软函数与 Collins--Soper 演化}

光样 Wilson 线带有无限长的快度积分。即使 UV 维数正规化已经固定，软胶子与共线胶子仍可能在 rapidity 方向产生额外发散。因此，光锥 TMD 不能只依赖 \(\mu\)，还必须引入 rapidity 标度 \(\zeta\)，并通过软函数完成减除\cite{Collins2011,JiReview2021}。物理量对 \emph{rapidity 方案} 的依赖应在硬核、软函数和 TMD 的组合中抵消。

记表示相关的软函数为 \(S_R(\bT;\mu,Y_1,Y_2)\)。在因子化适用的方案中，它可分解成 rapidity 无关的内禀部分和 Collins--Soper 指数：
\begin{equation}
  S_R(\bT;\mu,Y_1,Y_2)
  =S_{I,R}(\bT;\mu)\,
  \exp\left[K_R(\bT;\mu)\,\mathcal Y(Y_1,Y_2)\right],
  \label{eq:softfactor}
\end{equation}
其中 \(\mathcal Y\) 是依赖于正则化方案的快度组合，\(K_R\) 是表示 \(R\) 的 Collins--Soper 核。对胶子 TMD，Wilson 线处于伴随表示，软因子和 CS 核必须使用同一表示；夸克 TMD 中提取的数值不能在没有匹配关系的情况下直接替代胶子结果。

物理 TMD 的 rapidity 演化由
\begin{equation}
  2\zeta\frac{\partial}{\partial\zeta}
  \ln f^{R}(x,\bT;\mu,\zeta)
  =K_R(\bT;\mu)
  \label{eq:CSequation}
\end{equation}
描述。其 UV 演化满足
\begin{equation}
  \mu^2\frac{\partial}{\partial\mu^2}K_R(\bT;\mu)
  =-\Gamma^{R}_{\rm cusp}(\alpha_s),
  \label{eq:csrge}
\end{equation}
其中 \(\Gamma^{R}_{\rm cusp}\) 是表示相关的 cusp 反常维数。对胶子，领头项与 \(C_A\) 成正比；对夸克则为 \(C_F\)。这一表示依赖性是胶子 TMD 不能简单复用夸克软因子的根本原因之一。

在实际因子化公式中，常把未减除的 TMD 除以 \(\sqrt{S_R}\)。平方根来自软函数包含两条规范链，而单个 TMD 关联函数只含一条规范链。符号和幂次依赖于快度定义，但结构要求不变：\textbf{TMD、软函数、硬核和 CS 演化必须来自同一方案}。

\section{Lorentz 分解与胶子算符投影}

胶子算符含两个场强指标，因此比夸克双线性更容易发生 Lorentz 投影歧义。对核子自旋平均的矩阵元，最一般的分解可写成
\begin{align}
  M_{\mu\alpha;\lambda\beta}(z,p)
  &=\mathcal M_{pp}\,
  \big(g_{\mu\lambda}p_\alpha p_\beta
  -g_{\mu\beta}p_\alpha p_\lambda
  -g_{\alpha\lambda}p_\mu p_\beta
  +g_{\alpha\beta}p_\mu p_\lambda\big)\nonumber\\
  &\quad+\mathcal M_{zz}\,
  \big(g_{\mu\lambda}z_\alpha z_\beta
  -g_{\mu\beta}z_\alpha z_\lambda
  -g_{\alpha\lambda}z_\mu z_\beta
  +g_{\alpha\beta}z_\mu z_\lambda\big)\nonumber\\
  &\quad+\mathcal M_{zp}\,
  \big(g_{\mu\lambda}z_\alpha p_\beta
  -g_{\mu\beta}z_\alpha p_\lambda
  -g_{\alpha\lambda}z_\mu p_\beta
  +g_{\alpha\beta}z_\mu p_\lambda\big)\nonumber\\
  &\quad+\mathcal M_{pz}\,
  \big(g_{\mu\lambda}p_\alpha z_\beta
  -g_{\mu\beta}p_\alpha z_\lambda
  -g_{\alpha\lambda}p_\mu z_\beta
  +g_{\alpha\beta}p_\mu z_\lambda\big)\nonumber\\
  &\quad+\mathcal M_{ppzz}\,
  (p_\mu z_\alpha-p_\alpha z_\mu)
  (p_\lambda z_\beta-p_\beta z_\lambda)\nonumber\\
  &\quad+\mathcal M_{gg}\,
  \big(g_{\mu\lambda}g_{\alpha\beta}
  -g_{\mu\beta}g_{\alpha\lambda}\big),
  \label{eq:sixamplitudes}
\end{align}
其中六个不变振幅依赖于 \(z^2\) 与 Ioffe 时间 \(\nu=-p\cdot z\)\cite{Balitsky2020}。光锥胶子 PDF 由 \(\mathcal M_{pp}\) 决定：
\begin{equation}
  -\,\mathcal M_{pp}(\nu,0)
  =\frac{1}{2}\int_{-1}^{1}\dd x\,e^{-ix\nu}\,x f_g(x).
  \label{eq:Mpp}
\end{equation}
在 Euclidean 算符中，单独的 \(F^{zi}F^{zi}\) 或 \(F^{ti}F^{ti}\) 会混入 \(\mathcal M_{zz},\mathcal M_{gg}\) 或其它结构。为了提取 twist-2 的 \(\mathcal M_{pp}\)，需要选择线性组合，例如
\begin{equation}
  \mathcal M_{0i;i0}+\mathcal M_{ij;ij}
  =2p_0^2\,\mathcal M_{pp}
  +\text{高扭度/接触项},
  \label{eq:Mppcombination}
\end{equation}
其中横向指标按二维 Euclidean 空间缩并。该组合也是格点上常见的非极化胶子流投影之一\cite{Balitsky2020,FanQuasiGluon2018}。不同投影的 UV 反常量纲一般不同，因此“同一个胶子 PDF”不能由任意 Lorentz 分量等价地读出。

\section{胶子极化和 TMD 分类}

对非极化核子，最低阶的非极化胶子 TMD 为 \(f_1^g\)，它在 \(\bT\to0\) 时连接到胶子动量密度 \(xg(x)\)。线性极化胶子分布 \(h_1^{\perp g}\) 通过无迹横向张量投影提取，描述横向平面内的极化关联；在 \(k_T\) 空间，它表现为 \(k_T^ik_T^j\) 型的角度依赖\cite{Meissner2007}。核子极化后还可以定义螺旋度胶子、worm-gear 型和横向极化胶子 TMD；这些分布不能由 \(f_1^g\) 的标量倍数得到，因为它们对应不同的自旋和角度结构。

本文把非极化 \(f_1^g\) 作为主对象，原因是：
\begin{enumerate}
  \item 它不需要极化核子插值算符，先把 gluon loop、Wilson 线和 soft/rapidity 结构暴露出来；
  \item 它与共线胶子 PDF、动量求和和 QCD 演化有最直接的交叉检验；
  \item 它的 TMD 因子化和软函数在概念上最接近标准 Drell--Yan/SIDIS 结构，便于把纯理论误差与极化投影误差分开。
\end{enumerate}
极化胶子 TMD 的计算不是简单替换一个 Gamma 矩阵。场强指标、对偶张量、核子自旋投影和规范链方向会同时改变，且部分 T-odd 结构对过程方向敏感\cite{Egerer2022,JiReview2021}。

\section{过程依赖与规范链方向}

TMD 因子化中的 Wilson 线不是数学装饰，而是过程依赖的物理输入。Drell--Yan 的初态相互作用和 SIDIS 的末态相互作用对应相反的规范链方向。对非极化 \(f_1^g\)，在严格 leading twist 且采用兼容的软函数方案时，方向差异可以被吸收到统一的方案中；对 T-odd 或自旋相关胶子分布，方向变化可以给出符号或相位差异。因此，以下两种说法必须区分：
\begin{itemize}
  \item “同一个 TMD”是指在固定过程类和固定 rapidity 方案下的普适函数；
  \item “同一个裸格点算符”则不是过程无关的，路径方向、端点和软减除都改变数值。
\end{itemize}
在格点上模拟 \(+\hat z\) 与 \(-\hat z\) staple 并比较差值，可以作为过程依赖和 T-odd 结构的诊断；但只有配合相应的 soft/rapidity 方案和合法的 \(L_\parallel\to\infty\) 极限，差值才具有物理意义。

\section{物理极限与一致性检查}

一个胶子 TMD 结果至少应通过以下极限检查：
\begin{align}
  &\bT\to0:\qquad
  f_1^g(x,\bT;\mu,\zeta)\text{ 与短距 OPE 和胶子 PDF 的微扰匹配相容},
  \label{eq:limit-b0}\\
  &\int\dd^2\kT\,f_1^g(x,\kT;\mu,\zeta)
  =f_g(x,\mu)+\mathcal O(\text{scheme/higher twist}),
  \label{eq:limit-collinear}\\
  &\int_0^1\dd x\,x\,f_g(x,\mu)+\sum_q\int_0^1\dd x\,x\,f_q(x,\mu)=1,
  \label{eq:limit-momentum}\\
  &b_T\gg\Lambda_{\rm QCD}^{-1}:\qquad
  f_1^g\text{ 进入非微扰区，不能用固定阶微扰外推代替}.
  \label{eq:limit-largeb}
\end{align}
这些关系不是独立的数值拟合约束，而是方案一致性的物理要求。若某个外推在 \(\bT\to0\) 与共线 PDF 不一致，或在大 \(\bT\) 处仍使用未加控制的微扰核，就不能用“格点误差”掩盖方案错误。

\clearpage
\chapter{格点规范理论的基础对象}

\section{链接、平行移动与 Wilson 线}

格点规范场的自由度位于有向链接上。采用本节统一的显式耦合约定
\begin{equation}
  U_\mu(x)=\mathcal P\exp\!\left[i g_0\int_x^{x+a\hat\mu}\!\dd y^\nu A_\nu(y)\right],
  \qquad
  U_{-\mu}(x)=U_\mu(x-a\hat\mu)^\dagger ,
  \label{eq:link-definition}
\end{equation}
其中 \(A_\mu=A_\mu^aT^a\)，生成元满足 \(\Tr(T^aT^b)=\tfrac12\delta^{ab}\)。链接不是普通相位，而是把颜色参考系从 \(x\) 平行移动到 \(x+a\hat\mu\) 的群元素。若把 \(g_0A_\mu\) 合并记为 \(A_\mu\)，后续所有场强与 Clover 公式都应删去显式 \(g_0\)；两种约定不能混用。

在规范变换 \(\psi(x)\to\Omega(x)\psi(x)\) 下，
\begin{equation}
  U_\mu(x)\longrightarrow \Omega(x)U_\mu(x)\Omega^\dagger(x+a\hat\mu).
  \label{eq:link-gauge}
\end{equation}
沿路径的链接乘积因此只在两个端点变换。开 Wilson 线必须与端点场或场强收缩，闭合 Wilson 圈则取颜色迹，二者才构成规范不变量。这个简单的端点规则决定了准 PDF、准 TMD 和胶子非局域算符中路径方向、表示与指标顺序的全部物理含义。

连续 Wilson 线是路径编序指数
\begin{equation}
  W_\Gamma(y,x)=\mathcal P\exp\!\left[i g_0\int_\Gamma \dd z^\mu A_\mu(z)\right].
  \label{eq:wilson-line}
\end{equation}
格点上它由沿途链接的有序乘积代替。路径不是数值实现细节：改变两条长边的方向会改变软胶子重求和，改变颜色表示会改变 Casimir 与 cusp 反常维，交换端点和负方向则改变厄米共轭关系。

\section{Wilson 圈、静态势与有限温度}

设 \(W_{rT}\) 是空间距离为 \(r\)、时间跨度为 \(T\) 的矩形 Wilson 圈，并按颜色维数归一化。取其大 \(T\) 极限，
\begin{equation}
  \langle W_{rT}\rangle
  \sim \exp[-T V(r)]\left[1+\Ord(e^{-T\Delta E})\right],
  \label{eq:wilson-loop-potential}
\end{equation}
则 \(V(r)\) 是静态夸克--反夸克势。红外区域常用
\begin{equation}
  V(r)=A+\frac{B}{r}+\sigma r
  \label{eq:static-potential}
\end{equation}
描述自能常数、Coulomb 项和线性禁闭项。面积律 \(\langle W_C\rangle\propto e^{-\sigma\mathcal A_C}\) 是禁闭的格点判据之一；强耦合展开可给出面积律的首项，但不能代替连续极限下的非微扰测量。振幅、自能和常数项依赖归一化，只有斜率与差值的连续极限才具有直接的物理意义。

静态势也连接标度设定。Sommer 参数定义为
\begin{equation}
  r_0^2\left.\frac{\dd V}{\dd r}\right|_{r=r_0}=1.65.
  \label{eq:sommer}
\end{equation}
格点先计算无量纲 \(r_0/a\)；只有选定 \(r_0\) 的物理约定后，才能赋值 \(a\)。因此“\(r_0\simeq0.5~{\rm fm}\)”是外部标度输入，不是格点理论自身的预言。

在有限温度 \(T=1/(aN_T)\) 下，沿周期时间方向的闭合 Wilson 线给出 Polyakov 圈
\begin{equation}
  P(\mathbf x)=\Tr\prod_{j=0}^{N_T-1}U_4(\mathbf x,j).
  \label{eq:polyakov}
\end{equation}
纯规范理论具有中心对称性 \(P\to zP\)，\(z\in\mathbb Z_3\)；\(\langle P\rangle=0\) 区分禁闭相，\(\langle P\rangle\ne0\) 对应中心对称破缺。动力学费米子显式破缺中心对称性，因此在全 QCD 中 Polyakov 圈不是严格有序参量，禁闭--退禁闭应作为 crossover 处理。

\section{从 plaquette 到场强张量}

基本 plaquette 定义为
\begin{equation}
  P_{\mu\nu}(x)=U_\mu(x)U_\nu(x+\hat\mu)
  U_\mu^\dagger(x+\hat\nu)U_\nu^\dagger(x).
  \label{eq:b01-plaquette}
\end{equation}
对小 \(a\) 依次使用 Baker--Campbell--Hausdorff 公式和 \(A_\nu(x+a\hat\mu)=A_\nu(x)+a\partial_\mu A_\nu(x)+\cdots\)，得到
\begin{equation}
  P_{\mu\nu}(x)
  =\exp\!\left[i a^2g_0F_{\mu\nu}(x)+\Ord(a^3)\right],
  \qquad
  F_{\mu\nu}=\partial_\mu A_\nu-\partial_\nu A_\mu
  +ig_0[A_\mu,A_\nu].
  \label{eq:plaquette-field-strength}
\end{equation}
若在此推导中把 \(g_0\) 吸收进 \(A_\mu\)，指数中应写成 \(i a^2F_{\mu\nu}\)，场强对易项也相应地不含显式 \(g_0\)。

为减小单个 plaquette 的统计噪声并提高超立方对称性，取四个定向共面 plaquette 的对称和 \(Q_{\mu\nu}\)。所有定向都应使其领头项为 \(\mathbb 1+i a^2g_0F_{\mu\nu}\)。例如在常用定向约定下可写
\begin{equation}
  Q_{\mu\nu}=P_{\mu\nu}+P_{\nu,-\mu}
  +P_{-\mu,-\nu}+P_{-\nu,\mu}.
  \label{eq:clover-sum}
\end{equation}
于是可直接定义反厄米的格点曲率
\begin{equation}
  F^{c}_{\mu\nu}(x)
  =-\frac{i}{8g_0a^2}
  \left[Q_{\mu\nu}(x)-Q_{\mu\nu}^\dagger(x)\right].
  \label{eq:clover-field}
\end{equation}
在基本表示中还应去掉颜色迹：
\begin{equation}
  F_{\mu\nu}\longmapsto
  F_{\mu\nu}-\frac1{N_c}\Tr(F_{\mu\nu})\,\mathbb 1 .
  \label{eq:traceless-field}
\end{equation}
这一步应发生在场强与 Wilson 线收缩之前。式~\eqref{eq:clover-field} 的 \(g_0\) 与 \(a^2\) 因子是和式~\eqref{eq:link-definition} 绑定的一整套约定；若改用 \(U=e^{iaA}\) 或不同 \(Q\) 归一化，必须重新推导而不是只替换一个符号。

格点场强具有 \(\Ord(a^2)\) 的离散化误差。固定物理尺度并减小 \(a\) 时，它趋于连续场强；靠近截断时，单 plaquette 的高维算符、有限体积镜像和与平滑链接不一致的端点都会进入系统误差。对胶子算符而言，场强、Wilson 线、迹投影和 Fourier 约定属于同一个算符定义，不能分别“调参”。

\chapter{离散化误差、改进与标度设定}

\section{Symanzik 有效展开}

格点关联函数可在长距尺度上用连续 QCD 加局域高维算符展开：
\begin{equation}
  S_{\rm eff}=S_{\rm QCD}
  +a\int\!\dd^4x\,\mathcal L_1(x)
  +a^2\int\!\dd^4x\,\mathcal L_2(x)+\cdots .
  \label{eq:symanzik}
\end{equation}
\(\mathcal L_n\) 由连续场、连续导数和格点对称性允许的算符组成，量纲为 \(4+n\)。该展开说明离散误差可按 \(a\) 幂次系统分类；它不意味着在有限 \(a\) 上删除了 UV 自由度，也不意味着所有高阶系数都能由树级计算得到。

标准 Wilson 规范作用量从 \(\Ord(a^2)\) 偏差开始。树级 Symanzik 改进在 plaquette 之外加入 \(1\times2\) 矩形圈。在独立平面、同一方向归一化的常用约定中，树级系数为
\begin{equation}
  c_0=\frac53,\qquad c_1=-\frac1{12},
  \qquad c_0+8c_1=1.
  \label{eq:symanzik-coefficients}
\end{equation}
最后一个等式是检验系数与求和约定是否自洽的最简单条件；改变矩形圈的计数方式会改变组合因子，不能只引用 \(-1/12\) 而省略求和定义。

Wilson 费米子的主导误差为 \(\Ord(a)\)。Sheikholeslami--Wohlert 项
\begin{equation}
  S_f=S_W+a^5c_{\rm SW}\sum_x
  \bar\psi(x)\frac{i}{4}\sigma_{\mu\nu}F_{\mu\nu}(x)\psi(x)
  \label{eq:b01-clover-action}
\end{equation}
在适当调谐 \(c_{\rm SW}\) 后消除物理矩阵元中的 \(\Ord(a)\) 项。树级 \(c_{\rm SW}=1\)，单圈形式为 \(1+\Ord(g_0^2)\)，非微扰值可由 Schrödinger 泛函或等价的手征 Ward 条件确定。“使用 Clover”本身不是精度保证：残余 \(\Ord(a\alpha_s)\)、重夸克 \(\Ord(a^2m_Q^2)\)、算符混合和不完整的手征外推仍须逐项估计。

\section{Tadpole 改进与作用量一致性}

格点微扰论中的蝌蚪图使链接的真空期望值明显小于单位矩阵。一个常用非微扰量是
\begin{equation}
  u_0=\left\langle\frac1{N_c}\Re\Tr P_{\mu\nu}\right\rangle^{1/4},
  \qquad U_\mu\longmapsto U_\mu/u_0 .
  \label{eq:tadpole}
\end{equation}
重标链接后，作用量、Clover 系数、临界质量与算符归一化必须同步改变。例如在本节 Clover 归一化下，未重标链接的树级值为 \(c_{\rm SW}=1\)；显式采用 \(U/u_0\) 后，常用树级值为 \(u_0^{-3}\)，临界 hopping 参数也要按相应约定换算。把两种表述混在一起会制造伪装的 \(\Ord(a)\) 误差。

规范作用和费米子作用若都采用 tadpole 改进，应使用同一套自洽定义。不同 \(u_0\) 或只在费米子部分重标，会使规范场和费米子场在不同的有效归一化上“运作”。树级改进之后仍可能存在 \(\Ord(a\alpha_s)\) 交叉项；多格距拟合若只保留 \(a^2\)，该误差可能被错误吸收到物理常数中。

\section{标度设定不是物理极限}

格点计算首先给出无量纲比值。固定一个物理参考长度 \(\ell_{\rm ref}\) 后，
\begin{equation}
  a=\frac{\ell_{\rm ref}}{\ell_{\rm ref}/a}.
  \label{eq:b01-scale-setting}
\end{equation}
\(r_0\)、梯度流尺度 \(w_0\) 或强子质量都可充当 \(\ell_{\rm ref}\)，但它们的统计噪声、夸克质量依赖和外部输入不同。梯度流的常用定义是
\begin{equation}
  \frac{\dd B_\mu(x,t)}{\dd t}
  =-g_0^2\frac{\delta S_W[B]}{\delta B_\mu(x,t)},
  \qquad
  t^2\langle E(t)\rangle\big|_{t=t_0}=0.3,
  \label{eq:flow-scale}
\end{equation}
其中 \(w_0^2=t_0\) 是由同一条件定义的两种记法。流半径 \(r_f=\sqrt{8t}\) 控制平滑尺度；它不是一个可以任意选择后忽略的“数值参数”。

在固定格距族上，改进后的物理量常写成
\begin{equation}
  X(a,m_\pi)=X_0+d_a a^2+d_4a^4
  +d_\pi(m_\pi^2-m_{\pi,\rm phys}^2)+\cdots .
  \label{eq:continuum-fit}
\end{equation}
是否加入 \(a^4\)、交叉项 \(a^2P_z^2\) 或 \(a^2m_Q^2\) 取决于目标和格距覆盖，而不能由数据的 \(\chi^2\) 单独决定。至少两个格距只能检查一个格距趋势；要对多项同时拟合，需要足够多的独立格距和可控的相关误差。

\chapter{费米子谱、色散关系与组态采样}

\section{自由格点色散与加倍子}

最近邻自由标量场的传播子在 Euclidean 动量空间为
\begin{equation}
  G^{-1}(p)=\frac4{a^2}\sum_{\mu=1}^4\sin^2\!\left(\frac{ap_\mu}{2}\right)+m^2.
  \label{eq:scalar-propagator}
\end{equation}
把 \(p_4=iE(\mathbf p)\) 代入 \(G^{-1}=0\)，得到
\begin{equation}
  \cosh[aE(\mathbf p)]
  =\cosh(am)+\sum_{i=1}^3\left[1-\cos(ap_i)\right].
  \label{eq:lattice-dispersion}
\end{equation}
在固定物理动量并令 \(a\to0\) 时，
\begin{equation}
  E^2(\mathbf p)=m^2+\mathbf p^2+\Ord(a^2).
  \label{eq:continuum-dispersion}
\end{equation}
四维超立方对称性通常排除一般 \(\Ord(a)\) 项；若写 \(\Ord(ap)\)，应把它理解为对可用动量范围过宽的保守估计，或明确成 \(\Ord(a^2E^2)\)。在有限格距上，式~\eqref{eq:lattice-dispersion} 可用于拟合高 boost 强子；若数据偏离连续色散，必须先区分激发态、动量涂抹和格点色散，而不能直接归因于新物理。

naive Dirac 算符在动量空间含有 \(\sin(ap_\mu)\)。因此每个动量分量在 \(0\) 与 \(\pi/a\) 附近都可能贡献低能极点，四维共出现 \(2^4=16\) 个候选模。Wilson 项
\begin{equation}
  D_W(p)=m+\frac{i}{a}\sum_\mu\gamma_\mu\sin(ap_\mu)
  +\frac{r}{a}\sum_\mu\left[1-\cos(ap_\mu)\right]
  \label{eq:wilson-dirac}
\end{equation}
给布里渊区边界附近的模增加正比于 \(1/a\) 的有效质量，同时显式破坏标准手征对称性。对空间零动量、正能分支，自由场关系常可化为 \(e^{aE}=1+am\) 的特定形式；它只适用于该分支，不能当作任意动量的精确色散。

Nielsen--Ninomiya 约束表明：平移不变、局域、厄米并保持标准手征反对易关系的格点费米子不能同时避免加倍子。实际方案必须改变其中一个条件。Wilson/Clover 显式破缺手征对称性；Staggered/HISQ 保留 \(U(1)\) 子群但引入 taste；Domain-Wall/Overlap 利用第五维或 Ginsparg--Wilson 非局域性；Twisted Mass 通过最大扭曲实现自动 \(\Ord(a)\) 改进。表~\ref{tab:fermions} 给出选择时应同时检查的性质。

{\small
\begin{longtable}{L{2.4cm}L{1.8cm}L{2.1cm}L{2.4cm}L{2.8cm}}
\caption{主要格点费米子方案的结构性比较。表中“成本”是相对方法级判断，不是具体机器性能测量。}\label{tab:fermions}\\
\toprule
方案 & 手征性 & 主导离散误差 & 主要代价/风险 & 典型用途\\
\midrule
\endfirsthead
\toprule
方案 & 手征性 & 主导离散误差 & 主要代价/风险 & 典型用途\\
\midrule
\endhead
Wilson & 显式破缺 & \(\Ord(a)\) & 加法质量重整化、例外组态 & 基础谱学与大质量夸克\\
Clover & 显式破缺 & \(\Ord(a^2)\)（调谐后） & \(c_{\rm SW}\)、\(a\alpha_s\) 与算符混合 & 强子结构与部分子算符\\
Staggered\slash HISQ & \(U(1)\) 子群 & \(\Ord(a^2)\)，HISQ 进一步压低 & taste 分裂、根技巧与粗格距限制 & 大体积动力学系综\\
Domain-Wall & 有限 \(L_s\) 近似 & \(m_{\rm res}\) 控制后 \(\Ord(a^2)\) & 第五维成本、残余质量 & 手征极限与交叉检验\\
Twisted Mass & 显式破缺 & 最大扭曲下 \(\Ord(a^2)\) & 同位旋破缺、临界质量调谐 & 轻海夸克与谱学\\
Overlap & 精确 GW & \(\Ord(a^2)\) & 矩阵符号函数的成本、拓扑冻结 & 手征对称性基准与验证\\
\bottomrule
\end{longtable}
}

\section{Markov 链与 Metropolis 基础}

格点路径积分的采样目标可写成
\begin{equation}
  P(U)=\frac{e^{-S_g[U]}\det D[U]}{\int\mathcal DU'\,e^{-S_g[U']}\det D[U']}.
  \label{eq:target-weight}
\end{equation}
对 Markov 转移 \(T(U'|U)\)，细致平衡条件
\begin{equation}
  P(U)T(U'|U)=P(U')T(U|U')
  \label{eq:detailed-balance}
\end{equation}
保证 \(P\) 是平稳分布。对称提议的 Metropolis 接受概率为
\begin{equation}
  T_{\rm acc}(U'|U)=\min\!\left[1,e^{-(S[U']-S[U])}\right].
  \label{eq:metropolis}
\end{equation}
单链接扫描只改变与链接相邻的 plaquette；热浴直接从局部条件分布抽样，过松弛则用保作用量的确定性移动加速遍历。实际误差不能只写 \(\sigma/\sqrt N\)：链样本存在自相关，应报告积分自相关时间和有效样本数。

\section{赝费米子与 HMC}

对两味简并费米子，使用正定矩阵 \(D^\dagger D\)：
\begin{equation}
  \int\mathcal D\phi^\dagger\mathcal D\phi\,
  e^{-\phi^\dagger(D^\dagger D)^{-1}\phi}
  =\det(D^\dagger D).
  \label{eq:b01-pseudofermion}
\end{equation}
抽取 \(\chi\sim\exp(-\chi^\dagger\chi)\) 并令 \(\phi=D\chi\)，即可得到协方差为 \(D^\dagger D\) 的高斯场。单味或非偶味情形需要用有理/多项式近似（RHMC/PHMC），并分别控制近似阶数与力误差。

HMC 扩展相空间为 \((U,\Pi)\)。取
\begin{equation}
  H[U,\Pi]=\frac12\Tr\Pi^2+S_g[U]
  +\phi^\dagger(D^\dagger D)^{-1}\phi,
  \label{eq:hmc-hamiltonian}
\end{equation}
通过蛙跳积分得到候选终态。定义 \(F=-\delta S/\delta\omega\) 后，可写为
\begin{equation}
  \Pi_{1/2}=\Pi_0-\frac{\epsilon}{2}F[U_0],\qquad
  U_{k+1}=e^{i\epsilon\Pi_{k+1/2}}U_k,
  \qquad
  \Pi_{k+3/2}=\Pi_{k+1/2}-\epsilon F[U_{k+1}].
  \label{eq:leapfrog}
\end{equation}
末端以 \(\min(1,e^{-\Delta H})\) 接受。精确算术下蛙跳可逆并保持相空间体积，有限步长的能量偏差为 \(\Ord(\epsilon^2)\)，由全局接受步骤修正。

HMC 的生产质量至少需要报告：接受率、轨迹长度、步长、力、Dirac 真残差、拓扑荷历史、不同拓扑扇区的采样和自相关时间。拓扑冻结、近零模和力发散不能由“接受率看起来正常”排除。HMC 的产物是组态集合，不是物理矩阵元；后续两点/三点函数、谱学拟合和重正化仍各有独立误差。

\chapter{非定域算符与重正化}

\section{三类 UV 结构不能混为一谈}

含 Wilson 线的裸矩阵元通常表现为
\begin{equation}
  h_{\rm bare}(z,b_T,L;a)
  \sim Z_{\rm endpoint}(a,\mu)\,
  e^{-\delta m L/a}\,
  e^{-V(b_T)L}\,
  h_{\rm phys}(z,b_T;\mu),
  \label{eq:nonlocal-structure}
\end{equation}
其中第一项是端点对数与算符反常维，第二项是 Wilson 线自能的线性项，第三项是两条近平行长 Wilson 线之间交换胶子的 pinch-pole 结构。若 Wilson 线沿光锥方向延伸到无穷，还会出现 rapidity 发散；rapidity 不是普通 UV 发散，不能用 \(\delta m\) 或局部反常量纲替代。

\begin{longtable}{L{2.3cm}L{3.4cm}L{2.8cm}L{4.0cm}}
\caption{非定域算符中常见发散及所需控制方式。}\label{tab:b01-divergences}\\
\toprule
结构 & 物理来源 & 典型尺度 & 控制方式\\
\midrule
\endfirsthead
\toprule
结构 & 物理来源 & 典型尺度 & 控制方式\\
\midrule
\endhead
线性自能 & 开 Wilson 线的静态色源自能 & \(L/a\) 或 \(|z|/a\) & 质量抵消、同长度 Wilson 圈平方根、多格距自重整\\
端点对数 & 场强/夸克顶点与 Wilson 线连接 & \(\ln(1/a)\) & 反常量纲、RI/MOM 端点因子、短距匹配\\
pinch pole & 两条长 Wilson 线之间的有效势 & \(V(b_T)L\) & 同几何 Wilson 圈、\(L\to\infty\) 平台\\
cusp & 路径转角与有限夹角 & 转角反常维 & 路径匹配的 Wilson 圈减除\\
rapidity & 光锥方向软/共线积分 & \(L_\parallel\) 或 \(\zeta\) & soft function、rapidity regulator、CS 演化\\
\bottomrule
\end{longtable}

\section{局部 RI/MOM 与它的边界}

对局域夸克双线性算符 \(O_\Gamma=\bar u\Gamma d\)，RI/MOM 条件可概括为
\begin{equation}
  Z_\Gamma\frac1{12}\Tr\!\left[
  \langle p|O_\Gamma|p\rangle
  \langle p|O_\Gamma|p\rangle_0^{-1}
  \right]_{p^2=\mu^2}=1,
  \label{eq:ri-mom}
  \end{equation}
其中下标 \(0\) 表示树级或自由场归一化，\(\mu\) 必须落在微扰窗口
\begin{equation}
  \Lambda_{\rm QCD}^2\ll\mu^2\ll a^{-2}.
  \label{eq:perturbative-window}
\end{equation}
RI/MOM 可以确定局部反常量纲和有限混合矩阵，并通过微扰转换到 \(\MSbar\)。但离壳顶点条件不知道开 Wilson 线的总长度与同几何自能；把普通 RI/MOM 直接用于长 Wilson 线或 staple，通常会留下线性项。因此，RI/MOM 失效不表示所有动量减除都无用，而是说明必须补上非定域几何对应的减除。

\section{质量抵消、Wilson 圈与自重整}

辅助场方法把开 Wilson 线看成无动能静态色源的两点函数。若辅助场 \(Z\) 具有运动方程
\begin{equation}
  [\partial_z-ig_0A_z(z)]Z(z)=0,
  \qquad L(z,0)=Z(z)Z^\dagger(0),
  \label{eq:auxiliary-z}
\end{equation}
其线性自能可由质量抵消项吸收。对固定几何和表示，可把重整化结构写成
\begin{equation}
  L_R(z,0)=Z_Z^{-1}e^{-\delta m|z|}L(z,0),
  \label{eq:wilson-renorm}
\end{equation}
其中 \(Z_Z\) 仍包含端点对数。

对 staple 型准 TMD，同几何 Wilson 圈减除写成
\begin{equation}
  h_R(z,b_T,L;a)
  =\lim_{L\to\infty}
  \frac{h_{\rm bare}(z,b_T,L;a)}
  {\sqrt{Z_E(2L+z,b_T;a)}} ,
  \label{eq:wilson-loop-subtraction}
\end{equation}
其中 \(Z_E\) 的矩形圈必须与开 Wilson 线具有匹配的总长度、转角、横向分离和颜色表示。若线性项与 pinch 项都按长度成比例，开方的圈可抵消它们；残余端点对数再通过短距比率（SDR）或混合方案匹配到 \(\MSbar\)。

自重整化利用同一矩阵元的比值
\begin{equation}
  h_R(z,P_z)
  =\frac{h_{\rm bare}(z,P_z)}
  {\left[h_{\rm bare}(z_0,P_z=0)\right]^{z/z_0}},
  \label{eq:self-renorm}
\end{equation}
以短距点 \(z_0\) 标定共同指数因子。当 $z/z_0$ 为整数或通过内插定义时，它可消除与 \(z\) 成正比的共同线性项；但它是否同时消除端点对数、离散项和 rapidity 结构，必须用多格距稳定性与方案匹配证明。单格距下 \(z_0\) 与参数化模型往往不可辨识。

混合方案把短距与长距分开：
\begin{equation}
  h_R(z,P_z)=
  \begin{cases}
    h_{\rm bare}(z,P_z)/h_{\rm bare}(z,0), & z\le z_s,\\[2mm]
    h_{\rm bare}(z,P_z)/
    \left[Z_R(z,a,\mu)h_{\MSbar}(z_s,0,\mu)\right], & z>z_s,
  \end{cases}
  \label{eq:hybrid}
\end{equation}
切换点必须满足 \(a\ll z_s\ll\Lambda_{\rm QCD}^{-1}\)，并检查两侧在 \(z_s\) 处是否连续、最终结果是否依赖 \(z_s\)。若切换点变化超过目标误差，幂修正或微扰匹配尚未受控。

\section{梯度流与重整子}

梯度流将规范场向平滑配置演化。在固定流时下，规范不变的流场关联函数可以具有受控连续极限；其常用尺度窗口为
\begin{equation}
  a\ll r_f=\sqrt{8t}\ll
  \min\left(z,b_T,\Lambda_{\rm QCD}^{-1}\right).
  \label{eq:flow-window}
\end{equation}
左端保证 UV 平滑尺度大于格距，右端保证流没有抹平待测长距结构。流化算符通过短流时展开与普通重整化算符联系：
\begin{equation}
  O(t)=c_O(t,\mu)O_R(\mu)
  +\sum_i c_i(t,\mu)O_{i,R}(\mu)+\cdots ,
  \label{eq:flow-matching}
\end{equation}
因此流时不是普通 UV 标度，也不能替代 rapidity 重正化或算符混合矩阵。

在微扰展开中，Wilson 系数往往是渐近级数。红外重整子使大阶系数按
\begin{equation}
  c_n\sim n!\left(\frac{\beta_0}{2\pi}\right)^n
  \label{eq:renormalon-growth}
\end{equation}
增长，最小项误差对应更高 twist 的幂次模糊。若用主值（PV）或其他 Borel 重求和定义领头重整子，则必须把相应的 \(m_0z\) 线性项、短距 Wilson 系数和匹配阶数放在同一方案中。只对一个局部分量做重求和、却保持其他分量在另一方案，会人为制造方案差。

\chapter{有效场论连接：Symanzik、HQET 与 LaMET}

三种展开的共同骨架是“大标度 + 微扰匹配 + 幂次修正”，但大标度不同：
\begin{center}
\begin{tabularx}{\textwidth}{L{2.4cm}L{3.2cm}L{3.2cm}Y}
\toprule
框架 & 大标度 & 小参数 & 主要失效边界\\
\midrule
Symanzik & \(a^{-1}\) & \(a\Lambda_{\rm QCD}\) & 靠近截断、高阶改进系数和非局域长程算符\\
HQET & \(m_Q\) & \(\Lambda_{\rm QCD}/m_Q\) & 粲夸克、\(am_Q\) 非小、强末态与有限质量修正\\
LaMET & \(P_z\) & \(\Lambda_{\rm QCD}/P_z\) & 端点、大 \(z\)、格点色散和 \(aP_z\) 非小\\
\bottomrule
\end{tabularx}
\end{center}

HQET 从 QCD 场的速度分解开始：
\begin{equation}
  Q(x)=e^{-im_Qv\cdot x}\left[h_v(x)+H_v(x)\right],
  \qquad v^2=1,
  \label{eq:hqet-decomposition}
\end{equation}
其中 \(v\!\!/h_v=h_v\)，\(v\!\!/H_v=-H_v\)。积掉小分量后，
\begin{equation}
  \mathcal L_{\rm HQET}
  =\bar h_viv\cdot D h_v
  +\frac{1}{2m_Q}\bar h_v(iD_\perp)^2h_v
  +\frac{g_0}{4m_Q}\bar h_v\sigma_{\mu\nu}G^{\mu\nu}h_v
  +\Ord(m_Q^{-2}).
  \label{eq:hqet-lagrangian}
\end{equation}
领头项显示重夸克自旋--味对称性；动能和色磁项分别描述残余运动与超精细分裂。HQET 的匹配因子吸收高标度 UV 物理，非微扰矩阵元包含低能强子结构。对格点重味，必须同时报告 \(m_Qa\)、\(\Lambda_{\rm QCD}a\) 和重整化标度；只做形式上的 \(1/m_Q\) 展开而不控制 \(am_Q\)，不能得到可控连续极限。

LaMET 可以写成类似的有效展开：
\begin{equation}
  F(P_z/\Lambda)
  =Z\!\left(\frac{P_z}{\Lambda},\frac{\Lambda}{\mu}\right)f(\mu)
  +\Ord\!\left(\frac{\Lambda^2}{P_z^2},\frac{\Lambda}{P_z}\right).
  \label{eq:lamet-expansion}
\end{equation}
最低阶的类比很有用：\(P_z\leftrightarrow m_Q\)，光锥分布 \(f\leftrightarrow\) HQET 低能矩阵元。但 LaMET 的算符还带 Wilson 线几何、rapidity 和软函数，不能仅凭 HQET 类比推出胶子 TMD 的因子化。对重介子 LCDA，若
\begin{equation}
  P_z\gg m_H\gg\Lambda_{\rm QCD},
  \label{eq:bhqet-hierarchy}
\end{equation}
则可先匹配准分布到 QCD LCDA，再匹配到 boosted HQET LCDA。两个匹配核和两个幂修正参数必须独立记录；把 \(m_H/P_z\) 与 \(\Lambda_{\rm QCD}/m_H\) 合并成一个“小参数”会失去可检验性。

\chapter{综合验收流程}

从格点组态到物理结论，建议按下列顺序设置门控：
\begin{enumerate}
  \item \textbf{定义门：}固定链接方向、生成元、场强符号、Wilson 线表示、Fourier 约定、态归一化和算符端点。
  \item \textbf{代数门：}在恒等链接、纯规范变换和小矩阵自由场中检查协变性、厄米性、Clover 归一化和色迹投影。
  \item \textbf{采样门：}报告 HMC/RHMC 参数、接受率、力、真残差、拓扑荷与自相关；没有有效样本数就不报告统计误差。
  \item \textbf{谱学门：}用至少两个源/汇结构和多个动量检查基态平台、色散关系与激发态污染。
  \item \textbf{重正化门：}区分局部反常量纲、线性 Wilson 线自能、pinch/cusp、rapidity；验证同几何因子、短距窗口和多格距稳定性。
  \item \textbf{连续极限门：}分别扫描 \(a\)、\(P_z\)、\(L\)、流时和匹配标度；任何极限顺序变化都必须成为系统误差项。
  \item \textbf{物理论述门：}只有所有门均通过，才把矩阵元称为目标 PDF、TMD 或重味物理量；否则只能报告诊断量或中间矩阵元。
\end{enumerate}


\clearpage
\chapter{格点QCD的构造与数值基础}
\label{sec:lattice-foundations}

\section{从连续路径积分到有限格点}

格点QCD把连续时空替换成间距为 \(a\)、大小为 \(L^3\times T\) 的四维超立方格点。玻色场定义在格点或链接上，费米场定义在格点上。连续积分测度被有限维积分替代：
\begin{equation}
  \int\mathcal D A\,\mathcal D\bar\psi\,\mathcal D\psi
  \longrightarrow
  \prod_{x,\mu}\int_{\mathrm{SU}(3)}\dd U_\mu(x)
  \prod_{x}\int\dd\bar\psi(x)\dd\psi(x).
  \label{eq:lattice-measure}
\end{equation}
这一替换不是“把导数改成差分”这么简单。要保持规范不变性，必须把物质场之间的相位信息放到有限链接变量中；要保持费米子反对易关系，必须正确离散化Grassmann积分；要得到正确的连续极限，还必须调节裸参数并控制系统误差。

格点上的物理量通常以无量纲数计算。若代码中格距为 \(a\)，则物理能量、动量和长度分别由
\begin{equation}
  E_{\rm phys}=\frac{E_{\rm lat}}{a},
  \qquad
  p_{\rm phys}=\frac{p_{\rm lat}}{a},
  \qquad
  r_{\rm phys}=a r_{\rm lat}
  \label{eq:lattice-units}
\end{equation}
恢复。不同合作组可能报告 \(a^{-1}\) 或物理 \(a\)，也可能用梯度流标度 \(t_0\)、\(w_0\) 或强子质量确定尺度。比较数值前必须先统一单位和约定；把 \(a^{-1}\) 当作固定不变会掩盖尺度设定误差。

\section{链接的几何与Wilson规范作用量}

最基本的规范不变量是plaquette：
\begin{equation}
  P_{\mu\nu}(x)
  =\frac{1}{N_c}\Re\Tr\left[
  U_\mu(x)U_\nu(x+a\hat\mu)
  U_\mu^\dagger(x+a\hat\nu)U_\nu^\dagger(x)
  \right].
  \label{eq:lattice-plaquette}
\end{equation}
其小 \(a\) 展开为
\begin{equation}
  1-P_{\mu\nu}(x)
  =\frac{a^4}{2N_c}\Tr F_{\mu\nu}^2(x)
  +\mathcal O(a^6).
  \label{eq:plaquette-expansion}
\end{equation}
Wilson规范作用量取
\begin{equation}
  S_g[U]
  =\beta\sum_{x,\mu<\nu}
  \left(1-P_{\mu\nu}(x)\right),
  \qquad
  \beta=\frac{2N_c}{g_0^2}.
  \label{eq:wilson-gauge-action}
\end{equation}
它在 \(a\to0\) 时回到 \(\frac14F_{\mu\nu}^aF^{\mu\nu,a}\)。改良作用量通过加入更大环或其他表示的组合消除部分 \(a^2\) 修正，但会更贵，也可能改变拓扑性质。选择作用量时必须与生成系综的代码、尺度设定和重整化方案一起记录。

\section{格点费米子、倍增与Wilson项}

朴素中心差分Dirac算符为
\begin{equation}
  D_{\rm naive}(x,y)
  =m_0\delta_{x,y}
  +\frac{1}{2a}\sum_\mu\gamma_\mu
  \left[U_\mu(x)\delta_{x+\hat\mu,y}
  -U_{-\mu}(x)\delta_{x-\hat\mu,y}\right].
  \label{eq:naive-dirac}
\end{equation}
其自由谱在动量 \(p_\mu\) 处给出
\begin{equation}
  am_{\rm eff}(\mathbf p)
  =\sqrt{\sum_\mu\sin^2(ap_\mu)},
  \label{eq:naive-spectrum}
\end{equation}
因此在布里渊区边界附近的多个动量都可接近零，形成费米子倍增。Wilson项用二阶差分给倍增模增加正比于 \(1/a\) 的质量，代价是显式破坏手征对称性。

Clover作用量把Sheikholeslami--Wohlert项写成
\begin{equation}
  S_W
  =c_{\rm sw}\frac{iag_0}{4}
  \sum_x\bar\psi(x)\sigma_{\mu\nu}F_{\mu\nu}(x)\psi(x),
  \label{eq:clover-action}
\end{equation}
其中 \(F_{\mu\nu}\) 由四个plaquette对称组合。\(c_{\rm sw}\) 可在微扰或非微扰上标定。对TMD计算，Clover项既影响手征极限，也影响三点函数的接触项和重心化。代码中的clover系数、场强符号与传播子归一化不能各自独立“调参”。

\section{HMC采样与拓扑扇区}

系综不是由“随机生成每个链接”得到的，而是从路径积分权重
\begin{equation}
  \mathcal W[U]=\det D[U]\,e^{-S_g[U]}
  \label{eq:lattice-weight}
\end{equation}
采样。伪费米子把行列式写成高斯积分：
\begin{equation}
  \det D[U]
  =\int\mathcal D\phi^\dagger\mathcal D\phi\,
  e^{-\phi^\dagger D[U]^{-1}\phi},
  \label{eq:pseudofermion}
\end{equation}
实际用Hamiltonian Monte Carlo（HMC）引入共轭动量并做分子动力学演化。算法~\ref{alg:hmc} 给出概念化流程；真实代码还包含多时间步、质量预处理、Hasenbusch分解和多种hashing检查。

\begin{algorithm}[htbp]
\caption{HMC生成规范组态的概念流程。}
\label{alg:hmc}
\begin{algorithmic}[1]
\Require 当前规范场 \(U\)，步长 \(\epsilon\)，轨迹长度 \(N_{\rm md}\)，伪费米子 \(\phi\)
\Ensure 新规范场 \(U'\)，接受或拒绝标记
\State 从Gaussian分布抽取共轭动量 \(\Pi\)
\State 计算费米力 \(F_f=-\partial S_f[U]/\partial U\) 和规范力 \(F_g\)
\For{\(k=1,\ldots,N_{\rm md}\)}
  \State \(\Pi\gets\Pi+\frac{\epsilon}{2}(F_g+F_f)\)
  \State \(U\gets\exp(i\epsilon\Pi)U\)
  \State 重新计算 \(F_f\)，并检验Dirac求解收敛
  \State \(\Pi\gets\Pi+\frac{\epsilon}{2}(F_g+F_f)\)
\EndFor
\State 以 \(\min(1,e^{-\Delta H})\) 接受 \((U',\Pi')\)
\State 记录拓扑荷、残差、力、接受率和自相关时间
\end{algorithmic}
\end{algorithm}

拓扑荷在连续极限下量子化为整数，但有限格点上可能出现长拓扑隧道。HMC轨迹过短或步长过大时，采样可能停留在单一扇区，导致体积和连续极限外推被污染。仅检查接受率不低于某个阈值并不能排除拓扑冻结；应同时报告拓扑荷历史、自相关和不同拓扑扇区的采样。

\section{两点、三点与源汇结构}

基态谱学从两点函数开始。对插入动量 \(\mathbf p\) 的源 \(O_s\) 和汇 \(O_\Gamma\)，定义
\begin{equation}
  C_2(t;\mathbf p)
  =\sum_{\mathbf x}e^{-i\mathbf p\cdot\mathbf x}
  \left\langle O_\Gamma(\mathbf x,t)O_s^\dagger(\mathbf 0,0)\right\rangle .
  \label{eq:lattice-two-point}
\end{equation}
三点函数加入流或胶子算符 \(J\)：
\begin{equation}
  C_3(t,\tau;\mathbf p_f,\mathbf p_i)
  =\sum_{\mathbf x_f,\mathbf x_J}
  e^{-i\mathbf p_f\cdot\mathbf x_f}
  e^{i(\mathbf p_f-\mathbf p_i)\cdot\mathbf x_J}
  \left\langle O_f(\mathbf x_f,t)J(\mathbf x_J,\tau)
  O_i^\dagger(\mathbf 0,0)\right\rangle .
  \label{eq:lattice-three-point}
\end{equation}
源、汇与插入算符的时间顺序必须在 \(0<\tau<t\)，插入在边界时会出现接触项和不完整传播子。对于TMD，\(J\) 还携带横向分离 \(\bT\)、Wilson线路径和场强指标；三点函数不再只是一个标量，而是随 \(\bT\) 和投影变化的矩阵元族。

\section{传播子、涂抹与基态重叠}

点源传播子为
\begin{equation}
  S(x;y)=D^{-1}(x,y),
  \label{eq:propagator-def}
\end{equation}
但在大体积中直接点源会带来巨大统计噪声。常用源包括：
\begin{itemize}
  \item 点源：便宜、全动量信息，但高激发态重叠大；
  \item 高斯/动量涂抹源：增加基态重叠，适合低动量谱；
  \item 墙源与随机源：降低方差，代价是更多传播子或更复杂归一化；
  \item distillation源：把多个源压缩到低维基空间，适合三点收缩；
  \item 多源、全模式平均：抑制源位置和低模涨落。
\end{itemize}
源改变的是态投影与统计，不改变物理算符定义。必须对所有源使用同一插入算符、Wilson线和归一化，否则不同源之间的差异可能来自方案而非物理。

\section{尺度设定、连续极限与体积}

格点计算给出无量纲数，必须通过实验或理论输入定标。常用标度包括 \(t_0\)、\(w_0\)、\(\Omega\) 重子质量和 pion 衰变常数。尺度设定的方程可抽象为
\begin{equation}
  a^{-1}=F_{\rm ref}\big/\sqrt{t_0/a^2},
  \label{eq:scale-setting}
\end{equation}
但不同标度的 \(a\) 误差不完全相关。报告物理量时应说明先固定哪个标度、是否做联合拟合以及系统误差如何传播。TMD矩阵元本身依赖 \(P_z\) 和 \(\bT\)，把整个动量空间用单一 \(a^{-1}\) 缩放后，还要检查动量分辨率与流时分辨率。

连续极限通常写成
\begin{equation}
  O(a)=O(0)+c_1a+c_2a^2+\cdots,
  \label{eq:continuum-expansion}
\end{equation}
Wilson/Clover作用量的 \(\mathcal O(a)\) 项在改进系数足够好时可以减小，但不能仅凭“使用Clover”就写成 \(a^2\) 展开。手征极限、体积极限和大动量极限也需要联合处理。一个安全做法是同时对至少两个方向做外推，并检查路径依赖；只对唯一最细格距做线性外推，无法识别交叉项。

\section{梯度流的作用与边界}

梯度流把规范场沿流时 \(t_f\) 演化：
\begin{equation}
  \frac{\dd V_\mu(x,t_f)}{\dd t_f}
  =-g_0^2\frac{\delta S_g}{\delta V_\mu(x,t_f)}V_\mu(x,t_f),
  \qquad
  V_\mu(x,0)=U_\mu(x).
  \label{eq:flow-equation-intro}
\end{equation}
流动使高动量模被平滑，短流时下可通过微扰匹配把流场与\(\MSbar\)算符关联。流时并非越小越好：太小保留大UV噪声，太大会引入 \(t_f/a^2\) 离散误差和幂修正。定义无量纲流时 \(8t_f/a^2\)，并在固定 \(t_f\) 与固定物理尺度之间做取舍。

\section{系综、统计与计算成本}

一个格点系综由裸参数、作用量、体积、格距和生成轨迹共同定义。有效样本数不是“组态总数”，而是考虑自相关后的数量：
\begin{equation}
  N_{\rm eff}=\frac{N}{\tau_{\rm int}},
  \label{eq:lattice-neff}
\end{equation}
其中 \(\tau_{\rm int}\) 可用自相关函数估计。对于三点函数和大 \(z\) 矩阵元，自相关可能因拓扑和低模而更长。报告组态数、轨迹间隔和自相关估计，比只报告“用了多少GPU小时”更能支持可复现性。

\begin{checkpointbox}{格点基础自检}
\begin{enumerate}
  \item 作用量、链接方向、Clover系数和场强定义是否一致；
  \item 质量、尺度、体积和拓扑荷是否由同一系综报告；
  \item 源类型是否与传播子、涂抹和归一化一致；
  \item 连续极限是否包含至少两个格距，而非只改一个参数；
  \item 大 \(z\) 与周期边界是否分开检查；
  \item 流时间是否同时满足离散误差与物理幂修正窗口。
\end{enumerate}
\end{checkpointbox}

\clearpage
\chapter{Euclidean 格点表述}

\section{为什么不能直接在格点上模拟光锥}

格点 QCD 计算的是带正权重的 Euclidean 路径积分。光锥方向 \(n^2=0\) 在 Wick 转动后不具有实的 Euclidean 对应，实时 Wilson 线也不能直接作为有限格距算符计算。直接的困难包括：
\begin{itemize}
  \item 光锥分离和实时演化要求解析延拓，数值上不可控；
  \item 光样 Wilson 线产生 rapidity 发散和 cusp 发散；
  \item 物理 TMD 依赖初态/末态规范链方向，不能用一个无方向的局域算符替代；
  \item TMD 的 \(k_T\) 分布需要大体积和稳定的横向 Fourier 变换；
  \item 胶子算符是单态算符，统计噪声、真空减除和算符混合都比非单态夸克算符更复杂。
\end{itemize}

\section{准 TMD 算符与 staple 几何}

LaMET 的核心思想是：在强子具有大空间动量 \(P_z\) 时，等时、空间分离的 Euclidean 关联函数可以在大动量极限下匹配到光锥关联函数\cite{Ji2013,Ji2014,JiReview2021}。对胶子 TMD，定义准关联函数
\begin{equation}
  \Oq(z,\bT;P_z,L_\parallel)
  =
  F^{zi}_a(z\hat z+\bT)\,
  U^{ab}_{\Gamma_\pm}(z\hat z+\bT;0)\,
  F^{zj}_b(0),
  \label{eq:quasioperator}
\end{equation}
并将横向指标按光锥投影的 leading-twist 结构缩并。坐标空间矩阵元写成
\begin{equation}
  h_g(z,\bT;P_z,L_\parallel;1/a)
  =\langle P_z|\Oq(z,\bT;P_z,L_\parallel)|P_z\rangle .
  \label{eq:quasimat}
\end{equation}
在有限格距上，准 TMD 还含有 Wilson 线自能、cusp、pinch-pole 和端点对数发散。若直接对 \(z\) 做 Fourier 变换而不先减除这些结构，所得函数不能被解释为物理分布。

物理的 \(b_T\) 依赖通过
\begin{equation}
  \widetilde f^{g}(x,\bT;P_z;\mu_R)
  =\int\frac{\dd z}{2\pi}\,e^{ixP_z z}
  \frac{h_g(z,\bT;P_z,L_\parallel;1/a)}
  {\mathcal N(x,P_z,a,L_\parallel)}
  \label{eq:quasifourier}
\end{equation}
定义；\(\mathcal N\) 是归一化和重正化因子，依赖于所用方案。这里故意没有把 \(\mathcal N\) 写成一个普遍的指数因子：胶子 TMD 的 Wilson 线自能、cusp 角度、场强端点和 pinch-pole 结构不能用一个标量函数在没有假设的情况下全部吸收。

\section{准 TMD、软函数与 CS 核的关系}

在 LaMET 因子化成立的区域，重正化后的准 TMD 与光锥 TMD 的关系可概括为
\begin{equation}
  \frac{\widetilde f^{g}(x,\bT;P_z;\mu)}
  {\sqrt{S_{I,g}(\bT;\mu)}}
  =
  \left[H_g\otimes f_g\right](x,\bT;\mu;\zeta_z)
  \exp\left[\frac{1}{2}K_g(\bT;\mu)
  \ln\frac{\zeta_z}{\zeta}\right]
  +\text{高扭度修正},
  \label{eq:lametgluon}
\end{equation}
其中 \(\zeta_z\sim(2xP_z)^2\) 是准分布的自然快度标度，\(\otimes\) 概括硬匹配中可能出现的 \(x\) 卷积以及胶子-夸克单态混合。对非单态或近似忽略混合的计算，可把硬核取为对角形式；但在核子单态胶子通道中，不能把“忽略 \(g-q\) 混合”当作普适定理，只能作为有误差估计的近似\cite{ZhangGluonLaMET2019,JiReview2021}。

把抽象卷积写开，可以把“胶子来自胶子”和“胶子来自单态夸克”分成两个结构项：
\begin{keyformula}{大动量匹配的卷积展开}
\begin{align}
  \frac{\widetilde f^{g}(x,\bT;P_z;\mu)}
  {\sqrt{S_{I,g}(\bT;\mu)}\,e^{\frac12K_g\ln(\zeta_z/\zeta)}}
  &=\int_x^1\frac{\dd y}{y}\,
  H_{gg}\!\left(\frac{x}{y},\frac{\mu}{P_z}\right)
  f_g(y,\bT;\mu,\zeta)
  \nonumber\\
  &\quad+\int_x^1\frac{\dd y}{y}\,
  H_{gq}\!\left(\frac{x}{y},\frac{\mu}{P_z}\right)
  q_S(y,\bT;\mu,\zeta)
  +\mathcal O(\alpha_s^2,\text{高扭度}).
  \label{eq:lametconvolution}
\end{align}
\end{keyformula}

在只保留非单态或类对角近似时，可令 \(H_{gq}\to0\)，并把 \(H_{gg}\) 当作对角核；一旦保留单态演化，上式就是一个矩阵卷积，必须与非微扰重正化矩阵分开记录。硬核变量 \(x/y\) 还提醒我们：格点端点 \(x\to0,1\) 的误差会通过卷积扩散到整个 \(x\) 区间，不能只在单个动量点上检查。

高扭度修正至少包括
\begin{equation}
  \mathcal O\!\left(\frac{M_N^2}{P_z^2},
  \frac{\Lambda_{\rm QCD}^2}{P_z^2},
  \frac{1}{b_T^2P_z^2},
  \frac{\Lambda_{\rm QCD}}{P_z}\right),
  \label{eq:powercorrections}
\end{equation}
具体幂次取决于算符、投影和匹配方案。大 \(P_z\) 因此不是可选的数值优化，而是因子化定义的一部分；没有大动量窗口时，即使矩阵元本身有限，也不能把结果解释为光锥 TMD。

\section{准 PDF、赝 PDF 与 Ioffe-time 表述}

胶子 TMD 不是孤立的格点对象。理解它与准 PDF、赝 PDF 和 reduced Ioffe-time distribution（RITD）的关系，有助于把短距 OPE、Fourier 变换和 TMD 的 \(b_T\) 匹配分开。对共线胶子 PDF，准 PDF 定义为
\begin{equation}
  \tilde f_g(x,P_z)
  =\int\frac{\dd z}{2\pi xP_z}
  e^{ixP_z z}
  \langle P_z|\mathcal O_g(z,0)|P_z\rangle ,
  \label{eq:quasipdf}
\end{equation}
其中 \(\mathcal O_g(z,0)=F^{zi}_a(z)U_{ab}(z,0)F^{zj}_b(0)\) 的横向指标按 twist-2 投影缩并。这里的 \(1/(xP_z)\) 和 Fourier 相位必须与矩阵元归一化一起定义；等价地，可以使用动量密度 \(x\tilde f_g(x,P_z)\)，以避免端点处的额外歧义。

令 Ioffe 时间
\begin{equation}
  \nu=P_z z,
  \label{eq:Ioffe}
\end{equation}
并记坐标空间矩阵元为 \(M_g(\nu,z^2)\)。在短距离 \(z^2\ll\Lambda_{\rm QCD}^{-2}\) 区域，算符乘积展开给出
\begin{equation}
  M_g(\nu,z^2)
  =\sum_{n=0}^{\infty}\frac{(-i\nu)^n}{n!}
  C_{g,n}(z^2\mu^2)\,
  \langle x^n\rangle_g(\mu)
  +\mathcal O(z^2\Lambda_{\rm QCD}^2),
  \label{eq:OPEgluon}
\end{equation}
其中 \(C_{g,n}\) 是微扰 Wilson 系数，\(\langle x^n\rangle_g\) 是胶子 PDF 的 Mellin 矩。因子 \((-i\nu)^n/n!\) 来自 \(z\) 的 Taylor 展开，而 Wilson 系数包含 UV 重整化和 DGLAP 型对数。该式说明：准 PDF 不是直接等于 PDF，而是把光锥矩重新求和为空间关联；只有在 \(P_z\to\infty\)、固定 \(x\) 以及适当重正化后，才恢复光锥分布\cite{Ji2013,Ji2014,Izubuchi2018,JiReview2021}。

赝 PDF 方法从同一 Euclidean 矩阵元出发，先构造 reduced Ioffe-time distribution。一个常用定义是
\begin{equation}
  \mathscr M_g(\nu,z^2)
  =\frac{M_g(\nu,z^2)/M_g(0,0)}
  {M_g(0,z^2)/M_g(0,0)}
  =\frac{M_g(\nu,z^2)}{M_g(0,z^2)},
  \label{eq:RITDdef}
\end{equation}
等价地可以写成 \(M_g(\nu,z^2)/M_g(0,z^2)\) 的归一化形式。分子保留 boost 依赖，分母用零动量矩阵元吸收 Wilson 线的 UV 结构，因此 RITD 在固定 \(z^2\) 下可以在不同 \(P_z\) 之间比较。胶子 RITD 与光锥 PDF 的匹配写成
\begin{align}
  \mathscr M_g(\nu,z^2)
  &=\int_0^1\dd y\,\frac{y g(y,\mu)}{\langle x\rangle_g}
  R_{gg}(y\nu,z^2\mu^2)\nonumber\\
  &\quad+\frac{P_z}{P_0}\int_0^1\dd y\,
  \frac{y q_S(y,\mu)}{\langle x\rangle_g}
  R_{gq}(y\nu,z^2\mu^2)+\cdots,
  \label{eq:RITDmatching}
\end{align}
其中 \(q_S\) 概括单态夸克组合，\(P_0\) 是参考动量，\(R_{gg},R_{gq}\) 包含演化和方案转换。若忽略 \(R_{gq}\)，必须把误差单独列出；这不是“胶子只与胶子匹配”的严格陈述。

准 PDF 与赝 PDF 的主要区别在于幂修正的组织方式。准 PDF 的匹配以 \(P_z\) 为大标度，典型修正为 \(M_N^2/P_z^2\) 和 \(\Lambda_{\rm QCD}^2/P_z^2\)；赝 PDF 的匹配以 \(z^2\) 为短距标度，典型修正为 \(z^2\Lambda_{\rm QCD}^2\)。二者在短距 OPE 的公共区域可以相互转换，但不能把一种方法的端点外推直接套到另一种方法上\cite{Izubuchi2018,Cichy2019}。

\section{从共线准 PDF 到 quasi-TMD 的逻辑差异}

共线准 PDF 的 Wilson 线是一条直线；quasi-TMD 则需要 staple，因为 TMD 的横向分离与过程相关的软结构同时存在。因此两者的重正化问题并不相同：
\begin{itemize}
  \item 准 PDF 的直线 Wilson 线主要产生线性自能和端点对数，局部 Wilson 圈或质量减除常常足以组织重正化；
  \item quasi-TMD 的 staple 还包含两条平行长 Wilson 线之间的 pinch-pole、cusp 角度和横向分离，必须使用同几何 Wilson 圈、软函数和 CS 演化；
  \item 准 PDF 的匹配只在 \(x\) 与硬核之间进行，而 TMD 匹配还必须保留 \(\bT\)、\(\zeta\) 和过程相关的规范链。
\end{itemize}
因此，“胶子准 PDF 已经算到连续极限”不能推出“胶子 TMD 已完成”。前者是后者的重要构件，但 TMD 额外要求 \(b_T\) 分辨、rapidity 减除和表示相关软结构。

\section{准 TMD 波函数与软函数提取}

介子准 TMD 波函数提供了一条重要的辅助路线。对赝标量介子，定义带 staple 的准波函数，并把它与介子形状因子、软函数和硬核联系。该路线的优点是两点函数计算成本较低，且软函数和 CS 核可以通过不同强子态或不同动量的比值提取\cite{ZhangSoft2020,ChuCS2022}。它不能直接替代核子的胶子 TMD，但可以验证以下普适结构：
\begin{enumerate}
  \item 规范链自能是否按同几何的 Wilson 圈平方根抵消；
  \item 内禀软函数是否在多个 \(P_z\) 下稳定；
  \item CS 核是否与强子态无关；
  \item 匹配核是否在误差范围内消除硬贡献的动量依赖。
\end{enumerate}

\section{软函数提取的显式因子化链}

在大动量赝标量介子形状因子方法中，定义四夸克形状因子
\begin{equation}
  F(b_T,P_z)
  =\langle \pi(-P_z)|
  (\bar q_1\Gamma q_1)(\bT)\,
  (\bar q_2\Gamma q_2)(0)
  |\pi(P_z)\rangle_c ,
  \label{eq:mesonformfactor}
\end{equation}
其中不同味道和连通收缩用于避免不需要的 disconnected 拓扑。leading-power 因子化为
\begin{equation}
  \begin{aligned}
  F(b_T,P_z)
  &=S_{I,g}(b_T;\mu)
  \int_0^1\dd x\,\dd x'\,
  H_g(x,x',P_z;\mu)\\
  &\quad\times
  \Phi_g^\dagger(x',b_T,-P_z;\mu)\,
  \Phi_g(x,b_T,P_z;\mu)
  +\mathcal O(P_z^{-2}),
  \end{aligned}
  \label{eq:mesonfactorization}
\end{equation}
其中 \(\Phi_g\) 是 quasi-TMDWF，\(H_g\) 是硬核。若在领头阶把 \(H_g\) 近似为常数，并利用赝标量介子的宇称关系 \(\phi_g(0,b_T,-P_z)=\phi_g(0,b_T,P_z)\)，则
\begin{equation}
  S_{I,g}(b_T;\mu)
  =\frac{2N_c\,F(b_T,P_z)}
  {|\phi_g(0,b_T,P_z)|^2}
  +\mathcal O(\alpha_s,P_z^{-2}),
  \label{eq:SIextraction}
\end{equation}
其中 \(\phi_g(0,b_T,P_z)=\int_0^1\dd x\,\Phi_g(x,b_T,P_z)\)。在单圈精度下，必须把硬核 \(H_g\)、波函数归一化和 Fierz 重排一起纳入；否则 \(S_I\) 的 \(P_z\) 平台可能只是硬核误差被误当作软物理\cite{ZhangSoft2020,ChuSoft2023}。

CS 核可由不同动量下的准 TMDWF 比值提取：
\begin{equation}
  K_g(b_T;\mu)
  =\frac{1}{\ln(P_1^z/P_2^z)}
  \ln\left|
  \frac{H_g(xP_2^z;\mu)\,
  \widetilde\Phi_g(x,b_T,P_1^z;\mu)}
  {H_g(xP_1^z;\mu)\,
  \widetilde\Phi_g(x,b_T,P_2^z;\mu)}
  \right|
  +\mathcal O(P_z^{-2}).
  \label{eq:CScore}
\end{equation}
若使用 \(\Phi_g\) 的 \(x\) 积分或特定 \(x\) 点，必须验证所得 \(K_g\) 对 \(x\)、\(P_1^z/P_2^z\) 和硬核阶数不敏感。CS 核的普适性不是“任取两个动量相除”即可自动成立，而是因子化和幂修正控制后的结果。

\section{短距匹配与软因子不可交换}

把准 TMD 写成
\begin{equation}
  \widetilde f_g
  =H_g\otimes f_g\,
  \exp\left[\frac{1}{2}K_g\ln\frac{\zeta_z}{\zeta}\right]
  \frac{1}{\sqrt{S_{I,g}}}
  +\text{高扭度}
  \label{eq:matchingrepeat}
\end{equation}
时，四种常见错误是：
\begin{enumerate}
  \item 用裸准 TMD 直接除一个没有同几何的软函数；
  \item 用夸克表示或夸克通道拟合的 \(K\) 替代 \(K_g\)；
  \item 在匹配后乘一个标量 \(Z\) 吸收所有混合，却不报告 \(g-q\) 矩阵；
  \item 在 \(P_z\) 尚未进入因子化窗口时，把 \(1/P_z^2\) 修正吸收到软函数中。
\end{enumerate}
这些操作的共同问题是把不同尺度、不同表示和不同过程的贡献混在一个因子中。正确的顺序应当是：固定算符和几何，完成 UV 重正化，构造软/rapidity 减除，再匹配硬核，最后做大动量外推。每一步都应有独立检查和可回退的方案。

\clearpage
\chapter{因子化、软共线与匹配的深入展开}
\label{sec:factorization-details}

\section{因子化的物理图像}

因子化不是把一个复杂函数任意拆成几个因子的代数操作，而是利用尺度分离把“可微扰的短距碰撞”“沿运动方向的共线辐射”和“低能软辐射”分配到不同算符中。只要这些区域在幂次展开中互不重叠，它们的贡献就可以相乘或卷积；重叠区域造成的紫外与rapidity依赖则由重正化核吸收。TMD因子化的物理图像可以概括为一句话：硬过程选定部分子和标度，共线结构决定纵向动量分布，横向分离与软辐射共同生成TMD核。

考虑一个硬过程产生具有横向动量 \(\mathbf q_T\) 的末态。若存在层级
\begin{equation}
  \Lambda_{\rm QCD}\ll q_T\ll Q,
  \label{eq:fact-hierarchy}
\end{equation}
则可以先把 \(Q\) 对应的短距系数算出来，再把 \(q_T\) 对应的横向结构和软辐射保留为非微扰函数。TMD因子化把这个层级进一步组织成 \(\mu\)、\(\zeta\) 和 \(\bT\) 三个变量。若 \(q_T\sim Q\)，TMD与固定阶微扰展开重叠；若 \(q_T\sim\Lambda_{\rm QCD}\)，横向动力学进入非微扰区，必须使用TMD分布而不是简单的共线PDF。

在格点准TMD中，\(P_z\) 扮演制造大标度的角色。它同时提供接近光锥的 boost、控制幂修正，并把不同 \(x\) 分量的动量尺度推到较高的微扰区域。但 \(P_z\) 不是硬过程的 \(Q\)，也不能直接替代 \(\mu\) 或 \(\zeta\)。因此格点匹配必须保留三层标度的独立演化，不能把整个因子化链压成一个“除以软函数”的操作。

\begin{figure}[htbp]
\centering
\begin{tikzpicture}[
  region/.style={draw=qcdnavy!65!black,rounded corners=3pt,align=center,minimum width=3.3cm,minimum height=1.25cm,fill=qcdmist},
  soft/.style={draw=qcdrust!70!black,rounded corners=3pt,align=center,minimum width=3.3cm,minimum height=1.25cm,fill=qcdsand},
  arrow/.style={-{Latex[length=2.2mm]},thick,qcdnavy!75!black}
]
\node[region] (hard) at (0,0) {硬区域\\\(k\sim Q\)};
\node[region] (coll) at (4.2,0) {共线区域\\\(k^+\sim Q,\;k_T\sim q_T\)};
\node[soft] (soft) at (8.4,0) {软区域\\\(k\sim q_T\)};
\draw[arrow] (hard) -- node[above,font=\small]{短距匹配} (coll);
\draw[arrow] (coll) -- node[above,font=\small]{rapidity/soft} (soft);
\node[align=left,font=\small,below=0.65cm of coll] {因子化要求各区域贡献可分离，\\并由对应核吸收重叠发散};
\node[align=center,font=\small,below=1.65cm of hard] {微扰控制区域};
\draw[thick,qcdrust!70!black]
  ([xshift=0.2cm,yshift=0.35cm]hard.south west) -- ([xshift=0.2cm,yshift=0.35cm]coll.south east);
\end{tikzpicture}
\caption{TMD因子化的硬、共线与软区域示意。区域边界不是固定动量截断，而由幂计数与重正化方案定义。}
\label{fig:factorization-regions}
\end{figure}

\section{共线算符与动量分数}

对一个沿负 \(z\) 方向入射的强子，可以把共线动量写成
\begin{equation}
  k^\mu=xP^\mu+k_T^\mu+k_{\rm rem}^\mu,
  \qquad
  k^+\sim Q,
  \qquad
  k^-\sim\frac{\mathbf k_T^2}{Q},
  \label{eq:collinear-scaling}
\end{equation}
其中 \(x\) 是部分子动量分数，\(\mathbf k_T\) 是横向动量。共线算符提取 \(x\) 与横向结构，软算符只保留低能规范场。对胶子，横向指标必须投影到物理极化；对夸克，需要Dirac投影。这个分解本身不指定Wilson线方向，方向由硬过程决定。

共线分布的动量守恒体现在卷积上：
\begin{equation}
  f_{i/h}(x)=\int_x^1\frac{\dd y}{y}\,
  P_{ij}(x/y)f_{j/h}(y),
  \label{eq:collinear-convolution}
\end{equation}
其中 \(P_{ij}\) 是劈裂核。TMD版本把横向动量也保留下来，因此卷积从一维变成二维横向卷积；胶子TMD的劈裂核包含 \(g\to gg\) 与 \(g\to q\bar q\)，这解释了为什么单态匹配自然带有 \(2\times2\) 或更大矩阵结构。

\section{Wilson线方向与过程依赖}

TMD算符中的Wilson线不是装饰，而是把从部分子到硬过程的色流并行移动。若硬过程需要向未来方向的规范链，则使用
\begin{equation}
  W_{\Gamma}^{[R]}(x,y)
  =\mathcal P\exp\left[ig\int_\Gamma\dd z^\mu A_\mu^{[R]}(z)\right],
  \label{eq:wilson-general}
\end{equation}
未来指向与过去指向的交换涉及伴随表示的转置以及软函数的复共轭。对于胶子算符，Wilson线取伴随表示，横向端点还必须用场强张量连接；对于夸克TMD，Wilson线取基本表示。两种表示不能通过一个常数因子互换，因为它们的色Casimir和软反常维度不同。

在格点中，Wilson线的方向被离散成staple。一个长度为 \(L_\parallel\)、横向分离为 \(\bT\) 的staple可以写作
\begin{equation}
  W_{\rm staple}(z,\bT)
  =U_z(z)\,W_{\perp}(z+L_\parallel\hat z,\bT)\,
  U_z^\dagger(z+\bT),
  \label{eq:staple-order}
\end{equation}
其起点、终点、横向路径顺序和反向链接必须与算符端点一致。若代码对正负 \(z\) 使用不同的链接方向，或者把两端场强的伴随指标顺序颠倒，通常不会立即产生编译错误，却会在 \(P_z\to-P_z\) 时给出不满足厄米性的关联函数。因此检查 \(h(-z)=\pm h^*(z)\) 是必要的数值不变量。

\section{软函数与rapidity发散}

软函数来自低能胶子交换。一个典型的裸软函数可写成
\begin{equation}
  S_R(\bT)=\int\frac{\dd^2\mathbf k_T}{(2\pi)^2}
  \left\langle 0\left|
  W_{\Gamma_1}^{[R]\dagger}(\bT)
  W_{\Gamma_2}^{[R]}(0)
  \right|0\right\rangle ,
  \label{eq:bare-soft}
\end{equation}
其中两条Wilson线的方向由过程决定。由于Wilson线在光锥方向延伸，积分在 \(k^+\) 或 \(k^-\) 端可能出现rapidity发散：
\begin{equation}
  S_{\rm bare}
  =\left(\frac{\nu}{\delta}\right)^{\gamma_S}
  S_R(\nu),
  \label{eq:rapidity-divergence}
\end{equation}
\(\delta\) 是rapidity regulator，\(\nu\) 是相应的重正化标度。这里的关键不是把 \(\delta\to0\) 直接取掉，而是把 \(\nu\) 演化与硬核、TMD核的 \(\zeta\) 依赖组合成方案无关的物理量。

常用rapidity regulator包括：
\begin{itemize}
  \item 解析或指数型因子 \(e^{-\delta|k_\pm|}\)，便于微扰计算，但可能破坏四维协变性；
  \item 有限长度Wilson线 \(L_\parallel\)，便于格点实现，但会引入pinch-pole与 \(L_\parallel\to\infty\) 外推；
  \item 软函数比值或Collins构造，直接用几何抵消一部分发散；
  \item 横向动量或Wilson圈减除，保留格点可实现的短距参考点。
\end{itemize}
不同regulator给出的中间软函数可以差一个有限因子，但CS核与物理TMD必须一致。把某个regulator下的 \(\gamma_S\) 直接与另一个方案的匹配核相加，是方案错误的典型来源。

\section{从重正化不变性推导CS方程}

设光锥TMD在方案 \(S\) 下写成 \(f^{R}(x,\bT;\mu,\zeta)\)。物理截面可写成
\begin{equation}
  \sigma=\int \dd x\,\dd^2\mathbf q_T\,
  H_S(Q,\mu,\zeta)\,f^{R}_S(x,\bT;\mu,\zeta)
  +\text{高扭度}.
  \label{eq:factorization-physical}
\end{equation}
由于左边与 \(\mu,\zeta\) 无关，硬核的演化必须与TMD演化相消。若把硬核写成 \(H=H_0\exp[-\frac12K(\bT;\mu)\ln(\zeta/\zeta_0)]\)，则得到
\begin{equation}
  2\zeta\frac{\partial}{\partial\zeta}
  \ln f^{R}(x,\bT;\mu,\zeta)
  =K_R(\bT;\mu).
  \label{eq:cs-equation}
\end{equation}
进一步要求 \(K_R\) 与 \(\mu\) 的组合满足
\begin{equation}
  \mu\frac{\dd}{\dd\mu}K_R(\bT;\mu)
  =-\gamma_K(\alpha_s),
  \label{eq:cs-rge}
\end{equation}
因此CS方程可以积分成
\begin{equation}
  f^{R}(x,\bT;\mu,\zeta)
  =f^{R}(x,\bT;\mu,\zeta_0)
  \exp\left[
  \frac12K_R(\bT;\mu)
  \ln\frac{\zeta}{\zeta_0}
  \right].
  \label{eq:factor-cs-solution}
\end{equation}
若 \(K_R\) 包含 \(\mu\) 依赖，则应在积分中使用相应RGE解。实际拟合中常把 \(K_R\) 分解成微扰短距部分和非微扰长距部分；分解点 \(b_T^*\) 的选择必须与该标度下的匹配系数一起报告。

\section{软核的普适性与表示标签}

CS核应与强子态和具体硬过程无关，但可以和颜色表示、Wilson线几何以及rapidity方案有关。对夸克和胶子，分别记作 \(K_q\) 与 \(K_g\)。在微扰短距区域，它们满足由Casimir控制的单圈关系，但不能把这个关系外推到所有 \(b_T\)：
\begin{equation}
  K_i(\bT;\mu)
  =\Gamma_{\rm cusp}^i\,L(\bT;\mu)+K_i^{(1)}(\mu)
  +\mathcal O(\alpha_s^2),
  \qquad
  L=\ln\frac{b_T^2\mu^2e^{2\gamma_E}}4.
  \label{eq:cs-perturbative}
\end{equation}
其中 \(\Gamma_{\rm cusp}^i\) 与 \(C_i\) 成正比。若把胶子核简单写成 \(K_g=(C_A/C_F)K_q\)，只在单圈短距极限可成立；非微扰核与软函数的高阶结构必须独立验证。

\section{硬匹配的算符基底}

匹配不是把一个函数“乘一个系数”，而要先选择低能算符基底。对胶子TMD的单态部分，最低阶基底至少包括
\begin{equation}
  \mathcal B=\left\{f_g,f_{q,S},f_{\bar q,S}\right\},
  \label{eq:tmd-basis}
\end{equation}
其中 \(f_{q,S}\) 是夸克与反夸克单态组合。若只保留胶子与一种夸克组合，可以压缩为 \(2\times2\) 矩阵；若把不同味道独立保留，矩阵维数更大。Wilson系数满足
\begin{equation}
  \widetilde f_i
  =C_{ij}\otimes f_j
  \exp\left[\frac12K_j\ln\frac{\zeta_z}{\zeta}\right],
  \label{eq:matching-basis}
\end{equation}
这里的 \(C_{ij}\) 与UV重正化矩阵 \(Z_{ij}\) 不是同一对象。前者把准算符矩阵元映射到光锥基底，后者把一个UV标度下的算符矩阵元映射到另一个UV标度。二者都可能呈矩阵形式，但作用的空间和标度不同。

在核子胶子通道中，胶子准算符可以没有显式的夸克端点，仍然通过微扰匹配产生 \(g\leftrightarrow q\) 的味混合。判断是否需要 \(gq\) 核不能只看算符是否含夸克场，而要看低能基底是否允许混合。若忽略混合，应给出基于味道对称性或可辨识性的明确假设。

\section{\texorpdfstring{小 \(b_T\) OPE与长距外推}{小 bT OPE 与长距外推}}

在 \(b_T\ll1/\Lambda_{\rm QCD}\) 区域，TMD可以写成局部矩阵元的OPE：
\begin{equation}
  f_{g/h}(x,\bT;\mu,\zeta)
  =\sum_{n}\sum_{i}
  C_{g i,n}(b_T;\mu,\zeta)\otimes
  f_{i/h,n}(x;\mu)+\mathcal O(b_T^2\Lambda_{\rm QCD}^2),
  \label{eq:small-b-ope}
\end{equation}
其中 \(n\) 标记 twist 或自旋结构。小 \(b_T\) 是微扰区，但 \(b_T\to0\) 也会产生大对数，需要把匹配标度选择为 \(\mu_b\sim1/b_T\) 并进行RGE重求和。把拟合区间扩展到 \(b_T\gtrsim1/\Lambda_{\rm QCD}\) 时，OPE幂修正和软非微扰核可能同时重要，不能只用多项式拟合替代物理论证。

一个常见的处理是把 \(b_T\) 分解为
\begin{equation}
  b_T^*=\frac{b_T}{\sqrt{1+b_T^2/b_{\max}^2}},
  \label{eq:bstar}
\end{equation}
把 \(b_T^*\) 用于微扰核，把 \(b_T-b_T^*\) 的效应吸收到非微扰Sudakov因子。该处理不是从第一性原理唯一推出的，必须报告 \(b_{\max}\) 依赖、长距核的拟合形式和误差。若格点数据只在 \(b_T/a\le4\) 上通过流窗口，根本没有足够数据检验 \(b_T^*\) 的连续极限，不能把模型外推结果当作物理TMD。

\section{\texorpdfstring{有限 \(P_z\)、有限 \(a\) 与有限 \(L\)}{有限 Pz、有限 a 与有限 L}}

把因子化写成
\begin{equation}
  \widetilde f_g(x,\bT;P_z;a,L)
  =H_g\otimes f_g
  +\mathcal O\left(\frac{M_N^2}{P_z^2},
  \frac{\Lambda_{\rm QCD}^2}{P_z^2},
  \frac{\Lambda_{\rm QCD}^2}{b_T^2P_z^2},a\Lambda_{\rm QCD},
  \frac{z_{\max}}{L}\right)
  \label{eq:fact-errors}
\end{equation}
时，必须说明每一项如何检测。固定 \(P_z\) 只做 \(a\to0\) 无法区分 \(aP_z\) 固定带来的混合；固定 \(a\) 只做 \(P_z\to\infty\) 会把离散Wilson线误差当成幂修正。合理策略是做二维窗口扫描，并用 \(aP_z\) 同时增大的序列检验稳定性。\(L\to\infty\) 则检查大 \(z\)、软函数和低动量模是否受周期镜像影响。

\section{方案对照表}

表~\ref{tab:factorization-schemes} 汇总常见TMD方案的适用范围。它不是方案优劣排名，而是提醒每个中间量都必须和对应硬核成对出现。

\begin{table}[htbp]
\centering
\caption{TMD/准TMD常见方案及其配套对象。}
\label{tab:factorization-schemes}
\small
\begin{tabularx}{\textwidth}{L{3.0cm}L{4.2cm}Y}
\toprule
方案 & 主要减除对象 & 必须配套的报告项 \\
\midrule
\(\MSbar\) & UV端点对数与局域反常量纲 & 匹配系数阶数、\(\mu\) 初值、Wilson线表示 \\
Wilson圈 & 线性自能、pinch pole、部分 cusp & 圈几何、\(L_\parallel\) 扫描、短距参考点 \\
短距比率 & 零动量参考点与剩余UV & 参考点范围、\(z_0,b_{T,0}\)、连续极限 \\
Hybrid & 短距微扰与长距质量减除 & \(z_s\) 窗口、混合参数、短长距重叠区 \\
软函数减除 & rapidity/几何相关软发散 & regulator、\(\zeta,\nu\)、CS核和软表示 \\
流时匹配 & 流时UV与短距系数 & \(\tau/a^2\)、\(a\to0\)、流步长与单位 \\
\bottomrule
\end{tabularx}
\end{table}

\section{一个可执行的匹配工作流}

为避免从矩阵元直接跳到物理TMD，表~\ref{tab:matching-workflow} 把匹配拆成可以逐项验收的动作。每一步失败时，都应回到该步而不是通过后续拟合补偿。

\begin{table}[htbp]
\centering
\caption{TMD匹配工作流及停止条件。}
\label{tab:matching-workflow}
\small
\begin{tabularx}{\textwidth}{L{1.1cm}L{4.1cm}Y}
\toprule
步骤 & 动作 & 通过条件与失败信号 \\
\midrule
1 & 固定算符、路径、表示和Fourier约定 & 符号表闭合；\(P_z\) 反向与复共轭关系满足 \\
2 & 基于同几何Wilson圈完成UV减除 & 多 \(L_\parallel\) 或短距参考点给出稳定平台；否则不进入软减除 \\
3 & 构造软/rapidity减除并选CS标度 & 多个 \(P_z\) 下CS核稳定；否则检查regulator和表示 \\
4 & 计算硬核矩阵并纳入口味混合 & 硬核阶数与 \(\mu\) 依赖一致；混合矩阵未遗漏 \\
5 & 检查 \(1/P_z^2,aP_z,z_{\max}/L\) 窗口 & 二维外推方向一致；否则缩小发布区间 \\
6 & 输出TMD及误差预算 & 统计、系统、匹配阶数和方案误差均有证据 \\
\bottomrule
\end{tabularx}
\end{table}

\begin{physicalbox}{因式分解的判据}
一个“因子化结果”至少要同时满足：无未减除的UV/rapidity发散；不同 \(P_z\)、\(a\) 和 \(L\) 的窗口收敛；匹配核与软核使用同一表示和几何；幂修正的符号与量级可估计；高扭度项与目标精度相容。缺一项时，结果可以是有价值的中间矩阵元，但不能作为完整TMD发布。
\end{physicalbox}

\clearpage
\part{部分子分布、OPE 与匹配}
\chapter{共线部分子分布、OPE与DGLAP演化}

本章把强子内部分子的纵向动量信息放在共线因子化语言中讨论。重点不是把PDF当作
单一标度上的经验曲线，而是固定三件事：光锥算符的定义、短距离OPE的算符基底，
以及因子化标度改变时PDF的演化。只有这三件事同时一致，格点矩阵元、微扰匹配和
实验全局分析才能放在同一条证据链上。

\section{光锥坐标与PDF定义}

取
\begin{equation}
  v^\pm=\frac{v^0\pm v^3}{\sqrt2},\qquad
  v\cdot w=v^+w^-+v^-w^+-v_\perp\cdot w_\perp,
  \qquad g_{+-}=g_{-+}=1,
  \label{eq:b04-lightcone-metric}
\end{equation}
并令强子沿 $z$ 方向运动。沿纵向的boost只重标度 $v^+$ 与 $v^-$。因此，光锥
关联天然把大动量极限、纵向动量分数和横向平移分开处理。

非极化夸克与胶子光锥PDF定义为
\begin{align}
  q(x,\mu)&=\int\frac{\dd\xi^-}{4\pi}\,
  e^{-ixP^+\xi^-}
  \langle P|\bar\psi(\xi^-)\gamma^+
  W(\xi^-,0)\psi(0)|P\rangle_\mu,
  \label{eq:b04-lc-quark}
  \\
  xg(x,\mu)&=\int\frac{\dd z^-}{2\pi P^+}\,
  e^{-ixP^+z^-}
  \langle P|F^{a}_{+\mu}(z^-)W_{\rm adj}^{ab}(z^-,0)
  F^{b\mu}_{\phantom{b\mu}+}(0)|P\rangle_\mu .
  \label{eq:b04-lc-gluon}
\end{align}
其中Wilson线取路径排序指数，规范链的方向和端点约定必须与矩阵元中的场序一致。
$\mu$是UV重整化与因子化标度，不是强子动量。胶子定义中的$1/P^+$与左侧的$x$
共同固定动量求和规则；若改用$g(x)$而非$xg(x)$，所有矩和匹配核的归一化都要
同步改变。

在光锥规范$A^+=0$下，规范链退化为单位矩阵，正支$q(x)$可解释为夸克动量密度，
负支给出反夸克分布。物理支撑为$x\in[-1,1]$，胶子采用$x>0$的正支并在前向极限
写作$H^g(x,0,0)=xg(x)$。横动量已经积掉，故PDF不是三维分布，也不包含TMD软函数。

动量与价数守恒分别给出
\begin{equation}
  \int_0^1\dd x\,x\left[\sum_f q_f(x,\mu)+g(x,\mu)\right]=1,
  \qquad
  \int_0^1\dd x\left[q_f(x,\mu)-\bar q_f(x,\mu)\right]=N_f,
  \label{eq:b04-pdf-sumrules}
\end{equation}
其中$N_f$是该味的净价夸克数。第一条在微扰演化中不应漂移；第二条只对相应味
组合成立。数值实现应同时检查局部归一化、端点行为和这些守恒量，而不能只比较
中心曲线。

\section{DGLAP方程与劈裂核}

共线辐射产生$\alpha_s^n\ln^n(\mu^2/\mu_0^2)$型大对数。把这些对数吸收到PDF后，
非单态分布满足
\begin{equation}
  \frac{\dd f_i^{(-)}(x,\mu)}{\dd\ln\mu^2}
  =\frac{\alpha_s(\mu)}{2\pi}
  \int_x^1\frac{\dd y}{y}\,P_{qq}\!\left(\frac{x}{y}\right)
  f_i^{(-)}(y,\mu),
  \qquad f_i^{(-)}=f_{q_i}-f_{\bar q_i}.
  \label{eq:b04-dglap-nonsinglet}
\end{equation}
味单态$\Sigma=\sum_i(f_{q_i}+f_{\bar q_i})$与胶子耦合为
\begin{equation}
  \frac{\dd}{\dd\ln\mu^2}
  \begin{pmatrix}\Sigma(x,\mu)\\g(x,\mu)\end{pmatrix}
  =\frac{\alpha_s(\mu)}{2\pi}
  \int_x^1\frac{\dd y}{y}
  \begin{pmatrix}
    P_{qq}(x/y)&2n_fP_{qg}(x/y)\\
    P_{gq}(x/y)&P_{gg}(x/y)
  \end{pmatrix}
  \begin{pmatrix}\Sigma(y,\mu)\\g(y,\mu)\end{pmatrix}.
  \label{eq:b04-dglap-singlet}
\end{equation}
卷积变量$y>x$表示动量分数为$x$的部分子来自演化前分数更大的母部分子。$2n_f$
是味求和，不是新的颜色因子。

领头阶劈裂核为
\begin{align}
  P_{qq}(z)&=C_F\left[\frac{1+z^2}{(1-z)_+}
  +\frac32\delta(1-z)\right],\notag\\
  P_{qg}(z)&=T_F\left[z^2+(1-z)^2\right],\notag\\
  P_{gq}(z)&=C_F\frac{1+(1-z)^2}{z},\notag\\
  P_{gg}(z)&=2C_A\left[\frac{z}{(1-z)_+}
  +\frac{1-z}{z}+z(1-z)\right]
  +\frac{\beta_0}{2}\delta(1-z),
  \label{eq:b04-splitting}
\end{align}
其中
\begin{equation}
  C_F=\frac{N_c^2-1}{2N_c},\qquad T_F=\frac12,\qquad C_A=N_c,
  \qquad \beta_0=\frac{11}{3}C_A-\frac43T_Fn_f .
  \label{eq:b04-color}
\end{equation}
对$SU(3)$有$C_F=4/3$、$C_A=3$。式\eqref{eq:b04-splitting}中的$+$分布定义为
\begin{equation}
  \int_0^1\dd z\,[F(z)]_+\,G(z)
  =\int_0^1\dd z\,F(z)[G(z)-G(1)].
  \label{eq:b04-plus}
\end{equation}
它实现端点减除，保证劈裂概率在积分意义下有限。把$+$分布当作普通函数逐点相加，
会在$z=1$产生未减除的假发散。

\paragraph{物理意义.}
$P_{ab}(z)$不是任意多项式，它编码规范理论中$b\to a$的共线辐射概率密度。
$P_{qq}$与$P_{gq}$保持费米子线，$P_{qg}$描述胶子产生夸克对，$P_{gg}$包含
非Abel三胶子顶点。动量求和规则固定若干$\delta(1-z)$系数，并使单态演化矩阵在
$n=1$矩处具有零本征值。

\section{Mellin矩与OPE等价性}

定义
\begin{equation}
  M_n(\mu)=\int_0^1\dd x\,x^{n-1}f(x,\mu),
  \qquad
  \gamma_n=\int_0^1\dd z\,z^{n-1}P(z).
  \label{eq:b04-mellin}
\end{equation}
卷积在Mellin空间变成乘法：
\begin{equation}
  \frac{\dd M_n^{(-)}(\mu)}{\dd\ln\mu^2}
  =\frac{\alpha_s(\mu)}{2\pi}\gamma_{qq}^{(n)}M_n^{(-)}(\mu).
  \label{eq:b04-moment-dglap}
\end{equation}
对$P_{qq}$显式计算得到
\begin{equation}
  \gamma_{qq}^{(n)}
  =C_F\left[
  \frac32+\frac{1}{n(n+1)}
  -2\sum_{k=2}^{n}\frac1k\right],
  \qquad n\geq1.
  \label{eq:b04-gamma-qq}
\end{equation}
在固定耦合近似下，矩按幂律演化；在领头阶跑动耦合下，
\begin{equation}
  M_n(Q^2)=M_n(Q_0^2)
  \left[
  \frac{\ln(Q_0^2/\Lambda_{\rm QCD}^2)}
       {\ln(Q^2/\Lambda_{\rm QCD}^2)}
  \right]^{\gamma_n/(2\beta_0)} .
  \label{eq:b04-moment-solution}
\end{equation}
大$n$矩对应大$x$区域。该式说明强子PDF的大$x$部分随$Q^2$发生对数压低，但
$\Lambda_{\rm QCD}$和非微扰初值不由DGLAP本身产生。

同一个$\gamma_n$也可由twist-2局域算符的重整化得到。对相应的无迹对称化算符，
\begin{equation}
  \langle P|\mathcal O_n(\mu)|P\rangle
  =2a_n(\mu)\left(P^{\mu_1}\cdots P^{\mu_n}-\text{trace}\right),
  \qquad
  a_n(\mu)=M_n(\mu).
  \label{eq:b04-ope-moment}
\end{equation}
因此DGLAP劈裂核的Mellin矩等于相应twist-2算符的反常维数：
\begin{equation}
  \gamma_{n}^{\rm DGLAP}=\gamma_n^{\rm OPE}.
  \label{eq:b04-ope-dglap}
\end{equation}
这不是形式类比，而是“共线辐射重求和”和“局域算符重整化”描述同一短距离物理的
两种语言。它也为格点局域矩、准PDF匹配和DGLAP演化之间的对应提供理论基础。

\section{尺度窗口与失效条件}

DGLAP是因子化标度演化方程，不是完整的非微扰求解器。最低限度要求
\begin{equation}
  \Lambda_{\rm QCD}\ll\mu,\qquad
  \alpha_s(\mu)\ll1,\qquad
  \text{目标$x$处端点对数受控}.
  \label{eq:b04-dglap-window}
\end{equation}
在极小二$x$处，固定阶DGLAP可能遗漏BFKL型或饱和型重求和；在大$x\to1$处，
阈值修正、软辐射和高阶twist可能同时增强。若把$\mu$降到Landau极点附近，劈裂核
的微扰展开本身失去意义。任何“由DGLAP推出小$x$上升”或“由DGLAP精确预言端点”
的陈述，都必须先说明因子化窗口。

\begin{table}[htbp]
  \centering
  \small
  \caption{DGLAP对象与常见误用}
  \label{tab:b04-dglap-check}
  \begin{tabularx}{\textwidth}{@{}lXX@{}}
    \toprule
    对象 & 正确含义 & 常见误用 \\
    \midrule
    $P_{ab}(z)$ & 共线劈裂核，含端点减除 & 当作未减除的逐点概率 \\
    $\gamma_n$ & 劈裂核或twist-2算符的反常维数 & 与QCD $\beta$函数混同 \\
    PDF初值 & 非微扰输入，需实验或格点确定 & 声称由DGLAP唯一预言 \\
    $\mu$依赖 & 方案与标度依赖，可被系数函数补偿 & 与物理强子动量$P_z$混同 \\
    \bottomrule
  \end{tabularx}
\end{table}

\chapter{OPE、光线展开与幂次修正}

\section{短距离OPE与格点高矩困难}

两个算符在短距离下的乘积可展开为
\begin{equation}
  \mathcal O_1(z)\mathcal O_2(0)
  \xrightarrow[z\to0]{}
  \sum_n C_n(z;\mu)\mathcal O_n(0),
  \label{eq:b04-short-ope}
\end{equation}
其中Wilson系数$C_n$含短距离奇异性，局域算符$\mathcal O_n$的矩阵元编码非微扰
物理。量纲递增的算符被$(z\Lambda_{\rm QCD})$的幂次压低。把光锥双线性在
$z=0$附近Taylor展开，得到
\begin{equation}
  \bar\psi(z^-)\gamma^+W(z^-,0)\psi(0)
  =\sum_{n=0}^{\infty}\frac{(z^-)^n}{n!}
  \bar\psi(0)\gamma^+(iD^+)^n\psi(0)+\cdots ,
  \label{eq:b04-taylor}
\end{equation}
其核子矩阵元的Fourier变换给出Mellin矩。

传统格点路线直接计算这些局域算符，但在超立方群下高自旋算符会发生复杂混合，
高维算符还可能与低维算符产生$1/a^{\Delta d}$型幂次发散混合，统计噪声也随导数数
增加。因此实际可稳定控制的矩数有限，不能仅凭前几个矩无模型重建完整$x$依赖。
保持非局域形式可绕开高自旋局域算符基底，但不表示非局域算符没有UV发散，Wilson线
自能仍产生线性项。

\section{光线OPE与twist截断}

有限空间分离$z$不满足$z\Lambda_{\rm QCD}\ll1$的普通短距条件。光线OPE采用
\begin{equation}
  z=\mathcal O(\eta),\qquad P=\mathcal O(\eta^{-1}),\qquad
  zP=\mathcal O(1),\qquad z\Lambda_{\rm QCD}\to0,
  \label{eq:b04-lightray-counting}
\end{equation}
使Ioffe时间$zP$保持有限，但所有twist-2局域算符的贡献必须重求和。其结构为
\begin{equation}
  \bar q(zv)\slashed v[zv,0]q(0)
  =\int_0^1\dd\alpha\,H_\parallel(z,\alpha;\mu_F)
  \Pi_{\rm l.t.}\!\left[\bar q(\alpha zv)\slashed z q(0)\right]
  +\text{higher twist}.
  \label{eq:b04-lightray-ope}
\end{equation}
这里$\Pi_{\rm l.t.}$是leading-twist投影，$H_\parallel$是可微扰计算的系数。
对强子矩阵元取期望值后，hard与collinear部分进入匹配核，软矩阵元进入光锥PDF或
ITD。这就是准PDF、空间关联和赝PDF服从同一类因子化的OPE依据。

普通短距OPE在$\nu=zP\to0$时截断高维算符；光线OPE保持有限$\nu$而按twist截断。
“高自旋”与“高twist”不是同一个展开参数。一个算符可以自旋很高但仍是twist-2；
把所有高自旋项误称为高twist，会完全破坏准PDF的因子化计数。

\section{Wilson系数、重正子与精度层级}

准PDF与光锥PDF有相同的红外共线结构，差别吸收入UV Wilson系数。单圈裸系数同时
包含UV极点和共线IR极点。UV极点由局域抵消项减除，IR极点与光锥系数中的对应项
抵消，剩下有限对数与常数构成匹配核。若两边的IR结构不一致，就不能用“减去一个
微扰核”修补，必须重新检查算符、投影或因子化定理。

在$\overline{\rm MS}$中，匹配系数通常只是渐近级数。其大阶系数可写成
\begin{equation}
  h_k\sim k!\left(\frac{\beta_0}{2\pi}\right)^k,
  \qquad k\to\infty .
  \label{eq:b04-factorial}
\end{equation}
Borel变换
\begin{equation}
  \mathcal B[H](w)=\sum_{k}\frac{h_k}{k!}w^k
  \label{eq:b04-borel}
\end{equation}
因此出现极点。Wilson线自能关联的UV重正子在$w=1/2$附近，IR重正子位于整数
$w=1,2,\ldots$。Borel平面上的IR极点不是数值偶然，而是微扰系数与非微扰幂修正
共享同一渐近结构的信号。位于积分路径上的极点使Borel和方案依赖；用主值或其他
正规化时，应把相应模糊性显式吸收到高阶twist参数。

\begin{table}[htbp]
  \centering
  \small
  \caption{准PDF幂次修正的精度层级}
  \label{tab:b04-power-hierarchy}
  \begin{tabularx}{\textwidth}{@{}lXXX@{}}
    \toprule
    层级 & 典型尺度 & 可控制的物理 & 失配风险 \\
    \midrule
    twist-2 & $O(1)$ & PDF与leading-twist匹配 & 把有限$P_z$残差全部归于统计误差 \\
    包含twist-3 & 通常新增$O(\Lambda/P_z)$ & 领先幂次、线性Wilson线模糊 & 端点增强与重正子处方不统一 \\
    包含twist-4 & 通常新增$O(\Lambda^2/P_z^2)$ & 靶质量、常规$1/P_z^2$修正 & 仍有更高twist或方案转换遗漏 \\
    \bottomrule
  \end{tabularx}
\end{table}

端点增强需要单独说明。重正子分析给出的准PDF紧凑形式可出现
\begin{equation}
  Q_\parallel(x,P_z)=q(x)\left\{
  1+\mathcal O\!\left[
  \frac{\Lambda_{\rm QCD}^2}{P_z^2}\frac{1}{x^2(1-x)}
  \right]\right\},
  \label{eq:b04-qpdf-power}
\end{equation}
其$1/P_z^2$不能仅凭幂次反推为twist-4。Fourier变换、端点奇异和重正子重排会把
twist-3的非微扰信息组织成额外的$x$依赖；判断twist必须回到OPE算符，而不是数
最终公式中的$P_z$指数。相反，只保留twist-2时，常把剩余项简写为
$\mathcal O(\Lambda_{\rm QCD}^2/P_z^2,M^2/P_z^2)$。两种陈述对应不同的精度目标，
必须在同一段文字中先声明“是否显式计入twist-3”。

\chapter{LaMET、准PDF与赝PDF}

\section{Euclidean障碍与准PDF}

光锥关联依赖实时间，而格点Monte Carlo计算Euclidean关联。LaMET不把实时间直接
解析延拓到格点上，而是先计算空间分离的等时矩阵元，再利用大动量与短距离因子化。
准夸克PDF定义为
\begin{equation}
  \tilde q_\Gamma(x,P_z)
  =\int_{-\infty}^{\infty}\frac{\dd z}{4\pi}\,
  e^{ixP_z z}
  \langle P_z|\bar\psi(z)\Gamma U(z,0)\psi(0)|P_z\rangle_R,
  \label{eq:b04-quasi}
\end{equation}
其中$U$是沿$z$的直空间Wilson线。非极化常用$\Gamma=\gamma^t$或$\gamma^z$，
螺旋度用$\gamma^z\gamma_5$或$\gamma^t\gamma_5$，横向性用$\sigma^{z\mu}$。
准PDF可格点计算，但依赖$P_z$、定义域扩展到整个实轴、没有严格概率诠释。$|x|>1$
的贡献不是新物理自由度，而是有限动量和空间分离造成的非光锥结构，应在$P_z\to\infty$
与匹配后消失。

\section{因子化、极限次序与有效场论}

在leading-twist因子化成立时，
\begin{equation}
  \tilde q(x,P_z)
  =\int_{-1}^{1}\frac{\dd y}{|y|}\,
  C\!\left(\frac{x}{y},\frac{\mu}{|y|P_z}\right)q(y,\mu)
  +\mathcal O\!\left[
  \frac{\Lambda_{\rm QCD}^2}{x^2P_z^2},
  \frac{\Lambda_{\rm QCD}^2}{(1-x)^2P_z^2},
  \frac{M^2}{P_z^2}
  \right].
  \label{eq:b04-lamet-factorization}
\end{equation}
$C$从$\delta(1-x/y)$开始作微扰展开。hard模式进入$C$，collinear非微扰矩阵元
进入$q$。若需要显式twist-3精度，式\eqref{eq:b04-lamet-factorization}的剩余项
必须替换为相应高阶twist贡献，不能继续把误差只写成$1/P_z^2$。

LaMET的有效场论类比以$P_z$承担大标度：
\begin{equation}
  F(P_z/\Lambda_{\rm QCD})
  =Z(P_z/\Lambda_{\rm QCD},\Lambda_{\rm QCD}/\mu)f(\mu)
  +\text{higher powers}.
  \label{eq:b04-lamet-eft}
\end{equation}
它类似HQET把$m_Q$作为大标度，但相似性是结构性的，不是把胶子质量替换成$P_z$。
光锥定义先取$P_z\to\infty$再完成UV重整化；格点先有$1/a$截断，再取有限$P_z$。
两类极限次序的差异进入匹配核。可行的格点窗口至少要求
\begin{equation}
  \max(\Lambda_{\rm QCD},M)\ll P_z\ll a^{-1},
  \qquad \text{且}\qquad \alpha_s(\mu_{\rm match})\ll1 .
  \label{eq:b04-lamet-window}
\end{equation}
不同准算符只要具有同一光锥极限，就属于同一普适类，可选幂修正较小者；普适类
不意味着有限$P_z$下各算符的数值和系统误差相同。

\section{赝PDF、约化ITD与短距离因子化}

同一空间矩阵元也可保持在坐标空间：
\begin{equation}
  \mathcal M(z^2,z\cdot P)
  =\langle P|\bar\psi(z)\Gamma U(z,0)\psi(0)|P\rangle,
  \qquad \nu=zP_z .
  \label{eq:b04-itd}
\end{equation}
约化ITD用零动量或真空矩阵元作分母，
\begin{equation}
  \mathfrak M_{\rm red}(\nu,z^2)
  =\frac{\mathcal M(z^2,\nu)}
  {\mathcal M(z^2,0)},
  \label{eq:b04-ritd}
\end{equation}
以压低只依赖$z^2$的端点与自能结构。赝PDF是$\nu$的Fourier变换，并在$z^2\to0$
回到光锥PDF。其短距离因子化为
\begin{equation}
  \begin{aligned}
  \mathfrak M_{\rm red}(\nu,z^2)
  &=\int\dd w\,C(w,z^2\mu^2)
  I(w\nu,\mu^2)+\mathcal O(z^2\Lambda_{\rm QCD}^2),\\
  I(\nu,\mu^2)&=\int_{-1}^{1}\dd x\,e^{ix\nu}f(x,\mu^2).
  \end{aligned}
  \label{eq:b04-pseudo-factorization}
\end{equation}
准PDF与赝PDF的匹配核通过Fourier变换对应，但端点外推与幂修正的组织不同，不能
把一套数值尾部拟合直接搬到另一套。

\begin{table}[htbp]
  \centering
  \small
  \caption{准PDF与赝PDF的工作对象}
  \label{tab:b04-quasi-pseudo}
  \begin{tabularx}{\textwidth}{@{}lXX@{}}
    \toprule
    项目 & 准PDF & 赝PDF/ITD \\
    \midrule
    基本量 & $x\tilde q(x,P_z)$或$\tilde q$ & $\mathfrak M_{\rm red}(\nu,z^2)$ \\
    大标度 & $P_z\to\infty$，再做Fourier匹配 & $z\to0$，同时保持$\nu=zP_z$有限 \\
    端点问题 & 小$x$、$x\to1$及$|x|>1$的Fourier振荡 & 大$z$尾部与短距匹配窗口 \\
    典型幂修正 & 端点增强的$P_z^{-2}$组织，可含twist-3 & $z^2\Lambda_{\rm QCD}^2(1-x)$组织 \\
    主要验证 & 匹配核、RGR、端点稳定性 & 短距因子化、$\nu$依赖与$z^2$残差 \\
    \bottomrule
  \end{tabularx}
\end{table}

\section{内禀标度与重整化群重求和}

准PDF匹配核的NLO对数可写成
\begin{equation}
  C^{(1)}(x,P_z;\mu)
  \supset \frac{\alpha_sC_F}{2\pi}
  \ln\!\left(\frac{4x^2P_z^2}{\mu^2}\right),
  \label{eq:b04-matching-log}
\end{equation}
其自然标度为
\begin{equation}
  \mu_{\rm eff}\sim Q_{\rm eff}=2|x|P_z .
  \label{eq:b04-qeff}
\end{equation}
当$Q_{\rm eff}$远离参考标度$\mu$时，固定阶匹配含大对数。RGR先在
$\mu_0\sim Q_{\rm eff}$处做固定阶反匹配，再用DGLAP把PDF演化到统一标度。准PDF
匹配核的重整化群核与DGLAP劈裂核相同，这不是额外假设，而是二者共同源于共线
重整化的结果。

小$x$窗口由$2|x|P_z$不能进入非微扰区决定。数量级上
\begin{equation}
  x_{\min}\sim\frac{\Lambda_{\rm QCD}}{2P_z}.
  \label{eq:b04-xmin}
\end{equation}
在$P_z=2.5$~GeV、$\Lambda_{\rm QCD}\sim0.3$~GeV时，这个尺度约为$0.06$。这是
量级估计，不是普适阈值；系数、匹配阶数、方案和$P_z$窗口都会改变可接受范围。
当$Q_{\rm eff}$接近Landau极点时，继续作固定阶DGLAP或微扰匹配没有理论依据。

\section{端点幂修正与证据等级}

基于重正子分析的赝PDF紧凑形式常写作
\begin{equation}
  \mathcal P(x,z)=q(x)\left\{1+\mathcal O[z^2\Lambda_{\rm QCD}^2(1-x)]\right\},
  \label{eq:b04-ppdf-power}
\end{equation}
与式\eqref{eq:b04-qpdf-power}的端点增强结构不同。该差异说明qPDF与pPDF的优劣
依赖目标$x$区间、$P_z$、$z$窗口和重正化方案，不存在全区间无条件的“某一方案更好”。
靶质量修正可由$M^2/P_z^2$显式计算，但它没有式\eqref{eq:b04-qpdf-power}的$1/x^2$
增强，因而不能把运动学靶质量项误当作高阶twist证据。

\chapter{夸克与胶子算符的格点实现}

\section{Dirac分类、连续混合与RI/MOM}

连续夸克双线性按
\begin{equation}
  \Gamma\in\{I,\gamma_5,\gamma_\mu,\gamma_\mu\gamma_5,
  \sigma_{\mu\nu}=\tfrac{i}{2}[\gamma_\mu,\gamma_\nu]\}
  \label{eq:b04-bilinears}
\end{equation}
分为标量、赝标量、矢量、轴矢和张量。Euclidean约定取
\begin{equation}
  \{\gamma_\mu,\gamma_\nu\}=2\delta_{\mu\nu},
  \qquad
  \gamma_4=\gamma^0_M,\qquad
  \gamma_i=i\gamma^i_M,\quad i=1,2,3.
  \label{eq:b04-euclidean-gamma}
\end{equation}
因此Minkowski光锥投影不能未经约定检查就与$\gamma^t$或$\gamma^z$等同。

对非极化准PDF，源文献在所述非前行、非极化运动学下指出：$\Gamma=\gamma^z$可在
连续极限与twist-3标量型双线性发生$O(a^0)$混合，必须以$2\times2$重整化矩阵处理；
$\Gamma=\gamma^t$不与该标量算符混合，因而简化RI/MOM重整化。螺旋度通道有类似
结构，$\gamma^z\gamma_5$可能混入twist-3轴矢结构。该结论依赖算符、Wilson线方向、
投影和运动学，不能推广成“$\gamma^t$对所有GPD/TMD都无混合”的普遍定理。

RI/MOM条件可概括为
\begin{equation}
  Z_q^{-1}Z_O(z;\mu_R)
  \frac1{12}\operatorname{Tr}
  \!\left[V(p,z)V_{\rm Born}^{-1}(p,z)\right]_{p^2=\mu_R^2}
  =1,
  \label{eq:b04-ri-mom}
\end{equation}
其中$V$为截肢顶点，$Z_q$由夸克传播子确定。最小投影与$\slashed p$投影在特定
参考动量下可接近，但其差异随投影和动量增加。该差异必须进入匹配和系统误差，
不能只在最终图上选取看起来更稳定的投影。

\section{Wilson线与非局域重整化}

携带长Wilson线的矩阵元含有线性UV发散。辅助场表述
\begin{equation}
  Z_x=L(x,\infty),\qquad
  (\partial_x-igA_x)Z_x=0,
  \qquad L(x,y)=Z_xZ_y^\dagger
  \label{eq:b04-zfield}
\end{equation}
把Wilson线改写为辅助重夸克两点函数，从而得到
\begin{equation}
  L_{\rm ren}(x,y)=Z_L^{-1}e^{-\delta m|x-y|}L(x,y).
  \label{eq:b04-wilson-ren}
\end{equation}
$\delta m$吸收线性发散，$Z_L$处理剩余对数。短距ratio、hybrid与
self-renormalization是组织$\delta m$、端点因子和长距矩阵元的不同方案。普通
局域RI/MOM顶点不能自动减除整条Wilson线的自能；若$z_s$、参考点或$\delta m+m_0$
变化大于目标误差，就说明短长距分离尚未受控。

胶子算符有两个场强指标和伴随Wilson线。基础表示中的可乘性组合为
\begin{equation}
  \mathcal O_g(z)
  =M_{tx;tx}(z)+M_{ty;ty}(z)-2M_{xy;xy}(z),
  \label{eq:b04-gluon-combination}
\end{equation}
其中
\begin{equation}
  M_{\mu\lambda;\nu\rho}(z)
  =\sum_{\vec x}\operatorname{Tr}\!\left[
  F_{\mu\lambda}(\vec x+z\hat z)U(\vec x+z\hat z,\vec x)
  F_{\nu\rho}(\vec x)U(\vec x,\vec x+z\hat z)\right].
  \label{eq:b04-gluon-M}
\end{equation}
负号来自Minkowski与Euclidean时间-横向分量的转换。该组合使各构建块在单圈水平
具有匹配的UV反常量纲，但“存在一个可乘性组合”不等于所有Lorentz投影都可乘性化。
辅助重夸克分析识别出多个独立构建块，并给出不是可乘性组合的反例；实现时必须按
指标数、交叉结构和端点因子分类。

\section{Clover场强、蒸馏与三点函数}

Plaquette的短距展开为
\begin{equation}
  P_{\mu\nu}=1+ia^2F_{\mu\nu}(x)+\mathcal O(a^3),
  \label{eq:b04-plaquette}
\end{equation}
对称Clover组合
\begin{equation}
  F_{\mu\nu}(x)=-\frac{i}{8a^2}[Q_{\mu\nu}(x)-Q_{\nu\mu}(x)]
  \label{eq:b04-clover}
\end{equation}
减小统计涨落。Clover系数、场强符号、链接方向和场强指标的缩并必须统一；只检查
某个局域纯规范极限不足以验证非局域胶子算符。

蒸馏通过三维协变Laplace算符的低模构造投影
\begin{equation}
  \Box_{xy}(t)=\sum_{k=1}^{N_{\rm ev}}v_x^{(k)}(t)v_y^{(k)\dagger}(t)
  =V(t)V^\dagger(t).
  \label{eq:b04-distillation}
\end{equation}
传播子在蒸馏空间中的perambulator为
\begin{equation}
  \tau^{(k,\ell)}_{\alpha\beta}
  =V^{(k)\dagger}D^{-1}_{\alpha\beta}V^{(\ell)}.
  \label{eq:b04-perambulator}
\end{equation}
它只依赖规范组态和Dirac算符，可在插值算符基中复用。高动量必须使用相位蒸馏；
换相位、源型、场强投影或Wilson线几何时，缓存键必须相应改变。蒸馏降低的是
传播子与收缩成本，不自动改善算符UV结构或消除激发态污染。

胶子三点函数与核子线不直接连通，必须做算符真空减除和两点函数真空减除。常用的
基态主导拟合可写成
\begin{equation}
  R_{3{\rm pt}}(T,\tau)
  =\frac{C_{3{\rm pt}}(T,\tau)}{C_{2{\rm pt}}(T)}
  \simeq c_0+c_1e^{-\Delta E(T-\tau)}
  +c_1e^{-\Delta E\tau}.
  \label{eq:b04-gluon-ratio}
\end{equation}
当$T$不足、插入点靠近源或汇、或只使用单一$z$时，$c_0$可能仍是污染混合后的
有效平台。至少需要多个$T$、多个插入窗、正负$z$和独立两态或多态拟合来检验。

\section{匹配核与连续外推}

重整化胶子矩阵元先Fourier变换为准分布，再与光锥胶子PDF作卷积匹配。味单态通道
还必须在匹配层区分夸克与胶子：
\begin{equation}
  \begin{pmatrix}\tilde g\\ \tilde q_S\end{pmatrix}
  =\int\frac{\dd y}{y}
  \begin{pmatrix}C_{gg}&C_{gq}\\C_{qg}&C_{qq}\end{pmatrix}
  \begin{pmatrix}g\\q_S\end{pmatrix}
  +\text{higher twist}.
  \label{eq:b04-gluon-matching}
\end{equation}
这与UV重整化混合不是同一对象。匹配核阶数、$\mu$选择、RGR和大$\lambda=zP_z$
外推都会影响最终结果。

连续极限与无穷动量极限不能由单点结果替代。常用的联合拟合物为
\begin{equation}
  f(x,P_z,a)
  =f(x)+b_1(x)a^2+b_2(x)a^2P_z^2
  +b_3(x)P_z^{-2}+\cdots .
  \label{eq:b04-joint-extrap}
\end{equation}
至少需要多个格距和多个动量才能辨识主要项。若只有单一$a$或单一$P_z$，式
\eqref{eq:b04-joint-extrap}退化，结果只能称为该参数点的有效结果。

\chapter{广义部分子分布与全局结构}

\section{非前行运动学与定义}

GPD把PDF的同一非局域双线性夹在不同初末强子态之间。固定
\begin{equation}
  P=\frac{P_i+P_f}{2},\qquad
  \Delta=P_f-P_i,\qquad
  t=\Delta^2,\qquad
  \xi=\frac{P_i^+-P_f^+}{P_i^++P_f^+}
  =-\frac{\Delta^+}{2P^+},
  \label{eq:b04-gpd-kinematics}
\end{equation}
于是$P_i=P-\Delta/2$，$P_f=P+\Delta/2$。$x$与平均纵向动量相联系，$\xi$描述
纵向动量转移，$t$的Fourier共轭与横向位置相联系。定义中必须同时固定$\Delta$的
方向、Fourier相位、旋量约定和$E$项符号；只改其中一项会造成不可追踪的整体换号。

核子矢量GPD的光锥矩阵元分解为
\begin{equation}
  \begin{aligned}
  &\int\frac{\dd z^-}{4\pi}e^{ixP^+z^-}
  \langle P_f,S_f|
  \bar\psi(-z^-/2)\gamma^+
  W(-z^-/2,z^-/2)\psi(z^-/2)
  |P_i,S_i\rangle\\
  &=\bar u(P_f,S_f)\left[
  \gamma^+H^q(x,\xi,t)
  +\frac{i\sigma^{+\mu}\Delta_\mu}{2M}E^q(x,\xi,t)
  \right]u(P_i,S_i).
  \end{aligned}
  \label{eq:b04-vector-gpd}
\end{equation}
轴矢通道含有$\widetilde H^q$与$\widetilde E^q$，横向极化还给出四个张量GPD。
核子不可能由一个投影同时读出$H$与$E$；必须选择不同Dirac结构、初末自旋和动量
方向，构造投影矩阵后解线性方程。自旋零介子的最低阶矢量分解只需一个形状函数。

GPD的DGLAP区与ERBL区由$|x|$和$|\xi|$区分：
\begin{equation}
  |x|\geq|\xi|:\ \text{DGLAP区},\qquad
  |x|\leq|\xi|:\ \text{ERBL区},\qquad
  x\in[-1,1],\quad |\xi|\leq1.
  \label{eq:b04-gpd-regions}
\end{equation}
前向区可连续退化为概率型PDF或反夸克分布；ERBL区描述夸克-反夸克振幅，不能直接
称为反夸克概率。演化核也随区域改变，DGLAP核不能原样代入ERBL边界。

\section{前向极限、求和规则与多项式性}

当$\Delta=0$，即$\xi=t=0$时，
\begin{align}
  H^q(x,0,0)&=q(x),&
  H^q(-x,0,0)&=-\bar q(x),&x>0,\notag\\
  \widetilde H^q(x,0,0)&=\Delta q(x),&
  \widetilde H^q(-x,0,0)&=\Delta\bar q(x),&x>0,\notag\\
  H_T^q(x,0,0)&=\delta q(x),&
  H_T^q(-x,0,0)&=-\delta\bar q(x),&x>0.
  \label{eq:b04-gpd-forward}
\end{align}
胶子通道采用$x$加权约定，前向写作$H^g(x,0,0)=xg(x)$。若改变为未加权的$g(x)$，
矩、形状因子和匹配核的归一化必须一起改变。

对$x$积分消去部分子动量信息，留下局域流矩阵元：
\begin{equation}
  \int_{-1}^{1}\dd x\,H^q(x,\xi,t)=F_1^q(t),\qquad
  \int_{-1}^{1}\dd x\,E^q(x,\xi,t)=F_2^q(t).
  \label{eq:b04-gpd-formfactors}
\end{equation}
轴矢通道对应$G_A^q(t)$与$G_P^q(t)$。积分结果与$\xi$无关，是leading-twist
多项式性的约束，不是任意模型自动具有的性质。

定义Mellin矩
\begin{equation}
  M_n^q(\xi,t)=\int_{-1}^{1}\dd x\,x^{n-1}H^q(x,\xi,t),
  \label{eq:b04-gpd-mellin}
\end{equation}
leading-twist下应有
\begin{equation}
  M_n^q(\xi,t)=\sum_{k=0}^{n-1}\xi^kA_{n,k}^q(t).
  \label{eq:b04-gpd-polynomiality}
\end{equation}
宇称和交叉对称性可使部分系数为零。多项式性是检验格点外推、模型参数化和匹配核
端点行为的关键约束：一个在每一点都光滑的二维拟合仍可能违反GPD。

双分布表示自动实现多项式性：
\begin{equation}
  H^q(x,\xi,t)
  =\int_{-1}^{1}\dd\beta\int_{-1+|\beta|}^{1-|\beta|}\dd\alpha\,
  \delta(x-\beta-\xi\alpha)f^q(\beta,\alpha,t)
  +\theta(|\xi|-|x|)D^q\!\left(\frac{x}{\xi},t\right).
  \label{eq:b04-double-distribution}
\end{equation}
$f$生成普通PDF信息和$\xi$多项式矩，$D$-term只在ERBL区出现。$D$不贡献弹性
形状因子，却影响最高阶广义形状因子。缺少$D$-term的模型通常能拟合形状因子，
却不能重现完整多项式结构。

\section{角动量、层析与实验含义}

在$\xi=0,t=0$定义Ji角动量求和规则
\begin{equation}
  J_q=\frac12\left[A_{20}^q(0)+B_{20}^q(0)\right]
  =\frac12\int_{-1}^{1}\dd x\,x[H^q(x,0,0)+E^q(x,0,0)],
  \qquad
  \sum_qJ_q+J_g=\frac12 .
  \label{eq:b04-ji-sumrule}
\end{equation}
数值验证至少需要完整$E$分离、足够$x$覆盖和$t\to0$外推。只计算某个$x$窗口或
单一Dirac投影，不能宣称验证该求和规则。

在$\xi=0$时对横向动量转移作Fourier变换，可定义冲击参数分布
\begin{equation}
  q(x,\vec b_\perp)
  =\int\frac{\dd^2\Delta_\perp}{(2\pi)^2}
  e^{-i\vec\Delta_\perp\cdot\vec b_\perp}
  H^q(x,0,-\Delta_\perp^2).
  \label{eq:b04-impact-parameter}
\end{equation}
这里$\vec b_\perp$是部分子相对强子横向中心的共轭变量，不是TMD中的部分子横向
动量$\vec k_\perp$。两者可通过Wigner型关联联系，但不是同一测量量。

DVCS的Compton振幅是GPD与硬系数的卷积。实验通常先提取Compton Form Factor，
再结合束流/靶极化、方位角和不同过程分离GPD组合，不会逐点直接测出$H(x,\xi,t)$。
DVMP、TCS和双DVCS提供互补约束。格点QCD的价值是直接计算非微扰矩阵元，而不是
绕过因子化、幂修正和模型依赖。

\section{格点三点函数与准GPD}

格点使用空间分离的双局域算符
\begin{equation}
  H^\Gamma_{fi}(z,P_i,P_f)
  =\langle N(P_f,S_f)|\bar q(z_1)\Gamma W(z_1,z_2)q(z_2)
  |N(P_i,S_i)\rangle,
  \qquad
  z_{1,2}=y\mp\frac{z}{2}\hat e_z .
  \label{eq:b04-euclidean-gpd}
\end{equation}
常用投影包括$\gamma^t$、$\gamma^z\gamma_5$和$\gamma^t\gamma_\perp\gamma_5$。
$\gamma^t$在前向非极化准PDF中可避免某类连续混合，但GPD的非前行投影仍需逐项
分析，不能直接照搬。

三点函数的谱分解为
\begin{equation}
  C_3^{fi}(T,\tau;z)
  \xrightarrow[T\gg\tau\gg0]{}
  Z_fH^\Gamma_{fi}(z)e^{-E_f(T-\tau)}e^{-E_i\tau}\overline{Z_i}.
  \label{eq:b04-gpd-3pt}
\end{equation}
当初末态相同时，它退化为前向准PDF型三点函数。非前行GPD还要保留$P_f-P_i$、
初末自旋、Wilson线端点和非局域方向。胶子插入与核子线不连通时，必须进行算符和
两点函数的真空减除，信噪比通常更差。

一般比值可写成
\begin{equation}
  R(\tau)=\frac{C_3^{fi}(T,\tau)}{C_2^f(T)}
  \left[
  \frac{C_2^i(T-\tau)C_2^f(\tau)C_2^f(T)}
       {C_2^f(T-\tau)C_2^i(\tau)C_2^i(T)}
  \right]^{1/2},
  \label{eq:b04-gpd-ratio}
\end{equation}
其平台给出运动学因子乘裸矩阵元。若两点函数因周期边界或反向传播带负号，不能
对每个因子简单取绝对值；应先验证重叠因子的符号约定。至少需要多个$T$、多个
插入窗、不同动量和正负$z$来约束激发态。

非前行准GPD在坐标空间有两个端点Ioffe变量。伪GPD需要用双Fourier变换
\begin{equation}
  \mathcal P(\tau_1,\tau_2,z^2)
  =\int\frac{\dd\zeta_1}{2\pi}\frac{\dd\zeta_2}{2\pi}\,
  e^{i\tau_1\zeta_1+i\tau_2\zeta_2}
  \mathcal H(\zeta_1,\zeta_2,z^2),
  \qquad \zeta_i=z_iP_z,
  \label{eq:b04-pseudo-gpd}
\end{equation}
再通过非前行匹配得到$x,\xi$分布。只对一个端点作Fourier变换会丢失$\xi$与端点
相位信息。leading-twist非前行匹配是多变量双卷积；味非单态只有夸克核，味单态和
螺旋度需要夸克-胶子矩阵，横向性在手征极限不与胶子混合。

匹配实现必须通过以下极限：$\xi\to0$回到前向核；$\Delta\to0$回到前向三点函数；
$t\to0,\xi=0$同时满足形状因子和PDF归一化；$z\to0$回到局域流；$P_z\to\infty$
控制幂修正；$x\to\pm1$不产生非物理Fourier尖峰。任何一项失败都应先查符号、
端点和归一化，而不是提高匹配阶数掩盖低阶错误。

\section{全局拟合、格点输入与证据状态}

PDF全局分析在参考标度参数化非微扰初值，用DGLAP演化到实验标度，再最小化大量
截面与数据的$\chi^2$。快照资料列出的代表性分析为
\begin{center}
  \small
  \begin{tabularx}{0.95\textwidth}{@{}l l l X@{}}
    \toprule
    分析 & 阶数 & 方法 & 主要角色 \\
    \midrule
    CT18 & NNLO & Hessian & 多过程全局约束、误差线性化 \\
    NNPDF4.0 & NNLO & 神经网络+MC & 高灵活参数化与不确定性传播 \\
    MSHT20 & NNLO & Hessian & 与CT家族交叉检验 \\
    JAM24 & NLO/NNLO & Monte Carlo & 自旋与专属过程补充 \\
    \bottomrule
  \end{tabularx}
\end{center}
这些版本是来源文档写作时的历史快照，不应被表述为2026年的最新版本。实验输入覆盖
DIS、固定靶Drell--Yan、弱玻色子、喷注、顶对、孤立光子和粲产生；不同过程主要约束
不同味与$x$区间。大$x$、较大$x$胶子、奇异海和味分离仍存在模型与数据缺口。

格点计算可在大$x$、胶子分布和实验约束较弱的方向提供独立输入。但它与全局拟合的
结合要求：使用同一$\overline{\rm MS}$标度、同一重正化方案、可信的匹配核和完整的
误差相关矩阵。把一条格点中心曲线直接加入全局拟合，而不传播$P_z,a,\mu,z_s$、
拟合窗口和重正化方案误差，会人为夸大约束力。

\begin{table}[htbp]
  \centering
  \small
  \caption{来源文档报告的历史格点证据，本批未独立复算}
  \label{tab:b04-lattice-evidence}
  \begin{tabularx}{\textwidth}{@{}lXX@{}}
    \toprule
    对象 & 参数快照 & 可支持的有限陈述 \\
    \midrule
    同位旋矢量夸克PDF & CLS，$a=0.086$~fm，$M_\pi=356$~MeV，$P_z=2.3$~GeV & 匹配显著改变中等$x$形状；重整化和投影差异在最终PDF中减小 \\
    非极化胶子PDF & 三个间距$0.105,0.0897,0.0775$~fm，$M_\pi\simeq300$~MeV，$P_z\lesssim1.97$~GeV & 在报告误差内与部分全局拟合一致；两端仍有有限$a,P_z$限制 \\
    胶子试点关联 & `C11P29S`，$a=0.105$~fm，$M_\pi=290$~MeV，317组态 & 短$z$平台较稳定，长$z$噪声增加；属于试点而非连续极限 \\
    \bottomrule
  \end{tabularx}
\end{table}

表\ref{tab:b04-lattice-evidence}中的所有数字均来自来源文档的二次汇总。其作用是把
证据定位到具体系综和方法，而不是替代原始论文、组态检查或独立统计复核。专著若把
这些结果升级为“已由本库验证”，必须另行审计原始数据与代码。


\clearpage
\chapter{胶子算符、极化投影与重正化基底}
\label{sec:b05-gluon-operators}

本章把“胶子分布”拆成三个不能互相替代的层次：连续场论中的非局域算符、格点上的有限距离矩阵元，以及把格点量映射到光锥分布的方案链。胶子场属于伴随表示，场强又带两个反对称 Lorentz 指标，因此同一个物理分布在写进格点代码之前必须先固定色表示、指标缩并、Wilson 线方向与重正化基底。

\section{双胶子算符与 Lorentz 不变振幅}

\paragraph{定义.}
沿空间方向 \(z^\mu=z\hat z^\mu\) 连接两个场强的规范协变算符写成
\begin{equation}
  O_g^{\mu\nu\rho\sigma}(x;z)
  =F_a^{\mu\nu}(x+z\hat z)\,
   U_z^{ab}(x+z,x)\,
   F_b^{\rho\sigma}(x),
  \label{eq:b05-gluon-operator}
\end{equation}
其中 \(U_z\) 是基础表示链接构造的平行移动；在矩阵元中通过
\begin{equation}
  U_z^{ab}(x+z,x)
  =2\,\operatorname{Tr}\!\left[T^aU_z(x+z,x)T^bU_z(x,x+z)\right]
  \label{eq:b05-adjoint-link}
\end{equation}
解释其伴随指标。每个场强可选取 \(6\) 个独立 \(( \mu,\nu)\) 分量，因此双线性一般需要 \(6\times6=36\) 个 Lorentz 结构来审计。

\paragraph{公式：六个不变振幅.}
对自旋平均的外部态，定义矩阵元
\begin{equation}
  M_{\mu\alpha;\lambda\beta}(z,p)
  =\langle p|G_{\mu\alpha}(z)W[z,0]G_{\lambda\beta}(0)|p\rangle .
  \label{eq:b05-bare-matrix}
\end{equation}
Lorentz 协变性把它分解为
\begin{align}
  M_{\mu\alpha;\lambda\beta}
  &=(g_{\mu\lambda}p_\alpha p_\beta
    -g_{\mu\beta}p_\alpha p_\lambda
    -g_{\alpha\lambda}p_\mu p_\beta
    +g_{\alpha\beta}p_\mu p_\lambda)M_{pp}\nonumber\\
  &\quad+(g_{\mu\lambda}z_\alpha z_\beta
    -g_{\mu\beta}z_\alpha z_\lambda
    -g_{\alpha\lambda}z_\mu z_\beta
    +g_{\alpha\beta}z_\mu z_\lambda)M_{zz}\nonumber\\
  &\quad+(g_{\mu\lambda}z_\alpha p_\beta
    -g_{\mu\beta}z_\alpha p_\lambda
    -g_{\alpha\lambda}z_\mu p_\beta
    +g_{\alpha\beta}z_\mu p_\lambda)M_{zp}\nonumber\\
  &\quad+(g_{\mu\lambda}p_\alpha z_\beta
    -g_{\mu\beta}p_\alpha z_\lambda
    -g_{\alpha\lambda}p_\mu z_\beta
    +g_{\alpha\beta}p_\mu z_\lambda)M_{pz}\nonumber\\
  &\quad+(p_\mu z_\alpha-p_\alpha z_\mu)
         (p_\lambda z_\beta-p_\beta z_\lambda)M_{ppzz}\nonumber\\
  &\quad+(g_{\mu\lambda}g_{\alpha\beta}
    -g_{\mu\beta}g_{\alpha\lambda})M_{gg}.
  \label{eq:b05-six-amplitudes}
\end{align}
六个振幅依赖于 \(z^2\) 和 Ioffe 时间 \(\nu=-(p\cdot z)\)。

时间反演与场交换使 \(M_{pp},M_{zz},M_{gg},M_{ppzz}\) 以及
\(M_{pz}-M_{zp}\) 为 \(\nu\) 的偶函数，而 \(M_{pz}+M_{zp}\) 为奇函数。
这个奇偶结构是后面用 \(z\leftrightarrow -z\) 分离极化信号的场论根据。

\paragraph{公式：非极化投影.}
在光锥极限 \(z=z^-\hat n_-\) 下，最低扭度胶子分布来自 \(M_{pp}\)：
\begin{equation}
  -M_{pp}(\nu,z^2=0)
  =\frac12\int_{-1}^{1}dx\,e^{-ix\nu}x f_g(x),
  \qquad f_g(x)=g(x).
  \label{eq:b05-mpp-to-pdf}
\end{equation}
注意被积函数含 \(x\)：这里得到的是动量密度 \(xg(x)\) 的 Fourier 变换，而不是 \(g(x)\) 本身。实验和全局拟合若使用不同的 \(xg(x)\) 约定，必须统一后才能比较。

\section{非极化、螺旋度与 Ioffe-时间投影}

格点上不能直接令 \(z^2=0\)，只能选取一组在同圈具有相同反常量纲的等时分量。常用的非极化组合为
\begin{equation}
  \mathcal O_{\mathrm{unpol}}(z)
  =M_{tx;tx}(z)+M_{ty;ty}(z)-2M_{xy;xy}(z).
  \label{eq:b05-unpol-combination}
\end{equation}
在部分有效场论讨论中可写成
\begin{equation}
  M_{ti;it}+M_{ij;ji}=2p_0^2M_{pp}+M_{\mathrm{tr}}+\cdots ,
  \label{eq:b05-pp-combination}
\end{equation}
其中 \(M_{\mathrm{tr}}\) 概括 \(M_{gg}\)、迹缩并和高扭度接触项。\(M_{ij;ji}\) 并不能在所有约定下被当成 \(M_{pp}\) 的无污染替代；式 \eqref{eq:b05-pp-combination} 的余项必须在算子基底和匹配方案中保留。

\paragraph{极化投影.}
引入对偶场强
\begin{equation}
  \widetilde G^{\mu\nu}=\frac12\epsilon^{\mu\nu\rho\sigma}G_{\rho\sigma},
  \qquad \epsilon^{0123}=+1,
  \label{eq:b05-dual-field}
\end{equation}
并定义极化矩阵元
\begin{equation}
  \widetilde m_{\mu\alpha;\lambda\beta}(z,p,s)
  =\langle p,s|G_{\mu\alpha}(z)W[z,0]\widetilde G_{\lambda\beta}(0)|p,s\rangle .
  \label{eq:b05-polarized-matrix}
\end{equation}
自旋无关的背景由偶部进入；螺旋度相关部分用
\begin{equation}
  \widetilde M_{\mu\alpha;\lambda\beta}(z,p,s)
  =\widetilde m_{\mu\alpha;\lambda\beta}(z,p,s)
  -\widetilde m_{\mu\alpha;\lambda\beta}(-z,p,s)
  \label{eq:b05-zodd}
\end{equation}
提取。格点上应同时计算 \(+z\) 与 \(-z\)，并保持链接路径反向、端点场强顺序和组态样本完全相同。

极化 Ioffe-时间分布由
\begin{equation}
  -i\widetilde I_p(\nu,z^2=0)
  =\widetilde M_{(ps)}^{(+)}(\nu,0)-\nu\widetilde M_{(pp)}(\nu,0)
  =\frac{i}{2}\int_{-1}^{1}dx\,e^{-ix\nu}x\Delta g(x)
  \label{eq:b05-polarized-ioffe}
\end{equation}
给出。线性组合中减去 \(\nu\widetilde M_{(pp)}\) 的作用，是消去由外部动量与自旋结构产生的非螺旋度振幅，而不是对数据作经验性减背景。

\paragraph{约化赝 ITD.}
令同一极化投影的坐标矩阵元为 \(\widetilde M(z,p)\)。定义
\begin{equation}
  \mathfrak M(\nu,z^2)
  =\frac{\widetilde M(z,p)}
         {\widetilde M(0,p)},
  \label{eq:b05-reduced-itd}
\end{equation}
并与微扰核 \(K\) 匹配：
\begin{equation}
  \mathfrak M(\nu,z^2)
  =\int_{0}^{1}du\,K(u\nu,\mu^2z^2)I_p(u\nu,\mu)
  +\mathcal O(z^2\Lambda_{\mathrm{QCD}}^2).
  \label{eq:b05-pseudo-factorization}
\end{equation}
比值可消去只依赖 \(z\) 的乘性 UV 因子，但不能消去硬核、软函数、rapidity 方案或大 \(z\) 幂修正。因此“约化后无显式 \(Z\)”不等于“无需重正化论证”。

\begin{table}[H]
  \centering
  \caption{胶子投影的用途与主要污染。}
  \label{tab:b05-projections}
  \small
  \begin{tabularx}{\textwidth}{@{}p{0.23\textwidth}p{0.25\textwidth}X@{}}
    \toprule
    投影 & 目标量 & 必须报告的余项或几何 \\
    \midrule
    \(M_{tx;tx}+M_{ty;ty}-2M_{xy;xy}\) & 非极化 \(xg(x)\) & \(M_{gg}\)、迹项、Wilson 方向、匹配核 \\
    \(\widetilde M_{0i;0i}+\widetilde M_{ij;ij}\) & 螺旋度 \(\Delta g(x)\) & 对偶张量符号、\(z\)-奇组合、核子自旋投影 \\
    约化矩阵元 \(\mathfrak M\) & 赝 ITD & 归一化点、短距窗口、\(z^2\) 幂修正 \\
    \(F W F\) 的任意单项 & 诊断量 & 一般不是 twist-2 的本征算符，不能单独发布 PDF \\
    \bottomrule
  \end{tabularx}
\end{table}

\section{可乘性重正化、三角混合与构建块}

\paragraph{严格的纯准算符陈述.}
在规范不变的 UV 正规化中，纯准胶子算符的端点和 Wilson 线自能可由有限个因子组织。按四个场强指标中含 \(z\) 方向的个数 \(s\) 分类，可把结构写成
\begin{equation}
  O_{g,R}^{\mu\nu\rho\sigma}(z)
  =e^{-\delta m|z|}Z_{\mathrm{op}}^{-1}(s)
   O_{g,B}^{\mu\nu\rho\sigma}(z),
  \label{eq:b05-multiplicative}
\end{equation}
\begin{equation}
  Z_{\mathrm{op}}(s)
  =Z_{wg}Z_{vg1}^{s/2}Z_{vg2}^{(2-s)/2}.
  \label{eq:b05-zop}
\end{equation}
这里 \(s=0,1,2\)；\(\delta m\) 概括格点正规化下的线性自能，\(Z_{wg}\) 是 Wilson 线端点因子，\(Z_{vg1},Z_{vg2}\) 区分含 \(z\) 指标的分量。由于整体 \(Z\) 依赖 \(s\)，任意改变各分量权重的线性组合未必保持可乘性。

\paragraph{辅助场闭包中的三角混合.}
为了系统识别小基底内的重定义，可引入伴随表示辅助场
\begin{equation}
  \mathcal L_{\mathrm{aux}}
  =\mathcal L_{\mathrm{QCD}}
  +\overline Q_a(x)(in\cdot D-m)_{ab}Q_b(x).
  \label{eq:b05-auxiliary-lagrangian}
\end{equation}
在固定规范与维数正规化中，与场强--辅助场双线性可能重定义的有三组结构：
\begin{equation}
  \begin{pmatrix}J_{1,R}^{\mu\nu}\\J_{2,R}^{\mu\nu}\\J_{3,R}^{\mu\nu}\end{pmatrix}
  =
  \begin{pmatrix}
    Z_{11}&Z_{12}&Z_{13}\\
    0&Z_{22}&Z_{23}\\
    0&0&Z_{33}
  \end{pmatrix}
  \begin{pmatrix}J_1^{\mu\nu}\\J_2^{\mu\nu}\\J_3^{\mu\nu}\end{pmatrix}.
  \label{eq:b05-triangular-mixing}
\end{equation}
其中 \(J_2\) 是场强与辅助场组成的规范不变结构，\(J_3\) 比例于辅助场运动方程。物理矩阵元中 \(J_3\) 的接触项可消失，但在检查格点算子基底时不能把它的存在从代数分类中删掉。

\paragraph{四个构建块.}
积分掉辅助场后，可作为独立重正化单元的结构为
\begin{align}
  \mathcal O_R^1
  &=Z_{11}^2e^{\delta m\Delta z}
    F^{ti}W_{z}F^{ti},\label{eq:b05-block1}\\
  \mathcal O_R^2
  &=Z_{22}^2e^{\delta m\Delta z}
    F^{zi}W_{z}F^{zi},\label{eq:b05-block2}\\
  \mathcal O_R^3
  &=Z_{11}Z_{22}e^{\delta m\Delta z}
    F^{ti}W_{z}F^{zi},\label{eq:b05-block3}\\
  \mathcal O_R^4
  &=Z_{22}^2e^{\delta m\Delta z}
    F^{z\mu}W_{z}F^{z}_{\ \mu},
  \label{eq:b05-block4}
\end{align}
其中横向指标 \(i\) 需按声明的方式缩并，\(\Delta z\) 是两端纵向坐标差。指数因子表达的是自能结构；实际格点拟合仍要用多格距和同几何 Wilson 圈检验。

\begin{table}[H]
  \centering
  \caption{胶子算符重正化中两类“混合”的区别。}
  \label{tab:b05-two-mixings}
  \small
  \begin{tabularx}{\textwidth}{@{}p{0.22\textwidth}p{0.28\textwidth}X@{}}
    \toprule
    层次 & 数学对象 & 物理含义与边界 \\
    \midrule
    UV 算符基底 & \(O_{i,R}=Z_{ij}O_j\) & 同量子数、同量纲的格点算符重定义；决定能否用标量因子或必须用矩阵 \\
    辅助场闭包 & \(J_{1,2,3}\) 三角矩阵 & 用于分类端点和 Wilson 线结构，不表示任意物理投影都相互混合 \\
    味单态匹配 & \(C_{gg},C_{gq},C_{qg},C_{qq}\) & 准分布匹配到光锥 PDF 时的胶子--夸克卷积，不是 UV 自能 \\
    投影误差 & \(M_{\mathrm{tr}}\)、高扭度、接触项 & 由选错 Lorentz 组合或省略迹项产生，不能被重正化因子吸收 \\
    \bottomrule
  \end{tabularx}
\end{table}

一个安全的结论是：在指定纯准算符基底和规范不变正规化下，36 个分量可按 \(s\) 分类处理；在包含辅助场或不同物理投影的更大闭包中，仍可能出现三角或完整矩阵混合。把这两句话合并成“36 个分量在所有基底下都无混合”是错误的。

\section{Clover 场强、Wilson 链接与非连通矩阵元}

格点上用四个定向 plaquette 的 Clover 平均构造场强。取围绕同一点的定向 plaquette \(P_{\mu\nu}\)，定义
\begin{equation}
  Q_{\mu\nu}=P_{\mu\nu}+P_{\nu,-\mu}+P_{-\mu,-\nu}+P_{-\nu,\mu},
  \label{eq:b05-clover-sum}
\end{equation}
并以反厄米部分构造
\begin{equation}
  F_{\mu\nu}^{\mathrm{clov}}
  =-\frac{i}{8a^2}\left(Q_{\mu\nu}-Q_{\mu\nu}^{\dagger}\right).
  \label{eq:b05-clover-field}
\end{equation}
不同快照对四个 plaquette 的书写顺序和反向链接定向不完全一致；代码审计应以“同起点、同闭合方向、反厄米”作为归一化条件，而不能只看某一项的显式正负号。

基础表示的有限距离矩阵元写成
\begin{equation}
  M_{\mu\lambda;\nu\rho}(z)
  =\sum_{\vec x}\operatorname{Tr}\!\left[
    F_{\mu\lambda}(x+z\hat z)U_z(x+z,x)
    F_{\nu\rho}(x)U_z(x,x+z)\right],
  \label{eq:b05-lattice-correlator}
\end{equation}
其中第二条链接是返回路径。非极化和螺旋度流分别为
\begin{align}
  \mathcal O_{\mathrm{unpol}}^{(0)}(z)
  &=\left[-F^{0i}W_zF_{0i}+2F^{12}W_zF_{12}\right]_{\mathrm{E}} ,
  \label{eq:b05-lattice-unpol}\\
  \mathcal O_{\mathrm{hel}}^{(0)}(z)
  &=-\Bigl\{F^{01}W_zF^{23}-F^{02}W_zF^{13}
    +2F^{12}W_zF^{03}-(z\leftrightarrow -z)\Bigr\}_{\mathrm{E}} .
  \label{eq:b05-lattice-hel}
\end{align}
第一式中的相对负号来自时间型与空间型场强的 Wick 转动；第二式中的 \(z\)-奇反对称化是螺旋度投影，不是普通的统计对称化。

\paragraph{非连通三点函数.}
胶子流是纯规范算符，在 Wick 收缩后与核子夸克线拓扑分离。逐组态计算时使用
\begin{equation}
  C_{3}(z,P_z;t_s,t_i)
  =\left\langle
    \left(\mathcal O_g(z,t_i)-\langle\mathcal O_g(z,t_i)\rangle\right)
    \left(C_2(P_z;t_s,t_0)-\langle C_2(P_z;t_s,t_0)\rangle\right)
  \right\rangle ,
  \label{eq:b05-disconnected}
\end{equation}
再以 \(C_3/C_2\) 和源汇间隔拟合提取基态矩阵元。真空减除必须在组态层面完成；先对乘积求平均再减去两个均值会保留交叉项的偏置。

\paragraph{验收条件.}
本节的算符构造只有在以下条件同时满足时才可进入匹配阶段：色表示为伴随或显式转换为基础；场强为反厄米且四叶定向一致；Wilson 线方向、端点、长度在裸量和减除量中相同；投影含迹项处理；\(+z\) 与 \(-z\) 数据使用同一源汇与统计样本；非连通真空减除可复现。

\chapter{TMD、准 TMD 与软结构}
\label{sec:b05-tmd}

TMD 不是“在 PDF 上多写一个 \(k_T\)”。它把纵向动量、横向分离、UV 重正化标度和 rapidity 标度组织成同一个因子化对象。格点计算只能直接得到 Euclidean 矩阵元；软函数、CS 核、硬匹配和极限外推必须逐层闭合。

\section{定义、共线极限与 TMD 分类}

在 \(b_T\) 空间定义带 Dirac 投影的裸关联函数
\begin{equation}
  h_{\Gamma}(z,b_T,L,P_z;1/a)
  =\langle P|\bar\psi(0)\Gamma
    W_{\sqsubset}(0,z;b_T)\psi(z,b_T)|P\rangle ,
  \label{eq:b05-bare-tmd}
\end{equation}
其中 \(W_{\sqsubset}\) 是有限长度 staple，\(L\) 是纵向臂长。把 UV 减除后的矩阵元记为 \(h_R\)，准 TMD 的 Fourier 变换写成
\begin{equation}
  \widetilde f_{\Gamma}(x,b_T,P_z)
  =\mathcal N^{-1}\int\frac{dz}{2\pi}\,
    e^{-ixzP_z}h_R(z,b_T,P_z),
  \label{eq:b05-quasi-tmd-fourier}
\end{equation}
\(\mathcal N\) 固定场强归一化和 Fourier 约定。坐标空间截断、大 \(z\) 尾部和 \(x\) 空间非物理区都属于同一反演问题，不能用同一条光滑曲线同时“拟合掉”。

\paragraph{共线极限.}
在 leading twist 且 rapidity 方案固定时，
\begin{equation}
  f_{g/h}(x,\mu)
  =\int d^2k_T\,f_1^g(x,k_T;\mu,\zeta),
  \label{eq:b05-collinear-limit}
  \end{equation}
并满足总动量求和
\begin{equation}
  \int_0^1dx\,x\left[f_g(x,\mu)+\sum_qf_q(x,\mu)\right]=1.
  \label{eq:b05-momentum-sum}
\end{equation}
这两个关系是方案一致性检验。横向积分截断、单态混合、零模减除或重正化误差都可能让数值不满足式 \eqref{eq:b05-collinear-limit}，不能只把它们归入“统计误差”。

\begin{table}[H]
  \centering
  \caption{八种 leading-twist quark TMD 的分类与格点状态。表中状态是 B05 文档所述先例层级，不代表多格距物理闭合。}
  \label{tab:b05-quark-tmd}
  \footnotesize
  \begin{tabularx}{\textwidth}{@{}p{0.18\textwidth}p{0.11\textwidth}p{0.10\textwidth}X@{}}
    \toprule
    函数 & 核子/夸克极化 & T 性质 & 计算状态与关键风险 \\
    \midrule
    \(f_1\) & U/U & 偶 & 核子单系综先例；连续极限与软化未闭合 \\
    \(g_{1L}\) & L/L & 偶 & 需极化核子态；完整胶子对应尚未完成 \\
    \(g_{1T}\) & T/L & 偶 & worm-gear；自旋投影与非连通噪声 \\
    \(h_1\) & T/T & 偶 & 手征奇数；格点混合与端点行为 \\
    \(h_{1L}^{\perp}\) & L/T & 偶 & worm-gear L；需横向极化投影 \\
    \(h_{1T}^{\perp}\) & T/T & 偶 & pretzelosity；高扭度污染风险 \\
    \(f_{1T}^{\perp}\) & T/U & 奇 & Sivers；依赖 Wilson 线方向，仍属探索 \\
    \(h_1^{\perp}\) & U/T & 奇 & Boer--Mulders；核子与介子已有先例，胶子对应未完成 \\
    \bottomrule
  \end{tabularx}
\end{table}

胶子 TMD 不能把表 \ref{tab:b05-quark-tmd} 中的 Gamma 矩阵替换为场强后直接复用。胶子 Wilson 线在伴随表示，场强投影有无迹、对偶和线性极化结构，T-odd 与过程方向的关系也不同。

\section{Staple 几何、四类发散与 Wilson 圈减除}

一个常用的 Euclidean staple 路径为
\begin{equation}
  \Gamma:\quad 0\longrightarrow L\hat z
  \longrightarrow L\hat z+b_T
  \longrightarrow z\hat z+b_T .
  \label{eq:b05-staple-path}
\end{equation}
路径方向和长度是算符定义的一部分。非极化 leading-twist 分布在某些方案中可对方向差异重吸收，但 T-odd、自旋相关和软函数矩阵元不能把 \(+L\) 与 \(-L\) 的裸量混为一谈。

\begin{table}[H]
  \centering
  \caption{staple 关联函数的主要发散与错误处置。}
  \label{tab:b05-staple-divergences}
  \small
  \begin{tabularx}{\textwidth}{@{}p{0.20\textwidth}p{0.25\textwidth}X@{}}
    \toprule
    结构 & 典型形式 & 处置与常见错误 \\
    \midrule
    线性自能 & \(e^{-\delta m(2L+z)/a}\) & 用同长度 Wilson 圈平方根或质量重正化；不同表示不能互借系数 \\
    Cusp & 由 cusp 反常维数控制的对数 & 转角和表示必须同几何；不能用普适常数替代 \\
    Pinch pole & \(e^{-V(b_T)L}\) & 由两条平行长 Wilson 线交换产生；需要 \(L\) 平台或同几何圈减除 \\
    端点对数 & \(\ln a\) 型复合算符发散 & 由短距比率或匹配吸收；普通 RI/MOM 单除不够 \\
    Rapidity & 依赖光锥方向和软胶子 & 属于软函数/rapidity 方案，不是普通 UV 发散 \\
    \bottomrule
  \end{tabularx}
\end{table}

设矩形 Wilson 圈为 \(Z_E(2L+z,b_T;1/a)\)，其渐近结构可概括为
\begin{equation}
  \ln Z_E
  =-2\delta m L-2V(b_T)L
   -\Gamma_{\mathrm{cusp}}\ln(\mu L)
   +C_E(b_T;\mu)+\cdots .
  \label{eq:b05-wilson-loop-asymptotic}
\end{equation}
若开 Wilson 线的线性自能与 pinch 项按总长度增长，则
\begin{equation}
  h^{\mathrm{sub}}(z,b_T,P_z;1/a)
  =\lim_{L\to\infty}
   \frac{h^{\mathrm{bare}}(z,b_T,L,P_z;1/a)}
        {\sqrt{Z_E(2L+z,b_T;1/a)}}
  \label{eq:b05-wilson-loop-subtraction}
\end{equation}
可以抵消主项。式 \eqref{eq:b05-wilson-loop-asymptotic} 是拟合与诊断模板，不是对所有 \(L\) 的严格等式；只有同表示、同转角、同端点和同横向分离的数据才能进入同一平台拟合。

\paragraph{短距比率.}
选取短距参考点 \((z_0,b_{T,0})\)，把剩余端点对数匹配到 \(\overline{\mathrm{MS}}\)：
\begin{equation}
  h^{\mathrm{SDR}}(z,b_T,P_z;\mu)
  =\frac{h^{\mathrm{sub}}(z,b_T,P_z;1/a)}
         {h^{\mathrm{sub}}(z_0,b_{T,0},0;1/a)}
   \frac{h^{\overline{\mathrm{MS}}}(z_0,b_{T,0};\mu)}
        {h^{\overline{\mathrm{MS}}}(z,b_T;\mu)}.
  \label{eq:b05-sdr}
\end{equation}
参考点必须同时满足 \(a\ll z_0,b_{T,0}\ll 1/\Lambda_{\mathrm{QCD}}\)。固定参考点在不同格距上对应不同的物理距离会产生离散化偏差；固定物理参考点又会改变微扰精度，因此参考点误差必须通过扫描或匹配转换显式估计。

\paragraph{RI/MOM 的边界.}
RI/MOM 可以确定某些局域反常量纲和混合矩阵，却不自动消除非局域 Wilson 线的线性自能、pinch pole 或 cusp 结构。B05 文档记录的五格距检验表明，单独采用 RI/MOM 后仍有随 \(z\) 或 \(a\) 增强的格距依赖。可审计结论是“RI/MOM 单独不足，需要与同几何减除联合”，而不是“所有 TMD 只能采用一种方案”。

\section{软函数、Collins--Soper 演化与匹配}

表示 \(R\) 的软函数可写成
\begin{equation}
  S_R(b_T;\mu,Y_1,Y_2)
  =S_{I,R}(b_T;\mu)
   \exp\!\left[K_R(b_T;\mu)\mathcal Y(Y_1,Y_2)\right],
  \label{eq:b05-soft-factorization}
\end{equation}
其中 \(S_I\) 是 rapidity 无关的内禀部分，\(K_R\) 是 Collins--Soper 核，\(\mathcal Y\) 依赖正规化方案。对胶子，\(R=\mathrm{adj}\)，领头 cusp 系数与 \(C_A\) 成正比；quark 的 \(C_F\) 数值不能直接替换。

TMD 的 rapidity 演化和 CS 核的 UV 演化分别为
\begin{equation}
  2\zeta\frac{\partial}{\partial\zeta}\ln f^R(x,b_T;\mu,\zeta)
  =K_R(b_T;\mu),
  \qquad
  \mu^2\frac{\partial}{\partial\mu^2}K_R(b_T;\mu)
  =-\Gamma_{\mathrm{cusp}}^R(\alpha_s).
  \label{eq:b05-cs-evolution}
\end{equation}
软函数在因子化式中以平方根进入，是因为一条 TMD Wilson 线与软函数的两条 Wilson 线几何不同。该平方根、符号和常数项都属于方案，必须与匹配核同步。

\paragraph{结构因子化式.}
准 TMD 到光锥 TMD 的关系可组织为
\begin{equation}
  \begin{aligned}
  \frac{\widetilde f(x,b_T;\zeta_z,\mu)}
       {\sqrt{S_I(b_T;\mu)}}
  =H\!\left(\frac{\zeta_z}{\mu^2}\right)
   \exp\!\left[\frac12\ln\!\left(\frac{\zeta_z}{\zeta}\right)K(b_T;\mu)\right]
   f(x,b_T;\mu,\zeta)\\
  +\mathcal O\!\left(\frac{\Lambda_{\mathrm{QCD}}^2}{\zeta_z},
   \frac{M^2}{P_z^2},\frac{1}{b_T^2\zeta_z}\right),
  \end{aligned}
  \label{eq:b05-tmd-factorization}
\end{equation}
其中 \(\zeta_z\) 由准 TMD 的动量标度确定。不同文献可能使用 \(H\)、卷积或归一化不同的写法；在数值合并前必须统一到同一 Fourier 变换和场强归一化。

\paragraph{硬匹配与大对数.}
固定阶硬核含 \(\ln(\zeta_z/\mu^2)\) 的幂次。在 \(x\) 很小或动量很大时，\(\zeta_z\) 与 \(\mu\) 的比值可能远离 1，固定阶展开失效。重整化群重求和把硬核从自然标度 \(\mu_0\sim 2xP_z\) 演化到参考标度 \(\mu\)。当 \(\mu_0\) 进入 Landau 区时，有限 \(P_z\) 无法直接给出更小 \(x\) 的可控微扰匹配，必须报告 \(x_{\min}\) 和模型/外推误差。

\section{格点 TMD 的证据等级}

\begin{table}[H]
  \centering
  \caption{quark TMD 方法先例与直接 gluon TMD 的边界。}
  \label{tab:b05-tmd-evidence}
  \small
  \begin{tabularx}{\textwidth}{@{}p{0.22\textwidth}p{0.25\textwidth}X@{}}
    \toprule
    证据类型 & 已记录内容 & 可支持与不可支持的结论 \\
    \midrule
    Wilson 圈减除 & 五种格距、有限 \(L\) 平台、不同费米子作用量 & 支持 \(L\to\infty\) 自能减除；不支持完整物理 TMD 已闭合 \\
    \(\pi\) 介子软结构 & 大动量形状因子抽取 \(S_I\)，不同动量比抽取 \(K\) & 支持软核可计算及系综交叉检验；gluon 表示仍需独立匹配 \\
    核子 \(f_1\) & 单系综、NNLO+RGR、核子 TMD 分布先例 & 支持算法链；连续极限和物理体积尚未由同一数据闭合 \\
    Boer--Mulders & 核子与 pion 的手征奇数 T-odd 先例 & 支持 T-odd 方法；不能外推为 gluon T-odd 数值 \\
    直接 gluon TMD & B05 内未见中心值、协方差和多格距组合 & 当前结论只能是框架与工作流，不能发布物理分布 \\
    \bottomrule
  \end{tabularx}
\end{table}

把表 \ref{tab:b05-tmd-evidence} 的不同行相乘、相除或平均会产生错误证据：Wilson 圈五格距验证不能替代核子连续极限；quark 软核不能替代伴随表示软核；单系综 TMD 先例不能替代多格距胶子结果。

\section{梯度流方案与流时验收}

连续梯度流定义为
\begin{equation}
  \partial_tB_\mu(t,x)=D_\nu G_{\nu\mu}(t,x),
  \qquad B_\mu(0,x)=A_\mu(x),
  \label{eq:b05-flow-equation}
\end{equation}
格点上以 Wilson plaquette 作用量离散化。定义流半径
\begin{equation}
  r_f=\sqrt{8t},
  \label{eq:b05-flow-radius}
\end{equation}
并把 \(r_f\) 视为平滑尺度，而不是一个额外的普通重正化标度。固定正流时下，流场规范不变局域关联函数的 UV 行为可由标准 QCD 参数吸收；这不表示任何 Euclidean 非局域矩阵元自动等于光锥物理。

\paragraph{小流时展开.}
短流时下，流化算符按规范不变局域算符展开：
\begin{equation}
  O_R(t,x)
  =\sum_i c_i(t)O_{R,i}(x)
   +\mathcal O(t\Lambda_{\mathrm{QCD}}^2).
  \label{eq:b05-sftx}
\end{equation}
对非局域 TMD，Wilson 系数还必须依赖 \(z,b_T,P_z\)；固定外部物理分离后再对 \(t\) 展开的要求为
\begin{equation}
  a\ll r_f\ll\min\left(|z|,b_T,\Lambda_{\mathrm{QCD}}^{-1}\right),
  \qquad
  P_zr_f\ll1.
  \label{eq:b05-flow-window}
\end{equation}
如果 \(r_f\) 已经接近或超过 \(b_T\)，流会抹平待测的横向结构；如果 \(P_zr_f\) 不小，boost 会把流半径放大成硬尺度；如果 \(r_f\) 太小，格距效应和 SFTX 高阶项又不受控。

\paragraph{flow 不替代的对象.}
梯度流可改善 UV 行为和统计噪声，但不替代 rapidity 重正化、软函数、CS 核、Wilson 线方向或味单态匹配。一个经流化后有限、跨组态稳定的矩阵元仍可能因违反式 \eqref{eq:b05-flow-window}、缺少软减除或使用错误表示而只能作为诊断量。

\paragraph{验证协议.}
至少应做到：扫描多个满足窗口的 \(t\)；在每个 \(t\) 上重新完成 Wilson 圈、场强和准 TMD；用同一匹配核转换到共同方案；对 \(t\to0\) 作受控外推；比较不同外推窗口和匹配阶数；把流时误差、格距、体积、统计和匹配误差分开报告。单一流时的有限结果不足以支持物理分布的中心值。

\chapter{质子自旋、胶子极化与连续极限}
\label{sec:b05-spin-continuum}

自旋问题的困难不在“是否还有别的贡献”，而在不同分解、标度和算符定义下“贡献”指什么。连续极限的困难也不在“多拟合几个参数”，而在格距、动量、距离、流半径和匹配标度之间是否存在非空尺度窗口。

\section{自旋预算与两套分解}

\paragraph{历史量级.}
朴素非相对论夸克模型把质子写成三个 \(S\)-波夸克，容易给出 \(\Delta\Sigma\sim1\)。Ellis--Jaffe 利用 SU(3) 味对称和超子衰变给出
\begin{equation}
  \Delta\Sigma=3F-D\approx0.60\pm0.12,
  \label{eq:b05-ellis-jaffe}
\end{equation}
这已经表明“1”不是唯一自然预期。1987 年 EMC 的测量进一步给出很小的 \(\Delta\Sigma\)，从而把胶子和轨道角动量推到自旋预算的中心。历史数值依赖 QCD 标度、SU(3) 破缺和演化，不能把 EMC 的单一数值直接与现代全局分析逐项相减。

\paragraph{Jaffe--Manohar 分解.}
在无限动量或光锥规范语境中，
\begin{equation}
  \frac12
  =\frac12\Delta\Sigma+L_q+\Delta G+L_G,
  \label{eq:b05-jm}
\end{equation}
其中
\begin{equation}
  \Delta G(Q^2)=\int_0^1dx\,\Delta g(x,Q^2).
  \label{eq:b05-delta-g}
\end{equation}
这个分解能显式分离胶子螺旋度和胶子轨道项，但依赖坐标、规范与过程类。把 \eqref{eq:b05-jm} 中每一项当成与方案无关的局域算符矩阵元，会混淆不同定义。

\paragraph{Ji 分解.}
规范不变的角动量分解写成
\begin{equation}
  J=J_q+J_g,
  \qquad
  J_q=\frac12\Delta\Sigma+L_q,
  \label{eq:b05-ji}
\end{equation}
而 \(J_g\) 是胶子总角动量。\(J_g\) 不能在没有额外方案的情况下再被直接分为“胶子自旋”和“胶子轨道”两部分。两种分解中的 \(\frac12\Delta\Sigma\) 可对照，但 \(L_q,L_G,\Delta G,J_g\) 不能逐个等号连接。

\paragraph{轴反常的作用.}
单味单态轴矢流满足
\begin{equation}
  \partial^\mu A_\mu^0
  =N_f\frac{\alpha_s}{2\pi}\operatorname{Tr}F_{\mu\nu}\widetilde F^{\mu\nu}.
  \label{eq:b05-axial-anomaly}
\end{equation}
反常 Ward 恒等式可以在特定极限下把味单态轴荷与拓扑密度矩阵元联系起来，但反常流不能直接等同于胶子自旋密度。把 \(\Delta\Sigma\) 与 \(\Delta G\) 通过简单加法重复分配，是自旋预算中常见的错误来源。

\section{胶子螺旋度与赝 ITD 证据}

光锥螺旋度分布可由
\begin{equation}
  \Delta g(x,Q^2)
  =\frac{i}{2xP^+}\int\frac{d\xi^-}{2\pi}
    e^{-ixP^+\xi^-}
    \langle P,S|F_a^{+\alpha}(\xi^-)W^{ab}\widetilde F^+_{\alpha,b}(0)|P,S\rangle
  \label{eq:b05-helicity-pdf}
\end{equation}
定义。格点上的极化坐标矩阵元需先做 \(z\)-奇组合，再进入 Ioffe-时间或准分布。HadStruc 的赝 ITD 计算是一个单格距、较重 pion 质量的先例，表 \ref{tab:b05-spin-constraints} 所列数值只能作为来源文献快照。

\begin{table}[H]
  \centering
  \caption{B05 文档记录的 \(\Delta G\) 约束快照。数值未在本任务中独立重分析；比较前须统一标度、积分范围和正定性假设。}
  \label{tab:b05-spin-constraints}
  \small
  \begin{tabularx}{\textwidth}{@{}p{0.25\textwidth}p{0.22\textwidth}X@{}}
    \toprule
    方法 & 来源记录的中心值 & 边界与解读 \\
    \midrule
    DSSV14 & \(0.20^{+0.06}_{-0.07}\) & 全局分析快照，依赖数据集与参数化 \\
    NNPDFpol1.1 & \(0.13\pm0.14\) & 大误差，不能据此声明精确符号 \\
    JAM17 & \(0.06\pm0.21\) & 与零相容，显示约束仍弱 \\
    Ji-scheme 格点 & \(0.251\pm0.047\pm0.016\) & 局域算符方案，匹配系数系统误差较大 \\
    HadStruc 赝 ITD & 仅给出比较趋势 & 单格距探索，不提供完整 \(x\) 依赖或连续极限 \\
    JAM24 & 不要求正值 & 放松正定性后，正负号不能由旧共识外推 \\
    \bottomrule
  \end{tabularx}
\end{table}

对极化结果的验证应优先报告：核子自旋投影；对偶张量符号；\(z\)-奇组合；匹配核是否为 \(\Delta g\) 专用；单态胶子--夸克混合；赝 ITD 的归一化点和 \(z^2\) 范围；以及多格距、物理质量和体积是否闭合。

\section{非极化胶子 PDF 的三格距连续极限}

B05 连续极限主文档记录了三套 CLQCD \(2+1\) 味系综，算符为
\begin{equation}
  \mathcal O_g(z)=M_{tx;tx}+M_{ty;ty}-2M_{xy;xy}.
  \label{eq:b05-continuum-operator}
\end{equation}
Hybrid 方案在短距使用零动量比值，在长距使用自重整化因子：
\begin{equation}
  h_R(z,P_z)
  =\frac{h_B(z,P_z)}{h_B(z,0)}\theta(z_s-|z|)
   +\eta_s\frac{h_B(z,P_z)}{Z_R(z,1/a)}\theta(|z|-z_s).
  \label{eq:b05-hybrid}
\end{equation}
连续性系数 \(\eta_s\) 在 \(z_s\) 处固定重整化后的矩阵元。长距因子采用参数化
\begin{equation}
  Z_R=\exp\left[
    \frac{kz}{a}\ln(a\Lambda_{\mathrm{QCD}})+m_0z
    +\frac{5C_A}{3b_0}\ln\frac{\ln(1/a\Lambda_{\mathrm{QCD}})}
    {\ln(\mu/\Lambda_{\mathrm{QCD}})}
    +\frac12\ln\left(1+\frac{d}{\ln(a\Lambda_{\mathrm{QCD}})}\right)^2
  \right].
  \label{eq:b05-zr}
\end{equation}
这些参数强相关，单独看 \(k\) 或 \(m_0\) 的大误差会误导；应报告 \(Z_R(z)\) 本身及其对最终矩阵元的相关误差。

\begin{table}[H]
  \centering
  \caption{B05 连续极限文档记录的 CLQCD 系综。}
  \label{tab:b05-continuum-ensembles}
  \small
  \begin{tabularx}{\textwidth}{@{}lcccccX@{}}
    \toprule
    系综 & \(a/\mathrm{fm}\) & \(L^3\times T\) & \(m_\pi/\mathrm{MeV}\) & \(m_\pi L\) & \(N_{\mathrm{conf}}\) & 主要用途 \\
    \midrule
    C24P29 & 0.105 & \(24^3\times72\) & 292.3(1.0) & 3.746(13) & 879 & 粗格距与较大体积 \\
    E32P29 & 0.0897 & \(32^3\times64\) & 287.3(2.5) & 4.179(36) & 900 & 中间格距 \\
    F32P30 & 0.0775 & \(32^3\times96\) & 300.4(1.2) & 3.780(15) & 777 & 细格距与高动量检验 \\
    \bottomrule
  \end{tabularx}
\end{table}

表中三套系综的 pion 质量约为 300 MeV，并非物理 pion 点。连续外推不能自动替代手征外推；有限体积参数 \(m_\pi L\) 只覆盖约 3.7--4.2。

\paragraph{大距离外推.}
坐标空间矩阵元在 \(\lambda=zP_z\) 大时信噪比下降。B05 文档采用
\begin{equation}
  h_R(\lambda)=l_1\lambda^{-a_1}e^{-\lambda/\lambda_0}
  \label{eq:b05-lambda-tail}
\end{equation}
补全 Fourier 变换所需尾部。幂律项与 PDF 端点行为相关，指数项概括有限关联长度。该形式必须通过改变拟合区间检验；只用一组区间得到的光滑尾部不能作为独立验证。

\paragraph{联合极限.}
最终外推使用
\begin{equation}
  xg(x,P_z,a)
  =xg_0(x)+a^2f(x)+a^2P_z^2h(x)+\frac{d(x)}{P_z^2}.
  \label{eq:b05-joint-extrapolation}
\end{equation}
\(a^2f\) 是普通离散化误差，\(a^2P_z^2h\) 是 boost 与格距的交叉项，\(d/P_z^2\) 概括高扭度与靶质量修正。三项不能互相吸收；若只有一种 \(P_z\) 或一个格距，式 \eqref{eq:b05-joint-extrapolation} 的系数不可辨识。

\paragraph{结果解读.}
B05 文档报告最终 \(xg(x)/\langle x\rangle\) 在误差内与若干全局拟合一致，大 \(x\) 受压，\(x<0.2\) 与 \(x>0.8\) 被视为 LaMET 不可靠区。PDF 第 6 页的图只支持这一层级；任何具体点的中心值仍应回到原始数据表或发表表格，而不能从扫描图像估读。

\section{误差预算与尺度窗口}

\begin{table}[H]
  \centering
  \caption{胶子 PDF 连续极限的主要误差与可验证构造。}
  \label{tab:b05-continuum-errors}
  \small
  \begin{tabularx}{\textwidth}{@{}p{0.24\textwidth}p{0.25\textwidth}X@{}}
    \toprule
    误差来源 & B05 文档中的构造 & 验收判据 \\
    \midrule
    大 \(\lambda\) 尾部 & 改变拟合起点与终点 & 两组外推在联合误差内一致，且 Fourier 振铃受控 \\
    匹配与标度 & \(\mu=2\to4\ \mathrm{GeV}\) & 标度差进入系统误差，而非常数平移 \\
    连续/无穷动量 & 联合结果与最小格距、较高动量结果比较 & 差值不大于声明的外推误差 \\
    混合方案阈值 & \(z_s=0.3\to0.2\ \mathrm{fm}\) & 中心值变化与 \(z_s\) 相关性显式报告 \\
    统计与拟合 & bootstrap/jackknife、相关协方差 & 重采样整条分析链，不只重采样最后一点 \\
    手征与体积 & 多个 \(m_\pi\)、\(L\) & 本轮 CLQCD 数据只近似 300 MeV，不能宣称已闭合 \\
    \bottomrule
  \end{tabularx}
\end{table}

连续外推还应满足一组可操作窗口：
\begin{equation}
  a\ll |z|,\ b_T\ll \frac{L_s}{2},
  \qquad
  \Lambda_{\mathrm{QCD}}\ll P_z,\ \frac{1}{|z|},
  \qquad
  aP_z\ll1,
  \qquad
  m_\pi L\gtrsim4.
  \label{eq:b05-scale-windows}
\end{equation}
对 flow-TMD 再加式 \eqref{eq:b05-flow-window}。如果窗口为空，正确结果是报告“没有可发布区间”，而不是把窗口外的点纳入同一全局拟合。
\section{从中间量到物理结论的等级}

\begin{table}[H]
  \centering
  \caption{B05 结果的分级验收。}
  \label{tab:b05-claim-levels}
  \small
  \begin{tabularx}{\textwidth}{@{}p{0.21\textwidth}p{0.31\textwidth}X@{}}
    \toprule
    等级 & 已闭合内容 & 允许的表述 \\
    \midrule
    L1：裸格点量 & 场强、链接、矩阵元、真空减除、统计误差 & 可报告诊断量与关联函数 \\
    L2：方案量 & UV 减除、短距参考点、Wilson 圈、比值或 flow 匹配 & 可报告在明确方案下的有限矩阵元 \\
    L3：光锥分布 & soft/rapidity、CS、硬匹配、Fourier、大 \(z\) 尾部 & 可报告 TMD 或 PDF 候选分布及其适用范围 \\
    L4：物理闭合 & 多格距、大动量、手征、体积、单态混合与独立交叉检验 & 才可称为完整的第一性原理物理结果 \\
    \bottomrule
  \end{tabularx}
\end{table}

任何缺少一层的结果都不能用更高层的名称包装。例如，没有软减除的“准 TMD 矩阵元”不是 TMD；单流时的“有限关联函数”不是已匹配的 gluon TMD；单格距的“连续极限比较”不是连续极限。


\clearpage
\chapter{重正化与梯度流}

\section{发散分类}

胶子 TMD 关联函数的 UV 结构比局域胶子算符复杂。表~\ref{tab:divergences} 给出常见的发散类型及其物理来源。不同发散的减除必须使用同一 Wilson 线几何；否则可以出现“数值上看似抵消、实际上改变了 TMD 方案”的错误。

\begin{table}[htbp]
\centering
\caption{胶子 staple 关联函数中的主要 UV 与 rapidity 结构。}
\label{tab:divergences}
\footnotesize
\begin{tabularx}{\textwidth}{L{3.0cm}L{3.2cm}Y}
\toprule
结构 & 典型形式 & 物理来源与处理 \\
\midrule
Wilson 线自能 & \(e^{-\delta m\,L_\parallel/a}\) & 与路径长度和表示有关，需用同长度 Wilson 圈或质量重正化吸收 \\
cusp 发散 & \(\exp[-\Gamma_{\rm cusp}\ln(\mu L)]\) & 转角角度、表示和路径端点决定；不能用一个普适常数替代 \\
pinch pole & \(e^{-V(\bT)L_\parallel}\) & 两条长 Wilson 线之间的有效势；在 \(L_\parallel\to\infty\) 形成平台或需 Wilson 圈减除 \\
端点对数 & \(Z_O(\mu,1/a)\) & 场强与 Wilson 线连接处的局域反常量纲，可通过 RI/MOM、比率或匹配处理 \\
rapidity 发散 & \(f_{\rm bare}/\sqrt{S_R}\) & 与快度积分、软胶子和规范链方向有关，不能由普通 UV 重正化消除 \\
算符混合 & \(O_i=Z_{ij}O_j\) & 同量纲胶子算符、单态夸克算符和不同 Lorentz 投影可能混合 \\
\bottomrule
\end{tabularx}
\end{table}

\section{Wilson 圈减除与短距方案}

对长度为 \(L_\parallel\)、横向分离为 \(\bT\) 的 Wilson 圈取
\begin{equation}
  Z_{E,R}(\bT;L_\parallel;1/a)
  =\frac{1}{d_R}\left\langle
  \Tr_R\,U_{\square,R}(\bT;2L_\parallel)
  \right\rangle ,
  \label{eq:wilsonloop}
\end{equation}
其中 \(d_R\) 是表示维数。若开 Wilson 线的自能按总长度增长，则 \(\sqrt{Z_{E,R}}\) 可抵消相应的线性因子；Wilson 圈与 staple 具有相同的转角结构时，cusp 贡献也可部分抵消。该减除必须使用胶子的伴随表示，不能直接套用夸克的基本表示数值。

把 Wilson 圈的渐近结构写成结构式，可以同时看到三类必须区分的贡献：
\begin{equation}
  \ln Z_{E,R}(\bT;L_\parallel;1/a)
  =-2\delta m_R\,L_\parallel
  -2V_R(\bT)L_\parallel
  -\Gamma^{R}_{\rm cusp}\ln(\mu L_\parallel)
  +C_R(\bT;\mu)
  +\mathcal O(a/L_\parallel),
  \label{eq:wilsonloopasymptotics}
\end{equation}
其中第一项是线性自能，第二项是静态势导致的 pinch-pole 结构，第三项是 cusp 对数，\(C_R\) 概括有限端点项。实际计算不能把式~\eqref{eq:wilsonloopasymptotics} 当作严格定理，而应把它当作设计不同 \(L_\parallel\) 扫描和判据的物理模板：只有同一表示、同一转角和同一横向分离的数据才能放进同一拟合。

在消除线性项和 pinch-pole 之后，仍需处理剩余的对数 UV 发散。短距比率方案把非局域矩阵元除以零动量、短距离参考点上的矩阵元：
\begin{equation}
  h^{\rm ratio}_g(z,\bT;P_z)
  =
  \frac{h_g(z,\bT;P_z;1/a)}
  {h_g(z_0,\boldsymbol b_{T,0};0;1/a)}
  \frac{h_g^{\MSbar}(z_0,\boldsymbol b_{T,0};\mu)}
  {h_g^{\MSbar}(z,\bT;\mu)},
  \label{eq:ratioscheme}
\end{equation}
其中 \(z_0,\boldsymbol b_{T,0}\) 位于短距离微扰窗口。混合方案则在短距离使用比率、在长距离使用自重整化或质量重正化，以便同时控制 \(z\to0\) 与 \(a\to0\) 的非交换极限\cite{JiHybrid2021,Huo2021}。这些方案的共同要求是：短距离参考点、Wilson 线几何、场强投影和 Fourier 约定必须在同一计算中固定。

RI/MOM 方案可以确定某些局域反常量纲，但它不自动消除非局域 Wilson 线的线性发散，也不能自动处理 staple 的 pinch-pole。对于非手征费米子作用量，手征对称性破缺还可能引发额外的算符混合。因此，“RI/MOM 重正化后曲线随 \(z\) 变平”不是充分的物理判据；还必须检查连续极限、不同格距和不同几何下的稳定性\cite{Huo2021,Ishikawa2017}。

\section{辅助场、36 个分量与可乘性重正化}

胶子准算符的 UV 分类可以通过辅助伴随“重夸克”场系统化。引入
\begin{equation}
  \mathcal L_{\rm aux}
  =\mathcal L_{\rm QCD}
  +\overline Q_a(x)(in\cdot D-m)_{ab}Q_b(x),
  \label{eq:auxiliarylagrangian}
\end{equation}
其中 \(n^\mu\) 指向非局域分离方向。Wilson 线可以写成辅助场的两点函数，从而使非局域算符的端点和规范链自能转化为局域算符乘积的重正化问题\cite{ZhangGluonLaMET2019}。在维数正规化中，纯胶子准算符的 UV 发散局域在端点，并可按 Lorentz 指标中 \(z\) 分量的个数分类。对四个场强指标 \(\mu,\nu,\rho,\sigma\)，令 \(s\) 为其中等于 \(z\) 的指标数，则符号结构可概括为
\begin{equation}
  O^{\mu\nu\rho\sigma}_{g,\rm bare}
  =e^{-C_g|z|}
  Z_{wg}Z_{vg1}^{s/2}Z_{vg2}^{(2-s)/2}
  O^{\mu\nu\rho\sigma}_{g,\rm ren}
  +\text{contact/EOM},
  \label{eq:multiplicativegluon}
\end{equation}
其中 \(C_g\) 概括线性自能，\(Z_{wg}\) 是 Wilson 线端点因子，\(Z_{vg1},Z_{vg2}\) 依赖于 \(z\) 指标数。该结构说明：并非所有 Lorentz 组合都可用同一个标量 \(Z\) 重正化；但存在若干可乘性重正化的构建块，可用于构造准 PDF 或准 TMD 的物理投影\cite{LiMaQiu2019,ZhangGluonLaMET2019}。

对具有相同量子数的算符集合，可乘性重正化应写成矩阵而不是标量。令 \(\mathcal O_i^{R}\) 为第 \(i\) 个重正化算符，则
\begin{keyformula}{重正化矩阵与 RG 演化}
\begin{align}
  \mathcal O_i^{R}(\mu)
  &=Z_{ij}(\mu,\mu_0)\mathcal O_j^{R}(\mu_0),
  \label{eq:renormmatrix}\\
  \mu\frac{\dd}{\dd\mu}\mathcal O_i^{R}(\mu)
  &=\gamma_{ij}(\alpha_s)\mathcal O_j^{R}(\mu),
  \label{eq:renormrge}
\end{align}
其中重复指标求和。矩阵非对角元就是胶子-夸克或不同 Lorentz 投影之间的混合；若只保留对角元，必须有物理或对称性理由，并把它作为近似误差报告。
\end{keyformula}

对于物理 TMD staple，还需要把该局域分类与横向分离、长 Wilson 线夹角和软函数结合。可乘性重正化的“准 PDF 结论”不能被误写为“任意 staple 算符都可乘性重正化”。一个安全的表述是：在固定几何、固定表示和固定投影下，UV 发散可由有限个因子组织；但 pinch-pole 和 rapidity 结构仍需要专门减除。

\section{RI/MOM 条件及其适用边界}

RI/MOM 方案的出发点是对截肢 Green 函数施加离壳条件。若 \(\Gamma_O(p)\) 是算符 \(O\) 的截肢顶点，\(P_\Gamma\) 是投影算符，则一种典型定义为
\begin{equation}
  Z_O^{-1}(\mu_R,p_R)
  \left[
  \frac{P_\Gamma\Gamma_O(p)}
  {P_\Gamma\Gamma_O^{(0)}(p)}
  \right]_{p^2=-\mu_R^2,\;p_z=p_R^z}
  =1,
  \label{eq:ri-momcondition}
\end{equation}
其中 \(\Gamma_O^{(0)}\) 是树级顶点。对具有多个 Lorentz 结构的胶子算符，必须同时说明投影、色指标缩并和混合矩阵；单一标量条件通常不足以固定所有分量。

RI/MOM 的优点是系统、可非微扰计算，并且能与 \(\MSbar\) 通过微扰转换。它的局限在于：
\begin{itemize}
  \item 离壳 Green 函数包含规范依赖和 Goldstone 型 IR 结构，需检查规范固定和动量窗口；
  \item 非局域 Wilson 线的线性自能不是普通局域反常量纲，单靠 RI/MOM 除法不能证明消失；
  \item 物理矩阵元和离壳顶点对算子混合的投影可能不同；
  \item 对破坏手征对称性的作用量，标量-伪标量或张量型混合可能需要额外矩阵。
\end{itemize}
因此，RI/MOM 可以作为局域反常量纲和混合矩阵的工具，但对 staple 型 TMD，必须与 Wilson 圈、短距比率或自重整化联合使用，并在多个格距上做连续检验\cite{Huo2021,Zhang2022TMD}。

\section{Ratio、hybrid 与 self-renormalization}

Ratio 方案用零动量矩阵元构造比值：
\begin{equation}
  h_R^{\rm ratio}(z,P_z)
  =\frac{h_{\rm bare}(z,P_z;1/a)}
  {h_{\rm bare}(z,0;1/a)}
  \frac{h_{\rm bare}(0,0;1/a)}{h_{\rm bare}(0,P_z;1/a)} .
  \label{eq:ratiodetail}
\end{equation}
它的优点是消除了只依赖 \(z/a\) 的自能和部分端点因子；缺点是零动量分母在大 \(z\) 处可能进入非微扰区，且比值的局域极限会改变整体归一化。
更早的准 PDF 方案还常把非微扰重正化与短距比率组合成两步流程，以便同时控制端点混合和幂修正\cite{Alexandrou2017}。

Hybrid 方案把短距离和长距离分开：
\begin{equation}
  h_R^{\rm hyb}(z,P_z)=
  \begin{cases}
  \dfrac{h_{\rm bare}(0,0)}{h_{\rm bare}(0,P_z)}
  \dfrac{h_{\rm bare}(z,P_z)}{h_{\rm bare}(z,0)}, & |z|\le z_s,\\[1.0em]
  \dfrac{h_{\rm bare}(0,0)}{h_{\rm bare}(0,P_z)}
  \dfrac{h_{\rm bare}(z,P_z)}{h_{\rm bare}(z_s,0)}
  e^{(\delta m+m_0)(|z|-z_s)}, & |z|>z_s,
  \end{cases}
  \label{eq:hybriddetail}
\end{equation}
其中 \(z_s\) 是短距切换点，\(\delta m\) 是 Wilson 线线性发散系数，\(m_0\) 概括残余质量或重整子歧义。短距端可以用微扰 Wilson 系数匹配到 \(\MSbar\)，长距端则保留非微扰矩阵元的指数衰减。若 \(z_s\) 的选择改变最终结果超过统计误差，说明幂修正或微扰匹配尚未受控\cite{JiHybrid2021,GoodYaoLin2026}。

Self-renormalization 采用乘性比值
\begin{equation}
  h_R^{\rm self}(z,P_z)
  =\frac{h_{\rm bare}(z,P_z;1/a)}
  {\left[h_{\rm bare}(z_0,0;1/a)\right]^{z/z_0}},
  \label{eq:selfrenorm}
\end{equation}
或用零动量短距点 \(z_0\) 的拟合参数化线性项与对数项。它不依赖 Wilson 圈的显式测量，适合多格距联合分析；但其成功条件是多格距数据能区分 \(a^{-1}\) 线性和对数依赖。单格距时，不同 \(z_0\) 或参数化可能给出不可辨识的结果\cite{Huo2021}。

\section{规范固定、HYP 与平滑链接}

Wilson 线的信噪比和非微扰矩阵元对规范固定的精度敏感。对需要固定规范的计算，残差应远低于常用阈值，尤其在细格距和长路径上；规范固定误差若随路径长度累积，可能伪装成物理的 \(z\) 依赖。对规范不变实现，HYP、stout 或 Wilson flow 可以降低 UV 噪声，但会改变算符的 UV 结构，因此必须把平滑参数纳入重正化和匹配方案，而不能把不同平滑度的裸矩阵元直接比较\cite{ZhangGauge2024,Peardon2009,Bali2016}。

表~\ref{tab:schemematrix} 概括各方案的分工。一个完整 TMD 计算通常不是选一个方案，而是把 Wilson 圈、短距匹配、软/rapidity 和 flow 组合成一条有明确适用域的链。

\begin{table}[htbp]
\centering
\caption{胶子非局域算符重正化方案的适用域。}
\label{tab:schemematrix}
\footnotesize
\begin{tabularx}{\textwidth}{L{2.4cm}L{3.2cm}L{3.0cm}Y}
\toprule
方案 & 主要消除对象 & 优点 & 主要风险/缺口 \\
\midrule
Wilson 圈开方 & 线性自能、pinch-pole、部分 cusp & 同几何、物理直观 & 必须同表示/同路径；不能单独消除端点对数 \\
Ratio & \(z/a\) 型 UV、部分端点 & 实现简单，短距归一化明确 & 长距零模、端点归一化和幂修正 \\
Hybrid & 短距 ratio + 长距质量减除 & 兼顾微扰与非微扰 & \(z_s\) 依赖、\(\delta m+m_0\) 拟合和格距数 \\
Self-renorm & 线性项、重整子 & 多格距联合，非微扰 & 单格距不可辨识；参数化模型依赖 \\
RI/MOM & 局域反常量纲和混合矩阵 & 系统、可转换到 \(\MSbar\) & 非局域自能、规范依赖和混合投影 \\
Gradient flow & UV 平滑、短流时匹配 & 规范不变、降低噪声 & 流窗口、SFTX 截断和软/rapidity 不替代 \\
\bottomrule
\end{tabularx}
\end{table}

\section{梯度流作为 UV 正规化}

在 Euclidean 规范理论中定义流场 \(B_\mu(\tau,x)\)：
\begin{equation}
  \partial_\tau B_\mu(\tau,x)
  =D_\nu G_{\nu\mu}(\tau,x),
  \qquad
  B_\mu(0,x)=A_\mu(x),
  \label{eq:flowequation}
\end{equation}
其中 \(G_{\mu\nu}\) 是流场的场强，\(\tau\) 的量纲为长度平方。格点上的 Wilson flow 通过
\begin{equation}
  \dot V_\tau(x,\mu)
  =-g_0^2\left[\partial_{x,\mu}S_W(V_\tau)\right]V_\tau(x,\mu),
  \qquad
  V_0=U
  \label{eq:wilsonflow}
\end{equation}
实现\cite{Luscher2010}。Lüscher--Weisz 的微扰分析表明，在标准 QCD 参数已重正化的前提下，流场的规范不变局域关联函数在 \(d\le4\) 时是 UV 有限的\cite{LuscherWeisz2011}。这使流半径
\begin{equation}
  r_f=\sqrt{8\tau}
  \label{eq:flowradius}
\end{equation}
成为一个独立于格距的平滑尺度。

然而，流动并不等于完整 TMD 重正化。流场是 UV 调节器，短流时展开把流化算符表示为普通重正化算符的线性组合：
\begin{equation}
  O_R(\tau,x)
  \sim\sum_i c_i(\tau)\,\mathcal O_{R,i}(x)
  +\mathcal O(\tau\Lambda_{\rm QCD}^2),
  \label{eq:sftx}
\end{equation}
其中 Wilson 系数 \(c_i\) 与流方案和算符表示有关。对非局域 TMD 算符，短距展开还必须同时考虑 \(z,\bT\) 和 \(P_z\)。因此，流化只提供一个可控的短距离窗口，不能替代软函数、rapidity 方案和硬匹配。

\section{流的微扰图像与短流时系数}

在自由场极限中，把流场按 Fourier 模式展开，线性化流方程给出横向极化模式的高斯衰减：
\begin{equation}
  \widetilde B_\mu(\tau,k)
  =\left[e^{-k^2\tau}P^\perp_{\mu\nu}(k)
  +P^\parallel_{\mu\nu}(k)\right]\widetilde B_\nu(0,k),
  \label{eq:flowkernel}
\end{equation}
其中 \(P^\perp_{\mu\nu}=g_{\mu\nu}-k_\mu k_\nu/k^2\) 是横向投影，\(P^\parallel_{\mu\nu}=k_\mu k_\nu/k^2\) 是纵向投影。该式说明流半径 \(r_f=\sqrt{8\tau}\) 是扩散长度，不是简单的“平滑步数”。在非阿贝尔理论中，非线性项会改变高阶系数，但短流时展开仍按规范不变局域算符组织。

规范作用量密度的典型短流时展开为
\begin{equation}
  E(\tau,x)
  =\frac{1}{4}G^a_{\mu\nu}G^a_{\mu\nu}
  \sim \langle E(\tau,x)\rangle
  +c_E(\tau)\left[\frac{1}{4}F^a_{\mu\nu}F^a_{\mu\nu}\right]_R
  +\mathcal O(\tau),
  \label{eq:Esftx}
\end{equation}
其中 \(c_E(\tau)\) 的紫外行为由重整化群控制。更一般地，任何流化规范不变算符都可写成
\begin{equation}
  O_R(\tau,x)
  =\sum_i c_i(\tau)\mathcal O_{R,i}(x)
  +\mathcal O(\tau\Lambda_{\rm QCD}^2),
  \label{eq:genericsftx}
\end{equation}
系数 \(c_i\) 必须通过微扰匹配或非微扰方案确定。对胶子 TMD，非局域算符的 SFTX 可以看成在固定 \(z,\bT\) 下先对流场做局域展开，再对剩余非局域结构做算子匹配；因此流时间和外部分离的极限不可交换。

\section{流时间外推与匹配策略}

单一流时间的数值结果不能作为物理结果。推荐至少做到：
\begin{enumerate}
  \item 选取若干满足 \(a\ll r_f\ll b_T,z\) 的流时；
  \item 在每个流时上独立完成 Wilson 圈、场强和准 TMD 矩阵元计算；
  \item 用同一匹配核把结果转换到 \(\MSbar\) 或选定的连续方案；
  \item 对 \(\tau\to0\) 做线性或包含高阶项的受控外推；
  \item 检查不同外推窗口、不同 \(z_0\) 和不同 \(\alpha_s\) 阶数下的稳定性；
  \item 将流时间误差与格距、体积、统计和匹配误差分别报告。
  \end{enumerate}
如果只在 \(\tau=3a^2\) 上做计算，即使矩阵元有限，也无法区分短流时物理、格距效应和流化导致的非微扰污染。该点与第 6 节库内审计中的 \(r_f>b_T\) 问题直接相关。

\section{流时间窗口}

为了让流化 TMD 同时接近物理光锥关联和保持可微扰匹配，需要满足尺度层级
\begin{equation}
  a\ll r_f\ll \min\left(z,b_T,\Lambda_{\rm QCD}^{-1}\right),
  \qquad
  P_zr_f\ll1 .
  \label{eq:flowwindow}
\end{equation}
第一组不等式表示：流半径必须大于格距以平滑 UV，但小于非局域算符的物理分离，否则流会抹平待测的横向结构。第二式要求 boost 不把流半径放大成硬尺度。实际计算中还需加入
\begin{equation}
  b_T\ll L_s/2,
  \qquad
  L_\parallel\gg b_T,
  \label{eq:volumeconditions}
\end{equation}
以控制有限体积和 staple 平台。一个“流后矩阵元有限”的结果，如果违反这些窗口，仍然只能作为诊断。

\section{梯度流、混合与单态演化}

梯度流可以降低高维算符污染和统计噪声，但它不能消除所有算符混合。对非极化胶子，等时算符的不同 Lorentz 分量可能有不同的反常量纲；可乘性重正化的组合必须在场强指标和横向投影上固定\cite{ZhangGluonLaMET2019,LiMaQiu2019}。在核子单态通道中，胶子与夸克的物理演化还会在因子化层发生混合，表现为硬核矩阵 \(C_{gg},C_{gq},C_{qg},C_{qq}\)。这与 UV 重正化矩阵不是同一个概念，但二者都要求在最终误差预算中显式列出。

\clearpage
\chapter{梯度流重整化：从规范场演化到可匹配算符}
\label{sec:b06-flow-renormalization}

\section{对象、符号与三条不能混用的链}

梯度流在本项目中有三种经常被混称的角色。第一，它是一种规范协变的四维扩散，用正流时
\(\tau\) 压制高动量涨落；第二，它定义一个有限流时的重整化方案，使流场算符可以通过小流时
展开与无流算符匹配；第三，它可能作为准 PDF 或准 TMD 的 UV 正规化工具。三者共享流方程，
却不共享完整的物理含义。一个有限的流化矩阵元既不是无流 \(\MSbar\) 矩阵元，也不是光锥
分布函数。

\begin{table}[htbp]
\centering
\caption{B06 统一记号。流时和几何分离在物理极限中必须保持独立。}
\label{tab:b06-symbols}
\small
\begin{tabularx}{\textwidth}{@{}p{0.20\textwidth}p{0.28\textwidth}X@{}}
\toprule
符号 & 物理含义 & 使用边界 \\
\midrule
\(\tau\) 或 \(t\) & 物理流时，量纲为长度平方 & 格点无量纲量记为 \(\widehat\tau=\tau/a^2\)；两者不可混用 \\
\(r_f\) & \(r_f=\sqrt{8\tau}\) 的平滑半径 & 不是 smear 步数，也不是普通重整化标度 \\
\(z\) & 沿 Wilson line 的空间分离 & 非局域算符的短距匹配变量 \\
\(\bT\) 或 \(b_T\) & TMD 的横向坐标空间分离 & 只有 staple/TMD 几何才能承载其物理信息 \\
\(L_\parallel\) & staple 纵向臂长 & 必须检查平台或外推；不能默认等于物理无穷远 \\
\(P_z\) & 核子沿 \(z\) 的 Euclidean 大动量 & 控制 LaMET 幂修正和匹配窗口 \\
\(Z_\chi\) & 流夸克场的乘性重整化 & 含流夸克的非局域算符必须保留；纯流胶子算符没有这一因子 \\
\(C_i,\delta m_A\) & GF 到无流准算符的端点系数和线质量 & 依赖投影、表示、流方案和微扰阶数 \\
\(S_I^g,K_g\) & 胶子 soft 结构和 Collins--Soper 核 & 必须与 TMD Wilson-line 几何、伴随表示同一方案 \\
\bottomrule
\end{tabularx}
\end{table}

本章首先建立流场与 SFTX 的语言，然后说明哪些重整化因子确实不需要，哪些仍然必须保留。
下一章把固定流时连续极限与流时间窗口写成可执行顺序。第~\ref{sec:b06-tmd-continuum}
节再把同一逻辑用于胶子 TMD，特别区分直线 Wilson-line OPE 与 staple TMD。

\section{流方程、热核图像与格点实现}

在 \(D\) 维 Euclidean QCD 中，纯规范梯度流满足
\begin{equation}
  \partial_\tau B_\mu(\tau,x)
  =D_\nu G_{\nu\mu}(\tau,x),
  \qquad B_\mu(0,x)=A_\mu(x),
  \label{eq:b06-flow-equation}
\end{equation}
其中 \(G_{\mu\nu}=\partial_\mu B_\nu-\partial_\nu B_\mu+[B_\mu,B_\nu]\)，
\(D_\mu=\partial_\mu+[B_\mu,\cdot\,]\) 作用在伴随表示上。在线性化极限中，
横向模式满足
\begin{equation}
  \widetilde B_\mu(\tau,k)
  =\left[P^\perp_{\mu\nu}(k)e^{-\tau k^2}
  +P^\parallel_{\mu\nu}(k)\right]\widetilde B_\nu(0,k)
  +O(B^2),
  \label{eq:b06-heat-kernel}
\end{equation}
因此 \(\sqrt{8\tau}\) 是高斯平滑宽度的量级。这个结论说明为什么短流时可以接近无流算符，
也说明为什么大流时会抹掉待测的 \(z\) 或 \(b_T\) 结构。式~\eqref{eq:b06-heat-kernel}
不是全阶非微扰定理；非阿贝尔相互作用改变高阶结构，但不会取消正流时的 UV 平滑。

格点上实际演化的是链接 \(V_\tau(x,\mu)\)。以 Wilson 规范作用量为例，
\begin{equation}
  \dot V_\tau(x,\mu)
  =-g_0^2\left[\partial_{x,\mu}S_W(V_\tau)\right]V_\tau(x,\mu),
  \qquad V_0(x,\mu)=U(x,\mu).
  \label{eq:b06-wilson-flow}
\end{equation}
常用三阶 Runge--Kutta 积分。若时间步为 \(\epsilon\)，单步局部误差为
\(O(\epsilon^3)\)；最终报告还必须记录实际步数、最后一步截断和链接幺正性/行列式守卫。
式~\eqref{eq:b06-wilson-flow} 的连续极限并不自动给出目标流时的精确数值，离散积分误差应独立
检查。

物理上应监测 Wilson plaquette 作用量密度，而不是把所有 Clover 组合都当作逐点单调量：
\begin{equation}
  s_W(\tau,x)=\frac16\sum_{\mu<\nu}
  \left[1-\frac{\Re\Tr P_{\mu\nu}(\tau,x)}{N_c}\right].
  \label{eq:b06-plaquette-density}
\end{equation}
正流时通常降低平均作用量密度并改善短距噪声，但“作用量下降”只是实现诊断，不能替代流时
窗口、连续极限和算符匹配。

\section{\texorpdfstring{\(D+1\) 维表示与 UV 有限性}{D+1 维表示与 UV 有限性}}

纯胶子流场 UV 有限性的核心不是“平滑后什么都有限”，而是流相关函数可以重新组织为一个
\(D+1\) 维局部场论。体部场是 \(B_\mu(\tau,x)\)，流方程通过 Lagrange multiplier
\(L_\mu\) 强制。其体作用量可写成
\begin{equation}
  S^{(G)}_{\rm fl}
  =-2\int_0^\infty d\tau\int d^D x\,
  \Tr\left\{
  L_\mu\left[
  \partial_\tau B_\mu-D_\nu G_{\nu\mu}
  -\alpha_0D_\mu\partial_\nu B_\nu
  \right]\right\},
  \label{eq:b06-dplus1-action}
\end{equation}
其中 \(\alpha_0\) 项只改变流时间方向的规范固定。对 \(L_\mu\) 作泛函积分恢复流方程；体传播子
在 \(k^2\to\infty\) 时含 \(e^{-\tau k^2}\)，于是正流时的体图 UV 收敛。

\begin{quote}
\textbf{纯胶子定理的精确表述。} 在标准 QCD 参数已经重整化后，由正流时
\(B_\mu(\tau,x)\) 构成的规范不变局域乘积，其关联函数具有有限的连续极限；流规范场
\(B_\mu\) 本身不需要额外的波函数重整化。BRS 对称性和 \(D+1\) 维幂次计数把可能的抵消项
限制在 \(\tau=0\) 边界，即普通四维 QCD 的抵消项。
\end{quote}

不能把这个定理外推成“任意流场乘积均无额外因子”。若算符包含流夸克场
\(\chi(\tau,x)\)，Lüscher 的乘性重整化为
\begin{equation}
  \chi_R=Z_\chi^{-1/2}\chi,
  \qquad
  \bar\chi_R=\bar\chi\,Z_\chi^{-1/2}.
  \label{eq:b06-zchi-definition}
\end{equation}
在来源所用 MSbar 约定下，一圈系数写作
\begin{equation}
  Z_\chi
  =1+\frac{3C_F}{16\pi^2}\frac{g^2}{\epsilon}
  +O(g^4).
  \label{eq:b06-zchi-oneloop}
\end{equation}
含 \(n\) 个流夸克和 \(\bar n\) 个流反夸克的规范不变局域算符需要整体因子
\begin{equation}
  O_{R,\tau}
  =(Z_\chi)^{(n+\bar n)/2}O_\tau .
  \label{eq:b06-zchi-operator}
\end{equation}
ringed field 通过流夸克动能归一化给出不依赖微扰截断的 \(Z_\chi\) 定义。ringed 双线性在
端点处把两个场因子合并成一个整体归一化；普通未 ringed 的流夸克三点函数不能因为核子
two-point overlap 在 ratio 中相消就当作已经重整化。

\begin{quote}
\textbf{纠错边界。} “\(\tau>0\) 的流场 UV 有限”与“任意非局域 Wilson-line 算符已经完成
重正化和匹配”是不同命题。前者来自流场表示；后者还必须处理端点、线质量、cusp、
pinch-pole、rapidity、软因子以及可能的算符混合。
\end{quote}

\section{小流时展开与方案转换}

对纯胶子局域算符，SFTX 的典型形式是
\begin{equation}
  O(\tau,x)
  \sim \langle O(\tau,x)\rangle
  +\sum_i c_i(\tau,\mu)O_{R,i}(x;\mu)
  +O(\tau\Lambda_{\rm QCD}^2).
  \label{eq:b06-sftx-general}
\end{equation}
若多个算符具有相同量子数，系数应写成矩阵 \(c_{ij}\)，而不是把每个投影都当作可乘性
重整化。规范作用量密度给出最熟悉的例子：
\begin{equation}
  E(\tau,x)
  \sim \langle E(\tau,x)\rangle
  +c_E(\tau)\left[\frac14F^a_{\mu\nu}F^a_{\mu\nu}\right]_R
  +O(\tau).
  \label{eq:b06-sftx-energy}
\end{equation}
对于非局域胶子算符，\(z\)、\(\bT\) 和 \(P_z\) 必须固定在匹配的两端；不能先对 \(z\) 积分，
再用一个只对局域算符成立的系数。

本批次材料给出的胶子直线 Wilson-line 一圈端点为
\begin{align}
  c_{\parallel\perp}(\tau,\mu)
  &=1+O(\alpha_s^2),
  \label{eq:b06-c-parperp}\\
  c_{\perp\perp}(\tau,\mu)
  &=1+\frac{\alpha_s C_A}{4\pi}
  \ln\left(2\mu^2\tau e^{\gamma_E}\right)
  +O(\alpha_s^2),
  \label{eq:b06-c-perpperp}\\
  \delta m_A(\tau)
  &=-\frac{\alpha_s C_A}{4\pi}\sqrt{\frac{2\pi}{\tau}}
  +O(\alpha_s^2).
  \label{eq:b06-line-mass}
\end{align}
这里的 \(\parallel,\perp\) 是相对于 Wilson line 方向的场强指标分类，不是格点代码中
\(t,i\) 标签的简单替换。沿物理长度 \(|z|\) 的线质量因子为
\begin{equation}
  L(z;\tau)=e^{-\delta m_A(\tau)|z|}.
  \label{eq:b06-line-factor}
\end{equation}
对 U031 的两个生产投影，候选 GF-to-quasi 转换为
\begin{align}
  K_U(z;\tau,\mu)
  &=\kappa_F^2\,L(z;\tau)\,c_{\perp\perp}^{-2}(\tau,\mu),
  \label{eq:b06-KU}\\
  K_H(z;\tau,\mu)
  &=\kappa_F^2\,L(z;\tau)
  \left[c_{\parallel\perp}(\tau,\mu)c_{\perp\perp}(\tau,\mu)\right]^{-1}.
  \label{eq:b06-KH}
\end{align}
其中 \(U\) 指非极化 \(T_U-S_U\)，\(H\) 指 helicity 的候选 \(T_H-S_H\) 组合。
\(\kappa_F\) 是树图归一化约定；来源中的诊断取 \(\kappa_F=1\)，但没有独立证明它就是最终
Clover/color-projection 的完整归一化。式~\eqref{eq:b06-KH} 的 helicity 版本也是候选，
不能称为完整 partonic matching。

符号和重复扣除必须作为硬约束：
\begin{itemize}
  \item \(\tau\) 在 \(\ln(2\mu^2\tau e^{\gamma_E})\) 中必须是物理
        \(\mathrm{GeV}^{-2}\)，不能直接放入 \(\widehat\tau=\tau/a^2\)；
  \item \(L(z;\tau)\) 使用 \(|z|\)，因此正负方向共享同一线质量；
  \item 若 ratio/hybrid 已经吸收线质量，不得再次乘 \(L(z;\tau)\)；
  \item 端点因子应先作用在独立分量上，再重构 combined，避免分量与组合使用不同方案；
  \item \(K_U,K_H\) 仍属于 GF 到无流准算符的转换，不包含 TMD soft、rapidity 或 LaMET。
\end{itemize}

\section{推荐的可执行顺序}

在进入极限分析前，流化胶子直线 OPE 的最小顺序是：
\begin{enumerate}
  \item 从 thin link 开始作四维 Wilson flow，固定物理 \(\tau\)，保存实际流步和流后规范；
  \item 在同一 \(V_\tau\) 上构造 Clover、对偶场强和直线伴随 Wilson line；
  \item 保存正负 \(z\)、原始分量、无迹/未去迹两种 projection 以及所有线性恒等式；
  \item 对每个 source 作真空扣除和 disconnected \(C_3/C_2\)，保留共同 nopol 分母；
  \item 在逐 replica、逐组态或明确的 fit-level scalar 路径上施加 \(K_U,K_H\)，并记录两者差异；
  \item 在通过流窗口的 \((z,P_z,\tau)\) 上做连续极限和匹配后的 \(\tau\) 外推；
  \item 最后才进入有限 \(z\) Fourier、LaMET/pseudo-ITD 和 singlet mixing。
\end{enumerate}
任何一步只产生有限文件或通过形状测试时，状态仍应写成该步的对象，而不是“physical PDF”。

\chapter{流时窗口、固定流时连续极限与匹配次序}
\label{sec:b06-flow-window}

\section{三个极限不能合并成一个拟合参数}

令有限格距、有限流时的胶子矩阵元为
\begin{equation}
  \mathcal M^{\rm lat,GF}_{i}
  \bigl(z,\bT,P_z;\tau,a\bigr),
  \label{eq:b06-lattice-matrix-element}
\end{equation}
其中 \(i\) 标记 Lorentz 投影和 channel。固定物理 \(z,\bT,P_z,\tau\)，先定义连续
GF-scheme 矩阵元
\begin{equation}
  \mathcal M^{\rm GF}_{i}(z,\bT,P_z;\tau)
  =\lim_{a\to0}\mathcal M^{\rm lat,GF}_{i}(z,\bT,P_z;\tau,a).
  \label{eq:b06-fixed-tau-continuum}
\end{equation}
对 U032 明确分析的直线 OPE projection，再作 GF-to-MSbar 转换：
\begin{equation}
  \mathcal M^{\MSbar,{\rm q}}_{i}(\tau;\mu)
  =e^{-\delta m_A(\tau)|z|}
  C^{-1}_{ij}(\tau,\mu)\mathcal M^{\rm GF}_{j}(\tau).
  \label{eq:b06-gf-to-quasi}
\end{equation}
只有匹配后的对象才应分析残余流时依赖：
\begin{equation}
  \mathcal M^{\MSbar,{\rm q}}_{i}(\mu)
  =\lim_{\tau\to0}\mathcal M^{\MSbar,{\rm q}}_{i}(\tau;\mu).
  \label{eq:b06-matched-flow-limit}
\end{equation}
这三行不是同一个外推中的三个变量。未匹配的 GF 矩阵元在 \(\tau\to0\) 通常没有有限的
非零极限，因为线质量因子和端点系数本身按 \(1/\sqrt\tau\) 或对数发散。把
\(\mathcal M^{\rm lat,GF}\) 直接拟合成 \(\tau\to0\) 截距，会得到一个依赖匹配方案和
截断的数值。

\begin{table}[htbp]
\centering
\caption{三类对象、允许的极限和不能越过该层的宣称。}
\label{tab:b06-object-levels}
\small
\begin{tabularx}{\textwidth}{@{}p{0.20\textwidth}p{0.27\textwidth}p{0.25\textwidth}X@{}}
\toprule
对象 & 依赖量 & 允许操作 & 仍不能宣称 \\
\midrule
finite-flow bare & \(a,\tau,z,\bT,P_z,L_\parallel\) 与源汇结构 & ratio、谱拟合、真空扣除、协方差 & 重整化矩阵元或物理分布 \\
GF-scheme continuous & 固定物理 \(\tau,z,\bT,P_z\) & 多格距连续极限、SFTX/GF conversion & 光锥 TMD/PDF；TMD soft 也未处理 \\
MSbar quasi & \(\tau\) 已转换为 \(\mu\)，仍有限 \(P_z\) & 匹配后 \(\tau\to0\)、Fourier、LaMET & 无穷动量极限和完整 singlet 物理量 \\
light-cone & \(\mu,\zeta\) 与目标过程/表示 & 与实验或全局分析比较 & 若 soft、mixing 或极限缺失则仍无效 \\
\bottomrule
\end{tabularx}
\end{table}

\section{\texorpdfstring{\(a^2/\tau\) 与极限交换}{a2/tau 与极限交换}}

固定物理 \(\tau\) 时，流半径随格距缩小而变大，因此格点 cutoff 效应不一定从普通的
\(a^2\) 开始。基于 U032 所引局域 flowed PDF moment 结果，本批次把它作为诊断形式：
\begin{equation}
  \widehat{\mathcal M}^{\MSbar,{\rm q}}(\tau;a)
  =\mathcal M^{\MSbar,{\rm q}}
  +c_1\frac{a^2}{\tau}
  +c_2a^2P_z^2
  +c_3\tau\Lambda_{\rm QCD}^2
  +c_4\alpha_s^{N+1}
  +\cdots .
  \label{eq:b06-cutoff-ansatz}
\end{equation}
各项物理来源不同：\(a^2/\tau\) 表示流半径相对格距的离散误差，\(a^2P_z^2\) 表示大动量与
粗格距混合，\(\tau\Lambda_{\rm QCD}^2\) 表示短流时展开的高维修正，最后一项是匹配截断。
本式不是对所有非局域胶子算符已经证明的普适拟合公式；系数可以相关，高阶项也不必按
这里的顺序排列。它的用途是迫使分析显式报告极限交换和窗口，而不是给任意数据增加参数。

若 \(a/\sqrt\tau\) 不小时，先在固定 \(\tau/a^2\) 上比较不同格距会把不同物理流时混在一起；
若先取 \(\tau\to0\) 再取 \(a\to0\)，则会把 cutoff 增强项误当成流时物理。正确的数据组织是
在每个物理 \(\tau\) 上分别完成式~\eqref{eq:b06-fixed-tau-continuum}，再比较匹配后的
\(\tau\) 依赖。若没有多个格距，只能报告“固定格距、有限流时的参数点结果”。

\section{联合流时窗口与源汇 footprint}

对直线 Wilson-line 胶子 OPE，最小尺度层级是
\begin{equation}
  a\ll r_f=\sqrt{8\tau}\ll |z|,
  \qquad
  \sqrt{\tau}\,|P_z|\ll1,
  \qquad
  \tau\Lambda_{\rm QCD}^2\ll1,
  \qquad
  P_z^2\gg M_N^2 .
  \label{eq:b06-line-window}
\end{equation}
第一个不等式要求流化确实压低 UV，同时不能吞掉待测的纵向结构；第二式控制 boost 后的
流尺度；第三式保证 SFTX 的幂次展开；第四式保证 LaMET 的大动量条件。它们必须逐点满足，
不能在平均参数上满足后就保留所有点。

非局域插入还要求流半径不碰源和汇：
\begin{equation}
  r_f\ll t_{\rm insert}-t_{\rm source},
  \qquad
  r_f\ll t_{\rm sink}-t_{\rm insert}.
  \label{eq:b06-source-sink-footprint}
\end{equation}
对 TMD staple 还要加入
\begin{equation}
  r_f\ll b_T,
  \qquad
  b_T\ll L_s/2,
  \qquad
  L_\parallel\gg b_T,
  \label{eq:b06-tmd-window}
\end{equation}
并检查 \(L_\parallel\to\infty\) 的平台或受控外推。直线 OPE 没有物理 \(b_T\)，不能把
式~\eqref{eq:b06-tmd-window} 的第二、三式伪装成已经控制的直线 Wilson-line 条件。

窗口是交集，不是可逐点优选菜单。实践上应先按物理 \(\tau,z,b_T,P_z\) 建立 mask，再在
同一 mask 上做拟合，并报告被排除的点。若交集为空，结论只能是“当前格距/体积/动量下
没有可发布窗口”，不能通过增大误差棒制造物理平台。

\section{三类匹配与两类混合}

\begin{longtable}{@{}p{0.18\textwidth}p{0.27\textwidth}p{0.21\textwidth}p{0.22\textwidth}@{}}
\caption{B06 要求的匹配层级。每一行改变对象，不能只改名字。}\label{tab:b06-matching-layers}\\
\toprule
层级 & 输入到输出 & 主要依赖 & 典型错误 \\
\midrule
\endfirsthead
\toprule
层级 & 输入到输出 & 主要依赖 & 典型错误 \\
\midrule
\endhead
GF 方案转换 & GF-scheme flowed matrix element \(\to\) MSbar unflowed quasi & \(\tau,\mu,z,C_i,\delta m_A\) & 把 \(K_U/K_H\) 叫作 LaMET，或把线质量重复扣除 \\
准量到光锥 & finite-\(P_z\) quasi PDF/TMD \(\to\) light-cone distribution & \(P_z,\mu,\zeta\)；\(C_{\rm hard},S,K\)；\(\alpha_s\) 阶数 & 用有限 \(P_z\) 直接声称光锥 PDF \\
UV 算符混合 & \(O_i^{R}=Z_{ij}O_j\) & 同量子数算符、反常量纲矩阵 & 把 UV \(Z_{ij}\) 与 flavor-singlet 硬核混为一类 \\
因子化味混合 & \(g\leftrightarrow q\) 在准量到光锥 matching 中卷积 & \(C_{gg},C_{gq},C_{qg},C_{qq}\) & 关闭 mixing 却不报告系统误差 \\
\bottomrule
\end{longtable}

最重要的操作原则是：GF 流场不需要额外波函数重整化，并不表示流化 Wilson-line 算符没有
端点或线质量；TMD 有软因子和快度结构，也不表示 first matching 可以省略。相反，第二层
matching 越完整，越需要明确第一层的输入方案。

\section{来源中的数值状态与可验证性}

N406 是当前 GF 胶子产品中最新的统计快照，但它仍不是连续极限。来源记录的关键元数据为：
\begin{center}
\small
\begin{tabularx}{\textwidth}{@{}lX@{}}
\toprule
项目 & 来源记录 \\
\midrule
matched small-flow 输入 & \(N_{\rm conf}=406\)，\(N_{\rm boot}=1000\) \\
格距 & \(a=0.0775\,\mathrm{fm}\) \\
流时 & \(\tau/a^2=1.0,1.4,1.8,2.2,2.6,3.0,3.4,3.8\) \\
小流时模型 & \(M(\tau)=M_0+c\tau\)，replica-level AIC 权重 \\
动量/分离 & \(P_z=3,4,5\)，\(z/a=0,\ldots,24\) \\
产品边界 & \(P_z=3\) 有缺失外推点；\(P_z=5\) 有 1 个 \(z=0\) 归一化符号翻转风险 \\
\bottomrule
\end{tabularx}
\end{center}
这些数字来自来源 manifest 与交接记录，不是本批次重跑。它们证明已有统计产品覆盖
406 个组态和多个流时，不证明存在多个格距，也不证明匹配后结果已经收敛。

\begin{table}[htbp]
\centering
\caption{B06 证据的用途边界。}
\label{tab:b06-evidence-boundaries}
\small
\begin{tabularx}{\textwidth}{@{}p{0.18\textwidth}p{0.23\textwidth}p{0.27\textwidth}X@{}}
\toprule
来源 & 能证明 & 不能证明 & 章节用途 \\
\midrule
U061 理论专著 & 流方程、\(D+1\)/BRS 结构、\(Z_\chi\)、局域 SFTX & 非局域 TMD 的完整匹配或数值结果 & 提供梯度流重整化章节的定理和纠错边界 \\
U032 文献综合 & GF-SFTX-LaMET 的两阶段逻辑和固定流时顺序 & 单一文献已完成质子胶子全链 & 提供极限与窗口综合 \\
U031 算法报告 & 直线 OPE、ratio、端点转换和契约测试 & 真实生产能力、TMD soft 或连续极限 & 提供流时窗口章节的实现对象 \\
U030 算法报告 & staple/TMD 链、参数与失败模式 & 物理胶子 TMD；test9 仅 10 组态且窗口失败 & 提供胶子 TMD 验收与真实数据章节的对象和风险表 \\
U023 产品交接 & N406 元数据、缺失点和 warning & 新物理测量或连续外推 & 给出数据状态，不作为新结果 \\
\bottomrule
\end{tabularx}
\end{table}

\chapter{胶子 TMD、staple 匹配与连续极限验收}
\label{sec:b06-tmd-continuum}

\section{从准 PDF 的直线链到 TMD 的 staple 链}

胶子 TMD 的光锥定义同时依赖 \(x,b_T,\mu\) 和 rapidity 标度 \(\zeta\)。在示意图层面，
可写成
\begin{equation}
  f_1^g(x,b_T;\mu,\zeta)
  =\mathcal N_g\int\frac{dz^-}{2\pi}
  e^{-ixP^+z^-}
  \langle P|O_g^{+i;+i}(z^-,b_T;\mu,\zeta)|P\rangle,
  \label{eq:b06-tmd-definition}
\end{equation}
其中 \(O_g^{+i;+i}\) 含两条光样 Wilson line、横向连接和过程相关方向，\(\mathcal N_g\)
依场强归一化和指标投影而定。式~\eqref{eq:b06-tmd-definition} 不是格点可直接计算的
Euclidean 表达式；格点用有限 \(P_z\)、类空分离和有限长 staple 逼近它。

胶子准 PDF 的直线几何是
\begin{equation}
  O^{\rm line}_{\mu\nu;\rho\sigma}(z)
  =\sum_x\Tr\left[
  F_{\mu\nu}(x+z\hat z)
  W_{\rm adj}(x+z,x)
  F_{\rho\sigma}(x)
  \right],
  \label{eq:b06-line-operator}
\end{equation}
而 TMD 的几何必须多出横向分离：
\begin{equation}
  O^{\rm TMD}_{\mu\nu;\rho\sigma}(z,b_T;L_\parallel)
  =\sum_x\Tr\left[
  F_{\mu\nu}(x+z\hat z+b_T)
  W_{\rm st}(z,b_T;L_\parallel)
  F_{\rho\sigma}(x)W^\dagger_{\rm st}
  \right].
  \label{eq:b06-staple-operator}
\end{equation}
其中 \(W_{\rm st}\) 是沿三段路径有序指数化的有限臂长 staple。式~\eqref{eq:b06-line-operator}
的短距 OPE 只组织直线线的线质量、端点和局域算符；式~\eqref{eq:b06-staple-operator}
还包含两条近乎平行长线之间的 pinch-pole、cusp、软因子和快度结构。两者不能用同一个
\(K_U/K_H\) 或同一个 Wilson-loop subtraction 互相替代。

当 \(b_T\to0\) 时，staple 算符可以退化为某种直线几何，但只有 soft subtraction、端点
scheme 和 Fourier convention 同时定义后，才可与共线 PDF 比较。一个数值上的
\(b_T=0\) 点单独不能证明 TMD/PDF 归一化一致。

\section{Lorentz 投影与 disconnected 估计量}

沿用 U031 的 Euclidean 记号，定义
\begin{equation}
  T_i=M_{ti;ti},\qquad
  T_j=M_{tj;tj},\qquad
  S_0=M_{ij;ij},
  \label{eq:b06-ope-primitives}
\end{equation}
则非极化胶子组合为
\begin{equation}
  O_U=M_{titi}-M_{ijij}
  =T_i+T_j-2S_0 .
  \label{eq:b06-unpolarized-combination}
\end{equation}
helicity 目标为
\begin{equation}
  O_H^{\rm phys}
  =\left(M_{titi}-M_{ijij}\right)_{\rm odd},
  \label{eq:b06-helicity-combination}
\end{equation}
其中第二个场强使用对偶张量。历史文件中的 \(M_{titi}+M_{ijij}\) 只能作符号和方向诊断。
方向的偶/奇投影使用
\begin{equation}
  O_{\rm even}=O_{+z}+O_{-z},
  \qquad
  O_{\rm odd}=O_{+z}-O_{-z},
  \label{eq:b06-even-odd}
\end{equation}
不额外除以 2；任何归一化改变都必须同步进入 Fourier 和 matching。

对胶子 disconnected 插入，配置估计量为
\begin{align}
  C_{3,\chi}
  &=\left\langle O_\chi C_{2,\chi}\right\rangle
  -\left\langle O_\chi\right\rangle
  \left\langle C_{2,\chi}\right\rangle,
  \label{eq:b06-c3-disconnected}\\
  C_{2,\rm den}
  &=\left\langle C_{2,\rm nopol}\right\rangle,
  \qquad
  R_\chi=\frac{C_{3,\chi}}{C_{2,\rm den}}.
  \label{eq:b06-ratio-common-denom}
\end{align}
helicity 的 \(C_{3,H}\) 保留复数 \(pol35\) numerator，两个 channel 使用共同 nopol
分母。最终物理相位投影为
\begin{equation}
  \operatorname{Re}[-iR_H]=\operatorname{Im}R_H,
  \label{eq:b06-helicity-phase}
\end{equation}
但只能在逐 replica 的真空扣除、ratio、协方差和 jackknife 完成后进行。提前取实部/虚部
会破坏复数的配置相关性；用 \(pol35\) 均值作近零分母则同时改变物理量和统计稳定性。

\section{flow、soft、CS 与 hard matching 的对象图}

一种可执行的 TMD 实现把匹配组织为
\begin{equation}
  \begin{aligned}
  x\widetilde g(x,b_T;P_z)
  =\int\frac{dy}{|y|}\,
  C_{\rm hyb}\!\left(\frac{x}{y},\lambda_s,\frac{\mu}{yP_z}\right)
  \frac{y\,g(y,b_T;\mu)}{\sqrt{S_I^g(b_T;\mu)}}\\
  e^{\frac12\ln(\zeta_z/\zeta)K_g(b_T;\mu)}
  +\text{power corrections}.
  \end{aligned}
  \label{eq:b06-tmd-matching}
\end{equation}
这里 \(S_I^g\) 必须是伴随表示、同一 staple 几何和同一 rapidity 方案的 soft 结构，
\(K_g\) 是相应 CS 核。若 \(\sqrt{S_I^g}\)、Wilson-loop subtraction 或 staple 长度外推
缺失，式~\eqref{eq:b06-tmd-matching} 只能作为算法示意。离散实现中的矩阵
\(Z_{ij}\) 也必须作用在正确的 channel 和 \(b_T\) 网格上；它不是把 soft 因子替换成一个
标量。

从有限 \(P_z\) 准量到光锥单态分布还需要
\begin{equation}
  \begin{pmatrix}\Sigma(x,\mu)\\ g(x,\mu)\end{pmatrix}
  =\mathbf C^{\rm LaMET}\!\left(\frac{\mu}{P_z}\right)\otimes
  \begin{pmatrix}\widetilde\Sigma(x,P_z;\mu)\\ \widetilde g(x,P_z;\mu)\end{pmatrix}
  +\mathcal O\!\left(\frac{\Lambda_{\rm QCD}^2}{[xP_z]^2},
  \frac{\Lambda_{\rm QCD}^2}{[(1-x)P_z]^2}\right),
  \label{eq:b06-lamet-mixing}
\end{equation}
其中 \(\Sigma=\sum_f(q_f+\bar q_f)\)。C 的卷积方向取决于来源 convention；不能把
一个 quasi\(\to\)light-cone kernel 当成反向 kernel。对 helicity，应使用独立的
polarized hard matrix 和 soft/rapidity 配套，而不是把式~\eqref{eq:b06-lamet-mixing}
中的非极化 kernel 做符号替换。

\begin{table}[htbp]
\centering
\caption{B06 建议的胶子 TMD 物理链。每一层都有独立的失败关闭条件。}
\label{tab:b06-tmd-chain}
\small
\begin{tabularx}{\textwidth}{@{}p{0.05\textwidth}p{0.20\textwidth}p{0.28\textwidth}X@{}}
\toprule
层 & 对象 & 必须固定 & 失败关闭条件 \\
\midrule
T0 & 算符和路径 & 场强、对偶、生成元、伴随链、方向、Fourier 符号 & 端点规范恒等式或路径约定不闭合 \\
T1 & 有限流时矩阵元 & \(a,\tau,z,b_T,P_z,L_\parallel\)、源汇 & ratio 不稳定、mask 被零填充、单一 source \\
T2 & GF 方案转换 & \(C_i,\delta m_A,\kappa_F,\mu\) & 使用不匹配投影或重复线质量 \\
T3 & soft/rapidity & \(S_I^g,K_g,\zeta,L_\parallel\)、伴随表示 & 缺同几何 Wilson loop 或 \(L_\parallel\) 无平台 \\
T4 & hard matching & \(C_{gg},C_{gq}\)、阶数、\(\mu/P_z\) & 关闭 mixing 且未给误差、kernel 方向不清 \\
T5 & 物理极限 & 多 \(a\)、匹配后 \(\tau\to0\)、多 \(P_z\)、体积与手征 & 任一极限只有单点或单格距 \\
\bottomrule
\end{tabularx}
\end{table}

\section{来源数据能支持到什么程度}

U030 的 test9 链有可执行的 flow、staple、ratio、\(Z_R\)、Fourier、CS 和 matching
接口，但参数是 \(24^3\times72\)、\(a=0.1053\,\)fm、10 个组态、\(z=0,\ldots,12\)、
\(b_T=0,\ldots,4\)，且流时约为 \(3a^2\)。对应
\(r_f/a=\sqrt{24}\approx4.90\)，已经大于非零 \(b_T/a\le4\)。这不是轻微的数值误差，
而是流半径会抹平目标横向结构的窗口失败。PDF 第 17 页的归档还记录了
\(N_{\rm conf}=10\)、\(13\times5\) 网格和多处 \(\operatorname{cond}=\infty\)；因此 test9 的有限图和
matched curve 只能证明接口被调用。

U031 的直线 OPE 报告提供了 \(9/9\) 契约测试和 106 项回归的静态记录，PDF 第 14 页
展示相应数值表。它证明小数组下的轴、投影、共同分母和组合恒等式，不证明真实 ensemble
的流窗口和连续极限。U023 的 N406 信息增加了统计产品覆盖，但依然只有一个格距；多个
\(\tau/a^2\) 不是多个物理流时，多个 bootstrap 也不是多个格距。

因此本批次可以作出的最强结论是：胶子直线 OPE、TMD staple 几何和 GF-SFTX-LaMET
的对象关系已经可以写成可审计的算法链；完整物理胶子 TMD 仍需要同一表示/同一几何的
soft 和 CS、完整单态 mixing、多个格距以及匹配后流时外推。来源中的测试计数、N406
状态和归档图都不构成本批次的新数值测量。

\section{失败模式与最小修正}

\begin{longtable}{@{}p{0.24\textwidth}p{0.31\textwidth}p{0.37\textwidth}@{}}
\caption{梯度流胶子 TMD 的常见错误、症状和失败关闭规则。}\label{tab:b06-failures}\\
\toprule
错误 & 典型症状 & 最小修正 \\
\midrule
\endfirsthead
\toprule
错误 & 典型症状 & 最小修正 \\
\midrule
\endhead
把 finite-flow bare \(C_3/C_2\) 称为 PDF & 曲线可画但没有统一 UV/rapidity 方案 & 状态保持 bare/diagnostic；先完成方案转换 \\
把三维 flow 当成四维 Wilson flow & 无 \(D+1\) 定理和四维匹配系数可用 & 重新定义 scheme；不得复用四维 \(Z_\chi,C_i\) \\
跨格距固定 \(\tau/a^2\) & 看似有多个流时，实际比较不同物理 \(r_f\) & 固定物理 \(\tau\)、\(z,b_T,P_z\)，逐格距重算 \\
先作 \(\tau\to0\) 再 \(a\to0\) & \(\tau\to0\) 截距依赖 cutoff 和 \(z\) & 先固定 \(\tau\) 取 \(a\to0\)，匹配后再作流时外推 \\
\(r_f>b_T\) 或 \(r_f>|z|\) & 横向/纵向结构被流化抹平，平台变平滑 & 多流时全局 mask；交集为空就停止 \\
把直线 OPE 的 \(K_U/K_H\) 用于 staple & soft、pinch-pole 和 rapidity 被静默遗漏 & 直线和 TMD 分链；TMD 另算 \(S_I,K_g,Z_E\) \\
把 even/odd 除以 2 & Fourier 和 helicity 归一化偏移 & 采用来源约定；任何改动都进入 matching \\
提前取 helicity 实部/虚部 & 协方差丢失，虚部投影不可恢复 & 全程保留 complex128，最后取 \(\operatorname{Im}R_H\) \\
用 \(pol35\) 均值作分母 & 近零分母和 ratio 病态 & 两 channel 共享 nopol 分母，逐 replica 相除 \\
重复乘 \(\exp[-\delta m|z|]\) & ratio/hybrid 后出现人为 \(z\) 斜率 & 明确非微扰减除是否已含线质量 \\
令 \(\kappa_F=1\) 并称为已证明 & 整体归一化误差伪装成 matching & 把它列为未验证归一化，做独立树图/投影检查 \\
关闭 singlet mixing 又称光锥单态 & 胶子与 quark singlet 比较出现整体偏移 & 加入 \(C_{gq},C_{qg}\) 或给出严格忽略依据 \\
\bottomrule
\end{longtable}

\section{分级验收与本章结论}

建议把 B06 的胶子 TMD 结果分为五级，并禁止跨级宣称：
\begin{enumerate}
  \item \textbf{C0：定义级。} 流方程、场强、对偶、伴随 Wilson line、直线/staple 几何和 Fourier 符号均有可执行检查；
  \item \textbf{C1：矩阵元级。} 源分辨两点/三点函数、真空扣除、共同分母、谱拟合和协方差通过质量门；
  \item \textbf{C2：GF 匹配级。} 固定物理 \(\tau\) 有连续极限，\(C_i,\delta m_A,\kappa_F\) 和可能的混合矩阵在同一方案；
  \item \textbf{C3：TMD 方案级。} 伴随表示 soft、rapidity、CS、Wilson-loop subtraction、staple 长度和 hard kernel 闭合；
  \item \textbf{C4：物理极限级。} 多格距、匹配后 \(\tau\to0\)、多 \(P_z\)、体积和手征误差均受控。
\end{enumerate}

本批次材料足以支持 C0 的定义整理、C1/C2 的算法和契约层面检查，以及 C3/C4 的明确失败
条件；不足以支持完整物理胶子 TMD 的数值声称。最重要的纠错不是增加一个 matching
公式，而是坚持对象分层：纯胶子流场的 UV 定理不等于非局域 TMD 的完整重整化；直线
OPE 的端点转换不等于 staple TMD 的软结构；单格距下的多个流时也不等于连续极限。


\clearpage
\part{关联函数、算法与统计推断}
\chapter{从格点关联函数到 TMD}

\section{两点与三点关联函数}

核子两点函数取
\begin{equation}
  C_2(t;\mathbf P)
  =\sum_{\mathbf x}e^{i\mathbf P\cdot\mathbf x}
  \langle N(\mathbf x,t)\overline N(\mathbf 0,0)\rangle .
  \label{eq:two-point}
\end{equation}
胶子算符插入的三点函数为
\begin{equation}
  C_{3,g}(t,\tau;\mathbf P)
  =\sum_{\mathbf x}e^{i\mathbf P\cdot\mathbf x}
  \langle N(\mathbf x,t)\,
  \Oq[\tau,\bT;L_\parallel]\,
  \overline N(\mathbf 0,0)\rangle .
  \label{eq:three-point}
\end{equation}
对纯胶子插入，必须减去真空期望值，并把胶子回路与核子传播子正确收缩。实际的核子胶子矩阵元包含单态和断开型贡献；若忽略某些拓扑，必须证明其在该通道中被压低，而不能只凭图的数量判断。

在源汇时间间隔 \(t_{\rm sep}=t-t_0\) 足够大时，基态矩阵元可通过比值或双态拟合提取：
\begin{equation}
  R_g(t,\tau;t_{\rm sep})
  =\frac{C_{3,g}(t,\tau)}{C_2(t)}
  \longrightarrow
  \frac{\langle P|\Oq|P\rangle}{2E_P}
  \left[1+\mathcal O(e^{-\Delta Et_{\rm sep}})
  +\mathcal O(e^{-\Delta E\min(\tau,t_{\rm sep}-\tau)})\right].
  \label{eq:ratio}
\end{equation}
对高动量核子，基态能量间隙和信噪比都会恶化，因此必须同时报告单态拟合、双态拟合、不同 \(t_{\rm sep}\) 窗口和模型平均结果。仅报告一个 plateau 点不足以证明基态矩阵元。

\section{谱分解、四态结构与比值归一化}

为了看清三点函数中的污染，写出有限体积谱分解。两点函数为
\begin{equation}
  C_2(t;\mathbf P)
  =\sum_n |A_n(\mathbf P)|^2e^{-E_n(\mathbf P)t},
  \qquad
  A_n(\mathbf P)=\langle 0|\chi_N(\mathbf P)|n,\mathbf P\rangle ,
  \label{eq:c2spectral}
\end{equation}
其中 \(\chi_N\) 是核子插值算符。三点函数在 \(t_0<\tau<t_s\) 时为
\begin{equation}
  C_3(t_s,\tau;\mathbf P)
  =\sum_{m,n}
  A_m^*(\mathbf P)A_n(\mathbf P)\,
  \langle m,\mathbf P|\mathcal O_g(\tau)|n,\mathbf P\rangle\,
  e^{-E_m(\mathbf P)(t_s-\tau)}
  e^{-E_n(\mathbf P)\tau}.
  \label{eq:c3spectral}
\end{equation}
把中间态求和重写为谱密度，可以把“拟合了多少个态”转化为连续的正谱问题。定义
\begin{equation}
  \rho_{ij}(\omega;\mathbf P)
  =\sum_n
  \langle 0|\chi_i(\mathbf P)|n,\mathbf P\rangle
  \langle n,\mathbf P|\chi_j^\dagger(\mathbf P)|0\rangle
  \delta\!\left(\omega-E_n(\mathbf P)\right),
  \label{eq:spectraldensity}
\end{equation}
则两点函数写成 Laplace 型变换
\begin{equation}
  C_{ij}(t;\mathbf P)
  =\int_0^\infty\dd\omega\,e^{-\omega t}
  \rho_{ij}(\omega;\mathbf P).
  \label{eq:laplacespectral}
\end{equation}
式~\eqref{eq:laplacespectral} 说明“激发态污染”并不是一个拟合术语，而是有限格点上离散谱与多粒子阈值的物理内容。只有 \(\rho_{ij}\) 在目标能区被足够多的源、汇和 \(t_{\rm sep}\) 稳定约束时，单态外推才有意义。
若 \(t_s\) 不够大，会出现四类贡献：
\begin{enumerate}
  \item \(m=n=0\)：基态矩阵元，目标项；
  \item \(m=0,n>0\) 或 \(m>0,n=0\)：源端或汇端跃迁；
  \item \(m,n>0\)：激发态-激发态项；
  \item 多粒子阈值与有限体积能级，可能表现为非指数或近简并结构。
\end{enumerate}
简单的 \(C_3/C_2\) 比值只在基态主导时收敛。若用两点函数在源端的局部比值，常见的 tree-level 归一化形式是
\begin{equation}
  R_g(\tau,t_s;\mathbf P)
  =\frac{\sum_{\mathbf x}e^{i\mathbf P\cdot\mathbf x}
  \langle N(t_s,\mathbf x)\mathcal O_g(\tau)\overline N(0)\rangle}
  {\sum_{\mathbf x}e^{i\mathbf P\cdot\mathbf x}
  \langle N(t_s,\mathbf x)\overline N(0)\rangle},
  \label{eq:ratio-equivalent}
\end{equation}
若使用分段 Wilson 线或单独的源归一化，还应乘以由两点函数构造的端点因子。实际实现应根据源汇位置和算符插入归一化选择正确因子；关键是分子和分母必须来自同一组源、同一组规范组态和同一套 Wilson 线。任何“先对源平均再重建比值”的做法都可能丢失源间关联和归一化信息。

\section{双态拟合、求和法与 GEVP}

双态拟合通常把两点和三点函数联合参数化：
\begin{align}
  C_2(t)&=A_0^2e^{-E_0t}
  +A_1^2e^{-E_1t}+\cdots,
  \label{eq:two-state}\\
  C_3(t_s,\tau)&=A_0^2\langle0|O|0\rangle
  e^{-E_0t_s}\left[1+
  r_1e^{-\Delta E(t_s-\tau)}
  +r_1'e^{-\Delta E\tau}
  +\cdots\right].
  \label{eq:three-state}
\end{align}
其中 \(r_1,r_1'\) 概括跃迁矩阵元与振幅比。对胶子算符，跃迁项不一定具有简单的正比结构，故应把 \(r_1,r_1'\) 作为独立拟合参数并检查可辨识性。

求和法定义
\begin{equation}
  S_g(t_s)=\sum_{\tau=\tau_{\min}}^{t_s-\tau_{\min}}
  R_g(\tau,t_s)
  \simeq \text{const}
  +t_s\,\frac{\langle P|O_g|P\rangle}{2E_P}
  +\mathcal O(e^{-\Delta Et_s}),
  \label{eq:summation}
\end{equation}
其斜率给出基态矩阵元，但求和会放大源端和汇端污染，且对 \(t_s\) 窗口敏感。GEVP 则用算符基 \(i,j=1,\dots,N\) 构造矩阵
\begin{equation}
  C_{ij}(t)v_j^{(n)}(t,t_0)
  =\lambda_n(t,t_0)C_{ij}(t_0)v_j^{(n)}(t,t_0),
  \label{eq:gevp}
\end{equation}
并对角化得到有效能量和 overlap。GEVP 的优点是变分压低激发态，缺点是矩阵估计噪声、基算符线性相关和本征向量符号追踪；对纯胶子插入，扩展基还应包含不同 smear、不同场强投影和不同动量。

\section{顺序源与胶子回路的收缩}

标准三点函数可以用顺序源方法降低 Dirac 求逆次数。对核子传播子 \(S_q(x,y)=M_q^{-1}(x,y)\)，在汇点构造带有算符和 Gamma 结构的源，再向插入时间传播。胶子算符是纯规范算符，因此在 Wick 收缩中它与夸克传播子的连接可以组织为胶子回路与核子两点函数的乘积，再加上可能的连接拓扑。典型形式可写成
\begin{equation}
  C_{3,g}(t_s,\tau)
  =\left\langle
  \mathrm{Tr}\left[
  \mathcal G_g(\tau)\,\mathcal C_N(t_s)
  \right]\right\rangle
  -\left\langle\mathcal G_g(\tau)\right\rangle
  \left\langle\mathcal C_N(t_s)\right\rangle ,
  \label{eq:disconnected}
\end{equation}
其中 \(\mathcal G_g\) 是胶子场强、Wilson 线和颜色指标的闭合结构，\(\mathcal C_N\) 是核子两点收缩，第二项是真空减除。若胶子算符还通过显式夸克流连接，必须把连通与断开部分分别保留，不能用一个统一的 `disconnected` 标签覆盖所有拓扑。

非连通贡献的统计误差通常较大，且需要源分辨和全系综真空协方差。一个常见错误是只减除全局平均 \(\langle \mathcal G_g\rangle\)，却忽略每个源、每个时间片和每个组态的真空结构；这会在高动量下产生与物理信号同量级的偏置。

\section{Clover 场强与规范链实现}

在格点上，Clover 场强通常由四个 plaquette 的反对称组合构造：
\begin{equation}
  F_{\mu\nu}(x)
  =-\frac{i}{8a^2}\left[Q_{\mu\nu}(x)-Q_{\nu\mu}(x)\right],
  \label{eq:clover}
\end{equation}
其中 \(Q_{\mu\nu}\) 是围绕 \(x\) 的四个 plaquette 的适当和。对非阿贝尔场强，还需在色空间做无迹投影：
\begin{equation}
  F_{\mu\nu}\longrightarrow
  F_{\mu\nu}-\frac{1}{N_c}\Tr(F_{\mu\nu})\,\mathbf 1 .
  \label{eq:traceless}
\end{equation}
这一步在连续理论中对基本表示恒成立，但在有限格距和非手征作用量中是必要的数值稳定化。

Wilson 线应由同一套格点链接逐段相乘。对伴随表示，可先构造基本表示路径，再通过
\begin{equation}
  U^{(\mathrm{adj})}_{ab}
  =2\Tr\left(T_aU^{(\mathrm{fund})}T_bU^{(\mathrm{fund})\dagger}\right)
  \label{eq:adjointlink}
\end{equation}
做表示转换，或在伴随矩阵空间直接累积链接。两种实现的对易、归一化和反号必须通过小矩阵测试。若沿路径采用 HYP、stout 或 Wilson flow，场强和长 Wilson 线必须使用同一套平滑链接，否则端点处的局域 UV 结构不一致。

\section{传播子、涂抹与谱学}

高动量核子的信噪比可以通过 momentum smearing、distillation、Gaussian smearing 和 GEVP 改善\cite{Bali2016,Peardon2009}。这些技术改变的是插值算符的 overlap，不改变 TMD 算符的定义。一个合格的谱学分析应区分：
\begin{itemize}
  \item 插值算符的基态 overlap 和有效质量平台；
  \item 源分辨两点函数的平移与归一化；
  \item 三点函数的激发态污染；
  \item 胶子算符插入位置 \(\tau\) 的对称性；
  \item 不同源、汇和动量下的 jackknife/bootstrap 相关性。
\end{itemize}
如果两点函数先对源做平均，再把它当作完整时间序列参与比值，可能引入不可恢复的归一化错误。源分辨缓存应保留源轴和每个源的规范组态信息，统计重采样也应在同一组 replica 内完成。

\section{Dirac 方程的物理内容与规范协变性}

核子两点和三点函数的每一个 Wick 收缩最终都归结为在固定规范背景中求夸克传播子。把格点 Dirac 算符记为 \(D_q[U]\)，对给定源 \(\eta\) 求解
\begin{equation}
  \bigl[D_q[U]\psi_\eta\bigr](x,\alpha,a)=\eta(x,\alpha,a),
  \qquad
  S_q(x,y;U)=\bigl[D_q[U]^{-1}\bigr]_{x,y},
  \qquad
  \psi_\eta=D_q[U]^{-1}\eta ,
  \label{eq:diracsolve}
\end{equation}
其中 \(x\) 标记格点位置，\(\alpha\) 与 \(a\) 分别标记 Dirac 和颜色指标。源 \(\eta\) 不是一个与 QCD 无关的任意向量：它编码源的位置、动量投影、自旋-颜色结构、涂抹算子以及可能的 sequential-source 顶点。解 \(\psi_\eta\) 因而不是“某个经典轨迹”，而是该组态下所有满足协变 Dirac 方程的夸克传播振幅之和。单组态的 \(\psi_\eta\) 只是路径积分中的一个条件样本，只有对规范组态平均并按 Wick 收缩组合后，才形成可解释为强子矩阵元的非微扰传播。

\begin{figure}[htbp]
\centering
\begin{tikzpicture}[>=Stealth,font=\small]
  \tikzset{
    panel/.style={draw=blue!35!black,rounded corners,fill=blue!3,minimum width=6.25cm,minimum height=4.4cm},
    vec/.style={draw=blue!60!black,rounded corners,fill=blue!9,align=center,inner sep=4pt},
    op/.style={draw=orange!70!black,rounded corners,fill=orange!12,align=center,inner sep=4pt},
    check/.style={draw=red!65!black,rounded corners,fill=red!7,align=center,inner sep=3pt}
  }
  \node[panel] (left) at (-3.35,0) {};
  \node[panel] (right) at (3.35,0) {};
  \node[font=\bfseries] at (-3.35,1.78) {(a) 规范背景中的传播子};
  \node[font=\bfseries] at (3.35,1.78) {(b) 求解器所见的问题};

  \node[vec] (source) at (-5.15,0.55) {源 \(\eta\)\\[-1mm]{\scriptsize 位置、自旋、颜色}};
  \node[vec] (sink) at (-1.55,0.55) {汇 \(\mathcal O[\psi_\eta]\)\\[-1mm]{\scriptsize 两点/三点收缩}};
  \draw[->,line width=1.2pt,blue!65!black] (source.east) to[bend left=18] node[above] {\(D_q[U]^{-1}\)} (sink.west);
  \draw[->,line width=1.0pt,orange!75!black] (source.east) to[bend right=22] node[below,align=center] {协变传播路径\\{\scriptsize 由 \(U_\mu(x)\) 决定}} (sink.west);
  \draw[->,line width=1.0pt,teal!70!black] (source.east) to[bend left=36] (sink.west);
  \node[op,text width=5.3cm] at (-3.35,-1.15) {不同组态给出不同背景 \(U\)；平均后才形成强子矩阵元};

  \node[vec] (eta) at (1.05,1.05) {\(\eta\)};
  \node[op] (inv) at (3.35,1.05) {\(D_q[U]^{-1}\)};
  \node[vec] (psi) at (5.45,1.05) {\(\psi_i\)};
  \node[vec] (obs) at (5.45,-0.55) {\(\mathcal O[\psi_i]\)};
  \node[check,text width=2.65cm] (res) at (2.75,-0.55) {\(r_i\)、\(r_i^{\rm true}\)\\容差扫描};
  \draw[->] (eta)--(inv); \draw[->] (inv)--(psi); \draw[->] (psi)--(obs);
  \draw[->,dashed,red!65!black] (obs.west)--(res.east);
  \draw[->,dashed,red!65!black] (res.north)--(inv.south);
  \node[align=center,font=\scriptsize] at (4.15,-1.35) {求解失败或容差不足时\\回到 Krylov/预条件分支};
\end{tikzpicture}
\caption{格点 Dirac 求逆的双重图像。左图：源、传播子和汇都生活在同一规范背景中，不同路径由链接变量相干叠加；右图：数值求解只看到非 Hermitian 线性系统、Krylov 近似和残差验收。只有把两幅图对应起来，才能把“解出传播子”与“得到物理矩阵元”区分开。}
\label{fig:dirac-solver}
\end{figure}

图~\ref{fig:dirac-solver} 把这一层物理内容与求解器视角并列展示。左图强调源、链接和汇属于同一个规范场样本；右图强调算法只负责近似逆矩阵的一列。若把右图的小残差直接等同于左图的物理正确性，就会遗漏规范协变、谱条件和观测量敏感度三类问题。

规范变换 \(\Omega(x)\) 同时作用于算符、源和解：
\begin{equation}
  D_q[U^\Omega]=\Omega D_q[U]\Omega^\dagger,
  \qquad
  \eta^\Omega=\Omega\eta,
  \qquad
  \psi_\eta^\Omega=\Omega\psi_\eta .
  \label{eq:diracgaugecov}
\end{equation}
因此迭代求解必须保持颜色协变；如果把源颜色、链接方向和求解器内积的约定混用，残差可以形式上很小，却对应一个改变了规范链或顶点结构的物理量。对 Wilson 型算符，实质量下还满足 \(\gamma_5\)-Hermiticity：
\begin{equation}
  D_q^\dagger=\gamma_5D_q\gamma_5,
  \qquad
  (\gamma_5D_q)^\dagger=\gamma_5D_q .
  \label{eq:gamma5hermiticity}
\end{equation}
这保证 \(D_q\) 与 \(D_q^\dagger\) 通过 \(\gamma_5\) 相似，非实本征值成共轭对出现；但 \(D_q\) 本身通常仍不是 Hermitian 正定矩阵，不能直接套用共轭梯度法。更重要的是，\(D_q\) 的小本征值会放大源和舍入误差；近零模既可能与拓扑背景和近手征模式相关，也会使高动量、大体积传播子的迭代数急剧增加。数值收敛因此不是单纯性能指标，而是传播子谱条件是否足以支撑当前物理误差目标的组成部分。

\section{BiCGSTAB、谱条件与收敛验收}

BiCGSTAB 用非对称 Krylov 子空间处理一般非 Hermitian 的 \(D_q\)，以短递推同时构造双正交残差和局部残差极小化\cite{VanDerVorst1992}。对共轭梯度法可以直接处理的 Hermitian 正定系统，当前 \(D_q\) 并不满足前提；若改用正规方程 \(D_q^\dagger D_q\)，条件数会被平方，因此它不是一个默认更安全的替代。理解 BiCGSTAB 的关键不是记住系数，而是看清每一步在物理传播问题中承担的角色。搜索方向满足
\begin{equation}
  v_{i-1}=D_qp_{i-1},
  \qquad
  \alpha_i=\frac{(\hat r_0,r_{i-1})}{(\hat r_0,v_{i-1})},
  \qquad
  s=r_{i-1}-\alpha_i v_{i-1},
  \label{eq:bicg-alpha}
\end{equation}
其中 \(\alpha_i\) 由双正交约束 \((\hat r_0,s)=0\) 确定。随后令 \(t=Ds\)，局部稳定化系数由
\begin{equation}
  \omega_i
  =\arg\min_{\omega\in\mathbb C}\|s-\omega t\|_2^2
  =\frac{(t,s)}{(t,t)}
  \label{eq:bicg-omega}
\end{equation}
给出，并更新
\begin{equation}
  \psi_i=h+\omega_i s,
  \qquad
  r_i=s-\omega_i t,
  \qquad
  h=\psi_{i-1}+\alpha_i p_{i-1}.
  \label{eq:bicg-update}
\end{equation}
式~\eqref{eq:bicg-omega} 是一个局部最小二乘投影，不是物理时间演化，也不保证每一步都使全局残差单调下降；它只压低当前 \(s\) 在 \(D_qs\) 方向上的剩余分量。把这一点误读为“残差必然下降”，会导致只看迭代曲线而忽略真实矩阵元偏差。

\begin{algorithm}[htbp]
\footnotesize
\linespread{0.95}\selectfont
\caption{Unpreconditioned BiCGSTAB for the lattice Dirac equation \(D[U]\psi=\eta\)}
\label{alg:bicgstab}
\begin{algorithmic}[1]
\Require \(D\equiv D_q[U]\); source \(\eta\); initial guess \(\psi_0\)
\Require relative tolerance \(\epsilon_{\rm rel}>0\); breakdown tolerance \(\epsilon_{\rm br}>0\); \(i_{\max}\)
\Ensure approximate solution \(\psi_i\); residual \(r_i\); convergence status
\State \(r_0\gets \eta-D\psi_0\)
\State choose \(\hat r_0\) with \((\hat r_0,r_0)\ne0\); default \(\hat r_0\gets r_0\)
\State optionally randomize \(\hat r_0\) when the default shadow is unlucky
\State \(\rho_0\gets(\hat r_0,r_0)\); \(p_0\gets r_0\)
\State \(\texttt{status}\gets\texttt{max\_iterations}\)
\If{\(|\rho_0|\le\epsilon_{\rm br}\)}
  \State \(\texttt{status}\gets\texttt{breakdown}\); \textbf{break}
\EndIf
\For{\(i=1,2,\ldots,i_{\max}\)}
  \State \(v\gets Dp_{i-1}\)
  \If{\(|(\hat r_0,v)|\le\epsilon_{\rm br}\)}
    \State \(\texttt{status}\gets\texttt{breakdown}\); \textbf{break}
  \EndIf
  \State \(\alpha\gets\rho_{i-1}/(\hat r_0,v)\)
  \State \(h\gets\psi_{i-1}+\alpha p_{i-1}\)
  \State \(s\gets r_{i-1}-\alpha v\)
  \If{\(\|s\|_2/\|\eta\|_2\le\epsilon_{\rm rel}\)}
    \State \(\psi_i\gets h\); \(r_i\gets s\)
    \State \(\texttt{status}\gets\texttt{converged}\); \textbf{break}
  \EndIf
  \State \(t\gets Ds\)
  \If{\(|(t,t)|\le\epsilon_{\rm br}\)}
    \State \(\texttt{status}\gets\texttt{breakdown}\); \textbf{break}
  \EndIf
  \State \(\omega\gets(t,s)/(t,t)\)
  \State \(\psi_i\gets h+\omega s\)
  \State \(r_i\gets s-\omega t\)
  \If{\(\|r_i\|_2/\|\eta\|_2\le\epsilon_{\rm rel}\)}
    \State \(\texttt{status}\gets\texttt{converged}\); \textbf{break}
  \EndIf
  \State \(\rho_i\gets(\hat r_0,r_i)\)
  \If{\(|\rho_i|\le\epsilon_{\rm br}\) or \(|\omega|\le\epsilon_{\rm br}\)}
    \State \(\texttt{status}\gets\texttt{breakdown}\); \textbf{break}
  \EndIf
  \State \(\beta\gets(\rho_i/\rho_{i-1})(\alpha/\omega)\)
  \State \(p_i\gets r_i+\beta(p_{i-1}-\omega v)\)
\EndFor
\State \textbf{return} \(\psi_i,r_i,\) status
\end{algorithmic}
\end{algorithm}

\begin{figure}[htbp]
\centering
\begin{tikzpicture}[
  >=Stealth,font=\small,node distance=0.55cm,
  block/.style={draw=blue!60!black,rounded corners,fill=blue!8,align=center,text width=3.6cm,inner sep=4pt},
  decision/.style={draw=orange!75!black,fill=orange!12,diamond,aspect=1.75,align=center,inner sep=1pt,text width=2.75cm},
  output/.style={draw=green!55!black,rounded corners,fill=green!10,align=center,text width=2.1cm,inner sep=4pt},
  fail/.style={draw=red!65!black,rounded corners,fill=red!7,align=center,text width=2.5cm,inner sep=3pt}
]
  \node[block] (init) at (0,5.9) {初始化 \(r_0,\hat r_0,p_0,\rho_0\)};
  \node[block] (matvec) at (0,4.45) {计算 \(v=Dp_{i-1}\)};
  \node[block] (alpha) at (0,3.0) {计算 \(\alpha,h,s\)};
  \node[decision] (sconv) at (0,1.35) {\(\|s\|/\|\eta\|\)\\\(\le\epsilon_{\rm rel}\)?};
  \node[block] (tvec) at (0,-0.25) {计算 \(t=Ds\) 与 \(\omega\)};
  \node[block] (update) at (0,-1.9) {更新 \(\psi_i,r_i\)};
  \node[decision] (rconv) at (0,-3.55) {\(\|r_i\|/\|\eta\|\)\\\(\le\epsilon_{\rm rel}\)?};
  \node[block] (next) at (0,-5.2) {更新 \(\rho_i,\beta,p_i\)};
  \node[output] (ok) at (5.0,1.35) {输出 \(\psi_i,r_i\)\\状态：收敛};
  \node[fail] (bad) at (4.0,-5.1) {输出最后迭代\\状态：breakdown\\或达到 \(i_{\max}\)};

  \draw[->] (init)--(matvec);
  \draw[->] (matvec)--(alpha);
  \draw[->] (alpha)--(sconv);
  \draw[->] (sconv)--node[above]{是}(ok);
  \draw[->] (sconv)--node[right]{否}(tvec);
  \draw[->] (tvec)--(update);
  \draw[->] (update)--(rconv);
  \draw[->] (rconv.east) -| node[pos=0.25,above]{是} (ok.south);
  \draw[->] (rconv)--node[right]{否}(next);
  \draw[->,dashed,blue!65!black] (next.west) -- ++(-2.15,0) |- node[pos=0.72,left]{\(i\gets i+1\)} (matvec.west);
  \draw[->,dashed,red!65!black] (matvec.east) -| node[pos=0.18,above]{任一分母/\(\rho_i\)/\(\omega\) 触界} (bad.north);
\end{tikzpicture}
\caption{Unpreconditioned BiCGSTAB 的控制流。蓝色为主迭代路径，橙色菱形为收敛判据，红色虚线为 breakdown 与最大迭代数出口。图与算法~\ref{alg:bicgstab} 一一对应；生产实现可在 \(D\) 两侧加入预条件，但收敛与失败状态的含义不变。}
\label{fig:bicgstab-flow}
\end{figure}

算法~\ref{alg:bicgstab} 刻意写成 unpreconditioned 形式，图~\ref{fig:bicgstab-flow} 则把其中的收敛、回代和失败分支展开为控制流；二者共同保证“迭代步骤”与“物理验收出口”一一对应。生产计算通常还会使用 even-odd、低模 deflation、multigrid 或混合精度预条件。预条件只改变求解路径，不改变方程本身：若 \(D_qM^{-1}y=\eta\)，则必须先得到 \(y\)，再用 \(\psi=M^{-1}y\) 恢复同一物理传播子。伪代码中的 \(\epsilon_{\rm br}\) 只是概念性阈值；实际实现应把它换算成类似 \(|\rho_i|\le\epsilon_{\rm br}\|\hat r_0\|\|r_i\|\) 和 \(\|t\|\le\epsilon_{\rm br}\|s\|\) 的尺度化判据，否则算法会随源归一化或格点体积改变 breakdown 行为。选择 \(\hat r_0\) 的方式会改变双正交投影；默认 \(\hat r_0=r_0\) 简单且可复现，而在正交性偶然失效时随机化 shadow 可降低特定 breakdown 的风险\cite{Schoutrop2022}，但它不能修复病态谱，也不能替代真残差检查。

算法中有三类“小量”不能混为一谈。第一类是递归残差 \(r_i\) 的范数，它决定当前迭代是否满足内部停止条件；第二类是重启后重新计算的真实残差
\begin{equation}
  r_i^{\rm true}=\eta-D_q\psi_i,
  \qquad
  \rho_i^{\rm true}=\frac{\|r_i^{\rm true}\|_2}{\|\eta\|_2},
  \label{eq:trueresidual}
\end{equation}
它用于发现递推关系在长迭代和混合精度下积累的偏差；第三类是物理可观测量的偏差。对观测量 \(O[\psi]\)，在局部线性化意义上可写成
\begin{equation}
  |\Delta O_{\rm solve}|
  \lesssim \kappa_O\,\rho_i^{\rm true}\|\eta\|_2,
  \label{eq:observableerror}
\end{equation}
其中 \(\kappa_O\) 概括该观测量对传播子扰动的敏感度。式~\eqref{eq:observableerror} 说明，固定一个跨项目通用的“容差 \(10^{-n}\)”并不充分：小信号、大时间间隔和高动量三点函数的 \(\kappa_O\) 可以远大于两点函数。合格做法是在目标物理点对 \(\epsilon_{\rm rel}\) 做收敛扫描，并要求求解引起的中心值移动明显小于统计误差；同时比较 BiCGSTAB、预条件 BiCGSTAB 或其他独立求解器的结果。

近零模会使迭代数增加，也使不同源之间看似无关的残差在物理上表现为长程传播。此时应通过低模投影、deflation、multigrid、规范等变预条件或增加迭代预算控制，而不是简单放宽容差。对于单态胶子矩阵元，非连通拓扑对传播子和真空减除都敏感，求解偏差可能与统计涨落同量级；因此求解器状态必须随源、动量、smearing、流时和组态一起进入缓存与误差表。表~\ref{tab:solver-map} 汇总算法变量与物理验收之间的对应关系。

\begin{table}[htbp]
\centering
\caption{BiCGSTAB 变量与格点传播子物理含义、验收证据的对应。}
\label{tab:solver-map}
\small
\begin{tabularx}{\textwidth}{L{2.1cm}L{4.0cm}Y}
\toprule
量 & 数值含义 & 物理/验收解释 \\
\midrule
\(D[U]\), \(\eta\) & 固定背景 Dirac 算符与物理源 & 必须固定作用量、边界条件、Gamma/颜色结构、动量与 smearing \\
\(r_i\) & 当前源残差 & 内部停止量；不能单独证明传播子或矩阵元准确 \\
\(\hat r_0\) & 双正交 shadow 残差 & 影响 Krylov 投影与 breakdown；随机化需记录 seed 和复现策略 \\
\(p_i\) & 搜索方向 & 在规范背景中组合已有传播信息；不具独立物理意义 \\
\(s,\omega_i\) & 稳定化前残差与局部最优系数 & 局部残差压低不等于全局单调收敛；需检查 \(\omega\approx0\) 等 breakdown \\
\(\rho_i^{\rm true}\) & 重启后的真实相对残差 & 必报诊断；与内部 \(r_i\) 一起用于判断递推漂移 \\
求解容差扫描 & 对 \(\epsilon_{\rm rel}\)、预条件和求解器变体重复计算 & 通过可观测量稳定性而非迭代次数宣布质量门 \\
\bottomrule
\end{tabularx}
\end{table}

\section{Distillation、perambulator 与顺序源}

高动量核子的胶子三点函数需要对多个源、多个动量、多个 Wilson 线长度和多个流时重复计算，直接求解传播子的成本很高。Distillation 通过低模投影降低有效自由度。先在每个时间片求解协变 Laplace 算符的低本征模 \(v_x^{(k)}\)，定义蒸馏算符
\begin{equation}
  \Box_{xy}(t)
  =\sum_{k=1}^{N_v}v_x^{(k)}(t)v_y^{(k)\dagger}(t),
  \qquad
  \Box^2=\Box,
  \label{eq:distillation}
\end{equation}
再用
\begin{equation}
  \tau^{(p,p')}_{\alpha\beta}
  =V^{(p)\dagger}
  D^{-1}_{\alpha\beta}
  V^{(p')}
  \label{eq:perambulator}
\end{equation}
表示低模空间中的传播子。动量投影和 momentum smearing 可以写成
\begin{equation}
  \widetilde v_x^{(k)}(\mathbf P)
  =e^{i\mathbf P\cdot\mathbf x}v_x^{(k)},\qquad
  \widetilde f(\mathbf p)\mapsto\widetilde f(\mathbf p-\mathbf P),
  \label{eq:momentumsmearing}
\end{equation}
从而把高动量核子的 overlap 集中到目标动量附近\cite{Bali2016,Peardon2009}。

对三点函数，顺序源方法把汇点核子、插入算符和源端传播子组织成一次或少数几次额外求逆。若 \(S_q\) 是夸克传播子，汇点源可写成包含 \(\gamma_5\)、颜色和自旋投影的向量；把该源向插入时间传播后，再与胶子场强或 Wilson 线收缩。distillation 的广义 perambulator 则把三点函数分解成低模传播子、算符顶点和核子 elemental 的乘积。两种方法的共同物理要求是：
\begin{enumerate}
  \item 源和汇使用同一套规范链接和 smearing；
  \item 胶子算符插入不改变核子插值算符的动量投影；
  \item 低模截断 \(N_v\) 的收敛通过增加 \(N_v\) 或与直接传播子比较验证；
  \item 复用缓存时绑定组态哈希、流时、Wilson 线长度和算符版本；
  \item 不在不同 smearing/流时的矩阵元之间做未经匹配的插值。
\end{enumerate}

\section{统计重采样与协方差}

格点组态之间存在自相关，不能把每个时间片或每个源当作独立样本。设观测量 \(X\) 在 bootstrap replica \(b=1,\dots,N_b\) 上的值为 \(X_b\)，则样本协方差可取
\begin{equation}
  \mathrm{Cov}_{ij}
  =\frac{1}{N_b-1}\sum_{b=1}^{N_b}
  (X_i^b-\overline X_i)(X_j^b-\overline X_j),
  \label{eq:bootstrapcov}
\end{equation}
并用相关 \(\chi^2\) 拟合
\begin{equation}
  \chi^2
  =(\mathbf d-\mathbf m(\theta))^T
  \mathrm{Cov}^{-1}
  (\mathbf d-\mathbf m(\theta)).
  \label{eq:chi2}
\end{equation}
从中间量到最终观测量、再到拟合参数，协方差应沿同一条链传播。对观测量 \(O_a=O_a(X)\) 与参数模型 \(\mathbf d=X\theta+\boldsymbol\epsilon\)，有
\begin{keyformula}{误差传播与广义最小二乘}
\begin{align}
  \mathrm{Cov}_{ab}^{(O)}
  &=\sum_{i,j}
  \frac{\partial O_a}{\partial X_i}
  \mathrm{Cov}_{ij}^{(X)}
  \frac{\partial O_b}{\partial X_j},
  \label{eq:errorpropagation}\\
  \widehat{\theta}
  &=\left(X^T\mathrm{Cov}^{-1}X\right)^{-1}
  X^T\mathrm{Cov}^{-1}\mathbf d,
  \label{eq:gls}\\
  \mathrm{Cov}^{(\theta)}
  &=\left(X^T\mathrm{Cov}^{-1}X\right)^{-1}.
  \label{eq:glscov}
\end{align}
\end{keyformula}

式~\eqref{eq:errorpropagation} 的 Jacobian 必须包含重正化、Fourier 和匹配的导数；若先在中心值上完成这些变换、再只对最终点做线性误差传播，跨 replica 的相关性会丢失。式~\eqref{eq:gls} 与式~\eqref{eq:glscov} 在协方差非奇异时给出标准解；接近奇异时应改用伪逆或正则化，并把截断谱作为系统误差扫描。
当协方差矩阵秩不足或条件数过大时，应使用 SVD、正则化或先做 binning，并把正则化参数作为系统误差扫描。对胶子 TMD，尤其要避免以下三个问题：
\begin{itemize}
  \item 用单一 central fit 冒充 replica 误差；
  \item 对每个 replica 使用不同的数据窗口而不记录选择偏差；
  \item 把拟合失败或 fallback 的 replica 静默替换成中心值。
\end{itemize}
合格报告应同时给出 central fit、replica 分布、失败率、协方差条件数、拟合窗口和每个物理点的支持状态。

\section{坐标空间 Fourier 与大 \texorpdfstring{\(z\)}{z} 尾部}

准 TMD 的 \(x\) 依赖来自 \(z\) 的 Fourier 变换。格点数据只覆盖有限 \(|z|\le z_{\max}\)，直接截断会产生振铃和不物理的端点。常用的处理包括：
\begin{enumerate}
  \item 在短距离使用微扰匹配和比率方案；
  \item 在中等距离保留格点矩阵元；
  \item 对大距离使用满足已知渐近行为的模型或外推；
  \item 用多组截断和尾部模型估计系统误差；
  \item 明确报告 Fourier 约定、归一化和 \(x\) 支撑。
\end{enumerate}
大 \(z\) 外推不是纯粹的技术选择。若尾部模型包含了未测量的 soft/rapidity 或混合贡献，误差会以不可区分的函数形式进入最终 TMD。因此，尾部模型必须与短距离重正化方案分开记录。

\section{Fourier 反演、端点与正定性}

坐标空间 Fourier 反演本质上是不适问题：有限 \(z_{\max}\)、离散格点和统计噪声会在 \(x\) 空间产生高频振荡。一个实用的反演流程应区分三类信息：
\begin{enumerate}
  \item \textbf{被数据约束的区间}：\(|z|\lesssim z_{\rm trusted}\) 内的矩阵元；
  \item \textbf{被外推模型约束的区间}：\(z_{\rm trusted}<|z|\lesssim z_{\max}\)；
  \item \textbf{完全由先验或边界条件决定的区间}：\(|z|>z_{\max}\)。
\end{enumerate}
常用尾部形式包括指数衰减、幂次乘指数、物理渐近约束或在 \(x\) 空间做带先验的展开。无论采用哪种，都应报告：
\begin{itemize}
  \item Fourier 约定：\(e^{+ixP_z z}\) 或 \(e^{-ixP_z z}\)；
  \item 是否使用 \(xg(x)\) 或 \(g(x)\)；
  \item \(z\) 网格的端点和权重；
  \item \(x\to0\) 与 \(x\to1\) 的支撑和端点处理；
  \item 多组 \(z_{\max}\)、尾部模型和窗口选择的结果；
  \item 是否施加正定性、动量求和或 normalization 约束，以及这些约束的影响。
\end{itemize}
胶子分布的 \(x\to0\) 行为尤其敏感于大 \(z\) 尾部，而 \(x\to1\) 行为敏感于短距和端点匹配。把两个区域用同一个平滑模型强行拟合，容易掩盖不同物理来源的系统误差。更稳健的做法是分别报告 \(x\) 可信区间、外推区间和模型依赖区间。

\section{联合外推与误差传播}

胶子 TMD 至少有 \(a,M_\pi,P_z,b_T,\tau,L_\parallel\) 六个控制变量。全联合外推可以写成
\begin{equation}
  f_1^g(x,b_T)
  =f_0(x,b_T)
  +c_a a^2+c_{aP}a^2P_z^2+c_M\frac{M_\pi^2}{\Lambda^2}
  +c_P\frac{\Lambda^2}{P_z^2}
  +c_\tau\tau+c_L e^{-m_\pi L_s}
  +\cdots .
  \label{eq:fullextrap}
\end{equation}
把所有数据点排成向量 \(\mathbf f\), 把式~\eqref{eq:fullextrap} 的基函数排成设计矩阵 \(X\), 即可写成
\begin{equation}
  \mathbf f=X\boldsymbol\theta+\boldsymbol\epsilon,
  \qquad
  X_{r\alpha}=\partial f^{g}_{\,r}/\partial\theta_\alpha .
  \label{eq:extrapolationmatrix}
\end{equation}
式~\eqref{eq:extrapolationmatrix} 的价值不在于把拟合写成矩阵，而在于迫使每个列对应一个明确的物理极限、窗口依赖或可忽略项；如果两列在同一组数据上几乎线性相关，参数就不可辨识，不能靠增加拟合轮数解决。
这样写只是说明误差结构，不是宣称所有系数都可独立拟合。实际应：
\begin{enumerate}
  \item 先在不改变物理定义的变量上做重正化和匹配；
  \item 用 bootstrap 保留同一 replica 内的相关性；
  \item 按尺度层级逐步加入外推项，避免参数数目超过数据自由度；
  \item 对每种拟合模型给出参数协方差、残差和拟合质量；
  \item 采用模型平均或信息准则，而不是选择“看起来最平滑”的结果；
  \item 将统计误差、拟合模型选择和外推区间分开报告。
\end{enumerate}
若只在一个 \(a\)、一个 \(\tau\)、一个 \(P_z\) 和一个 \(m_\pi\) 上计算，式~\eqref{eq:fullextrap} 中的系统项无法辨识。此时应把最终对象称为“该参数点的有效结果”，而不是连续极限物理量。

\section{连续极限、手征极限与无穷动量极限}

最终物理结果要求至少控制三类极限：
\begin{equation}
  a\to0,\qquad M_\pi\to M_\pi^{\rm phys},\qquad P_z\to\infty .
  \label{eq:threelimits}
\end{equation}
在有效范围内，可参数化
\begin{equation}
  f^{g}(x,\bT;a,M_\pi,P_z)
  =f^{g}_{0}(x,\bT)
  +c_a\,a^2+c_M\,M_\pi^2+c_P\,\frac{\Lambda_{\rm QCD}^2}{P_z^2}
  +c_{aP}\,a^2P_z^2+\cdots .
  \label{eq:extrapolation}
\end{equation}
其中 \(a^2P_z^2\) 项反映大动量与粗格距的非交换误差，不能假定它自动消失。若只有单一格距、单一流时或单一价夸克质量，只能报告在该参数点的模型依赖结果，不能声称完成连续物理极限。

\clearpage
\chapter{关联函数、谱学与有限体积}\label{sec:b02-correlators}

本节把强子质量、跃迁矩阵元、激发态控制和有限体积散射组织为同一条可审计链条。核心对象不是某一个拟合函数，而是 Euclidean 关联函数及其在有限体积哈密顿量本征态上的谱表示。不同计算步骤改变的是算符基、源汇结构或统计估计；它们都不改变目标矩阵元的定义。

\section{两点函数、谱表示与有效质量}\label{subsec:b02-two-point}

设源算符 \(O_s\) 在时间 \(0\) 产生指定量子数的态，汇算符 \(O_\Gamma\) 在时间 \(t\) 湮灭该态，并投影到空间动量 \(\mathbf p\)。两点关联函数定义为
\begin{equation}
  C_2(t;\mathbf p)
  =\sum_{\mathbf x}e^{-i\mathbf p\cdot\mathbf x}
  \left\langle O_\Gamma(\mathbf x,t)O_s^\dagger(\mathbf 0,0)\right\rangle .
  \label{eq:b02-two-point}
\end{equation}
在有限体积中插入哈密顿量本征态完备集并取 \(t>0\)，得到
\begin{equation}
  C_2(t;\mathbf p)
  =\sum_n A_n(\mathbf p)e^{-E_n(\mathbf p)t},
  \qquad
  A_n(\mathbf p)
  =\langle 0|O_\Gamma(\mathbf p)|n,\mathbf p\rangle
  \langle n,\mathbf p|O_s^\dagger(\mathbf 0)|0\rangle
  \label{eq:b02-two-point-spectral}
\end{equation}
的前提是源和汇使用同一组态、同一动量投影以及彼此相容的归一化。若源汇相同，\(A_n\) 化为正定的模方；若源汇不同，则必须保留两个 overlap 的乘积。最低能级不是“某个算符的质量”，而是该量子数扇区中与所用算符有非零交叠的最低哈密顿量本征态。

当时间方向长度为 \(N_T\) 且采用周期边界条件时，前向与后向传播同时存在。对自共轭通道，基态贡献可写成
\begin{equation}
  C_2(t)\simeq A_0\left[
  e^{-E_0t}+e^{-E_0(N_T-t)}
  \right],
  \label{eq:b02-cosh}
\end{equation}
其中忽略的是高态和有限体积混合。相应的有效质量不能无条件使用开放边界公式。对周期数据，应在 \(\cosh\) 模型上拟合，或由
\begin{equation}
  \frac{C_2(t)}{C_2(t+1)}
  =
  \frac{\cosh\!\left[E_{\rm eff}(t-N_T/2)\right]}
       {\cosh\!\left[E_{\rm eff}(t+1-N_T/2)\right]}
  \label{eq:b02-periodic-meff}
\end{equation}
数值求解 \(E_{\rm eff}\)。在远离时间边界的平台区，开放边界近似
\begin{equation}
  E_{\rm eff}(t+\tfrac12)
  =\frac1a\ln\frac{C_2(t)}{C_2(t+1)}
  \label{eq:b02-meff}
\end{equation}
才可作为诊断，而不是最终拟合的替代品。平台只说明当前模型、当前源汇和当前误差下高态贡献较小；它不证明算符基已经覆盖了整个谱。

信噪比必须与信号衰减分开讨论。对零动量核子型两点函数，渐近信噪比常写成
\begin{equation}
  \frac{\text{signal}}{\text{noise}}
  \sim e^{-(m_N-\tfrac32m_\pi)t}.
  \label{eq:b02-nucleon-snr}
\end{equation}
这是由信号和规范噪声不同的指数行为造成的。对某些同位旋矢量介子道，相对噪声可远弱于式~\eqref{eq:b02-nucleon-snr}；不能仅凭“图是连通的”就断言信噪比按 \(e^{-m_\pi t}\) 衰减。高动量核子还会受到动量投影、源汇 overlap 和 excited-state contamination 的联合限制。

最简单的介子和核子内插算符分别是
\begin{equation}
  O_{\pi^+}=\bar d^a\gamma_5u^a,
  \qquad
  O_N=\epsilon^{abc}u_a^T C\gamma_5d_b\,P_+u_c,
  \qquad
  P_+=\frac{1+\gamma_4}{2}.
  \label{eq:b02-interpolators}
\end{equation}
色指标必须缩并为单态，Dirac 结构决定 \(J^{PC}\)，味组合决定同位旋和味单态性质。更一般的算符基还可包含不同涂抹半径、协变导数和立方群不可约表示投影。算符数量增加可以改善变分覆盖，但也会放大协方差估计误差；不能把“加入更多算符”当作单调改进。

\section{三点函数、Wick 收缩与矩阵元提取}\label{subsec:b02-three-point}

在源、流插入与汇满足 \(0<\tau<t_s\) 时，三点函数写为
\begin{equation}
  C_3(t_s,\tau;\mathbf p_f,\mathbf p_i)
  =
  \sum_{\mathbf x_f,\mathbf x_J}
  e^{-i\mathbf p_f\cdot\mathbf x_f}
  e^{i(\mathbf p_f-\mathbf p_i)\cdot\mathbf x_J}
  \left\langle
  O_f(\mathbf x_f,t_s)\,
  J(\mathbf x_J,\tau)\,
  O_i^\dagger(\mathbf 0,0)
  \right\rangle .
  \label{eq:b02-three-point}
\end{equation}
插入完备态后，严格的谱展开为
\begin{equation}
  C_3(t_s,\tau)
  =\sum_{m,n}
  \langle 0|O_f|m\rangle
  \langle m|J|n\rangle
  \langle n|O_i^\dagger|0\rangle
  e^{-E_m(t_s-\tau)}e^{-E_n\tau}.
  \label{eq:b02-three-spectral}
\end{equation}
当 \(t_s\to\infty\) 且插入点同时远离源和汇时，\(m=n=0\) 项主导。有限 \(t_s\) 下，最先残留的通常是 \(m=0,n>0\) 或 \(m>0,n=0\) 的跃迁项，而非 \(m,n>0\) 的双高态项；因此只把拟合窗口向大 \(t_s\) 平移而不检查插入点位置，不能保证污染受控。

对同位旋矢量介子，交换两个源汇算符后，连接收缩可以写成
\begin{equation}
  C_{2,\mathrm{conn}}(t)
  =-\operatorname{Tr}\!\left[
  \Gamma_{\rm snk}S_f(t,0)\,
  \Gamma_{\rm src}S_g(0,t)
  \right],
  \label{eq:b02-meson-connected}
\end{equation}
其中 \(S_f=D_f^{-1}\)，迹同时作用于颜色、Dirac 与空间指标。负号来自一个闭合费米子环。对味单态或同一味的双线性插入，还必须加入形如
\begin{equation}
  C_{2,\mathrm{disc}}(t)
  =\sum_q\kappa_q
  \operatorname{Tr}\!\left[\Gamma_{\rm snk}S_q(t,t)\right]
  \operatorname{Tr}\!\left[\Gamma_{\rm src}S_q(0,0)\right]
  \label{eq:b02-meson-disconnected}
\end{equation}
的断开项；系数 \(\kappa_q\) 由算符的味组合决定。精确 \(SU(2)\) 同位旋极限下，同位旋矢量组合中的 \(u,d\) 断开项相消，因此该类通道主要保留连接收缩。该结论不能自动推广到味单态、奇异夸克或胶子算符。

介子连接三点收缩的骨架为
\begin{equation}
  C_{3,\mathrm{conn}}
  =-\operatorname{Tr}\!\left[
  \Gamma_{\rm snk}S_q(t_s,\tau)\,
  \Gamma_JS_q(\tau,0)\,
  \Gamma_{\rm src}S_q(0,t_s)
  \right],
  \label{eq:b02-meson-three-connected}
\end{equation}
具体符号和 Gamma 顺序取决于算符、复共轭以及传播子方向。实现审计应通过小矩阵或已知自由场恒等式验证，而不能只凭代码中负号的位置判断。核子三点函数还需对三个夸克线和 \(\epsilon^{abc}\) 配对作完整反对称化；把核子当成三条互不相关的介子线会漏掉交换项。

胶子插入是一个规范算符。其三点函数通常写成连续测量与真空减除的组合：
\begin{equation}
  C_{3,g}
  =\left\langle
  \mathcal G_g(\tau)\,\mathcal C_N(t_s)
  \right\rangle
  -\left\langle\mathcal G_g(\tau)\right\rangle
  \left\langle\mathcal C_N(t_s)\right\rangle ,
  \label{eq:b02-glue-vacuum}
\end{equation}
其中 \(\mathcal G_g\) 包含场强、Wilson 线和颜色指标，\(\mathcal C_N\) 是核子两点收缩。减除必须与源、时间片和组态结构相容；只减去组态平均全局常数可能留下源相关偏置。

最常用的同动量比值为
\begin{equation}
  R(t_s,\tau)
  =\frac{C_3(t_s,\tau)}{C_2(t_s)}.
  \label{eq:b02-ratio-simple}
\end{equation}
对 \(p_f\neq p_i\) 的形状因子，应改用两端动量匹配的形式
\begin{equation}
  R(t_s,\tau;\mathbf p_f,\mathbf p_i)
  =\frac{C_3(t_s,\tau;\mathbf p_f,\mathbf p_i)}
  {\sqrt{C_2(t_s;\mathbf p_f)C_2(t_s;\mathbf p_i)}}\,
  \mathcal K^{-1},
  \label{eq:b02-ratio-transition}
\end{equation}
其中 \(\mathcal K\) 汇集已知的运动学、源汇重叠和算符归一化因子。若忽略 \(\mathcal K\)，平台值可能只差一个常数，却会被误读为物理 \(Q^2\) 依赖。

在基态加一个有效高态的参数化下，源端和汇端跃迁系数一般不同：
\begin{equation}
  R(t_s,\tau)
  =c_0
  +c_f e^{-\Delta E(t_s-\tau)}
  +c_i e^{-\Delta E\tau}
  +c_{ee}e^{-\Delta E t_s}
  +\cdots .
  \label{eq:b02-two-state-general}
\end{equation}
只有在特定对称性或模型假设下才有 \(c_f=c_i\)。对胶子算符、不同源汇 smear 或不同动量，强制两端共用一个系数会制造模型偏差。拟合应在每个 bootstrap replica 内完成，并扫描 \(\tau\) 排除数、\(t_s^{\min}\) 与高态数目。

顺序源方法把汇端 Dirac 结构与源传播子组织成一个新的汇点源。若 \(D_q\Sigma_q=\eta_{\rm seq}\)，则 \(\eta_{\rm seq}\) 在汇时间片非零，并包含汇动量、自旋投影以及源传播子。把 \(\Sigma_q\) 传回插入时间片后，再与局部流或胶子算符收缩。其成本可概括为
\begin{equation}
  N_{\rm solve}^{3{\rm pt}}
  \simeq N_{\rm source}
  +N_{\rm sink\;momentum}
  +N_{\rm extra},
  \label{eq:b02-sequential-cost}
\end{equation}
其中 \(N_{\rm extra}\) 取决于插入结构、算符基和是否需要独立顺序源。改变初态动量通常只改变投影相因子；改变末态动量、汇投影或汇算符通常要求新的顺序源。因而“顺序源把求逆次数从两次降为一次”只描述同一固定源汇结构的局部优化，不是跨所有动量共享的普遍成本公式。

\section{扩展算符基与广义本征值问题}\label{subsec:b02-gevp}

对 \(N\) 个同量子数算符构造交叉关联矩阵
\begin{equation}
  C_{ij}(t;\mathbf p)
  =\left\langle O_i(t,\mathbf p)O_j^\dagger(0,\mathbf p)\right\rangle
  =\sum_n Z_{in}Z_{jn}^*e^{-E_nt}.
  \label{eq:b02-correlation-matrix}
\end{equation}
GEVP 定义为
\begin{equation}
  C(t)v_n(t,t_0)
  =\lambda_n(t,t_0)C(t_0)v_n(t,t_0),
  \qquad n=1,\dots,N.
  \label{eq:b02-gevp}
\end{equation}
当 \(C(t_0)\) Hermitian 正定时，可作 Cholesky 分解 \(C(t_0)=LL^\dagger\)，并把问题化为标准 Hermitian 问题
\begin{equation}
  \widetilde C(t)
  =L^{-1}C(t)L^{\dagger\,-1},
  \qquad
  \widetilde C(t)u_n=\lambda_nu_n,
  \qquad
  v_n=L^{\dagger\,-1}u_n.
  \label{eq:b02-cholesky}
\end{equation}
这比对 \(C(t_0)\) 显式求逆更适合处理接近线性相关的算符基。有限统计下应先取 Hermitian 部分 \(C_H=(C+C^\dagger)/2\)，但若复关联矩阵的虚部具有物理含义，不能进一步丢弃为纯实矩阵。

若算符基恰好覆盖前 \(N\) 个态且 overlap 矩阵可逆，则
\begin{equation}
  C(t)=ZD(t)Z^\dagger,
  \qquad
  D(t)=\operatorname{diag}(e^{-E_1t},\dots,e^{-E_Nt}),
  \label{eq:b02-z-factorization}
\end{equation}
从而
\begin{equation}
  \lambda_n(t,t_0)=e^{-E_n(t-t_0)}
  \label{eq:b02-gevp-exact}
\end{equation}
没有高态修正。真实系统有无穷多个态，所谓“GEVP 误差”主要是算符基未覆盖的态泄漏。对固定 \(t-t_0\)，可以把主关联函数写成基态加残余项：
\begin{equation}
  \lambda_n(t,t_0)
  =(1-A_n)e^{-E_n(t-t_0)}
  +A_ne^{-E_n'(t-t_0)},
  \label{eq:b02-principal}
\end{equation}
并定义偏离函数
\begin{equation}
  \delta_n(t,t_0)
  =\lambda_n(t,t_0)e^{E_n(t-t_0)}-1
  =A_n\left[e^{-(E_n'-E_n)(t-t_0)}-1\right].
  \label{eq:b02-departure}
\end{equation}
在最低未表示态主导时，固定 \(t-t_0\) 的修正幅度近似按
\begin{equation}
  |\delta_n(t,t_0)|
  \sim
  \exp\!\left[-(E_N-E_n)t_0\right]
  \label{eq:b02-t0-suppression}
\end{equation}
压低，其中 \(E_N\) 表示算符基之外最低相关态的能量。该式给出选择 \(t_0\) 的系统误差下界，但更大的 \(t_0\) 会使 \(C(t_0)\) 的信号更小、相对统计噪声更大，并缩短可用于 \(t>t_0\) 的拟合窗口。

在变分意义上，Rayleigh 商
\begin{equation}
  R(v;t,t_0)
  =\frac{v^\dagger C(t)v}{v^\dagger C(t_0)v}
  \label{eq:b02-rayleigh}
\end{equation}
是各 \(e^{-E_k(t-t_0)}\) 的凸组合，因此最大本征值满足
\begin{equation}
  E_0^{\rm eff}(t,t_0)
  \equiv\frac1{t-t_0}\ln\frac1{\lambda_{\max}(t,t_0)}
  \ge E_0.
  \label{eq:b02-variational-bound}
\end{equation}
有限统计数据不一定显示严格单调下降；平台判据必须与协方差、\(t_0\) 扫描和基大小扫描联合使用。相邻时间片的有效质量采用同一 \(t_0\)：
\begin{equation}
  E_n^{\rm eff}(t,t_0)
  =\frac1a\ln
  \frac{\lambda_n(t,t_0)}{\lambda_n(t+1,t_0)}.
  \label{eq:b02-gevp-meff}
\end{equation}
不要用 \(t_0+1\) 替换分母中的时间片；否则得到的是混合了参考时间变化的量。

标准数值库返回满足 \(v_m^\dagger C(t_0)v_n=\delta_{mn}\) 的 \(B\)-正交归一本征向量。若下游代码又做欧氏归一化，则非对角性可以保持，但对角元一般不再是 1。把本征向量用于三点投影或态追踪前，应显式恢复
\begin{equation}
  \widehat v_n
  =\frac{v_n}{\sqrt{v_n^\dagger C(t_0)v_n}}.
  \label{eq:b02-b-normalization}
\end{equation}
近简并时，按本征值序号追踪态会跳态；可用跨时间或跨体积的重叠矩阵
\begin{equation}
  \mathcal O_{nm}
  =\left|\widehat v_n^\dagger(t)C(t_0)\widehat v_m(t+\Delta t)\right|
  \label{eq:b02-state-tracking}
\end{equation}
识别态的延续。若对角重叠不显著大于非对角项，应报告简并和追踪失败，而不应强行连线。

\paragraph{GEVP 操作顺序.}
输入相关矩阵后，先检查每个时间片和每个 replica 的 Hermitian 性、正定性与条件数；再在固定 \(t_0\) 下求解并排序；随后构造有效质量或主关联函数拟合；最后扫描 \(t_0\)、基大小和拟合窗口。条件化必须记录在案。对高条件数的 replica，可做 SVD 截断或正则化并把结果差异纳入系统误差；静默丢弃这些 replica 会引入选择偏差。

\paragraph{方法回归证据的边界.}
上游固定种子解析谱模型曾比较直接对角化与 GEVP，并得到较大 \(t_0\) 更早收敛的图件；该模型的谱精确性、\(t_0\) 压低和 Cholesky 稳定性可以作为算法回归，但不是真实系综上的物理验证。真实数据的最优 \(t_0\)、大基条件数和带噪态追踪仍需独立研究。

\section{蒸馏、perambulator 与计算结构}\label{subsec:b02-distillation}

蒸馏从三维协变 Laplace 算符的低模开始：
\begin{equation}
  -\nabla^2(t)v_k(t)=\lambda_k(t)v_k(t),
  \qquad
  v_k^\dagger v_l=\delta_{kl}.
  \label{eq:b02-laplace}
\end{equation}
保留 \(N_v\) 个最低模，定义
\begin{equation}
  \Box(t)=V(t)V^\dagger(t),
  \qquad
  V=(v_1,\dots,v_{N_v}),
  \qquad
  \Box^2=\Box.
  \label{eq:b02-distillation-projector}
\end{equation}
低模截断时 \(\Box\) 是投影而非完整酉变换；只有 \(N_v\) 等于空间—颜色—自旋全维数时，\(V\) 才可作为完整基，\(\Box\) 回到单位算符。因而“蒸馏是酉编码”只能在完整基极限下成立，不能用来暗示截断没有信息损失。

在低模空间中定义 perambulator
\begin{equation}
  \tau_{\alpha\beta}^{(p,p')}(t',t)
  =V^{(p)\dagger}(t')D_{\alpha\beta}^{-1}(t',t)V^{(p')}(t).
  \label{eq:b02-perambulator}
\end{equation}
质量、边界条件或规范组态一旦改变，\(\tau\) 必须重新计算。插值算符的几何信息则被压缩进 elemental \(\Phi\)。介子两端函数是 \(\Phi,\tau\) 的迹组合；核子包含三个 perambulator 与两个交换项；三点函数还需广义 perambulator
\begin{equation}
  \mathcal J_{\alpha\beta}(t_f,t,t_i)
  =\left[
  V^\dagger(t_f)D^{-1}(t_f,t)\Gamma(t)D^{-1}(t,t_i)V(t_i)
  \right]_{\alpha\beta}.
  \label{eq:b02-generalized-perambulator}
\end{equation}
它把源侧传播、流顶点和汇侧传播合成一个低模对象。对胶子流，顶点还携带场强、Wilson 线和颜色表示；不能把纯夸克双线性流的 Gamma 结构直接套用。

高动量核子可用带相位的本征矢
\begin{equation}
  \widetilde v_k(\mathbf x,t)
  =e^{i\boldsymbol\zeta\cdot\mathbf x}v_k(\mathbf x,t),
  \qquad
  \boldsymbol\zeta\in\frac{2\pi}{L}\mathbb Z^3
  \label{eq:b02-momentum-smearing}
\end{equation}
把低模基集中到目标动量附近。相位必须是格点允许动量的整数倍，并在源汇两端一致使用。只改变汇端相位会使动量投影和 overlap 不再对应同一个平移不变算符。

蒸馏的收益来自复用和因子化，而不是一个普遍成立的复杂度下降。传播子的求解次数随 \(N_v\)、Dirac 分量和源汇需求增长；收缩费用依两点、三点、核子交换图和非连通环拓扑表现为 \(N_v\) 的不同幂次。若 \(N_v\) 接近全空间维数，预计算、存储和收缩可能比直接方法更昂贵。合格报告至少应给出 \(N_v\)、算符基大小、源数、汇动量数、缓存策略以及与完整基或直接传播子的小体积对照。

\section{有限体积散射、共振与衰变}\label{subsec:b02-finite-volume}

单体强子质量在 \(m_\pi L\gg1\) 时的有限体积修正通常指数小；两体相互作用则可在阈值附近产生 \(O(L^{-3})\) 的能级移动。有限体积在这里不是只应被外推掉的误差，而是散射信息的载体。Maiani--Testa 定理说明，大体积 Euclidean 关联函数不能直接替代非零相对动量的 Minkowski 散射振幅；Lüscher 方法改为测量多个体积下的离散能级。

对质心系中的两个相同粒子，令
\begin{equation}
  q=\frac{kL}{2\pi},
  \qquad
  Z_{00}(s;q^2)
  =\frac{1}{\sqrt{4\pi}}
  \sum_{\mathbf n\in\mathbb Z^3}
  \frac{1}{(|\mathbf n|^2-q^2)^s},
  \label{eq:b02-zeta}
\end{equation}
其中 \(s=1\) 由解析延拓定义。在本书相位约定下，最低 \(A_1^+\) 道的 s 波量子化条件可写为
\begin{equation}
  \tan\delta_0(k)
  =-\frac{\pi^{3/2}q}{Z_{00}(1;q^2)},
  \qquad
  W_n=2\sqrt{m^2+k_n^2},
  \label{eq:b02-luscher-swave}
\end{equation}
并对反三角函数分支作连续展开。不同文献可能把负号写在 \(\phi\) 或 \(\delta\) 的定义中；比较相移前必须统一约定。连续展开不是数学装饰，共振附近若选错分支会把 \(\delta\) 的平台误读。

阈值附近，按本批次 \(a_0=\lim_{k\to0}\delta_0(k)/k\) 的约定，
\begin{equation}
  W
  =2M_\pi
  -\frac{4\pi a_0}{M_\pi L^3}
  \left[
  1+c_1\frac{a_0}{L}
  +c_2\frac{a_0^2}{L^2}
  +O(L^{-3})
  \right]
  +O(L^{-6}),
  \label{eq:b02-threshold}
\end{equation}
其中 \(c_1=-2.837297\)、\(c_2=6.375183\) 是立方几何常数。若采用相反散射长度符号，整项符号也随之改变。该展开要求近阈值、短程相互作用和单道主导；不能用它描述远离阈值的宽共振。

对 \(\rho\to\pi\pi\) 一类的 \(p\) 波共振，需要非零总动量和小群不可约表示来区分分波。典型流程是先用扩展算符基和 GEVP 提取 \(I=1\)、\(J^{PC}=1^{--}\) 道的 \(W_n(L)\)，再由运动框架 Lüscher 关系得到 \(\delta_1(k_n)\)，最后拟合
\begin{equation}
  \frac{k^3}{W}\cot\delta_1(k)
  =\frac{6\pi}{g_{\rho\pi\pi}^2}
  \left(m_\rho^2-W^2\right),
  \qquad
  W=2\sqrt{m_\pi^2+k^2}.
  \label{eq:b02-rho-bw}
\end{equation}
由截距和斜率可提取 \(m_\rho\) 与 \(g_{\rho\pi\pi}\)，再由
\begin{equation}
  \Gamma_\rho
  =\frac{g_{\rho\pi\pi}^2}{6\pi m_\rho^2}k_\rho^3,
  \qquad
  k_\rho=\sqrt{m_\rho^2/4-m_\pi^2}
  \label{eq:b02-rho-width}
\end{equation}
得到宽度。这里的 \(W\) 是双粒子不变能量，不是简单的单体有效质量平台。

弱衰变还多一层归一化问题。Lellouch--Lüscher 关系把有限体积矩阵元 \(A_L\) 与无穷体积振幅联系起来：
\begin{equation}
  \left|A_\infty(k)\right|^2
  =\mathcal F(k,L)\left|A_L\right|^2,
  \label{eq:b02-ll}
\end{equation}
其中 \(\mathcal F\) 由同一道的相移与 Lüscher 函数对动量的导数结构确定。精确表达式依赖 \(K\to\pi\pi\) 道、总动量、irrep、相空间和态归一化，不能把某个单道公式当作普适常数。

有限体积分析的主要失效边界是：超过非弹性阈值仍只用单道公式；宽共振使有限体积能级无法分离；\(m_\pi L\) 太小导致指数修正不可忽略；态追踪在 avoided crossing 处断线；以及相移分支或散射长度符号未统一。出现这些情况时，应降级为“有限体积能谱结果”或模型依赖的共振参数，而不是宣称已得到无穷体积振幅。

\section{形状因子、衰变常数与软函数}\label{subsec:b02-form-factors}

强子形状因子是三点函数最直接的物理出口。对 \(\pi^+\) 的电磁流，
\begin{equation}
  \left\langle\pi^+(p_f)|V_\mu|\pi^+(p_i)\right\rangle
  =(p_f+p_i)_\mu F_\pi(Q^2),
  \qquad
  Q^2=-(p_f-p_i)^2.
  \label{eq:b02-pion-ff}
\end{equation}
在 \(Q^2\to0\) 附近，
\begin{equation}
  F_\pi(Q^2)
  =1-\frac16\langle r_\pi^2\rangle Q^2+O(Q^4),
  \label{eq:b02-charge-radius}
\end{equation}
因此电荷半径依赖小 \(Q^2\) 端点数据的可靠外推。格点动量只能取离散值，靠近零点的插值必须与连续极限和有限体积误差一起报告。

对自旋 \(1/2\) 核子，电磁矩阵元可分解为
\begin{equation}
  \left\langle N(p_f)|V_\mu|N(p_i)\right\rangle
  =\bar u(p_f)\left[
  \gamma_\mu F_1(Q^2)
  +\frac{i\sigma_{\mu\nu}q^\nu}{2m_N}F_2(Q^2)
  \right]u(p_i),
  \label{eq:b02-nucleon-ff}
\end{equation}
并定义
\begin{equation}
  G_E(Q^2)=F_1(Q^2)-\frac{Q^2}{4m_N^2}F_2(Q^2),
  \qquad
  G_M(Q^2)=F_1(Q^2)+F_2(Q^2).
  \label{eq:b02-sachs}
\end{equation}
Sachs 公式中的相对符号由 \(q=p_f-p_i\)、\(\sigma_{\mu\nu}\) 和 metric 约定共同决定；换用相反度规时不能只改一个符号。轴矢结构可写为
\begin{equation}
  \left\langle N(p_f)|A_\mu^a|N(p_i)\right\rangle
  =\bar u(p_f)\left[
  \gamma_\mu\gamma_5G_A(Q^2)
  +\frac{q_\mu}{2m_N}\gamma_5G_P(Q^2)
  \right]\frac{\tau^a}{2}u(p_i).
  \label{eq:b02-axial-ff}
\end{equation}
\(G_A(0)\)、诱导赝标量项和 PCAC 关系必须使用同一套流重整化与手征约定。

赝标量介子衰变常数的两种常见归一化为
\begin{equation}
  \left\langle0|A_\mu^a(0)|\pi^b(P)\right\rangle
  =i\delta^{ab}F_\pi P_\mu,
  \qquad
  f_\pi=\sqrt2\,F_\pi,
  \label{eq:b02-decay-constant}
\end{equation}
其中 \(F_\pi\simeq92.4\) MeV 与 \(f_\pi\simeq130.4\) MeV 是同一物理量的不同约定。纯轻子宽度在采用 \(F_\pi\) 时写成
\begin{equation}
  \Gamma(\pi^+\to\ell^+\nu_\ell)
  =\frac{G_F^2}{8\pi}|V_{ud}|^2F_\pi^2
  m_\ell^2m_\pi
  \left(1-\frac{m_\ell^2}{m_\pi^2}\right)^2.
  \label{eq:b02-pion-width}
\end{equation}
如果不先统一 \(F_\pi/f_\pi\)，由宽度反推 CKM 元会系统性偏差 \(\sqrt2\)。

中性介子混合的 bag 参数可定义为
\begin{equation}
  \left\langle\bar K^0|O_{LL}|K^0\right\rangle
  =\frac83M_K^2F_K^2B_K(\mu),
  \qquad
  O_{LL}=(\bar s\gamma_\mu(1-\gamma_5)d)^2.
  \label{eq:b02-bk}
\end{equation}
\(B_K\) 是重整化方案和标度相关的矩阵元参数，不能与 scheme-independent 的物理振幅混写。四夸克算符混合、手征性压低的矩阵元和幂次发散减除都必须显式进入误差预算。

大动量介子形状因子提供另一条 TMD 软函数通道。四夸克连通矩阵元可因子化为
\begin{equation}
  F(b_\perp,P^z)
  =S_I(b_\perp)
  \int_0^1\!\mathrm dx\,\mathrm dx'\,
  H(x,x',P^z)
  \Phi^\dagger(x',b_\perp,-P^z)
  \Phi(x,b_\perp,P^z)
  +O\!\left(\frac1{(P^z)^2}\right),
  \label{eq:b02-soft-factorization}
\end{equation}
其中 \(S_I\) 是内禀软函数，\(H\) 是硬核。只有在指定 Fierz 基、波函数归一化和领头阶硬核时，才可写
\begin{equation}
  S_I(b_\perp)
  =2N_c\frac{F(b_\perp,P^z)}
  {\left|\phi(0,b_\perp,P^z)\right|^2}
  +O\!\left(\alpha_s,\frac1{(P^z)^2}\right),
  \qquad
  \phi(0,b_\perp,P^z)=\int_0^1\!\mathrm dx\,\Phi(x,b_\perp,P^z).
  \label{eq:b02-soft-lo}
\end{equation}
该式要求连通插入、特定算符组合和足够大的 \(P^z\)。把同一比值直接用于胶子软函数、非连通插入或不同 Fierz 结构会产生未控制的匹配误差。

\begin{table}[htbp]
\centering
\small
\caption{B02 形状因子与衰变观测量所需的最小条件。}
\label{tab:b02-form-factor-conditions}
\begin{tabularx}{\textwidth}{>{\raggedright\arraybackslash}p{2.7cm}>{\raggedright\arraybackslash}X>{\raggedright\arraybackslash}p{3.7cm}}
\toprule
对象 & 物理定义与关键输入 & 主要失效模式 \\
\midrule
\(\pi\) 电磁形状因子 & 守恒流、两态归一化、\(Q^2\to0\) 外推 & 源汇动量不匹配、流重整化错误、端点外推偏差 \\
核子 \(G_E,G_M\) & \(F_1,F_2\)、Sachs 变换、物理 \(m_\pi\) 极限 & 高 \(Q^2\) 信噪比、激发态污染、非连通奇异插入遗漏 \\
\(G_A,G_P\) & 轴矢流、PCAC、手征与重整化约定 & 流归一化、赝标量极点、基态重叠 \\
\(F_\pi,F_K\) & 真空—介子矩阵元、\(F_\pi/f_\pi\) 约定 & 约定混用、连续/手征外推、有限体积 \\
\(B_K\) & 四夸克算符混合、方案相关 bag 参数 & 幂次发散减除、手征性混合、标度混用 \\
\(S_I(b_\perp)\) & 四夸克连通形状因子、Fierz 基、硬核与波函数归一化 & 非连通插入、高扭度、\(P^z\) 平台不足 \\
\bottomrule
\end{tabularx}
\end{table}


\clearpage
\chapter{稀疏矩阵、Krylov算法与收缩复杂度}
\label{sec:krylov-algorithms}

\section{Dirac算符作为 matrix-free 对象}

格点Dirac算符 \(D[U]\) 的维数为
\begin{equation}
  N_D=N_cN_sN_{\rm vol},
  \qquad
  N_{\rm vol}=\frac{T}{a}\left(\frac{L}{a}\right)^3,
  \label{eq:dirac-dimension}
\end{equation}
其中 \(N_c=3\) 是色维数，\(N_s=4\) 是旋量维数。一个 \(32^3\times64\) 的格点已经有 \(N_D\simeq2.1\times10^7\)。虽然矩阵极端稀疏，显式存储仍不现实，因此实际求解器只提供矩阵乘向量（matvec）：
\begin{equation}
  y\gets D[U]x.
  \label{eq:matvec}
\end{equation}
这称为 matrix-free 实现。任何优化都不能把整个矩阵稠密化；算法设计、缓存布局、通信和预条件都必须围绕 matvec 成本展开。

\section{谱、条件数与误差放大}

对可逆矩阵 \(D\)，条件数定义为
\begin{equation}
  \kappa_p(D)=\|D\|_p\|D^{-1}\|_p,
  \label{eq:condition-number}
\end{equation}
若输入有相对误差 \(\epsilon_{\rm in}\)，解的最坏相对误差可被放大到
\begin{equation}
  \frac{\|\delta x\|}{\|x\|}
  \lesssim\kappa_p(D)\frac{\|\delta b\|}{\|b\|}.
  \label{eq:condition-amplification}
\end{equation}
轻夸克质量、非零化学势、粗格距和大体积都会让 \(\kappa\) 增大。近零模并不意味着矩阵不可逆，但会让求解时间和残差估计变得极不稳定。

对Wilson或Clover费米子，常用 \(\gamma_5\)-Hermiticity 把问题重写为 Hermitian 形式：
\begin{equation}
  A=\gamma_5 D,
  \qquad A^\dagger=A,
  \qquad A^2=D^\dagger D.
  \label{eq:gamma5-hermitian}
\end{equation}
若 \(D\) 在正质量区没有零模，则可用共轭梯度法直接求解 \(Ax=\gamma_5b\)。即使使用BiCGSTAB，也可利用 \(\gamma_5 D\) 的谱信息设计预条件。代码中的 Hermiticity 必须由 \(D\) 的实际边界条件和内积 \((x,y)=\sum x^\dagger y\) 决定，不能从连续公式直接假定。

\section{Krylov子空间与投影}

给定初始向量 \(x_0\) 和残差 \(r_0=b-Dx_0\)，第 \(m\) 个Krylov子空间为
\begin{equation}
  \mathcal K_m(D,r_0)
  =\operatorname{span}\{r_0,Dr_0,\ldots,D^{m-1}r_0\}.
  \label{eq:krylov-space}
\end{equation}
Krylov方法在 \(x_0+\mathcal K_m\) 中寻找近似解，不同方法只是选择不同投影条件：
\begin{align}
  \text{CG: }&r_m\perp\mathcal K_m,\\
  \text{GMRES: }&r_m\perp D\mathcal K_m,\\
  \text{BiCG: }&r_m\perp\mathcal K_m(D^\dagger,\tilde r_0).
\end{align}
理解投影条件比背递推公式更重要。对于Hermitian正定矩阵，CG最小化能量范数；对于非Hermitian矩阵，GMRES最小化残差范数，但存储所有基向量；BiCGSTAB用短递推降低存储，却不保证残差单调下降。

\section{共轭梯度法}

对Hermitian正定 \(A\)，CG从 \(x_0\) 出发，令 \(r_0=b-Ax_0\)、\(p_0=r_0\)。每一步满足
\begin{align}
  \alpha_k&=\frac{r_k^\dagger r_k}{p_k^\dagger Ap_k},
  \label{eq:cg-alpha}\\
  x_{k+1}&=x_k+\alpha_kp_k,\\
  r_{k+1}&=r_k-\alpha_kAp_k,\\
  \beta_k&=\frac{r_{k+1}^\dagger r_{k+1}}{r_k^\dagger r_k},\\
  p_{k+1}&=r_{k+1}+\beta_kp_k.
  \label{eq:cg-beta}
\end{align}
算法~\ref{alg:cg} 给出可执行结构。真正的停止条件不能只看递推残差；应周期性重算 \(b-Ax\)，以避免有限精度导致的假收敛。

\begin{algorithm}[htbp]
\caption{带显式真残差检查的共轭梯度法。}
\label{alg:cg}
\begin{algorithmic}[1]
\Require Hermitian正定算子 \(A\)，右端 \(b\)，初值 \(x_0\)，容差 \(\varepsilon\)
\Ensure 近似解 \(x\) 和真实残差 \(\|b-Ax\|_2\)
\State \(r_0\gets b-Ax_0\)，\(p_0\gets r_0\)，\(k\gets0\)
\While{\(\|r_k\|_2>\varepsilon\) 且未超过预算}
  \State \(Ap\gets A[p_k]\)
  \State \(\alpha_k\gets(r_k^\dagger r_k)/(p_k^\dagger Ap)\)
  \State \(x_{k+1}\gets x_k+\alpha_kp_k\)
  \State \(r_{k+1}\gets r_k-\alpha_kAp\)
  \State 每 \(m\) 步重算 \(r_{k+1}\gets b-Ax_{k+1}\)
  \State \(\beta_k\gets(r_{k+1}^\dagger r_{k+1})/(r_k^\dagger r_k)\)
  \State \(p_{k+1}\gets r_{k+1}+\beta_kp_k\)
  \State \(k\gets k+1\)
\EndWhile
\State 输出 \(x_k\)、matvec次数和真残差
\end{algorithmic}
\end{algorithm}

CG的收敛率在理想情况下受
\begin{equation}
  \frac{\|x_k-x\|_A}{\|x_0-x\|_A}
  \le2\left(
  \frac{\sqrt\kappa-1}{\sqrt\kappa+1}
  \right)^k
  \label{eq:cg-rate}
\end{equation}
控制。该界对谱聚集很悲观；若少数大本征值被deflation移除，剩余谱的有效条件数会显著降低。

\section{GMRES与Arnoldi基}

对非Hermitian \(D\)，GMRES构造正交基 \(v_1,\ldots,v_m\)，满足Arnoldi关系
\begin{equation}
  DV_m=V_{m+1}\overline H_m,
  \label{eq:arnoldi}
\end{equation}
其中 \(V_m=[v_1,\ldots,v_m]\)，\(\overline H_m\) 是上Hessenberg矩阵。近似解写成
\begin{equation}
  x_m=x_0+V_my_m,
  \qquad
  y_m=\arg\min_y\|\beta e_1-\overline H_my\|_2,
  \label{eq:gmres-small}
\end{equation}
其中 \(\beta=\|r_0\|_2\)。小Hessenberg最小二乘可以用Givens旋转递推更新。GMRES的残差单调不增，但每次迭代需保存基向量，内存为 \(\mathcal O(mN_D)\)；重启GMRES把 \(m\) 限制在较小值，却可能丢失收敛信息。

\begin{algorithm}[htbp]
\caption{重启GMRES的概念流程。}
\label{alg:gmres}
\begin{algorithmic}[1]
\Require 非Hermitian算子 \(D\)，右端 \(b\)，初值 \(x_0\)，重启长度 \(m\)
\Ensure 近似解 \(x\) 与真残差
\State \(r_0\gets b-Dx_0\)，\(\beta\gets\|r_0\|_2\)，\(v_1\gets r_0/\beta\)
\For{\(j=1,\ldots,m\)}
  \State \(w\gets Dv_j\)
  \For{\(i=1,\ldots,j\)}
    \State \(h_{ij}\gets v_i^\dagger w\)，\(w\gets w-h_{ij}v_i\)
  \EndFor
  \State \(h_{j+1,j}\gets\|w\|_2\)，\(v_{j+1}\gets w/h_{j+1,j}\)
\EndFor
\State 用最小二乘求 \(y_m\)，更新 \(x_m=x_0+V_my_m\)
\If{残差未收敛}
  \State 以 \(x_m\) 为新初值重启，并检查基向量是否退化
\EndIf
\State 输出解、内存占用和真残差
\end{algorithmic}
\end{algorithm}

\section{BiCGSTAB的推导与breakdown}

BiCG通过双Krylov子空间同时处理 \(D\) 与 \(D^\dagger\)，但需要两次矩阵乘且递推可能breakdown。BiCGSTAB把BiCG残差再用局部最小残差多项式光滑。给定影子残差 \(\hat r_0\)，常见迭代为
\begin{align}
  \rho_i&=\hat r_0^\dagger r_{i-1},\\
  \beta&=\frac{\rho_i}{\rho_{i-1}}\frac{\alpha}{\omega},\\
  p_i&=r_{i-1}+\beta(p_{i-1}-\omega v),\\
  v&=Dp_i,\\
  \alpha&=\frac{\rho_i}{\hat r_0^\dagger v},\\
  s&=r_{i-1}-\alpha v,\\
  t&=Ds,\\
  \omega&=\frac{t^\dagger s}{t^\dagger t}.
  \label{eq:bicgstab-kernel}
\end{align}
当 \(\hat r_0^\dagger v\approx0\)、\(\rho_i\approx0\) 或 \(t^\dagger t\approx0\) 时会发生breakdown。随机化 \(\hat r_0\) 可降低结构化breakdown概率，但不是对所有系统的保证。若 \(\omega\approx0\)，局部最小残差步骤失效，应切换到GMRES、改变预条件或重新启动。

BiCGSTAB每步通常需要两次matvec和少量全局内积。其成本低于完整GMRES，但对条件数、预条件质量和谱分布仍很敏感。残差曲线出现平台并不一定是算法“收敛慢”；应检查真残差、\(\rho\)、\(\omega\)、向量范数和NaN/Inf。

\section{预条件：把难点移到可解的方向}

预条件子 \(M^{-1}\) 求解等价系统
\begin{equation}
  M^{-1}Dx=M^{-1}b.
  \label{eq:preconditioned-system}
\end{equation}
理想预条件满足 \(M\approx D\) 且 \(M^{-1}\) 便宜。常用策略包括：
\begin{itemize}
  \item 偶奇预条件：利用时间/空间棋盘结构把部分耦合块对角化；
  \item 质量预条件：把轻夸克问题嵌入较重质量系统；
  \item 不完全分解或结构化近似：降低高频误差；
  \item 低模deflation：显式移除接近零的本征向量；
  \item multigrid：在不同格距层级处理低频误差；
  \item 混合精度：外层高精度、内层低精度。
\end{itemize}
预条件改变的是计算路径，不改变物理方程。比较不同预条件时必须报告matvec次数、全局同步、内存和最终真残差，而不能只比较墙钟时间。

\section{Deflation与近零模}

假设已知 \(k\) 个小本征值对应的正交基 \(U=[u_1,\ldots,u_k]\)，且 \(U^\dagger U=I\)。投影算子
\begin{equation}
  P_\perp=I-UU^\dagger
  \label{eq:deflation-projector}
\end{equation}
可把小模从迭代空间中移除。求解 \(Dx=b\) 可以拆成粗空间修正与投影后系统：
\begin{equation}
  x=Uy+P_\perp z,
  \qquad
  P_\perp D P_\perp z=P_\perp b.
  \label{eq:deflation-solve}
\end{equation}
粗空间系数通过
\begin{equation}
  (U^\dagger DU)y=U^\dagger b
  \label{eq:deflation-coarse}
\end{equation}
求解。实际实现需要对小 \(k\times k\) 矩阵做奇异值检查。deflation对TMD大 \(z\) 和近零模尤其重要，但低模基底必须由同一组态、同一Dirac算符和同一边界条件产生；借用量子数不匹配的低模会破坏规范协变性和矩阵元归一化。

\begin{algorithm}[htbp]
\caption{基于低模投影的deflation概念流程。}
\label{alg:deflation}
\begin{algorithmic}[1]
\Require 算子 \(D\)，右端 \(b\)，低模基 \(U\)，粗空间容差
\Ensure 近似解 \(x\)
\State 正交化 \(U\)，检查 \(\|U^\dagger U-I\|\)
\State 构造 \(D_{UU}\gets U^\dagger DU\)
\If{\(D_{UU}\) 奇异或条件数超限}
  \State 拒绝该低模基底并回退到未deflation求解
\EndIf
\State 求粗解 \(y\gets D_{UU}^{-1}U^\dagger b\)
\State 求投影残差 \(b_\perp\gets b-DUy\)
\State 在 \(P_\perp\) 空间迭代求解 \(P_\perp DP_\perp z=b_\perp\)
\State \(x\gets Uy+P_\perp z\)，并重算真残差
\end{algorithmic}
\end{algorithm}

\section{Multigrid与低频误差}

Krylov方法擅长消除高频误差，但低频光滑误差在细格点上传播很慢。multigrid通过粗化格点、限制残差、粗层求解和延拓修正来加速：
\begin{equation}
  r_H=Rr_h,
  \qquad
  D_He_H=r_H,
  \qquad
  x_h\gets x_h+Pe_H,
  \label{eq:multigrid-cycle}
\end{equation}
其中 \(R\) 是restriction，\(P\) 是prolongation。规范场上的自适应粗化必须尊重拓扑和规范协变性；不能简单把链接平均后当作粗层链接。multigrid设置不当会在粗层引入非物理零模，导致外层迭代振荡或提前收敛到错误解。

\begin{algorithm}[htbp]
\caption{两层multigrid V-cycle的概念结构。}
\label{alg:multigrid}
\begin{algorithmic}[1]
\Require 细层系统 \(D_hx_h=b_h\)，粗化算子 \(R,P\)，光滑步数 \(n_1,n_2\)
\Ensure 修正后的细层近似解
\State 执行 \(n_1\) 次细层光滑，得到 \(x_h\)
\State 计算细层残差 \(r_h\gets b_h-D_hx_h\)
\State 限制 \(r_H\gets Rr_h\)
\State 在粗层近似求解 \(D_He_H=r_H\)
\State 延拓并修正 \(x_h\gets x_h+Pe_H\)
\State 执行 \(n_2\) 次细层光滑并重算真残差
\end{algorithmic}
\end{algorithm}

\section{混合精度与全局同步}

GPU上的双精度矩阵乘通常受内存带宽和全局同步限制。混合精度求解器在低精度中执行大部分matvec和内积，在外层用高精度检查残差：
\begin{equation}
  \frac{\|\delta x\|}{\|x\|}
  \lesssim\kappa(D)\left(\epsilon_{\rm low}+
  \frac{\|r_{\rm true}\|}{\|b\|}\right).
  \label{eq:mixed-precision-bound}
\end{equation}
低精度并非只是“更快”；若条件数大，低精度递归会污染Krylov基。常见策略是使用补偿内积、周期性高精度重启或分层精度。性能报告至少包含精度格式、全局归约次数、matvec次数和真残差。

\section{随机估计与噪声源}

当不能计算所有传播子时，常用随机源近似迹或全模式平均。设零均值随机向量 \(\xi\) 满足
\begin{equation}
  \mathbb E[\xi_i]=0,
  \qquad
  \mathbb E[\xi_i^\dagger\xi_j]=N\delta_{ij},
  \label{eq:random-source}
\end{equation}
则
\begin{equation}
  \mathbb E[\xi^\dagger D^{-1}\xi]
  =N\,\Tr D^{-1}.
  \label{eq:random-trace}
\end{equation}
估计方差依赖源的谱和 \(D^{-1}\) 的非对角结构。Z2、U(1)、时空稀释和低模deflation可以降低方差，但任何一种源都必须与物理归一化、反粒子方向和Fourier投影一起验证。随机估计的误差是统计误差的一部分，不能混入系统误差后重复相加。

\section{顺序源、传播子复用与收缩图}

三点函数计算的核心是选择哪些传播子复用。若末态动量 \(\mathbf p_f\) 和源 \(\eta\) 已知，顺序源满足
\begin{equation}
  D\Sigma=e^{-i\mathbf p_f\cdot x}\gamma_5S,
  \label{eq:seq-source-krylov}
\end{equation}
其中汇点时间片以外的右端为零。每个新的末态动量通常需要新的顺序源，但多个初始动量可以通过改变投影相因子复用同一传播子。收缩图的成本可以抽象为
\begin{equation}
  C_{\rm contract}
  =\mathcal O\left(
  N_{\rm src}N_{\rm sink}N_{\rm op}N_v^p
  \right),
  \label{eq:contraction-cost}
\end{equation}
其中 \(p\) 对连通两点、非连通环和重子三点分别取不同值。报告应给出 \(N_v\)、源数、汇点动量数、算符插入数、每一步是否缓存，而不是只给一个总GPU时间。

\section{复杂度与验证矩阵}

\begin{table}[htbp]
\centering
\caption{常见Dirac求解器的粗粒度复杂度与适用域。}
\label{tab:krylov-comparison}
\small
\begin{tabularx}{\textwidth}{L{2.1cm}L{2.0cm}L{2.6cm}Y}
\toprule
方法 & 每次迭代matvec & 存储 & 适用与风险 \\
\midrule
CG & 1 & \(O(N_D)\) & Hermitian正定；条件数大时收敛慢 \\
BiCGSTAB & 2 & \(O(N_D)\) & 非Hermitian；可能breakdown或残差振荡 \\
GMRES(\(m\)) & 1 & \(O(mN_D)\) & 残差下降平滑；重启可能停滞 \\
Deflated Krylov & 1--2 & 基础 + \(O(kN_D)\) & 近零模明显；低模质量和成本需验证 \\
Multigrid & 层级matvec & 多层级向量 & 低频误差主导；粗化需保持规范结构 \\
\bottomrule
\end{tabularx}
\end{table}

\begin{checkpointbox}{求解器验收清单}
\begin{enumerate}
  \item 是否报告真实残差而不只是递推残差；
  \item 是否记录breakdown、重启、NaN和最大迭代；
  \item 是否检查 \(\gamma_5\)-Hermiticity、边界条件和源归一化；
  \item 是否区分求解误差、随机源误差和统计误差；
  \item 是否在小规模系统上与稠密解或已知解析解对照；
  \item 是否在目标GPU上报告matvec次数、内存、通信和精度。
\end{enumerate}
\end{checkpointbox}

\clearpage
\chapter{统计推断、极限外推与逆问题}
\label{sec:statistics-inverse}

\section{从组态到估计量}

格点数据由有限组态组成。一个观测量 \(O\) 的估计量记作 \(\hat O\)，它本身是随机变量。报告均值 \(\bar O\) 之前，应明确：
\begin{enumerate}
  \item 样本单位是组态、轨迹还是独立bootstrap replica；
  \item 是否对每个组态内部先做源平均或动量平均；
  \item 是否存在同一组态内不可忽略的相关性；
  \item 是否有选择性地丢弃了“坏组态”，选择准则是什么；
  \item 误差是否来自同一组态相关样本，而不是独立误差的平方和。
\end{enumerate}
若在组态之间先平均再对源误差独立处理，组态相关和源相关的交叉项常被低估。正确顺序是先定义可重采样的底层单位，再在该单位内执行完整分析。

\section{自相关、积分时间与有效样本}

对平稳序列 \(O_t\)，自相关函数为
\begin{equation}
  \rho_O(k)=\frac{\langle(O_t-\bar O)(O_{t+k}-\bar O)\rangle}
  {\langle(O_t-\bar O)^2\rangle}.
  \label{eq:autocorrelation}
\end{equation}
积分自相关时间为
\begin{equation}
  \tau_{\rm int}=\frac12+\sum_{k=1}^{\infty}\rho_O(k),
  \qquad
  N_{\rm eff}=\frac{N}{2\tau_{\rm int}}.
  \label{eq:tau-int}
\end{equation}
当自相关长度与可用样本数同量级时，简单截断求和会严重低估误差。应比较不同自相关窗口、binning大小和jackknife-block长度；若结果随block长度改变，说明误差模型尚未收敛。

\section{Jackknife与block bootstrap}

Jackknife通过删除一个block重新计算观测量，得到方差估计：
\begin{equation}
  \sigma_{\rm JK}^2
  =\frac{B-1}{B}\sum_{b=1}^{B}
  (\hat O_{-b}-\bar O_{\rm JK})^2,
  \label{eq:jackknife}
\end{equation}
其中 \(B\) 是block数。非线性的外推、比值和矩阵求逆适合用delete-one-block重算整条链，而不是先算每个block的中心值再套delta方法。Block bootstrap则从block中重采样并保留局部相关性；block长度应大于主要自相关时间。

\section{Bootstrap全链传播}

设第 \(r\) 个bootstrap replica 提供一组独立数据 \(D_r\)，完整分析链写成
\begin{equation}
  \mathcal A:D_r\longmapsto
  \{\widehat Z_r,\widehat C_r,\widehat{\mathbf M}_r,
  \widehat f_r,\widehat{\sigma}_r\},
  \label{eq:bootstrap-chain}
\end{equation}
其中每一步都使用同一replica的关联函数、协方差、重正化和匹配。最终协方差为
\begin{equation}
  \widehat{\rm Cov}(f)
  =\frac{1}{R-1}\sum_{r=1}^{R}
  (f_r-\bar f)(f_r-\bar f)^T.
  \label{eq:bootstrap-cov}
\end{equation}
若不同阶段分别重采样，或在重正化时混用其他replica的协方差，就会丢失相关误差。Bootstrap不是“多跑几次拟合”，而是一个可审计的随机复制实验。

\begin{algorithm}[htbp]
\caption{保留组态相关性的全链bootstrap。}
\label{alg:full-chain-bootstrap}
\begin{algorithmic}[1]
\Require 组态索引集合 \(S\)，block长度 \(L_b\)，完整分析函数 \(\mathcal A\)
\Ensure 中心值、协方差和失效replica数
\State 对原始数据运行 \(\mathcal A\)，得到中心结果 \(f_0\)
\For{\(r=1,\ldots,R\)}
  \State 从block索引中有放回抽取 \(S_r\)
  \State 用 \(S_r\) 重算所有关联函数和协方差
  \State 执行重正化、软减除、匹配和外推，得到 \(f_r\)
  \State 记录拟合状态、残差、矩阵条件数和失败原因
\EndFor
\State 用成功replica计算 \(\bar f\) 与 \(\widehat{\rm Cov}(f)\)
\State 报告失败率和失败是否依赖观测值
\end{algorithmic}
\end{algorithm}

\section{相关拟合与广义最小二乘}

若数据矢量 \(y\) 的协方差为 \(C\)，模型为 \(y_{\rm model}(\theta)\)，则广义最小二乘估计满足
\begin{equation}
  \chi^2(\theta)
  =(y-y_{\rm model})^TC^{-1}(y-y_{\rm model}),
  \qquad
  \hat\theta=\arg\min_\theta\chi^2(\theta).
  \label{eq:stats-gls}
\end{equation}
参数协方差近似为
\begin{equation}
  {\rm Cov}(\hat\theta)
  =\left(J^TC^{-1}J\right)^{-1},
  \qquad
  J_{i\alpha}=\frac{\partial y_{{\rm model},i}}{\partial\theta_\alpha}.
  \label{eq:gls-cov}
\end{equation}
估计 \(C\) 的样本数不足时，直接求逆会放大噪声。应使用shrinkage、因子模型、低秩加对角近似或bootstrap协方差，并检查 \(C\) 的条件数。若协方差随拟合参数变化，应在每个参数点重新计算，而不是固定中心协方差。

\section{模型选择与模型平均}

增加拟合参数通常会降低 \(\chi^2\)，但不一定提高预测能力。AIC与BIC分别定义为
\begin{align}
  {\rm AIC}&=\chi^2+2k,
  \label{eq:aic}\\
  {\rm BIC}&=\chi^2+k\ln N,
  \label{eq:bic}
\end{align}
其中 \(k\) 是参数数、\(N\) 是有效数据数。模型权重可写成
\begin{equation}
  w_m=\frac{e^{-\frac12{\rm IC}_m}}
  {\sum_{m'}e^{-\frac12{\rm IC}_{m'}}}.
  \label{eq:model-weight}
\end{equation}
信息准则不能替代物理可辨识性。若两个模型在拟合窗口内给出几乎相同的函数，但在目标区域外相差很大，应把模型差异作为系统误差，而不是只报告权重最高的模型。

\section{GEVP与激发态分离}

广义本征值问题（GEVP）通过矩阵相关函数分离能级：
\begin{equation}
  C(t)v_n(t,t_0)
  =\lambda_n(t,t_0)C(t_0)v_n(t,t_0).
  \label{eq:stats-gevp}
\end{equation}
在理想谱分解下，\(\lambda_n(t,t_0)\) 趋于 \(e^{-E_n(t-t_0)}\)。GEVP矩阵的维数由算符基底决定，基底不足时高态会污染低态，基底过多则协方差估计不稳。必须报告基底算符、\(t_0\)、本征值稳定性以及广义矩阵的条件数。

\begin{algorithm}[htbp]
\caption{GeVp谱学与稳定性扫描。}
\label{alg:gevp}
\begin{algorithmic}[1]
\Require 相关矩阵 \(C(t)\)，参考时间 \(t_0\)，算符基底标号
\Ensure 能级估计、本征向量和模型误差
\State 对每个bootstrap replica 正则化 \(C(t_0)\) 并求广义本征值
\State 对每个能级计算有效质量 \(E_n(t,t_0)\)
\State 扫描 \(t_0\)、基底大小和拟合窗口
\State 检查本征向量与相邻时间片的重叠稳定性
\State 丢弃条件数超过阈值的replica并记录失败率
\State 输出能量、拟合区间、基底和协方差
\end{algorithmic}
\end{algorithm}

\section{极限外推的实验设计}

连续极限、体积极限和无穷动量极限不是简单的“拟合到无穷”。每个极限至少需要两个可辨识参数点，并应记录：
\begin{itemize}
  \item 变量定义与目标物理标度；
  \item 有效展开参数的范围与高阶项估计；
  \item 数据窗口的上下边界；
  \item 交叉项 \(aP_z\)、\(aM_\pi\) 或 \(M_\pi L\) 是否被控制；
  \item 不同外推路径是否给出一致结果。
\end{itemize}
若只在一组参数上做多项式外推，系数不可辨识。此时可以报告“局部趋势”或“模型依赖估计”，但不能把它称为物理极限。

\section{坐标空间到动量空间的逆问题}

TMD的Fourier变换可写成积分算子
\begin{equation}
  \widetilde f(x)=\int\dd z\,K(x,z)h(z),
  \label{eq:inverse-integral}
\end{equation}
有限格点只给离散、带噪的 \(h(z_i)\)，而大 \(z\) 数据缺失。离散化后得到
\begin{equation}
  \widetilde f=K h+\epsilon,
  \qquad
  \epsilon\sim\mathcal N(0,C_\epsilon).
  \label{eq:inverse-discrete}
\end{equation}
即使 \(K\) 是方阵，其小奇异值也会放大噪声。条件数 \(\kappa(K)\) 很大时，最小二乘解可能剧烈振荡，端点与正定性也会破坏。

\begin{figure}[htbp]
\centering
\begin{tikzpicture}[
  stage/.style={draw=qcdnavy!70!black,rounded corners=2pt,fill=qcdmist,align=center,minimum width=3.1cm,minimum height=0.95cm},
  arrow/.style={-{Latex[length=2mm]},thick,qcdrust!75!black}
]
\node[stage] (data) {有限 \(z\) 数据\\和协方差};
\node[stage,right=0.8cm of data] (model) {尾部/先验\\模型};
\node[stage,right=0.8cm of model] (solve) {正则化反演\\或拟合};
\node[stage,below=0.8cm of solve] (valid) {正定性、端点\\与残差检查};
\node[stage,left=0.8cm of valid] (err) {系统误差\\与模型平均};
\draw[arrow] (data) -- (model);
\draw[arrow] (model) -- (solve);
\draw[arrow] (solve) -- (valid);
\draw[arrow] (valid) -- (err);
\end{tikzpicture}
\caption{从坐标空间数据到分布函数的逆问题链。正则化不是替代缺失数据，而是把模型假设显式化。}
\label{fig:inverse-pipeline}
\end{figure}

\section{\texorpdfstring{大 \(z\) 尾部与正则化}{大 z 尾部与正则化}}

常用尾部模型为
\begin{equation}
  h(z)\sim\sum_j c_jz^{-d_j}e^{-z/\lambda_j},
  \label{eq:tail-model}
\end{equation}
其中指数衰减来自能隙或有效质量，幂次来自端点与高扭度。尾部参数必须在 \(\lambda=zP_z\) 的大距离区域拟合，并通过不同 \(P_z\) 的联合约束检验。若尾部模型对目标 \(x\) 区间的贡献超过统计误差，应将其变化纳入系统误差。

导数方法把Fourier变换改写为
\begin{equation}
  \widetilde f_d(x)
  =\frac{i}{x}\int\dd z\,e^{ixP_zz}h'(z),
  \qquad x\neq0,
  \label{eq:derivative-method}
\end{equation}
以压低端点振铃；但 \(1/x\) 会放大小 \(x\) 误差，而且 \(x=0\) 需要单独处理。导数、尾部模型和正则化参数不能在同一组数据上任意调优；应区分训练窗口和验证窗口。

\section{Backus--Gilbert与Bayes视角}

Backus--Gilbert方法把目标量写成核函数 \(A(x,y)\) 的线性泛函，并通过最小化分辨率与噪声的权衡来构造：
\begin{equation}
  \hat f(x_0)=\int\dd x\,A(x_0,x)\widetilde f(x).
  \label{eq:backus-gilbert}
\end{equation}
Bayes方法则用后验
\begin{equation}
  p(f|h)\propto
  \exp\left[-\frac12(Kf-h)^TC_\epsilon^{-1}(Kf-h)\right]
  p(f),
  \label{eq:bayes-inverse}
\end{equation}
并把先验、正则化和模型选择显式记录。两种方法都产生带偏估计；报告必须同时给出分辨率尺度、先验敏感性和后验预测检查。

\section{系统误差、blinding与可复现性}

系统误差不是把所有变化相加，而是按物理来源分组：
\begin{equation}
  \sigma_{\rm sys}^2
  =\sum_s\sigma_s^2+2\sum_{s<s'}C_{ss'}.
  \label{eq:systematic-cov}
\end{equation}
若两个变化来自同一拟合窗口或同一重正化方案，其协方差可能非零。Blinding可以通过未知偏移、随机化分析选择或保留测试集实现，但必须在正式解盲前固定分析方案。所有随机种子、bootstrap索引、模型扫描空间和版本哈希应保存，以便重新生成结果图。

\begin{checkpointbox}{统计与逆问题验收}
\begin{enumerate}
  \item 是否定义独立样本单位和自相关时间；
  \item 是否在每个replica内传播完整相关链；
  \item 是否报告fit窗口、参数协方差、残差和模型权重；
  \item 是否控制外推变量与交叉项，而不是在单点上宣称极限；
  \item 是否报告反演核、条件数、正则化参数和分辨率；
  \item 是否把模型差异、blinding和随机种子纳入可复现记录。
\end{enumerate}
\end{checkpointbox}

\clearpage
\chapter{涂抹、传播子与源构造}
\label{sec:b03-smearing-source}

涂抹不是物理算符定义之外的“美化步骤”，而是对插值场、规范连接或两者施加一个
明确、可重复且参数化的线性或非线性变换。它主要改变源与目标强子态的 overlap、
短程噪声和有效算符的紫外结构；它不会因为关联函数误差变小就自动给出更接近物理
极限的结果。因此，任何涂抹都必须和源、汇、动量投影、规范链接、重整化以及误差
传播一起记录。本节把夸克场涂抹、动量涂抹、蒸馏和规范场平滑放入同一框架，并在
节末给出选择与失效条件。

\section{协变 Laplace 核与夸克场涂抹}

设 \(\Delta_3(t)\) 为固定时间片上的三维规范协变 Laplace 算符。对空间格点
\(\boldsymbol{x},\boldsymbol{y}\)，可用连接变量
\(\widetilde U_j(\boldsymbol{x},t)\) 写成
\begin{equation}
  -\Delta_3(\boldsymbol{x},\boldsymbol{y};t)
  =6\delta_{\boldsymbol{x}\boldsymbol{y}}
  -\sum_{j=1}^{3}\Bigl[
  \widetilde U_j(\boldsymbol{x},t)\delta_{\boldsymbol{x}+\hat j,\boldsymbol{y}}
  +\widetilde U_j^\dagger(\boldsymbol{x}-\hat j,t)
  \delta_{\boldsymbol{x}-\hat j,\boldsymbol{y}}\Bigr].
  \label{eq:b03-laplace}
\end{equation}
在规范变换 \(V(\boldsymbol{x},t)\) 下，连接变量作协变平移，故
\(-\Delta_3\) 的谱和由低模张成的投影均有明确规范含义。低本征模是空间缓变模，
高本征模是短波模；涂抹的作用正是重新加权这两类模。

Wuppertal/Jacobi 涂抹由反复作用
\begin{equation}
  J_{\sigma,n_\sigma}(t)
  =\left(1+\frac{\sigma\Delta_3(t)}{n_\sigma}\right)^{n_\sigma},
  \qquad \sigma>0,
  \label{eq:b03-jacobi}
\end{equation}
定义。若 \(-\Delta_3\) 的非负本征值为 \(\lambda_k\)，自由场中
\(\lambda(\boldsymbol p)=(4/a^2)\sum_j\sin^2(ap_j/2)\)，则该核的本征值为
\begin{equation}
  J_k=\left(1-\frac{\sigma\lambda_k}{n_\sigma}\right)^{n_\sigma}.
  \label{eq:b03-jacobi-spectrum}
\end{equation}
固定 \(n_\sigma\) 且 \(\sigma\to0\) 时，一阶展开给出扩散型权重；固定物理
\(\sigma\) 而 \(n_\sigma\to\infty\) 时
\(J_{\sigma,n_\sigma}\to e^{\sigma\Delta_3}\)。在自由场中其坐标空间宽度约为
\(\sqrt{2\sigma}\)。这个极限说明 Jacobi 与 Gaussian 涂抹的关系，也说明
“涂抹半径”不是严格的紧支撑截断：大本征模只是被压得更强，而不是严格置零。
数值实现必须保证 \(1-\sigma\lambda_{\max}/n_\sigma\ge0\) 或明确记录高模符号
翻转的风险；否则核的谱权重可能失去正性。

对源和汇分别使用 \(J_s\)、\(J_{\mathrm{snk}}\) 后，两点函数的有效源结构变成
\(J_s^\dagger J_s\) 与 \(J_{\mathrm{snk}}^\dagger J_{\mathrm{snk}}\) 的组合。
只涂抹源而不涂抹汇，或在不同分析阶段混用不同参数，都会改变 overlap 和接触项，
不能把由此产生的曲线差异解释为统计涨落。

\section{动量涂抹：相位中心与平移不变性}

对目标纵向动量 \(P_z\)，普通零动量涂抹核在高 boost 下与强子波函数的重叠迅速
下降。动量涂抹在源上加入空间相位，使涂抹核的动量中心从零移到参考动量
\(\boldsymbol{\zeta}\)。一种定义是
\begin{equation}
  q_{\boldsymbol{\zeta}}(\boldsymbol{x},t)
  =e^{i\boldsymbol{\zeta}\cdot\boldsymbol{x}}\phi(\boldsymbol{x},t),
  \qquad
  \Phi_{\boldsymbol p}
  =\sum_{\boldsymbol{x}} e^{-i(\boldsymbol p-\boldsymbol{\zeta})\cdot\boldsymbol{x}}
  q_{\boldsymbol{\zeta}}(\boldsymbol{x},t),
  \label{eq:b03-momentum-smear}
\end{equation}
其中第二个等号把相位吸收到投影动量中。直接以相位源执行的涂抹核不是严格的
平移不变核；在插值算符的源端取 \(e^{i\boldsymbol{\zeta}\cdot\boldsymbol{x}}\)、
汇端取相共轭 \(e^{-i\boldsymbol{\zeta}\cdot\boldsymbol{x}}\) 后，整体关联函数在
平移下保持不变。U045 所述四种构造中，只有这种单相位配对在源汇相消的意义上保持
平移不变；这是一项应在具体 perambulator 定义中复核的结构性质，而不是说其他
相位组合一定不能用于任何计算。

动量涂抹不改变 TMD 算符的 Wilson 线方向、横向分离或软函数，也不替代大动量外推。
它的作用是提高目标动量附近的 overlap。若 \(e^{i\zeta x}\) 与格点动量量子化、
Fourier 相位或源汇方向不一致，关联函数的实部/虚部关系会被破坏。典型的失败模式
包括：相位符号反置、只在一端加相位、把物理外动量与涂抹参考动量重复相加，以及
在复共轭传播子上漏掉相位共轭。

\section{蒸馏：低模投影与 perambulator}

蒸馏把 Jacobi 核的低模特征显式化。解每个时间片的本征值问题
\begin{equation}
  -\Delta_3(t)v^{(k)}(t)=\lambda_k v^{(k)}(t),
  \qquad k=1,\ldots,N_v,
  \label{eq:b03-distill-eigen}
\end{equation}
并定义秩为 \(N_v\) 的投影
\begin{equation}
  \Box(t)=\sum_{k=1}^{N_v}v^{(k)}(t)v^{(k)\dagger}(t),
  \qquad
  \Box^2=\Box,\qquad \Box\xrightarrow{N_v\to N_s}\mathbb 1 ,
  \label{eq:b03-distill-projector}
\end{equation}
其中 \(N_s\) 是单时间片的空间-颜色维数。它在给定截断下保留长波模，但会丢弃高模。
传播子在低模空间中的表示为
\begin{equation}
  \tau_{\alpha\beta}^{(p_f,p_i)}(t_f,t_i)
  =V^{(p_f)\dagger}(t_f)D_{\alpha\beta}^{-1}(t_f,t_i)V^{(p_i)}(t_i).
  \label{eq:b03-perambulator}
\end{equation}
动量涂抹可把相位附加到本征矢，而不必为每个目标动量重新求逆；这是
distillation 与顺序源在成本结构上的关键区别。

蒸馏的截断误差不能仅由“曲线变平滑”判断。可靠检验至少包括：增加 \(N_v\) 后
矩阵元是否稳定；与直接传播子或较大基计算比较；目标算符的量子数与低模基是否匹配；
perambulator 的缓存是否绑定同一组态、同一 Laplace 算符和同一规范链接。若基中
缺少与目标态耦合的重要模式，增大统计量不会修复由此产生的系统偏差。

\section{规范场平滑：APE、HYP、stout 与 Wilson flow}

规范场平滑把 link 的短程涨落替换为局域 staple 平均。APE 型更新为
\begin{equation}
  U'_\mu(x)=\mathcal P_{SU(3)}\left[
  (1-\alpha)U_\mu(x)+\frac{\alpha}{6}\sum_{\nu\ne\mu}C_{\mu\nu}(x)\right],
  \label{eq:b03-ape}
\end{equation}
其中
\(C_{\mu\nu}(x)=U_\nu(x)U_\mu(x+\hat\nu)U_\nu^\dagger(x+\hat\mu)\)，
\(\mathcal P_{SU(3)}\) 为群投影。APE 的投影使更新对 link 不再处处解析，因而不能
直接作为需要保持可微性的 HMC 变换。

HYP 平滑在超立方体内部递归构造更局域的 staple。可抽象写成
\begin{equation}
  V^{(r+1)}_\mu(x)=\mathcal P_{SU(3)}\left[
  (1-\alpha_r)V^{(r)}_\mu(x)+\alpha_r\mathcal S^{(r)}_\mu(x)\right],
  \qquad r=1,2,3,
  \label{eq:b03-hyp}
\end{equation}
其中 \(\mathcal S^{(r)}_\mu\) 是第 \(r\) 层超立方体 staple 的归一化平均，
\(V^{(0)}=U\)。U045 给出的标准权重 \((0.75,0.6,0.3)\) 是历史实现参数，不是
物理常数。HYP 的平滑范围较局域，常用于胶子场强和 Wilson 线算符短程噪声控制；
但层数、权重和归一化必须由具体实现定义，不能把不同权重的裸矩阵元直接比较。

stout 平滑用解析指数保持 \(SU(3)\)，避免显式群投影的非解析性：
\begin{equation}
  U'_\mu(x)=e^{iQ_\mu(x)}U_\mu(x),
  \qquad
  Q_\mu=\frac{i}{2}\left(\Omega_\mu^\dagger-\Omega_\mu\right)
  -\frac{i}{2N_c}\operatorname{Tr}
  \left(\Omega_\mu^\dagger-\Omega_\mu\right),
  \label{eq:b03-stout}
\end{equation}
其中 \(\Omega_\mu\) 由与 \(U_\mu\) 共面的 staple 加权构成。\(Q_\mu\) 是
反厄米无迹矩阵，故 \(e^{iQ_\mu}\in SU(3)\) 且对 link 解析可微。stout 可直接
用于 HMC 或作为 Laplace 本征矢前的规范场预处理。

Wilson flow 把离散平滑推广为连续演化：
\begin{equation}
  \frac{dB_\mu(x,t_f)}{dt_f}
  =-g_0^2\frac{\delta S_{\rm YM}[B]}{\delta B_\mu(x,t_f)},
  \qquad B_\mu(x,0)=U_\mu(x).
  \label{eq:b03-flow}
\end{equation}
在 plaquette 作用量下以 Runge--Kutta 步进；平滑半径通常取
\(r_f=\sqrt{8t_f}\)。小步 stout 与 flow 在连续极限下有对应关系，但 flow 的关键
性质是：在标准 QCD 参数重整化后，\(t_f>0\) 的流场具有受控的紫外有限性。这个
性质首先属于流场及其定义良好的复合量，不应无条件推广到任意裸非局域算符。

\begin{longtable}{@{}p{0.14\textwidth}p{0.25\textwidth}p{0.25\textwidth}p{0.26\textwidth}@{}}
\caption{涂抹与平滑方案的可审计比较}\label{tab:b03-smearing-comparison}\\
\toprule
方案 & 作用对象与核心变换 & 保持的性质 & 主要窗口与失效条件 \\
\midrule
\endfirsthead
\toprule
方案 & 作用对象与核心变换 & 保持的性质 & 主要窗口与失效条件 \\
\midrule
\endhead
Jacobi--Gaussian & 夸克场；Laplace 多项式核 & 规范协变；参数小时连续可微 & 过强时压掉物理短程；高模符号/归一化错误 \\
动量涂抹 & 夸克源；空间相位与投影 & 正确源汇配对下平移不变 & 相位符号或端点配对错误；不能替代大动量外推 \\
蒸馏 & 低 Laplace 模投影 & 规范协变；秩截断明确 & \(N_v\) 未收敛；基态模式遗漏；缓存键不一致 \\
APE & 规范 link；staple 平均后投影 & 保持 \(SU(3)\) 与规范协变 & 显式投影非解析；平滑过强；不能嵌入可微 HMC 更新 \\
HYP & 超立方体分层 staple & 保持 \(SU(3)\) 与局域性 & 权重/归一化依赖实现；平滑改变算符 UV 结构 \\
stout & 解析指数链接更新 & 保持 \(SU(3)\)、规范协变、可微性 & 步数和权重未收敛；与匹配方案混用 \\
Wilson flow & 连续规范场演化 & 规范协变；\(t_f>0\) 的受控 UV 性质 & 流窗口为空；\(r_f\) 与目标距离/格距不分离 \\
\bottomrule
\end{longtable}

涂抹参数属于分析定义。选择时应同时检查基态 overlap、目标算符的几何和矩阵元的
连续极限，而不应只最小化某一时间片的统计误差。若不同涂抹参数给出的中心值差异
大于统计误差，必须把它并入系统误差或明确判定窗口失败。

\chapter{传播子求解：从 Dirac 逆到可验证矩阵元}
\label{sec:b03-propagator-solver}

\section{Dirac 逆、传播子与 hopping 图像}

在固定规范背景 \(U\) 中，夸克传播子满足
\begin{equation}
  \sum_{y,\beta,b}D_{\alpha a;\beta b}(x,y;U)
  S_{\beta b;\gamma c}(y,z;U)
  =\delta_{\alpha\gamma}\delta_{ac}\delta_{xz}.
  \label{eq:b03-dirac-inverse}
\end{equation}
Dirac 算符 \(D\) 是稀疏矩阵，但 \(D^{-1}\) 通常稠密；因此实际计算是求一列或
少数几列 \(D\chi=\eta\)，而不是显式形成逆矩阵。这里的 \(\eta\) 是源向量，
携带位置、颜色、Dirac 结构、动量投影和涂抹信息。单个组态上的 \(\chi\) 只是该
规范背景下的条件样本，只有在规范组态平均并按 Wick 收缩组合后才有强子矩阵元的
意义。

对 Wilson 型算符，可把质量与对角项吸收到归一化后写成
\begin{equation}
  D=C(1-\kappa H),
  \qquad
  S=D^{-1}=C^{-1}\sum_{j=0}^{\infty}\kappa^jH^j .
  \label{eq:b03-hopping}
\end{equation}
\(H^j(x,y)\) 是长度为 \(j\) 的格点路径和，路径权重包含规范 link 乘积及 Dirac
因子。因子 \((1-\gamma_\mu)(1+\gamma_\mu)=0\) 排除了立即折返的路径，因此
hopping 展开给出“非回溯路径求和”的物理图像。大夸克质量使 \(\kappa\) 小、短路径
主导；趋近临界质量时 \(\kappa\) 接近 \(\kappa_c\)，级数收敛极慢。该展开用于理解
传播子的路径结构，不能替代轻夸克情形下的数值迭代。

实质量 Wilson 型算符满足
\begin{equation}
  \gamma_5D\gamma_5=D^\dagger,
  \qquad
  S(y,x)=\gamma_5S(x,y)^\dagger\gamma_5 .
  \label{eq:b03-gamma5}
\end{equation}
这可用于反向传播子和求解复用，但要求 Dirac 表示、颜色索引、源汇时间顺序和
复共轭约定完全一致。若实现使用不同的 \(\gamma_5\)-Hermiticity 定义，先在
小矩阵上验证式~\eqref{eq:b03-gamma5}，再进入生产分析。

\section{源的选择与成本}

源决定传播子投影到哪些动量、位置、自旋和颜色结构。点源只在一个格点注入一个
Dirac-颜色分量，若需形成该点的完整 \(4\times3\) 传播子，需要 12 次独立求逆；
它提供全动量信息，但基态 overlap 可能较差。Gaussian/Jacobi 源通过空间 Laplace
核扩展源，动量源在点源上乘 \(e^{i\boldsymbol p\cdot\boldsymbol x}\)。随机源以
满足 \(\langle\eta_i^*\eta_j\rangle=\delta_{ij}\) 的 \(Z_2/Z_4\) 向量估计迹：
\begin{equation}
  \operatorname{Tr}\!\left[\Gamma D^{-1}(x,x)\right]
  \simeq\frac{1}{N_r}\sum_{r=1}^{N_r}
  \eta^{(r)\dagger}(x)\Gamma[D^{-1}\eta^{(r)}](x).
  \label{eq:b03-stochastic-trace}
\end{equation}
随机误差按 \(1/\sqrt{N_r}\) 降低，低模噪声、源间相关和真空减除必须与统计误差
一起报告。

顺序源用于三点函数。若 \(S(z;0)\) 是已算好的源传播子，可将汇点结构、插入时间和
投影组合成辅助源
\begin{equation}
  \eta_{\rm seq}(z)=e^{-i\boldsymbol p_f\cdot\boldsymbol z}
  P_{\rm sink}(z)\gamma_5S(z;0)\,\delta_{z_0,t_s},
  \label{eq:b03-sequential-source}
\end{equation}
再求解一次 \(D\chi_{\rm seq}=\eta_{\rm seq}\)。通常每个新的末态动量、汇点投影
或插入结构都需要新的辅助源；改变初始动量则可能通过重新投影已有传播子复用。
因此成本不是“顺序源只多一次求逆”，而是与末态动量数和汇点算符数成线性关系的
额外求解。

\begin{table}[htbp]
\centering
\caption{常用夸克源的结构与主要风险}
\label{tab:b03-source-types}
\begin{tabularx}{\textwidth}{@{}>{\raggedright\arraybackslash}p{0.17\textwidth}X X@{}}
\toprule
源类型 & 主要用途与成本 & 失效模式 \\
\midrule
点源 & 12 次求逆得到单点完整传播子；全动量信息 & 基态 overlap 小；大体积统计噪声高 \\
Gaussian--Jacobi & 预处理空间低模，提高基态 overlap & 涂抹参数与 sink 不一致；过度平滑改变接触项 \\
动量源 & 注入确定空间动量，服务 RI/MOM 或投影 & 相位符号/归一化错误；高 boost overlap 仍可能差 \\
随机源 & 估计迹、all-to-all 或非连通贡献 & 低模未精确处理；源相关和真空减除误差 \\
顺序源 & 组合汇点与插入，减少三点收缩的独立求逆 & 每个末态动量/投影增加一次求解；中间态污染 \\
distillation 源 & 低秩 perambulator，支持多算符复用 & \(N_v\) 截断、基失配、缓存键错误 \\
\bottomrule
\end{tabularx}
\end{table}

\section{Krylov 求解与谱条件}

对 \(\gamma_5\)-Hermitian 的可逆 \(D\)，可求解 Hermitian 正定法方程
\begin{equation}
  A\chi=D^\dagger\eta,
  \qquad A=D^\dagger D.
  \label{eq:b03-normal-equation}
\end{equation}
若 \(\kappa(A)=\lambda_{\max}(A)/\lambda_{\min}(A)\)，CG 的理想误差界为
\begin{equation}
  \frac{\|e_k\|_A}{\|e_0\|_A}
  \leq 2\left(\frac{\sqrt{\kappa(A)}-1}{\sqrt{\kappa(A)}+1}\right)^k .
  \label{eq:b03-cg-bound}
\end{equation}
法方程会把原 Dirac 谱的条件数平方，因此轻夸克和近零模使收敛显著变慢。只有当
求解矩阵确实是 Hermitian 正定时才可直接使用 CG；\(\gamma_5\)-Hermitian 并不
等于 Hermitian，非 Hermitian 的 \(D\) 应使用适当的非对称 Krylov 方法或经过
明确谱分析的等价变换。

奇偶预条件把 Wilson 型 Dirac 矩阵分块为
\begin{equation}
  D=\begin{pmatrix}D_{ee}&D_{eo}\\D_{oe}&D_{oo}\end{pmatrix},
  \qquad
  S_{ee}=D_{ee}^{-1}+D_{ee}^{-1}D_{eo}
  \left(D_{oo}-D_{oe}D_{ee}^{-1}D_{eo}\right)^{-1}
  D_{oe}D_{ee}^{-1},
  \label{eq:b03-even-odd}
\end{equation}
实际常先消去一个宇称，再对约化算符迭代；对角块求逆只有在作用量结构和边界条件
允许时才应视为低成本操作。多移位 CG 利用同一 Krylov 空间同时求
\((D+\sigma_i)\chi_i=\eta\)，复用矩阵-向量乘法；对近似奇异移位仍需检查数值稳定
和停止条件。deflation 把低模基底 \(V_k\) 从残差中投影掉，
\begin{equation}
  r_0^{\perp}=r_0-V_kV_k^\dagger r_0,
  \label{eq:b03-deflation}
\end{equation}
但基底必须来自同一组态、同一 Dirac 算符、同一质量和同一边界条件。借用不匹配的
低模会加速一个错误的系统。

\section{多重网格与进阶求解}

多重网格把低频误差限制到粗层，在粗层求修正，再延拓回细层：
\begin{equation}
  x\leftarrow x+P\,A_c^{-1}R(b-Ax),
  \label{eq:b03-multigrid}
\end{equation}
其中 \(R\) 是 restriction、\(P\) 是 prolongation、\(A_c\) 是粗层算符。粗化必须
保持规范协变和拓扑结构；简单平均链接可能产生非物理零模。Domain-wall/overlap
算符需要专门的五维或符号函数预条件，overlap 的
\(\operatorname{sign}(\gamma_5D_W)\) 通常通过有理或 Chebyshev 近似实现，每次
近似都嵌套内层求逆，不能用普通 Wilson-CG 的成本模型估计。

\section{残差认证与失败条件}

求解器打印的递归残差可能因有限精度和重启动偏离真实残差。每次生产求解至少应保存
\begin{equation}
  r_{\rm rec}=b-Ax_k,
  \qquad
  \epsilon_{\rm true}=\frac{\|b-Ax_k\|_2}{\|b\|_2},
  \label{eq:b03-true-residual}
\end{equation}
并周期性用原算符重算 \(r_{\rm true}\)。停止容差应通过目标矩阵元的容差扫描确定：
把 \(\epsilon_{\rm true}\) 缩小一个数量级，若中心值和误差带不变，才能认为求解
误差低于统计误差。若不同求解器、预条件或重启策略给出不一致结果，应保留失败状态，
不得把失败 replica 静默替换为中心值。

\begin{table}[htbp]
\centering
\caption{传播子求解的常见失效模式与验收证据}
\label{tab:b03-solver-failures}
\begin{tabularx}{\textwidth}{@{}>{\raggedright\arraybackslash}p{0.17\textwidth}X X@{}}
\toprule
现象 & 物理/数值来源 & 验收证据 \\
\midrule
残差很小但矩阵元漂移 & 递归残差假收敛、近零模放大 & 真实残差、不同初值/求解器交叉检查 \\
迭代数随夸克变轻爆炸 & 谱条件数、近零模、拓扑背景 & 残差历史、条件数、deflation/multigrid 对照 \\
顺序源结果依赖末态动量 & 每个动量新辅助源未复用/归一化错误 & 源汇投影、额外求逆计数、动量扫描 \\
随机迹误差不随 \(N_r\) 下降 & 源间相关、低模未精确处理、真空减除偏置 & 不同噪声族、低模分解、块 bootstrap \\
涂抹后平台提前出现但极限不稳定 & 过强平滑压掉激发态和物理短程结构 & 多涂抹参数、sink 一致性、连续极限 \\
\bottomrule
\end{tabularx}
\end{table}

\chapter{Monte Carlo 数据的统计误差与重采样}
\label{sec:b03-monte-carlo-errors}

\section{从组态样本到估计量}

格点 Monte Carlo 给出的是按路径积分权重抽取的有限规范组态。一个物理结果是组态
样本上的估计量 \(\hat O\)，本身具有统计分布。分析开始前必须固定可重采样的底层
单位：组态、轨迹、block、源、replica，还是多个系综的联合样本。若同一组态内的
源、时间片或噪声源高度相关，先把它们独立平均会丢失交叉协方差；若不同的组态已经
按轨迹分组，重采样应在组态/block 层进行，而不是在单个时间片层进行。

若 \(X_1,\ldots,X_N\) 是独立同分布样本，则
\begin{equation}
  \bar X=\frac1N\sum_{i=1}^NX_i,
  \qquad
  \widehat\sigma_X^2=\frac{1}{N-1}\sum_{i=1}^N(X_i-\bar X)^2,
  \qquad
  \operatorname{Var}(\bar X)=\frac{\sigma_X^2}{N}.
  \label{eq:b03-iid-mean}
\end{equation}
这是 \(\sigma/\sqrt N\) 规则的来源，也是其适用边界。对 HMC 轨迹序列，相邻组态
通常不独立；将每个时间片或每个源都当作独立样本，会系统性地低估误差。

\section{自相关与有效样本数}

对平稳序列 \(X_t\)，定义归一化自相关函数
\begin{equation}
  \rho_X(t)=\frac{C_X(t)}{C_X(0)},
  \qquad
  C_X(t)=\left\langle (X_s-\langle X\rangle)(X_{s+t}-\langle X\rangle)\right\rangle .
  \label{eq:b03-autocorr}
\end{equation}
当自相关长度远小于样本长度时，
\begin{equation}
  \tau_{\rm int}=\frac12+\sum_{t=1}^{\infty}\rho_X(t),
  \qquad
  \operatorname{Var}(\bar X)\simeq\frac{\sigma_X^2}{N}\,2\tau_{\rm int},
  \qquad
  N_{\rm eff}=\frac{N}{2\tau_{\rm int}}.
  \label{eq:b03-tau-int}
\end{equation}
式~\eqref{eq:b03-tau-int} 的 \(N_{\rm eff}\) 定义必须与 \(\tau_{\rm int}\) 的
\(1/2\) 约定配套；若换用另一种积分时间约定，公式中的因子 2 会移位。有限长度
序列中，不能无界地向大 \(t\) 求和：应比较多个积分窗口、自相关函数的首次负值、
block 长度和 jackknife 方差，报告平台及窗口系统误差。

binning 是更直接的自相关诊断。把连续 \(K\) 个原始测量平均成
\(X_m^{(K)}=K^{-1}\sum_{i=(m-1)K+1}^{mK}X_i\)。若原始测量独立，
\begin{equation}
  \operatorname{Var}\!\left(X_m^{(K)}\right)=\frac{\sigma_X^2}{K}.
  \label{eq:b03-binning}
\end{equation}
因此可画 \(K\widehat\sigma^2_{(K)}\) 对 \(K\)，寻找平台。平台未出现时，block
仍未覆盖主要自相关；平台出现后继续增大 \(K\) 会使 block 数减少、方差估计本身
变差。常用的“先 binning、再 bootstrap/jackknife”只是把相关数据近似成近似
独立 block，不能自动修复长程相关或拓扑冻结。

\section{Bootstrap：全链重采样}

设有 \(B\) 个近似独立的 block。对第 \(r\) 个 bootstrap replica，从 block 索引中
有放回抽取 \(B\) 个，形成 \(D_r\)，并在该 replica 内重跑从关联函数、真空减除、
重正化、Fourier/匹配到拟合和外推的**同一分析链**，得到 \(\theta_r\)。估计方差用
\begin{equation}
  \bar\theta_{\rm boot}=\frac1R\sum_{r=1}^R\theta_r,
  \qquad
  \widehat{\operatorname{Cov}}(\theta)
  =\frac{1}{R-1}\sum_{r=1}^R
  (\theta_r-\bar\theta_{\rm boot})(\theta_r-\bar\theta_{\rm boot})^T .
  \label{eq:b03-bootstrap-cov}
\end{equation}
bootstrap 的中心值原则上仍用全样本估计 \(\hat\theta\)，而 \(\theta_r\) 的散布
用于误差和区间。若把 \(\bar\theta_{\rm boot}\) 直接替代中心值，应明确这是偏差
修正；更稳健的做法是报告 \(\bar\theta_{\rm boot}-\hat\theta\) 的偏差、百分位区间和
BCa 等区间作为敏感性检查。

重采样数量 \(R\) 不能只看“通常取多少”。应检查增加 \(R\)、改变随机种子和改变
block 长度后中心值、标准误和区间是否稳定。高度非线性或近阈值观测量可能产生重尾
分布，此时正态近似和偏差修正不一定可靠，需报告分位数、失败率和选择效应。
每个 replica 的成败状态必须保存；把失败 replica 静默替换成中心值会人为缩小误差。

\section{Jackknife 与 block delete-one}

对 \(B\) 个 block，逐一删除第 \(b\) 个 block 并重跑完整分析，得到
\(\theta_{-b}\)。标准 delete-one jackknife 方差为
\begin{equation}
  \bar\theta_{\rm JK}=\frac1B\sum_{b=1}^B\theta_{-b},
  \qquad
  \widehat\sigma_{\rm JK}^2
  =\frac{B-1}{B}\sum_{b=1}^B
  (\theta_{-b}-\bar\theta_{\rm JK})^2 .
  \label{eq:b03-jackknife}
\end{equation}
若只用全样本值 \(\hat\theta\) 替代 \(\bar\theta_{\rm JK}\)，得到的只是近似版本，
并不严格等于 jackknife 方差。Jackknife 偏差修正可写为
\begin{equation}
  \hat\theta_{\rm JK,bc}=B\hat\theta-(B-1)\bar\theta_{\rm JK}.
  \label{eq:b03-jackknife-bias}
\end{equation}
Jackknife 对平滑的均值、比值和拟合参数通常有效；对中位数、阈值选择、模型选择或
不连续统计量可能失效。是否优先使用 jackknife 或 bootstrap 取决于统计量的光滑性、
block 数和失败模式，而不存在仅由样本量决定的固定偏好。

\section{相关拟合与协方差归一化}

拟合向量 \(d\) 与模型 \(f(\theta)\) 时，广义最小二乘为
\begin{equation}
  \chi^2(\theta)=(d-f(\theta))^TC_d^{-1}(d-f(\theta)),
  \qquad
  \hat\theta=\arg\min_\theta\chi^2(\theta).
  \label{eq:b03-gls}
\end{equation}
\(C_d\) 必须是**被拟合估计量**的协方差。如果数据向量是 \(N\) 个组态的均值，
而 \(C_{\rm conf}\) 是单组态协方差，则
\begin{equation}
  C_d=\frac{C_{\rm conf}}{N}
  \label{eq:b03-covariance-of-mean}
\end{equation}
在组态独立且协方差估计无偏时成立；更稳妥的做法是从 jackknife/bootstrap 的均值
分布直接估计 \(C_d\)，这样可同时保留源、时间片、重正化和拟合链的相关性。若把
单组态协方差当作均值协方差直接求逆，拟合误差的尺度会错误。

对线性模型 \(f=X\theta\)，协方差近似为
\begin{equation}
  \operatorname{Cov}(\hat\theta)
  =\left(X^TC_d^{-1}X\right)^{-1},
  \qquad
  J_{i\alpha}=\frac{\partial f_i}{\partial\theta_\alpha},
  \label{eq:b03-gls-cov}
\end{equation}
非线性模型的雅可比 \(J\) 应在 \(\hat\theta\) 处计算；若重正化、Fourier 或匹配
也依赖参数，导数必须包含这些步骤。协方差秩不足或条件数过大时，SVD、收缩、低秩
加对角或伪逆可以作为诊断，但每种正则化都要扫描截断阈值并传播到最终结果。

\section{全链传播与系统误差预算}

统计相关可以由 bootstrap 全链传播；模型选择、窗口选择、匹配阶数和重整化方案等
“硬选择”不服从组态重采样，必须单独扫描。一个可审计的结果至少区分三类误差：
\begin{enumerate}
  \item 由组态抽样和源/噪声抽样产生的统计误差；
  \item 由同一组态内相关的拟合、重正化、匹配和外推传播的协方差；
  \item 由模型、窗口、截断和方案选择产生、不能由 bootstrap 自动覆盖的系统变化。
\end{enumerate}
若系统变化之间近似独立，可报告平方和根
\(\sigma_{\rm quad}^2=\sigma_{\rm stat}^2+\sum_s\sigma_s^2\)；若变化之间相关，
应使用
\begin{equation}
  \sigma_{\rm sys}^2=\sum_s\sigma_s^2+2\sum_{s<s'}C_{ss'},
  \label{eq:b03-system-cov}
\end{equation}
或直接在每个 bootstrap replica 内联合扫描。把同一个窗口变化重复计入多个误差项
会低估相关性；反过来，把明显相关的量强行平方相加会给出虚假的小误差。

\begin{table}[htbp]
\centering
\caption{误差传播路径及适用边界}
\label{tab:b03-error-paths}
\begin{tabularx}{\textwidth}{@{}>{\raggedright\arraybackslash}p{0.18\textwidth}X X@{}}
\toprule
路径 & 能传播的相关性 & 不能自动处理的问题 \\
\midrule
全链 bootstrap & 组态、源、拟合、重正化、匹配和联合外推的统计相关 & 模型函数、窗口、算符几何等硬选择 \\
delete-one block jackknife & 近似独立 block、比值和拟合参数 & 非光滑统计量、长程相关未覆盖的 block 长度 \\
独立变化平方和根 & 各系统扫描的中心值差异 & 扫描项之间的相关、重采样统计相关 \\
协方差联合拟合 & 已知并可靠估计的完整协方差 & 协方差秩不足、模型误差和非高斯尾部 \\
\bottomrule
\end{tabularx}
\end{table}

\chapter{极限外推、联合拟合与尺度窗口}
\label{sec:b03-extrapolation-windows}

\section{外推是带物理依据的渐近展开}

格点结果依赖格距 \(a\)、体积 \(L\)、介子质量 \(m_\pi\)、强子动量 \(P_z\)、
流时 \(t_f\)、Wilson 线长度和涂抹参数。外推的目标不是把最后一个数据点接一条
最长曲线，而是把有限参数结果嵌入一个可检验的渐近展开：
\begin{equation}
  F(x;a,L,m_\pi,P_z,t_f,\ell)
  =F_0(x)+\sum_j c_j(x)\,\Phi_j(a,L,m_\pi,P_z,t_f,\ell)
  +\mathcal R(x;\text{window}).
  \label{eq:b03-general-expansion}
\end{equation}
每个基函数 \(\Phi_j\) 应来自离散化、有限体积、手征对称性、因子化或重整化群，
而不是只在拟合窗口内“形状好看”。\(\mathcal R\) 是未纳入模型的高阶项，必须在
窗口边界和目标精度上被论证为小量；否则结果只能称为参数点的有效值。

\section{连续极限与联合拟合}

对 \(O(a^2)\) 改进的作用量，标量观测量常写成
\begin{equation}
  O(a)=O_0+c_1a^2+c_2a^4+\cdots .
  \label{eq:b03-continuum-expansion}
\end{equation}
只有两个格距时，保留 \(a^2\) 项是模型假设，无法从数据中检验 \(a^4\)；三个格距
才开始具有识别高阶项的自由度。Clover 改进也不自动保证所有算符都是纯
\(O(a^2)\)：改进系数、边界项、接触项和算符混合都可能留下线性或交叉误差。

对 LaMET/TMD 矩阵元，离散化误差与 boost 动量耦合。一个常用的联合参数化为
\begin{align}
  F(x;a,m_\pi,P_z)
  &=F_0(x)+c_a(x)a^2+c_{aP}(x)a^2P_z^2
  +c_P(x)\frac{1}{P_z^2}\notag\\
  &\quad+c_m(x)m_\pi^2+c_L(x)e^{-m_\pi L}+\cdots .
  \label{eq:b03-joint-expansion}
\end{align}
其中 \(c_{aP}\) 项描述 Wilson 线或准算符离散化与高动量的交叉效应。参数化的
具体函数形式取决于作用量、算符和投影；式~\eqref{eq:b03-joint-expansion} 是
组织窗口的工具，不是所有胶子 TMD 的普适定理。若只有单一 \(a\)、单一 \(P_z\)
或单一 \(m_\pi\)，对应系数和交叉项不可辨识，不能通过增加拟合轮数制造约束。

连续极限、手征极限和无穷动量极限在理论上不交换。格点上必须先有有限 UV 截断，
再讨论大动量；LaMET 匹配则要求
\begin{equation}
  \Lambda_{\rm QCD},M_N\ll P_z\ll a^{-1},
  \qquad
  \Lambda_{\rm QCD}\ll \frac1{b_T},
  \label{eq:b03-lamet-window}
\end{equation}
同时成立。若 \(aP_z\) 接近 1，格点色散关系和高扭度/离散 Wilson 线误差都进入
非微扰区；若 \(b_T\) 接近 \(L/2\)，周期镜像会污染横向结构。此时“多拟合一个
\(1/P_z^2\) 项”不能修复窗口失败。

\section{大 \(\lambda\) 尾部：模型，而不是定律}

准分布通常由坐标空间矩阵元的 Fourier 变换给出。采用本节统一约定
\begin{equation}
  \widetilde q(x,P_z)
  =\frac{P_z}{2\pi}\int_{-\infty}^{\infty}\mathrm{d}z\,
  e^{ixP_z z}h_R(z,P_z),
  \qquad
  \lambda=zP_z .
  \label{eq:b03-quasi-fourier}
\end{equation}
格点只覆盖有限 \(|z|\)，且大 \(z\) 的信噪比通常快速下降。截断 Fourier 变换会
把未测尾部转换成振铃和端点畸变。一个常用的尾部参数化是
\begin{equation}
  h_R(z,P_z)\simeq\sum_{j=1}^{N_t}c_jz^{-d_j}e^{-z/\lambda_j},
  \qquad \lambda>\lambda_L,
  \label{eq:b03-tail-model}
\end{equation}
但指数因子、幂次和项数都是待检验的模型选择。拟合时必须记录：使用的
\(\lambda_L\)、项数、参数相关矩阵、有效点数、不同 \(P_z\) 的交叉约束，以及
从尾部到目标 \(x\) 区间的积分贡献。若多个模型在数据窗口内拟合质量相近、但目标
区间外差异很大，模型差异必须进入系统误差，不能只选 \(\chi^2\) 最小者。

derivative 方法可减少截断振铃。对式~\eqref{eq:b03-quasi-fourier} 分部积分，
假定边界项已减除或在无穷远消失，则
\begin{equation}
  \widetilde q(x,P_z)
  =\frac{P_z}{2\pi}\int\mathrm{d}z\,e^{ixP_z z}h_R(z,P_z)
  =\frac{i}{2\pi x}\int\mathrm{d}z\,e^{ixP_z z}
  \frac{\partial h_R(z,P_z)}{\partial z},
  \qquad x\ne0.
  \label{eq:b03-derivative-method}
\end{equation}
与只写 \(i/x\) 的形式相比，\(1/(2\pi)\) 来自本节固定的 Fourier 归一化；若采用
另一种准分布定义，前因子会同步改变。导数压低的是常数和慢变边界项，不改变
\(h_R\) 本身的统计噪声；\(1/x\) 会放大 \(x\to0\) 的误差，\(x=0\) 需单独处理。

\section{无穷动量外推}

若已经对每个 \(P_z\) 完成与动量相关的重正化和匹配，则匹配后的量可写成
\begin{equation}
  q(x,P_z)=q_0(x)+\frac{c_1(x)}{P_z^2}
  +\frac{c_2(x)}{P_z^4}+\cdots .
  \label{eq:b03-pz-expansion}
\end{equation}
\(1/P_z^2\) 的起点由因子化中的幂修正和对称性决定；端点或特定算符可能还有额外
结构，不能仅凭“看起来像 \(1/P_z\)”改基函数。两个动量点只能拟合一个领先修正，
且结果高度依赖是否已进入渐近区；三个或更多点才可开始检查高阶项。单 \(P_z\)
结果最多是在假设修正足够小的情况下的有效值。

当 \(x\) 较小时，匹配核的硬标度可能由 \(2xP_z\) 控制。若该标度接近
\(\Lambda_{\rm QCD}\)，固定阶匹配和 RGR 都进入非微扰区；这时应缩小报告的
\(x\) 区间，而不是用大 \(z\) 外推填补。小 \(x\) 的大对数、尾部模型和统计噪声
往往同时增强，必须分别报告其贡献。

\section{拟合设计、可辨识性与模型误差}

把全部数据点排成向量 \(d\)，把外推基函数排成设计矩阵 \(X\)。即使写成
\begin{equation}
  d=X\theta+\varepsilon,
  \qquad
  \operatorname{Cov}(d)=C_d,
  \label{eq:b03-fit-design}
\end{equation}
也不表示所有列都可辨识。若两列在可用窗口内近似共线，参数协方差会爆炸；此时应
删除不可辨识项、合并为有效参数，或扩大 \(a,P_z,m_\pi,L\) 覆盖范围。报告至少
给出设计矩阵条件数、参数相关矩阵、残差模式、留一数据点或留一块检验。

模型平均可作为系统误差的展示，但不能替代物理判据。若模型 \(m\) 的权重为
\(w_m\)，则
\begin{equation}
  w_m=\frac{e^{-\frac12{\rm IC}_m}}
  {\sum_{m'}e^{-\frac12{\rm IC}_{m'}}},
  \label{eq:b03-model-weight}
\end{equation}
其中 IC 可为 AIC/BIC 或针对相关数据的推广。权重依赖候选模型集合；若候选集合
未包含真实渐近形式，权重再精确也不代表物理。建议同时报告“最优模型”“等高线
模型族”和“目标区间模型包络”，而不是只报单一中心值。

\begin{table}[htbp]
\centering
\caption{极限外推的最小验收矩阵}
\label{tab:b03-extrapolation-checks}
\begin{tabularx}{\textwidth}{@{}>{\raggedright\arraybackslash}p{0.18\textwidth}X X@{}}
\toprule
方向 & 必须扫描的量 & 失败判据 \\
\midrule
连续极限 & 多个 \(a\)、\(a^2\) 与 \(a^4\) 项、\(aP_z\)、作用量改进系数 & 单一格距、高阶项改变中心值、窗口外推主导 \\
无限体积 & 多个 \(m_\pi L\)、\(b_T<L/2\)、长 Wilson 线平台 & \(b_T\) 或 \(L_\parallel\) 与 \(L\) 不分离 \\
手征极限 & \(m_\pi\) 覆盖、ChPT 项数、价/海夸克质量 & 只有单点或手征模型不可辨识 \\
大动量 & 多个 \(P_z\)、\(aP_z\)、匹配标度与 RGR & \(P_z\) 接近格点截断；小 \(x\) 进入非微扰区 \\
大 \(\lambda\) 尾部 & \(\lambda_L\)、项数、尾部函数、\(z_{\max}\) & 模型包络超过目标精度；拟合数据不足 \\
误差传播 & bootstrap 全链、模型扫描、相关系统项 & 失败 replica 静默丢弃；相关误差平方相加 \\
\bottomrule
\end{tabularx}
\end{table}

最终报告应明确哪些区间由数据约束、哪些由模型外推、哪些没有通过物理窗口。一个
通过所有统计拟合但位于 \(P_z\gg a^{-1}\)、\(r_f\gtrsim b_T\) 或
\(m_\pi L\ll1\) 的结果，不是更强的物理证据，而是更清楚诊断出的窗口失败。


\clearpage
\part{实现、验证与研究前沿}
\chapter{\texorpdfstring{高级专题：端点、自旋、小 \(x\) 与过程依赖}{高级专题：端点、自旋、小 x 与过程依赖}}
\label{sec:advanced-topics}

\section{端点区域的因子化}

TMD分布在 \(x\to1\) 时进入端点区。此时部分子携带几乎全部纵向动量，剩余胶子辐射受到相空间压制，固定阶展开中出现大对数。令 \(\xi=1-x\)，端点对数可写成
\begin{equation}
  \alpha_s^n\frac{\ln^m\xi}{\xi},
  \qquad m\le2n-1,
  \label{eq:endpoint-logs}
\end{equation}
其中最奇异项来自软和共线辐射。端点区不能简单套用中等 \(x\) 的匹配核，因为阈值修正、高扭度和软函数可能同时增强。格点准TMD在 \(x\to1\) 还受到有限 \(P_z\) 与 Fourier截断影响；即使矩阵元在 \(x=1\) 附近有限，也不能自动得到正确的端点极限。

一个实用的处理顺序是：
\begin{enumerate}
  \item 先在中等 \(x\) 区域验证匹配核和RGE；
  \item 单独检查端点附近的尾部、Fourier振铃和正定性；
  \item 仅在因子化和微扰阶数有控制时发布端点结果；
  \item 否则给出端点区间和模型误差，而不是单点数值。
\end{enumerate}

\section{\texorpdfstring{小 \(x\) 与胶子饱和}{小 x 与胶子饱和}}

小 \(x\) 区域胶子密度增大，固定阶DGLAP演化中的对数需要重求和，非线性饱和效应可能重要。BFKL、BK和CGC等框架描述不同层次的演化。TMD胶子分布在小 \(x\) 处与横向动量、色偶极和Wilson线几何密切相关。格点大动量方法受
\begin{equation}
  x_{\min}\sim\frac{\Lambda_{\rm QCD}}{2P_z}
  \label{eq:advanced-xmin}
\end{equation}
限制；要进入更小 \(x\)，必须增大 \(P_z\)，同时控制 \(aP_z\)、信噪比和匹配标度。

小 \(x\) 不是“外推到最后一点”即可。需要检查：
\begin{itemize}
  \item 小 \(x\) 区域的匹配核是否仍在微扰区；
  \item 横向动量是否进入非微扰饱和区；
  \item 大 \(z\) 尾部与 \(xP_z\) 的关系；
  \item 不同 \(P_z\) 的 \(x\) 重叠区是否一致；
  \item 单态胶子-夸克混合是否被保留。
\end{itemize}

\section{TMD波函数与光锥波函数}

介子TMD波函数把部分子横向分离和纵向动量分数组织成振幅。与分布函数不同，波函数保留粒子数和自旋投影信息。一个示意性的分解为
\begin{equation}
  \Psi_g(x,\bT)
  =\sum_{n}\int\prod_i\dd x_i\,\dd^2\mathbf k_{T,i}\,
  \delta\left(1-\sum_ix_i\right)
  \delta^2\left(\sum_i\mathbf k_{T,i}\right)
  \psi_n(x_i,\mathbf k_{T,i}),
  \label{eq:tmd-wavefunction}
\end{equation}
其中 \(n\) 标记Fock扇区。格点准TMDWF通过类空Wilson线提供可计算的Euclidean对应，但必须经过软函数、rapidity和匹配。波函数的领头Fock态近似不能自动给出完整分布函数，因为高阶Fock态和多粒子关联可能在小 \(x\) 或大 \(b_T\) 处重要。

\section{胶子自旋与极化TMD}

胶子TMD按核子自旋和胶子极化分类。非极化、螺旋度和线偏振胶子分布对应不同的场强投影。横向极化核子还可能产生T-odd胶子分布。其算符需要插入自旋相关投影和适当的Wilson线，某些结构对过程方向敏感。

极化TMD的常见错误是把共线螺旋度PDF的投影直接搬到TMD。TMD多出的横向指标、规范链和软函数会改变混合与rapidity结构。对胶子，伴随表示还意味着极化投影必须与色单态收缩同时定义。任何极化结果都应报告：
\begin{enumerate}
  \item 核子自旋方向与投影；
  \item 胶子极化张量；
  \item Wilson线方向；
  \item 软函数和CS核的表示；
  \item T-odd项的符号约定。
\end{enumerate}

\section{T-odd、过程依赖与符号约定}

T-odd分布不能仅由局域算符的实部得到；它们与初/末态相互作用和Wilson线方向有关。对Drell--Yan与SIDIS，相应方向可能给出符号差异。格点计算如果只模拟一个staple方向，不能直接宣称得到过程无关的T-odd函数。

一个可检验的符号关系应写成
\begin{equation}
  f_{1T}^{\perp,g\,[{\rm DY}]}(x,\bT)
  =\pm f_{1T}^{\perp,g\,[{\rm SIDIS}]}(x,\bT)
  +\text{方案项},
  \label{eq:todd-sign}
\end{equation}
具体符号取决于Wilson线、软函数和自旋定义。论文必须把“裸矩阵元的符号变化”与“物理T-odd分布的符号关系”分开。没有同方案软减除时，裸量符号差异不构成物理结论。

\section{胶子TMD与EIC观测量}

电子-离子对撞机可通过双强子产生、背对背关联、重夸克和喷注等过程探测胶子TMD。实验观测量通常对 \(x\)、\(Q^2\)、\(q_T\)、快度和自旋不对称敏感。理论输入包括：
\begin{itemize}
  \item 胶子TMD及CS演化；
  \item 软函数和过程相关Wilson线；
  \item 共线PDF与 fragmentation functions；
  \item 微扰系数和重正化标度；
  \item 非微扰Sudakov与 \(b_T\) 模型。
\end{itemize}
格点结果若只给出 \(b_T/a\le4\) 的矩阵元，不能直接填充实验所需的全 \(b_T\) 区间。应明确哪些区间由格点约束、哪些由模型外推、哪些完全未验证。

\section{与实验闭环的最小数据产品}

为使格点结果可被现象学使用，至少应提供：
\begin{enumerate}
  \item 中心值与协方差矩阵；
  \item \(x\)、\(b_T\)、\(\mu\)、\(\zeta\) 的完整网格和单位；
  \item 重正化、软减除和匹配方案；
  \item 高扭度、连续极限和外推误差；
  \item 可插值格式与版本哈希；
  \item 明确的适用范围和禁止外推区间。
\end{enumerate}
只给一张中心曲线图不足以进行全局分析。

\section{高级专题的证据矩阵}

\begin{table}[htbp]
\centering
\caption{高级专题的理论状态与证据要求。}
\label{tab:advanced-status}
\small
\begin{tabularx}{\textwidth}{L{2.7cm}L{3.8cm}Y}
\toprule
专题 & 当前可控部分 & 仍需证据 \\
\midrule
端点 \(x\to1\) & 共线阈值重求和与部分匹配 & 格点Fourier端点、高扭度和软修正 \\
小 \(x\) & RGE与饱和框架 & 更大 \(P_z\)、横向非微扰核和混合 \\
TMDWF & 介子准TMDWF先例 & 核子胶子、软函数和高Fock态 \\
胶子自旋 & 共线矩与部分投影 & 完整TMD表示、T-odd与过程依赖 \\
EIC & 因子化与观测量框架 & 全 \(b_T\)、方案统一的格点输入 \\
\bottomrule
\end{tabularx}
\end{table}

\begin{checkpointbox}{高级专题发布门}
\begin{enumerate}
  \item 是否明确目标 \(x\) 和 \(b_T\) 区间；
  \item 是否区分共线、TMD和波函数对象；
  \item 是否保留胶子表示与Wilson线方向；
  \item 是否对端点和小 \(x\) 单独估计模型误差；
  \item 是否提供可被实验使用的协方差和适用域。
\end{enumerate}
\end{checkpointbox}

\clearpage
\chapter{重正化、匹配与胶子 TMD 的系统构造}

\section{从共线分布到横动量依赖分布}

共线胶子分布描述纵向动量份额，而 TMD-PDF 同时保留横向分离。采用伴随表示
Wilson 线的非极化胶子定义可以写成
\begin{equation}
  x\,g(x,\mu^2)
  =\frac{1}{2\pi P^+}\int \dd\xi^- e^{-ixP^+\xi^-}
  \langle P|F^{+i,a}(\xi^-)W^{ab}_{\mathrm{adj}}(\xi^-,0)
  F^{+i,b}(0)|P\rangle_{\mu},
  \label{eq:b10-lightcone-pdf}
\end{equation}
其中 $P^+$ 是光锥正方向动量，$i=1,2$ 是横向指标，$W_{\mathrm{adj}}$ 是
保证规范不变性的路径有序 Wilson 线。式~\eqref{eq:b10-lightcone-pdf} 的
物理意义是沿光锥分离的胶子场强关联；它定义在 Minkowski 光锥上，因此不能
直接替换为欧氏格点上的等时关联。若在光锥规范中令 $A^+=0$，Wilson 线化为
单位矩阵，但规范无关的物理结论不能依赖这一选择。

TMD 定义在同一纵向结构上增加横向位移 $\bperp$ 与 staple 规范链：
\begin{equation}
  x\,g(x,\bperp;\mu,\zeta)
  =\frac{1}{2\pi P^+}\int \dd\xi^- e^{-ixP^+\xi^-}
  \langle P|F^{+i}(\xi^-+\bperp)W_{\mathrm{staple}}
  F^{+i}(0)|P\rangle_{\mu,\zeta}.
  \label{eq:b10-tmd-def}
\end{equation}
这里 $\zeta$ 是 Collins--Soper 快度标度。与式~\eqref{eq:b10-lightcone-pdf}
相比，式~\eqref{eq:b10-tmd-def} 多出三个不可省略的结构：横向分离、同几何
软减除以及 Wilson 线方向造成的过程依赖。只有当 $q_\perp\ll Q$ 且横动量
因子化成立时，TMD 才能作为截面中的独立非微扰输入；把任意大 $q_\perp$ 的
数据写成 TMD 会越过因子化窗口。

在格点侧，常用准 TMD 矩阵元 $\tilde g(x,\bperp,\Pz)$ 逼近式
\eqref{eq:b10-tmd-def} 的光锥对象。其大动量展开应理解为
\begin{equation}
  \tilde g(x,\Pz)
  =\int_{-1}^{1}\frac{\dd y}{|y|}
  C\!\left(\frac{x}{y},\frac{\mu}{|y|\Pz}\right)g(y,\mu)
  +\Ocal\!\left(\frac{\LQCD^2}{\Pz^2},
  \frac{\LQCD^2}{x^2\Pz^2}\right),
  \label{eq:b10-lamet}
\end{equation}
其中 $C$ 是微扰匹配核，后两项分别是普通高扭度/靶质量修正和小 $x$ 增强。
该式是有效场论展开，不是对所有 $\Pz$ 的恒等式。实际分析至少要求
$\Pz\gtrsim1.5~\mathrm{GeV}$ 且 $x\Pz\gg\LQCD$；小 $x$ 端若
$x\Pz$ 落入非微扰区，匹配核与 Landau 极点问题不能被统计误差掩盖。

\section{欧氏算符、投影与 Wilson 线几何}

令 Wilson 线沿 $z$ 方向，横向方向为 $i,j\in\{x,y\}$。胶子双局域算符可写为
\begin{equation}
  M^{\mu\lambda;\nu\rho}(z,\bperp)
  =F^{\mu\lambda}(z\hat z+\bperp)U_{\mathrm{adj}}
  F^{\nu\rho}(0),
  \label{eq:b10-bilocal}
\end{equation}
其中 $U_{\mathrm{adj}}$ 是连接两端场强的伴随表示路径。非极化通道采用的
可乘组合是
\begin{equation}
  \Ocal_U(z,\bperp)
  =M^{ti;ti}+M^{tj;tj}-2M^{ij;ij},
  \qquad
  \Ocal_{U,\mathrm{even}}=\Ocal_U(+\!z)+\Ocal_U(-\!z).
  \label{eq:b10-unpolarized}
\end{equation}
若不除以 2，$z\to0$ 时奇数方向组合恒为零，可作为实现与数据布局的直接检查。
螺旋度通道的理论目标是
\begin{equation}
  \Ocal_{H,\mathrm{odd}}
  =\bigl(M^{ti;ti}_H+M^{tj;tj}_H-2M^{ij;ij}_H\bigr)_{+z}
  -\bigl(\cdots\bigr)_{-z}.
  \label{eq:b10-helicity}
\end{equation}
注意，存储格式中常见的 $M^{ti;ti}+M^{tj;tj}+2M^{ij;ij}$ 是诊断组合，
不能替代式~\eqref{eq:b10-helicity} 的物理螺旋度投影。

TMD staple 的格点几何由纵臂、横向翼和反向纵臂组成。若纵臂长度为 $L$，
并要求 $L$ 与横向位移保持在格点体积内，则 $\bperp\to0$ 时 staple 应连续
退化为直线 Wilson 线。这个极限提供三方面交叉检查：共线与 TMD 的公共归一、
$z=0$ 的局域性以及 even/odd 投影的一致性。若路径绕周期边界、纵臂与横翼
共线，或两端使用的方向指标不一致，则所有后续“规范不变”检查都不再对应同一
算符。

\section{三种尺度：UV、快度与动量}

胶子 TMD 同时含 $\mu$、$\zeta$ 与 $\Pz$。三者不能相互替代：$\mu$ 是 UV
重整化标度，$\zeta$ 是快度演化标度，$\Pz$ 是构造准光锥关联的动量。软函数
与 Collins--Soper 核的配对结构为
\begin{align}
  \frac{\dd}{\dd\ln\sqrt{\zeta}}\ln F(x,\bperp;\mu,\zeta)
  &=K(\bperp,\mu),
  \label{eq:b10-cs-evolution}\\
  \frac{\dd}{\dd\ln\mu}K(\bperp,\mu)
  &=-\gamma_\zeta(\alpha_s),
  \label{eq:b10-cs-mu}
\end{align}
其中 $F$ 是选定软减除方案后的 TMD，$K$ 是 Collins--Soper 核，
$\gamma_\zeta$ 是其 UV 反常量纲。式~\eqref{eq:b10-cs-evolution} 只在软函数、
TMD 和硬核使用同一快度方案时成立。若把一个方案的软因子除以另一个方案的
裸 TMD，或把 $\mu$ 演化与 $\zeta$ 演化混成一个拟合常数，可积性条件会被
人为破坏。

\section{梯度流重整化与小流时间展开}

有限流时间 $\tau$ 的 flowed 算符在连续极限中具有受控的 UV 性质，但它定义
的是带尺度 $\sqrt{\tau}$ 的新方案，并不自动等于 $\MSbar$。小流时间展开
写成
\begin{equation}
  \Ocal_{\mathrm{GF}}(\tau,z)
  =\sum_k C_k(\tau,\mu,z)\Ocal_k^{\MSbar}(z,\mu)
  +\Ocal(\tau\LQCD^2,\tau/z^2).
  \label{eq:b10-sfte}
\end{equation}
对 Wilson 线算符，线性自能可以组织为端点系数与指数线因子。按本批次来源的
一圈约定，
\begin{align}
  c_{\parallel\perp}(\tau,\mu)
  &=1+\Ocal(\alpha_s^2),
  \label{eq:b10-cpara}\\
  c_{\perp\perp}(\tau,\mu)
  &=1+\frac{\alpha_s(\mu)C_A}{4\pi}L_\tau(\mu)
  +\Ocal(\alpha_s^2),
  \label{eq:b10-cperp}\\
  \delta m_A(\tau)
  &=-\frac{\alpha_s(\mu)C_A}{4\pi}\frac{\sqrt{2\pi}}{\sqrt{\tau}}
  +\Ocal(\alpha_s^2),
  \label{eq:b10-dm}\\
  L_\tau(\mu)&=\ln\!\left(2\mu^2\tau e^{\gamma_E}\right).
  \label{eq:b10-Ltau}
\end{align}
于是双局域算符的一圈转换可概括为
\begin{equation}
  \Ocal_{ij}^{\MSbar}(z,\mu)
  =\frac{e^{-\delta m_A|z|}}{c_i c_j}
  \Ocal_{ij}^{\mathrm{GF},R}(\tau,z)+\Ocal(\tau),
  \qquad
  \begin{cases}
    c_i=c_j=c_{\perp\perp}, & U,\\
    \{c_i,c_j\}=\{c_{\parallel\perp},c_{\perp\perp}\}, & H.
  \end{cases}
  \label{eq:b10-endpoint}
\end{equation}
这里 $|z|$ 必须以 $\mathrm{GeV}^{-1}$ 计；若从格距单位换算，需显式乘以
$a[\mathrm{fm}]\times5.0677$。式~\eqref{eq:b10-endpoint} 只给出一圈端点
转换，不包含软函数、快度演化或连续极限。尤其不能把
\texttt{light\_cone\_matching} 的字符串标签当作匹配已经完成。

SFTE 的适用窗口不是“流时间越大越好”，而应同时满足
\begin{equation}
  \frac{a^2}{\tau}\ll1,\qquad
  \sqrt{8\tau}\ll|z|,\qquad
  \tau\Pz^2\ll1,\qquad
  \tau\LQCD^2\ll1.
  \label{eq:b10-sfte-window}
\end{equation}
第一项控制格点 UV 与流尺度的分离，第二项保证 Wilson 线仍远长于流化半径，
后两项要求流尺度低于强子与动量标度。若为了缩小误差而把大 $\tau$ 点纳入
线性外推，可能把 $\Ocal(\tau)$ 截断变成了中心值漂移；因此必须同时报告小流
窗口、大流诊断和窗口敏感性。

\section{非微扰重整化：自重整化、hybrid 与 RI/MOM}

自重整化利用零动量矩阵元的短距微扰形状和长距非微扰参数化。其思想可写成
\begin{equation}
  h_B(z,\Pz)\sim e^{-\delta m|z|}\,
  Z_{\mathrm{MS}}(z,\mu)\,B(z;\boldsymbol p),
  \label{eq:b10-zr-sketch}
\end{equation}
其中 $Z_{\mathrm{MS}}$ 在某些小 $z$ 节点由 NLO 匹配控制，$B$ 由大 $z$
行为拟合。具体实现采用含 $z/a$、$\ln(a\LQCD)$、$(m_0+m_2a^2)z$ 以及
$fa+f_2a^2$ 项的全局参数化，并通过逐 bootstrap/jackknife 样本重拟合传播
不确定度。该方案的优势是不依赖 Landau 规范固定；失效点是短距窗口、大 $z$
尾部模型和流时间不能同时受控时，参数化会以“吸收系统误差”的方式制造稳定
假象。

hybrid 方案在过渡尺度 $z_s$ 两侧拼接
\begin{equation}
  h_R(z,\Pz)=
  \begin{cases}
    \dfrac{h_B(z,\Pz)}{h_B(z,\Pz=0)}, & |z|\le z_s,\\[1.1em]
    \dfrac{h_B(z,\Pz)}{Z_R(z)}\,\eta_s, & |z|>z_s,
  \end{cases}
  \qquad
  \eta_s=\frac{Z_R(z_s)}{h_B(z_s,\Pz=0)}.
  \label{eq:b10-hybrid}
\end{equation}
短距比值消去与 $z$ 无关的 UV 因子和多数量子修正，长距 $Z_R$ 保留非微扰
信息。通常取 $z_s\sim0.2$--$0.4~\mathrm{fm}$，但 $z_s$ 是方案参数，不能
按点挑选“误差最小”的过渡位置。式~\eqref{eq:b10-hybrid} 在 $z_s$ 处连续
只说明拼接代数一致，不代表短距微扰和长距模型同时可信。

RI/MOM 通过 Landau 规范下的离壳夸克 Green 函数定义 $Z_O$，在窗口
$\LQCD\ll\mu_R\ll1/a$ 内把非微扰结果转换到 $\MSbar$。对非局域 Wilson
线，$z$ 与 $\mu_R$ 可比拟时幂次修正会污染匹配，因此通常只适合
$z\lesssim0.3~\mathrm{fm}$。B10 收录的 2022 年学位论文摘要报告：在约
$0.032$--$0.121~\mathrm{fm}$ 的多个格距和三种费米子作用量上，RI/MOM
重整化后仍可能残留 Wilson 线线性发散，且随格距减小而恶化。这个结论仅来自
摘要级阅读，本批次没有核对论文正文的投影、窗口和统计误差，不能单独据此
宣判现代 ratio、hybrid 或梯度流方案。

两种主线必须保持互斥：
\begin{align}
  \text{GF 链:}\quad
  &h^{\mathrm{GF,bare}}\longrightarrow h^{\MSbar}\longrightarrow
  \text{Fourier/matching},
  \label{eq:b10-chain-gf}\\
  \text{hybrid 链:}\quad
  &h_B\longrightarrow h_R^{(Z_R,\mathrm{hybrid})}\longrightarrow
  \text{Fourier/matching}.
  \label{eq:b10-chain-hybrid}
\end{align}
两者的输入可以是同一裸矩阵元，但输出方案不同。把 $\delta m_A$、比值因子、
$Z_R$ 或 Wilson 圈减除任意相乘，通常不及于任何定义良好的方案；分析 metadata
必须明确记录实际使用了哪一条链。

\section{Fourier 重构、NLO 卷积与软减除}

共线胶子准 PDF 的 sin 变换为
\begin{equation}
  \tilde g(x,\Pz)
  =\frac{2\Pz}{x}\int_0^{z_{\max}}\dd z\,
  h_R(z,\Pz)\sin(x\Pz z).
  \label{eq:b10-sin}
\end{equation}
准 TMD 则使用复相位与 Hermitian 延拓：
\begin{equation}
  x\,\tilde g(x,\bperp,\Pz)
  =\frac{1}{\mathcal N}\int\frac{\dd z}{2\pi\Pz}
  e^{-ix\Pz z}h_R(z,\bperp,\Pz),
  \qquad
  \mathcal N=\frac{\Pz^2}{P_t^2},
  \label{eq:b10-quasi-tmd}
\end{equation}
其中 $P_t=\sqrt{M_N^2+\Pz^2}$。在正半轴数据上采用
$h_R(-z)=h_R(z)^*$，等价于余弦乘实部、正弦乘虚部的组合。两种 Fourier
表示在有限 $z_{\max}$ 下都引入低通卷积与 Gibbs 振荡；$z_{\max}\to\infty$
并不是数值上可实现的极限，必须用大 $z$ 尾部模型或正则化估计缺失贡献。

NLO 胶子匹配的离散实现可组织为
\begin{equation}
  Z_{ij}=\delta_{ij}+\frac{\alpha_s C_A}{2\pi}
  \frac{x_i}{|x_i|}\Delta x
  \left[\frac{g_{ij}}{y_j}
  -\delta_{ij}y_i\sum_k\frac{g_{ik}}{y_k^2}\right],
  \qquad
  h^{\mathrm{PDF}}=Z^{-1}h_R^0.
  \label{eq:b10-zij}
\end{equation}
核 $g_{ij}$ 分为 $\xi=x_i/y_j<0$、$0<\xi<1$ 与 $\xi>1$ 三区，并在
$\xi=1$ 用 plus-distribution 或主值处理。矩阵求逆不是数学装饰：它同时
恢复端点 $\delta(1-\xi)$ 与软胶子求和规则。若只保留正半轴而丢掉 $\xi<0$
支路，会出现不随网格减小而消失的归一偏移；若 x 网格含精确零点，则必须停止
计算，而不是用任意有限值替代主值约定。

TMD 混合匹配把硬核、软减除和 CS 演化放在同一式中：
\begin{equation}
  x\tilde g(x,\bperp,\Pz)
  =\int\frac{\dd y}{|y|}
  C^{\mathrm{hybrid}}\!\left(\frac{x}{y},\lambda_s,
  \frac{\mu}{y\Pz}\right)
  \frac{y\,g(y,\bperp)}{\sqrt{S_I(\bperp,\mu)}}\,
  \exp\!\left[\frac12K(\bperp,\mu)
  \ln\frac{\zeta_z}{\zeta}\right].
  \label{eq:b10-tmd-match}
\end{equation}
式~\eqref{eq:b10-tmd-match} 的因子顺序对应物理层次，而不是任意约定：
$C^{\mathrm{hybrid}}$ 是微扰硬核，$S_I$ 是内禀软函数，$K$ 把参考快度演化
到目标 $\zeta$。只有三者来自同一颜色表示、同一 Wilson 线几何和同一快度
方案时，乘积才是方案无关的物理 TMD。工程上的推荐顺序是“归一、软减除、
CS 演化、NLO 卷积、联合外推”，并保存每一步的 metadata。

在短距区域，CS 核的一圈结构可写为
\begin{equation}
  K(\bperp,\mu)
  =-\frac{\alpha_s(\mu)C_A}{\pi}
  \ln\frac{\bperp^2\mu^2}{4}+\Ocal(\alpha_s^2),
  \label{eq:b10-cs-one-loop}
\end{equation}
常数项依赖方案。格点比值提取则使用两个动量：
\begin{equation}
  K(\bperp)\simeq
  \frac{\ln\!\left[h_R(z_0,\bperp,P_{z1})/
  h_R(z_0,\bperp,P_{z2})\right]}{\ln(P_{z1}/P_{z2})}.
  \label{eq:b10-cs-ratio}
\end{equation}
式~\eqref{eq:b10-cs-ratio} 只有在 $z_0$ 位于共同短距窗口、两个动量都满足
因子化且硬核差异受控时才有意义；否则它只是有效斜率，不是普适 CS 核。

\section{尚未闭合的归一化条件}

连续场强与代码 Clover 的场归一化之间可能存在树图因子 $\kappa_F$。若代码
输出没有显式包含 $g^2$，双局域算符应相应的 $\kappa_F^2$ 或等价转换吸收。
在本批次来源中，$\kappa_F$ 的微扰确定被明确列为未闭合链，而当前实现按
$\kappa_F=1$ 记录。因此，式~\eqref{eq:b10-endpoint} 的系数只有在场强归一、
Clover 投影和 Wilson 线表示全部对齐后才可用于物理比较。这个缺口不能由
ISO 误差或统计 bootstrap 自动修复。

\chapter{从组态到物理胶子 TMD 的工作流、验证与历史}

\section{端到端流程与停止条件}

一个可审计的格点胶子 TMD 流程应把数学对象、数据对象和物理结论分开。表
\ref{tab:b10-workflow} 给出最小流程；任一阶段的停止条件未满足，就只能保留
中间产物，不能把后续拟合结果称为物理 TMD。

\begin{longtable}{@{}p{0.06\textwidth}p{0.30\textwidth}p{0.50\textwidth}@{}}
\caption{B10 建议纳入第 18 章的端到端工作流。}\label{tab:b10-workflow}\\
\toprule
步 & 对象 & 进入下一步的条件 \\
\midrule
\endfirsthead
\toprule
步 & 对象 & 进入下一步的条件 \\
\midrule
\endhead
1 & ILDG thin links、几何与元数据 & 链接近 SU(3)、行列式、体积与边界语义通过；失败日志保留。\\
2 & 四维 Wilson flow & RK3 步长收敛；plaquette 单调诊断与流后有限性通过。\\
3 & flowed Clover、对偶场强与 Wilson 线 & traceless/legacy 投影明确；$\pm z$ 与 staple 几何合法；$z=0$ odd 为零。\\
4 & 六组件 OPE 与 schema 产物 & NPY/JSON 原子配对；轴、dtype、哈希和 \texttt{not\_applied} 清单完整。\\
5 & 两点函数与三点函数 & 源汇、宇称投影和真空扣除固定；$t_{\mathrm{sep}}$ 扫描有端点切割。\\
6 & ratio-v2 与 delete-one jackknife & 共同 nopol 分母、复数协方差、NaN 掩码和 selector 均通过。\\
7 & 基态矩阵元 $h(z,\bperp,\Pz)$ & plateau/多态拟合在预注册窗口内稳定；不按点挑窗口。\\
8 & 重整化方案 & 明确选择式~\eqref{eq:b10-chain-gf} 或式~\eqref{eq:b10-chain-hybrid}，方案因子不混用。\\
9 & Fourier、CS、软函数与 NLO 匹配 & 同一表示/几何/快度方案；$z_{\max}$ 与尾部模型误差可见。\\
10 & 联合外推与物理发布 & $a,\Pz,m_\pi,L$ 的窗口和相关性受控；所有未闭合项写入结果边界。\\
\bottomrule
\end{longtable}

\section{ratio-v2：共同分母、复数相位与掩码}

对每个通道 $\chi$，三点估计量应由配置/源协方差定义，而不是先对各组态取比值。
本批次新版采用的写法是
\begin{align}
  C_{3,\chi}
  &=\langle O_\chi C_{2,\chi}\rangle_{\mathrm{conf,src}}
  -\langle O_\chi\rangle_{\mathrm{conf,src}}
  \langle C_{2,\chi}\rangle_{\mathrm{conf,src}},
  \label{eq:b10-c3}\\
  C_{2,\mathrm{den}}
  &=\langle C_{2,\mathrm{nopol}}\rangle_{\mathrm{conf,src}},
  \qquad
  R_\chi=\frac{C_{3,\chi}}{C_{2,\mathrm{den}}}.
  \label{eq:b10-ratio}
\end{align}
非极化和螺旋度的 numerator 可以使用各自的通道，但分母必须共享 nopol。
尤其不能把 ensemble mean 的 pol35 作为近零分母。对螺旋度，完整复数必须先
经过协方差、真空扣除、ratio 和 jackknife，最后再做
\begin{equation}
  \operatorname{Re}[-\mathrm{i}R_H]=\operatorname{Im}R_H.
  \label{eq:b10-phase}
\end{equation}
若提前把 pol35 截成实部或虚部，式~\eqref{eq:b10-c3} 中的协方差项已经被
改变，最后再取虚部不能恢复。

delete-one jackknife 的每一样本应重新计算分子、共同分母和真空扣除：
\begin{equation}
  R_{(-k)}
  =\frac{C_{3,(-k)}}{C_{2,(-k),\mathrm{den}}},
  \qquad
  \sigma_{\mathrm{Re}}^2=\frac{N-1}{N}\sum_k
  \left(\operatorname{Re}R_{(-k)}
  -\overline{\operatorname{Re}R}\right)^2,
  \label{eq:b10-jk}
\end{equation}
虚部误差同理，实部与虚部不可合并成没有定义的单一标量。无效插入位置
$0\le\tau_{\mathrm{ins}}\le t_{\mathrm{sep}}$ 之外必须填 NaN，并保存
\texttt{valid\_insertion\_mask}。填零会把“没有观测”误当作“观测到零”，
从而系统性拉偏均值和误差。

\section{可复现验证与不能越过的边界}

新版工作流在来源中登记了三级证据：提交/标签说明、九项独立契约和全量回归。
其中可登记的断言包括：复杂 ratio 与参考估计器逐项一致、端点转换与参考一致、
随机局域规范变换下双局域 trace 不变、RK3 细步长收敛，以及全量回归通过。
这些数字是 U077 快照内的验证声明；B10 只做了文档和 PDF 元数据核对，没有
执行该快照中的安装、GPU、生产数据处理或测试命令。因此本草案把它们写成
“来源内登记证据”，而不是本次复跑结果。

同样的边界适用于物理结果。有限流时间的 flowed bare OPE 或 bare ratio 是
有意义的中间可观测量，但还不是物理胶子 PDF/TMD。当前材料明确把以下对象
列为未闭合或需真实生产接线：参考组态共同 manifest、单态 gluon--quark
mixing、与 schema-v2 ratio 相接的软函数/快度减除、完整 NLO 匹配、多格距
连续极限和场强树图归一 $\kappa_F$。只有在这些条件同时满足后，才可写出带
统计和系统误差的 $xg(x,\bperp;\mu,\zeta)$。

\section{历史谱系：从 RI/MOM 到梯度流 TMD}

本批次资料呈现了清晰但不等权的演进链。2022 年硕士学位论文以准 PDF 的 UV
线性发散为中心，摘要报告 RI/MOM 不能在所有格距上完全消除 Wilson 线发散；
2023 年硕士学位论文把 TMD-PDF、Wilson loop、HMC/Chroma 与 jackknife/
bootstrap 组织成格点工作流，但本批次只核对了题录和目录。两者提供历史问题
与工作流线索，不能替代 U077 中的显式公式。

旧版 \texttt{理论解析与工作流.tex} 采用共线胶子 PDF 主线，以 Eq.(20) 的
矩阵元、Eq.(25) 的非极化 OPE、$C_3/C_2$、Fourier 和 NLO matching 组织全文。
新版 \texttt{理论解析与工作流-v20260911.tex} 明确把它标为过时版本，并将主线
升级为梯度流重整化的胶子 TMD-PDF：新增 `SFTE` 端点转换、ratio-v2 复数统计、
CS 软结构、schema-v2 和验证边界。工作流版本因此应去重，而不是把两套公式
并列成两个物理结论。

从方法史看，可把演进概括为四步：先由 twist-2 定域算符计算低阶 Mellin 矩；
再由 LaMET/quasi-PDF 获得 $x$ 依赖；随后把直线 Wilson 线推广到 staple 并加入
软函数、CS 核和 TMD 因子化；最后用梯度流、系统误差窗口和多格距外推控制 UV
与连续极限。这个链条的每一步都增加新的标度和新的失效条件；后一步不会自动
修正前一步的错误。

\section{系统误差与失败模式}

\begin{tabularx}{\textwidth}{@{}p{0.22\textwidth}X@{}}
\toprule
来源 & 物理含义与最低报告要求 \\
\midrule
离散化 $a^2$ & Clover 改进后仍需至少三格距或明确说明只有单格距；流化半径必须
沿固定物理序列变化。 \\
流时间 SFTE & 同时检查式~\eqref{eq:b10-sfte-window}；小流截距与大流诊断分开，
不能以统计误差最小替代窗口控制。 \\
有限体积与激发态 & 报告 $M_\pi L$、$t_{\mathrm{sep}}$、端点切割和多态拟合；
disconnected 近似不自动消除所有污染。 \\
Fourier 截断 & 报告 $z_{\max}$、大 $z$ 尾部模型、Gibbs 振荡和逆问题正则化；
负值不能一律归因于统计。 \\
软函数/快度 & 报告 $S_I$、$\zeta$、Wilson 线方向与表示；缺少软减除时不能把
准 TMD 称为物理 TMD。 \\
匹配与混合 & 报告 $NLO$ 阶数、$Z_{ij}$、负半轴、$\mu$ 变异和单态混合；
正半轴截断或方案混用会造成不可忽略的伪偏移。 \\
方案与归一 & 固定 GF 或 hybrid 链；记录 $\kappa_F$、场强投影、Clover/对偶
约定和 metadata 中的未应用项。 \\
\bottomrule
\end{tabularx}


\clearpage
\chapter{当前理论进展与证据边界}

\section{PyQCD 文档中的实现层增补}

本节把 \path{refer/git-rep/PyQCD/docs/} 中的方法性笔记压缩成若干可检验的增量点。这些文档是二次综合与实现审计，不是新的实验测量；涉及具体物理结论时仍回到原始论文和格点数据。表~\ref{tab:pyqcd-docs-map} 给出选用文档与本文新增内容的对应关系，正文只保留经交叉核对后能够落到公式、窗口或状态门上的部分。

\begin{table}[htbp]
\centering
\caption{PyQCD 文档中与本文直接相关的增量材料及其使用边界。}
\label{tab:pyqcd-docs-map}
\small
\begin{tabularx}{\textwidth}{L{3.4cm}L{4.0cm}Y}
\toprule
文档主题 & 增量要点 & 本文中的落点与边界 \\
\midrule
\textit{大动量有效理论} & 极限交换、解析延拓、普适性类 & 补充 LaMET 的有效场论图像；不把类比当成证明 \\
\textit{胶子算符} & 36 分量、辅助场、三角混合、\(2\times2\) 匹配 & 区分 UV 可乘性重正化与匹配层味混合 \\
\textit{Wilson 线} & 线性/交叉/pinch-pole 结构与四类方案 & 形成重正化策略决策表 \\
\textit{3pt 构造} & 顺序源、广义 perambulator、多态拟合 & 补充成本、缓存和激发态控制 \\
\textit{外推} & 大 \(\lambda\) 外推、derivative 方法与联合外推 & 补充 Fourier 尾部与极限顺序 \\
\textit{误差统计} & 全协方差链与独立误差平方和 & 明确两条误差传播路径的适用域 \\
\textit{胶子 TMD 梯度流} & 从组态到 TMD 的十步流程与 SFTX & 用于实现层验收，不作为数值结果 \\
\textit{dev20 审计} & 当前 11 组态诊断链、fail-closed 状态 & 用于给出实现门控与未闭合项 \\
\bottomrule
\end{tabularx}
\end{table}

{\small\raggedright
表中完整文件名依次为
格点QCD中的大动量有效理论.tex、
格点QCD中的胶子算符.tex、
格点QCD中的Wilson线.tex、
格点QCD中的3pt构造.tex、
格点QCD中的外推.tex、
格点QCD中的误差统计.tex、
胶子TMD梯度流连续极限.tex 和
absorb\_dev20\_gluon\_tmd\_audit\_20261022.md，
均位于 \path{refer/git-rep/PyQCD/docs/}。\par}

\subsection{LaMET 的极限结构、普适类与准算符选择}

LaMET 可以在形式上类比为强子版本的重夸克有效理论：HQET 以 \(m_Q\to\infty\) 组织 \(1/m_Q\) 展开，LaMET 以 \(P_z\to\infty\) 组织 \(1/P_z\) 展开。最低阶匹配写成
\begin{equation}
  F(P_z/\Lambda)
  =Z\!\left(\frac{P_z}{\Lambda},\frac{\Lambda}{\mu}\right)f(\mu)
  +\mathcal O\!\left(\frac{\Lambda^2}{P_z^2}\right),
  \label{eq:lamet-eft}
\end{equation}
其中 \(F\) 是格点准观测量，\(f\) 是光锥有效理论中的物理量，\(Z\) 吸收两种极限顺序造成的 UV 差异。关键点是
\begin{equation}
  \lim_{P_z\to\infty}\lim_{a\to0}
  \;\neq\;
  \lim_{a\to0}\lim_{P_z\to\infty}.
  \label{eq:lamet-noncommute}
\end{equation}
格点计算必须先有 UV 截断再进行大动量外推，而光锥定义在概念上先取无穷动量再完成重整化。两种顺序的 IR 物理应相同，差别由匹配核和幂修正组织；因此必须同时满足
\begin{equation}
  \Lambda_{\rm QCD},M_N\ll P_z\ll a^{-1}.
  \label{eq:lamet-scales}
\end{equation}

同一光锥算符可以由多个 Euclidean 准算符逼近。把在 \(P_z\to\infty\) 下落到同一光锥极限的准算符集合称为普适类：
\begin{equation}
  \mathcal U_{\rm LC}
  =\left\{\mathcal O_i^{(E)}:
  \lim_{P_z\to\infty}Z_i^{-1}\mathcal O_i^{(E)}
  =\mathcal O^{\rm LC}\right\}.
  \label{eq:universalityclass}
\end{equation}
实际选择准算符时，应优先比较各成员在有限 \(P_z,a\) 下的幂修正、统计噪声和算符混合，而不是只选形式上最简单的一个。另一个常被忽视的要点是解析延拓：不能把 Minkowski 外动量直接代入 Euclidean 积分后宣布等于光锥结果；必须保留 Euclidean 极点结构，完成积分后再解析延拓。

\subsection{36 分量、三角混合与两类“混合”}

纯胶子准算符的 UV 结构可按四个 Lorentz 指标中 \(z\) 分量的个数 \(s\) 分类。令 \(C_g\) 表示线性自能，\(Z_{wg}\) 表示 Wilson 线端点因子，\(Z_{vg1},Z_{vg2}\) 表示两类 \(z\) 指标因子，则最小结构可写成
\begin{table}[htbp]
\centering
\caption{胶子准算符按 \(z\) 指标数的重正化因子结构。}
\label{tab:z-index-classes}
\small
\begin{tabularx}{\textwidth}{L{1.5cm}L{4.0cm}Y}
\toprule
\(s\) & 代表结构 & 必要的重正化因子 \\
\midrule
0 & 纯横向组合，如 \(F^{ti}F^{ti}\)、\(F^{ij}F^{ij}\) & \(e^{-C_g|z|}\)、\(Z_{wg}\) \\
1 & 一个纵向指标，如 \(F^{ti}F^{zi}\) & 另有 \(Z_{vg1}^{1/2}\) \\
2 & 两个纵向指标，如 \(F^{zi}F^{zi}\)、\(F^{z\mu}F^z_\mu\) & 另有 \(Z_{vg2}\) \\
\bottomrule
\end{tabularx}
\end{table}

辅助伴随重夸克场方法进一步表明，场强端点、Wilson 线和可能的辅助场算符可形成三角混合：
\begin{equation}
  \begin{pmatrix}J_{1,R}\\J_{2,R}\\J_{3,R}\end{pmatrix}
  =
  \begin{pmatrix}
    Z_{11} & Z_{12} & Z_{13}\\
    0 & Z_{22} & Z_{23}\\
    0 & 0 & Z_{33}
  \end{pmatrix}
  \begin{pmatrix}J_1\\J_2\\J_3\end{pmatrix},
  \label{eq:triangularmixing}
\end{equation}
其中 \(J_3\) 与辅助场运动方程相关，在物理矩阵元中可以消失。更重要的工程结论是：尽管全部 Lorentz 分量很多，真正需要独立追踪的可乘性构建块并不多；错误地用一个标量 \(Z\) 吸收所有端点结构，会把不同 \(s\) 分量的反常量纲混在一起\cite{LiMaQiu2019,ZhangGluonLaMET2019}。

UV 重正化阶段与匹配阶段的“混合”必须严格区分。前者的可乘性结构不要求胶子准算符与夸克准算符混合；后者在把准分布匹配到光锥单态 PDF 时出现味混合矩阵：
\begin{keyformula}{UV 可乘性与匹配层味混合}
\begin{align}
  \begin{pmatrix}\widetilde g(x,\mu)\\ \widetilde q_S(x,\mu)\end{pmatrix}
  &=\int_x^1\frac{\dd y}{y}
  \begin{pmatrix}
    C_{gg}\!\left(\frac{x}{y},\frac{\mu}{P_z}\right) &
    C_{gq}\!\left(\frac{x}{y},\frac{\mu}{P_z}\right)\\[2pt]
    C_{qg}\!\left(\frac{x}{y},\frac{\mu}{P_z}\right) &
    C_{qq}\!\left(\frac{x}{y},\frac{\mu}{P_z}\right)
  \end{pmatrix}
  \begin{pmatrix}g(y,\mu)\\ q_S(y,\mu)\end{pmatrix}
  \nonumber\\
  &\quad+\mathcal O\!\left(\frac{\Lambda_{\rm QCD}^2}{P_z^2}\right),
  \label{eq:matching2x2}\\
  \mathcal O_i^{R}(\mu)&=Z_{ij}^{\rm UV}(\mu,\mu_0)\mathcal O_j^{R}(\mu_0),
  \qquad
  \begin{pmatrix}\widetilde g\\ \widetilde q_S\end{pmatrix}
  \neq Z^{\rm UV}\begin{pmatrix}\widetilde g\\ \widetilde q_S\end{pmatrix}.
  \label{eq:uv-vs-matchingmixing}
\end{align}
\end{keyformula}
式~\eqref{eq:uv-vs-matchingmixing} 的第二行不是矩阵等式，而是强调不能把 \(Z^{\rm UV}_{gq}\) 与 \(C_{gq}\) 混成同一个对象。前者属于格点算符的 UV 重正化，后者属于光锥因子化的硬匹配。

\subsection{Wilson 线策略的决策边界}

库内文档把 Wilson 线发散的处理归纳为质量抵消、Wilson 圈、白重整化和 hybrid 四类。它们并非可以任意互换的“四种写法”，而是分别针对不同的发散组合。表~\ref{tab:wilson-strategies} 给出决策边界。

\begin{table}[htbp]
\centering
\caption{Wilson 线重正化策略的适用条件与失败信号。}
\label{tab:wilson-strategies}
\small
\begin{tabularx}{\textwidth}{L{2.5cm}L{4.1cm}L{3.2cm}Y}
\toprule
方案 & 主要消除对象 & 成立条件 & 失败信号 \\
\midrule
质量抵消/辅助 \(z\) 场 & 开放 Wilson 线的线性自能 & 同一路径长度与表示；残余对数另行处理 & 只看到线性项消失，但 \(L_\parallel\) 或 \(z\) 仍无平台 \\
Wilson 圈平方根 & 线性自能、pinch-pole、部分 cusp & 圈与 staple 同几何、同端点、同表示 & \(L_\parallel\) 扫描无平台或不同转角结果不一致 \\
白重整化 & 多格距的线性项与重整子模糊 & 至少多个 \(a\)、可控短距参考点 & 单格距时 \(z_0\) 和参数化不可辨识 \\
Hybrid & 短距 ratio 与长距质量减除 & \(z_s\) 窗口稳定，短距匹配与长距参数分离 & \(z_s\) 或 \(\delta m+m_0\) 变化超过目标误差 \\
RI/MOM & 局域反常量纲与有限混合 & 仅用于局域端点结构，需规范与投影固定 & 小 \(a\) 或大 \(z\) 的格距依赖反而增强 \\
\bottomrule
\end{tabularx}
\end{table}

一个直接可审计的 Wilson 圈比值为
\begin{equation}
  h_R(z,\bT;P_z)
  =\lim_{L_\parallel\to\infty}
  \frac{\tilde h(z,\bT,L_\parallel;P_z;1/a)}
  {\sqrt{Z_{E,R}(\bT;2L_\parallel+z;1/a)}} .
  \label{eq:wl-subtractionratio}
\end{equation}
文档中的五格距验证、Clover/Overlap 交叉检查和 \(L_\parallel\) 平台共同说明该减除可以同时组织线性、pinch-pole 与 cusp 结构；但残余对数仍需短距 ratio、hybrid 或自重整化处理。RI/MOM 的失败并不等于动量减除无意义，而是说明普通局域减除无法替代针对非局域 Wilson 线的同几何结构\cite{Zhang2022TMD,Huo2021,ZhangGauge2024}。

\subsection{顺序源、广义 perambulator 与真实成本}

三点函数的标准顺序源不是“少做一次求逆”的纯技巧，而是把汇点核子、算符插入和源端传播子重组为同一个协变对象。若 \(S=D^{-1}\eta\) 是源传播子，则辅助顺序传播子满足
\begin{equation}
  D\Sigma=e^{-i\mathbf p_f\cdot x}\gamma_5S,
  \label{eq:sequentialdirac}
\end{equation}
其中右手边只在汇点时间片非零。对每个末态动量 \(\mathbf p_f\) 通常需要一个新的顺序源，而改变初态动量 \(\mathbf p_i\) 往往只需改变投影相因子。因此三点谱学的成本近似由
\begin{equation}
  N_{\rm solve}^{3\rm pt}\simeq N_{\rm base}+N_{\mathbf p_f},
  \label{eq:sequentialcost}
\end{equation}
控制，不能把“顺序源减少求逆”误解为对所有动量同时免费。

蒸馏框架把传播子压缩到低模空间。广义 perambulator 写成
\begin{equation}
  J_{\alpha\beta}(t_f,t,t_i)
  =\left[V^\dagger(t_f)M^{-1}(t_f,t)\Gamma(t)
  M^{-1}(t,t_i)V(t_i)\right]_{\alpha\beta},
  \label{eq:generalizedperambulator}
\end{equation}
从而把源侧传播、流插入和汇侧传播编码成一个 \(N_v\times N_v\) 矩阵。其代价并不只是“\(N_v\) 很小”：传播子本身约按 \(N_v\) 增长，介子/重子连接两点函数可分别出现 \(N_v^3\) 和 \(N_v^4\) 的组合，非连通环通常按 \(N_v^3\) 组织。实际报告应给出 \(N_v\)、基算符数、源数、动量数和缓存命中，而不只给总 GPU 时间\cite{Peardon2009,Bali2016}。

\subsection{\texorpdfstring{大 \(\lambda\) 尾部、derivative 方法与 RGR 小 \(x\) 限制}{大 lambda 尾部、derivative 方法与 RGR 小 x 限制}}

坐标空间数据只能覆盖有限 \(z\)。为了做全区间 Fourier 变换，需要对 \(\lambda=zP_z\) 的大距离尾部建模。一个常用的结构是
\begin{equation}
  h_R(z,P_z)
  \sim c_1z^{-d_1}e^{-z/\lambda_0}
  +c_2z^{-d_2}e^{-z/\lambda_0}+\cdots ,
  \label{eq:large-lambda-tail}
\end{equation}
其中指数因子来自非微扰能隙，幂次来自高阶 twist 与端点行为。拟合窗口 \(\lambda>\lambda_L\)、尾部函数形式和 \(z_{\max}\) 的差异都必须进入系统误差。自重整化时通常先对每个格距做大 \(\lambda\) 外推，再做 \(a\to0\) 联合外推；顺序反过来会让大 \(z\) 噪声通过连续极限污染尾部参数。

derivative 方法用局域导数压制尾部截断：
\begin{equation}
  \widetilde q(x,P_z)
  =\frac{i}{x}\int\dd z\,e^{ixP_zz}
  \frac{\dd h_R(z,P_z)}{\dd z},
  \qquad x\neq0.
  \label{eq:derivativemethod}
\end{equation}
它在中等 \(x\) 可以减小截断振铃，但 \(x\to0\) 时出现 \(1/x\) 增强，且无法直接给出 \(x=0\) 的值；因此不能把 derivative 方法当成所有 \(x\) 区间的统一替代。

当 \(\zeta_z=(2xP_z)^2\) 与 \(\mu^2\) 相差很大时，固定阶硬核含大对数。重整化群重求和把硬核从自然标度 \(\mu_0=2xP_z\) 演化到参考标度：
\begin{equation}
  H\!\left(\frac{\zeta_z}{\mu^2}\right)
  =H\!\left(\frac{\zeta_z}{\mu_0^2}\right)
  \exp\left[
  \int_{\mu_0^2}^{\mu^2}\frac{\dd\bar\mu^2}{\bar\mu^2}
  \gamma_H(\alpha_s(\bar\mu))
  \right].
  \label{eq:rgr-hardkernel}
\end{equation}
这条链也解释了小 \(x\) 的根本限制：当
\begin{equation}
  x\lesssim x_{\min}\sim\frac{\Lambda_{\rm QCD}}{2P_z}
  \label{eq:xmin-lamet}
\end{equation}
时，\(\mu_0\) 进入非微扰区，固定阶匹配失去控制。大 \(P_z\) 不是单纯提高信噪比，而是扩大小 \(x\) 的微扰匹配窗口\cite{HeTMD2024,JiReview2021}。

\subsection{误差传播的两条路径与当前实现门控}

格点误差不能只用一个 bootstrap 标准差概括。文档给出了两条互补路径。路径 A 在每个 replica 内重跑完整分析链，自动保留拟合、重正化、匹配和外推之间的相关性；路径 B 先给中心值，再分别变化窗口、标度、拟合形式和尾部模型，最后按平方和组合：
\begin{equation}
  \sigma_{\rm quad}
  =\sqrt{\sigma_{\rm stat}^2+\sum_x\sigma_x^2}.
  \label{eq:quadratureerror}
\end{equation}
路径 B 透明、便宜，但假设各系统分量不相关；对于有限体积与手征外推、\(z_s\) 与匹配标度等明显相关的量，必须改用全协方差或直接扫描联合变化。合格报告应明确哪些误差来自同一 replica 相关样本，哪些来自“硬选择”的变化。

当前 PyQCD 诊断链实现层的关键状态可概括为表~\ref{tab:implementation-gates}。它补充的是“实现到底闭合到哪一步”，不是重复发表结果的证据矩阵。

\begin{table}[htbp]
\centering
\caption{由 PyQCD 文档审计整理出的实现层门控。}
\label{tab:implementation-gates}
\small
\begin{tabularx}{\textwidth}{L{3.2cm}L{4.4cm}Y}
\toprule
门控 & 当前证据 & 判定与后果 \\
\midrule
源分辨两点函数 & 保存 source 轴与 \((source,dt)\) shape；拒绝源平均缓存 & 通过；禁止把源间差异当重复测量 \\
非连通归一化 & 修正 \(corr3=None\) 路径的 \(N_{\rm source}N_{\rm conf}\) 归一化 & 通过；旧 ratio/c0/Sumratio 数值不得复用 \\
Fourier 约定 & quasi-TMD 与局域 spin-2 Ioffe-time 变换分离 & 通过；不能混用 \(1/(2\pi)\) 与 \(1/\pi\) 约定 \\
匹配支撑 & 输出 \(x\) 与卷积 \(y\) 支撑分开，端点区不发布 & 通过；\(x\to0,1\) 仍需额外误差 \\
流窗口 & \(\tau=3a^2\) 时 \(r_f\approx4.90a\)，非零 \(b_T/a\le4\) & 失败；只能作诊断，不能进入物理 \(b_T\) 结果 \\
流窗口扫描 & \(\tau=0.5a^2,1.0a^2\) 改善窗口，但部分动量仍有 fallback & 部分通过；两点不是 \(\tau\to0\) 外推 \\
流步长与单位 & 运行用步长 \(0.05\)，参考口径为 \(0.01\)；两套 fm-to-GeV 常数相差约 \(0.17\%\) & 未闭合；需步长收敛与统一常数 \\
soft/rapidity 与混合 & 无同几何、同伴随表示软因子和完整 \(Z_{ij}\) & 未闭合；禁止发布物理胶子 TMD-PDF \\
\bottomrule
\end{tabularx}
\end{table}

\begin{physicalbox}{文档证据的正确层级}
库内文档可以提供约定、实现路径、已知失败模式和可复现诊断，但不能把“代码能运行”升级为“物理结果已成立”。只有原始论文与独立格点数据支持的主张进入物理结论层；当前 11 组态诊断链即使给出有限的 ratio、c0 或 Sumratio，只要流窗口、soft/rapidity、混合或多参数极限未闭合，仍应保持 \texttt{diagnostic\_only} 和 \texttt{physical\_acceptance=false}。
\end{physicalbox}

\section{关键工作的证据矩阵}

表~\ref{tab:publication-matrix} 按“算符/理论、共线胶子、软结构、完整 TMD”四类整理库内最相关的格点工作。最后两列尤其重要：一个工作验证了方法，不等于它已经得到核子胶子 TMD。把“首次”“连续极限”“第一性原理”等措辞放回具体观测量和参数范围，是综述写作的基本要求。

\begin{footnotesize}
\begin{longtable}{L{2.6cm}L{2.0cm}L{2.7cm}L{2.6cm}L{2.7cm}}
\caption{库内胶子 TMD 相关工作的证据矩阵。}\label{tab:publication-matrix}\\
\toprule
工作 & 方法/对象 & 主要参数或范围 & 已建立 & 未建立/限制 \\
\midrule
\endfirsthead
\toprule
工作 & 方法/对象 & 主要参数或范围 & 已建立 & 未建立/限制 \\
\midrule
\endhead
Ji 2013, 2014\cite{Ji2013,Ji2014} & LaMET/准 PDF & 理论框架 & 空间关联连接光锥物理 & 非胶子专用；无格点数值 \\
Ji--Zhang--Zhao 2018；Izubuchi 等 2018\cite{JiZhangZhao2018,Izubuchi2018} & 重正化与因子化 & 理论 & 准 PDF/准 TMD 的因子化语言 & 胶子 staple 和软函数仍需专门定义 \\
Zhang 等 2019\cite{ZhangGluonLaMET2019} & 胶子准 PDF 算符 & 四个可乘性构建块 & 胶子算符分类、混合与匹配 & 共线 PDF，不是 TMD \\
Li--Ma--Qiu 2019\cite{LiMaQiu2019} & 准胶子算符 & 全阶 UV 结构 & 可乘性重正化的严格证明 & 不处理 TMD rapidity/pinch-pole \\
Balitsky 等 2020\cite{Balitsky2020} & 胶子赝分布 & 短距 ITD/RITD & 不变振幅、短距匹配 & 共线极限与伪分布 \\
Fan 等 2018\cite{FanQuasiGluon2018} & 首次胶子准 PDF & 格点 proof-of-principle & 胶子场强和 Wilson 线可算 & 早期算符/统计限制 \\
Fan--Lin 2021\cite{FanLinPion2021} & 介子胶子 PDF & 两 \(a\)，三 \(m_\pi\)，\(P_z\le2.29\) GeV & pseudo-PDF 与匹配链 & 介子；无 \(b_T\)/TMD \\
Salas-Chavira 等 2022\cite{SalasChavira2022} & K 介子胶子 PDF & 两 \(a\)，\(m_K\sim0.53\) GeV & 味道依赖与 Ioffe-time & 介子；无 TMD \\
Good--Yao--Lin 2026\cite{GoodYaoLin2026} & 核子胶子 PDF & \(a\simeq0.12\) fm，\(m_\pi^{\rm val}=310,690\) MeV，\(P_z\simeq1.71\) GeV & 核子 LaMET proof-of-principle、hybrid-ratio、单圈匹配 & 单格距、较重金属质量、非 TMD \\
Zhang 等 2020\cite{ZhangSoft2020} & TMD 软函数 & CLS \(a=.098\) fm，\(m_\pi^{\rm val}=547\) MeV & 形状因子/准 TMDWF 提取 \(S_I\)、CS 核先例 & 介子态、领头阶匹配、非胶子核子 TMD \\
Zhang 等 2022\cite{Zhang2022TMD} & 准 TMD 重正化 & 五个格距 \(0.032\)--\(0.121\) fm & Wilson 圈减除、RI/MOM 失效诊断 & 准 TMD 重正化，不是最终物理 TMD \\
Chu 等 2022\cite{ChuCS2022} & CS 核 & MILC \(a=.12\) fm，\(m_\pi^{\rm val}=670\) MeV，三动量 & 单圈匹配提取 \(K(b_T)\) & 单格距；胶子表示/核子直接计算未闭合 \\
Chu 等 2023\cite{ChuSoft2023} & \(S_I\) 与 CS 核 & MILC/CLS 交叉 & 多系综、Fierz、两种 TMD 方案对照 & 仍主要是介子辅助量 \\
Chu 等 2023\cite{ChuTMDWF2023} & \(\pi\) TMDWF & 双系综 & TMDWF 首算与模型比较 & 不是核子胶子 TMD \\
He 等 2024\cite{HeTMD2024} & 核子非极化 TMD & 单系综完整链 & 夸克 TMD 的 NNLO/RGR 先例 & 夸克通道，不能替代胶子 \\
Ma 等 2025\cite{MaBM2025} & Boer--Mulders & 核子与介子 & T-odd/自旋关联探索 & 夸克通道，胶子 counterpart 未完成 \\
Egerer 等 2022\cite{Egerer2022} & 胶子螺旋度 & 格点探索 & 胶子自旋结构的先例 & 共线/矩水平，非完整 TMD \\
\bottomrule
\end{longtable}
\end{footnotesize}

\section{共线胶子 PDF 的格点先例}

胶子 TMD 计算首先依赖胶子场强算符、胶子回路、重正化和 Fourier 匹配等基础技术。库内文献已经包含若干共线胶子 PDF 的格点研究：
\begin{itemize}
  \item 介子胶子 PDF 使用 pseudo-PDF 框架，在多个格距和介子质量上构造 reduced Ioffe-time distribution，并通过胶子-夸克匹配核提取 \(x\) 依赖\cite{FanLinPion2021}；
  \item 首个核子胶子 PDF 使用 MILC \(N_f=2+1+1\) HISQ 系综、单格距 \(a\approx0.12\)~fm、较重的价 \(\pi\) 质量、\(P_z\approx1.71\)~GeV 和 hybrid-ratio 重正化，给出了 LaMET 的 proof-of-principle\cite{GoodYaoLin2026}；
  \item 胶子准 PDF 的算符分类、乘性可重正化组合和一阶匹配核已经建立\cite{ZhangGluonLaMET2019,LiMaQiu2019,Balitsky2020}。
\end{itemize}
这些结果证明了胶子场强、Wilson 线和 Fourier 匹配的可行性，但没有提供 \(b_T\) 分辨或 rapidity 演化；因此不能直接改称为胶子 TMD-PDF。

\section{软函数与 Collins--Soper 核}

库内已发表的格点研究给出了 TMD 软函数和 CS 核的重要先例。Zhang 等在 CLS \(a=0.098\)~fm 系综上，通过大动量赝标量介子形状因子和准 TMD 波函数，在 leading-order 因子化下提取了内禀软函数，并研究了 CS 核\cite{ZhangSoft2020}。Chu 等在 MILC \(a\approx0.12\)~fm 系综上使用三个介子动量，在单圈匹配下把 CS 核提取到 \(b_T\lesssim0.6\)~fm，并系统讨论了算符混合、标度依赖和 \(1/P_z^2\) 修正\cite{ChuCS2022}。这些结果表明：
\begin{enumerate}
  \item 同几何 Wilson 圈可以显著消除长 Wilson 线的 UV 结构；
  \item 软函数和 CS 核在小 \(b_T\) 区域可以与微扰 QCD 对接；
  \item 大 \(b_T\) 区域包含不可由普通 PDF 拟合替代的非微扰信息；
  \item 当前结果仍主要基于单格距和特定的介子态，不能自动外推为核子胶子 TMD。
\end{enumerate}

\section{夸克 TMD 与直接胶子 TMD 的差距}

夸克 TMD 的格点计算为胶子通道提供了方法学模板：准 TMD 算符、Wilson 圈平方根、软函数、CS 核、硬匹配和大动量外推可以组合成完整链条。库内报告引用的核子非极化 TMD 与 Boer--Mulders 探索说明该方法在夸克通道已经进入全分布阶段\cite{HeTMD2024,MaBM2025}。然而，胶子通道不能直接照搬以下对象：
\begin{itemize}
  \item 基本表示 Wilson 线不能替代伴随表示 Wilson 线；
  \item 夸克端点发散、Dirac 投影和手征混合与胶子场强算符不同；
  \item 胶子的单态演化需要 \(g-q\) 混合矩阵，而不是单独一个非单态反常量纲；
  \item 胶子回路和真空减除的统计误差结构不同；
  \item 胶子 TMD 的极化分解包含两个场强指标，投影和规范链转角更复杂。
\end{itemize}
因此，把夸克软函数或夸克 CS 核数值直接代入胶子公式，至多得到诊断性估计。

\section{库内实现审计的结论}

本库近期的胶子 TMD 审计提供了一个重要的“失败关闭”案例。真实组态上的源分辨两点函数、Wilson flow 和 TMD OPE 已经能够运行，但在固定 \(\tau=3a^2\) 时，流半径 \(r_f=\sqrt{8\tau}\approx4.90a\)，而库内扫描的非零横向分离主要为 \(b_T/a\le4\)。这违反了 \(r_f<b_T\) 的短流时窗口，因此流化会污染待测的 \(b_T\) 结构。审计还发现：
\begin{itemize}
  \item 单组态或少量组态下，ratio/plateau 统计不可辨识；
  \item \(P_z=0\) 分母的部分点未通过严格质量门；
  \item 缺少同流时、同 staple 几何、同伴随表示的 soft/rapidity；
  \item 缺少完整胶子-夸克单态混合矩阵；
  \item 当前只有单格距、有限流时和有限组态数，不能构成连续极限。
\end{itemize}
因此，即使某些中间矩阵元有限，审计仍将链条标记为 \texttt{diagnostic\_only} 和 \texttt{physical\_acceptance=false}。这不是实现失败，而是物理验收边界被正确执行。它直接支持本文的核心结论：\textbf{可执行代码、有限矩阵元和物理 TMD 是三级不同的证据}。

\section{库内真实组态审计的定量快照}

为避免“实现能运行”被误读为“物理结果已成立”，这里把审计中最关键的状态独立列出。该快照不替代发表结果，只用于说明验收门如何工作。

\begin{table}[htbp]
\centering
\caption{库内胶子 TMD 诊断链的代表性状态。所有数字均来自本地审计记录，不是本文新产生的测量。}
\label{tab:audit-snapshot}
\small
\begin{tabularx}{\textwidth}{L{3.0cm}L{4.2cm}Y}
\toprule
项目 & 状态/数值 & 物理解释 \\
\midrule
系综与体积 & \(\beta=6.20\)，\(24^3\times72\)，11 个组态 & 单格距、单流时的小样本诊断，不是生产系综 \\
源分辨两点函数 & 保存 \((source,dt)\) 轴，shape \((72,72)\) & 避免把源平均量二次平移，保留源间关联 \\
动量标签 & \(P_0,P_{200},P_{400}\) & 覆盖零动量和两个非零动量，但未形成大动量连续扫描 \\
Sumratio 质量 & \(P_0=52/65\)，\(P_{200}=51/65\)，\(P_{400}=65/65\) & P400 中心质量较好，P0/P200 未通过严格覆盖门 \\
replica 统计 & 1000 bootstrap；fallback 非零 & fallback 不能当作中心物理值，必须保留状态 \\
流时间 & \(\tau=3a^2\)，\(r_f=\sqrt{24}\,a\approx4.90a\) & 只测一个流时，无法做 \(\tau\to0\) 外推 \\
横向窗口 & 非零扫描主要 \(b_T/a\le4\) & 多数点不满足 \(r_f<b_T\)，流会污染待测横向结构 \\
soft/rapidity & 无同几何、同伴随表示测量 & TMD 因子化链未闭合 \\
单态混合 & 无完整 \(Z_{ij}\) 矩阵和夸克输入 & 胶子-夸克 singlet mixing 未闭合 \\
链条状态 & \texttt{diagnostic\_only}；物理接受标记为 false & 只能报告诊断量，不应生成“物理 TMD”结论 \\
\bottomrule
\end{tabularx}
\end{table}

这组状态还揭示了一个常见的逻辑错误：把“矩阵元 finite/shape 正确”当作“矩阵元有物理意义”。有限的裸矩阵元可以完全由错误的真空归一化、源平均、流窗口或软因子导致。正确的判据是：数据在多个物理窗口、多个参数点和多个方案下稳定，并且每个中间量都有可复现的来源契约。

\begin{table}[htbp]
\centering
\caption{库内胶子 TMD 相关证据的分层。表中“可用于 TMD”仅表示方法学相关性，不表示已完成直接核子胶子 TMD 计算。}
\label{tab:evidence}
\small
\begin{tabularx}{\textwidth}{L{3.1cm}L{2.7cm}Y}
\toprule
证据类型 & 已建立内容 & 对核子胶子 TMD 的意义与缺口 \\
\midrule
胶子共线 PDF & 场强算符、ratio/hybrid、NLO 匹配、介子与核子先例 & 验证场强、Wilson 线与 Fourier 基础；缺少 \(b_T\) 和 rapidity \\
软函数/CS 核 & 同几何 Wilson 圈、内禀软函数、单圈 CS 提取 & 提供 nonperturbative 软结构先例；需胶子伴随表示和核子态闭合 \\
夸克 TMD & 准 TMD、Wilson 圈、NNLO/RGR、核子应用 & 提供算法模板；不能替代胶子表示和单态混合 \\
库内诊断链 & source-resolved correlators、flow、OPE、bootstrap、fail-closed & 证明实现可运行；当前不满足物理窗口和完整混合要求 \\
直接胶子 TMD & 尚需同方案 soft/rapidity、多格距、多流时、连续极限 & 本文核查的库内语料不足以支持物理数值声称 \\
\bottomrule
\end{tabularx}
\end{table}

\clearpage
\chapter{实现审计、调试与可复现实验}
\label{sec:implementation-audit}

\section{从物理公式到可执行规格}

一个可复现的TMD计算不能只包含公式和最终图。每个物理对象都应映射到明确的输入、输出、单位、形状、索引顺序、规范表示和状态门。可执行规格至少包括
\begin{equation}
  \mathcal S=
  (I,O,\mathcal D,\mathcal U,\mathcal C,\mathcal E),
  \label{eq:executable-spec}
\end{equation}
其中 \(I\) 是输入，\(O\) 是输出，\(\mathcal D\) 是数据布局，\(\mathcal U\) 是单位与约定，\(\mathcal C\) 是不变量与校验，\(\mathcal E\) 是失败状态。缺少任何一项，代码即使能运行也无法证明输出对应目标物理量。

以胶子两点矩阵元为例，输入至少包含规范组态标识、格距、体积、Dirac算符参数、源与汇类型、动量、横向分离、Wilson线方向和流时；输出应包含关联函数、真残差、归一化因子、缓存键和状态。只保存一个复数数组不足以复现结果。

\section{数据布局、索引与形状}

张量索引顺序是常见错误来源。一个推荐约定是
\begin{equation}
  C[\text{source},\text{sink},t,\mathbf p_f,\mathbf p_i,z,b_x,b_y,\text{operator}],
  \label{eq:data-layout}
\end{equation}
并在代码中显式命名每个轴。若代码把source轴平均掉，后续无法区分源间相关；若把 \(z\) 与时间轴混用，Fourier变换会错误地作用于时间片。保存shape只是最低要求，还应保存轴语义和单位。

缓存键必须包含所有影响结果的输入：
\begin{equation}
  K_{\rm cache}
  =\operatorname{Hash}\left(
  U,\;D_{\rm params},\;S,\;O,\;\bT,z,\tau,
  \text{conventions},\text{code version}
  \right).
  \label{eq:cache-key}
\end{equation}
若缓存键遗漏规范链方向或流时，改变物理方案后可能读到旧结果。缓存命中率必须与数值正确性分开报告。

\section{流水线与数据依赖图}

\begin{figure}[htbp]
\centering
\begin{tikzpicture}[
  node distance=0.65cm and 0.55cm,
  box/.style={draw=qcdnavy!70!black,rounded corners=2pt,fill=qcdmist,align=center,minimum width=2.7cm,minimum height=0.85cm},
  gate/.style={draw=qcdrust!75!black,rounded corners=2pt,fill=qcdsand,align=center,minimum width=2.7cm,minimum height=0.85cm},
  arrow/.style={-{Latex[length=2mm]},thick,qcdnavy!75!black}
]
\node[box] (cfg) {规范组态\\与元数据};
\node[box,right=of cfg] (prop) {传播子\\与残差};
\node[box,right=of prop] (corr) {两点/三点\\关联函数};
\node[gate,below=of corr] (ren) {UV与rapidity\\减除};
\node[gate,left=of ren] (match) {匹配与\\Fourier};
\node[gate,left=of match] (fit) {外推与\\误差传播};
\draw[arrow] (cfg) -- (prop);
\draw[arrow] (prop) -- (corr);
\draw[arrow] (corr) -- (ren);
\draw[arrow] (ren) -- (match);
\draw[arrow] (match) -- (fit);
\end{tikzpicture}
\caption{从组态到TMD的依赖图。每个箭头都应有缓存键、校验和失败状态。}
\label{fig:implementation-dag}
\end{figure}

图~\ref{fig:implementation-dag} 中的每一步都应有一个独立状态。若Dirac求解未收敛，后续关联函数不得标记为完成；若软函数缺失，匹配步骤应输出“未闭合”而不是静默继续。状态传播可以写成
\begin{equation}
  s_{i+1}=
  \begin{cases}
    \texttt{complete},&s_i=\texttt{complete}\ \text{且校验通过},\\
    \texttt{diagnostic\_only},&s_i\ \text{仅通过部分物理门},\\
    \texttt{failed},&\text{数值或约定校验失败}.
  \end{cases}
  \label{eq:state-propagation}
\end{equation}

\section{Fail-closed与证据门}

物理流水线应采用fail-closed：缺少关键输入时停止，而不是用默认值填补。典型硬门包括：
\begin{enumerate}
  \item 真残差超过阈值；
  \item source/sink/operator的shape不一致；
  \item 流时窗口为空；
  \item 软函数、Wilson圈或混合矩阵缺失；
  \item 缓存键与当前代码版本不匹配；
  \item 输出包含NaN、Inf或未定义Fourier相位。
\end{enumerate}
Fail-closed会增加失败率，但能防止把中间量误发布为物理结果。报告应区分“科学上未闭合”和“软件运行失败”：前者输出诊断状态，后者输出错误日志和恢复点。

\section{单元测试、回归测试与物理不变量}

单元测试应覆盖公式的最简单极限。例如：
\begin{itemize}
  \item 纯规范链接的plaquette等于1；
  \item \(P_z=0\) 或 \(z=0\) 的矩阵元满足定义归一化；
  \item 正负 \(z\) 的关联函数满足复共轭或厄米关系；
  \item Wilson圈在 \(L_\parallel=0\) 时退化为局域算符；
  \item Fourier变换对常数和纯相位输入有解析结果；
  \item 小规模Dirac矩阵与稠密解一致；
  \item bootstrap在完全相关极限下给出零比值误差。
\end{itemize}
回归测试还应固定随机种子并比较摘要哈希。物理不变量测试比“输出数组shape正确”更能发现约定错误。

\section{日志、哈希与证据链}

每次运行应保存：
\begin{equation}
  \mathcal R=
  \{\text{code hash},\text{input hash},\text{parameters},
  \text{environment},\text{output hash},\text{status}\}.
  \label{eq:run-record}
\end{equation}
日志不记录密钥或无关环境变量。对于大文件，记录SHA-256和生成命令；对于随机分析，记录种子、bootstrap索引和失败replica。这样可以在不复制数GB数据的情况下复核结果来源。

\section{性能模型与资源账本}

总时间可分解为
\begin{equation}
  T_{\rm total}
  =N_{\rm cfg}N_{\rm src}N_{\rm sink}N_{\rm op}
  \left(T_{\rm matvec}N_{\rm iter}+T_{\rm contract}\right)
  +T_{\rm IO},
  \label{eq:performance-model}
\end{equation}
其中 \(T_{\rm matvec}\) 包含全局同步和内存带宽，\(T_{\rm contract}\) 包含张量收缩，\(T_{\rm IO}\) 包含读写缓存。只报告GPU峰值浮点性能会低估实际瓶颈。优化前应先测量matvec、通信、内存和I/O占比。

\section{调试决策树}

\begin{table}[htbp]
\centering
\caption{TMD流水线的调试决策树。}
\label{tab:debug-tree}
\small
\begin{tabularx}{\textwidth}{L{3.0cm}L{4.2cm}Y}
\toprule
症状 & 优先检查 & 下一步 \\
\midrule
矩阵元NaN/Inf & 输入链接、Dirac残差、精度 & 在小子体积复现并检查最小系统 \\
结果依赖源位置 & 归一化、源平均、缓存键 & 分离source轴并重跑单源 \\
\(P_z\to-P_z\)不满足关系 & Fourier相位、Wilson方向、共轭 & 检查正负链接和端点指标 \\
重正化后仍强 \(z\) 依赖 & 同几何减除、混合、短距窗口 & 逐项关闭方案做对照 \\
大 \(z\) 误差爆炸 & 低模、随机源、周期镜像 & 检查信噪比和deflation \\
匹配后强 \(\mu\) 依赖 & 硬核阶数、方案配套、RGE & 用不同 \(\mu\) 检验相消 \\
拟合不稳定 & 协方差、模型参数、窗口 & 用bootstrap和交叉验证 \\
\bottomrule
\end{tabularx}
\end{table}

\section{复现包目录}

一个最小复现包可按以下结构组织：
\begin{verbatim}
run/
  manifest.json
  parameters.tsv
  cache/
  logs/
  outputs/
  plots/
  checksums.sha256
\end{verbatim}
其中manifest记录代码版本、命令和输入哈希；parameters记录物理参数；logs记录状态；checksums绑定输出。论文中的每个数值图都应能追溯到manifest和checksum。

\begin{checkpointbox}{实现审计验收}
\begin{enumerate}
  \item 每个物理对象是否有输入、输出、shape、单位和状态；
  \item 缓存键是否覆盖所有改变物理结果的参数；
  \item 失败是否fail-closed并保留可恢复信息；
  \item 是否有小系统解析解或稠密解对照；
  \item 是否记录代码、输入、输出和随机种子哈希；
  \item 是否把性能账本与物理误差账本分开。
\end{enumerate}
\end{checkpointbox}

\clearpage
\chapter{PyQCD 的代码架构与执行模型}

\section{从物理对象到模块边界}

PyQCD 的历史实现不是围绕单一“大程序”组织，而是把格点场、关联函数、重正化和极限外推拆成可替换的数组变换。其主数据流可以概括为
\begin{equation}
  U_\mu(x)
  \longrightarrow V_\mu(\tau;x)
  \longrightarrow \{F_{\mu\nu},\widetilde F_{\mu\nu},W_{\rm st}\}
  \longrightarrow O(z,b_\perp)
  \longrightarrow c_0
  \longrightarrow h_R
  \longrightarrow \tilde g
  \longrightarrow g,
  \label{eq:b07-dataflow}
\end{equation}
其中每一步改变的是物理对象或统计估计量，而不是只改变文件格式。代码层的模块边界应据此理解：格点和线性代数属于基础层，场强、Wilson 线和提升/投影属于算符层，传播子收缩与源分辨两点函数属于关联函数层，比值、重正化、匹配和 Fourier 变换属于提取层。把 \eqref{eq:b07-dataflow} 中的任一箭头省略，都不能由相邻模块“自动补齐”。

\begin{longtable}{>{\raggedright\arraybackslash}p{2.7cm}>{\raggedright\arraybackslash}p{5.3cm}>{\raggedright\arraybackslash}p{6.0cm}}
\caption{PyQCD 历史包结构与物理职责。}\label{tab:b07-architecture}\\
\toprule
模块 & 主要内容 & 物理解释与边界\\
\midrule
\endfirsthead
\toprule
模块 & 主要内容 & 物理解释与边界\\
\midrule
\endhead
\texttt{lattice} & 常数、格距/单位换算、Dirac 矩阵 & 固定表示与量纲；不替代算符归一化\\
\texttt{tools} & 后端选择、数组适配、HDF5 与旧格式读取 & 使算法与 numpy/cupy/torch 解耦；不保证物理约定一致\\
\texttt{vertex} & momentum phase、VdV/VVV、蒸馏顶点 & 提供几何与动量投影；单段 Wilson 线顶点不等于完整非局域算符\\
\texttt{contraction} & Wick 收缩、重子算符、动态收缩、顺序传播子 & 生成关联函数拓扑；必须保留源、汇、插入时刻和独立动量的轴语义\\
\texttt{operator} & Clover 场强、对偶场强、直线 Wilson 线、胶子 OPE/TMD 基元 & 固定场强符号、Lorentz 指标与路径；不自动给出核子矩阵元\\
\texttt{renorm} & Wilson flow、TMD 提取、$Z_R$、hybrid、NLO 匹配、外推 & 连接矩阵元与光锥量；其参数化与方案必须同时记录\\
\texttt{analysis} & Jackknife/bootstrap、meff、三点比值、不相连拟合 & 将组态样本变为估计量与协方差；不能修复上游归一化错误\\
\texttt{pipeline} & 配置、步骤调度、缓存和产物 & 固化执行顺序与依赖；不等于物理验收\\
\texttt{parallel} & 元任务规划、MPI 调度、设备绑定和释放 & 按组态/步骤并行；不改变单任务数值算法\\
\bottomrule
\end{longtable}

\section{实现链与物理证据链}

历史报告给出了两种容易混淆的“完整”含义。第一种是\emph{实现完整}：入口能够按顺序调用流、算符、比值、重正化、Fourier 和匹配函数，中间数组能够缓存。第二种是\emph{物理完整}：同一 Wilson 线几何下的 UV 重正化、soft/rapidity 减除、色表示、匹配核、多格距/多流时/大动量极限和统计误差全部闭合。前者是后者的必要条件，但远非充分条件。专著应把下表的两类状态分开。

\begin{center}
\small
\begin{tabularx}{\textwidth}{>{\raggedright\arraybackslash}p{3.2cm}X}
\toprule
问题 & 可接受的表述\\
\midrule
代码链是否连通？ & “历史实现提供了从规范场到匹配入口的函数链”，并列出快照版本和未执行步骤。\\
局部测试是否通过？ & “在该版本、该输入和该预算下通过”，附测试计数、输入形状和比较判据。\\
物理 TMD 是否成立？ & 只有在算符、soft/rapidity、重正化、匹配和所有极限共同一致时才成立。\\
\bottomrule
\end{tabularx}
\end{center}

本批次材料中的 U001、U006 和 U024 记录了历史包结构、九步管线、三后端和测试计数，U004 则给出了更细的物理公式链。由于这些报告对应不同日期与提交，正文不应把 \texttt{18/18}、\texttt{237/237} 和 \texttt{106 passed} 合并成一个“总通过数”；应把它们视为不同快照的版本化证据。

\section{张量布局与缓存语义}

历史实现把规范场写成
\begin{equation}
  U \in \mathbb C^{N_t\times N_z\times N_y\times N_x\times 4\times 3\times 3},
  \qquad U_{\mu,ab}(x)\in SU(3).
  \label{eq:b07-layout}
\end{equation}
这里的维序是物理约定的一部分：时间方向在最前，随后是三个空间轴，再是方向、色指标和色指标。任何 \texttt{roll}、转置或 \texttt{einsum} 都必须显式映射到该布局；否则同一个矩阵乘法的结果可能在数组上“成功”但在几何上错位。中间产物优先以 HDF5 保存，并允许旧 npy/npz 回退，是为了保留旧基线的可读性，不是自动证明新旧格式具有相同归一化。

对源分辨关联函数，缓存至少应保留组态标识、源时刻、汇时刻、插入时刻、动量标签、算符分离和规范场/流时标识。只缓存已经对源平均的数组会销毁源间相关性，也会使后续 bootstrap 或 jackknife 无法重建正确的协方差。

\chapter{梯度流重整化胶子 TMD 的计算链}

\section{Wilson flow：定义、RK3 与适用窗口}

连续流方程把规范场沿辅助流时 $t$ 演化：
\begin{equation}
  \partial_t B_\mu(t,x)=D_\nu G_{\nu\mu}(t,x),
  \qquad B_\mu(0,x)=A_\mu(x),
  \qquad G_{\nu\mu}=\partial_\nu B_\mu-\partial_\mu B_\nu+[B_\nu,B_\mu].
  \label{eq:b07-flow}
\end{equation}
在格点上，链接变量以无迹反厄米生成元演化。历史实现使用
\begin{equation}
  V(t+\varepsilon)=\exp\!\left[\varepsilon Z(V)\right]V(t),
  \qquad
  Z(V)=P_{\rm ah}\!\left[\Omega V^\dagger\right],
  \label{eq:b07-flowstep}
\end{equation}
其中
\begin{equation}
  P_{\rm ah}[M]=\frac{M-M^\dagger}{2}
  -\frac{\operatorname{Tr}(M-M^\dagger)}{2N_c}\mathbf 1,
  \label{eq:b07-ah}
\end{equation}
$\Omega$ 由六条与链接相邻的 staple 贡献组成。投影保证数值演化停留在 $SU(3)$；它不改变流方程在大流时可能造成的非微扰平滑。

历史实现采用 Lüscher 三阶 Runge--Kutta 结构：
\begin{align}
  W_1&=\exp\!\left[\frac{\varepsilon}{4}Z(W_0)\right]W_0,\nonumber\\
  W_2&=\exp\!\left[\frac{8\varepsilon}{9}Z(W_1)-\frac{17\varepsilon}{36}Z(W_0)\right]W_1,\nonumber\\
  V_{t+\varepsilon}&=\exp\!\left[\frac{3\varepsilon}{4}Z(W_2)
  -\frac{8\varepsilon}{9}Z(W_1)+\frac{17\varepsilon}{36}Z(W_0)\right]W_2.
  \label{eq:b07-rk3}
\end{align}
该格式的局部截断误差为 $\mathcal O(\varepsilon^4)$，单步演化的整体误差阶为 $\mathcal O(\varepsilon^3)$。历史报告对生产步长有 $0.05$ 与 $0.01$ 两种记载，故专著必须把“参考步长、生产步长、收敛检查所用步长”分别列出，不能用一个数字概括所有提交。

流半径定义为
\begin{equation}
  r_f=\sqrt{8t}.
  \label{eq:b07-flowradius}
\end{equation}
若目标是保留分离尺度 $z$ 或 $b_\perp$ 的短距结构，至少需要
\begin{equation}
  a\ll r_f\ll \min(z,b_\perp),
  \qquad P_z r_f\ll 1,
  \label{eq:b07-flowwindow}
\end{equation}
并扫描多个 $t$ 以检查 $t\to0$ 或短流时匹配。固定一个物理等价于格距相关值的流时，即使矩阵元有限，也不能区分短流时物理、格距伪迹和过度平滑。对于 $t=3a^2$，有 $r_f=\sqrt{24}\,a\approx4.90a$；当非零 $b_\perp$ 只覆盖到约 $4a$ 时，\eqref{eq:b07-flowwindow} 的右侧不等式并不成立。因此这类数据只能作为诊断，不能作为物理 $b_\perp$ 结构的最终证据。

\section{Clover 场强、Wilson 线与胶子算符}

Clover 场强取四个 plaquette 的对称反厄米组合：
\begin{equation}
  F_{\mu\nu}(x)=-\frac{i}{8}\sum_{k=1}^{4}\left(P_k-P_k^\dagger\right),
  \label{eq:b07-clover}
\end{equation}
四个方位为 $(\mu,\nu)$、$(-\mu,\nu)$、$(-\mu,-\nu)$、$(\mu,-\nu)$。对偶场强为
\begin{equation}
  \widetilde F_{\mu\nu}=\frac12\epsilon_{\mu\nu\rho\sigma}F^{\rho\sigma}.
  \label{eq:b07-dual}
\end{equation}
式 \eqref{eq:b07-clover} 是历史代码的裸 Clover 组合；是否额外出现 $1/(a^2g_0)$、色迹和规范耦合，取决于具体场强定义。专著必须从函数的输入输出和同一作用量约定出发检查，不能把两种归一化拼接。

胶子 TMD 算符需要同时保留纵向分离 $z$、横向分离 $b_\perp$ 和规范链几何。令 $W_{\rm st}$ 为连接两个场强端点的 staple Wilson 线，定义未求和色迹
\begin{equation}
  M^{\mu\lambda;\nu\rho}(z,b_\perp)
  =\operatorname{Tr}\!\left[
  F^{\mu\lambda}(x+z\hat z+b_\perp\hat b)\,
  W_{\rm st}\,
  F^{\nu\rho}(x)W_{\rm st}^\dagger
  \right],
  \label{eq:b07-M}
\end{equation}
历史实现采用组合
\begin{equation}
  O(z,b_\perp)
  =M^{tx;tx}(z,b_\perp)+M^{ty;ty}(z,b_\perp)
  -2M^{xy;xy}(z,b_\perp).
  \label{eq:b07-O}
\end{equation}
报告和代码注释称该组合在一圈水平具有乘性重正化性质，并在 $b_\perp\to0$ 时连接到胶子准 PDF 算符。这里应区分三个层次：组合的代数定义已经写定；一圈反常量纲相消是文献/代码声明；完整 nonperturbative 重正化、混合矩阵和 soft/rapidity 结构仍需独立验证。不能因为 \eqref{eq:b07-O} 是实数或能运行，就直接把它等同于物理 TMD。

\section{不相连三点函数与裸矩阵元}

胶子插入与核子价夸克线不连通时，历史链采用因子化近似
\begin{equation}
  C_3(dt,d\tau,z,b)=C_2(dt)\,\operatorname{OPE}(d\tau,z,b),
  \label{eq:b07-factorization}
\end{equation}
其中 $dt$ 是源汇时间分离、$d\tau$ 是插入相对源的时间。真空贡献需要显式扣除：
\begin{equation}
  C_3^{\rm disc}=C_3-C_2\langle \operatorname{OPE}\rangle,
  \qquad
  R(dt,d\tau,z,b)
  =\left\langle\frac{C_3^{\rm disc}}{C_2}\right\rangle_{t_i}.
  \label{eq:b07-vacuum}
\end{equation}
这一步不是装饰性的常数减除：$C_2$ 的归一化、源平移、组态对齐和插入时刻都必须使用同一套约定。

为了消除激发态污染，历史实现把比值拟合成
\begin{equation}
  R(dt,d\tau,z,b)=c_0(z,b)
  +c_1(z,b)e^{-\Delta E\,d\tau}
  +c_1(z,b)e^{-\Delta E(dt-d\tau)},
  \label{eq:b07-ratiofit}
\end{equation}
在 $dt$ 的中间窗口拟合，并要求 $d\tau$ 和 $dt-d\tau$ 同时远离边界。典型历史参数是 $dt\in[7,10]$、边界 cut 为 6、协方差奇异时使用 \texttt{svdcut=1e-6} 或 plateau 均值备份。式 \eqref{eq:b07-ratiofit} 的渐近极限给出 $R\to c_0$，但有限统计下还必须报告拟合窗口、协方差条件数和模型变体。

\section{自重整化与 hybrid 拼接}

在短距窗口，历史报告使用 NLO $\overline{\rm MS}$ 因子
\begin{equation}
  Z_{\overline{\rm MS}}(z,\mu)
  =1+\frac{\alpha_s C_A}{4\pi}
  \left[
  \frac53\ln\!\left(
  \frac{z^2\mu^2}{4e^{-2\gamma_E}}\right)+3
  \right].
  \label{eq:b07-zms}
\end{equation}
长距部分用零动量比值数据 $h_B(z,P_z=0)$ 参数化。历史模型在短距写
\begin{equation}
  \ln h_B(z)=\ln Z_{\overline{\rm MS}}(z,\mu)
  +(m_0+m_2a^2)z,
  \qquad z\le z_1,
  \label{eq:b07-hb}
\end{equation}
而在 $z>z_1$ 用一组拟合函数或节点描述。$z_1$ 是模型窗口边界，不是由格点自动确定的物理相变点。

历史实现的自重整化因子写成
\begin{align}
  \ln Z_R
  &=k\frac{z}{a\ln(a\Lambda)}
  +\frac12\ln\!\left(1+\frac{d}{\ln(a\Lambda)}\right)^2
  +\frac{5C_A}{3b_0}
  \ln\!\left[
  \frac{\ln(1/(a\Lambda))}{\ln(\mu/\Lambda)}
  \right]\nonumber\\
  &\quad +(m_0+m_2a^2)z+f_1a+f_2a^2,
  \label{eq:b07-zr}
\end{align}
其中 $b_0=11-2N_f/3$，$k,d,m_0,m_2,\Lambda$ 与离散项是全局拟合参数。该式是 PyQCD 采用的实现模型，不是所有自重整化方案的普适定义；尤其不能把线性指数项、对数项和格距修正分别解释成互相独立的物理观测量。

Hybrid 方案把短距比值和长距 $Z_R$ 拼接：
\begin{equation}
  h_R(z,P_z)=
  \begin{cases}
  \dfrac{h_B(z,P_z)}{h_B(z,0)},&z<z_s,\\[7pt]
  \dfrac{h_B(z,P_z)}{Z_R(z)}\eta_s,
  \quad \eta_s=\dfrac{Z_R(z_s)}{h_B(z_s,0)},&z\ge z_s.
  \end{cases}
  \label{eq:b07-hybrid}
\end{equation}
按定义，两支在 $z_s$ 处相等；这只证明拼接连续，不证明结果对 $z_s$ 不敏感。合格的系统误差扫描应改变 $z_s$、$Z_R$ 参数化和短距拟合窗口。

\section{Fourier 重构、Collins--Soper 核与 NLO 匹配}

历史 TMD 通道采用余弦重构
\begin{equation}
  x\tilde g(x,b_\perp,P_z)
  =\frac{1}{\mathcal N}\frac{1}{2\pi P_z}
  \int_0^{z_{\max}}\mathrm dz\,
  \cos(xP_zz)\,h_R(z,b_\perp,P_z),
  \qquad
  \mathcal N=\left(\frac{P_z}{P_t}\right)^2,
  \label{eq:b07-cos}
\end{equation}
其中 $P_t=\sqrt{M_N^2+P_z^2}$。共线反对称通道在旧代码中使用
\begin{equation}
  \tilde g(x,P_z)=\frac{2P_z}{x}
  \int_0^{z_{\max}}\mathrm dz\,
  h(z,P_z)\sin(xP_zz),
  \label{eq:b07-sin}
\end{equation}
并在 $x\to0$ 采用显式保护。式 \eqref{eq:b07-cos} 与 \eqref{eq:b07-sin} 服务不同的表示与对称性；它们的归一化不能互相替换，也不能只比较离散数组的最大值。

有限 $z_{\max}$ 造成 Fourier 截断。历史模型在数据区外用
\begin{equation}
  h_R(\lambda)\sim l_1\lambda^{-a_1}e^{-\lambda/\lambda_0},
  \qquad \lambda=zP_z,
  \label{eq:b07-tail}
\end{equation}
估计大 $\lambda$ 尾部。该 ansatz 只有在尾部拟合区间稳定、$\lambda_0$ 可辨识且对 $z_{\max}$ 不敏感时才可接受。

Collins--Soper 核可由两个动量的比值近似提取：
\begin{equation}
  K(b_\perp)=
  \frac{\ln[h_R(z_0,b_\perp,P_{z1})/h_R(z_0,b_\perp,P_{z2})]}
  {\ln(P_{z1}/P_{z2})}.
  \label{eq:b07-cs}
\end{equation}
这要求两个矩阵元已经使用同一套 UV 重正化、同一 staple 几何和兼容的软函数方案。$K$ 的强子无关性不是自动成立；有限动量、$b_\perp$ 截断和匹配阶数都会污染斜率。

把准分布映射到光锥 TMD 的硬匹配可概括为
\begin{equation}
  x\tilde g(x,b_\perp;\mu,\zeta)
  =\int\frac{\mathrm dy}{|y|}\,
  C\!\left(\frac xy,\lambda_s,\frac{\mu}{yP_z}\right)
  \frac{y\,g(y,b_\perp)}{\sqrt{S_I(b_\perp,\mu)}}
  \exp\!\left[
  \frac12\ln\!\left(\frac{\zeta_z}{\zeta}\right)K(b_\perp)
  \right],
  \label{eq:b07-tmdmatching}
\end{equation}
其中 $C$ 在历史实现中按 $x/y$ 分区，包含单圈胶子核与截断项。单态通道一般还需要夸克-胶子味混合，故式 \eqref{eq:b07-tmdmatching} 不能在没有说明的情况下当作完整单态公式。

离散实现把卷积写成矩阵反演或矩阵乘：
\begin{equation}
  h_R^{\rm PDF}=Z^{-1}h_R,
  \label{eq:b07-zmatrix}
\end{equation}
这里的 $Z$ 是硬匹配卷积矩阵，而不是 UV 重正化矩阵。主值奇点、端点支撑和 $Z^{-1}$ 的条件数都要报告；如果离散网格改变后系数不收敛，就不能把矩阵反演结果解释成物理 PDF。

\section{连续极限与尺度层级}

历史实现采用逐 $x$ 的联合参数化
\begin{align}
  f(x;a,P_z,m_\pi,L)
  &=xg_0(x)+f_xa^2+l_xa^4+h_xa^2P_z^2
  +d_xP_z^{-2}+b_xP_z^{-1}\nonumber\\
  &\quad +k_x(m_\pi^2-m_{\pi,\mathrm{phys}}^2)
  +c_xe^{-La m_\pi}.
  \label{eq:b07-extrap}
\end{align}
各项分别代表格距修正、动量-格距交叉项、大动量幂修正、手征修正和有限体积修正。式 \eqref{eq:b07-extrap} 是拟合 ansatz，不是由单系综数据自动证明的展开定理。若只有一个格距或一个 $P_z$，相应系数不可辨识；此时应删去该项并把目标降为诊断，而不是用先验把系数固定后宣称完成连续极限。

整个物理链至少应满足
\begin{equation}
  \Lambda_{\rm QCD},M_N\ll P_z\ll a^{-1},
  \qquad
  a\ll r_f\ll \min(z,b_\perp),
  \qquad
  |x|\gtrsim \frac{\Lambda_{\rm QCD}}{2P_z},
  \label{eq:b07-windows}
\end{equation}
其中最后一项是固定阶匹配控制小 $x$ 的粗略条件。$b_\perp$ 还必须远离体积的一半，以免周期图像和强子波函数绕回。若窗口交集为空，任何“平台”都只是中间量。

\begin{longtable}{>{\raggedright\arraybackslash}p{3.1cm}>{\raggedright\arraybackslash}p{3.5cm}>{\raggedright\arraybackslash}p{7.3cm}}
\caption{梯度流胶子 TMD 链的主要失败模式。}\label{tab:b07-failures}\\
\toprule
环节 & 失败信号 & 最小诊断\\
\midrule
\endfirsthead
\toprule
环节 & 失败信号 & 最小诊断\\
\midrule
\endhead
场强归一化 & 同态局域极限不匹配 & 固定规范链，比较 Clover 定义、色迹和 $a,g_0$ 因子\\
Wilson 线 & 随路径长度增长且无平台 & 同几何 Wilson 圈、多个 $L_\parallel$、同表示\\
flow 窗口 & $r_f$ 与 $b_\perp,z$ 同量级 & 多流时扫描并检查 $t\to0$\\
比值 & 对 $d\tau$ 或 $dt$ 无平台 & 双端点谱学拟合、源分辨缓存、协方差条件数\\
$Z_R$/hybrid & 对 $z_s$、拟合窗和初值敏感 & 多模型、多窗口、完整 bootstrap 重拟合\\
Fourier & 对 $z_{\max}$ 和尾部模型敏感 & 多个截断、尾部模型和可信 $x$ 区间\\
匹配 & 随网格、$\mu$ 或阶数不收敛 & 支撑域检查、矩阵条件数、NLO 变体\\
soft/rapidity & 与几何无关的软因子缺失 & 同几何 Wilson 圈与表示匹配\\
\bottomrule
\end{longtable}

\chapter{验证、并行工程与证据边界}

\section{质子两点函数与色散关系}

两点函数回归是检验比例归一化和谱学拟合的最小闭环。历史复现使用
\begin{equation}
  C(dt)=\frac{1}{N_t}\sum_{ti}
  \left|C^{\rm slice}(ti,ti+dt)\right|,
  \qquad
  aE_{\rm eff}(t)=\ln\frac{C(t)}{C(t+1)}.
  \label{eq:b07-corr2}
\end{equation}
绝对值投影适用于报告当时检查的通道，但会丢失复相位；若目标通道本身含相位或需要相位敏感干涉，必须先证明该投影不会改变物理结论。

平台拟合采用基态加激发态模型
\begin{equation}
  C(t)=c_0e^{-E_0t}\left[1+c_1e^{-\Delta E\,t}\right],
  \label{eq:b07-c2fit}
\end{equation}
历史窗口为 $t=3,\ldots,7$。在 $P_z=6$ 的样本中，历史报告给出的平均结果为
\begin{equation}
  E_0=1.640(5)a^{-1}=3.069(9)\ {\rm GeV},
  \label{eq:b07-e0}
\end{equation}
并与 $\sqrt{m_p^2+p^2}\approx3.14\ {\rm GeV}$ 的色散预期相差约 $2\%$。这组数字只属于对应报告的数据与拟合窗，不是本批次的新测量。

\begin{center}
\small
\begin{tabularx}{\textwidth}{lXXXX}
\toprule
方向 & $E_0/a^{-1}$ & $E_0/{\rm GeV}$ & $\chi^2/{\rm dof}$ & 条件数\\
\midrule
$x$ & 1.649(8) & 3.084(15) & 0.27 & $6.3\times10^7$\\
$y$ & 1.677(17) & 3.138(32) & 1.8 & $1.3\times10^8$\\
$z$ & 1.637(6) & 3.063(11) & 1.7 & $6.3\times10^3$\\
平均 & 1.640(5) & 3.069(9) & 2.4 & $1.4\times10^7$\\
\bottomrule
\end{tabularx}
\end{center}

该表的解释边界有三点。第一，三方向一致性与色散关系是必要检查，不是独立测量质子质量。第二，把平台窗改成 $[4,8]$ 时报告的平均值移动约 $3\%$，应计入系统误差。第三，$x/y$ 协方差的条件数达到 $10^7$--$10^8$，即使拟合收敛也必须检查窗口和正则化对中心值的影响。

\section{后端适配、持久化与 MPI 元任务}

历史 torch 适配层解决的是数组语义差异，而不是物理约定差异。复数输入按全局精度转换，numpy 风格函数经签名适配，二元运算补丁避免 \texttt{complex128} 与 \texttt{complex64} 混算破坏精度切换。其正确性判据应是同一输入、同一全局精度和同一后端能力下的最大差，而不是只看 dtype 或运行成功。

HDF5 工具每次调用独立打开文件，使多 rank/线程不必共享同一个句柄。元任务按 \texttt{(step, conf)} 划分，静态 round-robin 分配，每个任务写独立目录；rank $r$ 绑定设备
\begin{equation}
  \mathrm{device}(r)=\mathrm{cuda}:(r\bmod n),
  \label{eq:b07-gpumap}
\end{equation}
每任务结束后释放缓存。该方案没有 4D 网格 halo 交换，因而适合以组态为天然并行单位的蒸馏管线；代价是单个超大任务不能通过域分解进一步加速。

历史资源规划写成
\begin{equation}
  N a=n b,
  \label{eq:b07-resource}
\end{equation}
其中 $N$ 为总进程数，$a$ 为单任务显存，$n$ 为 GPU 数，$b$ 为单卡可预算显存。实际 $N$ 还要受 CPU 线程、可用内存、批次数和任务数封顶。资源不足时降至 $N=1$ 是正确行为，不应被理解为“并行失败”；真正需要验证的是任务映射是否无重复、无遗漏，以及串并行产物是否一致。

\section{版本化验证矩阵}

历史报告记录了多组不同层次的验证。它们不能横向求总和，应按快照、输入和验收目标分别登记。

\begin{longtable}{>{\raggedright\arraybackslash}p{2.4cm}>{\raggedright\arraybackslash}p{2.9cm}>{\raggedright\arraybackslash}p{4.1cm}>{\raggedright\arraybackslash}p{4.0cm}}
\caption{B07 历史验证证据与解释边界。}\label{tab:b07-evidence}\\
\toprule
时间/对象 & 证据 & 已证明 & 未证明\\
\midrule
\endfirsthead
\toprule
时间/对象 & 证据 & 已证明 & 未证明\\
\midrule
\endhead
2026-08-16 torch/MPI & 18/18 测试；12 个合成任务串并一致；后端差 $\sim10^{-15}$ & 适配层、合成调度和局部数值一致性 & 真实数据全步骤 MPI；生产规模资源行为\\
2026-08-18 全库报告 & A--E 零差异；237/237 PASS & 特定基线的可复算接口 & 当前提交、完整 soft/rapidity 和物理 TMD\\
2026-08-16 质子能量 & corr2 与参考最大差 0；三方向 $E_0$ 表 & 历史复现链和局部拟合 & 其他系综、其他动量、最终物理质量\\
2026-09-28 GPD 审计 & 106 tests PASS & 通用基元、前向 PDF/TMD 与几何测试 & GPD 专项链；测试数不能外推为 GPD 通过\\
固定流时诊断 & $r_f\approx4.90a$，$b_\perp/a\le4$ & 流窗口冲突可被识别 & 物理 $b_\perp$ 结果和连续极限\\
\bottomrule
\end{longtable}

\section{物理结果的 fail-closed 条件}

专著中的结果分级可压缩为四个门槛。第一，算符与代数门：场强、对偶场强、Wilson 路径、Lorentz 投影、Fourier 归一化和 $z=0$ 归一均可执行。第二，矩阵元门：源分辨两点函数、真空扣除、三点比值、谱学拟合和统计协方差通过。第三，方案门：同几何 soft/rapidity、UV 重正化、CS 核、NLO 匹配和单态混合处于同一方案。第四，极限门：多个 $a$、多个 $P_z$、多个流时和足够体积满足共同窗口。

只有四门同时通过，才能使用“核子胶子 TMD-PDF 的格点结果”这一表述。若只有算符和局部测试通过，应称为“算符原型”或“矩阵元诊断”；若矩阵元也通过但方案未闭合，应称为“候选实现”。固定 $t=3a^2$、有限 $b_\perp$ 且无伴随表示 soft/rapidity 的数据，即使代数结果有限，也不能跨过第三和第四门。


\clearpage
\chapter{历史工作流、审计与验证演进}
\label{sec:b08-history}

本节不把不同年代、不同输入和不同验收层级的报告压成一条“成功链”。相反，我们
按证据来源重建五种对象：蒸馏基线、分析脚本实战、梯度流 TMD 原型、逐层复现矩阵
以及最终失败关闭判定。贯穿全节的问题是：一项检查究竟证明了公式、文件格式、输出
投影、局部矩阵元，还是已经足以支持物理分布。只有最后一类需要同时满足算符、方案、
尺度窗口、统计和极限条件。

\section{证据对象与主张边界}

表~\ref{tab:b08-evidence-objects} 概括本批次材料的证据类型。表中的“历史结论”
表示当时报告中的解释，不表示本专著接受其物理有效性；尤其在源分辨归一化修复后，
早期 ratio、\(c_0\) 与 Sumratio 数值已经不能继续引用。

\begin{table}[htbp]
\centering
\caption{B08 历史材料的证据对象、允许提取的内容与不可升级的结论。}
\label{tab:b08-evidence-objects}
\small
\begin{tabularx}{\textwidth}{>{\raggedright\arraybackslash}p{2.5cm}>{\raggedright\arraybackslash}p{4.0cm}>{\raggedright\arraybackslash}X}
\toprule
材料 & 主要证据对象 & 主张边界 \\
\midrule
stab0 & 3/3/10 组态、\(N_{\rm ev}=50/100\)、VdV/VVV、连通与不相连比值 &
可用于说明管线演进和统计失败；不能把 10 组态的不相连拟合当作矩阵元测量 \\
stab1 & 真实数据上的 ratio、能量、FH、三方向分析及产物清单 &
45/45 是流程断言；没有闭合流时、混合、soft/rapidity 和连续极限 \\
dev5 系列 & Wilson flow、staple 算符、ratio、自重整化、Fourier、NLO/CS &
公式链与历史输出可审计；2026-10-22 审计后，旧的归一化数值必须撤回 \\
4150 矩阵 & 低层公式、HYP OPE、费米子收缩、momsmear 输出投影、effmass &
每一项只证明其明确列出的比较层；3pt/ratio/barematrix 仍无配对证据 \\
dev20 审计 & 源分辨真空减除、归一化反例、流窗口、严格质量门 &
当前只能支持诊断链和 fail-closed 状态，不能发布物理胶子 TMD-PDF \\
\bottomrule
\end{tabularx}
\end{table}

\section{stab0：蒸馏管线从可用计算到可检验统计}

\paragraph{蒸馏定义。}
蒸馏把三维协变 Laplace 算符的低本征模 \(v^{(k)}\) 作为夸克源与汇的压缩基。
传播子在低模空间中的矩阵元称为 perambulator，而介子和重子空间结构分别由
\begin{equation}
  V^{\rm meson}_{mn}(\mathbf p;t)
  =\sum_{\mathbf x}e^{-i\mathbf p\cdot\mathbf x}
  v_m^\dagger(\mathbf x;t)v_n(\mathbf x;t),
  \qquad
  V^{\rm baryon}_{mnl}(\mathbf p;t)
  =\sum_{\mathbf x}e^{-i\mathbf p\cdot\mathbf x}\,
  \varepsilon_{abc}v_m^a v_n^b v_l^c
  \label{eq:b08-vertices}
\end{equation}
给出。重子 VVV 的显式尺度为 \(N_{\rm ev}^3\)，这解释了为什么早期版本首先在
\(N_{\rm ev}\)、组态数和分析闭环之间做取舍。

\paragraph{谱学和不相连比值。}
源时间平均后的两点函数用有效质量平台提取能量：
\begin{equation}
  m_{\rm eff}^{\log}(t)=\frac1a\ln\frac{C(t)}{C(t+1)},
  \qquad
  m_{\rm eff}^{\cosh}(t)=\frac1a\operatorname{arccosh}
  \left[\frac{C(t-1)+C(t+1)}{2C(t)}\right].
  \label{eq:b08-meff}
\end{equation}
对不相连胶子贡献，历史报告采用
\begin{equation}
  C_3(t,\tau,z)=C_2(t)O(z,\tau),
  \qquad
  R(t,\tau,z)=\frac{C_3-C_2\langle O\rangle}{C_2},
  \label{eq:b08-disconnected-ratio}
\end{equation}
再逐 \(z\) 拟合两态模型
\begin{equation}
  R(\tau)=c_0+c_1e^{-dE\,\tau}+c_1e^{-dE\,(t-\tau)}.
  \label{eq:b08-two-state}
\end{equation}
式~\eqref{eq:b08-disconnected-ratio} 的关键风险不在代数形式，而在
\(C_2\)、\(O\) 与真空项如何按 source 和 configuration 归一化；这一点在
第~\ref{sec:b08-dev20} 小节会被后续审计重新定义。

\begin{table}[htbp]
\centering
\caption{stab0 三个历史版本的输入与功能面。所有数值均来自历史报告，不表示当前验收状态。}
\label{tab:b08-stab0-versions}
\small
\begin{tabular}{>{\raggedright\arraybackslash}p{2.2cm}ccc>{\raggedright\arraybackslash}p{7.0cm}}
\toprule
版本 & \(N_{\rm conf}\) & \(N_{\rm ev}\) & 胶子 OPE & 当时新增/不足 \\
\midrule
v20260803 & 3 & 50 & 否 &
建立集中配置、低模库和 LaTeX 报告；高动量统计不足 \\
v20260804 & 3 & 100 & 否 &
增加功能面，但质子无平台，pion \(P_2\) 与 \(P_0\) 退化 \\
v20260805 & 10 & 100 & 是 &
增加 \(pn\)、3pt/4pt 和自动报告；不相连拟合仍完全不受约束 \\
\bottomrule
\end{tabular}
\end{table}

历史报告把 v20260805 称为三者中唯一具有“计算--统计--报告”闭环的版本，但
这不等于完成了物理测量。10 组态下不相连参数出现
\(c_1\sim-9.8\times10^8\)、\(dE\sim8.9\times10^6\) 一类失控尺度，说明
中心值有限并不能证明参数可辨识。更严格的做法是同时报告协方差秩、通过质量门的
点数和失败 replica 比例，而不是只给出拟合曲线。

\paragraph{尺度元数据的版本分支。}
stab0 记录的格距为 \(a\simeq0.1053\,\mathrm{fm}\)，而 dev5 历史报告使用
\(a=0.0938\,\mathrm{fm}\)。本批次没有找到把两者统一到同一标定分支的原始元数据。
因此，后续章节必须把它们当作两个历史来源分支；质量、动量和流半径的数值不能
跨分支自动比较。类似地，stab0 中质量差异同时伴随 \(N_{\rm conf}\)、\(N_{\rm ev}\)
和拟合窗变化，不能把单一差异归因于某一个算法参数。

\section{stab1：流程断言与物理断言的分离}

stab1 把早期基线接上多个分析模块：02 比值、03 比值画图、04 两点能量、05
三方向相关、06 FH 变换。它记录了 10 组态、166 个文件、106 张图和
\(15\,\mathrm{MB}\) 产物，并在流程层执行 45/45 项检查。流程检查覆盖产物存在性、
有效质量偏差、拟合有限性和日志完整性，适合作为回归基线，但不检验 TMD 因子化
或软/快度结构。

其能量与聚合检查可写成
\begin{equation}
  E_0(\mathbf P)=\operatorname*{arg\,min}_{c_i,E_i}
  \chi^2\!\left[C_2;\sum_i c_i e^{-E_i t}\right],
  \qquad
  \mathrm{FH}(t)=\sum_\tau R(t+1,\tau)-\sum_\tau R(t,\tau),
  \label{eq:b08-energy-fh}
\end{equation}
其中前一项检验基态平台，后一项把 ratio 的常数部分与非零 \(dE\) 项的差异投影
出来。历史快照报告 \(m_P=1.097\,\mathrm{GeV}\)、
\(E(P_2)=1.533\,\mathrm{GeV}\)、\(E_0(P_2)=1.525\,\mathrm{GeV}\)，
但三方向/六方向平均因数据可得性退化为单一 \(z\)，10 组态下 \(c_0\)
仍有约 \(0.1\) 的窗口波动。\texttt{rank 6/33} 也说明协方差正则化参与了结果，而不是
所有方向都独立可辨识。

报告表中把 \(c_0(z=4)=0.36\) 列在 GeV 标题下。裸矩阵元在相邻材料中通常以
无量纲或代码单位记录，当前材料无法确认其量纲；专著只能报告数值与单位冲突，
不能擅自补上物理单位。

\section{dev5：一条可执行链路如何被重新验收}
\label{sec:b08-dev5}

\paragraph{流与场强。}
Wilson flow 通过热核型演化压低紫外模式：
\begin{equation}
  \partial_t B_\mu(t,x)=D_\nu G_{\nu\mu}(t,x),
  \qquad B_\mu(0,x)=A_\mu(x),
  \qquad r_f=\sqrt{8t}.
  \label{eq:b08-flow}
\end{equation}
历史实现采用六-staple 的 Wilson 作用量梯度、\(\mathfrak{su}(3)\) 反厄米投影
和截断指数 RK3；Clover 场强由四个 plaquette 的对称平均构造到
\(\mathcal O(a^2)\)。流半径不是“步数”，而是与 \(a,z,b_T\) 直接比较的物理长度。
因此，\(t=3a^2\) 带来的 \(r_f=\sqrt{24}\,a\simeq4.899a\) 必须与所有非零
\(b/a\le4\) 逐点比较，而不能只用平均值或最大 \(b\) 概括。

\paragraph{算符与比值。}
dev5 用 \(F W F W^\dagger\) 的色迹表示非局域胶子关联，并定义牵引表示组合
\begin{equation}
  O(z,b_T)=M^{tx;tx}(z,b_T)+M^{ty;ty}(z,b_T)-2M^{xy;xy}(z,b_T).
  \label{eq:b08-operator}
\end{equation}
其物理意图是消去不需要的局域混合并保留与胶子分布相关的张量结构；但“单圈可乘”
不是“任意 \(b_T\) 和任意几何均已重正化”。两态比值、傅里叶反演和 NLO/CS
匹配仍需独立的源、几何和流时一致性。

\begin{equation}
  \begin{aligned}
    C_3&=C_2O,\\
    R&=c_0+c_1e^{-dE\,\tau}+c_1e^{-dE\,(t-\tau)},\\
    h_R(z,b_T)&=\frac{h_B(z,b_T)}{h_B(0,b_T)},\\
    x\tilde g(x,b_T,P_z)&=\frac{1}{\mathcal N}
      \int\frac{dz}{2\pi P_z}e^{-ixzP_z}h_R(z,b_T,P_z).
  \end{aligned}
  \label{eq:b08-dev5-chain}
\end{equation}
这组公式是历史报告的主要可审计内容。其物理解释应停留在“给出准分布和中间
矩阵元的计算路径”；在 soft/rapidity、完整混合和多尺度极限缺失时，
\(xg(x,b_T)\) 与 \(K(b_T)\) 不能直接称为物理结果。

\paragraph{三次快照的差异。}
U009 是基线；U010 增加流协变、对偶张量、staple 规范和两态展开等推导框；
U011 再增加因果总览和图表的三段式解释。三者没有独立的新系综、新统计或新
重正化数据，故不能把页数增长计为证据增长。U011 中关于 \(N(1440)\)、CT18、
胶子半径和 \(b^*\) 饱和的具体叙述，本批次没有相应的原始计算，只能作为历史
解释保留。

\paragraph{后续审计改变了结论。}
U002 的 \texttt{corr3=None} 归一化修复使“归一化正确但数值很小”不再成立：在缺少
\(N_{\rm conf}\) 因子时，源和组态的乘积被多乘，旧 ratio、\(c_0\) 和
Sumratio 数值必须整体撤回。由此，dev5 的价值从“已经得到物理 PDF”转为
“保留了可重跑、可对照的实现链和失败样本”。这正是本章采用历史-only 解释的
原因。

\section{4150：逐层复现矩阵}
\label{sec:b08-4150}

4150 材料的优势是把“复现”拆成互不等价的层。表~\ref{tab:b08-4150-layers}
给出本批次能够支持的逐层陈述。表中所有相对差均为历史检查器输出，不能外推到
未检查的动量、时间片或算符几何。

{\small
\begin{longtable}{>{\raggedright\arraybackslash}p{2.6cm}>{\raggedright\arraybackslash}p{3.5cm}>{\raggedright\arraybackslash}p{2.5cm}>{\raggedright\arraybackslash}p{2.3cm}>{\raggedright\arraybackslash}p{2.0cm}}
\caption{4150 逐层对照矩阵。}\label{tab:b08-4150-layers}\\
\toprule
检查层 & 比较对象 & 基准 & 结果 & 允许主张 \\
\midrule
\endfirsthead
\toprule
检查层 & 比较对象 & 基准 & 结果 & 允许主张 \\
\midrule
\endhead
低层公式 & eigvec、phase、VdV、VVV、Clover、dual、OPE &
同场独立重写 & 最大 \(1.36\times10^{-15}\) & 公式/接口一致 \\
4D HYP OPE & x/y/z 三方向 9 通道 & 同场参考成品 & 9/9 PASS & 受控非极化通道一致 \\
3D HYP & 1/3/5 次、z 向 3 通道 & 无配对成品 & 运行成功，状态 unverified & 不可宣称通过 \\
费米子两点 & contract、nopol & 参考投影公式 & 多数动量 PASS & 中间/投影对象一致 \\
逐时间 VVV & 原始中间文件 & 缺失 & unverified & 不可替代为完整收缩 \\
momsmear & contract 到 \(P_+\)、边界、nopol & 输出投影关系 & 58/58 PASS & 输出关系自洽 \\
独立 smeared peram & 独立传播子重算 & 不存在 & unverified & 不可由普通 peram 推断 \\
极化输出 & 标准 \texttt{li} 与旧 \texttt{il} & 成品数组 & 74 PASS；41 仅旧轴序；1 diff；58 缺失 & 轴序未解决 \\
effmass & 7 个聚合资产、25 行 & raw/twoptall & 7/7 PASS & 4150 索引为 \(k=2\) \\
3pt, ratio,\newline barematrix & 同动量/同几何成品 & 缺失 & unverified & 暂不建立下游主张 \\
\bottomrule
\end{longtable}
}

低层和 4D HYP 的通过值得保留，但必须同时保留其基准说明：低层比较是同一规范场
上的独立重写公式，不是外部 raw VVV；momsmear 检查只复算输出关系，不重建
smearing 传播子。一个正确的“输出投影”不能证明输入 perambulator、动量相位和
核子收缩同时独立正确。最有信息量的反而是未验证项：3D HYP、逐时间 VVV、独立
smeared perambulator 和 3pt/ratio 下游均没有可配对基准。

\section{dev20：从可运行到失败关闭}
\label{sec:b08-dev20}

\paragraph{源分辨两点函数。}
正确的 ensemble ratio 不能在分析层丢掉 source 轴。令
\(C_{2,s}(t)\) 为 source \(s\) 的两点函数、\(O_s(t,\tau,z,b_T)\) 为同 source
的胶子结构，并在每个 replica 内取配置与 source 的平均，则可写成
\begin{equation}
  R_{\rm rep}(t,\tau)
  =\frac{\left\langle \overline{C_2O}
  -\overline{C_2}\,\overline{O}\right\rangle_{\rm conf}}
  {\left\langle \overline{\mathrm{Re}\,C_2}\right\rangle_{\rm conf}},
  \qquad
  \overline{X}=\frac{1}{N_s}\sum_{s=1}^{N_s}X_s.
  \label{eq:b08-source-ratio}
\end{equation}
分母和乘积项都必须保留两层归一化；先算源平均、再把一维时间序列当作绝对时间
中心二次平移，会产生不可恢复的相位和归一化误差。

\paragraph{解析反例。}
旧 \texttt{corr3=None} 路径在 \(N_{\rm conf}\) 维已经求和后只除了 \(N_s\)，随后又对
配置和 source 做双重平均。取 \(C_2=[1,3]\)、\(O=[2,6]\) 时，正确平均为
\begin{equation}
  \frac{\overline{C_2O}}{\overline{C_2}}=1,
  \label{eq:b08-correct-example}
\end{equation}
而缺少 \(N_{\rm conf}\) 的旧实现给出一与六的比值，即 6。这个反例不是数值
噪声，而是归一化错误；它直接解释了为什么旧 ratio、\(c_0\) 与 Sumratio 不能
通过“看起来有限”保留下来。

\paragraph{流窗口与诊断。}
在本批次数据中，\(\tau=3a^2\) 给出 \(r_f/a=\sqrt{24}=4.899\)，而非零
横向扫描为 \(b/a=0,\ldots,4\)。因此所有非零 \(b\) 点都不满足
\(r_f<b\)，小 \(z\) 点也不满足 \(r_f<z\)。这使流化算符的场强和 Wilson 线
可能抹去待测横向结构，故所有相关点只能是诊断点。受控扫描在
\(\tau=0.5a^2\) 与 \(1.0a^2\) 下观察到 \(P_{400}\) 中心质量较稳定，
但 P000/P200 仍存在质量或 flow-time 敏感性，且 replica fallback 非零。
两点扫描不构成 \(\tau\to0\) 外推。

\begin{table}[htbp]
\centering
\caption{dev20 诊断链的代表状态。表中仅复述历史审计记录，不是新测量。}
\label{tab:b08-dev20-status}
\small
\begin{tabularx}{\textwidth}{>{\raggedright\arraybackslash}p{3.1cm}>{\raggedright\arraybackslash}p{4.1cm}>{\raggedright\arraybackslash}X}
\toprule
项目 & 记录值/状态 & 验收含义 \\
\midrule
系综 & \(\beta=6.20\)，\(24^3\times72\)，11 组态 & 单格距、小样本；不是生产系综 \\
源轴 & shape \((72,72)\)，轴语义 source,dt & source 信息保留；禁止静默源平均 \\
Sumratio 质量 & P000 \(52/65\)，P200 \(51/65\)，P400 \(65/65\) & 严格 loader 要求 65/65 且无 replica fallback \\
fallback & 1.46\%、1.23\%、0.12\% & 即使 P400 全通过中心质量，仍不能忽略 replica 失败 \\
flow & \(\tau=3a^2\)，\(r_f=4.899a\) & 所有非零 \(b/a\le4\) 违反流窗口 \\
soft/mixing & 无同几何 soft/rapidity，无完整 \(Z_{ij}\) & TMD 因子化链未闭合 \\
链条状态 & diagnostic\_only；physical\_acceptance=false & 只能报告诊断量 \\
\bottomrule
\end{tabularx}
\end{table}

\paragraph{未解决项。}
当前数据没有同流时、同 staple 几何、同伴随表示的 soft/rapidity；没有完整
胶子-夸克单态混合矩阵；没有多个格距、多个流时的受控极限；Fourier 大
\(\lambda\) 尾部、端点 \(x\to0,1\)、流步长 \(0.05\) 与参考 \(0.01\) 的收敛，
以及统一 fm-to-GeV 常数仍未闭合。任何一项缺失都足以阻止“物理胶子 TMD-PDF”
的发布级主张。

\section{第16章应采用的综合判据}

把上述材料压缩成可写作的判据，可表示为
\begin{equation}
  \text{物理主张}
  \Longleftarrow
  \text{源分辨归一化}
  \wedge\text{流/尺度窗口}
  \wedge\text{soft/rapidity}
  \wedge\text{混合}
  \wedge\text{多尺度极限}.
  \label{eq:b08-claim-gate}
\end{equation}
式~\eqref{eq:b08-claim-gate} 中右侧每一个因子都必须有独立来源和独立失败模式。
历史报告的顺序大体是“先让整条链跑通，再补物理门”；当前专著应把顺序反过来：
定义验收对象、建立每层的 oracle 或对照、再决定哪些数值允许进入物理论述。
这样既保留了 dev5/stab1 的工程价值，也不会让早期有限结果重新获得已经撤回的
物理地位。


\clearpage
\chapter{真实数据复现与数值对照证据}
\label{sec:b09-real-data}

本节把 test6 至 test9 以及组态 4150 的运行记录压缩为一条可审计的证据链。这里的“复现”有三种不同强度：同一输入上的逐位回归、参考公式与 PyQCD 实现的数组级对照、以及在真实组态上完整跑通物理链。三者不能相互替代。一个有限矩阵元、一个合格形状或一次成功退出，只证明对应层级的契约成立；只有当算符几何、重正化、谱学投影、统计误差、软/快度减除、匹配和物理极限同时闭合时，才能把中间量升级为物理 TMD-PDF。

\section{证据分级与案例谱系}

本批次最重要的结论不是某个单一数值，而是证据随输入规模、链长和几何覆盖逐级增强，同时在几何与统计上仍保留明确断点。表~\ref{tab:b09-case-lineage} 给出各单元在专著中的位置。test7 和 test8 只有 1 个与 2 个组态，且报告存在缺图或内部矛盾，只作为失败关闭和流程演进的历史材料；test6 有 879 个组态，但目标是质子有效能量而非完整 TMD；test9 才把梯度流、胶子 OPE、不相连比值、自重整化和准 TMD 诊断连接起来；4150 对照则把重点放在真实输入、轴序、复数语义和逐层输出。

\begin{longtable}{@{}p{0.10\textwidth}p{0.22\textwidth}p{0.26\textwidth}p{0.32\textwidth}@{}}
\caption{B09 案例谱系与可支持的主张}\label{tab:b09-case-lineage}\\
\hline
单元 & 输入与目的 & 最强证据 & 专著中的结论等级 \\
\hline
\endfirsthead
\hline
单元 & 输入与目的 & 最强证据 & 专著中的结论等级 \\
\hline
\endhead
U017 & 879 组态，\(24^3\times72\)，\(P=(0,0,6)\)，提取质子 \(E_0\) & 四组 \texttt{corr2} 与参考数组最大绝对差为 0；平均方向 \(E_0=3.069(9)\) GeV & 真实统计链的数值回归，可用于验证实现等价性，不是新的能谱测量 \\
\hline
U018 & 1 组态 test7 冒烟测试 & 断链与缺图被记录；P0/P2 均错误写成 0.000 GeV & 历史失败案例，只进入限制说明 \\
\hline
U019 & 2 组态 test8 冒烟测试 & P0/P2 图与输出投影出现；无数值误差与拟合窗 & 流程升级证据，不采用其中心值 \\
\hline
U020 & 10 组态 test9 全链 & flow、OPE、ratio、self-renorm、quasi TMD、匹配和 CS 核全部落地 & 全链诊断，统计与几何不足以形成物理 TMD \\
\hline
U025 & 单组态 4150 功能对照 & 52 项中 51 项直接通过、1 项白名单；复 GEVP 残差约 \(10^{-15}\) & 实现语义和真实输入验收 \\
\hline
U027 & 4150 donghx 扩展现 & 低层 7/7、HYP 直线通道 9/9、正宇称投影 58/58、effmass 25/25 & 低层与 2pt 输出级的定量复现；完整 TMD 未闭环 \\
\hline
\bottomrule
\end{longtable}

表中的“通过”必须带对象和范围。例如，单组态功能对照通过，不等于多组态物理结果成立；直线 Wilson 线的 OPE 通过，不等于带横向 staple 的 TMD 算符通过；准 TMD 与匹配后曲线的数量级接近，也不等于匹配核、标度和误差预算已经闭合。

\section{test6：质子有效能量的逐位回归}

\paragraph{定义与算法。}
test6 读取每个组态、每个动量方向上的 \(72\times72\) 复数两点函数切片。实现先做平移不变的相对时间移位，再对源时间 \(t_i\) 平均，最后使用 delete-one jackknife 形成统计样本；输出数组形状为 \(879\times20\)。在固定平台窗 \(t=3,\ldots,7\) 内，对每个 jackknife 样本拟合双指数形式
\begin{equation}
  C(t)=c_0e^{-E_0t}\left(1+c_1e^{-dE\,t}\right).
  \label{eq:b09-test6-fit}
\end{equation}
这里 \(E_0\) 是拟合得到的有效能量参数，\(dE\) 概括第一激发态间隙。式~\eqref{eq:b09-test6-fit} 不是谱函数反演，不能据此声称已经分离全部激发态。有效质量
\begin{equation}
  m_{\mathrm{eff}}(t)=\log\left|\frac{C(t)}{C(t+1)}\right|
  \label{eq:b09-test6-meff}
\end{equation}
用于检查平台，但不能替代相关拟合与窗口稳定性分析。

\paragraph{结果与尺度。}
输入格距为 \(a=0.1053\) fm，动量 \(P=(0,0,6)\)。平均方向的中心结果为
\begin{equation}
  E_0^{\mathrm{ave}}=1.640(5)\,a^{-1}=3.069(9)\ \mathrm{GeV},
  \qquad \chi^2/\mathrm{dof}=2.4.
  \label{eq:b09-test6-result}
\end{equation}
x、y、z 三个方向分别给出 3.084(15)、3.138(32)、3.063(11) GeV。报告还用约 3.14 GeV 的色散预期作宽泛检查，但该比较没有给出独立的连续极限或体积外推，因此只能视为自洽性检查。

\paragraph{验证断言与失败边界。}
四个 \texttt{corr2} 数组与参考运行结果的形状、dtype 和 NaN 掩码一致，最大绝对差为零。这说明 PyQCD 的切片读取、相对时间、源平均、jackknife 索引和数组布局在该固定输入上与参考链等价。它不说明 879 个组态已经覆盖自相关、拓扑扇区或连续极限。报告还记录了两个必须保留的边界：前两个时间列因周期边界卷绕和零值除法退化；y 方向条件数约 \(1.3\times10^8\)，拟合最弱。源报告把 \(1.677(17)\) 的误差误写成“17\%”，专著采用括号中的绝对误差，不采用该百分比。

\section{test7 与 test8：从冒烟测试到可审计失败}
\label{subsec:b09-smoke}

test7 只有 1 个组态，test8 增至 2 个组态。它们覆盖了蒸馏两点函数、三点/两点比值、质子能量、FH 变换和三方向差异等入口，但没有形成可辨识的统计测量。test7 的多幅关键图缺失，并把 P0 质量与 P2 能量都写成 0.000 GeV，同时声称 P0 与 1.12 GeV 的已知结果一致；这是内部矛盾。test8 给出了约 1.060 GeV 的 P0 平台和 1.517 GeV 的 P2 平台，但没有报告平台窗、拟合函数、误差或重采样，两个组态也不足以估计系综误差。test8 的 PDF 还没有包含 TeX 末尾的结论段，说明构建产物与源文件不同步。

因此，这两个单元在专著中只承担两个作用。第一，它们展示了从 1 个到 2 个组态可以修复“图完全缺失”的运行问题，却没有修复统计问题。第二，它们提供了失败关闭规则：零信号、内部矛盾和缺图必须标为未验证，不能用“与基线一致”一句话覆盖。test7/test8 的具体中心值不进入第 16 章的匹配或 TMD 结论。

\section{test9：10 组态下的完整诊断链}
\label{subsec:b09-test9}

\paragraph{物理链。}
test9 使用 \(\beta=6.20\)、\(m_\mu=0.2770\)、\(m_s=0.2400\)、\(24^3\times72\) 的 10 个组态，流时取 \(\tau=3a^2\)。链为：蒸馏两点函数，Wilson flow，flowed gauge 上的胶子 TMD 矩阵元，不相连三点函数与真空扣除，逐 \((z,b_\perp)\) 拟合裸 \(c_0\)，用 \(h_R=c_0(z,b)/c_0(0,b)\) 自重整化，得到准 TMD，最后尝试 NLO 匹配和 Collins--Soper 演化。这个顺序不可交换：如果先对最终曲线做归一化，会掩盖 \(z,b_\perp\) 依赖；如果匹配前没有同一几何的软/快度减除，匹配核就没有与其配套的算符矩阵元。

\paragraph{Wilson flow。}
报告采用 RK3 积分
\begin{equation}
  \partial_t V(t,x)=-g_0^2\left[\partial_x S(V)\right]V(t,x),
  \qquad \varepsilon=0.05,\quad 60\ \text{steps},\quad \tau=3a^2.
  \label{eq:b09-flow}
\end{equation}
对 10 个组态，作用量密度从约 \(0.1495\)--\(0.1499\) 单调降到 \(0.0544\)--\(0.0548\)，最大非酉度约 \(5.7\times10^{-6}\)。这是流方程方向和数值近似保持 SU(3) 的强检查，但不是短流时匹配的证据。\(\tau=3a^2\) 对应
\begin{equation}
  r_f=\sqrt{8\tau}=\sqrt{24}\,a\approx4.90a.
  \label{eq:b09-flow-radius}
\end{equation}
若目标横向分离满足 \(b_\perp\lesssim r_f\)，流化会直接抹平待测结构；若只算一个流时，也无法做 \(\tau\to0\) 外推。故式~\eqref{eq:b09-flow} 的实现正确性不能替代流窗口。

\paragraph{胶子算符与矩阵元。}
报告使用的组合为
\begin{equation}
  O(z,b_\perp)=M^{tx;tx}+M^{ty;ty}-2M^{xy;xy},
  \qquad
  M^{\mu\lambda;\nu\rho}=\left\langle
  \operatorname{tr}\left(F_{\mu\lambda}WF_{\nu\rho}W^\dagger\right)
  \right\rangle,
  \label{eq:b09-tmd-operator}
\end{equation}
其中 \(F\) 为 Clover 场强，\(W\) 为 Wilson 线。该式是报告所采用的特定 Lorentz 投影，不能在没有证明的情况下替换为其他胶子 TMD 投影。单组态 \(O(0,0)\) 的时间平均约 9.9 而标准差约 31.7，已经显示矩阵元被大时间片涨落主导；这说明只看一个中心值不足以判断信号质量。

\paragraph{不相连比值。}
对两点函数 \(C_2(dt)\) 和胶子 OPE，报告构造
\begin{equation}
  C_3(dt,d\tau)=C_2(dt)\,\mathrm{OPE}(d\tau),
  \qquad
  R=\left\langle
  \frac{C_3-C_2\langle\mathrm{OPE}\rangle}{C_2}
  \right\rangle_{t_i},
  \label{eq:b09-ratio}
\end{equation}
并拟合
\begin{equation}
  R=c_0+c_1e^{-dE\,d\tau}+c_1e^{-dE(dt-d\tau)},
  \qquad t_{\mathrm{sep}}\in[7,10],\quad \mathrm{cut}=6.
  \label{eq:b09-ratio-fit}
\end{equation}
式~\eqref{eq:b09-ratio-fit} 中的两个指数使用同一个 \(c_1,dE\) 是报告采用的受限模型；专著应把它作为模型选择，而不是从第一性原理唯一推出的形式。报告得到 \(c_0=O(1)\)--\(O(10)\)、误差 \(O(\text{几})\)，统计噪声与信号同量级。

\paragraph{统计秩的关键纠正。}
报告称 10 组态不足以形成 14 个数据点的满秩 jackknife 协方差，因此改用 200 次 bootstrap 获得“满秩”协方差。这个措辞必须纠正。令中心化数据矩阵为 \(X\in\mathbb C^{N\times p}\)，其中 \(N=10\) 是独立组态数，则样本协方差
\begin{equation}
  \widehat C=\frac{1}{N-1}X^\dagger X
  \label{eq:b09-sample-cov}
\end{equation}
满足
\begin{equation}
  \operatorname{rank}\widehat C=\operatorname{rank}X\le \min(N-1,p)=9.
  \label{eq:b09-rank-bound}
\end{equation}
从 \(N\) 个样本有放回抽取 \(N_b\) 次只改变 replica 的数量，不改变由原始样本张成的空间。即使 \(N_b=200\)，未正则化的样本协方差秩仍至多 9；使矩阵可逆的 SVD 截断、shrinkage 或对角加载属于正则化，必须报告阈值及其偏差。正确表述是“200 次重采样配合正则化后数值可逆”，而不是“200 次 bootstrap 提供了满秩统计”。

\paragraph{准 TMD、匹配与 CS 核。}
报告采用的准 TMD 提取为
\begin{equation}
  x\widetilde g(x,b_\perp)
  =\frac{1}{2\pi P_z}\frac{P_z^2}{E^2}
  \int dz\,\cos(xzP_z)\,h_R(z,b_\perp),
  \qquad
  h_R(z,b)=\frac{c_0(z,b)}{c_0(0,b)}.
  \label{eq:b09-quasi}
\end{equation}
在 \(b_\perp=0\) 处得到 \(\langle x\rangle_{\mathrm{quasi}}\approx0.57\)，报告称匹配后约 0.59。该量级与胶子动量分数的常见范围相容，但报告没有给出所乘匹配核、UV 与快度标度、匹配阶数、尺度变化、误差传播或 Fourier 尾部模型。因此式~\eqref{eq:b09-quasi} 的结果只能标为量级诊断。

Collins--Soper 核采用两动量比值法，在 \(b=0.105\) fm 给出 \(K=1.6\)，并称 \(b\gtrsim0.2\) fm 后因噪声触发 clamp。clamp 是数值稳定装置，不是物理饱和。没有软函数、同几何 Wilson 圈、多个 \(P_z\) 的一致性和误差带时，不能把这一条数值写成核的确定测量。

\section{组态 4150：从功能对照到复 GEVP}
\label{subsec:b09-4150}

\paragraph{真实输入与张量契约。}
4150 的原始 Clover 规范场、本征矢量、light perambulator、HYP 场和若干两点函数成品均可读。规范场采用
\begin{equation}
  U\in\mathbb C^{N_t\times N_z\times N_y\times N_x\times4\times3\times3},
  \label{eq:b09-gauge-layout}
\end{equation}
方向 \(0,1,2,3\) 对应空间 \(x,y,z\) 与时间。真实链接的最大非酉度约 \(1.33\times10^{-15}\)。这个契约必须固定，否则后续的 \(z\) 位移、横向位移和时间投影会静默交换轴。

\paragraph{复 GEVP。}
PyQCD 与参考实现在 L20 用例上出现一个系统性差异。双方先构造 Hermitian 矩阵
\begin{equation}
  C_H(t)=\frac{C(t)+C^\dagger(t)}{2},
  \label{eq:b09-hermitize}
\end{equation}
再求解广义本征值问题
\begin{equation}
  C_H(t)v_n(t,t_0)=\lambda_n(t,t_0)C_H(t_0)v_n(t,t_0).
  \label{eq:b09-complex-gevp}
\end{equation}
参考路径进一步把 \(C_H\) 投影为 \(\operatorname{Re}C_H\)，而 PyQCD 保留复矩阵。Hermitian 矩阵的本征值是实数，但本征矢仍可为复向量；丢掉虚部会改变原方程。在随机复正定输入上，PyQCD 与参考对原复 GEVP 的残差分别约 \(1.02\times10^{-15}\) 和 \(1.17\times10^{-1}\)；纯实投影后两者逐位一致。另一个解析谱 \(C=Ue^{-Et}U^\dagger\) 上，本征值最大误差分别约 \(6.66\times10^{-16}\) 和 \(8.89\times10^{-2}\)。

据此，52 个功能用例中 51 个直接通过，L20 作为参考语义差异进入白名单，硬失败为零。白名单只改变验收分类，不改变 JSON 中的差异值。专著应明确：在实数子空间上逐位一致，并不能证明参考实现了原复问题。

\paragraph{低层与两点函数证据。}
4150 的 eigvec、phase、VdV、VVV、Clover、dual 和直线 OPE 共 7 个对象通过，最大相对误差约 \(1.36\times10^{-15}\)。4D10 HYP 的 \(x/y/z\) 三个方向共 9 个非极化直线通道通过，最大相对差约 \(3.31\times10^{-15}\)。两点函数输出包含 7 个根目录、58 个动量组，正宇称投影 58/58 通过；全体最大相对 \(L_2\) 误差约 \(8.29\times10^{-7}\)，最大绝对差约 \(1.98\times10^{-9}\)。effmass 聚合把 4150 定位为序列 \(4050+50k\) 的 \(k=2\)，7 个资产、25 个动量行通过，时间汇总最大相对误差约 \(2.58\times10^{-16}\)。

极化输出给出更细的异常清单：116 个实际存在的极化数组中，74 个与标准 \(li\) 投影一致，58 个文件缺失；41 个只与旧 \(il\) 转置轴序一致，1 个 `momsmear2x`、\(P=(0,0,6)\) 的 `Cg5g4/pol15` 对两种轴序都不一致。该异常槽位范数约 \(3.30\times10^{-9}\)，接近零，说明相对误差会因分母极小而失真；此时应同时报告绝对值、范数和上游 `contract` 状态。不能把 41 个轴序命中算作标准通过，也不能把 58 个缺失文件算作失败。

\paragraph{实现的证据等级。}
4150 证明了真实 HYP/LIME 读取、Clover/对偶张量、直线 Wilson-line OPE、蒸馏顶点、两点函数缩并、正宇称投影、momentum-smear 输出关系和 effmass 聚合的逐层一致性。它没有证明独立 smeared perambulator、逐时间 VVV、完整 3pt/ratio/barematrix 输入链或非零横向 staple。因此 4150 是 L0/L1 级实现证据，而不是 L2/L3 的完整 TMD 结果。

\section{本批次可接受的结论边界}

综合上述单元，可以把结果归纳为四条判据。第一，固定输入上的逐位一致是最强的实现等价性证据，但不自动说明统计或物理极限。第二，真实组态上的数组级对照必须同时报告容差、dtype、形状和状态语义；缺失文件既不是通过也不是失败。第三，少量组态只能验证流程与稳定性，协方差的秩由独立组态数限制，重采样次数不能制造信息。第四，低层直线 OPE 和两点函数通过后，仍需独立验证横向 staple、soft/rapidity、匹配和连续极限。

\chapter{准 TMD、匹配与真实组态上的几何边界}
\label{sec:b09-quasi-to-lightcone}

本节把第 14 章的案例转成第 16 章可用的物理接口。重点是说明“代码跑通”“准矩阵元可提取”“匹配到光锥”和“得到可发表的胶子 TMD-PDF”不是同一个命题。test9 与 4150 的证据足以支持前两层的部分内容，却不足以支持最后一层。

\section{从关联函数到准 TMD}

\paragraph{定义。}
设核子动量为 \(P=(0,0,P_z)\)，胶子算符沿空间方向分离 \(z\)，横向分离为 \(b_\perp\)。test9 使用式~\eqref{eq:b09-tmd-operator} 的特定 Lorentz 组合，先从三点/两点比值拟合裸矩阵元 \(c_0(z,b_\perp)\)，再构造自重整化矩阵元
\begin{equation}
  h_R(z,b_\perp)=
  \frac{c_0(z,b_\perp)}{c_0(0,b_\perp)}.
  \label{eq:b09-self-renorm}
\end{equation}
式~\eqref{eq:b09-self-renorm} 的物理依据是：若固定 \(b_\perp\) 下存在一个与 \(z\) 无关、可乘的 UV 因子，那么分子分母共同因子相消。它并不证明所有 \(z\) 依赖的 Wilson-line 自能、端点发散、算符混合和快度结构都被同一个因子吸收。因此，自重整化是一个需要独立 residual fitting 和方案比较的假设，不是“除以零分离就自动有限”的定理。

报告用式~\eqref{eq:b09-quasi} 的余弦变换得到准 TMD。这里 \(h_R\) 仍是 Euclidean 矩阵元，不是光锥分布函数。余弦变换把 \(z\) 空间映射到 \(xP_z\) 空间；有限 \(z_{\max}\)、离散 \(z/a\)、周期边界和大 \(z\) 噪声会产生振铃与尾部模型依赖。一个合理的提取至少要报告 \(z_{\max}\)、尾部截断、不同 \(P_z\) 的一致性和变换后的支持区间。

\paragraph{算法。}
本批次可以确认的最小算法链是：
\begin{enumerate}
  \item 用同一规范场构造 Clover 场强和 Wilson 线，计算 \(O(z,b_\perp)\)；
  \item 由同一源、汇和插入时间形成 \(C_2,C_3\)，在每个源时间上做真空扣除和比值；
  \item 在固定拟合窗内提取 \(c_0(z,b_\perp)\)，并保留每个组态的 replica；
  \item 在固定 \(b_\perp\) 上做式~\eqref{eq:b09-self-renorm}；
  \item 对 \(z\) 做式~\eqref{eq:b09-quasi} 的变换，得到准 TMD；
  \item 在同一几何和同一表示下接入 soft/rapidity 因子、CS 核和硬匹配；
  \item 对 \(a,\tau,P_z,b_\perp,z_{\max},L\) 做联合误差与极限分析。
\end{enumerate}
第 6 步不是可选的后处理。没有它，前面的曲线只能叫准矩阵元的 Fourier 像。第 7 步也不能由“单点看起来合理”替代；必须给出变化范围和误差传播。

\section{匹配与 Collins--Soper 演化的证据要求}

\paragraph{物理结构。}
准 TMD 到光锥 TMD 的匹配不是把一个函数整体乘以常数。它包含硬系数、卷积变量、可能的味混合和 soft/rapidity 方案转换。可以把结构写成
\begin{equation}
  \widetilde f_g(x,b_\perp;P_z,\mu,\zeta_z)
  =\int_x^1\frac{dy}{y}\,
  C_{g/g}\!\left(\frac{x}{y},b_\perp;\mu,P_z,\zeta_z\right)
  f_g(y,b_\perp;\mu,\zeta)
  +\text{更高扭度与混合项},
  \label{eq:b09-match-structure}
\end{equation}
其中 \(\widetilde f_g\) 是 Euclidean 准对象，\(f_g\) 是光锥 TMD，\(C_{g/g}\) 和混合项必须与算符表示、Path 几何和快度方案一致。式~\eqref{eq:b09-match-structure} 只表达接口结构，不提供未在源报告中给出的显式 NLO 系数。专著若使用具体核，必须另引相应文献并核对约定。

\paragraph{CS 核。}
Collins--Soper 核控制 TMD 对快度标度的演化。报告选择两动量比值法和 \(K(b_\perp)\) 记号，在 \(b=0.105\) fm 给出 \(K=1.6\)，并称 \(b\gtrsim0.2\) fm 后 clamp 饱和。这个结果的可支持表述是“在该输入、该流时、该几何下，两动量比值给出一个量级为 1 的核，并在高 \(b\) 被数值保护截断”。不支持“CS 核已非微扰确定”或“所有 \(b\) 上演化正确”。合格的 CS 证据需要：
\begin{itemize}
  \item 同一表示、同一 Wilson-line 几何和同一 soft/rapidity 方案；
  \item 至少两个 \(P_z\) 或两个动量组合，检验比值和平台；
  \item 去掉 clamp 前后的中心值、误差和失败率；
  \item 小 \(b\) 与微扰极限、大 \(b\) 与非微扰模型之间的连续接口；
  \item 将核的不确定度传播到最终 TMD，而不是只在核图上画误差带。
\end{itemize}

\paragraph{test9 的匹配诊断。}
test9 的 \(\langle x\rangle_{\mathrm{quasi}}\approx0.57\) 与“匹配后”约 0.59 可以用于检查数量级和代码路径，但源报告没有提供匹配核的显式参数、标度、阶数、卷积实现、误差或不同 \(P_z\) 的一致性。因此这两个数不能写入“已测得胶子 TMD 动量分数”的结论。它们最多说明匹配模块被调用、输出落在物理上不荒谬的范围内，以及后续需要用完整误差预算替换单点比较。

\section{Wilson 几何：直线 OPE 与横向 staple}

\paragraph{两种对象。}
4150 的低层对照主要是沿单个空间方向 $\hat d$ 的直线 OPE：
\begin{equation}
  O^{FF}_{\mu\nu}(z)=
  \sum_{\boldsymbol x}\operatorname{tr}\left[
  F_{\mu\nu}(x+z\hat d)W^\dagger(0\to z)F_{\mu\nu}(x)W(0\to z)
  \right].
  \label{eq:b09-straight-ope}
\end{equation}
test9 的目标对象则包含横向位移 \(b_\perp\) 和相应 staple 路径。两者共享场强与 Wilson-line 技术，但不是同一几何。直线零横向截面的通过只锁定了端点、方向、颜色缩并和场强归一化；它没有检验横向连接、转角、横向端点或非零 \(b_\perp\) 的 UV 结构。

\paragraph{几何判据。}
在把任何结果送入 ratio 之前，至少要满足：
\begin{equation}
  \begin{aligned}
  &\text{same gauge}\quad\land\quad
  \text{same path}\\
  &\quad\land\quad
  \text{same representation}\quad\land\quad
  \text{same soft/rapidity scheme}.
  \end{aligned}
  \label{eq:b09-geometry-gate}
\end{equation}
“same gauge”不仅指组态编号相同，还要求使用的 HYP/flow/stout 链接相同；“same path”包括正负方向、起点、纵向长度和横向端点；“same representation”包括端点场强与 Wilson line 的伴随表示；“same soft/rapidity scheme”包括软函数、Wilson 圈减除和快度标度。缺少任一项时，数值差异可能来自方案而不是物理。

\section{尺度窗口与失效条件}

\paragraph{必须同时满足的窗口。}
对本批次可用的最低限度，应报告
\begin{align}
  a&\ll r_f,\;1/\mu, \label{eq:b09-window-a}\\
  \Lambda_{\mathrm{QCD}}&\ll 1/b_\perp,\;P_z\ll 1/a, \label{eq:b09-window-ope}\\
  M_N,\Lambda_{\mathrm{QCD}}&\ll P_z\ll \pi/a, \label{eq:b09-window-lamet}\\
  a&\ll r_f\ll b_\perp,z,\qquad P_zr_f\ll1, \label{eq:b09-window-flow}\\
  b_\perp,L_\parallel&\ll L_s/2. \label{eq:b09-window-volume}
\end{align}
这些不是可任意挑选的“最优参数”，而是一个交集。test9 的 \(\tau=3a^2\) 给出 \(r_f\approx4.90a\)；如果某个目标的 \(b_\perp\) 小于或接近该值，式~\eqref{eq:b09-window-flow} 不成立。报告还应给出流时扫描；只有一个流时无法区分短流时匹配、离散效应和大 \(z\) 噪声。

\paragraph{几何与统计的空交集。}
如果 \(b_\perp\lesssim r_f\) 且统计量只有 10 个组态，那么即使中心曲线平滑，也没有同时满足流窗口和统计误差的物理点。此时正确的状态是“该参数点给出诊断上界”，而不是选择一个看起来稳定的 bin 拟合。若 \(b_\perp\gtrsim L_s/2\)，周期镜像又会污染矩阵元；若 \(P_zr_f\) 接近 1，boost 与流化尺度不可分离。专著应把这类点从物理结果表中明确剔除，并保留失败率。

\section{验收矩阵}

表~\ref{tab:b09-tmd-acceptance} 给出从已有证据到最终物理量的最小门控。它能帮助读者把第 14 章的“通过”转换为第 16 章的允许表述。

\begin{table}[htbp]
\centering
\caption{从已有证据到最终物理量的验收矩阵。}
\label{tab:b09-tmd-acceptance}
\begin{tabularx}{\textwidth}{@{}>{\raggedright\arraybackslash}p{0.17\textwidth}>{\raggedright\arraybackslash}p{0.25\textwidth}>{\raggedright\arraybackslash}X@{}}
\hline
层级 & 已有证据 & 允许的最强表述 \\
\hline
L0 算符与几何 & 4150 Clover/dual/直线 OPE；test9 的 \(O(z,b)\) 代码路径 & 场强、端点、轴序和一个特定直线投影已实现并做过数组级检查 \\
\hline
L1 矩阵元 & 4150 两点函数 58/58；test9 的 10 组态 ratio 与自重整化 & 固定输入下可生成中间矩阵元和诊断曲线，统计不确定度可能很大 \\
\hline
L2 方案一致性 & test9 尝试 CS 核与匹配；缺少同几何 soft\newline rapidity、显式匹配参数和误差链 & 只能称为方案接口原型，不能称为完整光锥 TMD \\
\hline
L3 物理极限 & 单一格距、单一流时、少量组态；无联合连续、手征和大动量外推 & 不支持“已测得物理核子胶子 TMD-PDF” \\
\hline
\end{tabularx}
\end{table}

\section{本章结论}
L0/L1 证据值得保留，因为它们排除了许多实现错误：真实 HYP 读取、复数 GEVP、两点函数轴序、正宇称投影、直线 OPE、Fourier 前的矩阵元形状都可以被逐项复核。L2/L3 仍为空，因为本批次没有给出同几何 soft/rapidity、完整匹配核、尺度扫描和多格距极限。第 16 章应把 test9 的准 TMD 和 CS 图作为“可复现的全链诊断”，而不是作为物理 TMD 的最终数值。


\clearpage
\chapter{系统误差与验证协议}

\section{误差预算}

表~\ref{tab:errors} 给出了胶子 TMD 计算的最小误差矩阵。每一项都必须有独立的扫描或解析边界，不能用总的 bootstrap 误差替代。

\begin{table}[htbp]
\centering
\caption{核子胶子 TMD-PDF 的系统误差矩阵。}
\label{tab:errors}
\small
\begin{tabularx}{\textwidth}{L{3.0cm}L{3.3cm}Y}
\toprule
误差源 & 控制变量 & 验收证据 \\
\midrule
格距与连续极限 & \(a\), \(aP_z\) & 至少多个格距，检查 \(a^2\) 和 \(a^2P_z^2\) 项 \\
有限体积 & \(L_s\), \(b_T\), \(L_\parallel\) & \(b_T<L_s/2\)，长 Wilson 线平台与体积扫描 \\
激发态 & \(t_{\rm sep}\), 拟合模型 & 单态/双态/GEVP 一致性与窗口稳定性 \\
重正化 & \(Z_L,Z_O,Z_{ij}\), \(\mu\) & 多方案交叉、连续极限、非零动量矩阵元稳定性 \\
rapidity/soft & \(S_R\), \(K_R\), 方案 & 同一表示和几何的软因子、多个 \(P_z\) 的 CS 检验 \\
梯度流 & \(\tau\), \(r_f\) & 多 \(\tau\) 外推和 \(a\ll r_f\ll b_T,z\) 窗口 \\
匹配核 & \(C_{gg},C_{gq}\), \(\alpha_s\) & 阶数变化、标度变化、单态混合误差 \\
高扭度 & \(P_z\), \(M_N\), \(b_T\) & 多个大动量，报告 \(1/P_z^2,1/(b_T^2P_z^2)\) 拟合 \\
Fourier 反演 & \(z_{\max}\), 尾部模型 & 多截断、多模型、正定性/端点检查 \\
Dirac 求逆 & \(\epsilon_{\rm rel}\), \(i_{\max}\), restart, 预条件 & 递归/真实残差、容差扫描、不同求解器与目标可观测量的稳定性 \\
统计 & 组态数、源数、bootstrap & 源分辨、相关性、不同 replica 的稳定性 \\
\bottomrule
\end{tabularx}
\end{table}

\section{统计误差、自相关与相关拟合}

格点数据的时间片、源和组态并不独立。若观测量 \(X_t\) 的时间序列具有自相关函数
\begin{equation}
  \rho_X(\Delta t)
  =\frac{\langle X_tX_{t+\Delta t}\rangle-\langle X_t\rangle^2}
  {\langle X_t^2\rangle-\langle X_t\rangle^2},
  \label{eq:autocorr}
\end{equation}
则有效样本数近似为
\begin{equation}
  N_{\rm eff}
  =\frac{N}{1+2\sum_{\Delta t=1}^{\Delta t_{\max}}\rho_X(\Delta t)},
  \label{eq:neff}
\end{equation}
并用 binning 检查误差是否随 bin 大小达到平台。对胶子回路、真空减除和高动量三点函数，单靠 \(\sigma/\sqrt{N}\) 往往会低估误差。

拟合时应在每个 bootstrap replica 内独立完成重正化、Fourier、匹配和外推，而不是先拟合中心值再对误差做线性传播。对于高度相关的数据，可先做 SVD 或采用带先验的协方差；任何正则化都必须报告奇异值谱、截断阈值和对最终结果的影响。第 5.8 节的协方差公式与本节共同构成统计验收门。

\section{尺度窗口与失败模式}

胶子 TMD 的物理结果必须位于一个互相兼容的多尺度窗口：
\begin{align}
  &\text{UV/连续：} & a&\ll r_f,\;1/\mu &&\ll \text{物理尺度},
  \label{eq:window-uv}\\
  &\text{OPE/匹配：} & \Lambda_{\rm QCD}&\ll 1/b_T,\;P_z &&\ll 1/a,
  \label{eq:window-ope}\\
  &\text{LaMET：} & M_N,\Lambda_{\rm QCD}&\ll P_z &&\ll \pi/a,
  \label{eq:window-lamet}\\
  &\text{流化：} & a&\ll r_f\ll b_T,z &&\text{且 }P_zr_f\ll1,
  \label{eq:window-flow}\\
  &\text{体积：} & b_T,L_\parallel&\ll L_s/2 &&\text{且 }t_s\gg1/\Delta E.
  \label{eq:window-volume}
\end{align}
图~\ref{fig:scale-windows} 把这些不等式组织成“几何长度”和“能量/动量”两条尺度轨道。它们必须同时成立，而不是分别选取各自最好看的参数点。

\begin{figure}[htbp]
\centering
\begin{tikzpicture}[>=Stealth,font=\small]
  \tikzset{
    lenbox/.style={draw=blue!55!black,rounded corners,fill=blue!8,minimum width=2.15cm,minimum height=0.72cm,align=center},
    ir/.style={draw=orange!65!black,rounded corners,fill=orange!10,minimum width=2.25cm,minimum height=0.72cm,align=center}
  }
  \node[font=\bfseries] at (-5.4,2.55) {几何尺度};
  \draw[->,line width=1.1pt,blue!60!black] (-6.1,2.0)--(6.1,2.0) node[right]{IR};
  \node[lenbox] at (-5.0,2.55) {格距 \(a\)};
  \node[lenbox] at (-2.0,2.55) {流半径 \(r_f\)};
  \node[lenbox] at (1.1,2.55) {\(b_T,z\)};
  \node[ir] at (4.6,2.55) {体积边界 \(L_s/2\)};
  \draw[<->,dashed,teal!70!black] (-2.9,3.25)--(2.25,3.25) node[midway,above] {可用的 \(r_f\ll b_T,z\) 窗口};

  \node[font=\bfseries] at (-5.4,-0.8) {能标/动量};
  \draw[->,line width=1.1pt,orange!65!black] (-6.1,-1.35)--(6.1,-1.35) node[right]{IR};
  \node[lenbox] at (-5.0,-0.8) {紫外截断 \(a^{-1}\)};
  \node[lenbox] at (-1.8,-0.8) {\(\mu,\ P_z\)};
  \node[lenbox] at (1.5,-0.8) {\(1/b_T\)};
  \node[ir] at (4.6,-0.8) {\(\Lambda_{\rm QCD}\)};
  \draw[<->,dashed,red!65!black] (-0.45,0.1)--(3.5,0.1) node[midway,above] {OPE/匹配窗口};
\end{tikzpicture}
\caption{胶子 TMD 的共同尺度窗口示意，不对应某组具体数值。上方为长度尺度，下方为能标或动量；只有两条轨道上的不等式同时成立，流化、短距匹配、大动量展开和有限体积控制才属于同一物理论述。}
\label{fig:scale-windows}
\end{figure}

这些窗口不是装饰性不等式。若 \(P_z r_f\sim1\)，流化会改变 boost 相关结构；若 \(1/b_T\lesssim\Lambda_{\rm QCD}\)，微扰匹配已进入非微扰区；若 \(b_T\gtrsim L_s/2\)，周期图像和核子波函数绕回会污染矩阵元；若 \(aP_z\) 接近 \(1\)，格点色散关系和 \(a^2P_z^2\) 误差不可控。

\begin{checkpointbox}{尺度窗口的最小审计}
报告结果前至少回答三个问题：(i) 每个数据点是否满足式~\eqref{eq:window-flow} 和式~\eqref{eq:window-volume}，而不是只在平均参数上满足；(ii) 改变 \(r_f,b_T,z_{\max}\) 或 \(P_z\) 后，中心值移动是否小于目标物理误差；(iii) 被窗口排除的点是否明确标为诊断点，而不是被静默丢弃。窗口交集为空时，任何“稳定平台”都不构成物理证据。
\end{checkpointbox}

表~\ref{tab:failure-modes} 汇总常见失败模式。每一项都应触发明确的“失败关闭”，而不是用更大的误差棒继续报结果。

\begin{table}[htbp]
\centering
\caption{胶子 TMD 计算的常见错误及其最小修正。}
\label{tab:failure-modes}
\small
\begin{tabularx}{\textwidth}{L{3.2cm}L{4.2cm}Y}
\toprule
失败模式 & 典型症状 & 最小修正 \\
\midrule
源平均后再比值 & 归一化随 \(t_s\) 漂移，ratio 出现非物理平台 & 保留 source 轴，按 replica 重算真空协方差 \\
两端 Wilson 线表示不一致 & 线性项残留，随 \(a,z\) 增强 & 同伴随表示、同路径长度、同 staple 几何 \\
Wilson 圈几何不匹配 & \(L_\parallel\) 无平台或 pinch-pole 残留 & 用同端点/同转角的矩形圈，检查 \(L\to\infty\) \\
把 RI/MOM 当万能减除 & 小 \(a\) 或大 \(z\) 格距依赖增强 & 联合 ratio/hybrid/self-renorm，检查连续极限 \\
流时过大 & \(b_T<r_f\)，横向结构被抹平 & 多 \(\tau\) 扫描，只保留满足窗口的点 \\
只用单流时 & 无 \(\tau\to0\) 外推 & 至少三个流时并报告截距 \\
匹配核阶数/标度不匹配 & \(\mu\) 变化导致系统漂移 & 同方案转换，扫描阶数与 \(\mu\) 区间 \\
忽略 \(g-q\) 混合 & 与全局 PDF 比较出现整体偏移 & 加入 \(Z_{ij}/C_{gq}/C_{qg}\) 或给出严格忽略依据 \\
Fourier 端点外推过强 & \(x\) 端点出现振铃/负值 & 多 \(z_{\max}\)、多尾部模型，报告可信区 \\
协方差秩不足 & 拟合参数不可辨识，结果依赖 seed & SVD/正则化、增加统计或减少参数 \\
\bottomrule
\end{tabularx}
\end{table}

\section{分级验收清单}


建议把胶子 TMD 结果分为四级，禁止跨级宣称：
\begin{enumerate}
  \item \textbf{L0：算符与代数。} 场强、伴随链接、staple 路径、投影、归一化和 Fourier 约定均有可执行检查；
  \item \textbf{L1：受控矩阵元。} 源分辨两点函数、Dirac 求逆的双残差门、三点比值、真空减除、谱学拟合和统计误差通过质量门；
  \item \textbf{L2：方案一致性。} 同几何 Wilson 圈、伴随软函数、CS 核、重正化矩阵和硬匹配核属于同一方案；
  \item \textbf{L3：物理极限。} 多个 \(a\)、多个 \(\tau\)、多个 \(P_z\) 和至少一个物理或可控手征点支持下，连续、手征和大动量极限同时闭合。
\end{enumerate}
只有达到 L3，才应使用“核子胶子 TMD-PDF 的格点结果”这一表述。L0/L1 可以称为算符原型或矩阵元诊断；L2 是完整理论链的候选，但仍可能受单一参数点限制。

\section{必须报告的约定}

为保证可复核，论文和代码至少应明确：
\begin{itemize}
  \item 光锥坐标、度规、场强符号和生成元归一化；
  \item 规范链的起点、终点、方向和表示；
  \item \(\mu\)、\(\zeta\)、\(\Pz\)、\(\tau\) 与 Fourier 变量的定义；
  \item 横向指标的缩并和无迹投影；
  \item Wilson 圈长度、参考点 \(z_0,\boldsymbol b_{T,0}\) 和短距窗口；
  \item 单态混合的保留或忽略方式及误差；
  \item Dirac 算符、初值、预条件、求解器、重启策略、递归残差与真实残差；
  \item 所有外推函数、数据截断、bootstrap seed 和来源哈希。
\end{itemize}

\clearpage
\chapter{展望}

\section{大动量、细格距与算符混合}

胶子 TMD 的下一步不是单一算法优化，而是多尺度计算。高 \(P_z\) 会增加 boost 和对偶 Wilson 线的重叠，要求更小 \(a\) 和更大体积；细格距又会增加 GPU、传播子和胶子回路的成本。可行路线包括：
\begin{itemize}
  \item 用自重整化和 hybrid 方案把短距离非微扰匹配与长距离 Wilson 线分开；
  \item 用多个流时和多个格距确定短流时展开系数；
  \item 采用导数流、stout/HYP 和动量涂抹改善高动量信号，同时保持几何一致；
  \item 通过 GEVP、源分辨缓存和机器学习重采样降低统计误差；
  \item 把胶子-夸克单态混合作为矩阵问题，而不是事后乘一个标量因子。
\end{itemize}

\section{胶子 TMD 的自旋与小 \texorpdfstring{\(x\)}{x} 物理}

非极化 \(\fone\) 是最先应该闭合的基准。随后可以扩展到线性极化胶子、螺旋度胶子和 T-odd 相关分布。胶子自旋和轨道角动量与 proton spin sum rule、胶子 helicity PDF 和 EIC 可观测量有关\cite{Egerer2022,JiReview2021}。在小 \(x\) 区域，TMD 演化、饱和效应和软胶子重叠会使固定阶匹配面临大对数；这要求把 CS 核、BFKL/CCFM 或偶极图像与格点输入的关系说清楚，而不是用单一 PDF 拟合外推。

\section{与实验的闭环}

格点胶子 TMD 的价值在于提供非微扰先验、检验因子化和约束全局拟合。与 EIC 的协同可形成一个闭环：
\begin{enumerate}
  \item 格点计算给出特定 \((\mu,\zeta,b_T)\) 方案下的非微扰输入；
  \item 微扰 QCD 将输入演化到实验标度；
  \item Drell--Yan、SIDIS 和喷注/重味过程检验因子化与方案；
  \item 实验数据反过来约束格点外推和软/rapidity 模型。
\end{enumerate}
在没有完整方案前，把格点曲线与全局 PDF 直接叠加只能作为定性比较，不能作为因子化证明。

\section{建议的分阶段研究计划}

一个可审计的胶子 TMD 项目可以按以下顺序推进。每一阶段都以可复现证据作为进入下一阶段的条件。

\begin{figure}[htbp]
\centering
\begin{tikzpicture}[
  >=Stealth,font=\small,
  stage/.style={draw=blue!55!black,rounded corners,fill=blue!8,align=center,text width=2.45cm,minimum height=0.9cm,inner sep=3pt},
  gate/.style={draw=orange!70!black,rounded corners,fill=orange!10,align=center,text width=2.45cm,minimum height=0.75cm,inner sep=2pt}
]
  \node[stage] (s1) at (-4.8,2.0) {规范组态\\\(U_\mu(x)\)};
  \node[stage] (s2) at (-1.6,2.0) {物理源\\\(\eta\)、smearing};
  \node[stage] (s3) at (1.6,2.0) {Dirac 求逆\\\(D_q[U]\psi=\eta\)};
  \node[stage] (s4) at (4.8,2.0) {两点/三点\\关联函数};
  \node[stage] (s5) at (4.8,-1.25) {重正化\\Wilson 圈/flow};
  \node[stage] (s6) at (1.6,-1.25) {soft/rapidity\\\(S_I^g,K_g\)};
  \node[stage] (s7) at (-1.6,-1.25) {硬匹配\\\(C_{gg},C_{gq}\)};
  \node[gate] (s8) at (-4.8,-1.25) {极限与外推\\L3 物理结果};
  \draw[->,line width=1pt] (s1)--(s2); \draw[->,line width=1pt] (s2)--(s3); \draw[->,line width=1pt] (s3)--(s4);
  \draw[->,line width=1pt] (s4)--(s5); \draw[->,line width=1pt] (s5)--(s6); \draw[->,line width=1pt] (s6)--(s7); \draw[->,line width=1pt] (s7)--(s8);
  \draw[->,dashed,red!65!black] (s3.south) -- ++(0,-0.55) node[below,align=center,font=\scriptsize]{双残差/容差扫描};
  \draw[->,dashed,teal!70!black] (s5.north) -- ++(0,0.45) node[above,align=center,font=\scriptsize]{方案一致性};
  \draw[->,dashed,orange!75!black] (s8.north) -- ++(0,0.5) node[above,align=center,font=\scriptsize]{失败关闭};
\end{tikzpicture}
\caption{从规范组态到 L3 胶子 TMD 结果的图形化验收链。算法求解、矩阵元、重正化、soft/rapidity、匹配和物理极限是串联门，不是可以独立宣称完成的平行任务；任一门的证据缺失都应回到对应阶段。}
\label{fig:acceptance-chain}
\end{figure}

图~\ref{fig:acceptance-chain} 把后文步骤组织成一条串行验收链。生产规模扩大不能越过失败门；相反，只有前一阶段的状态、方案和误差证据都可复现，才有资格进入下一阶段。
\begin{enumerate}
  \item \textbf{算符与几何阶段：}固定场强符号、生成元、伴随链接、staple 方向、横向投影和 Fourier 约定；建立小矩阵/自由场/恒等链接测试。
  \item \textbf{两点与三点原型：}在单系综上验证 source-resolved 两点函数、三点比值、真空减除、色散关系和激发态拟合；先以 \(b_T=0\) 或小 \(b_T\) 做闭环。
  \item \textbf{UV 与流窗口：}测量同几何 Wilson 圈，建立 ratio/hybrid/self-renorm 的交叉检验；扫描多个 \(\tau\)，只保留 \(a\ll r_f\ll b_T,z\) 的点。
  \item \textbf{soft/rapidity 阶段：}在同一表示和几何下计算 \(S_I^g\) 与 \(K_g\)，用多个 \(P_z\) 验证平台；与微扰小 \(b_T\) 极限对接。
  \item \textbf{匹配与混合阶段：}接入胶子硬核、\(g-q\) 混合矩阵和 \(\MSbar\) 转换；报告匹配阶数、标度区间和高阶误差。
  \item \textbf{多参数生产：}扩展到多个 \(a\)、多个 \(L_s\)、至少一个物理或可控手征质量；进行 bootstrap 联合分析。
  \item \textbf{极限与外推：}按第 7 节的窗口做 \(a\to0,\tau\to0,P_z\to\infty\) 联合外推，并给出每个 \(x,b_T\) 点的可信区。
  \item \textbf{发表与复现：}公开算符契约、缓存哈希、拟合窗口、协方差、失败 replica、Fourier 约定和来源 manifest。
\end{enumerate}

\section{开放问题清单}

按照物理重要性排序，当前最需要解决的问题包括：
\begin{enumerate}
  \item 胶子表示下的 soft/rapidity 因子如何与现有夸克 TMD 方案统一，并在多个系综上验证；
  \item 胶子 TMD 的单态混合矩阵能否在非微扰层面完整确定，而不是依赖 NLO 匹配和全局 PDF 拟合；
  \item 大 \(b_T\) 非微扰区的 Fourier 尾部如何与 CS 核、内禀软函数和格点体积联合约束；
  \item 高动量下 \(a^2P_z^2\)、流化半径和 Wilson line overlap 的联合控制；
  \item 线性极化胶子 TMD 的投影、符号和 T-odd 结构与 \(f_1^g\) 的分离；
  \item 小 \(x\) 区域能否在固定阶 LaMET 之外，与 BFKL/饱和物理建立可检验的桥接；
  \item 胶子 TMD 的格点先验如何进入 EIC 全局分析，而不引入重复计数或方案不一致；
  \item 以“失败关闭”为默认的报告标准如何推广到所有 TMD 计算，避免只发表通过质量门的子集。
\end{enumerate}

\clearpage
\chapter{结论}

使用格点 QCD 计算核子胶子 TMD-PDF 的核心并不是把一个 PDF 公式“翻译成 Euclidean 算符”，而是建立一条多尺度、方案一致的理论链：
\begin{equation}
  \text{光锥 TMD}
  \xleftarrow{\text{soft/rapidity+CS+匹配}}
  \text{准 TMD}
  \xleftarrow{\text{Wilson 圈+重正化+flow}}
  \text{格点矩阵元}.
  \label{eq:chain}
\end{equation}
在这条链上，胶子场强、伴随规范链、Wilson 圈减除、CS 核、硬匹配、单态混合和大动量极限缺一不可。梯度流和短流时展开可以显著改善 UV 控制，但它们只解决链中的一个尺度问题；rapidity、表示相关软因子和连续极限必须独立验证。

最底层的格点矩阵元也不是“由代码直接给出”的黑箱。每个夸克传播子都是固定规范背景中 \(D_q[U]\psi_\eta=\eta\) 的解；BiCGSTAB 的搜索方向、shadow 残差和局部稳定化系数只是构造这个逆矩阵列的方式。近零模和条件数会放大源、舍入与停止容差的影响，非连通胶子拓扑又会把求解偏差传播到真空减除和统计相关中。因此，合格的 TMD 报告不仅应说明用了哪一种 Krylov 算法，还应证明递归残差与真实残差、容差扫描、预条件选择和目标矩阵元对求解误差均处于统计误差之内。传播子求逆是物理链的一部分，不是可以隐藏在性能附录中的数值细节。

结合库内文献与实现审计，本文的结论是：共线胶子 PDF、TMD 软函数和 CS 核已经有重要格点先例；夸克 TMD 提供了完整算法模板；但本文所核查的库内证据尚不足以支持经过完整 soft/rapidity、混合、多格距和多流时验收的直接核子胶子 TMD-PDF 数值结果。现阶段最稳健的表述是“理论框架已建立，关键构件已有格点验证，直接胶子 TMD 物理计算仍在闭合中”。这一证据边界本身就是后续计算的出发点。

\appendix

\clearpage
\chapter{符号与约定速查}

\begin{longtable}{L{2.3cm}L{4.8cm}L{6.0cm}}
\toprule
符号 & 定义 & 备注 \\
\midrule
\endfirsthead
\toprule
符号 & 定义 & 备注 \\
\midrule
\endhead
\(F_{\mu\nu}\) & \(F=(i/g)[D_\mu,D_\nu]\)，\(D_\mu=\partial_\mu-igA_\mu\) & 与正文生成元约定一致 \\
\(T^a\) & \([T^a,T^b]=if^{abc}T^c\) & 伴随链与基本链的转换须显式验证 \\
\(\bT\) & 横向分离 & \(b_T=|\bT|\) \\
\(\zeta\) & rapidity 重正化标度 & 与 \(\zeta_z\sim(2xP_z)^2\) 区分 \\
\(K_R\) & 表示 \(R\) 的 Collins--Soper 核 & 胶子使用伴随表示 \\
\(S_{I,R}\) & 内禀软函数 & 与同方案 \(K_R\) 配套 \\
\(r_f\) & \(\sqrt{8\tau}\) & 流半径，不是格距 \\
\(L_\parallel\) & staple 纵向臂长 & 需报告平台/外推 \\
\bottomrule
\end{longtable}

\clearpage
\chapter{关键推导补充}

\section{从格点链接到 Clover 场强}

沿 \(\mu,\nu\) 平面的基本 plaquette 为
\begin{equation}
  U_{\mu\nu}(x)
  =U_\mu(x)U_\nu(x+\hat\mu)
  U_\mu^\dagger(x+\hat\nu)U_\nu^\dagger(x).
  \label{eq:plaquette}
\end{equation}
把每个链接用 \(U_\mu(x)=e^{iagA_\mu(x+\hat\mu/2)}\) 展开，Baker--Campbell--Hausdorff 公式给出
\begin{equation}
  U_{\mu\nu}(x)
  =\exp\left[ia^2F_{\mu\nu}(x)+\mathcal O(a^3)\right],
  \label{eq:plaquettebch}
\end{equation}
其中 \(F_{\mu\nu}\) 按正文约定包含非阿贝尔对易项。把四个方向的 plaquette 相加：
\begin{equation}
  Q_{\mu\nu}
  =U_{\mu\nu}+U_{-\nu,\mu}+U_{\nu,-\mu}+U_{-\mu,-\nu},
  \label{eq:cloversum}
\end{equation}
再取反对称组合
\begin{equation}
  F_{\mu\nu}^{\rm clover}
  =-\frac{i}{8a^2}\left(Q_{\mu\nu}-Q_{\nu\mu}\right)
  =\frac{1}{2}\left(F_{\mu\nu}-F_{\nu\mu}\right)+\mathcal O(a^2).
  \label{eq:cloverderive}
\end{equation}
对基本表示，连续场强无迹；有限格距下应做 \(F\to F-\frac1{N_c}\Tr F\,\mathbf 1\)。该投影必须在计算 Wilson 线收缩之前完成，否则迹部分会与颜色指标缩并混合。

\section{rapidity 发散的玩具积分}

为了看出 rapidity 发散不是普通 UV 极点，考虑一个简化的光锥积分
\begin{equation}
  I=\int\dd k^+\dd k^-\,
  \frac{f(k^+k^-)}{(k^+k^-)^{1+\epsilon}}.
  \label{eq:rapiditytoy}
\end{equation}
引入 \(m^2=k^+k^-\) 和快度变量 \(y=k^+/k^-\)，有
\begin{equation}
  I=\frac12\int\frac{\dd y}{y}
  \int\frac{\dd m^2\,f(m^2)}{m^{2+2\epsilon}}.
  \label{eq:rapiditytoy2}
\end{equation}
UV 维数正规化只控制 \(m^2\to\infty\)，而 \(y\to0\) 或 \(y\to\infty\) 的端点仍没有尺度截断。这正是光样 Wilson 线之间的软胶子交换产生的 rapidity 发散。TMD 软函数的作用是为这两个端点提供统一的快度减除，并使物理组合对快度方案不依赖。

\section{Collins--Soper 方程的解}

将 CS 方程写成
\begin{equation}
  \frac{\partial}{\partial\ln\zeta}
  \ln f(\zeta)=\frac12K(\mu),
  \label{eq:cslog}
\end{equation}
在固定 \(\mu\) 下积分得到
\begin{equation}
  f(\zeta)=f(\zeta_0)
  \exp\left[\frac12K(\mu)\ln\frac{\zeta}{\zeta_0}\right].
  \label{eq:cs-solution}
\end{equation}
另一方面，\(K\) 的 UV 演化满足
\begin{equation}
  \mu^2\frac{\partial K}{\partial\mu^2}
  =-\Gamma_{\rm cusp}(\alpha_s),
  \label{eq:k-rge-detail}
\end{equation}
因此可把非微扰信息放在参考标度 \(1/b_T\)：
\begin{equation}
  K(b_T;\mu)
  =K(b_T;1/b_T)
  -2\int_{1/b_T}^{\mu}\frac{\dd\mu'}{\mu'}
  \Gamma_{\rm cusp}(\alpha_s(\mu')).
  \label{eq:k-solution}
\end{equation}
对胶子，领头 cusp 系数与 \(C_A\) 成正比。这个表示清楚说明为什么 CS 核不能只看成一条经验曲线：它同时包含微扰 cusp 演化和大 \(b_T\) 的非微扰初值。

\section{Fourier 变换与小 \texorpdfstring{\(b_T\)}{bT} OPE}

坐标空间 TMD 和动量空间 TMD 的关系为
\begin{equation}
  f(x,\bT;\mu,\zeta)
  =\int\frac{\dd^2\kT}{(2\pi)^2}
  e^{i\kT\cdot\bT}f(x,\kT;\mu,\zeta).
  \label{eq:fourierTMD}
\end{equation}
在小 \(b_T\) 区域，可以把非局域算符做短距 OPE。对胶子，最低阶算符是 twist-2 场强双线性，其矩阵元给出胶子 PDF；更高维算符给出 \(b_T^2\) 幂修正。因此
\begin{equation}
  f_1^g(x,\bT;\mu,\zeta)
  =f_g(x,\mu)
  +b_T^2\,\mathcal C_2(x;\mu,\zeta)
  +\mathcal O(b_T^4\Lambda_{\rm QCD}^2),
  \label{eq:smallb}
\end{equation}
其中 \(\mathcal C_2\) 概括 twist-4 和方案相关贡献。该展开的有效范围是 \(b_T\Lambda_{\rm QCD}\ll1\)。在大 \(b_T\)，Wilson 线的非微扰势和软胶子使该展开失效，必须用格点数据或非微扰模型。

\section{梯度流热核与规范作用量}

在自由极限下，流方程对横向模式给出
\begin{equation}
  \partial_\tau\widetilde B_\mu(\tau,k)
  =-k^2P^\perp_{\mu\nu}\widetilde B_\nu(\tau,k).
  \label{eq:heatflow}
\end{equation}
解为
\begin{equation}
  \widetilde B_\mu(\tau,k)
  =P^\perp_{\mu\nu}e^{-k^2\tau}\widetilde B_\nu(0,k)
  +P^\parallel_{\mu\nu}\widetilde B_\nu(0,k).
  \label{eq:heatflowsolution}
\end{equation}
因此流半径 \(r_f=\sqrt{8\tau}\) 是高斯平滑宽度。规范作用量密度的领头短流时展开可写成
\begin{equation}
  E(\tau,x)
  =\frac14G^a_{\mu\nu}G^a_{\mu\nu}
  =\langle E\rangle
  +c_E(\tau)\left[\frac14F^a_{\mu\nu}F^a_{\mu\nu}\right]_R
  +\mathcal O(\tau),
  \label{eq:energyflow}
\end{equation}
其中 \(c_E\) 的 \(\tau\to0\) 对数行为由 \(\beta\) 函数控制。对非局域胶子 TMD，类似展开必须保留 \(z,\bT\)，并在固定外部分离后逐项匹配。

\section{Wilson 圈抵消的结构}

设开 Wilson 线矩阵元的自能和 pinch 结构近似为
\begin{equation}
  h_{\rm bare}(z,\bT;L)
  \sim
  \exp\left[-\delta m\,\frac{2L+|z|}{a}
  -V(\bT)L\right]\,
  h_{\rm phys}(z,\bT),
  \label{eq:barewilson}
\end{equation}
同几何矩形 Wilson 圈满足
\begin{equation}
  Z_E(\bT;2L)
  \sim
  \exp\left[-\delta m\,\frac{4L}{a}-2V(\bT)L\right]
  Z_{\rm finite}.
  \label{eq:zwilson}
\end{equation}
当路径长度和转角匹配时，\(\sqrt{Z_E}\) 可抵消开 Wilson 线中的线性项和 pinch 项，留下端点对数与物理矩阵元。这里的近似式用来说明结构，不应替代实际算符的严格重正化；特别是胶子伴随表示、转角角度和横向端点必须逐项核对。

\clearpage
\chapter{本库证据映射}

本文的文献和实现审阅覆盖以下库内入口：
\begin{itemize}
  \item \path{docs/pyqcd_physics/README.md}：当前物理导航和证据层级；
  \item \path{docs/pyqcd_reference/}：从固定提交吸收的格点 QCD、胶子 PDF/TMD、重整化和梯度流文档；
  \item \path{refer/papers/INDEX.md}：50 篇原文转排、双语目录和来源 arXiv 标识；
  \item 以下四个原文目录提供胶子算符、软函数和 CS 核原文：\\
  \path{refer/papers/Accessing_Gluon_PDF_LaMET_latex/}，\\
  \path{refer/papers/Gluon_Pseudo_Short_Distances_latex/}，\\
  \path{refer/papers/CS_Kernel_Quasi_TMDWF_latex/}，\\
  \path{refer/papers/TMD_Soft_Function_LaMET_latex/}；
  \item \path{refer/git-rep/PyQCD/pyqcd/operator/}、\path{renorm/}、\path{pipeline/}：场强、staple、flow、TMD 和匹配实现；
  \item \path{myqcd/formula_registry.py}、\path{pyqcd_theory_audit.py}：代数、量纲和有限维审计边界；
  \item \path{docs/pyqcd_reference/absorb_dev20_gluon_tmd_audit_20261022.md}：近期真实组态 smoke、fail-closed 状态和仍缺口。
\end{itemize}

\section{来源类别与用途}

\begin{longtable}{L{3.0cm}L{4.0cm}L{6.3cm}}
\toprule
来源类别 & 典型路径 & 本文使用方式 \\
\midrule
\endfirsthead
\toprule
来源类别 & 典型路径 & 本文使用方式 \\
\midrule
\endhead
上游源码 & \path{refer/git-rep/PyQCD/pyqcd/operator/} & 场强、staple、直线 Wilson 线和胶子 OPE 的接口语义 \\
上游源码 & \path{refer/git-rep/PyQCD/pyqcd/renorm/} & hybrid、self-renorm、flow、TMD 和 matching 的公式组织 \\
上游流水线 & \path{refer/git-rep/PyQCD/pyqcd/pipeline/} & 两点/三点、flow、TMD ratio、缓存和状态门 \\
原文转排 & \path{refer/papers/} & 原始公式、参数、因子化与文献边界 \\
书籍/教材 & \path{refer/books/} & Wilson 线、关联函数、统计和格点 QFT 背景 \\
吸收快照 & \path{docs/pyqcd_reference/} & 胶子 PDF/TMD 推导、代码解析、历史报告与证据快照 \\
公式审计 & \path{myqcd/} & SymPy 恒等式、量纲、有限维代理和来源行号绑定 \\
运行证据 & \path{logs/}、\path{data/runs/} & 真实 smoke、统计拟合、失败关闭和缓存哈希 \\
\bottomrule
\end{longtable}

\section{最小可复现报告清单}

\begin{table}[htbp]
\centering
\caption{胶子 TMD 数值结果的复现清单。}
\label{tab:repro-checklist}
\small
\begin{tabularx}{\textwidth}{L{3.3cm}Y}
\toprule
项目 & 必须记录的内容 \\
\midrule
规范与算符 & 场强约定、生成元、链接方向、伴随/基本表示、staple 路径和横向投影 \\
系综 & \(\beta,a,L_s,T_s,N_{\rm conf}\)、海/价夸克质量、作用量、边界条件、源数 \\
动量与流 & \(P_z\) 网格、核子色散关系、\(\tau\) 网格、\(r_f\)、\(P_zr_f\) \\
算符几何 & \(z/a\)、\(b_T/a\)、\(L_\parallel/a\)、路径端点、HYP/stout/flow 参数 \\
UV 重正化 & Wilson 圈、ratio/hybrid/self-renorm、短距点、\(\delta m+m_0\)、\(\mu\) \\
rapidity/soft & \(S_{I,g}\)、\(K_g\)、表示、方案定义、\(P_1/P_2\) 和 soft 提取方法 \\
匹配 & \(C_{gg},C_{gq}\)、\(\alpha_s\) 阶数、标度、RGR、是否忽略混合及误差 \\
统计 & 源分辨策略、binning、bootstrap seed、replica 数、协方差和 fallback 率 \\
谱学 & \(t_{\rm sep}\) 范围、拟合态数、GEVP 基、残差和模型平均 \\
Dirac 求逆 & \(D\) 的作用量与边界条件、\(\gamma_5\)-Hermiticity、初值、预条件、\(\epsilon_{\rm rel}\)、\(i_{\max}\)、重启、递归/真实残差和求解器交叉检查 \\
Fourier & 相位约定、\(z_{\max}\)、尾部模型、\(x\) 支撑、端点和正定性处理 \\
外推 & \(a\to0\)、\(m_\pi\to m_\pi^{\rm phys}\)、\(P_z\to\infty\)、\(\tau\to0\) 的函数和窗口 \\
状态 & 每个 \(x,b_T\) 点的可信区、系统误差、失败点和最终 evidence level \\
\bottomrule
\end{tabularx}
\end{table}

\clearpage
\chapter{数据与代码可用性}

本文不产生新的蒙卡系综或 GPU 数值。所引用的公式、文献转排、实现路径和审计日志均位于本库：
\begin{itemize}
  \item 文献来源索引：\texttt{refer/papers/INDEX.md}；
  \item 上游 PyQCD 固定提交：\texttt{refer/git-rep/PyQCD/}，当前提交标识由库内来源表记录；
  \item 公式审计：\texttt{python -m myqcd} 与 \texttt{myqcd/formula\_registry.py}；
  \item 历史吸收与质量审计：\texttt{docs/pyqcd\_reference/} 与 \texttt{logs/}；
  \item 真实数据、缓存和 build 产物：\texttt{data/}，按仓库规则不入库。
\end{itemize}
任何后续数值结果都应记录系综、规范组态哈希、格距、体积、价/海夸克质量、\(P_z\)、\(\tau\)、\(L_\parallel\)、\(b_T\)、重正化方案、soft/rapidity 方案、匹配阶数和重采样 seed。

\clearpage
\chapter{伦理、贡献、利益冲突与资助}

\noindent\textbf{伦理声明：}本文为理论综述，不涉及人类受试者、动物实验或个人隐私数据。

\noindent\textbf{作者贡献：}MyQCD 理论物理研究组负责文献范围界定、物理论证、公式约定统一、证据分层和文稿整理；任何真实系综或数值结果的使用均应另行署名和审计。

\noindent\textbf{利益冲突：}作者声明不存在与本文主题相关的利益冲突。

\noindent\textbf{资助：}本文整理工作未声明新的外部资助；库内原始文献和合作组结果的资助信息以其原文为准。

\noindent\textbf{AI 使用说明：}本文在 MyQCD 文献证据和本地源码约束下使用自动化研究写作工具辅助整理、交叉核对和 LaTeX 排版。所有物理结论、文献标识和验证状态均回溯到库内来源；未执行的真实系综、GPU、soft/rapidity 或连续极限验证没有作为结果陈述。

\clearpage
\chapter{约定与公式复核}

\section{链接、场强与 Clover}

\begin{longtable}{L{4.0cm}L{4.0cm}L{5.8cm}}
\caption{合稿时必须固定的约定及其检查。}\label{tab:conventions}\\
\toprule
对象 & 本节约定 & 替换约定时的检查\\
\midrule
\endfirsthead
\toprule
对象 & 本节约定 & 替换约定时的检查\\
\midrule
\endhead
生成元 & \(\Tr(T^aT^b)=\delta^{ab}/2\) & 若取 \(\Tr=\delta^{ab}\)，场强迹归一化和 \(C_F\) 同步改变\\
链接 & \(U_\mu=e^{iag_0A_\mu}\) & 若吸收 \(g_0\)，删去式~\eqref{eq:clover-field} 的 \(g_0\)\\
plaquette & \(P_{\mu\nu}=e^{ia^2g_0F_{\mu\nu}+O(a^3)}\) & 用自由场小 \(a\) 检查符号与指数因子\\
Clover 和 & 四个定向 plaquette 均给出同一 \(F_{\mu\nu}\) & 改变定向时重新检查 \(Q-Q^\dagger\) 的因子 8\\
Clover 系数 & \(c_{\rm SW}\) 与式~\eqref{eq:b01-clover-action} 绑定 & 若链接除以 \(u_0\)，使用 \(c_{\rm SW}=u_0^{-3}\) 的对应归一化\\
Wilson 圈 & 颜色迹按 \(N_c\) 归一化 & 不归一化会改变自能常数 \(A\)，但不改变势差和斜率\\
Fourier & 明确 \(e^{\pm iP_zz}\) 与 \(1/(2\pi)\) 位置 & \(z\to-z\)、复共轭和正负 \(x\) 同时变换\\
\bottomrule
\end{longtable}

\section{色散与 Wilson 费米子}

自由标量色散由式~\eqref{eq:lattice-dispersion} 给出。其小 \(a\) 展开的首个非连续项是 \(\Ord(a^2)\)；若采用改进导数或不同质量归一化，应重新展开传播子。Wilson 项中的 \(r\) 与 Clover 的 \(c_{\rm SW}\) 是两个独立参数：前者控制加倍模质量，后者抵消物理分支的 \(\Ord(a)\) 项。代码中把二者绑定为一个参数会改变手征极限或产生非物理的 doubler 残留。

HQET 和 LaMET 的类比只保留展开结构，不保留相同自由度。Symanzik 的 \(a\) 进入作用量与算符基，HQET 的 \(1/m_Q\) 进入重夸克有效拉氏量，LaMET 的 \(1/P_z\) 进入矩阵元匹配。三者都要求匹配核与目标量处于同一方案；跨框架借用数值前必须转换生成元、耦合和算符归一化。

\chapter{失败模式与最小诊断表}

\begin{longtable}{L{2.7cm}L{3.9cm}L{3.7cm}L{3.0cm}}
\caption{常见失败模式、可能原因与最小复核。}\label{tab:failures}\\
\toprule
现象 & 可能原因 & 最小复核 & 禁止的结论\\
\midrule
\endfirsthead
\toprule
现象 & 可能原因 & 最小复核 & 禁止的结论\\
\midrule
\endhead
Wilson 线结果随总长度指数衰减 & 未减除线性自能 & 同长度 Wilson 圈平方根；多格距斜率 & 把衰减解释为物理 PDF 尾部\\
RI/MOM 后小 \(a\) 或大 \(z\) 更不稳定 & 非局域线性项、cusp/pinch、端点混合 & 检查同几何减除、\(L\) 平台和表示 & 宣称 RI/MOM 已完全重正化\\
色散拟合斜率为 1 但结果仍依赖 \(P_z\) & 激发态、端点或匹配幂修正 & 多源/GEVP、不同 \(aP_z\) 序列 & 把单点斜率当全部误差受控\\
HMC 接受率正常但拓扑扇区不混合 & 轨迹短、步长/算法或拓扑冻结 & 拓扑历史、不同扇区、自相关 & 用接受率替代采样充分性\\
Clover 改进后仍有线性 \(a\) 趋势 & 系数/链接归一化不一致、算符未改进 & 自由场、Ward 恒等式、\(c_{\rm SW}\) 约定 & 称为纯统计波动\\
不同 \(u_0\) 或不同 Clover 符号给出不同连续极限 & 作用量与算符方案混用 & 固定生成元、\(U/u_0\) 和 \(F\) 归一化 & 混合方案中心值\\
自重整化结果依赖 \(z_0\) & 单格距不可辨识、端点对数未处理 & 多格距联合、不同短距点和 SDR/hybrid 交叉 & 宣称已经消除所有重整子\\
流化后短距结构被抹平 & \(r_f\) 接近或大于 \(z,b_T,\Lambda^{-1}\) & 扫描多个流时并检查重叠区 & 把流化当成万能量减除\\
\bottomrule
\end{longtable}

\section{证据等级}

\begin{enumerate}
  \item \textbf{定义级：}公式、算符和约定可由代数及规范变换直接检验。
  \item \textbf{自由场级：}链接、场强、传播子、加倍模和 Clover 系数可在恒等链接上复核。
  \item \textbf{算法级：}采样、残差、自相关和极限扫描给出可重复日志。
  \item \textbf{发表数值级：}系综参数与物理数来自来源报告；未重跑时只称“来源报告值”。
  \item \textbf{物理论述级：}只有定义、算法、重正化、匹配和极限门均通过，才允许把结果称为目标物理量。
\end{enumerate}

\noindent\textbf{结论.} 格点 QCD 的精度不是由某一个改进项或某一个求解器单独决定，而由以下链条共同决定：规范定义与场强归一化、离散化和对称性、标度设定、采样与谱学、非定域算符的 UV/rapidity 结构、连续与无穷动量极限。任何一环缺失，后续数值都只能作为诊断或中间结果。

\chapter{Wick 图拓扑、蒸馏与验证清单}\label{sec:b02-appendix}
本附录把正文用到的“费曼图”还原为 Wick 收缩拓扑。这里的线不是微扰展开中的虚粒子传播线，而是格点组态上的夸克传播子、闭合费米子环和 Wilson 线。图的有用之处在于保证算符、交换符号、闭合环和真空减除没有被遗漏；图本身不能替代代数式。

\section{连接、断开与海夸克}\label{subsec:b02-topology}

连接图在源算符与汇算符之间存在连续的费米子线，或至少存在把两者联系起来的收缩。断开图含有一个不参与该源汇连接的闭合费米子环，但它仍通过规范组态平均与其余部分相关联。这个定义是收缩层面的，不应与规范场的连通性、微扰 QCD 图或“是否出现在淬火近似中”混为一谈。

full QCD 的组态权重为
\begin{equation}
  Z=\int\mathcal DU\,e^{-S_G[U]}\prod_f\det D_f[U],
  \label{eq:b02-qcd-weight}
\end{equation}
其中行列式描述海夸克对规范场的反馈。淬火近似令
\begin{equation}
  \det D_f[U]\longrightarrow1,
  \label{eq:b02-quench}
\end{equation}
即关闭这种反馈。它是组态分布层面的近似，不是“把所有含闭合环的测量项手动设为零”的同义词。淬火数据中可以形式上计算某些闭合收缩，但该理论缺乏完整 QCD 的海夸克动力学和幺正解释。现代物理论文若仍使用部分淬火设置，应明确区分 \(m_{\rm sea}\) 与 \(m_{\rm val}\)，并说明手征外推方案。

对无法直接存储全部对角元 \(S_q(x,x)\) 的闭合环，可用随机源估计
\begin{equation}
  \operatorname{Tr}\!\left[\Gamma S_q\right]
  \simeq
  \frac1{N_r}\sum_{r=1}^{N_r}
  \eta_r^\dagger\Gamma S_q\eta_r,
  \qquad
  \left\langle\eta_{r,i}^*\eta_{r',j}\right\rangle
  =\delta_{rr'}\delta_{ij}.
  \label{eq:b02-stochastic-trace}
\end{equation}
迹包括颜色、Dirac 和空间指标。$Z_2$、$U(1)$、时空稀释和低模 deflation 只能降低方差；它们不改变估计的无偏性，也不消除源、时间片和组态之间的相关性。报告应给出随机源类型、稀释方式、求解容差和逐项方差，而不只是总 CPU/GPU 时间。

\begin{longtable}{>{\raggedright\arraybackslash}p{2.4cm}>{\raggedright\arraybackslash}p{3.5cm}>{\raggedright\arraybackslash}p{3.5cm}>{\raggedright\arraybackslash}p{3.8cm}}
\caption{B02 常见 Wick 拓扑及其审计点。}\label{tab:b02-topology-atlas}\\
\toprule
拓扑 & 代表对象 & 必须保留的结构 & 主要失败模式 \\
\midrule
\endfirsthead
\multicolumn{4}{l}{\small 续表：B02 常见 Wick 拓扑及其审计点。}\\
\toprule
拓扑 & 代表对象 & 必须保留的结构 & 主要失败模式 \\
\midrule
\endhead
\bottomrule
\endfoot
连接介子 & \(\pi^+\) 两点或电磁三点 & 两条味线的闭合迹、Gamma 顺序、交换符号 & 把同位旋矢量断开项或核子交换项漏掉 \\
连接重子 & 核子两点/三点 & 三条夸克线、\(\epsilon^{abc}\)、交换项 & 只实现一种线配对、自旋投影错位 \\
单态断开 & 味单态介子、奇异 PDF & 独立闭合环、味系数、随机迹方差 & 只做全局真空减除、忽略源/时间片结构 \\
胶子插入 & 胶子 PDF、胶子 TMD & 场强与 Wilson 线同一几何、真空减除 & 链路方向、指标顺序或表示归一化错误 \\
直接 Wilson 线 & 准 PDF & 端点、路径顺序、线性发散 & 端点链接反向、重整化几何不一致 \\
Staple Wilson 线 & 准 TMD & 三段路径、cusp/pinch-pole 结构、软函数因子 & 把 Wilson 圈减除当完整 TMD 重正化 \\
四夸克形状因子 & 介子软函数 & 两类双线型、横向分离、连通投影 & Fierz 基或波函数归一化混用 \\
蒸馏收缩 & perambulator/广义 perambulator & 源汇投影、流顶点、低模截断 & 视 \(VV^\dagger\) 为酉、缓存键缺项 \\
\end{longtable}

\section{Fourier、准分布与相位约定}\label{subsec:b02-fourier}

空间动量投影定义为
\begin{equation}
  O(\mathbf p,t)
  =\sum_{\mathbf x}e^{-i\mathbf p\cdot\mathbf x}O(\mathbf x,t),
  \qquad
  \mathbf p=\frac{2\pi}{L}\mathbf n.
  \label{eq:b02-momentum-projection}
\end{equation}
对有限区间上的准分布，本书采用
\begin{equation}
  \widetilde h_\Gamma(z,P^z)
  =\left\langle P^z\right|\bar q(z)\Gamma W(z,0)q(0)\left|P^z\right\rangle ,
  \label{eq:b02-quasi-correlator}
\end{equation}
\begin{equation}
  \widetilde q_\Gamma(x,P^z)
  =\mathcal N_\Gamma^{-1}
  \int\frac{\mathrm dz}{2\pi}
  e^{+ixP^zz}
  \widetilde h_\Gamma^R(z,P^z),
  \label{eq:b02-quasi-fourier}
\end{equation}
其中 \(\mathcal N_\Gamma\) 由算符和端点归一化固定。某些文献使用 \(e^{-ixP^zz}\)；这等价于把 \(x\) 换成 \(-x\)，或在取复共轭时改变相位。比较不同工作前必须统一这一约定。有限区间截断、离散 \(z\)、大 \(zP^z\) 噪声和尾部模型都属于 Fourier 逆问题的一部分；把截断数据直接做离散变换得到的是带窗函数的对象。

准 PDF 的大距离尾部建模应满足量纲、端点和重正化约束。一个常用的再生模型是
\begin{equation}
  \widetilde h^R(z,P^z)
  \sim c_1\lambda^{-d_1}e^{-\lambda/\lambda_0},
  \qquad
  \lambda=zP^z,
  \label{eq:b02-tail-model}
\end{equation}
但参数 \(c_1,d_1,\lambda_0\) 与拟合窗口、重正化方案和更高扭度修正相关。尾部外推应做模型平均或系统误差扫描，不能把外推区域当作直接测量平台。

\section{验证矩阵与可验证断言}\label{subsec:b02-validation}

\begin{longtable}{>{\raggedright\arraybackslash}p{2.5cm}>{\raggedright\arraybackslash}p{3.8cm}>{\raggedright\arraybackslash}p{3.8cm}>{\raggedright\arraybackslash}p{3.1cm}}
\caption{B02 最小验证矩阵。}\label{tab:b02-validation-matrix}\\
\toprule
对象 & 不变量或极限 & 失败信号 & 证据应如何记录 \\
\midrule
\endfirsthead
\multicolumn{4}{l}{\small 续表：B02 最小验证矩阵。}\\
\toprule
对象 & 不变量或极限 & 失败信号 & 证据应如何记录 \\
\midrule
\endhead
\bottomrule
\endfoot
两点谱 & 正谱、复共轭、周期 \(\cosh\) 对称 & 负有效质量、端点反号、平台随源汇消失 & \(t_{\min},t_{\max}\)、协方差、\(\chi^2\)、源汇类型 \\
Wick 收缩 & 每个算符的交换符号、同位旋极限 & 连通/断开项漏减或重复 & 小矩阵、自由场恒等式、拓扑清单 \\
三点比值 & \(t_s\to\infty\) 与 \(0\ll\tau\ll t_s\) 的稳定平台 & 接触项、两端污染不对称 & 两端独立系数、窗口扫描、replica 传播 \\
顺序源 & Dirac 方程真残差与汇点相位 & 动量混叠、汇投影错误 & 每源/动量的 matvec 数和真残差 \\
GEVP & \(B\)-正交归一、Cholesky 残差、态追踪 & 条件数爆炸、排序跳变、向量归一化漂移 & \(t_0\)、基大小、条件数、本征值与本征向量 \\
蒸馏 & \(N_v\to N_{\rm full}\) 回到直接传播子 & 低模截断伪平台、缓存串味 & 增加 \(N_v\) 扫描、基对照、缓存键和哈希 \\
Lüscher & 自由极限回到无相互作用量子化 & 分支选错、阈值外仍用单道式 & \(q\) 范围、分支、体积列表、精度 \\
形状因子 & \(Q^2\to0\) 的守恒荷与电荷半径 & 端点外推驱动半径、流重整化错误 & 流 \(Z\)、动量转移、连续/手征极限 \\
Fourier & \(z\to-z\) 的厄米性、归一化与端点 & 相位反号、窗函数未报、尾部主导 & 相位约定、窗函数、尾部模型与系统误差 \\
\end{longtable}

以下断言可由本批次资料直接转化为测试，而不需要真实系综：
\begin{enumerate}
  \item 对任意有限个正能态，两点谱必须为正、实、指数衰减；周期边界下的自共轭通道满足式~\eqref{eq:b02-cosh}。
  \item 同位旋矢量介子的连接收缩在精确 \(SU(2)\) 极限下不含双闭合环断开项；味单态组合一般含有。
  \item 当算符基与态空间完全匹配时，GEVP 本征值精确为式~\eqref{eq:b02-gevp-exact}；加入未覆盖高态后，偏差由 \(\delta_n\) 度量。
  \item Cholesky 化归、显式逆矩阵和标准 Hermitian 求解器在小矩阵上必须给出相同本征值；病态矩阵上应记录数值路径差异，而不是只比较默认容差。
  \item \(B\)-正交归一本征向量在欧氏归一化后一般不再满足 \(V^\dagger C(t_0)V=I\)；投影三点函数前必须显式修正。
  \item 蒸馏投影满足 \(\Box^2=\Box\)；当 \(N_v\) 等于完整基维数时，perambulator 收缩必须回到直接传播子结果。
  \item 自由场 \(Z_{00}\) 极限或已知散射长度玩具模型中，Lüscher 映射必须复现指定相移分支；改变 \(\delta\) 符号约定时整条映射应一致改号。
  \item 形状因子在零动量转移处必须回到对应守恒荷；\(F_\pi/f_\pi\)、Sachs 和流重整化约定在数值比较前必须统一。
  \item 准分布 Fourier 变换在 \(x\to-x\) 与复共轭下必须与所选相位约定一致；不同归一化只能通过 \(\mathcal N_\Gamma\) 或明确窗函数联系。
  \item 所有 bootstrap 重采样必须在组态层面完成，并在每个 replica 内重算关联函数、比值、拟合、重正化和 Fourier；只对最终中心值重采样会丢失相关性。
\end{enumerate}

\paragraph{合并后的最低验收门.}
一个 B02 结果只有同时具备以下证据才能进入物理结论：源、汇、插入时间和动量定义闭合；Wick 拓扑与真空减除可复核；拟合窗口、\(t_0\)、基大小或 \(N_v\) 有稳定性扫描；统计误差在 replica 内传播；连续极限、有限体积或有限 \(P^z\) 的误差有明确边界；并且明确区分方法回归、所引文献数值与 MyQCD 自己完成的验证。缺少任一项时，应把结果标记为中间矩阵元、诊断量或待验证项，而不是物理形状因子、散射参数或 TMD 软函数。

\chapter{B05 证据矩阵与验收表}\label{sec:b05-appendix}
\section{B05 单元覆盖}

\begin{longtable}{@{}>{\raggedright\arraybackslash}p{0.075\textwidth}>{\raggedright\arraybackslash}p{0.255\textwidth}>{\raggedright\arraybackslash}p{0.32\textwidth}>{\raggedright\arraybackslash}p{0.20\textwidth}@{}}
  \caption{B05 的 11 个内容单元、证据用途与处置。}\label{tab:b05-coverage}\\
  \toprule
  单元 & 主要 TeX 来源 & 进入正文的内容 & 处置与目标 \\
  \midrule
  \endfirsthead
  \toprule
  单元 & 主要 TeX 来源 & 进入正文的内容 & 处置与目标 \\
  \midrule
  \endhead
  \bottomrule
  \endfoot
  U035 & \texttt{构造胶子准算符.tex} & 六振幅、投影、重整化分类、格点流 & integrate；胶子算符、TMD 与自旋 \\
  U041 & \texttt{格点QCD中的TMD\_PDF.tex} & quark TMD、staple、软核、CS、匹配状态 & integrate；TMD 定义、软结构、匹配与验收 \\
  U063 & \texttt{格点QCD中的胶子极化.tex} & 轴反常、螺旋度算符、赝 ITD、系统误差 & integrate；胶子极化、TMD 与自旋 \\
  U064 & \texttt{格点QCD中的胶子算符.tex} & U035 扩展、hybrid/RI-MOM、应用边界 & integrate；胶子算符、TMD 与重正化 \\
  U076 & \texttt{格点上计算胶子准算符.tex} & 实现管道、蒸馏、非连通三点函数 & integrate；TMD、准 TMD 与软结构；本证据附录 \\
  U080 & \texttt{胶子PDF幻灯片.tex} & 因子化与准/赝 PDF 教学顺序 & deduplicate；U083/U086 \\
  U083 & \texttt{胶子PDF连续极限.tex} & 三格距、混合重整化、联合外推、误差 & integrate；TMD、准 TMD 与软结构 \\
  U084 & \texttt{胶子PDF连续极限幻灯片.tex} & 工作流、数据布局、GPU/MPI 说明 & deduplicate；U076/\allowbreak U083/\allowbreak U086 \\
  U085 & \texttt{胶子PDF连续极限论文.tex} & 与 U083 相同的正文、公式和补充材料 & deduplicate；U083 \\
  U086 & \texttt{胶子TMD梯度流连续极限.tex} & flow-TMD 推导、流时窗口、候选工作流 & integrate；梯度流 TMD、匹配与连续极限 \\
  U087 & \texttt{质子自旋危机解析.tex} & EMC 历史、两套自旋分解、证据争议 & integrate；质子自旋、胶子极化与连续极限 \\
\end{longtable}

\section{合并前的最低检查单}

\begin{enumerate}[leftmargin=*]
  \item \textbf{算符几何：}Wilson 线路径、方向、长度、端点和色表示是否与裸量、减除量及匹配核一致。
  \item \textbf{Lorentz 投影：}是否包含迹项、对偶张量符号、\(z\)-奇组合和核子自旋投影；是否把诊断单项误当 twist-2 本征算符。
  \item \textbf{UV 方案：}是否报告短距参考点、强子矩阵元归一化、\(Z_R\) 或 Wilson 圈、格距依赖和方案转换。
  \item \textbf{软与 rapidity：}TMD 是否除以正确的 \(S_I\) 平方根、使用同一表示 CS 核，并验证 \(\zeta\) 方案一致性。
  \item \textbf{匹配：}硬核阶数、Fourier 归一化、味单态混合、匹配标度和高阶误差是否写明。
  \item \textbf{极限：}是否同时报告大 \(z\) 尾部、多个 \(P_z\)、至少三格距的联合外推；flow 研究是否扫描多个流时。
  \item \textbf{误差：}统计、拟合、重正化、匹配、flow、手征、体积和外推误差是否分别给出并可传播到最终协方差。
  \item \textbf{证据等级：}结论是否停在表 \ref{tab:b05-claim-levels} 的已闭合层级，是否把 quark 先例或流程图误称为 gluon TMD 数值结果。
\end{enumerate}

\section{本批次保留的未解决项}

\begin{enumerate}[leftmargin=*]
  \item gluon TMD 的软函数和 CS 核如何以伴随表示与 quark 方案统一，尚无 B05 内的直接数据闭合。
  \item 大 \(b_T\) 区固定阶 OPE 失效后，格点数据、非微扰软核和现象学模型的误差如何分离。
  \item 味单态胶子--夸克混合矩阵的非微扰确定及其与 NLO/NNLO 匹配核的一致性。
  \item 梯度流与 rapidity/soft 方案的联合定义；流时外推和匹配阶数的相关误差尚缺多流时数值表。
  \item \(M_{ti;it}+M_{ij;ji}\) 中 \(M_{gg}\) 与接触项的精确约定，以及不同格点 Clover 定向的归一化一致性。
  \item 自旋预算的跨方案、跨标度协方差；现有预算区间不能直接作为闭合的物理结论。
\end{enumerate}

本附录不把“缺少一项”解释为整个框架无效。相反，B05 的可靠结论是：胶子算符、TMD 减除、软核、连续外推和自旋分解已经形成可审计的方法骨架；只有逐层满足上表条件，骨架才能升级为物理分布或自旋分量。

\chapter{GPD 扩展与旧胶子 PDF 代码}
\section{GPD 的运动学与算符定义}

广义部分子分布同时依赖纵向动量分数 $x$、偏度 $\xi$、动量转移平方 $t$ 和重正化标度 $\mu$。设初末核子动量为 $P_i,P_f$，采用
\begin{equation}
  P=\frac{P_i+P_f}{2},
  \qquad
  \Delta=P_f-P_i,
  \qquad
  t=\Delta^2,
  \qquad
  \xi=\frac{P_i^+-P_f^+}{P_i^++P_f^+}.
  \label{eq:b07-gpdkin}
\end{equation}
若文献采用 $\Delta=P_1-P_2$，必须同步交换初末态标签和旋量顺序，不能只改一个符号。

以夸克通道为例，非局域算符为
\begin{align}
  \mathcal O_q(z_1,z_2)
  &=\bar\psi(z_1)\Gamma[z_1,z_2]\psi(z_2),\nonumber\\
  [z_1,z_2]&=\mathcal P\exp\!\left[
  ig\int_0^1\mathrm dt\,(z_2-z_1)\cdot A(z_1+t(z_2-z_1))
  \right].
  \label{eq:b07-gpdop}
\end{align}
一个常用的自旋 $1/2$ 矢量分解写成
\begin{align}
  \frac{1}{2P^+}\langle P_f,S_f|
  &\bar q(-z^-/2)\gamma^+\nonumber\\
  &\quad\times W(-z^-/2,z^-/2)q(z^-/2)|P_i,S_i\rangle\nonumber\\
  &=\bar u(P_f,S_f)\left[
  \gamma^+H^q(x,\xi,t)
  +\frac{i\sigma^{+\mu}\Delta_\mu}{2M}E^q(x,\xi,t)\right]\nonumber\\
  &\quad\times u(P_i,S_i).
  \label{eq:b07-gpddecomp}
\end{align}
轴矢和横向通道分别给出带 tilde 与横向 GPD。换文献时，$E$、$\widetilde E$ 和横向 GPD 的符号约定必须整体换算。

夸克 GPD 的支撑区为 $x\in[-1,1]$。$|x|>|\xi|$ 是 DGLAP 区，在 $\xi=0$ 时退化到普通 PDF；$|x|<|\xi|$ 是 ERBL 区，对应夸克-反夸克振幅，不能直接解释成 PDF 的反夸克尾部。常见极限为
\begin{equation}
  \begin{aligned}
  H^q(x,0,0)=q(x),\quad
  H^q(-x,0,0)=-\bar q(x),\\
  \widetilde H^q(x,0,0)=\Delta q(x),\quad
  H_T^q(x,0,0)=\delta q(x).
  \end{aligned}
  \label{eq:b07-gpdlimits}
\end{equation}
对 $x$ 积分应回到相应形状因子；twist-2 GPD 的 Mellin 矩还满足多项式性。标准 Ji 关系可写成
\begin{equation}
  J_q=\frac12\left[A^q_{20}(0)+B^q_{20}(0)\right]
  =\frac12\int_{-1}^{1}\mathrm dx\,x
  \left[H^q(x,\xi,0)+E^q(x,\xi,0)\right].
  \label{eq:b07-ji}
\end{equation}
本批次只确认这是标准关系；没有证据表明 PyQCD 已实现该求和规则的数值提取。

\section{从准 GPD 到 GPD：格点链与匹配}

非前行三点函数需要独立的初末态动量。一个与历史审计一致的离散形式是
\begin{equation}
  C_3(t_f,\tau;z,P_f,P_i)
  =\sum_{\vec x,\vec y}
  e^{-i\vec P_f\cdot\vec x}
  e^{i(\vec P_f-\vec P_i)\cdot\vec y}
  \left\langle
  \mathcal N(\vec x,t_f)\,
  \mathcal O_\Gamma(z,\vec y,\tau)\,
  \bar{\mathcal N}(0,0)
  \right\rangle.
  \label{eq:b07-gpd3pt}
\end{equation}
源、汇和插入必须有独立轴与动量标签。仅仅把局部流替换成带空间 Wilson 线的 VdV，不会自动产生非前行运动学，也不会给出 GPD 的 Lorentz 分解。

准 GPD 与光锥 GPD 的一般匹配含味混合：
\begin{equation}
  \begin{pmatrix}\mathcal O_q\\ \mathcal O_g\end{pmatrix}
  =\begin{pmatrix}
  C_{qq}&C_{qg}\\ C_{gq}&C_{gg}
  \end{pmatrix}
  \otimes
  \begin{pmatrix}\mathcal O_q^{\rm LC}\\ \mathcal O_g^{\rm LC}\end{pmatrix}
  +\text{higher twist}.
  \label{eq:b07-gpdmatching}
\end{equation}
味非单态时只有 $C_{qq}$；核子单态必须保留胶子-夸克混合。非前行情形还涉及双 Ioffe 变量的 Fourier 结构，简单把单一 $z$ 变换套到 GPD 会隐藏 $\xi$ 和端点约定。

历史审计的实现状态可以概括为：胶子双场强、直线 Wilson 线、staple、通用 Wick、VdV/VVV 和允许初末态两点函数分离的比值接口存在；通用夸克双局域算符、独立 $P_i/P_f$ 的三点管线、GPD 自旋/洛伦兹投影、非前行 NLO 匹配和真实 GPD 数据不存在。现有前向 PDF/TMD 测试不能证明 GPD。

\begin{longtable}{>{\raggedright\arraybackslash}p{3.2cm}>{\raggedright\arraybackslash}p{4.2cm}>{\raggedright\arraybackslash}p{6.4cm}}
\caption{GPD 扩展的最低验收矩阵。}\label{tab:b07-gpdgates}\\
\toprule
门槛 & 通过条件 & 当前证据边界\\
\midrule
\endfirsthead
\toprule
门槛 & 通过条件 & 当前证据边界\\
\midrule
\endhead
约定 & 固定 $P_i,P_f,P,\Delta,\xi,t$、Fermi 相位、Wilson 方向和 Euclidean 基 & 理论可写，生产链未绑定\\
算符 & $z=0$ 回局域流；恒等链接回无链双线性；规范变换协变 & 有几何前驱，通用夸克算符缺失\\
收缩 & 独立动量标签；新增 GPD 图 oracle & 有通用 Wick/PJN oracle，非前行未见\\
比值 & $\Delta=0$ 或 $\xi=0$ 时回前向结果 & \texttt{ratio\_3pt} 接口可复用，上游数据缺失\\
重整化 & 同一 $Z$ 在不同动量态上一致，复相位不丢失 & 乘法重正化有理论依据，未接线\\
匹配 & $\xi=0$ 回前向核；PDF/DA 极限通过 & 只有前向 NLO 核，非前行核未实现\\
真实数据 & 组态对齐、平台、协方差和体积/格距误差 & 未发现 GPD 数据或验收\\
\bottomrule
\end{longtable}

\section{旧版 donghx 胶子 PDF 代码的结构}

旧代码把质子两点函数与胶子 OPE 分开计算，再用比值提取矩阵元。蒸馏部分用 Laplacian 低模投影压缩传播子：
\begin{equation}
  \Box_{xy}(t)=\sum_{k=1}^{N_{\rm ev}}v_x^{(k)}(t)v_y^{(k)\dagger}(t),
  \label{eq:b07-distillation}
\end{equation}
子空间传播子为
\begin{equation}
  \tau_{\alpha\beta}^{AB}(t_s,t_i)
  =v^{(A)\dagger}(t_s)D^{-1}_{\alpha\beta}(t_s,t_i)v^{(B)}(t_i).
  \label{eq:b07-perambulator}
\end{equation}
重子块通过 Levi-Civita 收缩三个本征矢量：
\begin{equation}
  \Phi_{abc}(P,t)
  =\sum_{\mathbf x}e^{-iP\cdot\mathbf x}
  \epsilon_{ijk}v_i^a(\mathbf x,t)v_j^b(\mathbf x,t)v_k^c(\mathbf x,t).
  \label{eq:b07-vvv}
\end{equation}
旧实现逐源时刻和源位置构造 $\Phi$，再以直接项与交换项收缩 quark perambulator，形成源分辨质子两点函数。该结构适合复用，但代码中的动量轴、相位符号、$C\gamma_5\gamma_\mu$ 型和宇称投影必须与 DR 基一起核对。

胶子算符由 Clover 场强、对偶场强和沿 $z$ 的 Wilson 线组成。历史代码依次把场强平移到 $z$、乘反方向链接回 0、插入对偶场强、再乘正方向链接回 $z$，最后取色迹并对空间求和。旧 TeX 把 disconnected 近似写成
\begin{equation}
  C_3(\Delta t,z)\propto C_2(\Delta t)\langle\mathcal O(z)\rangle.
  \label{eq:b07-legacyfactor}
\end{equation}
这只能视为因子化候选。实际提取应使用 \eqref{eq:b07-vacuum} 的真空扣除和同一源汇归一化；把 $\langle\mathcal O\rangle$ 直接当成不含真空结构的“OPE 真空期望”，会在非零真空贡献和有限源汇下出错。

\section{旧版 Fourier 与 LO 匹配的解释边界}

旧工作流把坐标空间矩阵元作 sine 变换：
\begin{equation}
  \tilde g(x,P_z)=\frac{2P_z}{x}
  \int_0^{z_{\max}}\mathrm dz\,
  h(z,P_z)\sin(xP_zz).
  \label{eq:b07-legacyfourier}
\end{equation}
代码只使用 $h(z,P_z)$ 的实部，并在 $x=0$ 置零。该做法依赖目标通道的实部/反对称性；若矩阵元有物理相位或正负 $z$ 耦合，必须先证明投影。匹配层只实现
\begin{equation}
  C_{gg}^{\rm LO}(\xi,\mu/P_z)=\delta(1-\xi),
  \label{eq:b07-legacyLO}
\end{equation}
因此在 LO 下 $\tilde g$ 与 $g$ 相同只是最低阶近似；不能把旧代码的 LO 输出当作 NLO 匹配结果。

旧代码的 Clover 公式与展示代码在 $1/(a^2g_0)$ 归一化上存在口径差异，PDF 中的希腊字母和上下标也有抽取损坏。专著应采用“代码返回的裸组合”与“连续场强定义”分开写的策略，并在唯一一个约定表中固定符号。旧报告的参考文献目录和代码片段是历史资产；它没有提供真实数据平台、连续极限或误差预算，因此只能作为附录中的代码考古，而不是物理结果的独立验证。

\section{本附录的状态总结}

GPD 部分应明确写成“理论定义与扩展蓝图已建立，生产链未启动”。旧胶子 PDF 代码部分应明确写成“前向 LaMET 代码考古，验证级别为历史实现”。两类材料都不能被当前胶子 TMD 综述中的一般理论覆盖，也不能反过来证明当前 PyQCD 已经完成 GPD 或物理胶子 TMD-PDF。

\chapter{历史证据与验证协议}\label{sec:b08-appendix-evidence}
本附录把历史日志转换为可复用的验证协议。协议不依赖某个版本的报告排版，也不
以“程序退出 0”代替物理正确性。每个阶段都必须同时记录输入身份、输出形状、数值
oracle、失败状态和证据等级。

\section{证据等级与允许措辞}

\begin{table}[htbp]
\centering
\caption{B08 使用的证据等级。}
\label{tab:b08-evidence-levels}
\small
\begin{tabularx}{\textwidth}{>{\raggedright\arraybackslash}p{1.4cm}>{\raggedright\arraybackslash}p{4.4cm}>{\raggedright\arraybackslash}X}
\toprule
等级 & 已证明的对象 & 允许措辞与禁止升级 \\
\midrule
E0 & 来源可定位、文件可读取 & “有历史材料”；不能据此发表物理结论 \\
E1 & 两遍编译、页框、缺字、渲染与安全区通过 & “报告可编译/可渲染”；不证明公式或算法 \\
E2 & 解析反例、低层公式 oracle、形状/归一化检查 & “该公式/接口在测试域内一致” \\
E3 & 中间对象的输出投影或文件关系一致 & “该输出关系通过”；不证明上游独立重算 \\
E4 & 真实数据的拟合、谱学、ratio 或矩阵元通过质量门 & “受控中间量”；仍缺 proof of physical TMD \\
E5 & 多格距、多流时、多动量、soft/rapidity、混合和极限闭合 & 才可使用“物理胶子 TMD-PDF 结果” \\
\bottomrule
\end{tabularx}
\end{table}

E1 到 E4 不可相互替代。一个 45/45 的流程检查、一个 58/58 的输出投影检查或
一个单组态有限矩阵元，都不能越过 E4 进入 E5。报告正文应把等级放在表格单元
或状态字段中，而不是只在附录文字中说明。

\section{数值 oracle、归一化与输出投影}

\paragraph{解析反例优先。}
任何 source/configuration 归一化都应有小数组反例。至少固定
\(C_2=[1,3]\)、\(O=[2,6]\)，分别验证 source 平均、组态平均、真空乘积减除和
最终 ratio；预期值、实际值和失败原因必须并排保存。只检查输出有限（NaN/Inf）
不足以发现因子 \(N_{\rm conf}\)、\(N_s\) 或重复平均。

\paragraph{低层公式对照。}
对 eigvec、phase、VdV、VVV、Clover、dual 和 OPE，可定义相对的截断误差
\begin{equation}
  \epsilon_{L2}(a,b)=\frac{\lVert a-b\rVert_2}
  {\max\{\lVert b\rVert_2,\epsilon_0\}},
  \label{eq:b08-relative-l2}
\end{equation}
其中 \(\epsilon_0\) 用于避免近零分母。报告必须同时给出最大绝对差和使用的
\(dtype\) 容差；若基准是同一规范场上的独立公式重写，证据等级只能记为 E2，
不能写成外部 raw 文件复现。4150 的 7 个低层对象、4D HYP 的 9/9 通道和
effmass 的 7 个聚合资产分别属于不同测试域，不能合并成单个“全部通过”。

\paragraph{输出投影。}
对 momentum-smear 文件，可检查
\begin{equation}
  \text{contract}\longrightarrow P_+\longrightarrow
  \text{anti-periodic boundary}\longrightarrow \text{nopol},
  \label{eq:b08-output-projection}
\end{equation}
并分别报告根目录、动量组、数组数、相对 L2 与缺失槽位。该检查证明输出文件
遵守投影和边界关系，不证明独立 smeared perambulator、相位生成或原始
VVV 一致。极化检查必须把标准 \texttt{li}、旧 \texttt{il}、两者均 diff 和文件缺失分成四个
互斥状态；近零分母导致的巨大相对误差应附绝对差，不能单独作为物理失败。

\section{真实数据与失败关闭}

真实数据链的最小状态集合为
\begin{equation}
  \mathcal S=
  (\text{success fraction},\text{coverage},\text{rank},
  \text{fallback},\text{window},\text{physical acceptance}).
  \label{eq:b08-state-vector}
\end{equation}
每个分量都必须有机器可读值。一个候选拟合成功不代表它可进入下一层；中心样本
通过也不代表 replica 无 fallback。对 dev20 诊断链，合法的状态是
\texttt{diagnostic\_only}，而不是用更大的误差棒掩盖流窗口或 soft/mixing 缺口。

协议至少执行以下检查：
\begin{enumerate}
  \item 严格读取 source 轴，拒绝一维源平均缓存；检查 shape、dtype 与 finite；
  \item 在每个 replica 内生成共享指数并重算真空协方差，保存 index SHA-256；
  \item 对每个 \(t,\tau,z,b_T\) 记录可辨识性、拟合质量、平台平坦度和支持状态；
  \item 逐点检查 \(a\ll r_f\ll b_T,z\) 与 \(P_zr_f\ll1\)，窗口外点只记诊断；
  \item 对 soft/rapidity、gluon-quark mixing、多格距和多流时缺失设置硬失败；
  \item 在最终 artifact 中同时保存 gauge/flow/OPE contract 链和 fallback 率。
\end{enumerate}

失败关闭不是把程序改成“总是报错”，而是把“缺少必要输入”和“数值计算失败”分开。
例如单组态下统计不可辨识，应输出 \texttt{statistically\_unidentifiable} 并停止 PDF
导出；这比生成一个任意中心值更符合物理验收。

\section{LaTeX/PDF 输出验收}

对每份报告至少运行两遍 XeLaTeX，并使用独立输出目录。通过条件包括：两遍均
退出 0；交叉引用稳定；\texttt{Overfull}、\texttt{Float too large}、
\texttt{Missing character} 和 \texttt{Underfull} 的适用计数为 0；页框数与实际页数相等。

随后用 \texttt{pdftoppm} 渲染每一页，检查
\begin{equation}
  N_{\rm frames}=N_{\rm pages}=N_{\rm rendered}
  \quad\wedge\quad
  N_{\rm occlusion}=N_{\rm clipped}=N_{\rm outside}=0.
  \label{eq:b08-render-acceptance}
\end{equation}
历史 pdfcheck 已给出 23 页和 27 页两个实例：前者两遍退出 0、无 overfull、
无 float-too-large、无缺字；后者进一步把交叉引用、underfull、纵横向 overfull
及安全区纳入检查。该协议只能给出 E1，报告中的算法和数值仍需各自的 E2--E4。

\section{历史时间线与交接点}

{\footnotesize
\begin{longtable}{>{\raggedright\arraybackslash}p{1.9cm}>{\raggedright\arraybackslash}p{2.8cm}>{\raggedright\arraybackslash}p{3.8cm}>{\raggedright\arraybackslash}p{4.6cm}}
\caption{从蒸馏基线到 fail-closed 审计的证据演进。}\label{tab:b08-timeline}\\
\toprule
时间/批次 & 输入规模 & 新增能力 & 未闭合点 \\
\midrule
\endfirsthead
\toprule
时间/批次 & 输入规模 & 新增能力 & 未闭合点 \\
\midrule
\endhead
stab0 v0803 & 3 组态，\(N_{\rm ev}=50\) & 蒸馏、2pt、3pt 原型 & 高动量统计、胶子 OPE \\
stab0 v0804 & 3 组态，\(N_{\rm ev}=100\) & 更大低模空间与更多关联目标 & 质子平台缺失、pion 动量退化 \\
stab0 v0805 & 10 组态，\(N_{\rm ev}=100\) & OPE、3pt/4pt、自动报告 & 不相连拟合不可辨识、无连续极限 \\
stab1 & 10 组态真实基线 & ratio、能量、FH、三方向模块 & 单 z、低秩协方差、无 soft/mixing \\
dev5 & 10 组态；\(a=0.0938\,\mathrm{fm}\) & flow、staple、ratio、hybrid、NLO/CS & 历史归一化、单格距/单流时、无 soft/mixing \\
4150 & 单组态逐层资产 & HYP、费米子、momsmear、effmass 对照 & 3D HYP、raw VVV、3pt/ratio 等未验证 \\
dev20 & 11 组态、1000 bootstrap & 源分辨归一化、严格质量门、fail-closed & 流窗口、soft/rapidity、混合、多尺度极限 \\
\bottomrule
\end{longtable}
}

\section{可复用的最小验收清单}

\begin{enumerate}
  \item 输入是否有稳定 ID、哈希、shape、dtype、边界条件和轴语义？
  \item 每个公式是否有解析反例或小矩阵 oracle，且误差阈值与 dtype 相符？
  \item 每个文件名是否与其实际投影、极化、源和动量语义一致？
  \item 每个 ratio 是否在 replica 内完成归一化，而不是先平均或后置覆盖？
  \item 每个 \((z,b_T)\) 点是否满足流窗口和 Fourier 支撑，不满足时是否标记诊断？
  \item 是否分开报告中心质量、replica fallback、协方差秩和物理状态？
  \item 是否把同几何 soft/rapidity、完整 mixing 和多格距列为硬门？
  \item 报告是否完成两遍编译、全页渲染、页框和安全区检查？
  \item 最终结论的措辞是否与证据等级表一致，未把 E1--E4 写成 E5？
\end{enumerate}

以上清单把历史工作流转化为专著可复核的最小规范。它既解释了
stab0、stab1 与 dev5
为何仍有方法学价值，也解释了 4150 的逐层对照为何不能外推，以及 dev20 为何在
真实数据可运行后仍必须保持失败关闭。

\chapter{关键公式与约定}

本附录给出主文反复使用的公式索引。页码不进入公式表，完整来源与读取深度见B04
批次的`coverage.tsv`和`notes.md`。

\begin{longtable}{@{}p{0.18\textwidth}p{0.40\textwidth}p{0.34\textwidth}@{}}
  \caption{B04关键公式与必须同时声明的约定}\label{tab:b04-formula-reference}\\
  \toprule
  对象 & 公式或定义 & 配套约定与限制 \\
  \midrule
  \endfirsthead
  \toprule
  对象 & 公式或定义 & 配套约定与限制 \\
  \midrule
  \endhead
  \bottomrule
  \endfoot
  光锥坐标 & $v^\pm=(v^0\pm v^3)/\sqrt2$，$g_{+-}=1$ & 与$P^\pm$、Fourier相位和Wilson线方向一起固定 \\
  夸克PDF & 式\eqref{eq:b04-lc-quark} & 指定场序、$+$规范链、$\overline{\rm MS}$或格点方案 \\
  胶子PDF & 式\eqref{eq:b04-lc-gluon} & 左侧为$xg$；$F_{+}^{\ \mu}$指标和伴随链接须一致 \\
  守恒 & 式\eqref{eq:b04-pdf-sumrules} & 动量规则含胶子；价数规则按味和反夸克定义 \\
  DGLAP & 式\eqref{eq:b04-dglap-nonsinglet}与\eqref{eq:b04-dglap-singlet} & $\mu$为因子化标度；端点采用$[\,]_+$减除 \\
  劈裂核 & 式\eqref{eq:b04-splitting} & LO；$n_f$为活跃轻味数，颜色因子按$SU(3)$ \\
  Mellin & 式\eqref{eq:b04-mellin}至\eqref{eq:b04-gamma-qq} & 矩指标$n$与算符自旋、多项式次数同步 \\
  OPE & 式\eqref{eq:b04-short-ope}与\eqref{eq:b04-lightray-ope} & 普通短距与light-ray的幂次计数不同 \\
  准PDF & 式\eqref{eq:b04-quasi}与\eqref{eq:b04-lamet-factorization} & 必须声明$\Gamma$、投影、$P_z$和匹配方案 \\
  ITD & 式\eqref{eq:b04-itd}与\eqref{eq:b04-ritd} & $\nu=zP_z$；分母为零动量/真空同几何矩阵元 \\
  胶子算符 & 式\eqref{eq:b04-gluon-combination} & 只适用于所列指标组合；不可外推到全部36分量 \\
  GPD & 式\eqref{eq:b04-gpd-kinematics}与\eqref{eq:b04-vector-gpd} & $\Delta=P_f-P_i$；旋量、相位和$E$项整体固定 \\
  GPD求和 & 式\eqref{eq:b04-gpd-formfactors}与\eqref{eq:b04-ji-sumrule} & leading twist；形状因子与$t\to0$外推不可省略 \\
\end{longtable}

\chapter{计算验收清单}

本清单把“代码可运行”和“物理量可发布”分开。建议每个结果至少标注
\texttt{diagnostic-only}、\texttt{candidate-physics}或\texttt{accepted}，并保存
对应证据。

\begin{enumerate}
  \item \textbf{约定：}固定光锥坐标、度规、Fourier相位、场强符号、生成元归一化、
  Euclidean gamma约定和反夸克定义。
  \item \textbf{算符：}声明Dirac结构、味组合、色表示、Wilson线端点与方向、
  胶子指标缩并和投影。$\gamma^t$与$\gamma^z$不得在未处理连续混合时互换。
  \item \textbf{归一化：}验证$x=0$、$z=0$、零动量和真空减除；检查动量求和、
  价数、PDF前向极限和GPD形状因子极限。
  \item \textbf{UV：}证明线性Wilson线发散被同几何减除，剩余对数进入正确方案；
  对混合通道给出完整矩阵，而不是一个标量$Z$。
  \item \textbf{匹配：}固定匹配阶数、$\mu$、投影、方案转换、味单态矩阵和
  $[\,]_+$端点。扫描$\mu$与阶数，检查残差是否进入误差预算。
  \item \textbf{幂修正：}先声明目标是twist-2、twist-3还是更高；分别报告端点增强、
  靶质量和重正子模糊。不得把式\eqref{eq:b04-qpdf-power}的$1/P_z^2$自动解释为
  twist-4。
  \item \textbf{Fourier与ITD：}报告$z_{\max}$、$\lambda=zP_z$范围、尾部模型、
  相位、归一化和端点权。大$z$外推结果必须与不同窗口和不同$P_z$交叉检验。
  \item \textbf{极限：}连续极限、无穷动量、手征、体积和激发态至少各有两个可辨识
  参数点；单点结果只能作为该点的有效结果。
  \item \textbf{误差：}传播统计协方差和系统误差相关性；$z_s$、$z_0$、$\mu$、
  $\lambda$窗口和匹配阶数都应作为独立系统源。
  \item \textbf{证据：}保留组态标识、作用量与流参数、源型、缓存键、真残差、
  原始或可重建矩阵元及哈希。中心值一致不能替代窗口和方案检查。
\end{enumerate}

\paragraph{物理接受条件.}
只有当算符与约定闭合、UV/rapidity及混合受控、匹配和幂修正处于声明精度、
所有极限在窗口内稳定、误差包含相关系统项，并与独立方案或独立数据交叉一致时，
才可把结果标记为\texttt{accepted}。若仅完成矩阵元、比值或单参数点，应标记为
\texttt{diagnostic-only}；若主要窗口已通过但连续极限或单态混合尚未闭合，应标记
为\texttt{candidate-physics}，不得改称最终物理PDF或GPD。

\chapter{胶子算符、关联函数与统计约定}\label{app:b09-conventions}
本附录只收集本批次证据链实际使用的约定。不同报告若把 \(z,b_\perp\)、场强符号、Wilson-line 方向或 Fourier 变量换掉，数值即使看似一致也不能直接相加或比较。

\section{低模、动量相位与颜色顶点}
\label{app:b09-vertices}

对空间体积 \(N_x^3\) 和低模 \(\phi_n^a(x)\)，定义
\begin{align}
  e_p(x)&=\exp\left[-\frac{2\pi i}{N_x}
  \left(p_zz+p_yy+p_xx\right)\right],
  \label{eq:app-phase}\\
  V_{mn}(p)&=\sum_{x,a}\phi_m^{a*}(x)e_p(x)\phi_n^a(x),
  \label{eq:app-vdv}\\
  V_{mnl}(p)&=\sum_x e_p(x)\epsilon_{abc}
  \phi_m^a(x)\phi_n^b(x)\phi_l^c(x).
  \label{eq:app-vvv}
\end{align}
\(V_{mn}\) 是低模基上的二体颜色/空间投影，\(V_{mnl}\) 把三个夸克颜色指标组合成 singlet。式~\eqref{eq:app-vvv} 的六置换符号来自 \(\epsilon_{abc}\)，在缓存和转置时必须保留。4150 的真实 eigvec 为 \((8,13824,3)\)，phase 为 \((13824,)\)，VdV 为 \((1,4,4)\)，VVV 为 \((4,4,4)\)。

\section{Clover 场强、对偶场强与直线 OPE}
\label{app:b09-field-strength}

本批次采用
\begin{align}
  F_{\mu\nu}(x)&=-\frac{i}{8}
  \left(C_{\mu\nu}(x)-C_{\mu\nu}^\dagger(x)\right),
  \label{eq:app-clover}\\
  \widetilde F_{\mu\nu}(x)&=\frac12
  \epsilon_{\mu\nu\rho\sigma}F_{\rho\sigma}(x),
  \label{eq:app-dual}
\end{align}
其中 \(C_{\mu\nu}\) 是四个正向和反向 plaquette 的 Clover 组合。直线 Wilson-line OPE 写成
\begin{equation}
  M_{\mu\nu;\rho\sigma}(z)=
  \sum_x\operatorname{tr}\left[
  F_{\mu\nu}(x+z\hat d)W^\dagger(0\to z)
  \widetilde F_{\rho\sigma}(x)W(0\to z)\right],
  \label{eq:app-line-ope}
\end{equation}
而 4150 的非极化 HYP 通道使用第二场强 \(F\) 而不是 \(\widetilde F\)。正负方向要使用各自独立的 link 链；取实部、绝对值或错误方向的链都不能被解释为同一矩阵元。

\section{胶子 TMD 组合与核子两点函数}
\label{app:b09-correlators}

本批次使用的胶子组合为
\begin{equation}
  O(z,b_\perp)=M^{tx;tx}+M^{ty;ty}-2M^{xy;xy}.
  \label{eq:app-tmd-operator}
\end{equation}
它是报告所采用的特定 Lorentz 结构。若改变流化、场强投影或 Wilson-line 几何，必须重新建立算符基和匹配。

核子两点函数写成
\begin{equation}
  C_2(\boldsymbol p;t_s,t_0)=
  \left\langle O_p(\boldsymbol p,t_s)\bar O_p(\boldsymbol p,t_0)\right\rangle .
  \label{eq:app-c2}
\end{equation}
正宇称投影采用 \(P_+=(1/2)(\gamma_0+\gamma_4)\)。对时间方向的反周期性，报告使用
\begin{equation}
  s_{++}(t_s,t_0)=
  \begin{cases}
    -1,&t_s<t_0,\\
    +1,&t_s\ge t_0.
  \end{cases}
  \label{eq:app-boundary-sign}
\end{equation}
这个符号必须在缩并阶段与投影一起作用。若先算 \(C\)、取绝对值或在统计平均后补符号，反向传播和周期镜像会被混合。

\section{复 GEVP}
\label{app:b09-gevp}

相关矩阵先 Hermitian 化：
\begin{equation}
  C_H(t)=\frac{C(t)+C^\dagger(t)}{2},
  \qquad
  C_H(t)v_n(t,t_0)=\lambda_n(t,t_0)C_H(t_0)v_n(t,t_0).
  \label{eq:app-gevp}
\end{equation}
复 Hermitian 矩阵的本征值是实数，但 \(v_n\) 可以是复向量。一个实现可以把输出投影到实部以符合旧接口，但该操作改变了原方程，必须在约定中显式标记。判别实现是否正确时，应报告原复 GEVP 的残差、广义正交性和解析谱误差，不能只报告投影后的逐位一致。

\section{数组级误差与统计秩}
\label{app:b09-error}

4150 对照采用
\begin{equation}
  \varepsilon_{\mathrm{rel}}(A_{\mathrm{PyQCD}},A_{\mathrm{ref}})
  =\frac{\left\|A_{\mathrm{PyQCD}}-A_{\mathrm{ref}}\right\|_2}
  {\max\left(\left\|A_{\mathrm{ref}}\right\|_2,10^{-300}\right)} .
  \label{eq:app-relerr}
\end{equation}
它是整数组的相对 \(L_2\) 指标，不是逐元素最大相对误差。对接近零的极化槽位，式~\eqref{eq:app-relerr} 的分母很小，可能把绝对差 \(10^{-4}\) 放大成 \(10^5\)；报告必须同时给出最大绝对差、范数和缺失状态。

对 \(N\) 个独立组态的中心化数据矩阵 \(X\)，样本协方差为
\begin{equation}
  \widehat C=\frac{1}{N-1}X^\dagger X,
  \qquad
  \operatorname{rank}\widehat C\le N-1.
  \label{eq:app-rank}
\end{equation}
Bootstrap 的 replica 数不是独立样本数。要从不足秩的协方差获得数值稳定的拟合，只能使用伪逆、SVD 截断、shrinkage 或对角加载，并把这些操作的阈值作为系统选择报告。否则“200 次 bootstrap 得到满秩”这一说法在数学上不成立。

\chapter{B09 复现证据矩阵与未决边界}
\label{app:b09-evidence}

本附录给出可直接回填来源表的九单元矩阵。表中“整合”表示正文采用；“去重”表示只保留其并入单元中的公式、数值或历史限制。PDF 仅用于页数、版式和独立结论核对，不与 TeX 重复计作两份物理证据。

\begingroup
\setlength{\tabcolsep}{3pt}
\begin{longtable}{@{}p{0.08\textwidth}p{0.18\textwidth}p{0.18\textwidth}p{0.30\textwidth}p{0.16\textwidth}@{}}
\caption{B09 单元、证据和处置（主证据为简称，完整路径见 \texttt{coverage.tsv}）}\label{tab:b09-evidence-matrix}\\
\hline
单元 & 主证据 & 关键数值或对象 & 本批次处置 & 未决边界 \\
\hline
\endfirsthead
\hline
单元 & 主证据 & 关键数值或对象 & 本批次处置 & 未决边界 \\
\hline
\endhead
U017 & test6 report 2026-08-16 & 879 组态；\(E_0=3.069(9)\) GeV；参考差 0 & 整合入真实数据复现与历史验证附录 & 单一格距；平台与自相关未扩展到物理极限 \\
\hline
U018 & test7 report & 1 组态；P0/P2 都写成 0.000 GeV；缺图 & 并入 U020 的历史限制 & 内部矛盾，不能作为物理结果 \\
\hline
U019 & test8 report & 2 组态；约 1.060/1.517 GeV；26/26 输出投影 & 并入 U020 的演进基线 & 无误差、无窗口；PDF 与 TeX 不同步 \\
\hline
U020 & test9 analysis & 10 组态；\(\tau=3a^2\)；准/匹配 0.57/0.59；\(K=1.6\) & 整合入准 TMD 几何、真实数据与历史验证章节 & 10 样本协方差、无完整 soft/rapidity、单流时、匹配参数缺失 \\
\hline
U025 & 4150 cmp1 report 2026-08-28 & 52 项中 51 直接通过；复 GEVP 残差 \(1.02\times10^{-15}\) & 整合入真实数据复现与验证附录 & 外部 4150 TMD/3pt 成品缺失；不能做端到端闭环 \\
\hline
U026 & donghx 4150 reproduction 2026-08-29 & 低层 7/7；HYP-OPE 9/9；2pt 与直线 OPE & 并入 U027 & 独立 smeared peram、逐时间 VVV、横向 staple 缺失 \\
\hline
U027 & donghx 4150 reproduction 2026-08-30 & 58/58 正宇称；effmass 25/25；极化 74/116 通过 & 整合入真实数据复现与验证附录 & 1 个极化异常；3pt, ratio, barematrix 未闭环 \\
\hline
U028 & donghx comparison summary & 统一 \(\varepsilon_{\mathrm{rel}}\)；五类通过项 & 并入 U027 与胶子算符约定附录 & 只列通过项，需与 U027 的失败清单同读 \\
\hline
U029 & donghx comparison table & 逐行 4150 对照表 & 并入 U027/U028 & 无独立物理结论，不复制本地路径长表 \\
\hline
\end{longtable}
\endgroup

\section{必须保留的未决项}

\begin{enumerate}
  \item test9 的 10 组态不能通过 200 次 bootstrap 获得秩 10 以上的样本协方差。任何拟合若使用了正则化，必须报告正则化方案、谱截断、失败 replica 和中心值移动。
  \item test9 的准 TMD 与匹配后数值缺少匹配核、标度、阶数、卷积实现和误差传播，只能作为量级诊断。
  \item test9 的 CS 核在高 \(b\) 触发 clamp；clamp 之前与之后的点属于不同统计质量，不能用一条平滑曲线掩盖。
  \item 4150 的直线 OPE 通过不覆盖横向 staple；非零 \(b_\perp\) 的 Wilson-line 自能、端点和软结构仍未验证。
  \item 4150 的 58 个极化文件缺失、41 个旧 \(il\) 轴序命中和 1 个近零异常槽位必须分栏保留，不能合并成单一通过率。
  \item 4150 的独立 smeared perambulator、逐时间 VVV、3pt/ratio/barematrix 和 Cluster Decomposition 仍未形成可配对闭环。
  \item test7/test8 的 1/2 组态和中心值只用于说明演进与失败关闭，不作为质量、能量或矩阵元测量。
  \item 现有专著中的“11 个组态”快照与本批次 test9 的“10 个组态”来源必须核对。数量不同可能来自不同运行窗口，专著不能在没有 manifest 对照时把两者混为同一条数据链。
\end{enumerate}

\section{建议写入主文的断言}

\begin{enumerate}
  \item 真实输入上的数组级一致性是必要的实现门，但不是物理接受门。
  \item GEVP 的复本征矢不能被默认实数化；必须报告原复问题的残差。
  \item 自重整化只在其乘性分解和几何前提成立时有效，不能替代 soft/rapidity 与算符混合。
  \item 直线 OPE、横向 staple TMD 和完整光锥 TMD 是三个不同层级的对象。
  \item 统计显著性由独立组态数、误差传播和窗口稳定性决定，重采样次数只描述数值误差估计的精度。
  \item 只要缺少同几何 soft/rapidity、显式匹配核或多尺度极限，结果就应保持 \texttt{diagnostic\_only}。
\end{enumerate}

\chapter{术语、符号与约定}
\section{对象与标度}

\begin{longtable}{@{}p{0.20\textwidth}p{0.34\textwidth}p{0.36\textwidth}@{}}
\caption{胶子 PDF/TMD 计算中最容易混淆的术语。}\label{tab:b10-terms}\\
\toprule
术语 & 定义 & 常见错误 \\
\midrule
\endfirsthead
\toprule
术语 & 定义 & 常见错误 \\
\midrule
\endhead
共线 PDF & 对横向动量积分后只保留纵向份额 $x$ 的分布；依赖 $\mu$。 & 把
quasi-PDF 直接当成共线 PDF，或把 TMD 的 $b_\perp$ 积分当成严格恒等式而不报告
幂修正。 \\
TMD-PDF & 同时依赖 $x,\bperp,\mu,\zeta$，含 staple 与软/快度结构。 & 只加一条
横向 Wilson 线而忽略软函数、过程方向和 CS 演化。 \\
quasi-PDF & 欧氏空间方向分离、由大 $\Pz$ 匹配到光锥 PDF 的准分布。 & 把有限
$\Pz$ 结果直接称为光锥结果。 \\
pseudo-ITD & 在 Ioffe 时间 $\nu=z\Pz$ 中表示同一矩阵元。 & 把坐标截断后的
振荡当成物理端点行为。 \\
RI/MOM & 在离壳夸克顶点和动量窗口内定义的非微扰 UV 方案。 & 对非局域 Wilson
线直接套用局域算符的减除，忽略 $z\mu$ 混合。 \\
hybrid & 短距 ratio 与长距自重整化/质量减除的拼接。 & 在 $z_s$ 两侧重复或遗漏
因子，或按点选择过渡尺度。 \\
CS 核 & 控制 TMD 对 $\zeta$ 演化的快度反常量纲。 & 把两动量比值的有效斜率
无条件外推到非微扰 $b_\perp$。 \\
SFTE & 正流时间算符对 $\MSbar$ 算符的小 $\tau$ 展开。 & 把流化后的 UV 有限性
误认为无需匹配，或把大流降噪结果用于小流外推。 \\
Wilson 圈减除 & 用同几何闭合圈吸收线性自能、pinch/cusp 结构。 & 用不同长度、
不同表示或不同转角几何的圈作减除。 \\
\bottomrule
\end{longtable}

\section{符号检查表}

\begin{tabularx}{\textwidth}{@{}p{0.15\textwidth}p{0.36\textwidth}X@{}}
\toprule
符号 & 含义 & 量纲/检查 \\
\midrule
$a$ & 格距 & fm 或 $\mathrm{GeV}^{-1}$；所有指数与对数必须同单位。 \\
$\tau$ & 流时间 & 物理量纲为长度平方；代码常用 $\tau/a^2$。 \\
$\sqrt{8\tau}$ & 流化半径 & 与 $|z|$、$b_\perp$ 比较时使用同一单位。 \\
$\Pz$ & 核子 boost 动量 & 不是 UV 标度 $\mu$，也不是快度标度 $\zeta$。 \\
$\nu=z\Pz$ & Ioffe 时间 & 无量纲；Fourier 共轭到 $x$。 \\
$Z_R$ & 自重整化因子 & 依赖 $z,a,\mu$；不能和硬匹配 $Z_{ij}$ 混为一谈。 \\
$K(\bperp,\mu)$ & Collins--Soper 核 & 依赖表示、Wilson 线几何和快度方案。 \\
$S_I(\bperp,\mu)$ & 内禀软函数 & 与完整软函数、$\zeta$ 方案分层记录。 \\
\bottomrule
\end{tabularx}

U055 的术语表适合作为附录索引，但它本身是综合性资料而非原始论文。其中弦张力
条目的数值与括号单位存在自相矛盾；专著应回到原始势能定义和自然单位换算，
不能原样复制该数字。U055 对 PDF、TMD、LaMET、RI/MOM、hybrid、CS 核和
Wilson 线仍可作为统一命名的有用底稿。

\chapter{来源导航、证据等级与版本谱系}

\section{六单元的使用方式}

\begin{longtable}{@{}p{0.13\textwidth}p{0.29\textwidth}p{0.48\textwidth}@{}}
\caption{B10 来源的证据定位。}\label{tab:b10-sources}\\
\toprule
单元 & 来源类型 & 专著中的用法 \\
\midrule
\endfirsthead
\toprule
单元 & 来源类型 & 专著中的用法 \\
\midrule
\endhead
U008 & 归档目录 README & 只证明旧 \texttt{examples/...} 路径和日志文件是历史快照；不能作为
现行运行入口或新物理证据。 \\
U033 & 2022 年准 PDF 学位论文 & 摘要级历史对照：RI/MOM 与 Wilson 线线性发散的
问题意识。正文数值与完整重整化条件未在本批次核对。 \\
U034 & 2023 年 TMD 学位论文 & 目录级工作流线索：TMD、Wilson loop、HMC/Chroma
和统计重采样。扫描 OCR 有字符损坏，未使用其公式。 \\
U055 & 格点 QCD 术语/缩写表 & 术语与符号附录的分类与索引；具体物理结论回溯到
U077 或原始文献，不把术语表当作独立推导。 \\
U078 & 旧版共线 PDF 工作流 & 历史版式与 Eq.(20)/Eq.(25) 叙事；已由 U077 取代，
只保留演进解释，不重复计入物理内容。 \\
U077 & 新版胶子 TMD 工作流 & 历史与复现章节的主要 TeX 证据：定义、SFTE、方案、
ratio-v2、匹配、验证边界和历史链。PDF 只用于页数与成对核对。 \\
\bottomrule
\end{longtable}

\section{行号、页码与 OCR 规则}

TeX/Markdown 来源使用“文件名:行号”锚定，例如 PDF/TMD 定义可回溯到
\texttt{理论解析与工作流-v20260911.tex:492--680}，算符与 TMD staple 可回溯到
\texttt{同文件:1220--1373}，SFTE 与窗口可回溯到 \texttt{同文件:2188--2300}。
这些行号指向固定快照，不随当前工作树的其他修改变化。

独立 PDF 使用物理 PDF 页码。U033 的第 2 页是中文摘要（印刷页 i），第 3 页
是英文摘要（印刷页 iii）；U034 的第 3 页是目录首页。扫描 PDF 的 OCR 出现
拆字、错码和字符重叠时，必须标为“公式/文本未可靠抽取”，不能根据残片补写
公式。PDF 页数与 TeX 自述不一致时，以来源文件的物理页数和 TeX 行为主记录，
并在 unresolved 中说明口径差异。

\section{证据等级}

\begin{enumerate}
  \item \textbf{A：定义或推导。} TeX 中给出对象、方程和适用条件的部分；
  仍需检查符号约定是否与专著主文一致。
  \item \textbf{B：来源内验证。} U077 登记的契约、回归、oracle 或参考诊断；
  未在 B10 重跑时只能写作“来源登记”。
  \item \textbf{C：摘要/目录级证据。} U033、U034 的题录、摘要和目录；可用于
  历史定位，不足以证明完整算法。
  \item \textbf{D：历史导航。} U008 的旧路径和 U078 的被替代版本；用于解释
  迁移与版本差异，不产生新的物理结论。
\end{enumerate}

任何数值进入正文前应回答四个问题：它来自哪个对象、在哪个尺度窗口、对应哪个
重整化/匹配方案、是否有独立复跑或只属于来源登记。若四个问题不能同时回答，
应把数值放入诊断表或未解决清单，而不是写成第一性原理结果。

\chapter{PyQCD 文档全量来源图谱}
\label{app:pyqcd-source-atlas}
\providecommand{\sourcepath}[1]{\seqsplit{#1}}
本附录把固定快照中的每一个文件映射到专著的来源单元和章节用途。它回答三个审计问题：文件是否存在且未被改写、知识内容进入哪一章、TeX/PDF 成对件与历史版本的重复关系是否被显式处理。This appendix is a machine-auditable ledger, not a bibliography by volume alone.
\begin{checkpointbox}{来源基线}
\small
\begin{itemize}
  \item 上游仓库：\url{https://gitee.com/zhangxin8069/PyQCD.git}。
  \item 固定提交：\texttt{0a2f85d30881fc8e}\allowbreak\texttt{ad60972c0afb9c02}\allowbreak\texttt{8abc348f}。
  \item 许可证：MIT，见 \sourcepath{docs/pyqcd\_reference/LICENSE.md}。
  \item 文件总数：166；内容文件：163；控制文件：3。
  \item 内容类型：77 个 TeX、80 个 PDF、6 个 Markdown；其中 76 对 TeX/PDF 成对保存，4 个 PDF 独立，1 个 TeX 独立。
  \item 文本规模：56\,850 行（TeX/Markdown）；PDF 规模：1\,566 页；原始字节：56\,165\,822。
  \item 完整性：163/163 个内容文件与吸收 manifest 的 SHA-256 一致；缺失文件 0，额外文件 0。
\end{itemize}
\end{checkpointbox}
\section{内容单元与章节映射}
成对 TeX/PDF 共享一个知识单元；PDF 仍保留为独立文件证据，并在逐文件账本中登记。每个单元在主批次中恰好出现一次。
\begin{description}[leftmargin=4.2em,labelwidth=3.7em,style=nextline,itemsep=0.35em]
\item[\texttt{U001}] \sourcepath{docs/pyqcd\_reference/PyQCD仓库结构与核心物理链解析.tex}
\par\noindent{\scriptsize
\textcolor{qcdsage}{伴生：}\sourcepath{docs/pyqcd\_reference/PyQCD仓库结构与核心物理链解析.pdf}。
类型：TeX + PDF；规模：166 行 / 4 页；目标：PyQCD 实现与复现。
}
\item[\texttt{U002}] \sourcepath{docs/pyqcd\_reference/absorb\_dev20\_gluon\_tmd\_audit\_20261022.md}
\par\noindent{\scriptsize
类型：Markdown；规模：378 行；目标：历史工作流与验证演进。
}
\item[\texttt{U003}] \sourcepath{docs/pyqcd\_reference/analy\_gpd\_20260928.tex}
\par\noindent{\scriptsize
\textcolor{qcdsage}{伴生：}\sourcepath{docs/pyqcd\_reference/analy\_gpd\_20260928.pdf}。
类型：TeX + PDF；规模：937 行 / 21 页；目标：PyQCD 实现与复现。
}
\item[\texttt{U004}] \sourcepath{docs/pyqcd\_reference/analy\_physics\_chain\_20260824.tex}
\par\noindent{\scriptsize
\textcolor{qcdsage}{伴生：}\sourcepath{docs/pyqcd\_reference/analy\_physics\_chain\_20260824.pdf}。
类型：TeX + PDF；规模：975 行 / 29 页；目标：PyQCD 实现与复现。
}
\item[\texttt{U005}] \sourcepath{docs/pyqcd\_reference/analy\_proton\_energy\_20260816.tex}
\par\noindent{\scriptsize
类型：独立 TeX；规模：300 行；目标：PyQCD 实现与复现。
}
\item[\texttt{U006}] \sourcepath{docs/pyqcd\_reference/analy\_pyqcd\_20260818.tex}
\par\noindent{\scriptsize
\textcolor{qcdsage}{伴生：}\sourcepath{docs/pyqcd\_reference/analy\_pyqcd\_20260818.pdf}。
类型：TeX + PDF；规模：324 行 / 11 页；目标：PyQCD 实现与复现。
}
\item[\texttt{U007}] \sourcepath{docs/pyqcd\_reference/analy\_torch\_parallel\_20260816.tex}
\par\noindent{\scriptsize
\textcolor{qcdsage}{伴生：}\sourcepath{docs/pyqcd\_reference/analy\_torch\_parallel\_20260816.pdf}。
类型：TeX + PDF；规模：442 行 / 9 页；目标：PyQCD 实现与复现。
}
\item[\texttt{U008}] \sourcepath{docs/pyqcd\_reference/from-logs/README.md}
\par\noindent{\scriptsize
类型：Markdown；规模：7 行；目标：学位论文综合与来源词库。
}
\item[\texttt{U009}] \sourcepath{docs/pyqcd\_reference/from-logs/logs/dev5/analy\_test9\_20260820.tex}
\par\noindent{\scriptsize
\textcolor{qcdsage}{伴生：}\sourcepath{docs/pyqcd\_reference/from-logs/logs/dev5/analy\_test9\_20260820.pdf}。
类型：TeX + PDF；规模：818 行 / 25 页；目标：历史工作流与验证演进。
}
\item[\texttt{U010}] \sourcepath{docs/pyqcd\_reference/from-logs/logs/dev5\_1/analy\_test9\_20260821.tex}
\par\noindent{\scriptsize
\textcolor{qcdsage}{伴生：}\sourcepath{docs/pyqcd\_reference/from-logs/logs/dev5\_1/analy\_test9\_20260821.pdf}。
类型：TeX + PDF；规模：863 行 / 33 页；目标：历史工作流与验证演进。
}
\item[\texttt{U011}] \sourcepath{docs/pyqcd\_reference/from-logs/logs/dev5\_2/analy\_test9\_20260822.tex}
\par\noindent{\scriptsize
\textcolor{qcdsage}{伴生：}\sourcepath{docs/pyqcd\_reference/from-logs/logs/dev5\_2/analy\_test9\_20260822.pdf}。
类型：TeX + PDF；规模：901 行 / 36 页；目标：历史工作流与验证演进。
}
\item[\texttt{U012}] \sourcepath{docs/pyqcd\_reference/from-logs/logs/donghx\_4150\_20260829/pdfcheck.md}
\par\noindent{\scriptsize
类型：Markdown；规模：70 行；目标：历史工作流与验证演进。
}
\item[\texttt{U013}] \sourcepath{docs/pyqcd\_reference/from-logs/logs/donghx\_4150\_20260829/summary.md}
\par\noindent{\scriptsize
类型：Markdown；规模：278 行；目标：历史工作流与验证演进。
}
\item[\texttt{U014}] \sourcepath{docs/pyqcd\_reference/from-logs/logs/stab0/stab0.tex}
\par\noindent{\scriptsize
\textcolor{qcdsage}{伴生：}\sourcepath{docs/pyqcd\_reference/from-logs/logs/stab0/stab0.pdf}。
类型：TeX + PDF；规模：589 行 / 13 页；目标：历史工作流与验证演进。
}
\item[\texttt{U015}] \sourcepath{docs/pyqcd\_reference/from-logs/logs/stab1/docs/stab1\_analysis.tex}
\par\noindent{\scriptsize
\textcolor{qcdsage}{伴生：}\sourcepath{docs/pyqcd\_reference/from-logs/logs/stab1/docs/stab1\_analysis.pdf}。
类型：TeX + PDF；规模：150 行 / 6 页；目标：历史工作流与验证演进。
}
\item[\texttt{U016}] \sourcepath{docs/pyqcd\_reference/from-logs/logs/stab1/v202608151614/stab1\_report.tex}
\par\noindent{\scriptsize
\textcolor{qcdsage}{伴生：}\sourcepath{docs/pyqcd\_reference/from-logs/logs/stab1/v202608151614/stab1\_report.pdf}。
类型：TeX + PDF；规模：74 行 / 6 页；目标：历史工作流与验证演进。
}
\item[\texttt{U017}] \sourcepath{docs/pyqcd\_reference/from-logs/logs/test6/docs/analy\_test6\_report\_20260816.tex}
\par\noindent{\scriptsize
\textcolor{qcdsage}{伴生：}\sourcepath{docs/pyqcd\_reference/from-logs/logs/test6/docs/analy\_test6\_report\_20260816.pdf}。
类型：TeX + PDF；规模：412 行 / 11 页；目标：真实数据复现与验证矩阵。
}
\item[\texttt{U018}] \sourcepath{docs/pyqcd\_reference/from-logs/logs/test7/v202608161007/test7\_report.tex}
\par\noindent{\scriptsize
\textcolor{qcdsage}{伴生：}\sourcepath{docs/pyqcd\_reference/from-logs/logs/test7/v202608161007/test7\_report.pdf}。
类型：TeX + PDF；规模：50 行 / 2 页；目标：真实数据复现与验证矩阵。
}
\item[\texttt{U019}] \sourcepath{docs/pyqcd\_reference/from-logs/logs/test8/v202608161117/test8\_report.tex}
\par\noindent{\scriptsize
\textcolor{qcdsage}{伴生：}\sourcepath{docs/pyqcd\_reference/from-logs/logs/test8/v202608161117/test8\_report.pdf}。
类型：TeX + PDF；规模：71 行 / 5 页；目标：真实数据复现与验证矩阵。
}
\item[\texttt{U020}] \sourcepath{docs/pyqcd\_reference/from-logs/logs/test9/test9\_analysis.tex}
\par\noindent{\scriptsize
\textcolor{qcdsage}{伴生：}\sourcepath{docs/pyqcd\_reference/from-logs/logs/test9/test9\_analysis.pdf}。
类型：TeX + PDF；规模：233 行 / 8 页；目标：真实数据复现与验证矩阵。
}
\item[\texttt{U021}] \sourcepath{docs/pyqcd\_reference/gevp\_convergence.pdf}
\par\noindent{\scriptsize
类型：独立 PDF；规模：1 页；目标：关联函数、谱学与有限体积。
}
\item[\texttt{U022}] \sourcepath{docs/pyqcd\_reference/gevp\_t0\_suppression.pdf}
\par\noindent{\scriptsize
类型：独立 PDF；规模：1 页；目标：关联函数、谱学与有限体积。
}
\item[\texttt{U023}] \sourcepath{docs/pyqcd\_reference/gradient\_flow.md}
\par\noindent{\scriptsize
类型：Markdown；规模：895 行；目标：梯度流、流时窗口与胶子 TMD。
}
\item[\texttt{U024}] \sourcepath{docs/pyqcd\_reference/pure\_pyqcd\_20260818.tex}
\par\noindent{\scriptsize
\textcolor{qcdsage}{伴生：}\sourcepath{docs/pyqcd\_reference/pure\_pyqcd\_20260818.pdf}。
类型：TeX + PDF；规模：304 行 / 10 页；目标：PyQCD 实现与复现。
}
\item[\texttt{U025}] \sourcepath{docs/pyqcd\_reference/report\_cmp1\_4150\_20260828.tex}
\par\noindent{\scriptsize
\textcolor{qcdsage}{伴生：}\sourcepath{docs/pyqcd\_reference/report\_cmp1\_4150\_20260828.pdf}。
类型：TeX + PDF；规模：507 行 / 16 页；目标：真实数据复现与验证矩阵。
}
\item[\texttt{U026}] \sourcepath{docs/pyqcd\_reference/report\_donghx\_4150\_reproduction\_20260829.tex}
\par\noindent{\scriptsize
\textcolor{qcdsage}{伴生：}\sourcepath{docs/pyqcd\_reference/report\_donghx\_4150\_reproduction\_20260829.pdf}。
类型：TeX + PDF；规模：617 行 / 23 页；目标：真实数据复现与验证矩阵。
}
\item[\texttt{U027}] \sourcepath{docs/pyqcd\_reference/report\_donghx\_4150\_reproduction\_20260830.tex}
\par\noindent{\scriptsize
\textcolor{qcdsage}{伴生：}\sourcepath{docs/pyqcd\_reference/report\_donghx\_4150\_reproduction\_20260830.pdf}。
类型：TeX + PDF；规模：732 行 / 27 页；目标：真实数据复现与验证矩阵。
}
\item[\texttt{U028}] \sourcepath{docs/pyqcd\_reference/report\_donghx\_pyqcd\_compare\_20260830.tex}
\par\noindent{\scriptsize
\textcolor{qcdsage}{伴生：}\sourcepath{docs/pyqcd\_reference/report\_donghx\_pyqcd\_compare\_20260830.pdf}。
类型：TeX + PDF；规模：361 行 / 12 页；目标：真实数据复现与验证矩阵。
}
\item[\texttt{U029}] \sourcepath{docs/pyqcd\_reference/report\_donghx\_pyqcd\_compare\_20260830\_2.tex}
\par\noindent{\scriptsize
\textcolor{qcdsage}{伴生：}\sourcepath{docs/pyqcd\_reference/report\_donghx\_pyqcd\_compare\_20260830\_2.pdf}。
类型：TeX + PDF；规模：351 行 / 10 页；目标：真实数据复现与验证矩阵。
}
\item[\texttt{U030}] \sourcepath{docs/pyqcd\_reference/report\_gluon\_tmd\_gradient\_flow\_20260828.tex}
\par\noindent{\scriptsize
\textcolor{qcdsage}{伴生：}\sourcepath{docs/pyqcd\_reference/report\_gluon\_tmd\_gradient\_flow\_20260828.pdf}。
类型：TeX + PDF；规模：789 行 / 22 页；目标：梯度流、流时窗口与胶子 TMD。
}
\item[\texttt{U031}] \sourcepath{docs/pyqcd\_reference/report\_gradient\_flow\_gluonpdf\_pyqcd\_20260909.tex}
\par\noindent{\scriptsize
\textcolor{qcdsage}{伴生：}\sourcepath{docs/pyqcd\_reference/report\_gradient\_flow\_gluonpdf\_pyqcd\_20260909.pdf}。
类型：TeX + PDF；规模：1029 行 / 17 页；目标：梯度流、流时窗口与胶子 TMD。
}
\item[\texttt{U032}] \sourcepath{docs/pyqcd\_reference/基于Gradient\_Flow的质子胶子PDF格点计算整体流程.md}
\par\noindent{\scriptsize
类型：Markdown；规模：661 行；目标：梯度流、流时窗口与胶子 TMD。
}
\item[\texttt{U033}] \sourcepath{docs/pyqcd\_reference/基于格点QCD对部分子分布函数的研究.pdf}
\par\noindent{\scriptsize
类型：独立 PDF；规模：56 页；目标：学位论文综合与来源词库。
}
\item[\texttt{U034}] \sourcepath{docs/pyqcd\_reference/基于格点QCD的横向动量依赖的部分子分布函数的研究.pdf}
\par\noindent{\scriptsize
类型：独立 PDF；规模：82 页；目标：学位论文综合与来源词库。
}
\item[\texttt{U035}] \sourcepath{docs/pyqcd\_reference/构造胶子准算符.tex}
\par\noindent{\scriptsize
\textcolor{qcdsage}{伴生：}\sourcepath{docs/pyqcd\_reference/构造胶子准算符.pdf}。
类型：TeX + PDF；规模：562 行 / 14 页；目标：胶子算符、TMD 与自旋。
}
\item[\texttt{U036}] \sourcepath{docs/pyqcd\_reference/格点QCD中的3pt构造.tex}
\par\noindent{\scriptsize
\textcolor{qcdsage}{伴生：}\sourcepath{docs/pyqcd\_reference/格点QCD中的3pt构造.pdf}。
类型：TeX + PDF；规模：783 行 / 20 页；目标：关联函数、谱学与有限体积。
}
\item[\texttt{U037}] \sourcepath{docs/pyqcd\_reference/格点QCD中的DGLAP演化方程.tex}
\par\noindent{\scriptsize
\textcolor{qcdsage}{伴生：}\sourcepath{docs/pyqcd\_reference/格点QCD中的DGLAP演化方程.pdf}。
类型：TeX + PDF；规模：688 行 / 18 页；目标：PDF、OPE、LaMET 与 GPD。
}
\item[\texttt{U038}] \sourcepath{docs/pyqcd\_reference/格点QCD中的GEVP.tex}
\par\noindent{\scriptsize
\textcolor{qcdsage}{伴生：}\sourcepath{docs/pyqcd\_reference/格点QCD中的GEVP.pdf}。
类型：TeX + PDF；规模：924 行 / 24 页；目标：关联函数、谱学与有限体积。
}
\item[\texttt{U039}] \sourcepath{docs/pyqcd\_reference/格点QCD中的OPE算符.tex}
\par\noindent{\scriptsize
\textcolor{qcdsage}{伴生：}\sourcepath{docs/pyqcd\_reference/格点QCD中的OPE算符.pdf}。
类型：TeX + PDF；规模：526 行 / 14 页；目标：PDF、OPE、LaMET 与 GPD。
}
\item[\texttt{U040}] \sourcepath{docs/pyqcd\_reference/格点QCD中的Symanzik有效理论.tex}
\par\noindent{\scriptsize
\textcolor{qcdsage}{伴生：}\sourcepath{docs/pyqcd\_reference/格点QCD中的Symanzik有效理论.pdf}。
类型：TeX + PDF；规模：449 行 / 13 页；目标：格点规范理论基础与系统误差。
}
\item[\texttt{U041}] \sourcepath{docs/pyqcd\_reference/格点QCD中的TMD\_PDF.tex}
\par\noindent{\scriptsize
\textcolor{qcdsage}{伴生：}\sourcepath{docs/pyqcd\_reference/格点QCD中的TMD\_PDF.pdf}。
类型：TeX + PDF；规模：937 行 / 30 页；目标：胶子算符、TMD 与自旋。
}
\item[\texttt{U042}] \sourcepath{docs/pyqcd\_reference/格点QCD中的UV涨落.tex}
\par\noindent{\scriptsize
\textcolor{qcdsage}{伴生：}\sourcepath{docs/pyqcd\_reference/格点QCD中的UV涨落.pdf}。
类型：TeX + PDF；规模：472 行 / 13 页；目标：格点规范理论基础与系统误差。
}
\item[\texttt{U043}] \sourcepath{docs/pyqcd\_reference/格点QCD中的Wilson线.tex}
\par\noindent{\scriptsize
\textcolor{qcdsage}{伴生：}\sourcepath{docs/pyqcd\_reference/格点QCD中的Wilson线.pdf}。
类型：TeX + PDF；规模：712 行 / 21 页；目标：格点规范理论基础与系统误差。
}
\item[\texttt{U044}] \sourcepath{docs/pyqcd\_reference/格点QCD中的luscher参数化.tex}
\par\noindent{\scriptsize
\textcolor{qcdsage}{伴生：}\sourcepath{docs/pyqcd\_reference/格点QCD中的luscher参数化.pdf}。
类型：TeX + PDF；规模：603 行 / 17 页；目标：关联函数、谱学与有限体积。
}
\item[\texttt{U045}] \sourcepath{docs/pyqcd\_reference/格点QCD中的smear算法.tex}
\par\noindent{\scriptsize
\textcolor{qcdsage}{伴生：}\sourcepath{docs/pyqcd\_reference/格点QCD中的smear算法.pdf}。
类型：TeX + PDF；规模：630 行 / 16 页；目标：数值算法、统计推断与外推。
}
\item[\texttt{U046}] \sourcepath{docs/pyqcd\_reference/格点QCD中的传播子求解.tex}
\par\noindent{\scriptsize
\textcolor{qcdsage}{伴生：}\sourcepath{docs/pyqcd\_reference/格点QCD中的传播子求解.pdf}。
类型：TeX + PDF；规模：269 行 / 10 页；目标：数值算法、统计推断与外推。
}
\item[\texttt{U047}] \sourcepath{docs/pyqcd\_reference/格点QCD中的光锥PDF、准PDF、赝PDF.tex}
\par\noindent{\scriptsize
\textcolor{qcdsage}{伴生：}\sourcepath{docs/pyqcd\_reference/格点QCD中的光锥PDF、准PDF、赝PDF.pdf}。
类型：TeX + PDF；规模：705 行 / 20 页；目标：PDF、OPE、LaMET 与 GPD。
}
\item[\texttt{U048}] \sourcepath{docs/pyqcd\_reference/格点QCD中的光锥PDF与quasi\_PDF.tex}
\par\noindent{\scriptsize
\textcolor{qcdsage}{伴生：}\sourcepath{docs/pyqcd\_reference/格点QCD中的光锥PDF与quasi\_PDF.pdf}。
类型：TeX + PDF；规模：413 行 / 12 页；目标：PDF、OPE、LaMET 与 GPD。
}
\item[\texttt{U049}] \sourcepath{docs/pyqcd\_reference/格点QCD中的光锥与光前.tex}
\par\noindent{\scriptsize
\textcolor{qcdsage}{伴生：}\sourcepath{docs/pyqcd\_reference/格点QCD中的光锥与光前.pdf}。
类型：TeX + PDF；规模：783 行 / 18 页；目标：PDF、OPE、LaMET 与 GPD。
}
\item[\texttt{U050}] \sourcepath{docs/pyqcd\_reference/格点QCD中的关联函数.tex}
\par\noindent{\scriptsize
\textcolor{qcdsage}{伴生：}\sourcepath{docs/pyqcd\_reference/格点QCD中的关联函数.pdf}。
类型：TeX + PDF；规模：604 行 / 15 页；目标：关联函数、谱学与有限体积。
}
\item[\texttt{U051}] \sourcepath{docs/pyqcd\_reference/格点QCD中的场强张量.tex}
\par\noindent{\scriptsize
\textcolor{qcdsage}{伴生：}\sourcepath{docs/pyqcd\_reference/格点QCD中的场强张量.pdf}。
类型：TeX + PDF；规模：576 行 / 14 页；目标：格点规范理论基础与系统误差。
}
\item[\texttt{U052}] \sourcepath{docs/pyqcd\_reference/格点QCD中的外推.tex}
\par\noindent{\scriptsize
\textcolor{qcdsage}{伴生：}\sourcepath{docs/pyqcd\_reference/格点QCD中的外推.pdf}。
类型：TeX + PDF；规模：567 行 / 16 页；目标：数值算法、统计推断与外推。
}
\item[\texttt{U053}] \sourcepath{docs/pyqcd\_reference/格点QCD中的大动量有效理论.tex}
\par\noindent{\scriptsize
\textcolor{qcdsage}{伴生：}\sourcepath{docs/pyqcd\_reference/格点QCD中的大动量有效理论.pdf}。
类型：TeX + PDF；规模：498 行 / 13 页；目标：PDF、OPE、LaMET 与 GPD。
}
\item[\texttt{U054}] \sourcepath{docs/pyqcd\_reference/格点QCD中的夸克算符.tex}
\par\noindent{\scriptsize
\textcolor{qcdsage}{伴生：}\sourcepath{docs/pyqcd\_reference/格点QCD中的夸克算符.pdf}。
类型：TeX + PDF；规模：457 行 / 13 页；目标：PDF、OPE、LaMET 与 GPD。
}
\item[\texttt{U055}] \sourcepath{docs/pyqcd\_reference/格点QCD中的常见名词短语与缩写.tex}
\par\noindent{\scriptsize
\textcolor{qcdsage}{伴生：}\sourcepath{docs/pyqcd\_reference/格点QCD中的常见名词短语与缩写.pdf}。
类型：TeX + PDF；规模：414 行 / 26 页；目标：学位论文综合与来源词库。
}
\item[\texttt{U056}] \sourcepath{docs/pyqcd\_reference/格点QCD中的常见费曼图.tex}
\par\noindent{\scriptsize
\textcolor{qcdsage}{伴生：}\sourcepath{docs/pyqcd\_reference/格点QCD中的常见费曼图.pdf}。
类型：TeX + PDF；规模：495 行 / 10 页；目标：关联函数、谱学与有限体积。
}
\item[\texttt{U057}] \sourcepath{docs/pyqcd\_reference/格点QCD中的广义部分子分布.tex}
\par\noindent{\scriptsize
\textcolor{qcdsage}{伴生：}\sourcepath{docs/pyqcd\_reference/格点QCD中的广义部分子分布.pdf}。
类型：TeX + PDF；规模：1354 行 / 25 页；目标：PDF、OPE、LaMET 与 GPD。
}
\item[\texttt{U058}] \sourcepath{docs/pyqcd\_reference/格点QCD中的强子谱学.tex}
\par\noindent{\scriptsize
\textcolor{qcdsage}{伴生：}\sourcepath{docs/pyqcd\_reference/格点QCD中的强子谱学.pdf}。
类型：TeX + PDF；规模：353 行 / 11 页；目标：关联函数、谱学与有限体积。
}
\item[\texttt{U059}] \sourcepath{docs/pyqcd\_reference/格点QCD中的形状因子.tex}
\par\noindent{\scriptsize
\textcolor{qcdsage}{伴生：}\sourcepath{docs/pyqcd\_reference/格点QCD中的形状因子.pdf}。
类型：TeX + PDF；规模：499 行 / 14 页；目标：关联函数、谱学与有限体积。
}
\item[\texttt{U060}] \sourcepath{docs/pyqcd\_reference/格点QCD中的标度参数.tex}
\par\noindent{\scriptsize
\textcolor{qcdsage}{伴生：}\sourcepath{docs/pyqcd\_reference/格点QCD中的标度参数.pdf}。
类型：TeX + PDF；规模：456 行 / 13 页；目标：格点规范理论基础与系统误差。
}
\item[\texttt{U061}] \sourcepath{docs/pyqcd\_reference/格点QCD中的梯度流重整化.tex}
\par\noindent{\scriptsize
\textcolor{qcdsage}{伴生：}\sourcepath{docs/pyqcd\_reference/格点QCD中的梯度流重整化.pdf}。
类型：TeX + PDF；规模：626 行 / 15 页；目标：梯度流、流时窗口与胶子 TMD。
}
\item[\texttt{U062}] \sourcepath{docs/pyqcd\_reference/格点QCD中的维克收缩与傅里叶变换.tex}
\par\noindent{\scriptsize
\textcolor{qcdsage}{伴生：}\sourcepath{docs/pyqcd\_reference/格点QCD中的维克收缩与傅里叶变换.pdf}。
类型：TeX + PDF；规模：469 行 / 13 页；目标：关联函数、谱学与有限体积。
}
\item[\texttt{U063}] \sourcepath{docs/pyqcd\_reference/格点QCD中的胶子极化.tex}
\par\noindent{\scriptsize
\textcolor{qcdsage}{伴生：}\sourcepath{docs/pyqcd\_reference/格点QCD中的胶子极化.pdf}。
类型：TeX + PDF；规模：694 行 / 16 页；目标：胶子算符、TMD 与自旋。
}
\item[\texttt{U064}] \sourcepath{docs/pyqcd\_reference/格点QCD中的胶子算符.tex}
\par\noindent{\scriptsize
\textcolor{qcdsage}{伴生：}\sourcepath{docs/pyqcd\_reference/格点QCD中的胶子算符.pdf}。
类型：TeX + PDF；规模：618 行 / 15 页；目标：胶子算符、TMD 与自旋。
}
\item[\texttt{U065}] \sourcepath{docs/pyqcd\_reference/格点QCD中的色散关系.tex}
\par\noindent{\scriptsize
\textcolor{qcdsage}{伴生：}\sourcepath{docs/pyqcd\_reference/格点QCD中的色散关系.pdf}。
类型：TeX + PDF；规模：496 行 / 15 页；目标：格点规范理论基础与系统误差。
}
\item[\texttt{U066}] \sourcepath{docs/pyqcd\_reference/格点QCD中的蒙卡方法.tex}
\par\noindent{\scriptsize
\textcolor{qcdsage}{伴生：}\sourcepath{docs/pyqcd\_reference/格点QCD中的蒙卡方法.pdf}。
类型：TeX + PDF；规模：347 行 / 11 页；目标：格点规范理论基础与系统误差。
}
\item[\texttt{U067}] \sourcepath{docs/pyqcd\_reference/格点QCD中的衰变道.tex}
\par\noindent{\scriptsize
\textcolor{qcdsage}{伴生：}\sourcepath{docs/pyqcd\_reference/格点QCD中的衰变道.pdf}。
类型：TeX + PDF；规模：474 行 / 15 页；目标：关联函数、谱学与有限体积。
}
\item[\texttt{U068}] \sourcepath{docs/pyqcd\_reference/格点QCD中的误差统计.tex}
\par\noindent{\scriptsize
\textcolor{qcdsage}{伴生：}\sourcepath{docs/pyqcd\_reference/格点QCD中的误差统计.pdf}。
类型：TeX + PDF；规模：472 行 / 15 页；目标：数值算法、统计推断与外推。
}
\item[\texttt{U069}] \sourcepath{docs/pyqcd\_reference/格点QCD中的费米子方案.tex}
\par\noindent{\scriptsize
\textcolor{qcdsage}{伴生：}\sourcepath{docs/pyqcd\_reference/格点QCD中的费米子方案.pdf}。
类型：TeX + PDF；规模：999 行 / 24 页；目标：格点规范理论基础与系统误差。
}
\item[\texttt{U070}] \sourcepath{docs/pyqcd\_reference/格点QCD中的连通图与非连通图.tex}
\par\noindent{\scriptsize
\textcolor{qcdsage}{伴生：}\sourcepath{docs/pyqcd\_reference/格点QCD中的连通图与非连通图.pdf}。
类型：TeX + PDF；规模：531 行 / 16 页；目标：关联函数、谱学与有限体积。
}
\item[\texttt{U071}] \sourcepath{docs/pyqcd\_reference/格点QCD中的部分子分布函数.tex}
\par\noindent{\scriptsize
\textcolor{qcdsage}{伴生：}\sourcepath{docs/pyqcd\_reference/格点QCD中的部分子分布函数.pdf}。
类型：TeX + PDF；规模：576 行 / 16 页；目标：PDF、OPE、LaMET 与 GPD。
}
\item[\texttt{U072}] \sourcepath{docs/pyqcd\_reference/格点QCD中的重夸克有效理论.tex}
\par\noindent{\scriptsize
\textcolor{qcdsage}{伴生：}\sourcepath{docs/pyqcd\_reference/格点QCD中的重夸克有效理论.pdf}。
类型：TeX + PDF；规模：426 行 / 12 页；目标：格点规范理论基础与系统误差。
}
\item[\texttt{U073}] \sourcepath{docs/pyqcd\_reference/格点QCD中的重整化.tex}
\par\noindent{\scriptsize
\textcolor{qcdsage}{伴生：}\sourcepath{docs/pyqcd\_reference/格点QCD中的重整化.pdf}。
类型：TeX + PDF；规模：741 行 / 21 页；目标：格点规范理论基础与系统误差。
}
\item[\texttt{U074}] \sourcepath{docs/pyqcd\_reference/格点QCD中的重采样方法.tex}
\par\noindent{\scriptsize
\textcolor{qcdsage}{伴生：}\sourcepath{docs/pyqcd\_reference/格点QCD中的重采样方法.pdf}。
类型：TeX + PDF；规模：507 行 / 16 页；目标：数值算法、统计推断与外推。
}
\item[\texttt{U075}] \sourcepath{docs/pyqcd\_reference/格点QCD蒸馏方法解析.tex}
\par\noindent{\scriptsize
\textcolor{qcdsage}{伴生：}\sourcepath{docs/pyqcd\_reference/格点QCD蒸馏方法解析.pdf}。
类型：TeX + PDF；规模：586 行 / 14 页；目标：关联函数、谱学与有限体积。
}
\item[\texttt{U076}] \sourcepath{docs/pyqcd\_reference/格点上计算胶子准算符.tex}
\par\noindent{\scriptsize
\textcolor{qcdsage}{伴生：}\sourcepath{docs/pyqcd\_reference/格点上计算胶子准算符.pdf}。
类型：TeX + PDF；规模：947 行 / 22 页；目标：胶子算符、TMD 与自旋。
}
\item[\texttt{U077}] \sourcepath{docs/pyqcd\_reference/理论解析与工作流-v20260911.tex}
\par\noindent{\scriptsize
\textcolor{qcdsage}{伴生：}\sourcepath{docs/pyqcd\_reference/理论解析与工作流-v20260911.pdf}。
类型：TeX + PDF；规模：5957 行 / 100 页；目标：学位论文综合与来源词库。
}
\item[\texttt{U078}] \sourcepath{docs/pyqcd\_reference/理论解析与工作流.tex}
\par\noindent{\scriptsize
\textcolor{qcdsage}{伴生：}\sourcepath{docs/pyqcd\_reference/理论解析与工作流.pdf}。
类型：TeX + PDF；规模：3837 行 / 89 页；目标：学位论文综合与来源词库。
}
\item[\texttt{U079}] \sourcepath{docs/pyqcd\_reference/胶子PDF完整推导.tex}
\par\noindent{\scriptsize
\textcolor{qcdsage}{伴生：}\sourcepath{docs/pyqcd\_reference/胶子PDF完整推导.pdf}。
类型：TeX + PDF；规模：742 行 / 15 页；目标：PDF、OPE、LaMET 与 GPD。
}
\item[\texttt{U080}] \sourcepath{docs/pyqcd\_reference/胶子PDF幻灯片.tex}
\par\noindent{\scriptsize
\textcolor{qcdsage}{伴生：}\sourcepath{docs/pyqcd\_reference/胶子PDF幻灯片.pdf}。
类型：TeX + PDF；规模：357 行 / 10 页；目标：胶子算符、TMD 与自旋。
}
\item[\texttt{U081}] \sourcepath{docs/pyqcd\_reference/胶子PDF笔记.tex}
\par\noindent{\scriptsize
\textcolor{qcdsage}{伴生：}\sourcepath{docs/pyqcd\_reference/胶子PDF笔记.pdf}。
类型：TeX + PDF；规模：445 行 / 9 页；目标：PDF、OPE、LaMET 与 GPD。
}
\item[\texttt{U082}] \sourcepath{docs/pyqcd\_reference/胶子PDF计算代码解析.tex}
\par\noindent{\scriptsize
\textcolor{qcdsage}{伴生：}\sourcepath{docs/pyqcd\_reference/胶子PDF计算代码解析.pdf}。
类型：TeX + PDF；规模：1218 行 / 30 页；目标：PyQCD 实现与复现。
}
\item[\texttt{U083}] \sourcepath{docs/pyqcd\_reference/胶子PDF连续极限.tex}
\par\noindent{\scriptsize
\textcolor{qcdsage}{伴生：}\sourcepath{docs/pyqcd\_reference/胶子PDF连续极限.pdf}。
类型：TeX + PDF；规模：578 行 / 16 页；目标：胶子算符、TMD 与自旋。
}
\item[\texttt{U084}] \sourcepath{docs/pyqcd\_reference/胶子PDF连续极限幻灯片.tex}
\par\noindent{\scriptsize
\textcolor{qcdsage}{伴生：}\sourcepath{docs/pyqcd\_reference/胶子PDF连续极限幻灯片.pdf}。
类型：TeX + PDF；规模：1576 行 / 50 页；目标：胶子算符、TMD 与自旋。
}
\item[\texttt{U085}] \sourcepath{docs/pyqcd\_reference/胶子PDF连续极限论文.tex}
\par\noindent{\scriptsize
\textcolor{qcdsage}{伴生：}\sourcepath{docs/pyqcd\_reference/胶子PDF连续极限论文.pdf}。
类型：TeX + PDF；规模：573 行 / 16 页；目标：胶子算符、TMD 与自旋。
}
\item[\texttt{U086}] \sourcepath{docs/pyqcd\_reference/胶子TMD梯度流连续极限.tex}
\par\noindent{\scriptsize
\textcolor{qcdsage}{伴生：}\sourcepath{docs/pyqcd\_reference/胶子TMD梯度流连续极限.pdf}。
类型：TeX + PDF；规模：1481 行 / 35 页；目标：胶子算符、TMD 与自旋。
}
\item[\texttt{U087}] \sourcepath{docs/pyqcd\_reference/质子自旋危机解析.tex}
\par\noindent{\scriptsize
\textcolor{qcdsage}{伴生：}\sourcepath{docs/pyqcd\_reference/质子自旋危机解析.pdf}。
类型：TeX + PDF；规模：564 行 / 14 页；目标：胶子算符、TMD 与自旋。
}
\end{description}
\section{逐文件哈希与规模账本}
本账本以文件而不是知识单元为粒度，确保 PDF 见证件、历史版本和 Markdown 工作记录没有被去重后从审计记录中消失。
\begin{description}[leftmargin=4.2em,labelwidth=3.7em,style=nextline,itemsep=0.18em]
\item[\texttt{001}] \sourcepath{docs/pyqcd\_reference/PyQCD仓库结构与核心物理链解析.pdf}
\par\noindent{\scriptsize PDF 见证；规模：4 页；字节：206\,963；SHA-256：\texttt{2d09cac620d75bf8}\allowbreak\texttt{7b01967fe155f35e}\allowbreak\texttt{65906959656fbf1a}\allowbreak\texttt{f8693cec45cc3318}。}
\item[\texttt{002}] \sourcepath{docs/pyqcd\_reference/PyQCD仓库结构与核心物理链解析.tex}
\par\noindent{\scriptsize TeX 正文；规模：166 行；字节：8\,336；SHA-256：\texttt{134ce0210d6793fc}\allowbreak\texttt{feda59b9fce7fdb4}\allowbreak\texttt{797a5031e300fd4d}\allowbreak\texttt{6905c116e960eef6}。}
\item[\texttt{003}] \sourcepath{docs/pyqcd\_reference/absorb\_dev20\_gluon\_tmd\_audit\_20261022.md}
\par\noindent{\scriptsize Markdown；规模：378 行；字节：18\,531；SHA-256：\texttt{4ee378fef6d34d8a}\allowbreak\texttt{54d91ded0172113c}\allowbreak\texttt{0b0e18d5ccbc02fb}\allowbreak\texttt{35b4049162c550ca}。}
\item[\texttt{004}] \sourcepath{docs/pyqcd\_reference/analy\_gpd\_20260928.pdf}
\par\noindent{\scriptsize PDF 见证；规模：21 页；字节：635\,062；SHA-256：\texttt{2b4c554ef47f5022}\allowbreak\texttt{f8f710058abd0f2f}\allowbreak\texttt{9a82d24dd5097613}\allowbreak\texttt{cf03b5e35d8dd14e}。}
\item[\texttt{005}] \sourcepath{docs/pyqcd\_reference/analy\_gpd\_20260928.tex}
\par\noindent{\scriptsize TeX 正文；规模：937 行；字节：45\,915；SHA-256：\texttt{ad2e68e844b6cf05}\allowbreak\texttt{26e0c2253ed0fb38}\allowbreak\texttt{0818763d6d576d00}\allowbreak\texttt{9719a6120799e5d1}。}
\item[\texttt{006}] \sourcepath{docs/pyqcd\_reference/analy\_physics\_chain\_20260824.pdf}
\par\noindent{\scriptsize PDF 见证；规模：29 页；字节：808\,442；SHA-256：\texttt{5a55abe4a7f4e1b0}\allowbreak\texttt{4fc26b934d7a68e6}\allowbreak\texttt{f81d564c704c080d}\allowbreak\texttt{e14d48e245cc2b71}。}
\item[\texttt{007}] \sourcepath{docs/pyqcd\_reference/analy\_physics\_chain\_20260824.tex}
\par\noindent{\scriptsize TeX 正文；规模：975 行；字节：56\,818；SHA-256：\texttt{9e5e3a1ff9ae08c5}\allowbreak\texttt{93222bc7827953eb}\allowbreak\texttt{23455a5f0f81e653}\allowbreak\texttt{65d1439824c7362b}。}
\item[\texttt{008}] \sourcepath{docs/pyqcd\_reference/analy\_proton\_energy\_20260816.tex}
\par\noindent{\scriptsize 独立 TeX；规模：300 行；字节：17\,301；SHA-256：\texttt{1c1be3bbfa90e622}\allowbreak\texttt{75b3f5d90a5d6ab1}\allowbreak\texttt{0ea7f71d2951dfb2}\allowbreak\texttt{068489135e7527d6}。}
\item[\texttt{009}] \sourcepath{docs/pyqcd\_reference/analy\_pyqcd\_20260818.pdf}
\par\noindent{\scriptsize PDF 见证；规模：11 页；字节：379\,821；SHA-256：\texttt{9b0dbc554d95d51e}\allowbreak\texttt{f24c90980c96e17c}\allowbreak\texttt{bd2c62ca7f4e7132}\allowbreak\texttt{cba1c6dd25b2526a}。}
\item[\texttt{010}] \sourcepath{docs/pyqcd\_reference/analy\_pyqcd\_20260818.tex}
\par\noindent{\scriptsize TeX 正文；规模：324 行；字节：18\,759；SHA-256：\texttt{fdacbc4e78a71fc9}\allowbreak\texttt{739a894fd69f7519}\allowbreak\texttt{186ce5504c92f9be}\allowbreak\texttt{ed10d082bfacd383}。}
\item[\texttt{011}] \sourcepath{docs/pyqcd\_reference/analy\_torch\_parallel\_20260816.pdf}
\par\noindent{\scriptsize PDF 见证；规模：9 页；字节：318\,869；SHA-256：\texttt{b88ab30252bf826f}\allowbreak\texttt{ea418721e4eeadf5}\allowbreak\texttt{176b1e22a910ec30}\allowbreak\texttt{7dde07ebbc6757c3}。}
\item[\texttt{012}] \sourcepath{docs/pyqcd\_reference/analy\_torch\_parallel\_20260816.tex}
\par\noindent{\scriptsize TeX 正文；规模：442 行；字节：21\,987；SHA-256：\texttt{cae2a64f7616b320}\allowbreak\texttt{cf4f213d8d938e37}\allowbreak\texttt{c48971f7ade562d8}\allowbreak\texttt{10126b79158340c2}。}
\item[\texttt{013}] \sourcepath{docs/pyqcd\_reference/from-logs/README.md}
\par\noindent{\scriptsize Markdown；规模：7 行；字节：327；SHA-256：\texttt{fc0b384080783923}\allowbreak\texttt{644288eb7864dcbf}\allowbreak\texttt{416cd9493f76e430}\allowbreak\texttt{bf11c125a1a08bd2}。}
\item[\texttt{014}] \sourcepath{docs/pyqcd\_reference/from-logs/logs/dev5/analy\_test9\_20260820.pdf}
\par\noindent{\scriptsize PDF 见证；规模：25 页；字节：1\,202\,601；SHA-256：\texttt{9e5d8f06527bace0}\allowbreak\texttt{2adc3195ab3a4a71}\allowbreak\texttt{4822e6183330abe3}\allowbreak\texttt{984d345786e9a534}。}
\item[\texttt{015}] \sourcepath{docs/pyqcd\_reference/from-logs/logs/dev5/analy\_test9\_20260820.tex}
\par\noindent{\scriptsize TeX 正文；规模：818 行；字节：64\,359；SHA-256：\texttt{52bd8bf496a7a4b0}\allowbreak\texttt{9738c1ec5259c265}\allowbreak\texttt{6712fca9b7746f8f}\allowbreak\texttt{e490bec1589d93ca}。}
\item[\texttt{016}] \sourcepath{docs/pyqcd\_reference/from-logs/logs/dev5\_1/analy\_test9\_20260821.pdf}
\par\noindent{\scriptsize PDF 见证；规模：33 页；字节：1\,245\,025；SHA-256：\texttt{a1572ea6320000a6}\allowbreak\texttt{ff8a6a6254078f52}\allowbreak\texttt{d89df71bc0c96453}\allowbreak\texttt{a8d0ae4989461572}。}
\item[\texttt{017}] \sourcepath{docs/pyqcd\_reference/from-logs/logs/dev5\_1/analy\_test9\_20260821.tex}
\par\noindent{\scriptsize TeX 正文；规模：863 行；字节：69\,355；SHA-256：\texttt{8bcb6539891ea2a2}\allowbreak\texttt{ae7e457520483476}\allowbreak\texttt{62ed1f96b7dd0ac0}\allowbreak\texttt{9bda309b145bba59}。}
\item[\texttt{018}] \sourcepath{docs/pyqcd\_reference/from-logs/logs/dev5\_2/analy\_test9\_20260822.pdf}
\par\noindent{\scriptsize PDF 见证；规模：36 页；字节：1\,280\,233；SHA-256：\texttt{d15cb8be273601f7}\allowbreak\texttt{129f99cbcba19e25}\allowbreak\texttt{988e81ce11b760ef}\allowbreak\texttt{c12785fff2b845cc}。}
\item[\texttt{019}] \sourcepath{docs/pyqcd\_reference/from-logs/logs/dev5\_2/analy\_test9\_20260822.tex}
\par\noindent{\scriptsize TeX 正文；规模：901 行；字节：74\,317；SHA-256：\texttt{8d0dc42c9031ed40}\allowbreak\texttt{3e2c4c739bb0a135}\allowbreak\texttt{7a03a4f528090fb5}\allowbreak\texttt{b2a1bc2f34b6d306}。}
\item[\texttt{020}] \sourcepath{docs/pyqcd\_reference/from-logs/logs/donghx\_4150\_20260829/pdfcheck.md}
\par\noindent{\scriptsize Markdown；规模：70 行；字节：2\,517；SHA-256：\texttt{c084c01cb5831f69}\allowbreak\texttt{dff5a0bed62685a9}\allowbreak\texttt{15c610d735b53371}\allowbreak\texttt{7679b2304e94595e}。}
\item[\texttt{021}] \sourcepath{docs/pyqcd\_reference/from-logs/logs/donghx\_4150\_20260829/summary.md}
\par\noindent{\scriptsize Markdown；规模：278 行；字节：14\,427；SHA-256：\texttt{1a21af7865602318}\allowbreak\texttt{f86c2b296f668bf7}\allowbreak\texttt{6b36b77876defb17}\allowbreak\texttt{341b34d284eda38b}。}
\item[\texttt{022}] \sourcepath{docs/pyqcd\_reference/from-logs/logs/stab0/stab0.pdf}
\par\noindent{\scriptsize PDF 见证；规模：13 页；字节：2\,463\,407；SHA-256：\texttt{1c82b836b99db690}\allowbreak\texttt{685e481b61fcaff6}\allowbreak\texttt{6138ae68d5c4cc62}\allowbreak\texttt{f9d8309af9d3dc7a}。}
\item[\texttt{023}] \sourcepath{docs/pyqcd\_reference/from-logs/logs/stab0/stab0.tex}
\par\noindent{\scriptsize TeX 正文；规模：589 行；字节：28\,489；SHA-256：\texttt{1b127b5b56560cf1}\allowbreak\texttt{d355f8d4fff0c4fd}\allowbreak\texttt{fe3c79498de0ab17}\allowbreak\texttt{280b489d7478fec0}。}
\item[\texttt{024}] \sourcepath{docs/pyqcd\_reference/from-logs/logs/stab1/docs/stab1\_analysis.pdf}
\par\noindent{\scriptsize PDF 见证；规模：6 页；字节：497\,785；SHA-256：\texttt{fb98c1ef65a170e9}\allowbreak\texttt{362f04bb991d86a9}\allowbreak\texttt{70eb68c1352160ef}\allowbreak\texttt{bcc3a5fcd6c04eec}。}
\item[\texttt{025}] \sourcepath{docs/pyqcd\_reference/from-logs/logs/stab1/docs/stab1\_analysis.tex}
\par\noindent{\scriptsize TeX 正文；规模：150 行；字节：7\,900；SHA-256：\texttt{08914a4fd86d1c6c}\allowbreak\texttt{8c425ac074846344}\allowbreak\texttt{61ae2ea8236813a8}\allowbreak\texttt{31f85df07066602f}。}
\item[\texttt{026}] \sourcepath{docs/pyqcd\_reference/from-logs/logs/stab1/v202608151614/stab1\_report.pdf}
\par\noindent{\scriptsize PDF 见证；规模：6 页；字节：538\,176；SHA-256：\texttt{7d5a713e6a56c8f3}\allowbreak\texttt{f4e02b94e59ad3ef}\allowbreak\texttt{a5fb4e6498665865}\allowbreak\texttt{eead8d964af53281}。}
\item[\texttt{027}] \sourcepath{docs/pyqcd\_reference/from-logs/logs/stab1/v202608151614/stab1\_report.tex}
\par\noindent{\scriptsize TeX 正文；规模：74 行；字节：2\,855；SHA-256：\texttt{3bc2467fa08b5f42}\allowbreak\texttt{1e5e3a2fb072306a}\allowbreak\texttt{984ccf0fe2687fe6}\allowbreak\texttt{d7f12aa1f30eeef9}。}
\item[\texttt{028}] \sourcepath{docs/pyqcd\_reference/from-logs/logs/test6/docs/analy\_test6\_report\_20260816.pdf}
\par\noindent{\scriptsize PDF 见证；规模：11 页；字节：674\,241；SHA-256：\texttt{a40e002a496987b6}\allowbreak\texttt{f7e1951a55cc975e}\allowbreak\texttt{9fba74e731585deb}\allowbreak\texttt{1f71a49ddc5c5925}。}
\item[\texttt{029}] \sourcepath{docs/pyqcd\_reference/from-logs/logs/test6/docs/analy\_test6\_report\_20260816.tex}
\par\noindent{\scriptsize TeX 正文；规模：412 行；字节：20\,924；SHA-256：\texttt{a36d4b185f94cc4c}\allowbreak\texttt{3b32b27472bc3bf3}\allowbreak\texttt{92d3abf531d43540}\allowbreak\texttt{a02ef7d01d175172}。}
\item[\texttt{030}] \sourcepath{docs/pyqcd\_reference/from-logs/logs/test7/v202608161007/test7\_report.pdf}
\par\noindent{\scriptsize PDF 见证；规模：2 页；字节：78\,209；SHA-256：\texttt{506531296ca5f609}\allowbreak\texttt{1118dfa9e9c53c29}\allowbreak\texttt{75e76ef019edb651}\allowbreak\texttt{d274ff0630c0beb6}。}
\item[\texttt{031}] \sourcepath{docs/pyqcd\_reference/from-logs/logs/test7/v202608161007/test7\_report.tex}
\par\noindent{\scriptsize TeX 正文；规模：50 行；字节：1\,736；SHA-256：\texttt{a463c7e2b3f105f8}\allowbreak\texttt{944a6459aeb88bc5}\allowbreak\texttt{4013ee281973f12b}\allowbreak\texttt{61e091880475fda0}。}
\item[\texttt{032}] \sourcepath{docs/pyqcd\_reference/from-logs/logs/test8/v202608161117/test8\_report.pdf}
\par\noindent{\scriptsize PDF 见证；规模：5 页；字节：439\,765；SHA-256：\texttt{b7aeac1d21c5d047}\allowbreak\texttt{67fe1580d9473996}\allowbreak\texttt{3add55bd1b0b8304}\allowbreak\texttt{fbce8e8679d330b9}。}
\item[\texttt{033}] \sourcepath{docs/pyqcd\_reference/from-logs/logs/test8/v202608161117/test8\_report.tex}
\par\noindent{\scriptsize TeX 正文；规模：71 行；字节：2\,724；SHA-256：\texttt{ea0f78463994d769}\allowbreak\texttt{d151528811d833fb}\allowbreak\texttt{acd922e249054fcc}\allowbreak\texttt{94b68c5dbba1f115}。}
\item[\texttt{034}] \sourcepath{docs/pyqcd\_reference/from-logs/logs/test9/test9\_analysis.pdf}
\par\noindent{\scriptsize PDF 见证；规模：8 页；字节：692\,938；SHA-256：\texttt{58fba31e5ef000eb}\allowbreak\texttt{a8eebac8e71e7bc9}\allowbreak\texttt{bdd35a4cac8a9f7f}\allowbreak\texttt{e325ca16acc4959a}。}
\item[\texttt{035}] \sourcepath{docs/pyqcd\_reference/from-logs/logs/test9/test9\_analysis.tex}
\par\noindent{\scriptsize TeX 正文；规模：233 行；字节：11\,279；SHA-256：\texttt{74cdf8723a9d7f39}\allowbreak\texttt{45253c0c439b0b96}\allowbreak\texttt{7c56021ad65b7abd}\allowbreak\texttt{909da7a1b439b710}。}
\item[\texttt{036}] \sourcepath{docs/pyqcd\_reference/gevp\_convergence.pdf}
\par\noindent{\scriptsize 独立 PDF；规模：1 页；字节：17\,355；SHA-256：\texttt{32940b99cbd21df5}\allowbreak\texttt{8aa0d34a788accca}\allowbreak\texttt{4c44e8518f1ea83b}\allowbreak\texttt{04dcab5f3e5729c2}。}
\item[\texttt{037}] \sourcepath{docs/pyqcd\_reference/gevp\_t0\_suppression.pdf}
\par\noindent{\scriptsize 独立 PDF；规模：1 页；字节：19\,451；SHA-256：\texttt{3405d3166c006178}\allowbreak\texttt{561532d2dc3c4f5d}\allowbreak\texttt{5b686a140618402f}\allowbreak\texttt{190614409a7ee609}。}
\item[\texttt{038}] \sourcepath{docs/pyqcd\_reference/gradient\_flow.md}
\par\noindent{\scriptsize Markdown；规模：895 行；字节：32\,498；SHA-256：\texttt{b7d8a3e7dc42cd30}\allowbreak\texttt{e89d07c590a19ab0}\allowbreak\texttt{c6063cfb26addb61}\allowbreak\texttt{0dc853f1c8f8a506}。}
\item[\texttt{039}] \sourcepath{docs/pyqcd\_reference/pure\_pyqcd\_20260818.pdf}
\par\noindent{\scriptsize PDF 见证；规模：10 页；字节：340\,582；SHA-256：\texttt{2c0ec28e4a5ec5bb}\allowbreak\texttt{4522897d4d875dc4}\allowbreak\texttt{a0b53bf4698e29d9}\allowbreak\texttt{1055c2c54e5e9d08}。}
\item[\texttt{040}] \sourcepath{docs/pyqcd\_reference/pure\_pyqcd\_20260818.tex}
\par\noindent{\scriptsize TeX 正文；规模：304 行；字节：17\,527；SHA-256：\texttt{cd538e7a7dedd8bb}\allowbreak\texttt{e7a79e57f0ef6908}\allowbreak\texttt{9b21e91f2a1552d5}\allowbreak\texttt{256a53cec084dce9}。}
\item[\texttt{041}] \sourcepath{docs/pyqcd\_reference/report\_cmp1\_4150\_20260828.pdf}
\par\noindent{\scriptsize PDF 见证；规模：16 页；字节：330\,758；SHA-256：\texttt{e03b57a097cca751}\allowbreak\texttt{bb56b50f5e52f687}\allowbreak\texttt{0517a9e4eea3dee3}\allowbreak\texttt{32f3209eeda46fff}。}
\item[\texttt{042}] \sourcepath{docs/pyqcd\_reference/report\_cmp1\_4150\_20260828.tex}
\par\noindent{\scriptsize TeX 正文；规模：507 行；字节：23\,921；SHA-256：\texttt{1cbc3d1c5af60faa}\allowbreak\texttt{03f58139bdb4669e}\allowbreak\texttt{3c73ab1ff165ec50}\allowbreak\texttt{d70922f01fd4af4f}。}
\item[\texttt{043}] \sourcepath{docs/pyqcd\_reference/report\_donghx\_4150\_reproduction\_20260829.pdf}
\par\noindent{\scriptsize PDF 见证；规模：23 页；字节：362\,452；SHA-256：\texttt{ef8d98a11e3e7c98}\allowbreak\texttt{ec0b39f4918bad78}\allowbreak\texttt{dfe3fa4fae7827cf}\allowbreak\texttt{26b026b57bee8ce1}。}
\item[\texttt{044}] \sourcepath{docs/pyqcd\_reference/report\_donghx\_4150\_reproduction\_20260829.tex}
\par\noindent{\scriptsize TeX 正文；规模：617 行；字节：32\,972；SHA-256：\texttt{e235ca69d223bd7f}\allowbreak\texttt{d67e27081e2ef59d}\allowbreak\texttt{d33225a8ac16e4b9}\allowbreak\texttt{ec7ef24509e0dec3}。}
\item[\texttt{045}] \sourcepath{docs/pyqcd\_reference/report\_donghx\_4150\_reproduction\_20260830.pdf}
\par\noindent{\scriptsize PDF 见证；规模：27 页；字节：391\,299；SHA-256：\texttt{f2a0cd8c2f708f97}\allowbreak\texttt{e8986cd07c975968}\allowbreak\texttt{3632f8a0573924e6}\allowbreak\texttt{b23cbc659d3ce4e4}。}
\item[\texttt{046}] \sourcepath{docs/pyqcd\_reference/report\_donghx\_4150\_reproduction\_20260830.tex}
\par\noindent{\scriptsize TeX 正文；规模：732 行；字节：39\,736；SHA-256：\texttt{39117f534b551b71}\allowbreak\texttt{1442cd45070fa354}\allowbreak\texttt{7b68be308c68f0e2}\allowbreak\texttt{d69ffba29611c9a7}。}
\item[\texttt{047}] \sourcepath{docs/pyqcd\_reference/report\_donghx\_pyqcd\_compare\_20260830.pdf}
\par\noindent{\scriptsize PDF 见证；规模：12 页；字节：256\,862；SHA-256：\texttt{8a93ababfe9de0de}\allowbreak\texttt{c1aa16d63c3a6543}\allowbreak\texttt{4045368880996148}\allowbreak\texttt{9414f3300f96ea7d}。}
\item[\texttt{048}] \sourcepath{docs/pyqcd\_reference/report\_donghx\_pyqcd\_compare\_20260830.tex}
\par\noindent{\scriptsize TeX 正文；规模：361 行；字节：19\,419；SHA-256：\texttt{33d203ebe5314e32}\allowbreak\texttt{7975903226a0c62b}\allowbreak\texttt{f44b9af6c0962815}\allowbreak\texttt{03677f15215e2c48}。}
\item[\texttt{049}] \sourcepath{docs/pyqcd\_reference/report\_donghx\_pyqcd\_compare\_20260830\_2.pdf}
\par\noindent{\scriptsize PDF 见证；规模：10 页；字节：151\,746；SHA-256：\texttt{ad1b560a8be2bf97}\allowbreak\texttt{7468a42b6cd9baa2}\allowbreak\texttt{39064d22c9fcc9a2}\allowbreak\texttt{eb8921b646fde207}。}
\item[\texttt{050}] \sourcepath{docs/pyqcd\_reference/report\_donghx\_pyqcd\_compare\_20260830\_2.tex}
\par\noindent{\scriptsize TeX 正文；规模：351 行；字节：20\,056；SHA-256：\texttt{789ceb3160b1a881}\allowbreak\texttt{d6a4ff8313cc50c7}\allowbreak\texttt{f037385ded8ca1ce}\allowbreak\texttt{423714b9671b873f}。}
\item[\texttt{051}] \sourcepath{docs/pyqcd\_reference/report\_gluon\_tmd\_gradient\_flow\_20260828.pdf}
\par\noindent{\scriptsize PDF 见证；规模：22 页；字节：657\,830；SHA-256：\texttt{47038e186fce679f}\allowbreak\texttt{f7b2a098f24ffffc}\allowbreak\texttt{3b4fa34ce303afab}\allowbreak\texttt{d8d4a5672c525c87}。}
\item[\texttt{052}] \sourcepath{docs/pyqcd\_reference/report\_gluon\_tmd\_gradient\_flow\_20260828.tex}
\par\noindent{\scriptsize TeX 正文；规模：789 行；字节：40\,334；SHA-256：\texttt{c8201a1fff3d9bb4}\allowbreak\texttt{54af0c3776b2ad6b}\allowbreak\texttt{e53b8f9f0d5dd312}\allowbreak\texttt{c54209d56fac0713}。}
\item[\texttt{053}] \sourcepath{docs/pyqcd\_reference/report\_gradient\_flow\_gluonpdf\_pyqcd\_20260909.pdf}
\par\noindent{\scriptsize PDF 见证；规模：17 页；字节：450\,305；SHA-256：\texttt{82507699c084eed4}\allowbreak\texttt{588834550cbd2335}\allowbreak\texttt{e586fbfc5c4ab97f}\allowbreak\texttt{114ab490f155dc13}。}
\item[\texttt{054}] \sourcepath{docs/pyqcd\_reference/report\_gradient\_flow\_gluonpdf\_pyqcd\_20260909.tex}
\par\noindent{\scriptsize TeX 正文；规模：1029 行；字节：37\,255；SHA-256：\texttt{a05bc2ba629d23da}\allowbreak\texttt{750264fc5d970aa2}\allowbreak\texttt{9e8a99332d8b1302}\allowbreak\texttt{2a51e69ed6f33ff8}。}
\item[\texttt{055}] \sourcepath{docs/pyqcd\_reference/基于Gradient\_Flow的质子胶子PDF格点计算整体流程.md}
\par\noindent{\scriptsize Markdown；规模：661 行；字节：26\,650；SHA-256：\texttt{a5c1ac4cbe14a543}\allowbreak\texttt{78c2803e9340d4d0}\allowbreak\texttt{704d6ee05f2a4196}\allowbreak\texttt{f7ce644ed6470924}。}
\item[\texttt{056}] \sourcepath{docs/pyqcd\_reference/基于格点QCD对部分子分布函数的研究.pdf}
\par\noindent{\scriptsize 独立 PDF；规模：56 页；字节：2\,466\,379；SHA-256：\texttt{cac44eac6456c007}\allowbreak\texttt{77cc5c4ccaee70a8}\allowbreak\texttt{aa63cbe467fec51e}\allowbreak\texttt{bebb5653fbbd7941}。}
\item[\texttt{057}] \sourcepath{docs/pyqcd\_reference/基于格点QCD的横向动量依赖的部分子分布函数的研究.pdf}
\par\noindent{\scriptsize 独立 PDF；规模：82 页；字节：9\,783\,095；SHA-256：\texttt{dac0a88df96c0d09}\allowbreak\texttt{03e903b319032ca7}\allowbreak\texttt{a270277fb44ef6ac}\allowbreak\texttt{3373813a1b28d2ea}。}
\item[\texttt{058}] \sourcepath{docs/pyqcd\_reference/构造胶子准算符.pdf}
\par\noindent{\scriptsize PDF 见证；规模：14 页；字节：391\,990；SHA-256：\texttt{fb02814754486d5e}\allowbreak\texttt{ea54f70bc88e5d27}\allowbreak\texttt{64563f12f0990427}\allowbreak\texttt{16c6341f00108f30}。}
\item[\texttt{059}] \sourcepath{docs/pyqcd\_reference/构造胶子准算符.tex}
\par\noindent{\scriptsize TeX 正文；规模：562 行；字节：31\,624；SHA-256：\texttt{a160f6f87cba567f}\allowbreak\texttt{f988b7f795477d64}\allowbreak\texttt{68d1bd4934d9bcf7}\allowbreak\texttt{a3b89d37eac51de6}。}
\item[\texttt{060}] \sourcepath{docs/pyqcd\_reference/格点QCD中的3pt构造.pdf}
\par\noindent{\scriptsize PDF 见证；规模：20 页；字节：619\,134；SHA-256：\texttt{d80a8423eb44d646}\allowbreak\texttt{b9f7a066ba77770d}\allowbreak\texttt{d7adf48cd983fdfb}\allowbreak\texttt{c9d2949d5a730997}。}
\item[\texttt{061}] \sourcepath{docs/pyqcd\_reference/格点QCD中的3pt构造.tex}
\par\noindent{\scriptsize TeX 正文；规模：783 行；字节：50\,957；SHA-256：\texttt{573884a330024d4e}\allowbreak\texttt{cafb8d8011a2f2c2}\allowbreak\texttt{3b1f4baef6b616f8}\allowbreak\texttt{e001a5d37306331a}。}
\item[\texttt{062}] \sourcepath{docs/pyqcd\_reference/格点QCD中的DGLAP演化方程.pdf}
\par\noindent{\scriptsize PDF 见证；规模：18 页；字节：528\,085；SHA-256：\texttt{e2c0bdcebf5eaffe}\allowbreak\texttt{b05064734235516e}\allowbreak\texttt{42ea3388847d7337}\allowbreak\texttt{eebf5247cfc2198a}。}
\item[\texttt{063}] \sourcepath{docs/pyqcd\_reference/格点QCD中的DGLAP演化方程.tex}
\par\noindent{\scriptsize TeX 正文；规模：688 行；字节：44\,598；SHA-256：\texttt{a43ca9182e9820bd}\allowbreak\texttt{b8a30558661daa9c}\allowbreak\texttt{26903bd516267f5d}\allowbreak\texttt{b442c554ada8d05e}。}
\item[\texttt{064}] \sourcepath{docs/pyqcd\_reference/格点QCD中的GEVP.pdf}
\par\noindent{\scriptsize PDF 见证；规模：24 页；字节：700\,272；SHA-256：\texttt{ab8efd3a047856fe}\allowbreak\texttt{6a04899417b33739}\allowbreak\texttt{cfdddd57c67a958e}\allowbreak\texttt{a7e4afbac6054c45}。}
\item[\texttt{065}] \sourcepath{docs/pyqcd\_reference/格点QCD中的GEVP.tex}
\par\noindent{\scriptsize TeX 正文；规模：924 行；字节：63\,489；SHA-256：\texttt{6de83102f1088059}\allowbreak\texttt{2f3bc75266dcc2d4}\allowbreak\texttt{8dae6c8ba1e8e4aa}\allowbreak\texttt{8cb8d53047b54d35}。}
\item[\texttt{066}] \sourcepath{docs/pyqcd\_reference/格点QCD中的OPE算符.pdf}
\par\noindent{\scriptsize PDF 见证；规模：14 页；字节：439\,579；SHA-256：\texttt{5c2f643305d95005}\allowbreak\texttt{41489e1e1353a1ac}\allowbreak\texttt{80ea9a48fd2b9018}\allowbreak\texttt{c28ee3c700579271}。}
\item[\texttt{067}] \sourcepath{docs/pyqcd\_reference/格点QCD中的OPE算符.tex}
\par\noindent{\scriptsize TeX 正文；规模：526 行；字节：32\,875；SHA-256：\texttt{3f80d0732fe8e181}\allowbreak\texttt{44b1f3a497f43d75}\allowbreak\texttt{dec7891cc6ac2092}\allowbreak\texttt{3a2f1eb7d094d56b}。}
\item[\texttt{068}] \sourcepath{docs/pyqcd\_reference/格点QCD中的Symanzik有效理论.pdf}
\par\noindent{\scriptsize PDF 见证；规模：13 页；字节：415\,796；SHA-256：\texttt{732f087cd6f3ff57}\allowbreak\texttt{e9a726c0a3b6ba6a}\allowbreak\texttt{12c87905d5e67fa3}\allowbreak\texttt{6b8ba79883892d5f}。}
\item[\texttt{069}] \sourcepath{docs/pyqcd\_reference/格点QCD中的Symanzik有效理论.tex}
\par\noindent{\scriptsize TeX 正文；规模：449 行；字节：29\,755；SHA-256：\texttt{f8abdc1c4aa3f7da}\allowbreak\texttt{3960ea55b3af5a27}\allowbreak\texttt{49f5080f45eff859}\allowbreak\texttt{d6573c06c5fd43e5}。}
\item[\texttt{070}] \sourcepath{docs/pyqcd\_reference/格点QCD中的TMD\_PDF.pdf}
\par\noindent{\scriptsize PDF 见证；规模：30 页；字节：645\,661；SHA-256：\texttt{832ab29c59f2fec2}\allowbreak\texttt{e28f57954cb73ae0}\allowbreak\texttt{c80bd6016865b544}\allowbreak\texttt{2b364d3624af4f3d}。}
\item[\texttt{071}] \sourcepath{docs/pyqcd\_reference/格点QCD中的TMD\_PDF.tex}
\par\noindent{\scriptsize TeX 正文；规模：937 行；字节：72\,101；SHA-256：\texttt{1a785c56d9dd6bad}\allowbreak\texttt{e13b6e5b390d0c08}\allowbreak\texttt{ae7fc5c891973cee}\allowbreak\texttt{a1b89a27982f02d0}。}
\item[\texttt{072}] \sourcepath{docs/pyqcd\_reference/格点QCD中的UV涨落.pdf}
\par\noindent{\scriptsize PDF 见证；规模：13 页；字节：438\,110；SHA-256：\texttt{03fdf94337ff02d1}\allowbreak\texttt{8835fdc5d8e62b51}\allowbreak\texttt{3f3e40449adef1f9}\allowbreak\texttt{8e50375250702c5c}。}
\item[\texttt{073}] \sourcepath{docs/pyqcd\_reference/格点QCD中的UV涨落.tex}
\par\noindent{\scriptsize TeX 正文；规模：472 行；字节：30\,337；SHA-256：\texttt{00a7feb7a3ef32a1}\allowbreak\texttt{b5d7f7b3c893df78}\allowbreak\texttt{3dd640adf28ed587}\allowbreak\texttt{6cb5cfa00e0cb7fe}。}
\item[\texttt{074}] \sourcepath{docs/pyqcd\_reference/格点QCD中的Wilson线.pdf}
\par\noindent{\scriptsize PDF 见证；规模：21 页；字节：559\,057；SHA-256：\texttt{ed84ee969e9f1697}\allowbreak\texttt{8f258cd5e3046cb2}\allowbreak\texttt{ba512f2c6b3eb164}\allowbreak\texttt{a3755a88109a83de}。}
\item[\texttt{075}] \sourcepath{docs/pyqcd\_reference/格点QCD中的Wilson线.tex}
\par\noindent{\scriptsize TeX 正文；规模：712 行；字节：51\,580；SHA-256：\texttt{6ebeae97d4d34fe0}\allowbreak\texttt{8b83f694deb96cab}\allowbreak\texttt{240069bf05d0cadc}\allowbreak\texttt{e27e3d9f4c98831f}。}
\item[\texttt{076}] \sourcepath{docs/pyqcd\_reference/格点QCD中的luscher参数化.pdf}
\par\noindent{\scriptsize PDF 见证；规模：17 页；字节：564\,536；SHA-256：\texttt{f409cc585cef82b8}\allowbreak\texttt{a7008c981c3d5f85}\allowbreak\texttt{fcf9734774e1c9c9}\allowbreak\texttt{122f3f0738259912}。}
\item[\texttt{077}] \sourcepath{docs/pyqcd\_reference/格点QCD中的luscher参数化.tex}
\par\noindent{\scriptsize TeX 正文；规模：603 行；字节：44\,514；SHA-256：\texttt{cea5b0e909fc4a62}\allowbreak\texttt{3bf298b8114e2aad}\allowbreak\texttt{71abc7cfe2141fbe}\allowbreak\texttt{5fdde5ed2ab69357}。}
\item[\texttt{078}] \sourcepath{docs/pyqcd\_reference/格点QCD中的smear算法.pdf}
\par\noindent{\scriptsize PDF 见证；规模：16 页；字节：455\,906；SHA-256：\texttt{2cbd9b2016fd246d}\allowbreak\texttt{11b3a93ac7420bac}\allowbreak\texttt{84bcc21230f1264b}\allowbreak\texttt{896ec4496c2354fe}。}
\item[\texttt{079}] \sourcepath{docs/pyqcd\_reference/格点QCD中的smear算法.tex}
\par\noindent{\scriptsize TeX 正文；规模：630 行；字节：36\,878；SHA-256：\texttt{30685c00813a42b7}\allowbreak\texttt{d4546db135448a0d}\allowbreak\texttt{326033eb68ed4b8f}\allowbreak\texttt{0b08673d31ae8e15}。}
\item[\texttt{080}] \sourcepath{docs/pyqcd\_reference/格点QCD中的传播子求解.pdf}
\par\noindent{\scriptsize PDF 见证；规模：10 页；字节：433\,173；SHA-256：\texttt{7b9d3cb555be3b6d}\allowbreak\texttt{0e584b39931cd4ff}\allowbreak\texttt{7683a4a136d9ab40}\allowbreak\texttt{b473c69a592e762d}。}
\item[\texttt{081}] \sourcepath{docs/pyqcd\_reference/格点QCD中的传播子求解.tex}
\par\noindent{\scriptsize TeX 正文；规模：269 行；字节：22\,693；SHA-256：\texttt{a4c82b384c16bdd9}\allowbreak\texttt{64dcbe7cb9583157}\allowbreak\texttt{6333aa2c8c3bff8a}\allowbreak\texttt{39020b752125af31}。}
\item[\texttt{082}] \sourcepath{docs/pyqcd\_reference/格点QCD中的光锥PDF、准PDF、赝PDF.pdf}
\par\noindent{\scriptsize PDF 见证；规模：20 页；字节：552\,380；SHA-256：\texttt{fb38862b59eb2d42}\allowbreak\texttt{b3aee9aae7de8530}\allowbreak\texttt{74fbbb30a2c66654}\allowbreak\texttt{9c4067688c4da387}。}
\item[\texttt{083}] \sourcepath{docs/pyqcd\_reference/格点QCD中的光锥PDF、准PDF、赝PDF.tex}
\par\noindent{\scriptsize TeX 正文；规模：705 行；字节：48\,477；SHA-256：\texttt{e45e4fb129745406}\allowbreak\texttt{5c04fd8d71096636}\allowbreak\texttt{c5aa7bb52319042f}\allowbreak\texttt{1cc4d5b3f8cd0d23}。}
\item[\texttt{084}] \sourcepath{docs/pyqcd\_reference/格点QCD中的光锥PDF与quasi\_PDF.pdf}
\par\noindent{\scriptsize PDF 见证；规模：12 页；字节：367\,168；SHA-256：\texttt{5d39e651bb6a7400}\allowbreak\texttt{ca2282032e1aafd3}\allowbreak\texttt{690a61fbd861dc0d}\allowbreak\texttt{0f88668f9f5bae11}。}
\item[\texttt{085}] \sourcepath{docs/pyqcd\_reference/格点QCD中的光锥PDF与quasi\_PDF.tex}
\par\noindent{\scriptsize TeX 正文；规模：413 行；字节：25\,659；SHA-256：\texttt{c6cd58cbda714e4c}\allowbreak\texttt{e0221b12caf92958}\allowbreak\texttt{ad66bb21a07de9f5}\allowbreak\texttt{0d06a64b13a8dd8a}。}
\item[\texttt{086}] \sourcepath{docs/pyqcd\_reference/格点QCD中的光锥与光前.pdf}
\par\noindent{\scriptsize PDF 见证；规模：18 页；字节：426\,135；SHA-256：\texttt{cd59c642241b34d0}\allowbreak\texttt{969540dcf636a10d}\allowbreak\texttt{bc1f473704ab9d31}\allowbreak\texttt{38bc125806ef1303}。}
\item[\texttt{087}] \sourcepath{docs/pyqcd\_reference/格点QCD中的光锥与光前.tex}
\par\noindent{\scriptsize TeX 正文；规模：783 行；字节：40\,923；SHA-256：\texttt{d6a20d7aadeabadf}\allowbreak\texttt{bdb5fb1edbd872bf}\allowbreak\texttt{f78871770d44c889}\allowbreak\texttt{cb808f500954f2e5}。}
\item[\texttt{088}] \sourcepath{docs/pyqcd\_reference/格点QCD中的关联函数.pdf}
\par\noindent{\scriptsize PDF 见证；规模：15 页；字节：449\,097；SHA-256：\texttt{5fe68e8c29cd0e61}\allowbreak\texttt{43b31f5046158f3f}\allowbreak\texttt{f35cd40058e2a7e3}\allowbreak\texttt{84e127f07d361231}。}
\item[\texttt{089}] \sourcepath{docs/pyqcd\_reference/格点QCD中的关联函数.tex}
\par\noindent{\scriptsize TeX 正文；规模：604 行；字节：34\,896；SHA-256：\texttt{b761f4a05ff3a6c2}\allowbreak\texttt{0284f0841ad63849}\allowbreak\texttt{8701319ecc52bd36}\allowbreak\texttt{d210e31f30df0ac9}。}
\item[\texttt{090}] \sourcepath{docs/pyqcd\_reference/格点QCD中的场强张量.pdf}
\par\noindent{\scriptsize PDF 见证；规模：14 页；字节：382\,178；SHA-256：\texttt{34ea3d20c45be2e1}\allowbreak\texttt{5717bcb6d4b09d48}\allowbreak\texttt{5eb07dc21d32deeb}\allowbreak\texttt{2e17acffb3d225a7}。}
\item[\texttt{091}] \sourcepath{docs/pyqcd\_reference/格点QCD中的场强张量.tex}
\par\noindent{\scriptsize TeX 正文；规模：576 行；字节：31\,156；SHA-256：\texttt{51c5846c5c8345bf}\allowbreak\texttt{d62d528596b856c8}\allowbreak\texttt{a45c16d0a973efb8}\allowbreak\texttt{8f2e79c18a444fc9}。}
\item[\texttt{092}] \sourcepath{docs/pyqcd\_reference/格点QCD中的外推.pdf}
\par\noindent{\scriptsize PDF 见证；规模：16 页；字节：491\,115；SHA-256：\texttt{c7968256a79eb20e}\allowbreak\texttt{d5b45dcc10c043b5}\allowbreak\texttt{0b54675a2780bb4f}\allowbreak\texttt{e369324bbfb2ccee}。}
\item[\texttt{093}] \sourcepath{docs/pyqcd\_reference/格点QCD中的外推.tex}
\par\noindent{\scriptsize TeX 正文；规模：567 行；字节：40\,966；SHA-256：\texttt{7e0d0e54965cf069}\allowbreak\texttt{c6c51346d8a617d3}\allowbreak\texttt{e1a91af1ef1e5817}\allowbreak\texttt{d4a8ca59ead7b443}。}
\item[\texttt{094}] \sourcepath{docs/pyqcd\_reference/格点QCD中的大动量有效理论.pdf}
\par\noindent{\scriptsize PDF 见证；规模：13 页；字节：415\,295；SHA-256：\texttt{a07c6a93b3c2c71a}\allowbreak\texttt{a1e312be0d5baeb0}\allowbreak\texttt{56ec08c00e421291}\allowbreak\texttt{9d078d389b241226}。}
\item[\texttt{095}] \sourcepath{docs/pyqcd\_reference/格点QCD中的大动量有效理论.tex}
\par\noindent{\scriptsize TeX 正文；规模：498 行；字节：29\,104；SHA-256：\texttt{ece25f53cba99342}\allowbreak\texttt{9fe85c19f418a750}\allowbreak\texttt{e66e0b982a800069}\allowbreak\texttt{309e9bb1ff4983aa}。}
\item[\texttt{096}] \sourcepath{docs/pyqcd\_reference/格点QCD中的夸克算符.pdf}
\par\noindent{\scriptsize PDF 见证；规模：13 页；字节：402\,811；SHA-256：\texttt{b848bf53dbe2df2c}\allowbreak\texttt{22f12ff8d43ed962}\allowbreak\texttt{c3902c83837ae544}\allowbreak\texttt{40e925dec1038e8b}。}
\item[\texttt{097}] \sourcepath{docs/pyqcd\_reference/格点QCD中的夸克算符.tex}
\par\noindent{\scriptsize TeX 正文；规模：457 行；字节：29\,049；SHA-256：\texttt{fc960e0f9fe18a12}\allowbreak\texttt{3e4523b7f1261b59}\allowbreak\texttt{9eb40d870d24420b}\allowbreak\texttt{122fe0dda2311e08}。}
\item[\texttt{098}] \sourcepath{docs/pyqcd\_reference/格点QCD中的常见名词短语与缩写.pdf}
\par\noindent{\scriptsize PDF 见证；规模：26 页；字节：532\,740；SHA-256：\texttt{23e63f872e382306}\allowbreak\texttt{f1d27756d20e953e}\allowbreak\texttt{4352e39f3c4af9da}\allowbreak\texttt{229eecc51ef607d4}。}
\item[\texttt{099}] \sourcepath{docs/pyqcd\_reference/格点QCD中的常见名词短语与缩写.tex}
\par\noindent{\scriptsize TeX 正文；规模：414 行；字节：55\,807；SHA-256：\texttt{5923a71cd886bf19}\allowbreak\texttt{cfb2e64913178985}\allowbreak\texttt{437df285a997a192}\allowbreak\texttt{887eb2e9606ab5ad}。}
\item[\texttt{100}] \sourcepath{docs/pyqcd\_reference/格点QCD中的常见费曼图.pdf}
\par\noindent{\scriptsize PDF 见证；规模：10 页；字节：349\,511；SHA-256：\texttt{37ca0f40e748ac16}\allowbreak\texttt{bd530576ba24d93d}\allowbreak\texttt{f0e6a2674c6cde1f}\allowbreak\texttt{9d543e42e9966347}。}
\item[\texttt{101}] \sourcepath{docs/pyqcd\_reference/格点QCD中的常见费曼图.tex}
\par\noindent{\scriptsize TeX 正文；规模：495 行；字节：26\,316；SHA-256：\texttt{bfce4234f827be27}\allowbreak\texttt{c4aa1ab9c0fba9c3}\allowbreak\texttt{677d4a7acd664bbe}\allowbreak\texttt{f912a57cbd60eda5}。}
\item[\texttt{102}] \sourcepath{docs/pyqcd\_reference/格点QCD中的广义部分子分布.pdf}
\par\noindent{\scriptsize PDF 见证；规模：25 页；字节：545\,329；SHA-256：\texttt{60486936257c8161}\allowbreak\texttt{a209160f5cd473c3}\allowbreak\texttt{070207c60d5cc039}\allowbreak\texttt{b73ab89462e22959}。}
\item[\texttt{103}] \sourcepath{docs/pyqcd\_reference/格点QCD中的广义部分子分布.tex}
\par\noindent{\scriptsize TeX 正文；规模：1354 行；字节：55\,674；SHA-256：\texttt{6436bbca893f8242}\allowbreak\texttt{46b9dd295c237b0d}\allowbreak\texttt{60c6760fcc32dd6f}\allowbreak\texttt{5935821b95c71fd2}。}
\item[\texttt{104}] \sourcepath{docs/pyqcd\_reference/格点QCD中的强子谱学.pdf}
\par\noindent{\scriptsize PDF 见证；规模：11 页；字节：439\,197；SHA-256：\texttt{b8872bb3cf9fabaa}\allowbreak\texttt{9ae5a80e43a9949d}\allowbreak\texttt{646376fb1d6388a9}\allowbreak\texttt{02f41e865146d1dd}。}
\item[\texttt{105}] \sourcepath{docs/pyqcd\_reference/格点QCD中的强子谱学.tex}
\par\noindent{\scriptsize TeX 正文；规模：353 行；字节：25\,165；SHA-256：\texttt{973ad5dcbd50830e}\allowbreak\texttt{d82df520be6dd76b}\allowbreak\texttt{d3ab64e783255ca3}\allowbreak\texttt{76238dc56f5cb0a3}。}
\item[\texttt{106}] \sourcepath{docs/pyqcd\_reference/格点QCD中的形状因子.pdf}
\par\noindent{\scriptsize PDF 见证；规模：14 页；字节：511\,026；SHA-256：\texttt{91da8f2e52c3ca24}\allowbreak\texttt{3d156f81ffd7955c}\allowbreak\texttt{2386418967d49a63}\allowbreak\texttt{ac9245fe7d2d75ac}。}
\item[\texttt{107}] \sourcepath{docs/pyqcd\_reference/格点QCD中的形状因子.tex}
\par\noindent{\scriptsize TeX 正文；规模：499 行；字节：36\,680；SHA-256：\texttt{4abc0fb2b9904a77}\allowbreak\texttt{7524058f9561bc3e}\allowbreak\texttt{89fe809a1a431b29}\allowbreak\texttt{4d1d325231f82e88}。}
\item[\texttt{108}] \sourcepath{docs/pyqcd\_reference/格点QCD中的标度参数.pdf}
\par\noindent{\scriptsize PDF 见证；规模：13 页；字节：415\,188；SHA-256：\texttt{a1f22b621bd77a14}\allowbreak\texttt{bbfa2554bd6a3e7f}\allowbreak\texttt{bc6a37d2e3269d2c}\allowbreak\texttt{73508b5046c5ceba}。}
\item[\texttt{109}] \sourcepath{docs/pyqcd\_reference/格点QCD中的标度参数.tex}
\par\noindent{\scriptsize TeX 正文；规模：456 行；字节：28\,707；SHA-256：\texttt{3523d6f1ff1fa017}\allowbreak\texttt{3079a6970592544d}\allowbreak\texttt{620989424a8e7130}\allowbreak\texttt{c87aabee537bb6e4}。}
\item[\texttt{110}] \sourcepath{docs/pyqcd\_reference/格点QCD中的梯度流重整化.pdf}
\par\noindent{\scriptsize PDF 见证；规模：15 页；字节：429\,035；SHA-256：\texttt{586a27b15c8e3bd1}\allowbreak\texttt{4330da6dbefd2c55}\allowbreak\texttt{90d49f6d1c1dd0b2}\allowbreak\texttt{11c518e447ad8803}。}
\item[\texttt{111}] \sourcepath{docs/pyqcd\_reference/格点QCD中的梯度流重整化.tex}
\par\noindent{\scriptsize TeX 正文；规模：626 行；字节：37\,208；SHA-256：\texttt{8717b2a859f7deec}\allowbreak\texttt{8714828e4325a730}\allowbreak\texttt{48ed9f71c0d5df5f}\allowbreak\texttt{b66404db56db44a8}。}
\item[\texttt{112}] \sourcepath{docs/pyqcd\_reference/格点QCD中的维克收缩与傅里叶变换.pdf}
\par\noindent{\scriptsize PDF 见证；规模：13 页；字节：425\,972；SHA-256：\texttt{046e5cc683b39afd}\allowbreak\texttt{08523201a018b738}\allowbreak\texttt{e7ecad5247f2c8a5}\allowbreak\texttt{c354a51a31dd4005}。}
\item[\texttt{113}] \sourcepath{docs/pyqcd\_reference/格点QCD中的维克收缩与傅里叶变换.tex}
\par\noindent{\scriptsize TeX 正文；规模：469 行；字节：29\,279；SHA-256：\texttt{28a3d50b1f48d74f}\allowbreak\texttt{b51f23b32c919f6b}\allowbreak\texttt{d8d1af3f7e857a1f}\allowbreak\texttt{621e388f9425709b}。}
\item[\texttt{114}] \sourcepath{docs/pyqcd\_reference/格点QCD中的胶子极化.pdf}
\par\noindent{\scriptsize PDF 见证；规模：16 页；字节：443\,418；SHA-256：\texttt{db5bd54dd07d031d}\allowbreak\texttt{f400d2f2438e88a9}\allowbreak\texttt{ee6e0cc41ef67f2f}\allowbreak\texttt{51d7a22e3792b67b}。}
\item[\texttt{115}] \sourcepath{docs/pyqcd\_reference/格点QCD中的胶子极化.tex}
\par\noindent{\scriptsize TeX 正文；规模：694 行；字节：38\,655；SHA-256：\texttt{1cef9c8e9fae72a1}\allowbreak\texttt{a43620e4e54affe8}\allowbreak\texttt{bc4d51348ef7936b}\allowbreak\texttt{9398d272636752f7}。}
\item[\texttt{116}] \sourcepath{docs/pyqcd\_reference/格点QCD中的胶子算符.pdf}
\par\noindent{\scriptsize PDF 见证；规模：15 页；字节：396\,500；SHA-256：\texttt{970ef249969bd25c}\allowbreak\texttt{3787a03d86105a2e}\allowbreak\texttt{eef736e6e08dcec6}\allowbreak\texttt{db7a2e6a6b183743}。}
\item[\texttt{117}] \sourcepath{docs/pyqcd\_reference/格点QCD中的胶子算符.tex}
\par\noindent{\scriptsize TeX 正文；规模：618 行；字节：34\,339；SHA-256：\texttt{644f1066ef957e02}\allowbreak\texttt{7216845aaa775430}\allowbreak\texttt{9e5b804ba32ff06f}\allowbreak\texttt{82b730674bcaf173}。}
\item[\texttt{118}] \sourcepath{docs/pyqcd\_reference/格点QCD中的色散关系.pdf}
\par\noindent{\scriptsize PDF 见证；规模：15 页；字节：504\,778；SHA-256：\texttt{d03a804df50ea2f2}\allowbreak\texttt{81cc8da217d49511}\allowbreak\texttt{a47a2471f4b60dfc}\allowbreak\texttt{61e99c8a78526456}。}
\item[\texttt{119}] \sourcepath{docs/pyqcd\_reference/格点QCD中的色散关系.tex}
\par\noindent{\scriptsize TeX 正文；规模：496 行；字节：37\,730；SHA-256：\texttt{0737abf85734aa31}\allowbreak\texttt{27bb8b9d4046b0f1}\allowbreak\texttt{5d3f5dd248b6d2f0}\allowbreak\texttt{43a21b7c210a53d6}。}
\item[\texttt{120}] \sourcepath{docs/pyqcd\_reference/格点QCD中的蒙卡方法.pdf}
\par\noindent{\scriptsize PDF 见证；规模：11 页；字节：462\,990；SHA-256：\texttt{444777355609bc8f}\allowbreak\texttt{6fbd17787e08833e}\allowbreak\texttt{0a2052859870f3f0}\allowbreak\texttt{cd3d810aebf0ab65}。}
\item[\texttt{121}] \sourcepath{docs/pyqcd\_reference/格点QCD中的蒙卡方法.tex}
\par\noindent{\scriptsize TeX 正文；规模：347 行；字节：28\,178；SHA-256：\texttt{a33d4f7d9e67cb93}\allowbreak\texttt{a1b2337bf4d1fc31}\allowbreak\texttt{e5c736a582ae6d96}\allowbreak\texttt{fe7a26d418dfc1ee}。}
\item[\texttt{122}] \sourcepath{docs/pyqcd\_reference/格点QCD中的衰变道.pdf}
\par\noindent{\scriptsize PDF 见证；规模：15 页；字节：512\,657；SHA-256：\texttt{373c1dc85d9ae551}\allowbreak\texttt{8be0df30128d0832}\allowbreak\texttt{d48545fa02b33b5a}\allowbreak\texttt{ddd5030f424582b6}。}
\item[\texttt{123}] \sourcepath{docs/pyqcd\_reference/格点QCD中的衰变道.tex}
\par\noindent{\scriptsize TeX 正文；规模：474 行；字节：37\,402；SHA-256：\texttt{a72d61b1ef1c9217}\allowbreak\texttt{bfd58efad717c822}\allowbreak\texttt{10ee05fd19b731e2}\allowbreak\texttt{e6d7ada7634a0aed}。}
\item[\texttt{124}] \sourcepath{docs/pyqcd\_reference/格点QCD中的误差统计.pdf}
\par\noindent{\scriptsize PDF 见证；规模：15 页；字节：490\,086；SHA-256：\texttt{e7d90bb41b7d7145}\allowbreak\texttt{304de93244ca5570}\allowbreak\texttt{813d3885aa73f597}\allowbreak\texttt{66d28f4e0edb8f2a}。}
\item[\texttt{125}] \sourcepath{docs/pyqcd\_reference/格点QCD中的误差统计.tex}
\par\noindent{\scriptsize TeX 正文；规模：472 行；字节：36\,291；SHA-256：\texttt{017f3acd57cdc449}\allowbreak\texttt{6ea24c08da1bf620}\allowbreak\texttt{028aae08fb2a3a42}\allowbreak\texttt{4f2d9423dd14422f}。}
\item[\texttt{126}] \sourcepath{docs/pyqcd\_reference/格点QCD中的费米子方案.pdf}
\par\noindent{\scriptsize PDF 见证；规模：24 页；字节：576\,602；SHA-256：\texttt{832e3f7f46ea0f83}\allowbreak\texttt{33bb4211d2a60053}\allowbreak\texttt{ebe537ea0a8c836e}\allowbreak\texttt{a84b731e9736e376}。}
\item[\texttt{127}] \sourcepath{docs/pyqcd\_reference/格点QCD中的费米子方案.tex}
\par\noindent{\scriptsize TeX 正文；规模：999 行；字节：61\,284；SHA-256：\texttt{a675522d85691d24}\allowbreak\texttt{5694fcfe6ddf5e82}\allowbreak\texttt{2c57aa52c26eb142}\allowbreak\texttt{3a191c77a27f6018}。}
\item[\texttt{128}] \sourcepath{docs/pyqcd\_reference/格点QCD中的连通图与非连通图.pdf}
\par\noindent{\scriptsize PDF 见证；规模：16 页；字节：537\,089；SHA-256：\texttt{46cb1b5de3a56184}\allowbreak\texttt{e78f635dbba24a30}\allowbreak\texttt{96a5fe1e0b7921f2}\allowbreak\texttt{72bfcb947f1d2158}。}
\item[\texttt{129}] \sourcepath{docs/pyqcd\_reference/格点QCD中的连通图与非连通图.tex}
\par\noindent{\scriptsize TeX 正文；规模：531 行；字节：42\,019；SHA-256：\texttt{4120d271fe16fdb9}\allowbreak\texttt{3d2deb344f93d908}\allowbreak\texttt{724b1070b0edc56c}\allowbreak\texttt{c5322b0585a94fa8}。}
\item[\texttt{130}] \sourcepath{docs/pyqcd\_reference/格点QCD中的部分子分布函数.pdf}
\par\noindent{\scriptsize PDF 见证；规模：16 页；字节：454\,948；SHA-256：\texttt{ec25264bf8ed9910}\allowbreak\texttt{72f1618c6870104f}\allowbreak\texttt{696c0c78df0f5904}\allowbreak\texttt{4a911c7f5a13746f}。}
\item[\texttt{131}] \sourcepath{docs/pyqcd\_reference/格点QCD中的部分子分布函数.tex}
\par\noindent{\scriptsize TeX 正文；规模：576 行；字节：34\,642；SHA-256：\texttt{ab83c2e090ab9639}\allowbreak\texttt{053ed7c0d75156d0}\allowbreak\texttt{3433a7328ca5bc73}\allowbreak\texttt{72e4b35493996091}。}
\item[\texttt{132}] \sourcepath{docs/pyqcd\_reference/格点QCD中的重夸克有效理论.pdf}
\par\noindent{\scriptsize PDF 见证；规模：12 页；字节：420\,832；SHA-256：\texttt{88682031dcc17573}\allowbreak\texttt{6c37133975604316}\allowbreak\texttt{cbfb658710595177}\allowbreak\texttt{5463c65d478515a5}。}
\item[\texttt{133}] \sourcepath{docs/pyqcd\_reference/格点QCD中的重夸克有效理论.tex}
\par\noindent{\scriptsize TeX 正文；规模：426 行；字节：27\,156；SHA-256：\texttt{7a7e801e6a3ae90a}\allowbreak\texttt{b7a1b0a199cdee71}\allowbreak\texttt{d961eac22a1d2581}\allowbreak\texttt{8e86d4aaf4bbff83}。}
\item[\texttt{134}] \sourcepath{docs/pyqcd\_reference/格点QCD中的重整化.pdf}
\par\noindent{\scriptsize PDF 见证；规模：21 页；字节：567\,098；SHA-256：\texttt{2ad98239dc1d640d}\allowbreak\texttt{86de9b1080f929d0}\allowbreak\texttt{53f05191dac715bb}\allowbreak\texttt{c90b09d058d27251}。}
\item[\texttt{135}] \sourcepath{docs/pyqcd\_reference/格点QCD中的重整化.tex}
\par\noindent{\scriptsize TeX 正文；规模：741 行；字节：52\,997；SHA-256：\texttt{59bb49dd32d10e4f}\allowbreak\texttt{9ebb580e7edf8be9}\allowbreak\texttt{ad7248f4379d7263}\allowbreak\texttt{6420d87f0f0b01da}。}
\item[\texttt{136}] \sourcepath{docs/pyqcd\_reference/格点QCD中的重采样方法.pdf}
\par\noindent{\scriptsize PDF 见证；规模：16 页；字节：517\,180；SHA-256：\texttt{aa90d94361af845c}\allowbreak\texttt{a2b0831eaf56e294}\allowbreak\texttt{620b311ace5a332b}\allowbreak\texttt{3a42f843153da3cf}。}
\item[\texttt{137}] \sourcepath{docs/pyqcd\_reference/格点QCD中的重采样方法.tex}
\par\noindent{\scriptsize TeX 正文；规模：507 行；字节：40\,132；SHA-256：\texttt{2dbc83ca816afaf4}\allowbreak\texttt{eea0e062f08bb583}\allowbreak\texttt{357c3b34599a8e64}\allowbreak\texttt{e054f653668be336}。}
\item[\texttt{138}] \sourcepath{docs/pyqcd\_reference/格点QCD蒸馏方法解析.pdf}
\par\noindent{\scriptsize PDF 见证；规模：14 页；字节：411\,181；SHA-256：\texttt{c1d3d4e36e46778b}\allowbreak\texttt{fe41036e73433f7d}\allowbreak\texttt{c448da662aa917a8}\allowbreak\texttt{5f8f52b095e9cdde}。}
\item[\texttt{139}] \sourcepath{docs/pyqcd\_reference/格点QCD蒸馏方法解析.tex}
\par\noindent{\scriptsize TeX 正文；规模：586 行；字节：31\,130；SHA-256：\texttt{ffe1ef95188fa8c6}\allowbreak\texttt{d2576b7a554046e7}\allowbreak\texttt{c14dfb0a455bff0c}\allowbreak\texttt{7faec4207db584b7}。}
\item[\texttt{140}] \sourcepath{docs/pyqcd\_reference/格点上计算胶子准算符.pdf}
\par\noindent{\scriptsize PDF 见证；规模：22 页；字节：555\,215；SHA-256：\texttt{a3bf39abebee5cfc}\allowbreak\texttt{66a5c211aca4f434}\allowbreak\texttt{97155e59ca985c8f}\allowbreak\texttt{4ed0c83c16f6baeb}。}
\item[\texttt{141}] \sourcepath{docs/pyqcd\_reference/格点上计算胶子准算符.tex}
\par\noindent{\scriptsize TeX 正文；规模：947 行；字节：49\,800；SHA-256：\texttt{4a141f1dd1b20aaf}\allowbreak\texttt{eeb73c94dd2467ad}\allowbreak\texttt{25da919f1b8bd900}\allowbreak\texttt{dd07c8785c1bac5f}。}
\item[\texttt{142}] \sourcepath{docs/pyqcd\_reference/理论解析与工作流-v20260911.pdf}
\par\noindent{\scriptsize PDF 见证；规模：100 页；字节：1\,373\,004；SHA-256：\texttt{ee275800365b5193}\allowbreak\texttt{0539d3a742f9312e}\allowbreak\texttt{4220c33841a42831}\allowbreak\texttt{3e0b8b141b16b34a}。}
\item[\texttt{143}] \sourcepath{docs/pyqcd\_reference/理论解析与工作流-v20260911.tex}
\par\noindent{\scriptsize TeX 正文；规模：5957 行；字节：300\,841；SHA-256：\texttt{699e9f3060838952}\allowbreak\texttt{1b0d2f82a8047f2d}\allowbreak\texttt{08ef9ed40bc1a60d}\allowbreak\texttt{9715e44a7796f1d3}。}
\item[\texttt{144}] \sourcepath{docs/pyqcd\_reference/理论解析与工作流.pdf}
\par\noindent{\scriptsize PDF 见证；规模：89 页；字节：1\,142\,493；SHA-256：\texttt{16601e44ccd3e852}\allowbreak\texttt{dc2763e405a8bbae}\allowbreak\texttt{a9e3898583dba10c}\allowbreak\texttt{d5cc7e0a864e8925}。}
\item[\texttt{145}] \sourcepath{docs/pyqcd\_reference/理论解析与工作流.tex}
\par\noindent{\scriptsize TeX 正文；规模：3837 行；字节：220\,708；SHA-256：\texttt{c9e5e0d25c50bee4}\allowbreak\texttt{f4c5220e380e0fe6}\allowbreak\texttt{835b947e771e35ab}\allowbreak\texttt{8d17046245963e0c}。}
\item[\texttt{146}] \sourcepath{docs/pyqcd\_reference/胶子PDF完整推导.pdf}
\par\noindent{\scriptsize PDF 见证；规模：15 页；字节：348\,290；SHA-256：\texttt{9e660d0139702b1c}\allowbreak\texttt{77b0b907987fd926}\allowbreak\texttt{a7f5b2b9cc078218}\allowbreak\texttt{0019d7d9bee22d77}。}
\item[\texttt{147}] \sourcepath{docs/pyqcd\_reference/胶子PDF完整推导.tex}
\par\noindent{\scriptsize TeX 正文；规模：742 行；字节：37\,897；SHA-256：\texttt{087e590d5818e495}\allowbreak\texttt{f0ac08aa5931bbd4}\allowbreak\texttt{e9bf2a5cb0436efa}\allowbreak\texttt{aef608e2094e355e}。}
\item[\texttt{148}] \sourcepath{docs/pyqcd\_reference/胶子PDF幻灯片.pdf}
\par\noindent{\scriptsize PDF 见证；规模：10 页；字节：219\,085；SHA-256：\texttt{1eb005812e1fd9e1}\allowbreak\texttt{233f9cabb0491251}\allowbreak\texttt{682f47575be78c2d}\allowbreak\texttt{830775c89947eed9}。}
\item[\texttt{149}] \sourcepath{docs/pyqcd\_reference/胶子PDF幻灯片.tex}
\par\noindent{\scriptsize TeX 正文；规模：357 行；字节：16\,128；SHA-256：\texttt{d2adf3c955782d99}\allowbreak\texttt{1ca1b142757d4308}\allowbreak\texttt{1d5e9d216cd7c799}\allowbreak\texttt{ac2b2dbb8dea7a6e}。}
\item[\texttt{150}] \sourcepath{docs/pyqcd\_reference/胶子PDF笔记.pdf}
\par\noindent{\scriptsize PDF 见证；规模：9 页；字节：283\,294；SHA-256：\texttt{7504e1c11dfd3412}\allowbreak\texttt{86a1333b572713ba}\allowbreak\texttt{02e03cf58ccac829}\allowbreak\texttt{9a3e24543574ebb8}。}
\item[\texttt{151}] \sourcepath{docs/pyqcd\_reference/胶子PDF笔记.tex}
\par\noindent{\scriptsize TeX 正文；规模：445 行；字节：22\,729；SHA-256：\texttt{f96be544e01665f9}\allowbreak\texttt{90019414e61d713e}\allowbreak\texttt{e587496fb1359a04}\allowbreak\texttt{ec718d7b52c04612}。}
\item[\texttt{152}] \sourcepath{docs/pyqcd\_reference/胶子PDF计算代码解析.pdf}
\par\noindent{\scriptsize PDF 见证；规模：30 页；字节：462\,296；SHA-256：\texttt{2ad33f213b61daae}\allowbreak\texttt{ba5426c428d58c58}\allowbreak\texttt{3a08e54e42559846}\allowbreak\texttt{8336bb945115c479}。}
\item[\texttt{153}] \sourcepath{docs/pyqcd\_reference/胶子PDF计算代码解析.tex}
\par\noindent{\scriptsize TeX 正文；规模：1218 行；字节：51\,688；SHA-256：\texttt{d12460de77c87586}\allowbreak\texttt{7bd2bad98eb1ec18}\allowbreak\texttt{19f4c8633ffc3aaf}\allowbreak\texttt{5c4e22d69a23151d}。}
\item[\texttt{154}] \sourcepath{docs/pyqcd\_reference/胶子PDF连续极限.pdf}
\par\noindent{\scriptsize PDF 见证；规模：16 页；字节：372\,848；SHA-256：\texttt{2061cd46250653c3}\allowbreak\texttt{843c51bb769d817c}\allowbreak\texttt{6401799a19ebe83e}\allowbreak\texttt{14354121402e05a3}。}
\item[\texttt{155}] \sourcepath{docs/pyqcd\_reference/胶子PDF连续极限.tex}
\par\noindent{\scriptsize TeX 正文；规模：578 行；字节：37\,759；SHA-256：\texttt{9533deecbd993ee2}\allowbreak\texttt{fc1ebba4bc43c3b4}\allowbreak\texttt{5f44663e3e4176f5}\allowbreak\texttt{e2bc89c6764e4515}。}
\item[\texttt{156}] \sourcepath{docs/pyqcd\_reference/胶子PDF连续极限幻灯片.pdf}
\par\noindent{\scriptsize PDF 见证；规模：50 页；字节：405\,678；SHA-256：\texttt{4c25a83f40752f64}\allowbreak\texttt{740a747ee56b6c3f}\allowbreak\texttt{cd3737732a36d94c}\allowbreak\texttt{23d8f6c8aac0d483}。}
\item[\texttt{157}] \sourcepath{docs/pyqcd\_reference/胶子PDF连续极限幻灯片.tex}
\par\noindent{\scriptsize TeX 正文；规模：1576 行；字节：64\,393；SHA-256：\texttt{1cbb0909185b1ccb}\allowbreak\texttt{fc5f7e6a653db68c}\allowbreak\texttt{515486bdbf1318b1}\allowbreak\texttt{c125c6930d29752b}。}
\item[\texttt{158}] \sourcepath{docs/pyqcd\_reference/胶子PDF连续极限论文.pdf}
\par\noindent{\scriptsize PDF 见证；规模：16 页；字节：371\,424；SHA-256：\texttt{565122646ad7873b}\allowbreak\texttt{50e2dc333f1e399c}\allowbreak\texttt{10bcfe7319e56e75}\allowbreak\texttt{103ea97790333ac6}。}
\item[\texttt{159}] \sourcepath{docs/pyqcd\_reference/胶子PDF连续极限论文.tex}
\par\noindent{\scriptsize TeX 正文；规模：573 行；字节：37\,479；SHA-256：\texttt{86160991a9ef06e3}\allowbreak\texttt{d699f7341ee9a0ad}\allowbreak\texttt{b2473342de2a5f9d}\allowbreak\texttt{695221577f0ecf88}。}
\item[\texttt{160}] \sourcepath{docs/pyqcd\_reference/胶子TMD梯度流连续极限.pdf}
\par\noindent{\scriptsize PDF 见证；规模：35 页；字节：585\,221；SHA-256：\texttt{fc507ef05992916e}\allowbreak\texttt{db6249977e702e76}\allowbreak\texttt{247119ed5fb2d0c6}\allowbreak\texttt{20c1e6dafcf443ae}。}
\item[\texttt{161}] \sourcepath{docs/pyqcd\_reference/胶子TMD梯度流连续极限.tex}
\par\noindent{\scriptsize TeX 正文；规模：1481 行；字节：77\,181；SHA-256：\texttt{34f56b3b909edd95}\allowbreak\texttt{64e4aa37416f97dd}\allowbreak\texttt{7d40b090c2b398b1}\allowbreak\texttt{e160951d0798f726}。}
\item[\texttt{162}] \sourcepath{docs/pyqcd\_reference/质子自旋危机解析.pdf}
\par\noindent{\scriptsize PDF 见证；规模：14 页；字节：416\,167；SHA-256：\texttt{229f03499407ad3e}\allowbreak\texttt{fb65322fe163cc8e}\allowbreak\texttt{5a651d961481f1b9}\allowbreak\texttt{8f6d64da42a35612}。}
\item[\texttt{163}] \sourcepath{docs/pyqcd\_reference/质子自旋危机解析.tex}
\par\noindent{\scriptsize TeX 正文；规模：564 行；字节：30\,901；SHA-256：\texttt{ba016d3810d81819}\allowbreak\texttt{1a41cfe3f1ea7dc2}\allowbreak\texttt{f8bfe5de09142feb}\allowbreak\texttt{7c8b426168be22b6}。}
\end{description}
\section{控制文件}
以下文件不进入物理内容计数，但决定快照边界和许可证解释。
\begin{description}[leftmargin=4.2em,labelwidth=3.7em,style=nextline,itemsep=0.25em]
\item[\texttt{C01}] \sourcepath{docs/pyqcd\_reference/AGENTS.md}
\par\noindent{\scriptsize 控制文件；规模：14 行；字节：857；SHA-256：\texttt{896c08fb62b3b348}\allowbreak\texttt{432bbc404da427e3}\allowbreak\texttt{50301de9724c2f2f}\allowbreak\texttt{4b5cf1c4682298b1}。}
\item[\texttt{C02}] \sourcepath{docs/pyqcd\_reference/LICENSE.md}
\par\noindent{\scriptsize 控制文件；规模：21 行；字节：1\,071；SHA-256：\texttt{edf5f3be1ced1214}\allowbreak\texttt{81c53c774c321a95}\allowbreak\texttt{5c8c6eca1f188f1e}\allowbreak\texttt{6cdeac0d00216776}。}
\item[\texttt{C03}] \sourcepath{docs/pyqcd\_reference/README.md}
\par\noindent{\scriptsize 控制文件；规模：60 行；字节：3\,243；SHA-256：\texttt{71a3f3e53ed91608}\allowbreak\texttt{de8fa72412cd712f}\allowbreak\texttt{9540b337dd669595}\allowbreak\texttt{7933fa17212a6d7b}。}
\end{description}
\section{覆盖结论与边界}
\begin{checkpointbox}{全量覆盖判据}
\small
\begin{enumerate}
  \item 163 个内容文件均有 SHA-256 和章节归宿；76 对 TeX/PDF 以同一知识单元计，但两个文件均登记。
  \item 历史重复版本不会被当作新物理结论重复堆叠；其差异用于验证演进、失败修复和实现状态变化。
  \item PDF 中的数值结果只按快照语境引用；除非专著另行登记当前源码、运行输入和验证输出，否则不升级为新的物理验收结果。
  \item 固定的上游历史路径、旧接口和集群命令仅属于引用语境，不构成本库当前执行入口。
\end{enumerate}
\end{checkpointbox}


\clearpage
\chapter{教程篇 I：从规范结构到Dirac方程}
\label{sec:tutorial-gauge-dirac}

\section{教程 1：从规范协变性推导Wilson线}

\paragraph{目的.} Wilson线并非额外假设，而是要求两个不同地点的场在同一规范变换下能够比较时自然出现的平行移动。设物质场按
\begin{equation}
  \psi(x)\to\Omega(x)\psi(x),
  \qquad
  \bar\psi(x)\to\bar\psi(x)\Omega^\dagger(x)
  \label{eq:tutorial-matter-gauge}
\end{equation}
变换。局域变换下，普通导数产生额外项：
\begin{equation}
  \partial_\mu[\Omega\psi]
  =\Omega\partial_\mu\psi+(\partial_\mu\Omega)\psi.
  \label{eq:tutorial-derivative-extra}
\end{equation}
引入协变导数
\begin{equation}
  D_\mu\psi=(\partial_\mu-igA_\mu)\psi,
  \qquad
  A_\mu\to\Omega A_\mu\Omega^\dagger+\frac{i}{g}\Omega\partial_\mu\Omega^\dagger,
  \label{eq:tutorial-covariant-gauge}
\end{equation}
后额外项抵消，故 \(D_\mu\psi\to\Omega D_\mu\psi\)。

\paragraph{沿路径的平行移动.} 把有限路径分成 \(N\) 段，每段长度 \(\Delta s\)。小位移下
\begin{equation}
  \psi(x+\Delta s)-U(x+\Delta s,x)\psi(x)
  \simeq\Delta s\left[\partial_s-igA_s(x)\right]\psi(x).
  \label{eq:tutorial-small-transport}
\end{equation}
若要求 \(U(x+\Delta s,x)\to\Omega(x+\Delta s)U(x+\Delta s,x)\Omega^\dagger(x)\)，则
\begin{equation}
  U(x+\Delta s,x)
  =1+igA_s(x)\Delta s+\mathcal O(\Delta s^2).
  \label{eq:tutorial-link-infinitesimal}
\end{equation}
把 \(N\) 段相乘并令 \(\Delta s\to0\)，得到
\begin{equation}
  U_\Gamma(y,x)
  =\lim_{N\to\infty}\prod_{k=0}^{N-1}
  \left[1+igA_s(s_k)\Delta s\right]
  =\mathcal P\exp\left[ig\int_\Gamma A_s\,\dd s\right].
  \label{eq:tutorial-wilson-product}
\end{equation}
路径有序来自矩阵 \(A_s(s_k)\) 在不同位置不对易；若忽略路径排序，就会得到错误的色结构。

\paragraph{格点版本与检查.} 在有限格距上，链接就是有限平行移动：
\begin{equation}
  U_\mu(x)\simeq\exp\left[iagA_\mu(x+a\hat\mu/2)\right],
  \label{eq:tutorial-link-midpoint}
\end{equation}
其中midpoint约定使离散误差为 \(\mathcal O(a^2)\)。检查方法包括：规范变换后链接两端分别乘 \(\Omega\) 与 \(\Omega^\dagger\)、闭合环的迹规范不变、反向链接等于Hermitian共轭。若代码只存储 \(U_\mu\) 而按 \(\exp(igA_\mu)\) 更新，却没有同步 \(U_{-\mu}\)，闭合plaquette会失去幺正性。

\section{教程 2：Plaquette展开与场强张量}

\paragraph{目的.} 说明Wilson规范作用量为什么以plaquette为最小构件。取 \(1\times1\) 闭环
\begin{equation}
  U_{\mu\nu}(x)
  =U_\mu(x)U_\nu(x+a\hat\mu)
  U_\mu^\dagger(x+a\hat\nu)U_\nu^\dagger(x).
  \label{eq:tutorial-plaquette-loop}
\end{equation}
把链接展开到 \(a^2\)：
\begin{align}
  U_\mu(x)
  &=1+iagA_\mu-\frac{a^2g^2}{2}A_\mu^2+\mathcal O(a^3),\\
  U_\nu(x+a\hat\mu)
  &=1+iagA_\nu+\frac{a^2g}{2}\partial_\mu A_\nu
  -\frac{a^2g^2}{2}A_\nu^2+\mathcal O(a^3).
\end{align}
逐项相乘并使用
\begin{equation}
  F_{\mu\nu}=\partial_\mu A_\nu-\partial_\nu A_\mu-ig[A_\mu,A_\nu],
  \label{eq:tutorial-field-strength}
\end{equation}
可得
\begin{equation}
  U_{\mu\nu}
  =1+iag^2F_{\mu\nu}
  -\frac{a^2g^2}{2}F_{\mu\nu}^2+\cdots .
  \label{eq:tutorial-plaquette-expanded}
\end{equation}
取实部并利用 \(\Tr F_{\mu\nu}=0\)：
\begin{equation}
  1-\frac1{N_c}\Re\Tr U_{\mu\nu}
  =\frac{a^4g^2}{2N_c}\Tr F_{\mu\nu}^2
  +\mathcal O(a^6).
  \label{eq:tutorial-plaquette-action-density}
\end{equation}
因此
\begin{equation}
  S_g=\frac{2N_c}{g^2}\sum_{x,\mu<\nu}
  \left(1-\frac1{N_c}\Re\Tr U_{\mu\nu}\right)
  \longrightarrow\frac14\int\dd^4xF_{\mu\nu}^aF^{\mu\nu,a}.
  \label{eq:tutorial-gauge-continuum}
\end{equation}

\paragraph{物理图像.} Plaquette测量一个最小方格上的“曲率”：若规范场是纯规范，链接可以完全由端点规范变换消去，闭环回到单位矩阵；非零曲率使闭环偏离单位矩阵。胶子动力学正是这些最小曲率的平方和。

\paragraph{退化检查.} 当 \(A_\mu=0\) 时 \(U_{\mu\nu}=1\)，作用量为零；当 \(g\to0\) 时 \(\beta\to\infty\)，量子涨落被抑制。若代码把plaquette写成逆序、漏掉某个链接共轭，则在纯规范场测试中会得到非单位闭环，这是比“编译通过”更有效的单元测试。

\section{教程 3：Wilson费米子色散与倍增}

\paragraph{目的.} 解释为什么朴素格点费米子会出现16个低能模。自由WilsonDirac算符写成
\begin{align}
  D_W(x,y)
  &=m_0\delta_{x,y}
  +\frac{1}{2a}\sum_\mu\gamma_\mu
  \left[\delta_{x+\hat\mu,y}-\delta_{x-\hat\mu,y}\right]\nonumber\\
  &\quad-\frac{r}{2a}\sum_\mu
  \left[2\delta_{x,y}
  -\delta_{x+\hat\mu,y}-\delta_{x-\hat\mu,y}\right].
  \label{eq:tutorial-wilson-dirac}
\end{align}
对平面波 \(\psi(x)=e^{ipx}u(p)\)，得到
\begin{equation}
  D_W(p)
  =m_0+\frac{i}{a}\sum_\mu\gamma_\mu\sin(ap_\mu)
  +\frac{2r}{a}\sum_\mu\sin^2\left(\frac{ap_\mu}{2}\right).
  \label{eq:tutorial-wilson-symbol}
\end{equation}
令 \(r=0\) 得到朴素谱
\begin{equation}
  m_{\rm naive}^2(p)
  =m_0^2+\frac1{a^2}\sum_\mu\sin^2(ap_\mu).
  \label{eq:tutorial-naive-spectrum}
\end{equation}
当任意一个 \(ap_\mu\simeq\pi\) 时，\(\sin(ap_\mu)\simeq0\)，因此除 \(p\simeq0\) 外还有多个低能模。四维中超立方体有 \(2^4=16\) 个角，故朴素费米子产生16个味道。

Wilson项在 \(p_\mu\simeq\pi\) 时给出
\begin{equation}
  m_W(p)=m_0+\frac{2r}{a}\sum_\mu\sin^2\frac{ap_\mu}{2}.
  \label{eq:tutorial-wilson-mass}
\end{equation}
对 \(p_\mu\simeq\pi\) 的模，质量为 \(m_0+2rn_\pi/a\)，在 \(a\to0\) 时发散；只有 \(p\simeq0\) 的模保留。代价是手征对称性在有限 \(a\) 下被显式破坏。

\paragraph{极限检查.} 当 \(a\to0\) 且 \(p\) 固定时，\(D_W(p)\to i\gamma_\mu p^\mu+m_0\)，Wilson项为 \(\mathcal O(a)\)。当 \(r=1\) 时，加倍模被推到截断；当 \(r<1\) 时，部分模可能仍较轻，不能忽略。代码中若把 \(r\) 与Clover系数混为一个参数，会改变手征极限。

\section{教程 4：用二乘二模型理解CG}

\paragraph{目的.} 用一个可手算的Hermitian正定系统展示Krylov迭代与条件数。取
\begin{equation}
  A=\begin{pmatrix}1&0\\0&100\end{pmatrix},
  \qquad
  b=\begin{pmatrix}1\\100\end{pmatrix}.
  \label{eq:tutorial-cg-system}
\end{equation}
精确解为 \(x=(1,1)^T\)。从 \(x_0=0\) 开始，\(r_0=b\)，沿 \(p_0=b\) 做一次精确线搜索：
\begin{equation}
  \alpha_0=\frac{b^Tb}{b^TAb}
  =\frac{10001}{1000001}\simeq0.0100.
  \label{eq:tutorial-cg-alpha}
\end{equation}
更新后残差为
\begin{equation}
  r_1=b-\alpha_0Ab
  =\begin{pmatrix}0.9900\\99\end{pmatrix}-\begin{pmatrix}0.0100\\100\end{pmatrix}
  \simeq\begin{pmatrix}0.98\\-1\end{pmatrix}.
  \label{eq:tutorial-cg-residual}
\end{equation}
第二次迭代沿与第一个方向A-共轭的方向消除剩余误差。理想CG至多两步给出精确解。

\paragraph{为什么真实格点不是两步.} 二乘二系统只有两个不同本征值；格点Dirac矩阵有大量本征值。CG的收敛速度由谱分布决定，轻夸克使小本征值密集，导致长期迭代。预条件的作用是把本征值聚集；deflation则把最靠近零的少数方向显式移除。

\paragraph{数值检查.} 在每一步记录递推残差与真残差：
\begin{equation}
  r_k^{\rm rec}=r_{k-1}-\alpha_{k-1}Ap_{k-1},
  \qquad
  r_k^{\rm true}=b-Ax_k.
  \label{eq:tutorial-cg-residuals}
\end{equation}
两者在双精度中可能相差多个数量级。若只打印递推残差，可能在真实误差仍大时误判收敛。

\section{教程 5：构造一个BiCGSTAB breakdown}

\paragraph{目的.} 说明BiCGSTAB的失败不是抽象警告。考虑
\begin{equation}
  D=\begin{pmatrix}0&1\\0&0\end{pmatrix},
  \qquad
  b=\begin{pmatrix}1\\0\end{pmatrix}.
  \label{eq:tutorial-breakdown-system}
\end{equation}
该矩阵不可逆，精确求解应在算法前被拒绝。若错误地继续迭代，取 \(x_0=0\)、\(r_0=b\)、影子残差 \(\hat r_0=(0,1)^T\)，则
\begin{equation}
  \rho_0=\hat r_0^\dagger r_0=0,
  \label{eq:tutorial-breakdown-rho}
\end{equation}
第一步就发生 \(\rho\)-breakdown。它提示的不是“增加迭代次数”，而是检查矩阵是否可逆、初值是否退化以及影子残差是否与右端正交。

\paragraph{一个可逆但接近breakdown的例子.} 取
\begin{equation}
  D_\epsilon=\begin{pmatrix}\epsilon&1\\0&\epsilon\end{pmatrix},
  \qquad 0<\epsilon\ll1.
  \label{eq:tutorial-near-breakdown}
\end{equation}
此时条件数约 \(\epsilon^{-1}\)，BiCGSTAB的 \(\omega\) 与 \(\rho\) 会因有限精度剧烈变化。可靠实现应：
\begin{enumerate}
  \item 在 \(\rho,\omega\) 接近机器精度时触发重启或切换GMRES；
  \item 对右端和影子残差做尺度归一化；
  \item 定期计算真残差并检查NaN/Inf；
  \item 保存失败迭代的状态，而不是只输出最后一次残差。
\end{enumerate}

\paragraph{物理对应.} QCD中近零模、非零化学势和粗格距可使 \(D\) 的有效谱接近这种结构。Breakdown不是“BiCGSTAB不适合QCD”的普遍结论，而是预条件、deflation和精度策略不足的信号。选择算法前应先理解谱，而不是先比较墙钟时间。

\section{教程 6：相关比值与bootstrap误差}

\paragraph{目的.} 展示为什么比值、拟合和矩阵求逆必须在同一replica内完成。设两个相关观测量 \(A,B\)，目标为
\begin{equation}
  R=\frac{A}{B}.
  \label{eq:tutorial-ratio}
\end{equation}
若用线性误差传播，
\begin{equation}
  \sigma_R^2
  =\left(\frac{\partial R}{\partial A}\right)^2\sigma_A^2
  +\left(\frac{\partial R}{\partial B}\right)^2\sigma_B^2
  +2\frac{\partial R}{\partial A}\frac{\partial R}{\partial B}C_{AB}.
  \label{eq:tutorial-ratio-delta}
\end{equation}
其中
\begin{equation}
  \frac{\partial R}{\partial A}=\frac1B,
  \qquad
  \frac{\partial R}{\partial B}=-\frac{A}{B^2}.
  \label{eq:tutorial-ratio-derivatives}
\end{equation}
如果 \(A\) 与 \(B\) 来自同一组态，\(C_{AB}\) 往往不可忽略；把两个误差直接平方相加会高估或低估。

\paragraph{Bootstrap做法.} 对每个组态集合 \(S_r\) 同时计算 \(A_r,B_r\)，再形成
\begin{equation}
  R_r=\frac{A_r}{B_r}.
  \label{eq:tutorial-ratio-bootstrap}
\end{equation}
最终误差为 \(R_r\) 的标准差。这样自动包含 \(A,B\) 的相关性，也能检测 \(B_r\) 接近零的非线性尾部。若先在中心样本上算 \(R_0\)，再单独bootstrap \(A\) 与 \(B\)，就丢失了相关性。

\paragraph{检查.} 当 \(A\) 与 \(B\) 完全相关且 \(A=cB\) 时，\(R=c\)，误差应为零。使用式~\eqref{eq:tutorial-ratio-delta} 只有在 \(C_{AB}=\sigma_A\sigma_B\) 时才给出零；使用全链bootstrap则会自动得到接近零的散布。这个极限是检验分析脚本的强测试。

\section{教程 7：指数尾部的Fourier变换}

\paragraph{目的.} 说明大 \(z\) 尾部如何影响小 \(x\) 与端点。取一个归一化玩具尾部
\begin{equation}
  h(z)=e^{-\lambda z},
  \qquad z\ge0,
  \label{eq:tutorial-exp-tail}
\end{equation}
其Fourier变换为
\begin{equation}
  \widetilde h(P_z,x)
  =\int_0^\infty\frac{\dd z}{2\pi}
  e^{ixP_zz}e^{-\lambda z}
  =\frac{1}{2\pi}\frac{1}{\lambda-ixP_z}.
  \label{eq:tutorial-exp-fourier}
\end{equation}
当 \(xP_z\ll\lambda\) 时，\(\widetilde h\simeq1/(2\pi\lambda)\)；当 \(xP_z\gg\lambda\) 时，幅度按 \(1/(xP_z)\) 衰减。若格点只覆盖 \(z\le z_{\max}\)，截断积分为
\begin{equation}
  \widetilde h_{z_{\max}}
  =\frac{1-e^{-(\lambda-ixP_z)z_{\max}}}
  {2\pi(\lambda-ixP_z)}.
  \label{eq:tutorial-exp-truncated}
\end{equation}
误差
\begin{equation}
  \delta\widetilde h
  =\widetilde h_{z_{\max}}-\widetilde h
  =-\frac{e^{-(\lambda-ixP_z)z_{\max}}}
  {2\pi(\lambda-ixP_z)}
  \label{eq:tutorial-exp-error}
\end{equation}
按 \(e^{-\lambda z_{\max}}\) 衰减，但分母含 \(xP_z\)，因此小 \(x\) 的绝对误差相对较大。

\paragraph{物理结论.} 大 \(z\) 尾部不能只凭“看起来趋于零”忽略。应检查 \(z_{\max}P_z\)、尾部拟合窗口和目标 \(x\) 区间，并把尾部模型变化纳入误差。指数尾部是理想化模型；真实QCD还可能有幂次、振荡和高扭度贡献。

\section{教程 8：标量玩具TMD因子化}

\paragraph{目的.} 用一个标量模型展示硬核、软核和TMD分布的乘法结构。设一个标量“部分子”与外部源耦合，领头区域贡献为
\begin{equation}
  \mathcal M(P_z,b_T)
  =\int\frac{\dd k_T^2}{(2\pi)^2}
  H(P_z,k_T)\,S(b_T,k_T)\,f(k_T).
  \label{eq:tutorial-scalar-factorization}
\end{equation}
若在leading power中 \(k_T\) 依赖可分离，
\begin{equation}
  S(b_T,k_T)\simeq S(b_T)\,\phi(k_T),
  \qquad
  f(k_T)\phi(k_T)\ \text{进入TMD波函数},
  \label{eq:tutorial-scalar-separation}
\end{equation}
则
\begin{equation}
  \mathcal M(P_z,b_T)
  =H(P_z)\,S(b_T)\,f(b_T)
  +\mathcal O\left(\frac{\Lambda^2}{P_z^2}\right).
  \label{eq:tutorial-scalar-result}
\end{equation}

\paragraph{为什么真实QCD更难.} 胶子有颜色和自旋，软函数带Wilson线表示，硬核还可能有 \(x\) 卷积和味混合。标量模型只保留“硬、软、共线”三因子结构，不能据此断言胶子TMD没有rapidity发散或单态混合。它的价值是帮助识别因子化逻辑，而不是替代QCD证明。

\paragraph{退化检查.} 当 \(b_T\to0\) 时，\(S(b_T)\) 应回到短距极限；当 \(P_z\to\infty\) 时，高扭度项消失；当软耦合关闭时，\(S=1\)，结果回到普通共线分布。任何数值实现都应至少通过这些玩具极限。

\section{教程 9：Collins--Soper一致性玩具}

\paragraph{目的.} 展示“物理量与方案无关”如何给出CS方程。设裸TMD与硬核分别为
\begin{align}
  f_{\rm bare}(x,\zeta)
  &=Z_f^{-1}(\zeta/\zeta_0)f_R(x,\zeta_0),\\
  H_{\rm bare}(Q,\zeta)
  &=Z_H(\zeta/\zeta_0)H_R(Q,\zeta_0).
\end{align}
物理乘积
\begin{equation}
  \sigma=H_{\rm bare}f_{\rm bare}
  =H_Rf_R
  \label{eq:tutorial-cs-product}
\end{equation}
必须与 \(\zeta\) 无关。因此
\begin{equation}
  Z_HZ_f^{-1}=1
  \quad\Longrightarrow\quad
  \zeta\frac{\partial}{\partial\zeta}\ln Z_H
  =\zeta\frac{\partial}{\partial\zeta}\ln Z_f.
  \label{eq:tutorial-cs-z}
\end{equation}
若把 \(Z_f\) 的指数写成
\begin{equation}
  Z_f(\zeta/\zeta_0)
  =\exp\left[\frac12K\ln\frac{\zeta}{\zeta_0}\right],
  \label{eq:tutorial-cs-zf}
\end{equation}
则立即得到
\begin{equation}
  2\zeta\frac{\partial}{\partial\zeta}\ln f_R
  =K.
  \label{eq:tutorial-cs-result}
\end{equation}
这个玩具推导说明CS方程来自因子化两边演化的相消，而不是额外的动力学假设。

\paragraph{方案边界.} 真实TMD的 \(Z_f\) 还可能含软函数、rapidity regulator和Wilson线几何。若硬核使用另一种方案，式~\eqref{eq:tutorial-cs-z} 中的 \(Z_H\) 必须同步改变。把不同方案的中心值混合，会人为破坏 \(\zeta\) 无关性。

\section{教程 10：梯度流热核与量纲匹配}

\paragraph{目的.} 说明流时 \(t_f\) 为什么能作为短距标度。在线性化近似中，规范场按热方程演化：
\begin{equation}
  \partial_{t_f}B_\mu=\partial^2B_\mu,
  \qquad
  B_\mu(t_f)=e^{t_f\partial^2}B_\mu(0).
  \label{eq:tutorial-flow-heat}
\end{equation}
动量空间解为
\begin{equation}
  B_\mu(k,t_f)=e^{-t_fk^2}B_\mu(k,0).
  \label{eq:tutorial-flow-momentum}
\end{equation}
因此高动量 \(k^2\gg1/t_f\) 被指数压低。若流场与\(\MSbar\)算符通过
\begin{equation}
  O(t_f)=C(t_f,\mu)O_R(\mu)
  \label{eq:tutorial-flow-matching}
\end{equation}
匹配，则 \(C\) 的微扰展开包含 \(\alpha_s\) 与 \(\ln(t_f\mu^2)\)。有效流时窗口要求
\begin{equation}
  a\ll\sqrt{8t_f}\ll\frac1{\Lambda_{\rm QCD}},
  \label{eq:tutorial-flow-window}
\end{equation}
左边控制离散误差，右边控制非微扰幂修正。

\paragraph{物理图像.} 梯度流像对规范场做扩散：短波涨落先被平滑，长波结构保留。它不改变规范协变性，但会改变算符的标度和有限 \(a\) 结构。流时匹配不是把“时间”当作任意参数，而是在短距展开内把 \(t_f\) 与 \(\mu\) 关联。

\paragraph{检查.} 当 \(t_f\to0\) 时 \(C\to1\)；当 \(a\to0\) 且 \(t_f\) 固定物理值时，流场应趋于连续流场；当 \(t_f\) 接近 \(L^2\) 时，有限体积和边界效应会污染匹配。若流窗口为空，应报告“没有可发布区间”，而不是在窗口外强行拟合。

\section{教程 11：Perambulator与收缩复杂度}

\paragraph{目的.} 说明distillation为什么能降低三点收缩成本。把传播子投影到低维空间 \(V\)，广义perambulator可写成
\begin{equation}
  J_{\alpha\beta}(t_f,t,t_i)
  =\left[V^\dagger(t_f)M^{-1}(t_f,t)
  \Gamma(t)M^{-1}(t,t_i)V(t_i)\right]_{\alpha\beta}.
  \label{eq:tutorial-perambulator}
\end{equation}
若不使用蒸馏，直接收缩的复杂度近似为
\begin{equation}
  C_{\rm direct}\sim N_{\rm src}N_{\rm sink}N_x^3,
  \label{eq:tutorial-direct-cost}
\end{equation}
其中 \(N_x\) 是空间格点数。使用维数 \(N_v\) 的蒸馏基后，复杂度变为
\begin{equation}
  C_{\rm dist}\sim N_{\rm src}N_{\rm sink}N_v^p,
  \qquad p=2,3,4,
  \label{eq:tutorial-dist-cost}
\end{equation}
具体指数取决于两点、介子三点、重子三点或非连通环。

\paragraph{成本不是单调下降.} 若 \(N_v\) 接近 \(N_x\)，蒸馏的预计算和存储可能比直接方法更贵。实际收益取决于：
\begin{enumerate}
  \item 低模基对目标强子态的覆盖；
  \item 需要复用的源和汇数量；
  \item 算符插入的Lorentz/颜色结构；
  \item 传播子缓存和GPU内存；
  \item 是否需要非连通图。
\end{enumerate}
报告应给出 \(N_v\)、基算符数、源/汇数、matvec数、缓存命中和最终误差，而不是只给“加速了多少倍”。

\paragraph{检查.} 当 \(N_v\to N_x\) 时，蒸馏结果应趋近完整收缩；当基算符只保留一个态时，结果应接近单态近似但不应被当作完整谱学。任何“加速”都必须与完整基或独立方法在小子体积上对照。

\clearpage
\chapter{教程篇 II：专业术语百科}
\label{sec:glossary}

本百科按使用场景而非字母顺序编排。每条术语给出“是什么、为什么重要、最容易错在哪里”三层信息；公式只保留最能区分概念的部分。阅读正文遇到不熟悉的名词时，可先在此定位其对象空间与误差来源。

\section{场论、QCD与因子化术语}

{\small
\begin{longtable}{L{2.8cm}L{1.9cm}L{8.7cm}}
\caption{场论、QCD与因子化核心术语。}\label{tab:glossary-theory}\\
\toprule
术语 & 类别 & 物理含义与常见错误 \\
\midrule
\endfirsthead
\toprule
术语 & 类别 & 物理含义与常见错误 \\
\midrule
\endhead
\bottomrule
\endfoot
规范链接 & 格点规范场 & 连接相邻格点的平行移动。它把局域规范变换分配到路径两端；漏掉反向链接或共轭会使闭环不幺正。 \\
Wilson线 & 非局域算符 & 路径有序指数 \(W=\mathcal P e^{ig\int A}\)。路径方向决定过程依赖，不能只当作标量相位。 \\
Staple & Wilson几何 & 两条长平行Wilson线加横向连接，用于TMD横向分离。与普通直线Wilson线的重正化不同。 \\
Cusp & Wilson端点 & 路径转角造成的额外UV发散。其系数依赖转角、表示和端点；不可用普适常数替代。 \\
Pinch pole & 长Wilson线 & 两条几乎平行Wilson线之间的有效势导致随长度增长的贡献。需要同几何Wilson圈或长度外推处理。 \\
Rapidity发散 & 软/共线 & 沿光锥方向积分产生的发散。普通UV重正化不能消除，需要rapidity regulator和软函数。 \\
Soft function & 软区域 & 低能胶子交换形成的规范不变核。依赖Wilson线几何、表示和rapidity方案。 \\
Collins--Soper核 & 演化核 & 控制TMD对 \(\zeta\) 的演化，常记 \(K(\bT;\mu)\)。它应强子无关，但可与表示和方案有关。 \\
TMD & 分布函数 & transverse-momentum-dependent分布，同时保留 \(x,\bT,\mu,\zeta\)。与共线PDF的差别不是多一个拟合参数。 \\
Quasi-TMD & Euclidean对象 & 用大 \(P_z\) 的类空分离逼近光锥TMD。必须经过UV重正化、软减除、匹配和外推。 \\
LaMET & 大动量方法 & large-momentum effective theory，用 \(1/P_z\) 展开连接准量与光锥量。极限顺序不可交换。 \\
Factorization & 因子化 & 将硬、共线、软区域分配到不同因子并由演化核吸收重叠。不是任意代数分解。 \\
Hard kernel & 微扰核 & 短距Wilson系数，通常含 \(\alpha_s\)、对数和大 \(P_z\)。其阶数必须与软/重正化方案配套。 \\
Matching & 匹配 & 把准算符或格点方案映射到光锥方案。可能含卷积、味混合和有限常数。 \\
Operator mixing & 算符混合 & 同量纲算符在UV下互相重正化，\(O_{i,R}=Z_{ij}O_j\)。与硬匹配混合不是同一层。 \\
Anomalous dimension & 反常量纲 & 算符随 \(\mu\) 演化的指数。标量、向量和矩阵反常量纲不能混用。 \\
Twist & 幂次分类 & 量纲减自旋，控制高扭度修正。leading twist 不是“所有修正都为零”。 \\
OPE & 算符乘积展开 & 短距下把非局域乘积展开为局域算符与Wilson系数。剩余项由更高 twist 或更高维算符组成。 \\
Wilson coefficient & 短距系数 & OPE或匹配中的微扰系数，依赖标度比和耦合。不能与格点重正化因子等同。 \\
DGLAP & 共线演化 & 控制PDF随 \(\mu\) 的劈裂卷积。TMD的CS演化与DGLAP是不同方向。 \\
CSS & TMD演化框架 & Collins--Soper--Sterman组织rapidity与软函数。不同regulator下中间量可不同，物理量应一致。 \\
SIDIS & 过程 & semi-inclusive deep inelastic scattering。Wilson线方向与Drell--Yan相反。 \\
Drell--Yan & 过程 & 轻子对产生过程。其TMD软结构与SIDIS过程方向不同，T-odd项尤其敏感。 \\
\end{longtable}
}

\section{格点、重正化与传播子术语}

{\small
\begin{longtable}{L{2.8cm}L{1.9cm}L{8.7cm}}
\caption{格点与重正化核心术语。}\label{tab:glossary-lattice}\\
\toprule
术语 & 类别 & 物理含义与常见错误 \\
\midrule
\endfirsthead
\toprule
术语 & 类别 & 物理含义与常见错误 \\
\midrule
\endhead
\bottomrule
\endfoot
Plaquette & 规范作用量 & 最小规范闭环，测量场强曲率。逆序、共轭或归一化错误会破坏纯规范测试。 \\
Clover & 费米作用量 & Sheikholeslami--Wohlert项，用 \(\sigma_{\mu\nu}F^{\mu\nu}\) 改善手征行为。系数和场强符号必须固定。 \\
HYP & 链接平滑 & 通过嵌套链接平均降低UV噪声，同时保持规范协变。它改变短距算符，必须与匹配一起处理。 \\
Gradient flow & UV正规化 & 沿流时演化规范场，抑制高动量模。短流时匹配与长流时非微扰修正必须同时控制。 \\
Flow time & 流参数 & 常用 \(t_f\) 或 \(8t_f\)。流时不是普通重正化标度，窗口需满足 \(a\ll\sqrt{8t_f}\ll1/\Lambda\)。 \\
RI/MOM & 重正化方案 & 在固定动量插入条件下定义投影反常量纲。不能自动消除非局域Wilson线线性发散。 \\
Ratio scheme & 重正化方案 & 用零动量短距参考点相除。参考点和短距窗口必须报告。 \\
Hybrid & 混合方案 & 短距用ratio、长距用质量减除。关键是短长距重叠区和 \(z_s\) 稳定性。 \\
Self-renormalization & 重正化策略 & 用矩阵元自身构造减除因子。需证明残留UV和rapidity结构可被同一因子吸收。 \\
Wilson loop & 规范不变量 & 闭合矩形环的真空期望。同几何平方根可减除部分自能与pinch pole。 \\
Linear divergence & UV发散 & 常表现为 \(e^{-\delta m L}\)。系数依赖表示和路径长度，不可用普通 \(\MSbar\) 常数替代。 \\
Endpoint divergence & 端点UV & 场强与Wilson线连接处的局域发散。通常进入算符反常量纲或短距系数。 \\
Contact term & 三点函数 & 插入与源/汇重合或近邻时的非传播贡献。必须与物理矩阵元分离。 \\
Source & 源算符 & 产生初始态的插值场。源类型影响基态重叠和统计，不改变物理定义。 \\
Sink & 汇算符 & 湮灭末态的插值场。源汇必须使用同一态投影和动量约定。 \\
Sequential source & 顺序源 & 把汇点传播子整体求解，减少三点函数求逆次数。每个末态动量通常需要新源。 \\
Distillation & 低模谱学 & 用低维基压缩传播子。收益取决于基覆盖、复用数和算符结构。 \\
Perambulator & 蒸馏对象 & 源侧、流插入和汇侧传播子组成的低模矩阵。成本随 \(N_v\) 的多项式增长。 \\
GEVP & 谱学拟合 & 广义本征值问题分离能级。需报告基底、\(t_0\) 和条件数。 \\
Effective mass & 谱学诊断 & \(m_{\rm eff}(t)=\ln[C(t)/C(t+1)]\)。平台不唯一，必须与谱模型和源汇稳定性联合判断。 \\
\end{longtable}
}

\section{数值算法与统计术语}

{\small
\begin{longtable}{L{2.8cm}L{1.9cm}L{8.7cm}}
\caption{数值算法、逆问题与统计术语。}\label{tab:glossary-numerics}\\
\toprule
术语 & 类别 & 物理含义与常见错误 \\
\midrule
\endfirsthead
\toprule
术语 & 类别 & 物理含义与常见错误 \\
\midrule
\endhead
\bottomrule
\endfoot
Condition number & 线性代数 & \(\kappa=\|D\|\|D^{-1}\|\)，衡量误差放大。大条件数使迭代和反演不稳定。 \\
Residual & 求解器 & \(r=b-Ax\)。递推残差可能与真残差差很多，必须周期重算。 \\
CG & Krylov & 用于Hermitian正定系统的短递推方法。依赖正定性和良好预条件。 \\
GMRES & Krylov & 对非Hermitian系统最小化残差，存储随迭代增长。重启可能停滞。 \\
BiCGSTAB & Krylov & 双共轭梯度加局部最小残差，短存储但可能breakdown或振荡。 \\
Preconditioner & 加速器 & 近似 \(M^{-1}\)，把谱聚集。必须报告matvec、内存和真残差。 \\
Deflation & 谱方法 & 显式移除低模。低模基底必须来自同一算符和边界条件。 \\
Multigrid & 多层求解 & 在粗层消除低频误差。粗化必须尊重规范结构和拓扑。 \\
Mixed precision & 精度策略 & 低精度内层、高精度外层。条件数大时可能污染Krylov基。 \\
Random source & 随机估计 & 用随机向量估计迹或全模式。方差取决于源和传播子结构。 \\
Bootstrap & 重采样 & 从底层样本有放回抽取并重跑完整分析。必须保留组态相关性。 \\
Jackknife & 重采样 & 删除block估计方差。Block长度应覆盖主要自相关。 \\
Autocorrelation & 时间序列 & 组态间相关性。有效样本数 \(N_{\rm eff}=N/(2\tau_{\rm int})\)。 \\
GLS & 拟合 & 广义最小二乘使用协方差逆。协方差奇异或样本不足时需正则化。 \\
AIC/BIC & 模型选择 & 在拟合优度与参数数之间折中。不能替代物理可辨识性。 \\
Model averaging & 系统误差 & 按模型权重或扫描范围组合结果。模型外推差异必须保留。 \\
Inverse problem & 逆问题 & 从有限带噪数据恢复分布。小奇异值会放大噪声，需要正则化或先验。 \\
Backus--Gilbert & 反演 & 在分辨率与方差之间权衡。输出带偏，需报告分辨率核。 \\
Regularization & 反演 & Tikhonov、先验或尾部模型。正则化参数不能在同一数据上任意调优。 \\
Blinding & 防偏 & 在正式解盲前隐藏分析选择或偏移。随机种子和方案必须预注册。 \\
\end{longtable}
}

\clearpage
\chapter{教程篇 III：公式与算法手册}
\label{sec:formula-algorithm-handbook}

\section{色代数与规范结构}

\begin{keyformula}{QCD基本代数}
\begin{align}
  [T^a,T^b]&=if^{abc}T^c,
  & \Tr(T^aT^b)&=\frac12\delta^{ab},
  \label{eq:handbook-color}\\
  C_F&=\frac{N_c^2-1}{2N_c},
  & C_A&=N_c,
  \label{eq:handbook-casimir}\\
  U_\mu(x)&\to\Omega(x)U_\mu(x)\Omega^\dagger(x+a\hat\mu),
  \label{eq:handbook-link}\\
  F_{\mu\nu}&=\partial_\mu A_\nu-\partial_\nu A_\mu-ig[A_\mu,A_\nu].
  \label{eq:handbook-fieldstrength}
\end{align}
\end{keyformula}

色代数计算的三条原则是：表示必须匹配、指标必须缩并、闭合算符必须规范不变。若中间量依赖规范参数，最终物理量只有在与相应Wilson线、软函数或投影组合后才可能规范无关。

\section{光锥、TMD与演化}

\begin{keyformula}{TMD定义与演化}
\begin{align}
  f_{g/N}(x,\bT;\mu,\zeta)
  &=\frac12\delta_T^{ij}\Phi_{g,ij}(x,\bT;\mu,\zeta),
  \label{eq:handbook-tmd}\\
  2\zeta\frac{\partial}{\partial\zeta}\ln f
  &=K(\bT;\mu),
  \label{eq:handbook-cs}\\
  \mu\frac{\dd}{\dd\mu}K
  &=-\gamma_K(\alpha_s),
  \label{eq:handbook-cs-rge}\\
  \widetilde f_g
  &=H_{gg}\otimes f_g+H_{gq}\otimes q_S
  +\mathcal O(\Lambda^2/P_z^2).
  \label{eq:handbook-matching}
\end{align}
\end{keyformula}

必须区分三个标度：\(\mu\) 控制UV，\(\zeta\) 控制rapidity，\(P_z\) 控制Euclidean大动量。三个标度的演化方程不能相互替代。

\section{格点关联函数与流}

\begin{keyformula}{格点与梯度流}
\begin{align}
  C_2(t,\mathbf p)
  &=\sum_{\mathbf x}e^{-i\mathbf p\cdot\mathbf x}
  \langle O_\Gamma(\mathbf x,t)O_s^\dagger(0)\rangle,
  \label{eq:handbook-c2}\\
  C_3(t,\tau,\mathbf p_f,\mathbf p_i)
  &=\sum_{\mathbf x_f,\mathbf x_J}
  e^{-i\mathbf p_f\cdot\mathbf x_f}
  e^{i(\mathbf p_f-\mathbf p_i)\cdot\mathbf x_J}
  \langle O_fJ O_i^\dagger\rangle,
  \label{eq:handbook-c3}\\
  \dot V_\mu(t_f)
  &=-g_0^2\frac{\delta S_g}{\delta V_\mu(t_f)}V_\mu(t_f),
  \label{eq:handbook-flow}\\
  a&\ll\sqrt{8t_f}\ll\Lambda_{\rm QCD}^{-1}.
  \label{eq:handbook-flow-window}
\end{align}
\end{keyformula}

三点函数的源、汇、插入顺序和接触项必须先固定，再讨论拟合。流时窗口是物理幂修正与离散误差的交集，不是可由单个最优值替代的调参范围。

\section{重正化、匹配与误差}

\begin{keyformula}{重正化与误差传播}
\begin{align}
  O_{i,R}(\mu)&=Z_{ij}(\mu,\mu_0)O_{j,R}(\mu_0),
  \label{eq:handbook-renorm}\\
  h_R(z,\bT;P_z)
  &=\lim_{L_\parallel\to\infty}
  \frac{h(z,\bT,L_\parallel;P_z)}
  {\sqrt{Z_E(\bT;2L_\parallel+z)}},
  \label{eq:handbook-wilson-ratio}\\
  \chi^2(\theta)
  &=(y-y_{\rm model})^TC^{-1}(y-y_{\rm model}),
  \label{eq:handbook-chi2}\\
  \sigma_{\rm sys}^2
  &=\sum_s\sigma_s^2+2\sum_{s<s'}C_{ss'}.
  \label{eq:handbook-sys}
\end{align}
\end{keyformula}

重正化因子、Wilson系数、软函数和统计协方差属于四个不同对象。把它们的误差或标度依赖合并成一个标量，会使物理检验失去意义。

\section{算法索引}

\begin{longtable}{L{2.2cm}L{2.8cm}L{2.7cm}L{5.2cm}}
\caption{本文算法索引、复杂度与失败信号。}\label{tab:algorithm-index}\\
\toprule
算法 & 输入/输出 & 粗复杂度 & 主要失败信号 \\
\midrule
\endfirsthead
\toprule
算法 & 输入/输出 & 粗复杂度 & 主要失败信号 \\
\midrule
\endhead
\bottomrule
\endfoot
HMC & 规范场到组态 & 每轨迹多次Dirac解 & 接受率异常、拓扑冻结、力发散 \\
CG & Hermitian正定 \(A,b\) & \(O(N_{\rm iter}N_D)\) & 真残差平台、正定性破坏 \\
GMRES & 非Hermitian \(D,b\) & \(O(mN_D)\)存储 & 重启停滞、基退化、内存超限 \\
BiCGSTAB & 非Hermitian \(D,b\) & 每步2次matvec & \(\rho,\omega\) breakdown、残差振荡 \\
Deflation & \(D,b,U\) & 粗空间加投影Krylov & 粗矩阵奇异、低模基底不匹配 \\
Multigrid & 多层系统 & 层级matvec & 粗化非物理、粗层零模 \\
Bootstrap & 组态block & \(O(RN_{\rm analysis})\) & replica失败依赖观测值 \\
GEVP & 相关矩阵 & \(O(N_b^3)\)每时间片 & 基底不足、条件数过大 \\
Tail fit & 大 \(z\) 数据 & 非线性拟合 & 尾部模型不可辨识、窗口依赖 \\
Derivative FT & \(h(z)\) 与导数 & 每次 \(O(N_zN_x)\) & 小 \(x\) 放大噪声、端点缺失 \\
Backus--Gilbert & 核与协方差 & 每目标点解线性系统 & 分辨率过粗、先验敏感 \\
Tikhonov & 离散逆问题 & SVD或Krylov & 正则化参数过拟合 \\
Sequential source & 源与汇动量 & 每次新汇点求逆 & 相位或汇投影错误 \\
Distillation & 低模基与传播子 & \(O(N_v^p)\) & 基截断、缓存错误、态遗漏 \\
Random trace & 随机源与 \(D^{-1}\) & \(O(N_r)\)次求解 & 方差大、源未稀释 \\
\end{longtable}

\begin{checkpointbox}{公式手册使用规则}
\begin{enumerate}
  \item 先确认公式属于哪个对象空间，再代入数值；
  \item 不跨方案混用 \(Z\)、\(H\)、\(S\) 和 \(K\)；
  \item 每次变换后做量纲、极限和复共轭检查；
  \item 算法结果必须附真残差、条件数和失败率；
  \item 同一公式在不同论文中记号不同，引用时同时记录定义和方案。
\end{enumerate}
\end{checkpointbox}

\chapter{教程篇 IV：复现实验与练习}
\label{sec:reproduction-labs}

本节给出六个最小实验。每个实验都设计成可以在小子体积、解析模型或合成数据上完成，目的是在进入生产系综前验证公式、索引、状态和误差传播。

\section{实验 1：纯规范Plaquette}

\paragraph{目标.} 验证链接方向、共轭、矩阵乘法和规范变换。

\paragraph{输入.} 随机生成规范变换 \(\Omega(x)\in\mathrm{SU}(3)\)，取纯规范链接
\begin{equation}
  U_\mu(x)=\Omega(x)\Omega^\dagger(x+a\hat\mu).
  \label{eq:lab-pure-gauge}
\end{equation}

\paragraph{步骤.} 对每个 \(x,\mu<\nu\) 构造plaquette，计算 \(\Re\Tr U_{\mu\nu}/N_c\)，再做一次随机规范变换并重复。

\paragraph{预期输出.} 所有plaquette在机器精度内等于1，规范变换前后不变。若出现偏离，优先检查 \(U_{-\mu}\)、链接共轭和矩阵乘法顺序。

\paragraph{失败信号.} 只有部分平面失败通常是指标顺序错误；所有平面随机偏离通常是链接未幺正化；规范变换后改变说明两端变换未配平。

\section{实验 2：二乘二CG与真残差}

\paragraph{目标.} 验证CG递推、停止条件和误差界。

\paragraph{输入.} 使用式~\eqref{eq:tutorial-cg-system} 的矩阵和右端，以及一个条件数更大的对角矩阵。

\paragraph{步骤.} 实现CG，分别记录递推残差、真残差、A-范数误差和matvec次数。用不同容差重复。

\paragraph{预期输出.} 二乘二系统至多两步收敛；大条件数系统迭代次数增加。真残差最终应低于容差，而不是只让递推残差下降。

\paragraph{失败信号.} 真残差停滞、A-范数误差不下降或出现负的 \(p^\dagger Ap\)，说明正定性、内积或更新顺序有问题。

\section{实验 3：相关比值Bootstrap}

\paragraph{目标.} 验证相关误差传播。

\paragraph{输入.} 生成相关高斯样本
\begin{equation}
  (A_r,B_r)\sim\mathcal N(0,C),
  \qquad
  C=\begin{pmatrix}\sigma_A^2&\rho\sigma_A\sigma_B\\
  \rho\sigma_A\sigma_B&\sigma_B^2\end{pmatrix}.
  \label{eq:lab-correlated-sample}
\end{equation}

\paragraph{步骤.} 对每个replica计算 \(R_r=A_r/B_r\)，并比较bootstrap标准差、delta方法和完全相关极限。

\paragraph{预期输出.} 当 \(\rho\to1\) 且 \(A=cB\) 时，比值误差趋于零；当 \(B\) 接近零时，分布出现非线性长尾。

\paragraph{失败信号.} Bootstrap误差与 \(C\) 的相关系数无关，或独立重采样 \(A,B\) 后得到错误零误差。

\section{实验 4：指数尾部Fourier}

\paragraph{目标.} 验证Fourier相位、截断误差和尾部模型。

\paragraph{输入.} 解析函数 \(h(z)=e^{-\lambda z}\)，在 \(0\le z\le z_{\max}\) 上采样。

\paragraph{步骤.} 实现式~\eqref{eq:tutorial-exp-truncated} 的数值积分，与解析结果比较；改变 \(P_z,z_{\max},\lambda\)。

\paragraph{预期输出.} 截断误差按 \(e^{-\lambda z_{\max}}\) 衰减，但在小 \(x\) 处相对误差增大；正负相位约定给出复共轭。

\paragraph{失败信号.} 常数尾部变换出现振荡、误差不随 \(z_{\max}\) 减小或实部不收敛，说明相位、端点权重或归一化错误。

\section{实验 5：匹配方案一致性}

\paragraph{目标.} 验证 \(\mu\) 与 \(\zeta\) 的演化相消。

\paragraph{输入.} 一个玩具硬核 \(H(\mu)\)、CS核 \(K(\mu)\) 和TMD分布 \(f(\mu,\zeta)\)，满足式~\eqref{eq:tutorial-cs-product}。

\paragraph{步骤.} 在 \(\mu_1,\mu_2\) 和 \(\zeta_1,\zeta_2\) 上计算 \(Hf\)，检查物理组合是否在数值误差内不变。故意混用不同方案，观察残余依赖。

\paragraph{预期输出.} 同方案组合与标度无关；跨方案组合出现人为 \(\mu,\zeta\) 依赖。

\paragraph{失败信号.} 同方案也不变，可能是硬核和分布都被人为固定；跨方案不变，可能是方案标签未真正进入公式。

\section{实验 6：完整误差预算}

\paragraph{目标.} 把小系统误差组织成可审计表格。

\paragraph{输入.} 为某个玩具观测量设定统计、格距、体积、流时、匹配阶数和尾部模型六类变化。

\paragraph{步骤.} 分别扫描每类变化，记录中心值、协方差和变化范围；区分相关与独立误差，生成总误差预算。

\paragraph{预期输出.} 每个误差来源都有可复现命令、参数范围和物理含义；总误差不是未经验证的平方和。

\paragraph{失败信号.} 误差表只有数值没有来源，或不同分析阶段重复计算同一组态误差。

\begin{table}[htbp]
\centering
\caption{六个复现实验的最低验收标准。}
\label{tab:lab-acceptance}
\small
\begin{tabularx}{\textwidth}{L{2.2cm}L{4.4cm}Y}
\toprule
实验 & 成功证据 & 不得宣称的内容 \\
\midrule
Plaquette & 纯规范闭环等于1且规范不变 & 真实胶子动力学已正确 \\
CG & 真残差低于容差且次数可复现 & 大生产系综性能已达标 \\
Bootstrap & 相关极限与解析误差一致 & 系统误差已经完整 \\
Fourier & 截断误差与解析趋势一致 & 物理TMD尾部已确定 \\
匹配 & 同方案标度相消、跨方案不消 & 完整单态匹配已实现 \\
误差预算 & 每项有命令、窗口和协方差 & 物理结果已通过全部验收 \\
\bottomrule
\end{tabularx}
\end{table}

\begin{checkpointbox}{进入生产计算前的门槛}
\begin{enumerate}
  \item 六个实验至少全部通过一次并保存日志；
  \item 每个实验都有失败注入测试；
  \item 公式约定、数据布局和缓存键已固定；
  \item 小系统结果与解析或稠密解一致；
  \item 生产流水线复用同一经过验证的组件，而不是另写一套实现。
\end{enumerate}
\end{checkpointbox}

\backmatter
\begin{thebibliography}{99}

\bibitem{Collins2011}
J. Collins,
\newblock \textit{Foundations of Perturbative QCD},
\newblock Cambridge University Press, Cambridge, 2011.
\newblock DOI: 10.1017/CBO9780511622656.

\bibitem{Meissner2007}
S. Meissner, A. Metz, and K. Goeke,
\newblock ``Relations between generalized and transverse momentum dependent parton distributions,''
\newblock Phys. Rev. D \textbf{76}, 034002 (2007).
\newblock DOI: 10.1103/PhysRevD.76.034002; arXiv:hep-ph/0703176.

\bibitem{Ji2013}
X. Ji,
\newblock ``Parton physics on a Euclidean lattice,''
\newblock Phys. Rev. Lett. \textbf{110}, 262002 (2013).
\newblock DOI: 10.1103/PhysRevLett.110.262002; arXiv:1305.1539.

\bibitem{Ji2014}
X. Ji,
\newblock ``Parton physics from large-momentum effective field theory,''
\newblock Sci. China Phys. Mech. Astron. \textbf{57}, 1407 (2014).
\newblock DOI: 10.1007/s11433-014-5492-3; arXiv:1404.6680.

\bibitem{JiZhangZhao2018}
X. Ji, J.-H. Zhang, and Y. Zhao,
\newblock ``Renormalization in large momentum effective theory of parton physics,''
\newblock Phys. Rev. Lett. \textbf{120}, 112001 (2018).
\newblock DOI: 10.1103/PhysRevLett.120.112001; arXiv:1706.08962.

\bibitem{Izubuchi2018}
T. Izubuchi, X. Ji, L. Jin, I. W. Stewart, and Y. Zhao,
\newblock ``Factorization theorem relating Euclidean and light-cone parton distributions,''
\newblock Phys. Rev. D \textbf{98}, 056004 (2018).
\newblock DOI: 10.1103/PhysRevD.98.056004; arXiv:1801.03917.

\bibitem{JiReview2021}
X. Ji, Y. Liu, Y.-S. Liu, J.-H. Zhang, and Y. Zhao,
\newblock ``Large-momentum effective theory,''
\newblock Rev. Mod. Phys. \textbf{93}, 035005 (2021).
\newblock DOI: 10.1103/RevModPhys.93.035005; arXiv:2004.03543.

\bibitem{ZhangGluonLaMET2019}
J.-H. Zhang, X. Ji, A. Schäfer, W. Wang, and S. Zhao,
\newblock ``Accessing gluon parton distributions in large momentum effective theory,''
\newblock Phys. Rev. Lett. \textbf{122}, 142001 (2019).
\newblock DOI: 10.1103/PhysRevLett.122.142001; arXiv:1808.10824.

\bibitem{LiMaQiu2019}
Z.-Y. Li, Y.-Q. Ma, and J.-W. Qiu,
\newblock ``Multiplicative renormalizability of quasi-parton operators,''
\newblock Phys. Rev. Lett. \textbf{122}, 062002 (2019).
\newblock DOI: 10.1103/PhysRevLett.122.062002; arXiv:1809.01836.

\bibitem{Balitsky2020}
I. Balitsky, W. Morris, and A. Radyushkin,
\newblock ``Gluon pseudo-distributions at short distances: forward case,''
\newblock Phys. Lett. B \textbf{808}, 135621 (2020); arXiv:1910.13963.

\bibitem{FanQuasiGluon2018}
Z.-Y. Fan, Y.-B. Yang, A. Anthony, H.-W. Lin, and K.-F. Liu,
\newblock ``Gluon quasi-parton-distribution functions from lattice QCD,''
\newblock Phys. Rev. Lett. \textbf{121}, 242001 (2018).
\newblock DOI: 10.1103/PhysRevLett.121.242001.

\bibitem{FanLinPion2021}
Z. Fan and H.-W. Lin,
\newblock ``Gluon parton distribution of the pion from lattice QCD,''
\newblock Phys. Lett. B \textbf{823}, 136778 (2021).
\newblock DOI: 10.1016/j.physletb.2021.136778; arXiv:2104.06372.

\bibitem{SalasChavira2022}
A. Salas-Chavira, Z. Fan, and H.-W. Lin,
\newblock ``First glimpse into the kaon gluon parton distribution using lattice QCD,''
\newblock Phys. Rev. D \textbf{106}, 094510 (2022).
\newblock DOI: 10.1103/PhysRevD.106.094510; arXiv:2112.03124.

\bibitem{GoodYaoLin2026}
W. Good, F. Yao, and H.-W. Lin,
\newblock ``First nucleon gluon PDF from large momentum effective theory,''
\newblock arXiv:2505.13321v3 [hep-lat] (2026).

\bibitem{ZhangSoft2020}
Q.-A. Zhang \textit{et al.} (Lattice Parton Collaboration),
\newblock ``Lattice-QCD calculations of TMD soft function through large-momentum effective theory,''
\newblock Phys. Rev. Lett. \textbf{125}, 192001 (2020).
\newblock DOI: 10.1103/PhysRevLett.125.192001; arXiv:2005.14572.

\bibitem{ChuCS2022}
M.-H. Chu \textit{et al.} (Lattice Parton Collaboration),
\newblock ``Nonperturbative determination of Collins--Soper kernel from quasi transverse-momentum dependent wave functions,''
\newblock Phys. Rev. D \textbf{106}, 034509 (2022).
\newblock DOI: 10.1103/PhysRevD.106.034509; arXiv:2204.00200.

\bibitem{Zhang2022TMD}
K. Zhang \textit{et al.} (Lattice Parton Collaboration),
\newblock ``Renormalization of transverse-momentum-dependent parton distribution on the lattice,''
\newblock Phys. Rev. Lett. \textbf{129}, 082002 (2022).
\newblock DOI: 10.1103/PhysRevLett.129.082002; arXiv:2205.13402.

\bibitem{ChuSoft2023}
M.-H. Chu \textit{et al.} (Lattice Parton Collaboration),
\newblock ``Lattice calculation of the intrinsic soft function and the Collins--Soper kernel,''
\newblock JHEP \textbf{08}, 172 (2023).
\newblock DOI: 10.1007/JHEP08(2023)172; arXiv:2306.06488.

\bibitem{ChuTMDWF2023}
M.-H. Chu \textit{et al.} (Lattice Parton Collaboration),
\newblock ``Transverse-momentum-dependent wave functions of pion from lattice QCD,''
\newblock Phys. Rev. D \textbf{109}, L091503 (2024).
\newblock DOI: 10.1103/PhysRevD.109.L091503; arXiv:2302.09961.

\bibitem{HeTMD2024}
J.-C. He \textit{et al.} (Lattice Parton Collaboration),
\newblock ``Unpolarized transverse-momentum-dependent parton distributions of the nucleon from lattice QCD,''
\newblock Phys. Rev. D \textbf{109}, 114513 (2024).
\newblock DOI: 10.1103/PhysRevD.109.114513; arXiv:2211.02340.

\bibitem{MaBM2025}
L. Ma \textit{et al.} (Lattice Parton Collaboration),
\newblock ``Quark transverse spin-momentum correlation of the nucleon from lattice QCD: the Boer--Mulders function,''
\newblock arXiv:2502.11807 [hep-lat] (2025).

\bibitem{Luscher2010}
M. Lüscher,
\newblock ``Properties and uses of the Wilson flow in lattice QCD,''
\newblock JHEP \textbf{08}, 071 (2010).
\newblock DOI: 10.1007/JHEP08(2010)071; arXiv:1006.4518.

\bibitem{LuscherWeisz2011}
M. Lüscher and P. Weisz,
\newblock ``Perturbative analysis of the gradient flow in non-Abelian gauge theories,''
\newblock JHEP \textbf{02}, 051 (2011).
\newblock DOI: 10.1007/JHEP02(2011)051; arXiv:1101.0963.

\bibitem{MonahanOrginos2017}
C. Monahan and K. Orginos,
\newblock ``Quasi parton distributions and the gradient flow,''
\newblock JHEP \textbf{03}, 116 (2017).
\newblock DOI: 10.1007/JHEP03(2017)116; arXiv:1612.01584.

\bibitem{Monahan2018}
C. Monahan,
\newblock ``Smeared quasidistributions in perturbation theory,''
\newblock Phys. Rev. D \textbf{97}, 054507 (2018); Erratum: Phys. Rev. D \textbf{110}, 059902 (2024).
\newblock DOI: 10.1103/PhysRevD.97.054507; arXiv:1710.04607.

\bibitem{Huo2021}
Y.-K. Huo \textit{et al.},
\newblock ``Self-renormalization of quasi-light-front correlators on the lattice,''
\newblock Nucl. Phys. B \textbf{969}, 115443 (2021).
\newblock DOI: 10.1016/j.nuclphysb.2021.115443; arXiv:2103.02965.

\bibitem{ZhangGauge2024}
K. Zhang \textit{et al.} (Lattice Parton Collaboration),
\newblock ``Impact of gauge fixing precision on the continuum limit of non-local quark-bilinear lattice operators,''
\newblock Phys. Rev. D \textbf{110}, 074505 (2024); arXiv:2405.14097.

\bibitem{JiHybrid2021}
X. Ji, Y. Liu, A. Schäfer, W. Wang, Y.-B. Yang, J.-H. Zhang, and Y. Zhao,
\newblock ``A hybrid renormalization scheme for quasi light-front correlations in large-momentum effective theory,''
\newblock Nucl. Phys. B \textbf{964}, 115311 (2021).
\newblock DOI: 10.1016/j.nuclphysb.2021.115311; arXiv:2008.03886.

\bibitem{Ishikawa2017}
T. Ishikawa, Y.-Q. Ma, J.-W. Qiu, and S. Yoshida,
\newblock ``Renormalizability of quasiparton distribution functions,''
\newblock Phys. Rev. D \textbf{96}, 094019 (2017).
\newblock DOI: 10.1103/PhysRevD.96.094019; arXiv:1707.03107.

\bibitem{Alexandrou2017}
C. Alexandrou \textit{et al.},
\newblock ``A complete non-perturbative renormalization prescription for quasi-PDFs,''
\newblock Nucl. Phys. B \textbf{923}, 394 (2017).
\newblock DOI: 10.1016/j.nuclphysb.2017.08.012; arXiv:1706.00265.

\bibitem{Cichy2019}
K. Cichy and M. Constantinou,
\newblock ``A guide to light-cone PDFs from lattice QCD: an overview of approaches, techniques and results,''
\newblock Adv. High Energy Phys. \textbf{2019}, 3036904.
\newblock DOI: 10.1155/2019/3036904; arXiv:1811.07248.

\bibitem{Peardon2009}
M. Peardon \textit{et al.},
\newblock ``Novel quark-field creation operator construction for hadronic physics in lattice QCD,''
\newblock Phys. Rev. D \textbf{80}, 054506 (2009).
\newblock DOI: 10.1103/PhysRevD.80.054506; arXiv:0905.2160.

\bibitem{Bali2016}
G. S. Bali, B. Lang, B. U. Musch, and A. Schäfer,
\newblock ``Novel quark smearing for hadrons with high momenta in lattice QCD,''
\newblock Phys. Rev. D \textbf{93}, 094515 (2016).
\newblock DOI: 10.1103/PhysRevD.93.094515; arXiv:1602.05525.

\bibitem{Egerer2022}
C. Egerer \textit{et al.} (HadStruc Collaboration),
\newblock ``Toward the determination of the gluon helicity distribution in the nucleon from lattice quantum chromodynamics,''
\newblock Phys. Rev. D \textbf{106}, 094511 (2022).
\newblock DOI: 10.1103/PhysRevD.106.094511; arXiv:2207.08733.

\bibitem{VanDerVorst1992}
H. A. van der Vorst,
\newblock ``Bi-CGSTAB: A fast and smoothly converging variant of Bi-CG for the solution of nonsymmetric linear systems,''
\newblock SIAM J. Sci. Stat. Comput. \textbf{13}, 631--644 (1992).
\newblock DOI: 10.1137/0913035.

\bibitem{Schoutrop2022}
C. Schoutrop, J. ten Thije Boonkkamp, and J. van Dijk,
\newblock ``Reliability investigation of BiCGStab and IDR solvers for the advection-diffusion-reaction equation,''
\newblock Commun. Comput. Phys. \textbf{32}, 156--188 (2022).
\newblock DOI: 10.4208/cicp.OA-2021-0182.

\end{thebibliography}

\end{document}
